% 本文件由 LaTeX 排版 Agent 根据有效 translation.json 自动生成。
% author=Gregory Nazianzen; work_id=3102; title=选演讲集

\chapter{选演讲集}

% structure: p0001 source=h1 target=omitted reason=page-title-already-rendered-by-parent

\Emph{演讲一}\newline{}\Emph{演讲二}\newline{}\Emph{演讲三}\newline{}\Emph{演讲七}\newline{}\Emph{演讲八}\newline{}\Emph{演讲十二}\newline{}\Emph{演讲十六}\newline{}\Emph{演讲十八}\newline{}\Emph{演讲二十一}\newline{}\Emph{演讲二十七} 第一篇神学演讲\newline{}\Emph{演讲二十八} 第二篇神学演讲\newline{}\Emph{演讲二十九} 第三篇神学演讲\newline{}\Emph{演讲三十} 第四篇神学演讲\newline{}\Emph{演讲三十一} 第五篇神学演讲\newline{}\Emph{演讲三十三}\newline{}\Emph{演讲三十四}\newline{}\Emph{演讲三十七}\newline{}\Emph{演讲三十八}\newline{}\Emph{演讲三十九}\newline{}\Emph{演讲四十}\newline{}\Emph{演讲四十一}\newline{}\Emph{演讲四十二}\newline{}\Emph{演讲四十三}\newline{}\Emph{演讲四十五}\par

\SourceRecord{Translated by Charles Gordon Browne and James Edward Swallow.；Nicene and Post-Nicene Fathers, Second Series；Vol. 7.；Edited by Philip Schaff and Henry Wace.；Buffalo, NY: Christian Literature Publishing Co.,；1894.；<http://www.newadvent.org/fathers/3102.htm>.}

% structure: p0001 source=h1 target=section reason=page-title

\section{第一篇讲道}

\Emph{论复活节及其迟延}\par

一、今天是复活之日，我的开始有良好的预兆。让我们隆重地庆祝这庆节，\BibleRef{依 66:5} 让我们彼此拥抱。让我们称那些憎恨我们的人为弟兄；更应如此对待那些因爱我们而做过或受过任何事的人。为了复活的原故，让我们宽恕一切冒犯：让我们彼此宽恕，我因我所受的高贵迫害（因为我现在可以称之为高贵）而宽恕你们；你们这些行使迫害的人，若你们有理由责备我的迟延，也请宽恕；因为或许在天主眼中，这迟延比别人的匆忙更为宝贵。因为暂时疏远天主也是好事，如同古时的伟大梅瑟 \BibleRef{出 4:10} 和后来的耶肋米亚 \BibleRef{耶 1:6} 所做的那样；然后，当祂召唤时，又欣然跑向祂，如同亚郎 \BibleRef{出 4:27} 和依撒意亚 \BibleRef{依 1:6} 所做的那样；只要两者都以虔诚的精神去做——前者出于对自己的软弱之感，后者出于对召唤者之能的敬畏。\par

二、一个奥迹傅油了我；我因一个奥迹退隐片刻，正如需要省察自己那样；现在我带着一个奥迹前来，带着这日子作为我怯懦和软弱的良好辩护者；为使今天从死者中复活的那一位也藉祂的圣神更新我；并穿上新人，将赐给祂的新创造，就是那些按天主的旨意重生的人，作为基督的良善模塑者和教师，甘愿与祂同死，也与祂同复活。\par

三、昨天羔羊被宰杀，门框被涂血，埃及为她的长子哀哭，毁灭者越过我们，印记可畏而可敬，我们被宝血围护。今天我们已经完全逃离埃及和法郎；没有人能阻止我们向主我们的天主举行庆节——我们出离的庆节；或庆祝那庆节，不用旧酵的恶意和邪恶，而用真诚和真实的无酵饼，\BibleRef{格前 5:8} 不带任何不敬和埃及的酵。\par

四、昨天我与祂同钉十字架；今天我与祂同受荣耀；昨天我与祂同死；今天我与祂同活；昨天我与祂同葬；今天我与祂同复活。但让我们向那为我们受苦并复活的主献上——你或许以为我打算说金银，或编织品，或透明昂贵的石头，那些只是地上易逝的物质，留在这下面，且大多数时候被恶人——世界和世界之王的奴仆——所拥有。让我们献上我们\Emph{自己}，这是天主最珍贵、最相称的财产；让我们把按肖像所造的归还给肖像。让我们认识我们的尊严；让我们尊崇我们的原型；让我们认识这奥迹的权能，以及基督为何而死。\par

五、让我们成为像基督一样，因为基督成为了像我们一样。让我们为祂的缘故成为天主的，因为祂为了我们成为了人。祂取了那较差的，为要把更好的赐给我们；祂成了贫穷的，为使我们藉着祂的贫穷而成为富足的；\BibleRef{格后 8:9} 祂取了奴仆的形体，为使我们重新获得自由；祂降下，为使我们被高举；祂受试探，为使我们得胜；祂受侮辱，为要荣耀我们；祂死了，为要拯救我们；祂升天，为要把我们这些因罪的堕落而低垂的人吸引到祂那里。让我们把\Emph{一切}献给祂，把\Emph{一切}奉献给祂，祂为我们自己献出了赎价和和好。但没有人能献出像自己一样的，即理解这奥迹，并为祂的缘故成为祂为我们所成的一切。\par

六、正如你们所见，祂赐给你们一位牧者；因为你们的善牧，那为自己的羊舍命的，正怀着希望和祈祷，并向你们这些他的属民提出要求；他将他自己双倍而非单倍地赐给你们，使他暮年的杖成为你们精神的杖。他又给无生命的圣殿加上一个有生命的；在那极其美丽和属天的圣所旁，加上这贫穷而微小的圣所，然而对他来说却是极有价值的，并且是用了许多汗水和辛劳建造的。但愿我能说它配得上他的辛劳。他将自己的一切（哦，伟大的慷慨！——或更确切地说，哦，父爱！）供你们使用：他的白发，他的青春，圣殿，大司祭，立遗嘱者，继承人，以及你们渴望的讲道；并且这些讲道不是虚浮的、向空气说出、只达到外耳的那种；而是圣神所写、刻在石版或肉版上的，不是刻在表面、容易磨掉的，而是印记极深，不是用墨，而是用恩宠。\par

七、这些是这位尊贵的亚巴郎，这位可敬可畏的首领，这位圣祖，这位一切美善的安息所，这位德行的标准，这位司祭职的圆满，今天向主献上他甘愿的祭品——他的独生子，那应许之子的人所赐给你们的礼物。你们方面，当向天主并向我们献上对牧者的服从，居住在青草地上，被平静的水滋养；认识你们的牧者，也被他认识；\BibleRef{若 10:14} 当他像牧者一样坦然地通过门呼唤你们时，你们跟随他；但不跟随一个像盗贼和叛徒一样爬进羊圈的外人；也不听从陌生声音，那声音会偷偷把你们从真理中带走，散落在山岭、\BibleRef{则 34:6} 旷野、陷阱和主不眷顾的地方；并会引诱你们离开对父、子和圣神的纯正信仰——同一权能和天主性——我的羊总是听祂的声音（但愿他们永远听），却用欺骗和败坏的话把他们从真牧者那里夺走。愿我们众人，牧者和羊群，都远离那如同毒草和死亡牧场的一切；引导和被引导远离它，使我们众人都在我们的主耶稣基督内合而为一，从今世直到天上的安息。愿荣耀和权能归于祂，直到永远。阿们。\par

\SourceRecord{Translated by Charles Gordon Browne and James Edward Swallow.；Nicene and Post-Nicene Fathers, Second Series；Vol. 7.；Edited by Philip Schaff and Henry Wace.；Buffalo, NY: Christian Literature Publishing Co.,；1894.；<http://www.newadvent.org/fathers/310201.htm>.}

% structure: p0001 source=h1 target=section reason=page-title

\section{\WorkTitle{第二篇讲道}}

\EditorialNote{普遍认为此讲道并非旨在口头宣讲。其目的是解释并辩护圣额我略近期的行为，该行为曾受到他在纳齐盎的朋友们的严厉批评。他大概是公元361年被父亲从庞图斯召回，他在那里与朋友圣巴西略一起隐修多年。他的父亲并不满足于儿子在家作为晚年依靠，而且确信他适合圣职，于是在公元361年圣诞节，不顾额我略的勉为其难，在会众高声赞同下，祝圣他为司铎。圣额我略甚至在多年之后仍称其祝圣为暴行，当时他几乎像被牛虻刺中的公牛一样狂怒，再次逃往庞图斯，在一位密友的深切同情下，将他受伤的情感埋葬于宜人的孤寂中。不久责任感重新抬头，他在公元362年复活节前回到父亲身边。在复活节那天，他向会众发表了第一篇讲道，会众的稀少表明纳齐盎人民对他的行为的不满。因此他在此讲道中向他们全面解释导致他隐退的动机。同时，如讲道的副标题所示，他阐述司铎职分的义务与尊严，后来所有论述此主题的作者都借鉴了它。圣金口若望在其著名论文中，圣大额我略在其\WorkTitle{牧灵指南}中，以及波舒哀在其颂扬圣保禄的颂词中，都不过是总结了这篇雄辩而精细的论文中的材料或发展其思考。}\par

\Emph{在辩护他逃往庞图斯及其返回，在他被祝圣为司铎之后，并阐述司铎职分的特征。}\par

1. 我被打败了，我承认自己的失败。我顺服了主，向祂祈祷。愿最蒙福的达味提供我的开场白，或者更确切地说，愿那位在达味内发言，甚至如今仍藉着他发言的那一位来提供。因为确实，开始每一言语和行动的最佳次序，就是从天主开始，并在天主内结束。至于原因，无论是关于我最初的叛逃和怯懦，在那期间我远远逃开，并离开你们一段时间，这在那些想念我的人看来或许很长；还是关于现在的温顺和心意的改变，在这改变中我再次将自己交付给你们，人们可能会根据对我的仇恨或爱戴，以不同方式思考和谈论：一方面拒绝宣告我无罪，另一方面则热诚欢迎我。因为没有什么比谈论他人的事更令人愉快，尤其是在情感或仇恨的影响下，这常常几乎完全使我们看不清真理。然而，我自己将无所畏惧地陈述真理，并在控告我们和勇敢辩护我们的两方之间公正裁决：一方面控告我自己，另一方面为自己辩护。\par

2. 因此，为使我的讲话按适当次序进行，我首先处理怯懦的问题。因为我不能容忍那些关注我努力之成败的人中，有任何人为我而感到困惑，既然天主乐意让我们的事情对基督徒有些重要，所以我将通过辩护来解除，如果有这样的人，那些已经受苦的人；因为尽可能并且如理性所允许，避免因我们的罪或猜疑而引起对团体的任何冒犯或绊脚石，乃是好的；因为我们知道，凡使这些小子中的一个跌倒，\BibleRef{玛 18:6}，那些使人跌倒的人，必然会从不说谎的那一位那里受到最严厉的惩罚。\par

3. 我现在的处境，诸位善人，并非出于缺乏经验和无知，不，让我稍微自夸一下，\BibleRef{格后 11:16}，也不是出于对神圣法律和规定的藐视。正如在身体中，有一个肢体统治并可以说主持，而另一个被统治并服从；同样在教会中，天主按照要么是平等的法律（它承认功过次序），要么是眷顾的法律（祂藉此将一切结合在一起），规定那些对此种对待有益的人，应受牧灵关怀和统治，并在言语和行为上被引导走上责任之路；而其他人则应成为牧者和教师，\BibleRef{弗 4:11}，为成全教会，这些人在德行和接近天主方面超越大多数，履行灵魂在身体中的功能，以及理智在灵魂中的功能；以便两者如此结合并连在一起，即使一个缺乏而另一个卓越，也能像我们身体的肢体一样，藉着圣神的和谐被结合并连在一起，形成一个完美的身体，真正配得上基督自己——我们的头。\BibleRef{弗 4:15}\par

4. 那么，我意识到无政府状态和无序，无论对于其他受造物还是对于人类，都不能比秩序和统治更有益；不，对于人类来说，这更是千真万确，因为在他们的情况下利害关系更大；因为即使他们达不到最高目的——即无罪——至少达到次要的目的——即从罪中恢复，这对他们来说是一件大事。既然这似乎是正确和公正的，那么我认为，所有人都希望统治，而没有人愿意接受统治，同样错误和无序。因为如果所有人都逃避这一职务——无论它应被称为服务还是领导——那么教会美丽的圆满\BibleRef{弗 1:23}就会受到极大妨碍，实际上不再美丽。而且，如果没有君王，没有总督，没有司祭，没有祭献，\BibleRef{欧 3:4}，也没有所有那些最高的职务（丧失这些职务，因他们的重罪，从前因不顺从而被判受此惩罚），那么在我们中间，天主将如何、由谁在这些神秘而提升的礼仪中受到敬拜呢？这些礼仪是我们最大最珍贵的特权。\par

5. 对于大多数致力于研究神圣事物的人来说，从被统治上升到统治，这既不是奇怪也不矛盾，也没有越出哲学所设定的界限，也不涉及耻辱；这正如一位优秀的水手成为瞭望员，而一位成功监视风向的瞭望员被委以舵柄；或者，如果你愿意，一位勇敢的士兵被任命为队长，一位好队长成为将军，并将整个战役的指挥权交给他。同样，我也不是因为渴望更高的等级而对这一级别的地位感到羞耻，正如也许一些荒唐而令人厌烦的人所猜想的那样，他们以自己的感受来判断他人的感受。我并非如此无知于这一等级的神圣伟大或人类卑微，以至于认为受造的本性，无论以多么微小的程度接近天主，都是一件小事；天主是唯一最光荣最卓越的，并在纯洁上超越每一个本性，包括物质的和非物质的。\par

6. 当时我的感受如何，以及我不服从的原因是什么？因为在大多数人看来，我那时似乎与自己的过去不一致，也不像他们所认识的我，而是有所退步，表现出比应有的更大的抗拒和固执。这些原因你们早已渴望听我道来。首先，也是最重要的，我被所发生之事的出乎意料所震惊，如同人被突如其来的声响吓到一样；我失去了理智的控制，而此前一直控制着我的自尊也随之瓦解。其次，我心中涌起一种对宁静与退隐之福的热切渴望——这种渴望，我自开始求学以来就比任何其他文人更为热衷，并且在最大、最危险的险境中我曾向天主许诺过，我也曾如此深切地体验过它，以至于我当时已处在它的门槛上，而这种渴望也因此被大大点燃，以致我无法忍受因一项专横的压迫行为而被强行投入喧嚣的生活，并被暴力从这种生活的神圣避难所中拖离。\par

7. 因为在我看来，没有什么比关上感官之门、逃离肉身与世界、收敛于自身之内、除绝对必要外不再与人间事务有更多联系、对自己和天主说话，\BibleRef{格前 14:28} 生活在可见事物之上、不断在我内保持神圣的印象纯净而未被此下界的错误标记所混淆、既是且不断成长为一个真正无瑕的天主和神圣之事的明镜、如同光上加光、原本黑暗的逐渐明亮、已藉着望德预享来世之福、与天使一同漫游、现今已因离弃尘世而超越大地、藉着圣神立于高处更为甘美。如果你们中有人曾拥有这种渴望，他就明白我的意思，并会同情我当时的情感。因为，也许我不应期望凭我的话说服大多数人，他们不幸地因自己的轻率或因不值得这应许之人的影响而嘲笑此类事情，那些人给善事冠以恶名，称哲学为胡言，加之嫉妒和民众的邪恶倾向——这些人总是趋向堕落；以致他们不停地纠缠于两种罪过之一：要么作恶，要么诋毁善事。\par

8. 此外，还有一种情感影响了我——我不知道这是卑劣还是高尚，但我要向你们坦白我所有的秘密。我为所有那些人感到羞耻，他们并不比普通人好，甚至，他们若不是更坏，就已经是件好事了——他们洗手未净，\BibleRef{谷 7:5} 正如俗语所说，灵魂未受启迪，却闯入最神圣的职务；在配得上走近圣殿之前，他们就觊觎圣所，在圣桌周围推挤，仿佛以为这个圣职是一种谋生手段，而非美德的典范，或是绝对权威，而非我们必须为之交代的服务。事实上，他们的人数几乎与所管辖的人相当；在虔敬上可怜，在尊位上不幸；以至于在我看来，随着时间和这种邪恶一同发展，他们将无人可统治，当人人都成为教师——而不是像应许所说的那样，由天主教导——并且人人预言，以致连\Quote{撒乌耳也在先知之中}，根据古代历史和谚语。因为无论在现今还是往昔，在各种发展兴衰之中，从未有过像现在基督徒中这样多的这类耻辱和滥用。如果阻止这股潮流非我力所能及，那么至少，见证我们对它的憎恶和羞愧，是宗教并非最不重要的职责。\par

9. 最后，还有一件比我已提及的任何事都更严重的事——我现在要说到问题的终点了：我不会欺骗你们，因为这在涉及如此重大主题时是不允许的。我过去和现在都不认为自己有资格管理一群羊或牛群，或有权管辖人的灵魂。因为在牛羊的情况中，只要使牛群或羊群尽可能肥壮就够了；为此目的，牛倌和牧人会寻找水草丰美的牧场，将他的牲畜从一处牧场赶到另一处，让它们休息、唤醒或召回，有时用他的牧杖，更常用他的笛子；除了偶尔与狼搏斗或照料病畜，他的大部分时间都消磨在橡树、阴凉和笛子中，他躺在美丽的草地上、清凉的水边，在有风的地方铺下卧榻，不时唱着情歌，身旁放着杯子，对着他的小公牛或羊群说话，其中最肥壮的那些供应他的宴席或报酬。但从未有人考虑过牛羊的美德；因为它们有什么美德呢？或者谁曾认为它们的利益比自己的享乐更重要呢？\par

10. 但对于人来说，尽管学会服从统治很难，但学会统治人似乎更难，而其中最难的是我们的这种统治——它藉着天主法律、朝向天主引导人——因为在有思想的人看来，它的风险与其高度和尊严相称。因为，首先，他必须像银子或金子一样，虽然在各种时节和事务中普遍流通，却永不发出虚假或掺杂的声音，也不显示任何需要进一步用火精炼的劣质材料的迹象；\BibleRef{格前 3:12} 否则，他的统治范围越广，他所造成的邪恶就越大。因为波及许多人的伤害，大于仅限于一个人的伤害。\par

11. 染一块布，或使接近之物沾染气味（无论香臭）都不容易；致命的毒气（它被恰当地称为瘟疫）感染空气，再通过空气进入生物体内，也不如一位上级的恶习如此迅速地占据其下属，而且比其反面——德行——要容易得多。因为邪恶在这方面尤其胜过良善，而最令我痛心的是，我察觉到恶习具有吸引力和现成性，没有什么比成为邪恶更容易，即使没有人引导我们走向它；而德行的获得却是罕见且困难的，即使有很多东西吸引和鼓励我们。我认为，最有福的\Person{哈盖}在他那奇妙而极其真实的比喻中正是看到了这一点：——\Quote{你们请问司铎关于法律说：若圣肉触到衣服，再触到肉、饮料或器皿，它能圣化所接触之物吗？他们回答说：不能。再问：若这些东西中任何一件触到不洁之物，它岂不立即沾染污秽吗？他们必定会告诉你们，它确实沾染了污秽，并不因接触而保持洁净。}\par

12. 这是什么意思？依我看，良善难以在人性中扎根，如同火在湿木上一样；而大多数人却随时准备并倾向于加入邪恶，如同干草，我的意思是，它们随时准备迎接火花和风，因其干燥而容易点燃和烧毁。因为任何人在轻微诱因下，都会更快地完全参与邪恶，胜过面对完全呈现在面前的良善而只达到轻微程度。事实上，一点点苦艾能最快地把苦味渗入蜂蜜；而即使双倍的蜂蜜也难以把甜味渗入苦艾；一小块卵石的移开就能使整条河流改道，而最强固的堤坝却难以约束或阻止它的奔流。\par

13. 因此，在我们所说的话中，第一点是我们应当提防的，即：不要成为德行魅力的拙劣画师，更甚者，或许成为劣质画师的模样，劣质模样给百姓，或者勉强逃脱那句谚语：我们着手医治别人\BibleRef{路 4:23}，而自己却满身疮痍。\par

14. 其次，即使一个人已经极其纯洁地远离了罪恶，我也不认为这对于将要教导他人德行的人来说是足够的。因为接受这职责的人，不仅需要远离邪恶（因为在大多数受他照管的人眼中，邪恶是最可耻的），还需要在善行上卓越，按照命令：\Quote{远离邪恶，行善。}他不仅要从灵魂中抹去恶习的痕迹，还要铭刻更美好的德行，以便在德行上超越他人，正如他在尊严上超越他们一样。他应当不知道良善或灵性进步的界限，应当挂念未达到的而非已获得的，并视脚下之物为通往下一步的阶梯；不应以为超越一般人就是大成就，而应把达不到当行之事视为损失；应以诫命而非邻人（无论他们是邪恶还是在一定程度上修德）来衡量自己的成功；不应以微小的天平衡量德行，因为它是归于至高者的，\BibleQuote{万有出于祂，万有归于祂。}\BibleRef{罗 11:35}\par

15. 他也不应认为同样的事物适用于所有人，正如并非所有人都有同样的身高，面容的特征、动物的本性、土地的特质、星辰的美观和大小也不尽相同；相反，他必须视恶行为普通人的过错，当在法律的严格管制下受惩罚；而对于统治者或领袖，过错在于未能达到最高的德行并不断进步，如果确实要以其高度的德行，通过劝说的影响力而非权威的强迫，带领人民达到一般的德行水平。因为非自愿的行为，除了出于压迫之外，既没有功德也不持久。被迫之事，如同被手强行拉开的植物，一旦放手就恢复原状；而出于选择之事则最合法且持久，因为它为善意的纽带所维系。所以，我们的法律和立法者最严格地命令我们：\BibleQuote{你们要牧放天主的羊群，不是出于不得已，而是出于自愿。}\BibleRef{伯前 5:2}\par

16. 但即便一个人远离恶习，达到了最高程度的德行，我也不知道有什么知识或能力能使他敢于承担这一职分。因为引导人——这个最变化多端、最复杂的受造物——在我看来确实是一门艺术之艺术、科学之科学。任何人可以通过比较灵魂医师的工作与身体的治疗来认识这一点；注意后者虽然费力，但我们的工作更费力，也更重要，这是由其主题的性质、科学的力量以及实践的目的所决定的。一个工作处理身体和必朽坏、有缺陷的物质，这物质最终必须分解并遭遇其命运\BibleRef{创 3:19}，即使在此刻，借助艺术，它可以克服自身内部的紊乱，但因疾病或时间而分解，服从自然法则，因为它无法超越自身的限制。\par

17. 另一种则关涉来自天主、具有神性且分有天上高贵性的灵魂，它即使与较低的本性相连，也仍奋力向上攀升。或许还有别的原因，只有将二者结合的天主以及那些在天主内受教于这些奥秘的人才能知晓，但就我和像我这样的人所能感知的，有两个原因：其一，为使灵魂通过与下方事物角力与摔跤（\BibleRef{弗 6:12}），如同金子经火试炼（\BibleRef{伯前 1:7}），从而承继天上的荣耀，将来所望之物作为德行的奖赏而获得，而非仅作为天主的恩赐。这实在是至善者之意愿，要使善甚至成为我们自己的——不仅因为播种在我们的本性中，更因我们自己的选择以及自由意志的运作而培养。第二个原因是，为使灵魂能逐步将较低的本性从粗重中释放，将其拉向并提升至天上，从而使灵魂之于身体如同天主之于灵魂，灵魂亲自引领那服务于它的质料，并使之如同同仆一般与天主结合。\par

18. 地点、时间、年龄、季节及诸如此类，是医生审查的对象；他会开具药物和饮食，防范有害之物，以免病人的欲望妨碍他的医术。有时在某些情况下，他会使用烧灼、刀割或更剧烈的疗法；但所有这些，无论看似如何费力与艰难，都比不上诊断与治疗我们的习惯、情欲、生活、意志以及内在的一切，要从我们复合的本性中驱逐一切残暴与凶恶，引入并确立温和与中悦天主的事物，并在灵魂与身体之间公平裁决：不允许卓越者被低劣者压倒（此乃最大不公），而是使那自然居于第二位的服从于统治与领导的力量——正如神圣律法所命，这律法极好地施加于祂的全部受造物，无论可见或超越我们的视野。\par

19. 另一点我也并非不知：医生所关注的那些对象的本性保持不变，并不产生任何狡猾的对抗或与其医术之器械为敌的诡计；相反，是治疗改变其对象，除非病人有轻微的违抗，但阻止或约束它并不困难。然而在我们这种情况中，人的思虑、自私、缺乏训练及不愿顺从，对德行的进步构成极大障碍，几乎成了对想帮助我们之人的武装抵抗。而本应坦露病症给灵性医生的那种热切，我们却用来逃避这种治疗，并以对抗自身利益来显示勇敢，以规避健康为技巧。\par

20. 因为我们要么隐藏自己的罪，将其深埋于灵魂深处，如同某种溃烂恶疾，以为能躲过人的眼目便能逃避天主的巨眼与公义；要么在罪中找借口，为我们的失足编造辩护理由；要么如同聋虺堵住耳朵，闭紧双耳，顽固拒绝听那召唤者的声音，拒绝用智慧的药物医治——灵性疾病正是由此被治愈；最后，我们中最胆大妄为、固执己见者，在那些本可治愈我们的人面前无耻地硬顶着罪，如俗语所说，昂首阔步踏入各种过犯。哦，何其癫狂——若不对此心态有更合适的称呼！我们本应爱戴为恩人者，反视之如仇敌，恨恶在城门口责备的人，厌恶正直之言（\BibleRef{亚 5:10}）；我们以为，通过与善待我们的人作战，就能成功，实则尽己所能伤害自己，如同那些以为自己吞噬他人血肉却咬住自己的人。\par

21. 为此我断言，我们医师的职责远比局限于身体的医师更为辛劳，也因此更有价值；此外，后者主要涉及表面，仅略微探究深层原因。而我们的全部治疗与努力都关乎心中的隐秘之人（\BibleRef{伯前 3:4}），我们的争战是针对那内在的仇敌与敌人，它使用我们自己作为武器来攻击我们自己，最可怕的是，它使我们陷于罪的死亡。因此，为对抗这些仇敌，我们需要宏大而完美的信德，更需要天主的协力，而依我所信，也需要我们自己不轻微的反制，这反制须通过言与行彰显，好使我们最宝贵的所有——我们自己，得到妥善的照料、洁净，并尽可能变得配得上。\par

22. 然而，且让我们转而思考这两种医治方式的目的（因这一点尚待考量）：前者保存肉身的健康良好状态（若已存在），或在其缺失时恢复之；尽管尚不清楚这些是否会对拥有它们的人有利，因为其反面往往带来更大益处，如同贫穷与财富、声誉或羞辱、低微或显赫的地位以及所有其他本质上中立、不倾向一方的境况，其好坏取决于经历者的意志和使用方式。而我们医术的目标，是为灵魂安上翅膀，将其从世界解救出来，交予天主；守护那按天主肖像（\BibleRef{创 1:26}）所造的，若它存留，则扶持之，若它处于危险中，则引领之，若它毁坏，则恢复之；使基督因圣神居于心中（\BibleRef{弗 3:17}）；简言之，使那属于天上军团的人得以神化，并获得天上的福乐。\par

23. 这正是我们的启蒙师\BibleRef{迦 3:24}法律的意愿，也是介于基督与法律之间的先知们的意愿，更是成全并终结\BibleRef{希 12:2}属神法律的基督的意愿；是空虚自己的天主性\BibleRef{斐 2:7}的意愿，是所取的血肉\BibleRef{希 2:14}的意愿，是那新的天主与人的结合——由二者组成之一个、而二者合为一——的意愿。为此，天主藉着灵魂与肉体结合，如此分开的本性藉着居中媒介与每一方的亲缘而联结在一起：所以一切成为一，是为了万有，也是为了我们的一位始祖——灵魂是为了那不顺服的灵魂，肉体是为了那与其合作并同受定罪的肉体，基督——祂超越并远离罪恶——是为了那成了罪恶奴仆的亚当。\par

24. 因此，新的代替了旧的\BibleRef{希 8:8}，那曾受苦的为受苦而复活，祂那超越我们的每一属性与我们的每一属性互换，新的奥秘代替了旧的安排，出于仁爱对待那因不顺服而堕落者。这就是诞生与贞女的缘由，马槽与白冷的缘由；诞生是为受造界，贞女是为女人\BibleRef{创 2:7}，白冷\BibleRef{路 2:7}是为伊甸，马槽是为乐园，微小可见的事物是为伟大隐密的事物。因此天使\BibleRef{路 2:14}先荣耀天上的，后荣耀地上的\BibleRef{格前 15:49}；牧羊人看见羔羊与牧者身上的荣耀；星辰引领贤士前来朝拜并献礼\BibleRef{玛 2:9}，为摧毁偶像崇拜。因此耶稣受洗，并获得从上而来的见证；祂禁食\BibleRef{玛 4:2}、受试探，并战胜了那曾战胜者。因此魔鬼被驱逐\BibleRef{玛 10:7}、疾病得医治，那大能的宣讲托付给卑微的人，并由他们成功传扬。\par

25. 因此，外邦人喧嚷，万民妄想空虚之事；因此，树对抗树\BibleRef{若 19:17}，手对抗手——一手放纵伸展\BibleRef{创 3:6}，另一手慷慨伸出；一手不受约束，另一手被钉穿\BibleRef{玛 27:35}；一手驱逐亚当，另一手使地极和好。这就是高举以补赎堕落的原因，苦胆是为尝味，荆棘冠是为恶的统治，死亡是为死亡，黑暗是为光明，埋葬是为归于尘土，复活是为复活。这一切都是天主对我们的训练，对我们软弱的医治，使旧亚当重返他堕落之处，引领我们到达生命树——那因不顺时不当食用而使我们疏远的知善恶树。\par

26. 这医治的施行者与同工，正是我们这些管理他人的人；对我们而言，认出并医治自己的激情与疾病已是大事——或者更确切地说，并非大事，只是此等阶层的多数人堕落使我如此说——但更大的事是能医治并巧妙洁净他人的激情与疾病，这既有利于需要医治的人，也有利于负责医治的人。\par

27. 再者，医治我们身体的人，他们有自己的辛劳、守望和忧虑——这我们都知道——并会从他人的痛苦中为自己收获痛苦，如同他们中一位智者所言；他们会将自己劳苦研究的成果以及从他人处借得的资料，提供给需要的人使用。他们认为自己所发现或未能发现的任何最微小细节，对健康或危险都具有重要影响。这一切的目的是什么？是使人能在世上多活几天，尽管他可能并非善人，而是最堕落的人之一——或许因其邪恶，更早死去对他更好，以摆脱那最严重的疾病——邪恶。但假如他是善人，他能活多久？永远吗？他从这生命中能获得什么？对清醒明智的人而言，从这生命中获得的最大福分，难道不是寻求解脱吗？\par

28. 但我们的努力所系的是灵魂的救恩——一个蒙福、不朽、注定因恶行或善行而受永罚或永赏的存有——我们应当怎样奋斗，需要多大的技巧来妥善治疗或使人得到治疗，改变他们的生命，将泥土交付给圣神？因为男人和女人、青年和老年、富人和穷人、乐观者和沮丧者、病人和健康人、统治者和被统治者、智者和愚者、懦夫和勇士、愤怒者和温柔者、成功者和失败者，所需的教导和鼓励并不相同。\par

29. 若你更仔细省察，已婚与未婚者之间的区别何其大；在未婚者中，独修者与团体共居者之间、在默观中精进者与勉强持守正道者之间、城镇居民与乡野村夫之间、质朴者与诡诈者之间、忙碌者与闲逸者之间、遭遇挫折者与不知不幸的顺遂者之间，区别又有多大。因为这些群体彼此在欲望和激情上的差异，有时比他们在身体特征上的差异更大；或者如果你愿意，比构成我们元素的混合与交融的差异更大，因此，规范他们并非易事。\par

30. 正如同一药物和同一食物并非对每个人的身体都适用，而是根据他们健康或虚弱的程度有所区别；同样，灵魂也以不同的教导和引导来对待。那些有过这种经验的人可以为这种治疗作证。有些人被教义引导，另一些人被榜样训练；有些人需要刺棒，另一些人需要缰绳；有些人懒惰，难以被唤醒行善，必须用言语击打来激励；另一些人在精神上过于热切，冲动难以约束，如同纯种小马，跑离转弯标杆，要改善他们，言语必须有约束和抑制的作用。\par

31. 有些人受益于赞扬，另一些人受益于责备，两者都应在适当的时候施用；如果不合时宜或不合理，则有害；有些人被鼓励纠正，另一些人被斥责纠正；有些人在公开受责时改正，另一些人在私下受责时改正。因为有些人惯于藐视私下的劝诫，但被众人的谴责唤回理智；而另一些人，若公开受责会变得鲁莽，却因私下受责而接受，并因同情而服从。\par

32. 有些人需要在最细微的细节上严密监视，因为他们如果以为未被察觉（他们会设法如此），就会因自以为有智慧而自高自大。对另一些人，最好不予理会，而是视而不见、听而不闻，正如俗语所说，以免我们因一再的责备而使他们陷入绝望，最终变得完全鲁莽，当他们失去了自尊——劝服力的源泉时。在某些情况下，我们甚至必须生气，却不感到生气；或对他们表现出我们并未感到的轻蔑；或表现出绝望，尽管我们并未真正对他们绝望，这要根据他们本性的需要。又有些人，我们必须以迁就和谦卑对待，帮助他们轻易地生出对更美好事物的希望。有时，征服某些人——或被别人征服，以及赞扬或贬抑——对某些人来说，在一种情况下是财富和权力，在另一种情况下是贫穷和失败，这往往更为有利。\par

33. 因为我们的治疗并不对应德行和恶行，其中一个在任何时候和任何情况下都是最卓越和有益的，而另一个则是极其邪恶和有害的；而且，我们同一种药物——无论是严厉还是温和，或是我们列举的其他任何药方——并不总是一成不变地在同样的情况下证明是最有益或最危险的；在某些情况下，它证明是良好和有用的，在另一些情况下，它又有相反的效果，我想，这取决于时间和情况以及病人的性情。现在要向你展示这一切之间的区别，并给你一个完全精确的看法，以便你能够简要地理解医疗技艺，即使对于最细心和最有技巧的人来说，也是完全不可能的；但实际的体验和实践对于形成一个医疗体系和一位医生是必要的。\par

34. 然而，我认为这一点是普遍公认的——正如在高空走钢丝的人向任何一边倾斜都不安全，因为即使倾斜似乎很轻微，后果却不轻微，他们的安全有赖于完全的平衡；同样，在我们中间，如果一个人向任何一边倾斜，无论是因恶习还是无知，对于他自己和那些被他带领的人，都会陷入跌倒犯罪的不小危险。但我们实在必须在君王的道路上行走，\BibleRef{户 20:17}，并小心不要向左或向右偏离，\BibleRef{箴 4:27}，正如箴言所说。因为我们的情欲正是如此，而在这件事上，善牧的任务就是如此，如果他想要恰当地认识他羊群的灵魂，并按照正确和公正的牧灵方法，以及配得上我们真正牧人的方法来引导他们。\par

35. 关于圣言的分配——最后提到我们的首要职责——那神圣而崇高的圣言，现在人人都准备谈论它；如果有人大胆地承担它，并认为它是在每个人智力范围之内，我对他的智慧（且不说愚蠢）感到惊讶。对我而言，这似乎不是一项小任务，并且需要不小的属灵能力，才能按时分给每个人应得的圣言之分，\BibleRef{路 12:42}，并以判断来调节我们意见的真理，这些意见涉及诸如世界或诸世界、物质、灵魂、理智、更好的或更坏的本性、眷顾——它维系和引导宇宙，并且在我们对它的经验中，似乎是按照某种原则被治理，但这一原则与地上和人类的经验相矛盾。\par

36. 此外，它们还涉及我们的原始构造和最终恢复，真理的预像，盟约，基督的第一次和第二次来临，祂的降生成人、苦难和分解，以及复活、末日、审判和报偿，无论是悲伤的还是荣耀的；最后，还有我们当如何看待原始的、有福的\OriginalTerm{天主圣三}{Trinity}。这对那些负责光照他人的人来说包含着极大的危险，如果他们为了避免多神论而将教义收缩为单一位格，从而只留给我们空洞的术语——如果我们认为圣父、圣子和圣神只是同一个位格；或者，另一方面，将祂分割为三个，或相异、或混乱、无原则、或可以说是敌对的神明，从而从反对方面跌入同样危险的错误：就像一株扭曲的植物，如果向相反方向弯曲过甚。\par

37. 因为在关于神学的三种软弱——无神论、犹太教和多神教——中，其中一种得到利比亚人\Person{撒伯流}的支持，另一种得到\Place{亚历山大里亚}人\Person{亚略}的支持，第三种得到我们中间某些极端正统派的支持，那么我的立场是什么？我能否避免这三种立场中有害的东西，而保持在虔诚的界限内？既不因新的分析与综合而被引入撒伯流的无神论，以致断言万事并非一体而是各自虚无（因为那些相互转化、彼此过渡的事物，就不再是它们各自所是，这样我们就有一个非自然的复合神祇，如同那些神话中的生物，是生动想象的产物）；也不因主张分离的诸多性体——正如亚略那名副其实的疯狂——而陷入犹太式的贫乏，并将嫉妒引入天主性中，将天主性局限于\OriginalTerm{无始者}{Unbegotten}独一，仿佛害怕我们的天主会消失，如果祂是一个同等本性真天主的父亲；也不因排列三个相互对立或联合的本原，而引入我们已逃脱的外邦人多本原论？\par

38. 既不可过分尊崇圣父以致剥夺祂的父性（因为如果圣子被分离并与祂疏远，并被归入受造物，祂又怎能是父呢？——因为一个异质的存在，或一个与其父混合混淆、使其混乱的存在，就不是儿子）；也不可过分尊崇基督，以致忽略保存祂的子性（因为如果祂的本原不归于父，祂又是谁的儿子呢？）以及父作为本原的地位，因为祂是父和生者；因为祂将是卑微无价值者的本原，或者更确切地说，\Quote{本原}一词将被用于卑微无价值的意义，如果祂不是那在圣子和圣神内被默观的天主性和美善的本原：前者是圣子和圣言，后者是那发出而不分离的圣神。因为必须既保存天主性的合一，又宣认三位，每一位有其自身的属性。\par

39. 对这一问题恰当而可敬的理解和阐述，需要比现在场合甚至我们一生所能允许的更长的讨论，而且，更重要的，无论现在还是任何时候，都需要圣神的援助，唯有藉着祂，我们才能领悟、阐述或把握有关天主的真理。因为只有纯洁者才能领悟那纯洁的、与自己性情相同的祂；我如今已简要地谈及这个问题，以表明讨论如此重大的问题有多困难，尤其是在由各种年龄和身份的人组成的广大听众面前，他们如同一个多弦的乐器，需要用多种方式弹奏；或者找到任何一种能够使他们都得到建造并用知识之光启明他们的言辞形式。因为不仅危险出自三个方面：理解、言语和听觉，以致在一方面（如果不是全部）的失败是必然的——因为要么心智未被光照，要么言语无力，要么听觉未经洁净而无法理解，因此，在一方面或所有方面，真理必然受损；而且，使那些自称教授任何其他科目者的指导如此容易和可接受的原因——即听众的虔诚——在这个问题上却涉及困难和危险。\par

40. 因为他们既然承担了为天主至高存在、为救恩、为我们众人首要希望辩争的责任，他们在信德上越热忱，就越敌视所说的话，以为顺从的态度不是虔诚，而是对真理的背叛，因此他们宁愿牺牲一切，也不愿放弃自己的私人信念和所受教导的惯常教义。我现在指的是那些性情温和、并非彻底败坏的人，他们若在真理上犯了错，也是出于虔诚而犯错，他们有热忱，但不合乎真知（见：\BibleRef{罗 10:2}），他们或许属于那些不过分被定罪、不被多鞭打的人（见：\BibleRef{路 12:47}），因为他们不是在恶习或堕落中未能遵行主的旨意。这些人或许会被说服，抛弃那引起他们敌意的虔诚观点，如果某种理由——或来自他们自己的心智，或来自他人——抓住了他们，并在关键时刻，如同燧石取铁一样，从一颗怀胎且值得光的心智中迸出火花，因为这样一点火花就会迅速点燃其中的真理火炬。\par

41. 对于那些出于虚荣或傲慢而妄论至高者，以\Person{雅乃乃}和\Person{杨布勒}的傲慢武装自己（见：\BibleRef{弟后 3:8}），不是反对梅瑟，而是反对真理，并起来反对健全教义的人，又该说什么呢？或者对于第三类人，他们因无知及其后果——鲁莽，以猪一般的行径冲向一切形式的教义，并将真理美丽的珍珠践踏在脚下？\par

42. 又那些对天主没有个人见解或措辞（不论好坏），却听取各种教义和教师，意图从所有中选出最好最安全的，除了自己之外，不相信有更好的真理判断者，又该如何呢？结果，他们被一个又一个似乎有理的观念吹动、转来转去，在被各种教义淹没和蹂躏之后（见：\BibleRef{弗 4:14}），并在接连更换教师和信条（他们像灰尘一样轻易抛弃）之后，最终耳朵和心智都疲倦了，而且，何等愚昧！他们对一切形式的教义都同样厌恶，并承担起可悲的角色：嘲笑和鄙视我们的信德，认为它不稳定不健全；在无知中从教师转到教义：就好像一个眼睛有病或耳朵受伤的人，抱怨太阳昏暗不发光，或抱怨声音不和谐微弱。\par

43. 因此，趁灵魂尚新鲜如未受印封的蜡时，将真理印在其上，比将虔诚的教训——我指错误的教训和信条——刻在已有的铭文之上更为容易，其结果是后者将前者扰乱并打乱。的确，走一条平坦且常有人行的路，比走一条未经践踏且崎岖的路更好；耕种一块常被犁头劈开翻松的土地，比耕种荒芜之地更好。但需要书写的灵魂，应当没有有害教训或罪恶深深刻痕的铭刻；否则，虔诚的书写者将面临双重任务：擦除先前的印象，并代之以更卓越、更值得持守的其他印象。那么，那些不仅因屈从于私欲偏情而显出其邪恶，甚至因言论而显出其邪恶的人何其多！邪恶的形式与特征何其多！受托教育和管理灵魂的职分者所面对的任务何其艰巨！实际上，我略去了大多数细节，免得我的言论不必要地冗长。\par

44. 若有人承担驯服和调教一个多形多状、由大小不同、温顺与野性程度各异的多种动物组合而成的动物，其主要任务——涉及一场相当艰苦的搏斗——将是管理这种异常异质的本性，因为构成它的每一种动物，按其本性或习惯，对同样的言语、食物、手抚、口哨或其他处理方式，会以不同的方式感受喜乐、愉悦或厌恶。那么，这动物的主人若不尽显其知识的多方多样，对每一种动物施以适合的处理，以便成功引领并保全这兽，他又能做什么呢？既然教会的共同身体由许多不同的性格和心智组成，如同一个由不协和部分组合而成的单一动物，那么它的治理者绝对必须一方面在一切事情上保持正直的单纯，另一方面在对待个人时尽可能多方多样、变化多端，并以适当和适宜的方式对待所有人。\par

45. 因为有些人需要用最简单、最初步的教训——即那些在习惯上如婴孩、可以说是新生、不能承受圣言的成人食物之人——的奶来喂养；不仅如此，若向他们提供超出其力量的食物，他们很可能会因心灵的无能而被压倒和压抑，如同我们物质身体的情况一样，无法消化和吸收所提供的东西，甚至会失去原有的力量。另一些人则需要那些在成全者中所讲的智慧（\BibleRef{格前 2:6}）以及更高、更坚实的食物，因为他们的官能已经充分操练，足以分辨（\BibleRef{希 5:14}）真理与虚假；若叫他们喝奶、吃病人蔬菜般的食物（\BibleRef{罗 14:2}），他们会感到厌烦。这也有正当理由：因为他们不会按照基督得着力量，也不会做出那值得称许的增长——圣言在正确喂养的人身上产生这种增长，使他成为完人，达到基督圆满年龄的程度（\BibleRef{弗 4:13}）。\par

46. 谁能胜任这些事呢？因为我们不像许多人那样，能败坏（\BibleRef{格后 2:16}）真理的圣言，将那使人欢畅心灵（\EditorialNote{依撒意亚 1:22}）的酒搀水——即用普通、廉价、低劣、陈腐、无味的东西掺杂我们的教训，以便从掺假中获利，并迎合那些遇见我们的人，讨好每个人，变成腹语者和喋喋不休者，他们用出自尘土又归于尘土的话语服务于自己的私欲，为博取多数人的特别好感，而极其伤害、甚至毁灭我们自己，流无辜单纯灵魂的血，这血将向我们要索。\par

47. 此外，我们知道，将自己缰绳交给比我们更有技巧的人，比在自己毫无经验时引导别人的路程更好；宁可善意地聆听，也不要鼓动未经训练的舌头。在与我认为颇有价值、至少对我们怀有好意的人商议这些事之后，我们宁愿学习那些我们不知道的言语与行动准则，而不是在无知中贸然教导。因为来到长者那里，甚至直到年老，接受长者的推理，是令人愉快的，因为它能帮助新近虔诚的灵魂。因此，在自己尚未充分受训之前就承担训练他人，如同人们所说的\Quote{在酒瓮上学陶艺}——即以他人灵魂为代价来操练虔诚——在我看来似乎是极度的愚昧或极度的鲁莽：愚昧，如果我们甚至不知道自己的无知；鲁莽，如果我们明知如此仍贸然从事。\par

48. 不仅如此，希伯来人中更具智慧者告诉我们，古代希伯来人中有一条极好且值得称赞的法律：并非每个年龄的人都可受托全部圣经，因为这对所有人并非更有益处，其全体并非人人都能立即理解，而其更深奥的部分若按其字面意义会对大多数人造成极大伤害。因此，有些部分，因其外表无可指摘，一开始就允许并普遍分享；另一些部分只托付给已满二十五岁的人——即那些将自身奥秘之美藏在朴素外表之下，作为勤勉和卓越生活的赏报；只向心灵已净化的人闪耀并呈现，理由是只有这个年龄才能超越肉体，并正确地从字面上升到精神。\par

49. 然而在我们中间，施教与受教之间并没有界限，如同古时约旦河内外支派之间的石头一般；也不是某部分托付给一些人，另一部分给另一些人；也没有关于经验等级的任何规则；这件事已经如此混乱不堪，以至于我们大多数人——不敢说全部——几乎在我们还未失去童稚的卷发和口齿不清以前，在我们还未进入天主的殿宇以前，在我们甚至还不知道圣经各卷的名称以前，在我们还未学到旧约和新约的特点和作者以前（因为我目前所指的，并非我们未能洗净因邪恶而沾染的灵性羞耻的泥泞和痕迹），如果我们，我说，仅仅通过道听途说而非研读，为自己预备了两三个虔诚作者的说法；如果我们对\Person{达味}有了短暂的经验，或者正当地穿上了一件小袍，或者至少系上了哲学家的腰带，或者在我们周围束上了某种形式的\OriginalTerm{虔诚}{piety}的外表——呸！我们何等轻易地坐上椅位，展示我们的精神！\Person{撒慕尔}即使在襁褓中也是圣洁的。\BibleRef{撒上 2:11}我们立刻成为智慧的教师，在神圣事物上地位崇高，成为经师和法学士中的佼佼者；我们自封为天上之人，并寻求被人称为\OriginalTerm{辣彼}{Rabbi}。\BibleRef{玛 23:7}字句无处可寻，一切都要按属灵的意义来理解，而我们的梦完全是胡言乱语，如果我们不被过分赞扬，我们就会恼怒。这是我们中较好较单纯的人的情况：那些更属灵、更高贵的人呢？他们常常谴责我们，视我们为无用之人，然后离弃我们，厌恶与我们这样不虔敬的人交往。\par

50. 现在，如果我们温和地对其中一人说话，逐步前进，如下论证：\Quote{告诉我说，我的好朋友，你称舞蹈和吹笛子算什么呢？} \Quote{当然算。} 他们可能会说。\Quote{那么智慧和有智慧呢？我们敢将其定义为对神圣和人类事物的知识。} 他们也会承认。\Quote{那么这些技艺比智慧和智慧更好更优越，还是智慧远胜于这些？} \Quote{甚至胜过一切，} 我知道他们会说。到此为止他们还算明智。\Quote{那么舞蹈和吹笛子需要教学和学习，这个过程需要时间，许多辛劳和汗水，有时还要付费，恳求开门，长期离家，以及为获得经验所必须做和忍受的一切：但是\OriginalTerm{智慧}{wisdom}，它是万物的首要，怀抱一切美善，甚至天主自己也更喜欢这个称号胜于祂被称为的一切名字；难道我们要认为它是如此微不足道和如此容易获得，以至于我们只要愿望就能有智慧吗？} \Quote{这样做是极其愚蠢的。} 如果我们或任何有学问和谨慎的人对他们这样说，并逐步尝试使他们从错误中洁净，那将是撒在石头上，\BibleRef{路 8:6} 对不愿听的人说话。\BibleRef{德 25:9} 他们离连自己的无知都认识不到的智慧如此遥远。我认为我们可以恰当地对他们引用\Person{撒罗满}的话：我看见在太阳之下有一种邪恶，\BibleRef{训 10:5} 一个自以为有智慧的人；\BibleRef{箴 26:12} 更大的邪恶是让一个连自己的无知都不知道的人去教导别人。\par

51. 这种心态尤其需要我们的眼泪和叹息，常常激起我的怜悯，因为我确信想象在很大程度上剥夺了我们现实，而虚浮的荣耀是人们在达到\OriginalTerm{德行}{virtue}道路上的巨大障碍。要医治和制止这种疾病，需要有\Person{伯多禄}或\Person{保禄}，那些\Person{基督}的伟大门徒，他们除了在言语和行为上的引导外，还领受了他们的\OriginalTerm{恩宠}{grace}，成为一切人中的一切，为要赢得一切人。\BibleRef{格前 9:22} 但对我们其他这样的人来说，被那些负责纠正和整顿此类事情的人正确引导和带领，就是一件大事。\par

52. 然而，既然我提到了\Person{保禄}和他那样的人，请允许我略过所有其他作为立法者、先知、领袖或任何类似职务的首要人物——例如\Person{梅瑟}、\Person{亚郎}、\Person{若苏厄}、\Person{厄里亚}、\Person{厄里叟}、民长、\Person{撒慕尔}、\Person{达味}、先知团体、\Person{若翰}、\Person{十二宗徒}以及他们的继任者，他们以许多辛劳和劳动在各自的时期内行使了权柄；所有这些我都略过，以便提出\Person{保禄}作为我主张的见证人，并让我们通过他的榜样来思考灵魂的牧养是何等重要的事情，以及它是否需要轻微的关注和少量的判断。但为了使我们认识并察觉到这一点，让我们听听\Person{保禄}自己关于\Person{保禄}所说的话。\par

53. 我不提他的劳苦、警醒、在饥渴、寒冷和赤身中的受苦、外在的攻击者、内在的反对者。我略过迫害、议会、监狱、锁链、控告者、审判台、每日每时的死亡、筐子、石头、棍打、旅行、陆地和海上的危险、深渊、船难、河流的危险、强盗的危险、来自同胞的危险、在假弟兄中的危险、亲手劳作、白白传福音、成为天使和人类注视的焦点、立于天主与人之间为祂辩护、将他们结合于祂、使他们成为祂自己的子民，\BibleRef{铎 2:14} 除了那些外面的事。\BibleRef{格后 11:28} 因为谁能详尽地述说这些事呢：每日的压力、个别的挂虑、对众教会的关怀、普遍的同情、兄弟之爱？有人跌倒，\Person{保禄}也软弱；有人受诽谤，\Person{保禄}就焦急如焚。\par

54. 他的教导之辛劳如何？他职务的多样性如何？他的慈爱如何？另一方面他的严厉如何？两者的结合与调和如何？这样他的温和不致削弱人，他的严厉也不致激怒人？他为奴隶与主人\BibleRef{弗 6:5}、统治者与被统治者\BibleRef{罗 13:1}、丈夫与妻子、父母与子女\BibleRef{弗 6:1}、婚姻与守贞、自律与放纵、智慧与无知、割礼与不割礼、基督与世界、肉身与圣神制定律法\BibleRef{迦 5:16}。他为一派人感恩，另一派人他责备。有人他称为自己的喜乐与冠冕\BibleRef{斐 4:1}，另一派人他斥为愚昧\BibleRef{迦 3:1}。有人持守正道他同道同行，分担他们的热忱；有人偏离正路他加以拦阻。时而施以绝罚\BibleRef{格前 5:5}，时而重申他的爱\BibleRef{格后 2:8}；时而忧伤，时而喜乐；时而以奶喂养，时而论及奥秘；时而屈尊俯就，时而提携至己身同列；时而以棍棒威胁，时而展现温柔之灵；时而对高傲者傲慢，时而对卑微者谦卑。忽而他是宗徒中最小的，忽而他证明基督在他内说话\BibleRef{格后 13:3}；忽而他渴望离世并已被奠祭，忽而他认为为他们的缘故存留肉身更为必要。因为他寻求的不是自己的利益，而是他儿女的利益\BibleRef{格前 10:33}，就是他在基督内藉福音所生的。这是他全部神权之目的：凡事不顾自己而顾及他人益处。\par

55. 他以自己的软弱和患难为荣。他以耶稣的死为乐，仿佛那是某种装饰。他在属血肉的事上高傲，在属神的事上喜乐；他在知识上并非粗野\BibleRef{格后 11:6}，并声称如同在镜中观看，模糊不清\BibleRef{格前 13:12}。他在心神上勇敢，并痛击自己的身体，将它当作对手抛掷。他以此要给我们什么教训和指导？不要为属世的事物骄傲，不要因知识自大，不要激动肉身反抗圣神。他为众人奋斗，为众人祈祷，为众人热诚，为众人燃烧，无论他们不在法律之下，或在法律之下；外邦人的宣讲者\BibleRef{弟后 1:11}，犹太人的保护者。他甚至为同族的弟兄极其勇敢——若我可以冒昧这样说——在他充满爱德的祈祷中，愿他们能代替他被带到基督面前\BibleRef{罗 9:3}。何等宽宏！何等热忱！他效法基督，基督为我们成了诅咒\BibleRef{迦 3:13}；他承担了我们的软弱，背负了我们的疾病\BibleRef{玛 8:17}；或借用更审慎的说法，他准备在基督之后，为众人忍受一切，甚至如同不敬神的人一样，只要他们得救。\par

56. 我何必细数？他活着不是为自己，而是为基督和他的宣讲。他将世界钉在十字架上归给自己\BibleRef{迦 6:14}，而他自己也被钉死在世界上和可见的事物上，他视万事为微小\BibleRef{斐 3:8}，且微小到不值得渴望；尽管从耶路撒冷及其周边直到依里黎康\BibleRef{罗 15:19}，他已圆满宣讲了福音，尽管他曾被提前提到三重天上，并看见了乐园的异象，听到了不可言说的话语。保禄如此，凡与他同灵性的人也如此。但我们担心，与他们相比，我们可能成为左安愚蠢的王子\BibleRef{依 19:11}，或是勒索者，榨取地里的出产，或虚假地祝福人民；并且进一步自娱自乐，扰乱你们的行径，或是统治你们的嘲弄者，或是掌权的孩童，心智不成熟，甚至连足以治理任何人的衣食都没有；或是教导谎言的先知，或是反叛的王子，因饥荒的困迫而应分享他们长老的耻辱，或是远未能对耶路撒冷说安慰话的司祭，依照依撒意亚所发出的责备与抗议，他曾被色辣芬用火炭洁净。\par

57. 既然如此，对于一颗敏感而忧伤的心，这项任务是如此沉重而艰巨吗？——对于一个明理之人，这简直是骨中的朽烂\BibleRef{箴 14:30}：而危险却很轻微，跌倒也不值一提吗？不，有福的欧瑟亚使我深感警醒，他说审判属于我们司铎和统治者\BibleRef{欧 5:1}，因为我们成了瞭望台的罗网；如同在大博尔山上铺开的网，由猎取人灵魂的猎人牢牢固定，他威胁要剪除邪恶的先知，用火吞灭他们的审判官，并暂时停止膏立君王和王子，因为他们为自己统治，而非为祂。\par

58. 因此，神圣的米该亚不能容忍以血建造熙雍——无论你如何解释这句话——也不能容忍以不义建造耶路撒冷，而其中的首领因贿赂审判，司铎为酬劳教导，先知为金钱占卜——他说这会有什么结果？熙雍必被耕耘如同田地，耶路撒冷必成为园中的茅舍，圣殿的山必成为丛林中的空地\BibleRef{米 3:10}。他也哀叹正直人的稀少，几乎没有一根茎或一串残余的葡萄留下，因为王子索取，审判官讨好，以至于他的话语几乎与伟大的达味相同：\Quote{上主，求祢拯救我，因为虔诚人断绝了}，并且说因此他们的祝福必逝去，如同被虫蛀蚀。\par

59. 岳厄尔再次召唤我们哀哭，并要求祭台的侍从在饥荒面前哀悼：他绝不容许我们在他人的不幸中狂欢；而且，在祝圣禁食、召集圣会、聚集老人、孩童和幼童之后，我们自身还必须披麻蒙灰\BibleRef{依 58:5}，出没于圣殿，极其谦卑地俯伏在地，因为田地荒废，素祭和奠祭从上的殿中断绝，直到我们因谦卑而赢得怜悯。\par

60. 哈巴谷如何呢？他发出更激烈的话语，甚至对天主自己不耐烦，并且谴责——仿佛在谴责——我们的善主，因为审判官的不义。他说：\Quote{上主，我呼求，祢不垂听，要到几时呢？我向祢喊说\InnerQuote{残暴}，祢却不拯救？祢为何让我看见辛劳与劳苦，使我注视乖谬与不虔？审判已对我不利，审判官成了掠夺者。因此法律松弛，公义从未显明。}随后便是宣告，以及随之而来的一切：\Quote{看，你们这些轻慢的人，要观看，要惊奇，要诧异，并且消失吧，因为我正在做一件事。}但我何必引用整个宣告呢？然而，稍后一点——因为我认为最好将这一点加在已说的话上——在谴责和哀叹许多在某方面不义或败坏的人之后，他谴责邪恶的首领和教师，将恶习斥为污秽的紊乱、醉酒和心灵错乱；指控他们给邻人喝酒，以便观看他们灵魂的黑暗，以及爬虫和野兽的洞穴，即邪恶思想的居所。它们的确是这样，并且他们与我们讨论这样的教训。\par

61. 怎能合适地略过玛拉基亚呢？他时而严厉控告司铎，指责他们轻视上主的名（\BibleRef{拉 1:6}），并说明他们如何这样做：在祭台上奉献污秽的饼和不是初果的肉——这些是他们不会向自己的一位总督奉献的，或者如果他们献了，也会被羞辱——然而却将这些作为还愿祭献给宇宙之王，即瘸腿的、有病的和畸形的，这些是全然不洁和可憎的。他又提醒他们天主的盟约，即与肋未子孙所立的生命与平安的盟约，以及他们应敬畏事奉祂，并因祂圣名的显现而战栗。他说：\Quote{真理之律在他口中，不义没有在他唇上找到；他以平安正直与我同行，并使许多人转离不义：因为司铎的唇应保存知识，人应从他的口寻求法律。}这原因何等可敬，同时又何等可畏！因为他是全能上主的使者。虽然我略过以下的咒骂，因它们措辞强烈，但我惧怕它们的真实。然而，这一点可以无碍地引用，以利于我们。他说：\Quote{难道应从你们手中接纳你们的祭品，并善意地接受吗？}就好像他极其愤怒，因他们的邪恶而拒绝他们的职分。\par

62. 每当我想起匝加利亚，我就因那收割的镰刀而战栗，同样也因他对司铎的见证、他对那位著名的耶叔亚大司铎的暗示而战栗——他描述他脱去污秽不当的衣服，然后穿上华美的司铎袍。至于天使对耶叔亚所说的话和吩咐，应当默然敬重，因为它们或许指向一个比那些众多司铎更高更大的对象（\BibleRef{希 7:23}）；但在他右边站着魔鬼，要抵挡他。在我看来，这个事实意义重大，需要不小的畏惧和警醒。\par

63. 谁的心如此刚硬如金刚石，以至对故意对余下牧者提出的控告和指责不颤抖、不羞愧呢？他说：\Quote{有牧者哀号的声音，因为他们的荣耀被摧毁了。有狮子咆哮的声音}（\BibleRef{匝 11:3}），因为这事临到了他们。他岂非几乎听到仿佛近在眼前的哀哭，并自己与受苦者一同哀哭吗？稍后一点，有一段更引人注目、更充满激情的说话：\Quote{喂养那被宰杀的羊群，它们的占有者宰杀它们而不悔改，卖它们的人说：\InnerQuote{愿上主受赞美，因为我们富足了！}而它们自己的牧者对它们没有感觉。因此，我将不再怜惜这地的居民。}这是上主全能者说的（\BibleRef{匝 11:5}）。又说：\Quote{刀剑，醒来吧，攻击牧者，击打牧者，使羊群四散，我将向牧者伸手；我的怒气向牧者发作，我将眷顾羔羊。}在威胁中又加上那些统治人民的人。他如此辛勤地致力于他的任务，以致难以轻易从谴责中脱身，我担心，如果我引用全部系列，会使你们失去耐心。有关匝加利亚，这就够了。\par

64. 略过达尼尔书中的长老们——因为最好略过他们，连同上主关于他们的公义判决与声明：邪恶从巴比伦而来，来自那些看似治理人民的古老审判官——我们如何被厄则克耳感动呢？他是伟大奥秘与异象的目睹者和阐释者。他命令守望者（\BibleRef{则 33:2}）不要对恶习及悬在其上的刀剑缄默不语——那样既不会有利于他们自己，也不会有利于罪人——而是宁愿保持警醒并预先警告，从而至少有益于那些发出警告的人，如果不同时有益于说话者和听者的话。\par

65. 他对牧者的进一步抨击又如何？\Quote{祸患将临于祸患，谣言将接续谣言，那时他们将向先知寻求异象，但法律必从司铎那里消失，计谋必从长老那里消失。}又这样说：\Quote{人子啊，你要对她说：你是未被浇灌的土地，在忿怒之日雨没有降在你身上；她中间的王子好像咆哮的狮子，撕裂猎物，以强力吞噬灵魂。}稍后一点：\Quote{她的司铎违反了我的法律，亵渎了我的圣物，他们不分别圣与俗，一切都对他们一样，他们掩眼不顾我的安息日，我在他们中间受了亵渎。}他威胁要毁灭那墙和那粉刷墙的人，即那些犯罪的人和那些掩盖他们罪过的人——正如恶的统治者和司铎所行的，他们使以色列家随从自己的心意而走入歧途，这些心意在他们欲求中疏远了（\BibleRef{则 14:5}）。\par

我也避免进入他对那些牧人的讨论：他们牧养自己，吞吃奶，穿羊毛，宰杀肥壮的，却不牧养羊群；没有扶助软弱的，没有包扎受伤的，没有领回迷失的，没有寻找失落的，没有看管强壮的，反而以严厉辖制他们，以压力毁坏他们；以致因为没有牧人，羊群四散在全平原和山上，成了所有飞鸟和走兽的食物，因为没有人寻找他们，领他们回来。结果如何？\Quote{主说：我指着我的永生起誓，因为这样的事发生了，我的羊群成了掠物，看，我必与牧人为敌，我必从他们手中追讨我的羊，并要聚集他们，使他们成为我的羊；但牧人必遭受如此这般的事，如恶牧人所应得的。}\par

然而，为了避免不合理的延长我的讲词，逐一列举所有先知和他们的话，我只再提一位，他在成形以前就被认识，并从母胎中被圣化（见：\BibleRef{耶 1:5}），那就是\Person{耶肋米亚}；其余的都略过。他渴望水浇他的头，泪泉流他的眼，好使他能充分地哀哭以色列；同样他也哀叹其统治者的堕落。\par

天主向他发言，斥责司祭：\Quote{司祭没有说：主在哪里？掌管法律的也不认识我；牧人也背叛了我。}祂又对他说：\Quote{牧人变得愚昧，没有寻求主，因此他们所有的羊群都不明白，都四散了。}又说：\Quote{许多牧人毁坏我的葡萄园，玷污了我所喜悦的产业，直至它成为没有路径的荒野。}祂又进一步斥责牧人：\Quote{祸哉，那牧养我的羊群，却毁坏和驱散它们的牧人！因此主论到喂养我子民的人如此说：你们驱散我的羊群，赶走它们，没有眷顾它们；看，我必因你们的行为降祸与你们。}此外，祂吩咐牧人哀号，羊群中的公羊哀哭，因为他们被宰杀的日子已经成就了。\par

我何必谈论古代的事？谁能用保禄为主教和长老所定的规则和标准考验自己，即他们要节制、清醒、不贪酒、不斗殴、善于教导、在一切事上无可指摘、不为恶人所及，而不发现自己在规则的正线上有相当大的偏离呢？耶稣对祂的门徒的吩咐又如何，当祂派遣他们去宣讲时？这些吩咐的主要目的——不细述细节——是使他们具有如此的美德，如此纯朴、谦逊，总之如此属天，以致福音得以进展，不亚于藉着他们的品行，也藉着他们的宣讲。\par

我被法利塞人的指责、经师的定罪所惊吓。因为我们本该大大超过他们，正如我们被命令的，如果我们渴望天国，却发现自己更深地陷入恶习，这是可耻的：以致我们该被称为蛇、毒蛇的种类、瞎眼的向导，滤出蚊虫却吞下骆驼，或是内里污秽的坟墓，尽管外表光鲜，或是外面清洁的盘子，以及他们所是或所受控诉的一切。\par

我昼夜被这些念头占据：它们消耗我的骨髓，吞噬我的肉体，不让我自信或昂首。它们压抑我的灵魂，贬抑我的心思，束缚我的舌头，使我思考的不是主教之位，或引导和指导他人——这远超我的能力——而是我自己如何逃避将来的忿怒，并从自己身上刮去一些恶习的铁锈。一个人必须先洁净自己，然后才能洁净他人；自己先成为智慧，然后才能使人有智慧；自己先成为光，然后才能发光；自己先亲近天主，然后才能带领他人亲近；自己先被圣化，然后才能圣化他人；自己有手，才能用手引导他人；自己有智慧，才能给人建议。\par

这事何时发生？那些凡事敏捷却不稳妥、随时建造又随时拆毁的人说。灯何时放在灯台上？塔冷通在哪里？因为他们如此称呼恩宠。这样说的人在友谊上比在敬畏上更热心。你们这些极其勇敢的人问我，这些事何时发生，我如何交代？即使是最高的年老也不算是过长的期限。因为白发加上审慎胜过无经验的青年，深思熟虑的犹豫胜过轻率的急促，短暂的统治胜过漫长的暴政：正如一份以荣誉赢得的小份额胜过庞大而不名誉且不确定的财产，一点金子胜过大量的铅，一点光明胜过许多黑暗。\par

但这种匆忙，因其不可靠和过于急促，有危险像落在石头上的种子（见：\BibleRef{路 8:6}），因为它们没有深厚的土壤（见：\BibleRef{玛 13:5}），立即发芽，却不能忍受太阳初次的热；或者像建造在沙土上的基础，连轻微的雨和风都不能抵挡。祸哉，你这城，其君王是孩童（见：\BibleRef{训 10:16}），撒罗满说。不要在言语上急躁（见：\BibleRef{箴 29:20}），撒罗满又说，声称言语的急躁不及行动的热烈严重。然而，有谁不顾这一切，要求匆忙胜过安全和效益呢？谁能像在一天内塑造泥像那样塑造真理的捍卫者？他将与天使一同站立，与总领天使一同光荣，使祭献上升到高处的祭台，并分享基督的司祭职，更新受造物，刻划肖像，为上界创造居民，而且——最伟大的——成为天主，并使他人成为天主？\par

74. 我知道我们是何等的仆役，身处何地，引导何人。我知道天主的崇高，人的软弱，以及相反，人的能力。天高，地深；凡被罪恶倾覆的，谁能升上？凡尚被此下界的幽暗和肉身的粗重所环绕的，谁能以全心纯然凝视那全心，并在不稳定与可见的事物中与稳定而不可见者交往？因为，即使是那些最受过净化的人，也很难在此看见\Quote{美}的形象，如同人在水中看见太阳。谁曾用手量过水，用掌测过天，用秤称过大地？谁曾用天平称过山岳，用秤称过丘陵？\EditorialNote{依撒意亚 40:12}祂安息之处何在？又谁能与祂相比？\par

75. 是谁以自己的圣言创造万有，以自己的智慧塑造了人，将分散的聚拢为一，将尘土与神气混合，组合成一个可见又不可见、暂时又不朽、属土又属天的生灵，能接近天主却不能领悟祂，既靠近又遥远？我曾说：\Quote{我要成为智慧的}——撒罗满说——\Quote{她却远超过我所能及的}；\BibleRef{训 7:24}又说：的确，增加知识，增加忧愁。因为我们所发现的喜乐，并不多于我们所遗漏的痛苦；我想，这痛苦就像那些被从水中拖出、仍感干渴，或无法留住自以为抓住的东西，或被闪电骤然抛入黑暗的人一样。\par

76. 这些事使我沮丧，使我谦卑，并说服我：听赞美的声音，胜于解释超出我能力之外的真理。祂的威严、崇高、尊贵，以及那些几乎不能容纳天主光辉的纯洁本性——深渊掩盖祂，祂的隐秘处是黑暗，因为祂是最纯粹的光明，\BibleRef{弟前 6:16}是多数人不能接近的；祂既在此宇宙之中，又超越宇宙；祂是全部的美善，\BibleRef{出 33:19}又超越全部的美善；祂光照理性，却又逃脱理性的敏捷与高度，总在被领悟时退去，借着逃离和溜走，将爱慕祂的人引向更高之处。\par

77. 我们所热切渴望的对象是如此伟大，而那位预备并引导灵魂赴婚筵的人也当如此。我害怕被捆起手脚扔出去，\BibleRef{玛 22:13}因为我未穿婚宴礼服，就冒然闯入那些坐席者中间。然而，我自幼就被邀请——若我可谈及多数人不知道的事——从母胎中就被交付给祂，因我母亲的许诺，后来在危险时刻得以确认。我的渴望随之增长，理性与之认同，我将自己的一切——财产、名声、健康、甚至言语——献给了我那位赢得并拯救我的主，我只从中获益，即能轻视它们，并拥有一些使我能以基督为优先的东西。天主的言语对我如蜜房甘甜，我呼求知识，扬声寻求智慧。\BibleRef{箴 2:3}此外，还有抑制愤怒、约束舌头、节制眼目、管束肚腹、践踏那附于地上的荣耀。我说话愚昧，\BibleRef{格后 11:23}但不得不说：在这些追求上，我或许不逊于许多人。\par

78. 然而，有一种哲学过于高深，非我所能及，那就是引导和治理灵魂的职责——在我尚未学会正确服从牧者，或灵魂尚未得到适当净化之前，就承担牧养羊群的重任。尤其是在这样的时代：当人看见众人混乱地东奔西跑，宁愿逃离混战，在隐退的庇护中躲避恶者的风暴与阴霾；当肢体彼此相争，那曾存留的少许爱德已逝去，司铎只是一个空名，因为，如经上所记，轻蔑倾注在首领身上。\par

79. 但愿它仅仅是空名！如今，愿他们的亵渎归于不敬虔者的头上！一切敬畏已从灵魂中被驱逐，无耻取而代之，而知识和圣神的深奥事理\BibleRef{格前 2:10}向任何愿意者开放；我们全都因谴责他人的不敬虔而成为虔诚者；我们求助于不敬虔的审判官，将圣物丢给狗，把珍珠扔在猪前，\BibleRef{玛 7:6}在亵渎的灵魂面前宣扬神圣之事；我们这些可怜虫，仔细履行敌人的祈祷，毫不羞耻地以自己的发明行淫。摩阿布人和阿孟人，甚至不被允许进入上主的教会，\BibleRef{申 23:3}却常参与我们至圣的礼仪。我们向所有人敞开的，不是正义之门，而是毁谤和党派傲慢之门；在我们中间，首位的不是那因敬畏天主而连闲话都避免的人，而是那最能流利地辱骂邻人的人，无论是明说还是暗讽；他在舌头下滚动着祸害和不义，或者更准确地说，是毒蛇的毒素。\par

80. 我们观察彼此的罪过，不是为了哀哭，而是为了加以责备；不是为了医治，而是为了加重；并借邻人的伤口为自己恶行开脱。好人与坏人的区分，不是根据个人品行，而是根据他们是否与我们意见一致或友好。我们今日赞美的，明日就诋毁；从别人那里受谴责，反成为我们钦佩的通行证。我们在恶行上竟如此宽宏大量，以至不敬虔的一切都被坦然赦免了。\par

81. 一切都回到了世界以其现有美好秩序与形态形成之前的原始状态，\newline{}\BibleRef{创 1:2}。普遍的混乱与失序呼求某种组织之手与能力。或者，若你愿意，它好比一场月淡星稀之夜的交战，无人能辨认敌友的面容；又好比波涛上的海战，狂风劲吹，浪花沸腾，波涛撞击，船只粉碎，长竿戳刺，水手长笛鸣响，倒下者的呻吟，我们在喧嚣之上提高嗓音，不知如何是好，且可叹没有机会展示我们的勇武，彼此攻击，倒在彼此手中。\par

82. 事实上，民众与司铎之间也无区别：在我看来，这不过是古代咒诅的简单应验：\BibleQuote{百姓如何，司铎也如何。}再者，伟大显赫之人与大众并无二致；不，他们公开与司铎为敌，而他们的虔诚反而助长了其说服力。的确，只要是为信仰，为最高且最重要的问题，让他们如此行事，我并不责备；不，说实话，我甚至赞美并祝贺他们。是的！我宁愿自己是那为真理而争战并招致仇恨者之一：或者说，我可以夸耀自己就是其中之一。因为一场可称颂的战争，胜过使人远离天主的和平：因此圣神装备那温和的战士，作为能在正义事业中争战的人。\par

83. 但如今，有些人甚至为小事且无目的地争战，并以极大的无知和冒失，接受任何人为其恶行的同伙。于是每个人都以信仰为托辞，这可敬的名号被拖入他们的私斗中。结果，正如所料，我们甚至在外邦人中遭恨，更难受的是，我们不能说这毫无道理。不，就连我们中最优秀的人也感到困扰，而在那些不乐于接受任何善事的群众中，这一结果并不奇怪。\par

84. 罪人在我们背上谋划；我们彼此算计之事，他们转过来对付我们全体：我们成了新的景象，不仅是对天使和世人，\BibleRef{格前 4:9}正如那最勇敢的运动员\Person{保禄}在与率领者和掌权者\BibleRef{弗 6:12}争战时所说，而且几乎对所有恶人，在每一时每一地，在广场上、在宴饮中、在庆典中、在哀伤时。不，我们已经——我几乎含泪而言——被搬上舞台，在放荡者的笑声中，最受欢迎的对白与场景便是对基督徒的丑化。\par

85. 这些便是我们内战的后果，以及我们过于急切地争夺良善与温和，还有我们对天主不合时宜的过分挚爱。摔跤或任何其他竞技，只允许按照固定规则进行；若有人破坏摔跤规则，或在任何其他竞赛中违反竞赛制定的规则行事，无论他如何有能有力，都将被喝倒彩、蒙羞并失去胜利；那么，若有人以非基督的方式为基督争战，却因以它的名义非法作战而取悦和平，又岂能蒙悦纳？\par

86. 是的，即使在现今，当基督被呼求时，魔鬼战栗，\BibleRef{雅 2:19}甚至我们的恶行也未熄灭这名的能力，而我们却不以侮辱如此可敬的事业和名为耻；几乎在公开场合，每日呼喊它并被人回喊；因为\BibleQuote{我的名因你们在外邦人中受了亵渎。}\par

87. 我并不惧怕外在的战争，也不惧怕那如今兴起攻击教会的野兽与邪恶的满溢，尽管他可能威胁以火、剑、野兽、悬崖、深渊；尽管他可能比以往所有狂人更不人道，并发现更严酷的新酷刑。我对这一切只有一个疗方，一条胜利之路：我要在基督内夸耀，\BibleRef{斐 3:3}即为了基督而死。\par

88. 然而，对于我自己的争战，我茫然不知该走何路，该求何同盟，该用何智慧之言，该设想何恩宠，该以何全副武装自卫，以对抗恶者的诡计。\BibleRef{弗 6:11}是\Person{梅瑟}通过在山上的举手\BibleRef{出 17:11}来征服他，以便那如此预表与预象的十字架得胜？还是\Person{若苏厄}作为他的继任者，与上主军队的统帅\BibleRef{苏 5:14}一同列阵？是\Person{达味}，或以弹琴，或以甩石争战，并被天主束腰备战，手指受训作战？是\Person{撒慕尔}为百姓祈祷\BibleRef{撒上 7:5}并献祭，膏立那能得胜者为王？是\Person{耶肋米亚}为以色列书写哀歌，应当哀叹这些事？\par

89. 谁会大声呼喊：\BibleQuote{上主，求祢宽恕祢的百姓，不要让祢的产业受羞辱，让外邦人统治他们？}\BibleRef{岳 2:17}什么\Person{诺厄}、\Person{约伯}和\Person{达内尔}，他们被一同算为祈祷之人，会为我们祈祷，使我们暂得喘息，恢复过来，一时彼此认出，不再作为联合的\Place{以色列}，而是\Place{犹大}和\Place{以色列}，\Person{勒哈贝罕}和\Person{雅洛贝罕}，\Place{耶路撒冷}和\Place{撒玛黎雅}，因我们的罪依次被交付，又依次被哀悼。\par

90. 因为我承认我太软弱，无法胜任这场争战，因此转身，在溃败中掩面，独自静坐，\BibleRef{哀 3:28}因我充满苦楚，寻求沉默，知道这是邪恶的时代，\BibleRef{米 2:3}蒙爱者已踢跳，\BibleRef{申 32:15}我们成了叛离的子女，\BibleRef{耶 3:14}是茂盛的葡萄树，\BibleRef{欧 10:1}真葡萄树，全都结果，全都美丽，从天降的雨水中蓬勃生长。因为美丽的冠冕，\EditorialNote{依撒意亚 62:3}荣耀的印记，\BibleRef{则 28:12}威荣的王冠，对我已变为羞耻；若有人面对这些事仍勇敢无畏，我祝福他的勇敢无畏。\par

91. 关于我们自身内部的战争，以及我们昼夜与那使我们卑微的身体（\BibleRef{斐 3:21}）争战的私欲，无论明暗，以及那借感官和此生其他快乐之源将我们抛来抛去的浪潮；关于我们被陷入的泥泞；关于那与心神之法律交战并力图摧毁我们内天主的肖像以及赐予我们的一切神圣流露的罪恶法律（\BibleRef{罗 7:23}），我尚未言及。因此，任何人若未经过长期的哲学训练，逐步将灵魂中高贵而光明的部分与卑下、与黑暗结合的部分分离，或未借天主的仁慈，或两者并用，并恒常举目仰望，以克服物质的压抑力量，则很难胜过这一切。而在人尚未尽可能获得此优胜、充分净化自己的理智、并在亲近天主方面远超同伴之前，我认为将管理灵魂的权责或天人之间的中保职务（因为我认为司铎正是如此）托付给他，是不安全的。\par

92. 使我心生此畏惧的原因是什么——使你们不致以为我无端恐惧，反而称赞我的远见？我从梅瑟亲口所述听到：当天主对他说话时，虽然许多人被召上山，其中甚至包括亚郎和他的两个司铎儿子，以及七十位长老，其余的人却奉命远远朝拜，只有梅瑟一人可亲近，百姓不得与他同去（\BibleRef{出 24:1}）。因为并非人人都可亲近天主，只有那如同梅瑟一样能承受天主荣耀的人才可以。此外，在此之前，当法律初次颁布时，号角声、闪电、雷声、黑暗、全山的烟雾，以及甚至连牲畜触摸山也要被石头砸死的可怕警告（\BibleRef{希 12:18}）和其他类似的威慑，使其余百姓退却；对他们而言，经过仔细净化后，仅能听到天主的声音已是极大的恩典。但梅瑟却上山进入云中，被付予法律，领受了两块石版——对大众而言是文字，对超脱大众的人而言是心神（\BibleRef{格后 3:6}）。\par

93. 我又听说纳达布和阿彼胡仅仅因为奉献了凡火而被异火烧灭（\BibleRef{肋 10:1}），他们不敬的器具成了惩罚的工具，在亵渎的当场即时毁灭；甚至他们的父亲亚郎——在天主眼中仅次于梅瑟——也不能救他们。我也知道司铎厄里，以及稍后的乌匝；前者为他儿子的过犯付出代价——他们胆敢用不当方式强索鼎中的初熟之物以玷污祭献，尽管厄里并未纵容他们不敬，反而屡次责备；后者只因用手扶住被牛倾翻的约柜（\BibleRef{撒下 6:6}），虽保存了约柜，自己却因天主的嫉妒（为维护约柜当受的敬畏）而被消灭。\par

94. 我也知道，无论是司铎身上的残疾（\BibleRef{肋 21:17}）还是祭品上的瑕疵，都不能被忽略；法律要求完美的祭品须由完美的人奉献——我认为这象征灵魂的完整。并非人人都可触摸司铎的衣袍或任何圣器；祭品本身除适当的人、在适当的时间和地点外，不得取用；不可仿制傅油膏或配制的香（\BibleRef{出 30:33}）；凡进入圣殿者，必须在灵魂和身体上毫无细微的不洁；至圣所远离妄加接近，以至于只许一人一年一次进入；帷幔、赎罪盖、约柜和革鲁宾，都远避众人的目光和触摸。\par

95. 既然我知道这些事，并且知道任何人若未先将自己奉献给天主，作为生活的、圣洁的祭品，并呈上合理的、中悦天主的敬拜（\BibleRef{罗 12:1}），以赞颂的祭献和痛悔的心灵向天主献祭——这正是那赐予万有的主向我们索要的唯一祭品——我又怎敢向祂献上外在的祭品（那伟大奥迹的预象），或披上司铎的外衣和名号？除非我的手已因圣善的作为而被祝圣；我的眼已习惯于只以惊叹注视受造物，只归于造物主，而不伤害受造物；我的耳已被主的教训充分开启，祂开了我的耳，使我毫无沉重地聆听，并将金耳环配以宝贵的红玉髓——即智慧人的言语置于顺服之耳中；我的口已被打开以吸取圣神，张开阔大以充满讲述奥迹和道理的神恩（\BibleRef{格前 14:2}），我的嘴唇被神圣的知识所约束，以便使用智慧之言——并且，我要加上一句，在适当的时候解开；我的舌头充满了欢乐，成为神圣旋律的乐器，清晨醒来警醒，勤劳直到贴于上颚；我的脚立于磐石，如同母鹿的蹄，我的步履被导向属神的样式，不致滑跌，甚至几近滑跌也不致滑跌；我的一切肢体都成了正义的武器（\BibleRef{罗 6:13}），所有的可朽性都被脱去，被生命吞灭（\BibleRef{格后 5:4}），且已屈服于圣神？\par

96. 有哪一个人，心从未因翻开圣经而燃烧，\BibleRef{路 24:32}，因那些在炉中炼净的天主纯言；有哪一个人未曾藉着将这些话三度写在宽广的心版上，而获得基督的心意，\BibleRef{格前 2:16}，也未曾被引入那对多数人仍隐藏、隐秘、幽暗的宝藏，以凝视其中的财富？\EditorialNote{依撒意亚 45:3}并能以属灵的事比较属灵的事，而使人富足。\BibleRef{格前 2:13}\par

97. 有哪一个人，从未依照我们的本分瞻仰上主的美善，未曾进入祂的圣殿——或者更确切地说，成为天主的宫殿\BibleRef{格后 6:16}和基督在圣神内的居所？\BibleRef{弗 2:22}有哪一个人，未曾认识预象与真理之间的关联与区别，从而脱离前者、紧贴后者，藉此逃离字面的旧约、服侍精神的新约，能洁净地从法律过渡到恩宠——法律在身体消解中得以精神性地满全？\BibleRef{罗 6:6}\par

98. 有哪一个人，未曾借着经验和默观，遍历基督的全部称号和权能的系列——包括祂原本所拥有的更高超的，以及后来为了我们而取用的更卑微的：即天主、圣子、肖像、圣言、智慧、真理、光明、生命、大能、蒸气、流露、光辉、造物主、君王、元首、法律、道路、门、根基、磐石、珍珠、和平、正义、圣化、救赎、人、仆人、牧人、羔羊、大司祭、祭品、在受造物中之首生者、从死者中之首生者、复活？有哪一个人，听见这些满含实意的名号却不留意，从未与圣言在任何由祂所负的每个名号所标志的真实关系中相交，也未曾有分于祂？\par

99. 最后，有哪一个人，虽然从未致力于或学会讲论天主奥秘中的隐藏智慧，\BibleRef{格前 2:17}虽然仍是婴孩，仍以奶为食，仍属于未在以色列中被计数的人，\BibleRef{户 1:3}也未登记在天主的军队中，虽然他尚未能像男子汉一样背起基督的十字架，虽然可能仍是较为不尊贵的肢体，却会喜乐而热切地接受被任命为基督满全的首领？\BibleRef{弗 1:23}没有一个人，若他听从我的判断并接受我的劝告！这是众事中最可畏惧的，在每一个明了成功之重大与失败之彻底毁灭的人眼中，这是最极端的危险。\par

100. 我常说，让别人为贸易而航行，横渡广阔的海洋，不断与风浪搏斗，以博取巨大财富——若命运如此——并在航海的渴望中冒大风险；而我，则远更喜欢留在岸上，犁一短而愉快的犁沟，远远地向大海及其收益致敬，靠贫乏而微薄的大麦饼生活，在安全与平静中拖曳我的生命，而不是为了巨大的收益而将自己暴露于如此长久与巨大的风险中。\par

101. 因为居高位者，若未能进一步前进和更广泛地传播德行，却满足于微小的成果，便招致惩罚，如同以巨大的灯盏照亮小屋，或在孩童的四肢上围上成人的全副铠甲。相反，出身低微者可以安全地承担轻担，并避免因尝试超出能力之事而招致嘲笑和更大风险的风险。因为我们听说过，一个人建造塔楼，除非他有足够的钱完成，否则是不相称的。\BibleRef{路 14:28}\par

102. 这就是我能够为我的逃跑所作的辩护——或许过长了些。这就是那些原因——令我痛苦，也可能令你们痛苦——将我携离你们，我的朋友和弟兄；然而，正如我那时看来，是不可抗拒的力量。我对你们的渴望，以及你们对我渴望的感觉，比任何事物都更促使我归回，因为没有什么比相互的爱更能使我们倾向爱。\par

103. 其次，还有我的责任、我的义务、我圣洁双亲的白发与衰弱——他们因我而受到的痛苦比因年迈更甚——这是亚巴郎族长（我尊崇其人，并被列于天使之中）和撒辣（她借着真理的教导，在属灵上生育了我）的事。如今，我已特别承诺要成为他们老年的支柱和衰弱的支持——这一承诺，我按自己的能力已经履行，甚至以牺牲哲学本身（对我而言最珍贵的财产和称号）为代价；或者更确切地说，虽然我哲学的首要目标似乎是显得没有哲学，但我不能忍受因单一目标而浪费我的劳力，也不能失去那份祝福——据记载，古代的一位圣人曾以欺诈的方式从他父亲那里偷得了祝福，他以献上的食物和伪装的毛皮欺骗了父亲，从而以可耻的诡计达到了善良的目的。\BibleRef{创 27:21}这就是我顺服和听话的两个原因。或许我的理由向两者屈服和顺服并非不合理，因为有时得胜是适时的，正如我也认为凡事都有定时，\BibleRef{训 3:1}光荣地被征服，胜过赢得危险而不法的胜利。\par

104. 我还要提及第三个极为重要的原因，然后便搁置其余。我追念往日，回溯一段古老历史，从中为自己的当前行为寻得劝勉；因为我们不可认为这些事件是无目的地被记载，也不可认为它们只是为取悦听众而搜集的言辞与事迹的集合，如同一种耳目的饵食，仅以快感为目的。让我们把这种戏谑留给传说和那些不大重视真理、以虚构之魅力和文辞之雅致来迷惑耳目的希腊人吧。\par

105. 然而我们，既将圣神的精确性延伸到最小的笔划和字头（\BibleRef{玛 5:18}），就绝不会接受那种不虔敬的主张，即认为最微小的事件也是由记载它们并因此流传至今的人随意处理的；相反，他们的目的是为我们在类似情况下（若我们遭遇）提供纪念和训诲，使过去的榜样能成为规则和模型，供我们警诫和效仿。\par

106. 那么，这故事是什么？它的应用何在？因为，叙述它或许对公众安全不无裨益。约纳也曾逃离天主的面容（\BibleRef{纳 1:3}），或者更确切地说，他自以为在逃离；但他被大海、风暴、签、鲸腹和三天埋葬所追上——这是更大奥秘的预象。他因不愿向尼尼微人宣布那可怕而严厉的消息，并因城市若因悔改而得救，自己随后会被证明是假先知而逃离；并非他不喜罪人得救，而是他耻于成为谎言的工具，且极其热忱于预言的声誉——这声誉有可能在他身上被摧毁，因为多数人无法洞察天主在这种情形下的深奥安排。\par

107. 然而，正如我从一位精通这些主题、能把握先知深意的人所学到的，通过对历史上似乎不合理之处给予合理解释，并非此原因导致约纳逃走，并带他到约培，又从约培往塔尔史士，当他将偷来的自己托付给大海时（\BibleRef{纳 1:3}）；因为这样一位先知不可能不知道天主的计划——即通过威胁，按照祂伟大的智慧、不可探究的判断和无法追踪的道路（\BibleRef{罗 11:33}），使尼尼微人逃脱所威胁的厄运；也不可设想，若他知道这一点，却拒绝与天主合作使用祂为救他们而设计的方法。此外，想象约纳希望在海中藏匿，并藉逃跑躲避天主的大眼，这绝对是荒谬愚蠢、不足置信的——不仅对一个先知而言，甚至对任何略有天主知觉、意识到祂全能之力的明智之人也是如此。\par

108. 相反，如同我的导师所说，也如我自己所深信，约纳比任何人都更清楚他向尼尼微人所传使命的目的，并且他在策划逃跑时，虽改变地点，却并未逃脱天主。这对任何人来说都是不可能的，无论他藏身于地心、海底，还是乘翅高飞（若有此法），或升入空中，或居于地狱最深处，或把自己裹入密云，或使用任何其他逃脱的计谋。因为唯独天主是无法逃脱或与之抗争的；若祂愿意擒拿并使人服从祂的手，祂必超越迅捷者、智取智者、推翻强者、贬抑高傲者、制服鲁莽者、抑制权力者。\par

109. 因此，约纳并非不知天主的大能之手——他曾用这手威胁他人；他也不认为能完全逃脱天主的权能——这是我们不能相信的。当他看见以色列的堕落，察觉预言恩宠转向外邦人时——这正是他隐退不传道、延迟执行命令的原因；于是他离开了欢乐的瞭望塔（这是约培在希伯来语中的含义），即他先前的尊严和声誉，投身于悲伤的深渊。因此他遭受风暴，陷入沉睡，沉船，被唤醒，被签所定，承认逃跑，被抛入海，被鲸吞下却未被毁灭；在那里他呼求天主，奇妙的是，第三天他如同基督被释放。但对此主题的处理必须暂缓，若天主允许，不久将更审慎地阐述。\par

110. 如今，回到我最初的观点，我心中浮现这样的思考和疑问：虽然约纳若因我提到的原因不情愿预言，或许能获得某种宽容——但在我自己的情况下，若我继续倔强、拒绝那已置于我肩上的服务之轭（我不知该称它为轻还是重），我能说什么？我能作何辩护？\par

111. 因为，就算我们承认——且此事中唯一可断言的正是这一点——我们远不能在天主面前履行司铎之职，只有在我们配得上教会之后才能配得上圣所，配得上主席之位之后才能配得上圣所，但也许其他人不会因不顺从的罪名而宣告我们无罪。如今，不顺从的恐吓是可怕的，随之而来的惩罚也是可怕的；同样，另一边也有可怕的后果，如果我们不是像撒乌耳那样躲在他父亲的行李中（\BibleRef{撒上 10:22}）——虽然他被召为君王只是短暂一时——而是轻易地走上前来，仿佛这是一件轻微又最容易的任务，而事实上连辞去它都不安全，也无法以第二次考虑来修正第一次。\par

112. 为此，我费尽心思要发现自己的责任，置身于两种恐惧之间：一种阻止我，另一种催逼我。我长时间在它们之间不知所措，左右摇摆，如同被不定的风吹动的水流，时而倾向这边，时而倾向那边，最终我向较强的一方屈服，对不顺服的恐惧胜过了我，将我带走。请留意，我是何等精确公正地衡量这两种恐惧：既不贪求未经授予我的职务，也不拒绝已授予我的职务。前者表明卤莽，后者表明不顺服，两者皆表明缺乏纪律。我的立场介于过于大胆与过于胆怯之间：比那些冲向每一种职务的人更胆怯，比那些回避所有职务的人更大胆。这就是我对这件事的判断。\par

113. 此外，为了更清楚地区分两者，我们对于职务的恐惧，可以在服从的法律中找到可能的帮助，因为天主按其良善赏报我们的信德，使信靠祂并将一切希望寄托于祂的人成为完美的统治者；但面对不顺服的危险，我不知道有什么可以帮助我们，也没有任何理由鼓励我们的信心。因为恐怕我们将要听到关于那些托付给我们的人的话：\Quote{我要从你手中追讨他们的灵魂；\BibleRef{则 3:18}并且，因为你弃绝了我，没有作我子民的首领和统治者，我也要弃绝你，使我不再作你的君王；并且，正如你们拒绝听从我的声音，转过倔强的背，不顺服，同样，当你们呼求我时，我必不顾念，也不垂听你们的祈祷。}但愿这些话语不会从公义的审判者临到我们，因为当我们歌唱祂的仁慈时，我们也必须歌唱祂的审判。\par

114. 我再次诉诸历史，思考古时最负盛名的人，他们曾因恩宠被优先选为统治者或先知，我发现有些人欣然回应召唤，另一些人却拒绝这份恩赐，并且那些退缩的并未因胆怯而受责备，那些挺身而出的也并未因热切而受责备。前者敬畏职务的伟大，后者信靠地服从了召唤他们的那一位。亚郎是热切的，但梅瑟抗拒；依撒意亚欣然顺服，但耶肋米亚害怕他的年轻，\EditorialNote{依撒意亚 6:8}直到他从天主那里得到超越他年龄的应许和能力才敢说预言。\BibleRef{耶 1:6}\par

115. 借着这些论点，我陶醉了自己，我的灵魂逐渐松弛，变得如铁般可塑，时间也助长我的论点，而天主的证言——我把自己的一生都托付给了它——成了我的顾问。因此，我没有悖逆，也没有转身退后，\EditorialNote{依撒意亚 50:6}我的主说，尽管祂没有被称为统治，却被牵到宰杀之地，如同羊羔；但我俯伏在天主大能的手下，\BibleRef{伯前 5:6}并为我以前的懒惰和不顺服祈求宽恕，如果这被归咎于我的话。我保持沉默，\EditorialNote{依撒意亚 42:14}但我不会永远沉默：我退后片刻，直到我审视自己、安慰我的忧伤；但现在我受命在会众中高举祂，在长者的座位上赞美祂。如果我以前的行为该受责备，我现在的行为则值得宽恕。\par

116. 何须更多言语？我在这里，我各位牧者及同牧者，我在这里，你们神圣的羊群，堪当基督——总司牧，\BibleRef{伯前 5:4}我在这里，我的父亲，完全被征服，按基督的法律而非按地上的法律服从你们：这是我的服从，以你们的祝福报偿它。以你们的祈祷引领我，以你们的言语指导我，以你们的精神坚定我。父亲的祝福坚立子女的房屋，\BibleRef{德 3:9}愿我和这属神的房屋都被坚立，我所渴慕的房屋，我祈祷它成为我永远的安息，当我从此处的教会转至彼处的教会，就是那些名字写在诸天之上的首生者的总集会。\BibleRef{希 12:23}\par

117. 这就是我的辩护：我已陈明其合理性。愿平安的天主，\BibleRef{希 13:20}祂使两者合而为一，\BibleRef{弗 2:14}并使我们彼此和好；祂使君王登上宝座，从灰尘中提拔穷乏人，从粪堆中举起乞丐；祂拣选了祂的仆人达味，将他从羊圈中带来，尽管他是叶瑟儿子中最小的和最年轻的；\BibleRef{撒上 17:14}祂将言语赐给那些以大能传福音、为福音的成全之人；——愿祂亲自用右手握住我，以祂的谋略引导我，以光荣收纳我；祂是牧者的牧者，\BibleRef{则 34:12}引导者的引导：使我们能以知识牧养祂的羊群，\BibleRef{耶 3:15}而不是用愚昧牧人的器械，\BibleRef{匝 11:15}按照祝福，而非按照从前对古时之人所宣告的诅咒；愿祂将力量和能力赐给祂的子民，并亲自将祂的羊群呈献给自己，\BibleRef{弗 5:27}羊群光彩夺目、毫无玷污，堪当天上的羊圈，在喜乐之人的居所，在圣徒的光辉中；以致在祂的圣殿中，每一个人，包括羊群和牧者都同声说：光荣，在基督耶稣我们的主内，愿光荣归于祂，直到永远。阿们。\par

\SourceRecord{Translated by Charles Gordon Browne and James Edward Swallow.；Nicene and Post-Nicene Fathers, Second Series；Vol. 7.；Edited by Philip Schaff and Henry Wace.；Buffalo, NY: Christian Literature Publishing Co.,；1894.；<http://www.newadvent.org/fathers/310202.htm>.}

% structure: p0001 source=h1 target=section reason=page-title

\section{第三篇演讲}

\Emph{致那些邀请了我却没有前来迎接我的人}\par

\Emph{（约公元362年复活节）}\par

一. 我的朋友们和弟兄们，你们来听我说话是多么迟缓呀，尽管你们在强逼我、把我从我的堡垒——隐修生活中拉出来时是那么迅速；这隐修生活是我所拥抱的，远超其他一切，视之为神圣攀登的助手和母亲，使人神化，我特别钦佩它，并把它当作我整个生活的指引。如今你们得到了它，怎么又如此鄙视你们那么渴望得到的东西呢？你们似乎更能渴望不在眼前的事物，却不会享受当前的事物；好像你们更愿意\Emph{拥有}我的教导，而不愿\Emph{从中获益}。是的，我甚至可以对你们这样说：\Quote{我在你们尝到我、或试验我以前，就使你们厌烦}\BibleRef{依 1:14}——这真是再奇怪不过了。\par

二. 你们既没有款待我作客，也不容我（如果我可以作一个更富有怜悯心的评论）款待你们，即使没有其他理由，也要尊重这条命令；你们没有在我开始一项新任务时扶持我，也没有鼓励我的胆怯，也没有安慰我所遭受的暴力；反而——我不愿说，却又不得不说——你们使我的节日不成节日，毫无喜庆地接待我；你们用悲伤混淆了庄严的节日，因为它缺少最能促使它快乐的因素，就是你们这些征服我的人在场，因为称你们为爱我的人是不恰当。凡是容易征服的东西，多么容易被人轻看；骄傲的人受到关注，而在天主面前谦卑的人却被忽略。\par

三. 你们要怎样？是被你们审判，还是我来审判你们？是宣判，还是接受判决？因为我若被审判，就希望被宣告无罪；我若宣判，就希望公正地判你们有罪？对你们的控诉是：你们没有用同等的爱回报我的爱，也没有用尊荣回报我的服从，也没有用你们现在的热忱向我保证未来——即使你们这样做了，我也几乎难以相信。但你们每个人都有某些事物比新旧牧者都更偏爱，既不尊重前者的白发，也不激发后者的青春精神。\par

四. 福音中有一个婚宴 \BibleRef{路 14:16}，和好客的主人及朋友们；这婚宴极其欢乐，因为是他儿子的婚礼。他召唤他们，他们却不来：他发怒了，而且——我因怕不祥之兆而跳过中间之事——但温和地说，他用别人填满了婚宴。愿天主不让这事临到你们；但是你们对待我（我该怎样温和地表达呢？）却像那些被邀请赴宴却起来反对并侮辱主人的人一样傲慢或鲁莽；因为你们，并不在圈外之人或受邀参加婚礼之列，而是你们自己邀请了我，把我绑到了祭台前，向我展示了洞房的荣耀，然后却抛弃了我（这是你们最光彩的一点）——一个去了自己的田地，另一个去了新买的牛轭，另一个去了刚娶的妻子，另一个去了其他琐事；你们都四散分开，毫不关心洞房和新郎。\BibleRef{玛 22:10}\par

五. 因此我充满沮丧和困惑——因为我不愿对自己所受的痛苦保持沉默——我差点就要扣下原本打算当作结婚礼物赐予的讲道，那是我所有礼物中最美好最珍贵的；我几乎要把它向你们释放，你们这些在我遭受暴力之后我非常渴望见到的人：因为我想从中得到一个绝佳的主题，而且我的爱磨利了我的舌头——这爱极其炽热，当它因某种意外忽视而产生悲伤而激起嫉妒时，就准备好指责。如果你们中有人曾被爱情刺伤，感到自己被忽视，他明白这种感受，也会原谅这样受苦的人，因为他自己也曾接近这种疯狂。\par

六. 但此刻我不可对你们说任何责备的话；愿天主永不让我这样做。甚至现在，我或许已经过分地斥责了你们——神圣的羊群，基督可赞的乳儿，天主的产业；天主，因着这些，祢是富有的，即使祢在其他方面贫乏。我以为，这些话正适合祢：\Quote{祢的产业落在佳美之处：是的，祢获得了最美好的基业。}我也不允许最人口众多的城市或最广阔的羊群比我们更有优势，我们是以色列众支派中最小的一个的小子们，是犹大千万中最少的 \BibleRef{撒上 23:23}，是从小城白冷中来的 \BibleRef{米 5:2}，基督在那里出生，从起初就众所周知并被朝拜；在那些人中，圣父被高举，圣子被视为与祂同等，圣神与他们一同受光荣：我们同心一体，思念同一件事，绝不损害天主圣三，既不偏爱一位过于另一位，也不排除任何一位；不像那些不良的裁判和丈量神性的人，他们过分抬高一位，就贬损并侮辱了全体。\par

七. 但你们，如果你们对我怀有善意——你们是我的田地，我的葡萄园，我自己的心肠，或者更确切地说，是属于那一位的，祂是我们共同的父，因为在基督内祂藉着福音生了你们——也要向我们表示一些尊重。这是公平的，因为我们尊崇你们超过一切：你们是我的见证人，你们和那些将我们手中之物交付给我们的人——我该说是\Emph{权柄}还是\Emph{服务}？如果最爱的人应得，我该如何衡量这份爱，我以自己的爱使你们成为欠我债的人？但不如尊重你们自己，尊重托付给你们的肖像（见：\BibleRef{创 1:27}），尊重托付这肖像的那一位，尊重基督的苦难以及你们由此而生的希望；持守你们所领受、在其中受培养、并藉以得救的信德，并相信能拯救别人（因为你们确知，没有多少人能夸耀你们所拥有的）；虔诚不在于多谈论天主，而在于大多时候保持沉默，因为舌头若不节制，对人是危险之物。要相信，听总是比说更安全，正如学习天主比教导人更令人愉快。将这些问题的更精细探究留给那些圣言的管家；至于你们自己，用言语少量敬拜，但更多以行动，并藉遵守法律而非仰慕立法者来表示爱祂；藉逃避邪恶、追求美德、活在圣神内、随从圣神行事、从祂汲取知识、在信德的基础上建造，不是用易朽的木、草、禾秸（见：\BibleRef{格前 3:12}），因为当火考验或毁灭我们的工程时，这些材料软弱且易耗；而是用金、银、宝石，这些存留且屹立。\par

八. 愿你们如此行事，并如此尊敬我们，无论我们在场或不在场，无论你们参与我们的讲道，还是宁愿做别的事。愿你们成为天主的子女，纯洁无瑕，在弯曲悖谬的世代中（见：\BibleRef{斐 2:15}）；愿你们永不陷入环绕四周的恶者的圈套，或被自己罪的锁链捆绑。愿在你们内的圣言永不因今生的忧虑而窒息，以致你们不结果实；但愿你们行走在君王的道路上，不偏左也不偏右，而是被圣神引导穿过窄门。这样，我们的一切事务都将顺利，无论是现在还是在彼处的审讯，因我们的主耶稣基督，愿光荣归于祂，直到永远。阿们。\par

\SourceRecord{Translated by Charles Gordon Browne and James Edward Swallow.；Nicene and Post-Nicene Fathers, Second Series；Vol. 7.；Edited by Philip Schaff and Henry Wace.；Buffalo, NY: Christian Literature Publishing Co.,；1894.；<http://www.newadvent.org/fathers/310203.htm>.}

% structure: p0001 source=h1 target=section reason=page-title

\section{第七篇演说}

\Emph{论其弟圣凯撒留的颂词。}\par

\EditorialNote{这篇演说的日期大概是公元369年春天。圣热罗尼莫将它排在圣额我略演说集的首位。凯撒留是圣人的弟弟，大概生于公元330年左右。他在家接受早期教育，后来在亚历山大里亚的各学校学习，在数学、天文学，尤其是医学方面取得了极大的成就。从亚历山大里亚回来后，应公众请求，君士坦西乌皇帝向他提供了一个在拜占庭的尊贵且高薪的职位，但额我略说服了他一同返回纳齐盎。过了一段时间，他回到拜占庭，在尤利安登基后，他被迫保留宫廷职务，尽管额我略责备他，他还是这样做，直到尤利安长期试图使他放弃基督教后，最终邀请他进行公开辩论。凯撒留尽管面对皇帝似是而非的论证，但取得了胜利，但既然他已明确自称为基督徒，便不能再留在宫廷。尤利安死后，他受到继任皇帝的尊重和提拔，直到他从瓦伦斯手中接受了比提尼亚财务官的职务。这个职务的确切性质和等级至今仍未被历史学者确定，其中一些人将圣额我略未曾提及的其他职务归于他，这些职务很可能是由同名者担任的。公元368年10月11日，尼西亚城几乎被地震完全摧毁，凯撒留奇迹般地幸免于难。这次脱险令他深受触动，他领受了圣洗圣事，并计划辞职退休，似乎要过一种苦修生活，但这些计划被他突然的早逝所打断。}\par

1. 我的朋友们、弟兄们、父老们（你们无论在实际上还是在名义上都是我所亲爱的）或许认为，我即将为离去的人献上悲伤的悼词，是急于承担这项任务，并且会像大多数人喜欢做的那样，长篇大论，辞藻华丽。因此，你们中一些曾承受过同样悲伤的人，准备加入我的哀悼和悲叹，以便在我的悲痛中哀悼你们自己的忧伤，并学会因朋友的苦难而感到痛苦；而另一些人则期待以聆听我言辞的方式满足耳朵的享受。因为他们认为，我必定会像从前那样，当我拥有大量世俗之物，且最渴望获得演说声望时——在我仰望那真实而最高的\OriginalTerm{圣言}{Word}，将一切交于天主（万物由祂而来），并以天主为万有之先——把我的不幸当作炫耀的机会。现在，请你们不要这样想我，如果你们愿意正确看待我的话。因为我既不想为他惋惜超过限度（正如我也不应赞同别人这样做），也不想过分赞美他。尽管语言对于一位有语言天赋的人是一种珍贵而特别的致敬，赞美对于一位非常喜爱我言辞的人更是一种敬意——是的，不仅是致敬，更是一种债务，是所有债务中最正当的。但即使在眼泪和赞美中，我也必须尊重关于这些事的法律；这并不违背我们的哲学，因为他说：\Quote{义人的纪念伴随赞美，}又说：\Quote{让眼泪落在死者身上，开始哀悼，好像自己受了重伤一样：\BibleRef{德 38:16}}使我们远离麻木和无节制。因此，我将不仅彰显人性的软弱，还要提醒你们灵魂的尊严，并给予忧伤者应有的安慰，将我们的悲伤从关乎肉体和暂时的事转移到属灵和永恒的事上。\par

2. 凯撒留的父母，首先提及我最应提及的一点，是你们都知道的。你们热切地注意他们的卓越，并怀着钦佩之心听闻，而且乐意向任何不知晓的人（如果有这样的人）讲述；因为没有一个人能完全做到，这超出了任何勤奋和热忱的单一舌头的能力。在他们值得歌颂的众多伟大优点中（希望我不会因赞美自己家族而显得过分），最大的一点，也是最突出其品格的，是虔诚。他们因白发而值得尊敬，但他们的德性并不亚于年龄；虽然他们的身体因年岁而弯曲，但他们的灵魂在天主内焕发青春。\par

3. 他的父亲从野橄榄枝被巧妙地嫁接在好橄榄树上，并如此分享其肥汁，以致被托付去嫁接他人，并负责牧养灵魂，以符合其高位的方式治理这群百姓，如同第二个亚伦或梅瑟，被吩咐亲自走近天主（\BibleRef{出 24:1}），并向站在远处的其他人传达天主的声音；温和、柔顺、面容平静、心灵热切，外表英俊，但内在更为富足。但我为何要描述你们已经认识的人呢？因为即使我说很多，也不能达到他应得的赞美，或你们每个人所知道并期望听到的。因此，更好的办法是让你们自己去想象他，而不是用我的言语割裂你们所景仰的对象。\par

4. 他的母亲因出身圣洁家族而奉献于天主，并拥有虔诚作为必要的遗产，不仅为她自己，也为她的孩子们——她确实是圣洁的初果所成的一团圣洁面团（\BibleRef{罗 11:16}）。她如此增长和扩大这一遗产，以致有些人（尽管这话大胆，我仍要说出来）已经相信并说过，她丈夫的成全唯由她而来；何等奇妙！她自己作为虔诚的报酬，获得了更大更完全的虔诚。他们两人都爱他们的子女并爱基督，但最特别的是，他们爱基督远胜于爱他们的子女；是的，他们甚至对子女的唯一享受就是子女被基督承认并称呼，他们衡量子女幸福的唯一标准是德性和与\OriginalTerm{至善}{the Chief Good}的密切结合。富有同情心、怜悯心，从蛀虫、强盗和这\OriginalTerm{今世之王}{prince of this world}（\BibleRef{若 14:30}）手中夺取许多财宝，将其从此世的旅居转移到〔真正的〕住所，为他们的子女积攒（\BibleRef{弟前 6:19}）天上的光辉作为最大的遗产。由此，他们达到了美好的老年，因德性和岁月同样可敬，充满日子——既有永存的，也有消逝的；每人在地上未能获得首奖仅仅是因为被对方所超越；是的，他们已满足了每一福乐的尺度，除了这最后的考验或锻炼（无论我们应将其称为何）——我是指他们不得不将那因年龄更易跌倒的孩子先送走，从而在安全中结束生命，并带着全家被提升到上面的领域。\par

5. 我详述这些细节，并非出于歌颂他们的愿望——因我深知，即使我将整篇讲辞用于赞美他们，也难以相称地做到——而是为了陈明\Person{凯撒瑞}从他父母继承的美德，使你们不致惊讶或怀疑：出身于如此先辈之人，竟当得起这样的赞誉；反之，倘若他视他人为典范，却忽视自己亲属家中榜样，那才真是奇怪。他的早年生活恰如真正出身高贵、命定善生者之所为。他人显而易见的品质——他的美貌、他的身材、他多样的优雅，以及嗓音中流露的和谐气质——我不多言；因为赞颂这类品质并非我的职责，无论它们在他人眼中何其重要。我要继续言说那些即使我想避开也难以略过的要点。\par

6. 在此等影响下成长并受养育，我们在此地所提供的教育中得到了充分训练。论及他凭才智的敏捷与广博远超我们大多数人的程度，无人能尽述。我如何能不流泪地追忆那些日子？这流泪表明，与我承诺相反，情感胜过了我哲学的克制。离家的时刻到来，我们首次分离：我在当时盛行的\Place{巴勒斯坦}学校研习修辞；他则去了\Place{亚历山大里亚}，那在古今均被视为一切学问之故乡。他的何种品质我应置于首位最优先？何种我可以略去而不损及我的描述？谁比他更忠于师长？谁比他更友善待同窗？谁比他更谨慎远离堕落者的交往和陪伴？谁比他更亲近最卓越之人，尤其是他那最受敬重、最光耀的同乡？因他知道，我们受同伴影响，或行善或为恶，影响甚深。正因这一切，谁比他更获当局尊敬？整个城市（虽然因庞大而个体隐匿其间）有谁因他的审慎而更受推崇，因他的才智而更被视为杰出？\par

7. 什么学问他未曾精通？甚或，在何种研究领域，他未超越那些以此为一学习之人？他容许何人甚至接近他，不仅在同龄人，即便在投入更多岁月研习的长者中？他研究一切学问如研究一门，而又如同不知其他，将每门学问研习得同样透彻。才智出众者，他于勤奋上超越；勤勉者，他于领悟迅捷上超越；不，他甚至以迅捷超越敏捷者，以勤奋超越刻苦者，并在两方面超越两者俱优之人。从几何与天文学——这门对他人的危险学问——他汲取了一切有益者（我意指他藉天体之和谐与秩序，引申而敬畏其造物主），并避开了有害者；他不像那些将受造物——这同侪——高举以反抗造物主之人，将一切存在或发生之事归于星辰影响，而是合理地将其运动及一切其他之物归于天主。在算术与数学，以及在奇妙的医学艺术——就其处理生理学、气质与病因，以拔除病根并随之摧毁其子嗣而言——有谁如此无知或好争，竟以为自己逊于他，不以被列于其下、得二等奖为乐？这并非无根据之论；东方与西方，以及他后来造访的每一处，都如同铭刻他学问的柱石。\par

8. 但当他将每一种卓越和知识如同大商人搜集各类货物般聚集于单一灵魂，驶向本城，欲与他人分享其学识之美货时，一件奇妙之事发生了。此事之述说对我最为振奋，亦或使你们喜悦，我须简要道来。我们的母亲，以对儿女的慈爱，曾献上祈祷：既然她一同送我们出门，也能一同见我们回家。因为我们——至少在母亲眼中，若不在他人看来——形成一对在同时出现时堪受她祈祷与目光的伴侣；但如今，唉，我们的联系已被切断。天主垂听义人的祈祷，并尊重父母对好子女的慈爱，如此安排：我们一方自\Place{亚历山大里亚}、一方自\Place{希腊}，一方经海路、一方经陆路，毫无计划或约定，同时到达同一城市。这城便是\Place{拜占庭}，现今统御\Place{欧罗巴}。在此，\Person{凯撒瑞}在短短时日后获得如此声誉，以致公共荣誉、与显赫家族联姻、以及国务会议席位均被提供给他；且由公共决议派遣使团至皇帝，恳求让这诸城之首城由学者之首来妆点和尊荣（若他关心它如何才是真正的首城、名实相符的话）；并使其所有其他卓越头衔之上再加一项——以\Person{凯撒瑞}为其医士与居民而增辉，虽然其辉煌早已因众多哲学及其他学问伟人之群集而有保障。但言尽于此。此时发生之事，在他人看来像是无因无由的偶然，如同今日屡屡自然发生之事，但对虔敬心灵而言，这已足够显明为敬畏天主的父母祈祷应验之结果——其子从海陆同时抵达。\par

9. 在\Person{凯撒里乌斯}的高贵品格中，有一项我们不可不提；在他人眼中，它或许微不足道，不值一提，但就我看来，当时以及后来，它都具有至高的重要性——如果手足之爱确实是值得称赞的美德。我在叙述他生平时，将永不停止地把它置于首位。尽管都城竭力以我所提及的荣誉来挽留他，并宣布绝不放他走，但我的影响力——他在一切场合都极为珍视——说服了他听从父母的祈祷，满足祖国的需要，并成全我自己的心愿。当他这样随我一同回家时，他不仅将我置于城市和民众之上，不仅置于荣誉和收入之上——这些荣誉和收入部分已从多方大量涌向他，部分则触手可及——甚至置于皇帝本人及其皇命之上。从那时起，我摆脱了一切野心，视之如严酷的主人、痛苦的紊乱，决心修德，使自己适应那更高的生活；更确切地说，这愿望早已诞生，而生活则随后而来。然而，我的兄弟，他将自己学问的初果献给了祖国，赢得了与他的努力相称的钦佩，随后却被名声的欲望所驱使——如他所说，也是出于成为城市守护者的愿望——以致投身朝廷，这确实不符合我的愿望或判断；因为我愿向你们承认，我认为与天主共处最卑微的等级，比赢得尘世君王的首位更为美好和伟大。尽管如此，我不能责备他，因为修德既是最高尚的职业，也是最困难的职业，只有少数人能从事，只有那些蒙天主宽宏大度所呼召的人——祂向蒙祂拣选而受尊荣的人伸出援手。然而，若有人选择了较低形式的生活，却追随良善，将天主和自身的救恩看得比尘世的光彩更重要，那也绝不是小事；他视尘世光彩如同舞台，或是在上演这世界戏剧时的多种短暂面具，而自己却按那肖像——他知道自己从天主领受了这肖像，且必须归还给赐予者——为天主而生活。我们清楚地知道，这确实是\Person{凯撒里乌斯}的意愿。\par

10. 在医生中，他毫不费力地获得了最前列的地位，仅凭展示他的能力——更确切地说，是他能力的些许样本，便立即被列入皇帝的朋友之列，并享有了最高的荣誉。但他将医术中的人道职能无偿地交给当局支配，深知没有什么比美德和善行的名誉更能带来进一步的升迁；因此，他在名声上远超过那些地位比他高的人。他的谦逊如此赢得了所有人的爱戴，以至于他们将自己宝贵的病人托付给他照看，而不必让他以\Person{希波克拉底}之名起誓，因为\Person{克拉特斯}的质朴与他相比也相形见绌：总体上，他赢得了超越其地位的尊敬；因为除了他当时被认为公正赢得的声誉之外，皇帝本人以及所有身居他们最近职位的人还预期他会有更大的声誉。但最重要的是，无论是名声还是周围的奢华，都没有腐蚀他高贵的灵魂；因为在他众多的荣誉中，他自己最关心的是成为一个基督徒，并被人知道是基督徒；与此相比，其他一切在他看来都不过是儿戏。因为那些东西属于我们在舞台上向他人扮演的角色，这舞台很快搭建又很快拆除——或许破坏比搭建更快，正如我们从人生的多变和命运的起伏中所见；而唯一真实且稳固持久的善是虔诚。\par

11. 这就是\Person{凯撒里乌斯}的修德之道，甚至身在朝廷也是如此：他所生活和死亡于其中的就是这些观念，在隐秘的人中，他向天主发现并呈献了一种比公开可见的更深的虔诚。如果我必须略过其他一切——他对困苦中亲人的保护，对傲慢的蔑视，对朋友的不自负，对有权势者的勇敢，他为真理与许多人进行的无数争论和辩论，不仅为了争辩，而是出于深厚的虔诚和热忱——我至少必须提及一点，特别值得注意。那位记忆不祥的皇帝正对我们发怒，他拒绝基督的疯狂，在使自己成为第一个牺牲品后，如今已令他无法容忍他人；然而他并不像其他反对基督的人那样，堂皇地自列于不虔之列，而是将迫害掩盖在公正的外表之下；并且，被占据他灵魂的弯曲之蛇所支配，以各种诡计将可怜的人拖入自己的坑中。他的第一个诡计和手段是，剥夺我们斗争的荣誉（因为，他这位高贵的人，却对基督徒吝啬这一点），使我们这些因身为基督徒而受苦的人被当作恶人受罚：第二个是，称这个过程为说服，而非暴政，以便那些选择站在不虔一方的人的耻辱比他们的危险更大。他有的用金钱收买，有的用职位，有的用承诺，有的用各种荣誉——他并非以君王的方式，而是以十足的奴仆姿态，当众赏赐这些，而每个人都受他言语的迷惑和自己榜样的影响。最后，他攻击了\Person{凯撒里乌斯}。何等的疯狂与愚蠢，竟妄想将像\Person{凯撒里乌斯}这样的人——我的兄弟，我们这样父母的儿子——作为猎物！\par

12. 但让我稍稍停驻于此，沉浸于我的故事中，如同那些亲历奇妙事件的人们一样。那位高尚的人，以基督的标记坚固自己，并以祂的圣言保卫自己，进入战场，对抗一位在兵器上和论辩技巧上都老练的对手。他丝毫未因景象而羞愧，也未因谄媚而退缩于崇高目标，他是一位斗士，在言与行上都准备好与旗鼓相当的对手较量。那么，这就是竞技场，而这位斗士就是虔诚的捍卫者。裁判是基督，以祂自己的苦难来武装斗士；另一边是可怕的暴君，以其论辩技巧说服人，并以权威重压人。作为旁观者，两边既有仍站在虔诚一方的人，也有被暴君夺去的人，他们注视着胜利倒向哪方，比斗士自己更焦虑于胜负。\par

13. 你们岂不担心凯撒留，怕他遭遇任何有损于他热忱的事吗？不，你们要鼓起勇气。因为胜利属于基督，祂战胜了世界。\BibleRef{若 16:33} 就我而言，请你们确信，我很愿意详细叙述那场合中使用的论点和指控的细节，的确，那讨论包含某些精彩和优雅之处，我怀着不小的乐趣沉浸其中；但这对像现在这样的场合和讲词来说很不合适。当凯撒留将他对手的一切诡辩撕成碎片，并将每一次或明或暗的攻击都当作儿戏般推开后，他以清晰响亮的声音宣称他是、并且仍然是一位基督徒——即使如此，他最终也未被释放。因为皇帝怀有强烈的渴望，想要享受他的文化并借此出名，然后在众人面前说出了他那著名的话——\Quote{哦，幸福的父亲！哦，不幸的儿子！}他这样屈尊来荣耀我——他曾因我在雅典的文化和虔诚而认识我——却分享凯撒留的屈辱，后者被押回再受审判（因为正义正适时地武装皇帝去对抗波斯人），而我们在凯撒留幸运逃脱和不见血的胜利之后迎接他，他因屈辱比因名声更为显赫。\par

14. 我认为这场胜利远比皇帝强大的权柄、辉煌的紫袍和昂贵的冠冕更为崇高和荣耀。我在描述它时比他从皇帝那里赢取一半帝国更为兴奋。在邪恶的日子里，他过着隐居生活，听从我们基督徒的律法，\BibleRef{玛 10:23} 这律法吩咐我们，有机会时要为真理冒险，不可因胆怯而背叛我们的信仰；但也要尽可能避免冲入危险，或出于对自己灵魂的恐惧，或出于对带来危险之人的顾惜。但当阴霾散去，公义的判决在异地被宣布，闪亮的刀剑击倒了不虔敬者，权柄回到基督徒手中时，无需多言他是带着何等荣耀和尊荣，以及多少伟大的见证——仿佛是在赐予而非接受恩惠——重新受到宫廷欢迎；他新的尊荣接续了旧日的荣光；时间更替了皇帝，但凯撒留的声誉和影响力与他们相处却毫不受损，甚至他们相互竞争，力求将他最紧密地拉拢到自己一边，并被人认为是他的特殊朋友和相识。这就是凯撒留的虔诚及其果实。愿所有人，无论长幼，都倾听并沿着同样的德行奔向同样的荣耀，因为光荣是善工的果实，\BibleRef{智 3:15} 倘若他们认为这是值得追求的，并且是真福的一部分。\par

15. 关于他的另一奇事，是他父母虔诚和他本人虔诚的强有力证据。他当时住在\Place{比提尼亚}，担任皇帝一项不轻的职务，即掌管他的税收并管理国库：因为皇帝已将之作为最高官职的预兆指派给他。不久前的\Place{尼凯亚}地震（据说是人们记忆中最严重的一次，几乎将全城居民和城市之美一同摧毁）中，他独自一人，或与极少数显贵一起，逃过了危险，被倒塌的废墟本身所庇护，以难以置信的方式幸免于难，只留下轻微的危险痕迹；然而他让恐惧引领他走向更重要的救恩，因为他将自己完全交托于至高天主眷顾；他放弃了转瞬即逝之物的服务，投身于另一个宫廷。他不仅自己有意如此，还以此为我与他共同热切祈祷的对象，通过书信邀请我，我当时抓住这个机会向他提出警告——因为我从未停止过这样做，为他的伟大天性被用在低于它的事务上而感到痛苦，也为一个同样适合哲学的灵魂如同太阳被云遮蔽一样，在公共生活的漩涡中黯然失色而感到痛苦。虽然他在地震中安然无恙，但他未能抵御疾病，因为他只是凡人。他的幸免于难是他独有的；他的死亡是全人类共通的：一个是其虔诚的标记，另一个是其本性的结果。前者为安慰我们而在他的命运之前发生，因此，虽然他的死使我们动摇，我们却因他得救的非凡方式而欢欣。如今，我们卓越的凯撒留已归还给我们，他那受尊敬的遗体和著名的尸身在一连串赞美诗和公开颂词中护送回家后，由他父母的圣洁之手所尊奉；他的母亲以宗教的节日礼服代替哀悼的装饰，以她的哲学战胜了眼泪，以咏唱圣诗平息了哀歌，因为她的儿子享有与他新重生灵魂相称的尊荣，这个灵魂是通过水，由圣神所转化的。\par

16. 这，\Person{凯撒留}，是我献给您的葬礼祭品，是我言语的初果，您常责备我迟迟不献上；然而若已献上，您却会将它弃去。我用这装饰来装点您，我深知，这装饰远比其他一切为您所珍爱，尽管它不是轻柔飘逸的丝绸织物——您生前以德行独为装点，并不像众人那样喜爱它；也不是透明的细麻布，也不是昂贵的香膏倾洒——您早已将这些让给了闺阁中的美人，它们的芬芳仅持续一日；也不是任何其他渺小之物，为渺小之心所珍视，而这一切今天都已被这苦涩的石头与您俊美的形象一同掩埋。愿希腊人的竞技与故事远离，那是不幸青年的荣誉，他们为微小的竞赛赢得微小的奖品；也愿一切奠酒、初果、花环和新摘的花朵远离，人们用这些来尊崇逝者，是出于古老的习俗和无理智的悲伤，而非理性。我的礼物是一篇演说，或许未来的时光会从我的手中接过它，并永远传扬，使之不让他——那位离我们而去的人——完全从大地消失，而是永远使我们所尊崇的人留在人们的耳中和心中，因为它向他们展示出我们哀悼之人的形象，比肖像画更清晰。\par

17. 这便是我的献礼；若它微薄且不配他的功德，天主喜爱那按我们能力所献的。我们礼物的部分现已完成，余下的我们如今将偿还，每年（我们这些仍存留的人）献上我们的尊崇与纪念。而现在，为你，神圣而圣洁的灵魂，我们祈求进入天堂；愿你享有\Person{亚巴郎}怀抱中的安息，愿你能看见天使的队列，以及圣徒们的荣耀与光辉；是的，愿你加入那队列，分享他们的喜乐，从高处俯视此地的一切——人们所谓的财富、可鄙的尊位、骗人的荣誉、我们感官的错谬、此生的纠缠及其混乱与无知，仿佛我们在黑暗中搏斗；而你则侍立在大王身旁，充满自祂发出的光明：愿我们日后能从那里领受一条不如此微弱的溪流——这溪流是我们在这镜子与谜语的日子里所想望的——而达到至善本身的源泉，以纯洁的心灵凝视纯净的真理，并因我们在此下为善的热切辛劳而获得报偿，即更完美地拥有和瞻仰至高之善：这是我们的圣书和导师所预言的，我们神圣奥秘的历程将引领我们到达的终点。\par

18. 如今还剩下什么？将圣言的治疗带给那些悲伤的人。对哀悼者而言，同情是一剂有力的良药，因为受苦者最能被那些承受同样苦难的人所安慰。因此，我特别向那些人说话，我为他们感到羞愧，如果他们在所有其他德行之外，没有展现出忍耐的要素。因为即使他们在爱子女方面超越所有人，愿他们同样在爱智慧、爱基督，以及特别在默想我们离世之事上超越众人，也将这默想印刻在子女身上，甚至使他们的整个生命成为死亡的准备。但若你的不幸仍蒙蔽你的理性，如同湿润模糊了我们的眼睛，使你无法清楚看见你的责任，那么，来吧，年长者，请接受一个年轻人的安慰；你们为父的，接受一个孩童的安慰，这本该受到与你们同龄的人、曾劝诫过许多人并从漫长岁月中汲取经验的长者的劝诫。然而，不要惊奇，若我以年少劝诫年长；若我在某些事上比白发者看得更清楚，我将它献给诸位。我们还有多久要活，你们可敬的年长者，如此接近天主？我们还要在此受苦多久？即便人的整个生命，与天主性体的永恒相比，也不长久，更不用说生命的残余，以及我所谓的我们人气息的分离，我们脆弱存在的终结。\Person{凯撒留}比我们领先了多少？我们还要为他的离世悲伤多久？我们不是正在奔向同一居所吗？我们不是很快就被同一块石头覆盖吗？我们不是不久就化为同样的尘土吗？在这些短暂的日子里，我们将得到什么益处呢？除非在我们看到、经历或甚至做了更多恶事之后，我们必须偿还那共同而不可避免的、对自然法则的赋税，即跟随一些人、在另一些人之前进入坟墓，哀悼这些人，被那些人哀哭，并从某些人那里得到我们曾向他人付出的眼泪之报偿。\par

19. 我的弟兄们，这就是我们度过这易逝生命的存在，这就是我们在世上的消遣：我们从无中生有，存在后又消散。我们是虚幻的梦，无影的幻象，\BibleRef{约 20:8} 如飞鸟的掠影，如海上不留痕迹的船只，如尘埃，如蒸气，如晨露，如速开速谢的花朵。至于人，他的日子如同草，如同田野的花，他茂盛一时。受感的\Person{大卫}曾言及我们的脆弱，又说过：\Quote{求你让我知道我日子的短少；} 他将人的日子定为\Quote{仅一拃之长。} 而你对\Person{耶肋米亚}作何感想？他因自己的诞生而哀怨母亲，且是为他人的过错。\BibleRef{耶 15:10} 我已见到一切，\BibleRef{训 1:14} 正如训道者所说：我在思想中回顾了所有人事——财富、享乐、权力、不稳的荣耀、逃避我们而非被获得的智慧；然后又是享乐、智慧，常常转圜相同的事物：感官的欢愉、果园、众多奴隶、财富储备、男女仆役、男女歌者、武备、长枪手、臣服的民族、收集的贡物、君王的骄傲，以及生命中一切必需与多余的事物，在这一切上我超越了我之前的众王。说了这一切之后，他还有什么可说？虚而又虚，万象皆虚，都是追风，这可能是指灵魂某种不合理的渴求，以及人因原初堕落而受罚的分心。但请听他所说的万事的结论：敬畏天主。这是他困惑中的支柱，也是你在此世下唯一的收获：从可见且摇动之物的混乱中，被引导到那稳固不动的事物那里。\BibleRef{格后 4:18} \BibleRef{希 12:27}\par

20. 那么，我们不要哀悼\Person{凯撒利乌斯}，而要哀悼我们自己，因为我们知道他所逃脱的、我们仍陷其中的灾祸，也知道我们将积蓄何种宝藏——除非我们紧贴天主，超越易逝之物，奔赴天上的生命，在仍在地上时便离开尘世，热切追随那提升我们的圣神。这对怯懦者是痛苦的，对勇敢之人却微不足道。让我们这样思考：\Person{凯撒利乌斯}将不再统治，反而会被他人统治；他不会恐吓任何人，但也不会惧怕任何凶恶的主人，尽管自己甚至常常不配做臣仆。他不再积聚财富，但也不会招致嫉妒，不为失败而痛苦，不再企图使收益翻倍。因为财富之症就是如此，对更多的欲望没有限度，不断地以饮来解渴。他不再炫耀口才，却因其言谈而被钦佩；他不再论述\Person{希波克拉底}和\Person{伽伦}的学说及其对头，但也不再为疾病所扰，不为他人的不幸而痛苦；他不再阐述\Person{欧几里得}、\Person{托勒密}和\Person{希罗}的原理，但也不再因无知者的浮夸之言而痛苦；他不再展示\Person{柏拉图}、\Person{亚里士多德}、\Person{皮浪}的学说，以及\Person{德谟克利特}、\Person{赫拉克利特}、\Person{阿那克萨戈拉}、\Person{克莱安特斯}、\Person{伊壁鸠鲁}等人，以及可敬的斯多葛学派和学园派全体成员的名号，但他也不再为解他们狡猾的三段论而烦恼。还需要更多细节吗？然而这些是众人所尊崇或渴望的：他身边不再有妻子或孩子，但他免于为他们或他们为他哀悼，免于将他们留给他人，或自己作为不幸的纪念被留下。他不继承财产，但他有那些最有用的继承人，正如他自己所愿，从而他离开时是个富人，带着他的一切。这是何等志向！何等新颖的安慰！他的遗嘱执行人何等宽大！一个宣言已被听到，值得众人倾听，母亲的悲痛已被一个美好而神圣的承诺所消除：完全将她儿子的财产作为对他的殡葬祭献，什么也不留给那些期待的人。\par

21. 这些安慰还不够吗？我要加上一个更有效的药方。我相信智者的言论：每一个美好而蒙天主喜爱的灵魂，一旦脱离身体的束缚，离开此世，便立刻享有对未来福乐的感受和知觉，因为那遮蔽它的东西已被净化或弃置——我不知道还能怎样称呼——它感到奇妙的喜乐和欢欣，欢欢喜喜地去迎接它的主，好像从今世痛苦的毒药中逃脱，抖落了束缚它、压制心灵翅膀的桎梏，从而开始享受为它预备的福乐，对此它已有所领悟。然后，稍后一些，它从大地——那给予并受托保管它的大地——领受那与它同类的肉身，这肉身曾与它一同追求天上的事物，并按照天主结合与分解的奥秘方式，与它一同进入那光荣的遗产。正如因紧密的结合而分享其苦难，同样也将一部分喜乐赋予肉身，将之完全收归己有，与它成为一体，在心神、理智和天主内，使那可朽与可变的被永生所吞没。至少请听听受默感的先知\Person{厄则克耳}如何讲述骨头与筋腱的接合，以及之后圣\Person{保禄}宗徒如何谈论地上的帐棚和天上非人手所造的房屋，一个要拆毁，另一个存在天上，声称离开身体便是与主同在，并哀叹在世的生活如同流放，因而渴望并急切地获得解脱。我为何在希望中气馁？为何像个朝生暮死的受造物？我等待总领天使的声音\BibleRef{得前 4:16}，最后的号角\BibleRef{格前 15:52}，诸天的转变，大地的变形，元素的解放，宇宙的更新\BibleRef{伯后 3:10}。那时我将看见\Person{凯撒留}本人，不再流亡，不再躺在棺架上，不再令人哀悼和怜悯，而是光明、光荣、属天，正如我在梦中常常看见你，最亲爱的兄弟，我的渴望如此描绘你，即使不完全是真理。\par

22. 但现在，放下哀叹，我将审视自己，省察我的情感，免得我不知不觉中内心有可哀叹之事。人子啊，因为这话是对你们说的，你们心硬且愚蒙要到几时呢？你们为何喜爱虚妄，寻找谎言，以为今世的生活是伟大的，这些日子是众多的，并且畏惧这离别——它本是可喜而愉悦的——仿佛它真是可悲可惧的？我们难道不认识自己吗？我们难道不抛弃可见之物吗？我们难道不仰望那不可见的吗？即使我们有些悲伤，难道不反而要为这漫长的居留而痛苦，如同圣\Person{达味}那样，称这里为黑暗的帐幕、苦难之地、深淤泥和死影？因为我们滞留于我们所携带的坟墓中，虽然我们是神，却像凡人一样死于罪恶。这就是我的恐惧，日夜伴随着我，不让我喘息：一边是光荣，另一边是惩罚之地；我渴望前者，直到能说：\Quote{我的灵魂因你的救恩而疲惫}；我颤栗地退缩于后者；然而我并不害怕我的身体在分解与朽坏中彻底消亡；而是害怕天主光荣的受造物——因为如果正直，它便是光荣的，正如如果有罪，它便是可耻的——在其中有理性和道德与希望，却要与畜生同受羞辱，死后并不更好；这样的命运是邪恶者所渴望的，他们配得那边的烈火。\par

23. 但愿我能致死我在地上的肢体\BibleRef{哥 3:5}，但愿我能将一切用于心神，走在窄而少人走的路上，而不是宽而容易的路\BibleRef{玛 7:13}。因为其后果是光荣而伟大的，我们的希望大于我们的功德。人算什么，你竟顾念他？这涉及我的新奥秘是什么？我渺小而伟大，卑微而崇高，可朽而永生，属世而属天。我与下界共享一种状态，与天主共享另一种；与肉身共享一种，与心神共享另一种。我必须与基督同葬，与基督同复活，与基督同作嗣子，成为天主的儿子，甚至成为天主本身。看，我们的论述进展到了何处。我几乎承认自己应感谢这灾难，它启发我如此的思想，使我更加热爱离开此世。这就是为我们而设的伟大奥秘的目的。这就是天主为我们而设的目的，祂为我们成了人，成了贫穷的\BibleRef{格后 8:9}，为复活我们的肉身\BibleRef{罗 8:11}，恢复祂的肖像，重塑人\BibleRef{哥 3:10}，使我们众人在基督内合而为一\BibleRef{迦 3:28}，祂在众人内完全成为祂所是的一切\BibleRef{格前 15:28}，使我们不再分男与女、化外人、西古提人、为奴的或自主的\BibleRef{哥 3:11}，这些都是肉身的标识，而只在我们内携带天主的印记，我们是由祂并为了祂而受造的\BibleRef{罗 11:36}，并且我们从祂那里领受了我们的形式和模型，以至于我们唯独被这印记所认识。\par

24. 是的，但愿我们所盼望的，按我们慷慨的天主的仁慈成就，祂要求少而赐予多，在现在和将来，都赐给那些真心爱祂的人；\BibleRef{格前 2:9} 为对他们的爱和希望，忍受一切，承行一切，事事感谢\BibleRef{得前 5:18}，无论顺逆：我指的是愉快和痛苦，因为理智知道连这些也常是救恩的工具；将自己的灵魂\BibleRef{伯前 4:19}和那些同行的旅伴——他们更准备好，已先我们安息——的灵魂托付给祂。现在，我们做完这事后，让我们停止我们的讲论，你们也停止你们的眼泪，因为你们现在正赶往你们的坟墓，那是你们给予\Person{凯撒留}的悲伤而长久的礼物，为他的父母在老年时早已预备妥当，如今却出人意料地赐给了他们年轻的儿子，但在祂——安排我们事务的那位——眼中并非无理。啊，万物的主和造物主，特别是这身体的主和造物主！啊，天主和父，以及属于祢之人的领航者！啊，生命与死亡的主！啊，我们灵魂的审判者和恩主！啊，万物按时由祢智慧的圣言——根据祢智慧与眷顾的深奥知识——所创造和改变者！求祢现在收纳\Person{凯撒留}，我们朝圣之旅的初果；如果最后的成了最先的，我们俯伏在祢治理万物的圣言前；也求祢日后在我们可寻见的时候收纳我们，按我们的益处让我们在肉身中生活多久；是的，收纳我们，准备好而不因祢的恐惧不安，不在我们末日离开祢，也不像贪恋世界和肉体的灵魂那样被强行夺走，而是充满对那在\Person{耶稣基督}我们的主内的有福而永恒生命的渴望，愿光荣归于祂，万世无疆。阿们。\par

\SourceRecord{Translated by Charles Gordon Browne and James Edward Swallow.；Nicene and Post-Nicene Fathers, Second Series；Vol. 7.；Edited by Philip Schaff and Henry Wace.；Buffalo, NY: Christian Literature Publishing Co.,；1894.；<http://www.newadvent.org/fathers/310207.htm>.}

% structure: p0001 source=h1 target=section reason=page-title

\section{Oration 8}

\Emph{论其姐果果尼亚。}\par

\EditorialNote{本篇讲道的确切年代不详。可以肯定（§23）是在凯撒留去世后（主历369年），并在其父去世前（主历374年）。我们从讲道本身以及一些作者对圣额我略诗歌的引述中所知不多（《诗·史》I.1.108,227）。讲道的地点几乎无疑是其结婚后所居住的城市。听众熟悉她婚后生活的公开细节。果果尼亚的父母和讲者本人虽为人所知，但提及的方式并不像在已知于纳西盎所发表的讲道中那样亲密。果果尼亚的神师和告解者也在场，显然处于权威地位，很可能坐在主教座上。果果尼亚的丈夫（《墓志铭》24）名叫阿利比乌。据克莱芒塞和伯努瓦一致引证埃利亚的权威，其家在依科尼雍，当时该城主教为福斯蒂诺。果果尼亚两个儿子的名字不详。埃利亚称他们二人都成了主教。圣额我略在其遗嘱中提到她的三个女儿：阿利比亚娜、欧根尼亚和农娜。这篇讲道以雄辩、虔诚和温情为特点，堪与论凯撒留的那篇媲美。}\par

\Emph{论其姐果果尼亚的葬礼演说。}\par

1. 在赞美我姐姐时，我是在尊敬我家族中的一员；然而我的赞美并非虚假，因为是因亲属而发，而是因真实而值得称赞；其真实不仅基于公正，而且基于众所周知的事实。因为，即使我愿意，也不允许偏袒；凡听我讲者，都像一个精明的评判者，站在我的讲辞与真理之间，以制止夸张，但若他是一个公正的人，便会要求真正应得的。所以，我担心的不是超过真理，而是相反，不及真理，因我极度不充分的颂词而减损了她应得的名誉；因为用适当的行动和言辞来匹敌她的优点实属难事。那么，我们切莫如此不公，去赞美他人的每一特质，而贬低我们自己的可贵品质，以致一些人因与己无关而得利，另一些人却因亲属关系而受损。因为无论赞美一方还是忽视另一方，都同样违背正义；若我们以真理为准绳和法则，只仰望她，无视凡俗和卑劣的一切，我们将根据每件事的价值去赞美或略过。\par

2. 然而，最不合理的是：既然我们拒绝认为欺骗、侮辱、控告或以任何方式——无论大小——不公正地对待亲属是正义的事，并认为对最亲近者行不义是最大的恶；我们竟会认为剥夺他们如此一篇讲辞（这最是善人应得的），而对恶人用更多言辞并乞求宽容，对待卓越者仅要求应得之物，这竟是正义之举。因为，若我们不被阻止（更公正的做法是）去赞美自己圈子以外的人，因为我们不认识且不能亲身见证其功德，那么，我们是否因友谊或众人嫉妒而被阻止去赞美我们所认识的人呢？尤其是那些已离世的人，为时已晚，不能再讨好他们，因为他们已逃脱，除其他一切外，已超脱赞美或责备的范围。\par

3. 既然已对这些作了充分辩护，并表明我作为讲者是何等必要，那么，让我继续我的颂词，摒弃一切雕琢和文辞的优雅（因我们所赞美者质朴，而质朴对她即是美），同时作为最不可推卸的债务，尽到所有应得的葬礼之礼，并且进一步教导众人热切效法同样的美德，因我每一言行之目的，在于促进托付于我者之成全。赞美我们逝者之故乡和家族的任务，我留给另一个人，他更严格遵循颂词的规则；他也不会缺少许多美好话题，若他愿以外在装饰来装扮她，如同人以黄金、宝石和工匠的艺术巧思去装扮一具光彩照人的美丽形体；这些装饰虽因对比而凸显丑陋，却无法为超越它们的美丽增添魅力。至于我，我只在提及我们共同的父母时稍循此规，因为忽略有如此父母和导师的大福是不恭敬的，然后我迅速转向她本人，不再多耗那些急于了解她是怎样的女性之人的耐心。\par

4. 有谁不知道我们末世的亚巴郎和撒辣，即额我略和他的妻子农娜呢？因为不提他们的名字作为对美德的激励是不妥的。他因信德而成义，她与忠信者同居；他超乎希望地成为多国之父（\BibleRef{罗 4:18}），她则灵性上为他们生产；他摆脱了父辈神祇的束缚，她是自由者的女儿和母亲；他为应许之地离开了亲族和家乡，而她是他流徙的原因；因为仅在这一方面，我冒昧为她求取高于撒辣的荣誉；他开始了如此崇高的朝圣之旅，她欣然与他同担劳苦；他将自己献于主，她既称丈夫为主，也以此看待他，并因此而成义；谁领受了应许，从他们按其所能生下了依撒格，谁就是恩赐。\par

5. 这位善牧是其妻子祈祷和引导的结果，他从她那里学到了善牧生活的理想。他慷慨地逃离了他的偶像，此后甚至驱逐了魔鬼；他从不与拜偶像者同席共食；以同等的尊荣、同心、同灵结合在一起，关注美德和与天主的共融如同关注肉体；寿数和白发相等，智慧和明达相等，彼此争胜，超越众人，在肉体上少有牵挂，在精神上已被提升，甚至在离世之前；不拥有世界，却拥有世界，因既轻视又善用它；舍弃财富，却因高贵的追求而富有；弃绝此地之物，换取彼处之物；拥有此生残存的少许，是从虔诚中所留，却拥有他们为之劳苦的丰盛和长久生命。关于他们，我再说一句话：他们被正确而公平地分派于两性；他是男子的光荣，她是女子的光荣，不仅是光荣，而且是美德之模范。\par

6. 戈尔戈尼亚的身世与声誉都源自他们；他们在她内播下了虔诚的种子，是她美好生活之源，也是她带着更美好的希望安然离世之因。这些是珍贵的特权，是许多以高贵出身自诩、以祖先为傲的人难以企及的。然而，若我必须用一种更富哲学与崇高精神的笔触来论及她的情况，那么，戈尔戈尼亚的故乡乃是天上的耶路撒冷（\BibleRef{希 12:22}），不是肉眼可见而是心神默观的对象；我们的家园就在那里，我们也正向那里奔赴；基督是那里的公民，与她同伍的是那些登记在天上的首生者的集会与教会，她们围绕着她伟大的奠基者默观祂的荣耀，参与那永恒的庆典。她的高贵在于持守了\OriginalTerm{肖像}{Image}，并与\OriginalTerm{原型}{Archetype}的完全相似，这相似由理性、德行与纯洁的渴望产生，在天主的事上越来越符合那些真正受教于天上奥秘的人；也在于知道我们从何而来、有何性质、为何受造。\par

7. 关于这些，我所知如此；因此我既发觉并断言，她的灵魂比东方人更高贵（\BibleRef{约 1:3}），乃是依据一种更优于常理的高贵与卑贱之分，其区别不在于血统，而在于品格；它不按家族来划分它所赞扬或责备的人，而是按个人。但既然我是在认识她的人中间谈论她的美德，但愿人人都贡献一己之见，助我一言；因为一人不可能面面俱到，无论他多么善于观察与领悟。\par

8. 在端庄方面，她如此出众，远远超越同时代的人，更不用说古时以端庄著称的人，以至于就所有人生命的两部分——即已婚与未婚状态——而言，后者更高超、更神圣，但也更艰难、更危险，前者则更卑微、更安全，她却能避免两者的弊端，并选取并融合两者中一切最好的：即后者的崇高与前者的安稳，从而成为端庄而不骄傲，将已婚与未婚的美德结合，并证明两者既不能绝对地将人捆绑于世界，也不能绝对地将人从世界分离（以至于一种因其本性必须完全避绝，另一种则应全然赞美）；而是理智高尚地主宰着婚姻与贞洁，将其作为德行的原料，在理性的巧手之下安排与塑造。因为虽然她进入了肉身的结合，却并未因此与精神分离；虽然丈夫是她的头，却并未忽略她的最初之头；而是按肉体法律的意志，或更确切地说，按赐予肉体这些法律的那位的意志，履行了少数属于世界与本性的事奉后，她将自己完全奉献给天主。而最卓越、最可敬的是，她还赢得了丈夫的认同，使他成为一位善良的同仆，而非无理的主人。不仅如此，她还使身体的果实——她的子女与孙辈——成为精神的果实，不仅将她自己的灵魂，而是将整个家庭与家族都奉献给天主，并因她在婚姻中的可亲可爱以及由此收获的美好果实，使婚姻变得光辉；她活着时，以自身为子女树立了一切美善的榜样；被召离世时，留下遗嘱作为对家庭的无声劝勉。\par

9. 神圣的撒罗满在他富有教益的智慧中，我是指他的\WorkTitle{箴言}，称赞那位顾家爱夫的妇女（\BibleRef{箴 31:10}），将她与那流连在外、不受约束、不端庄、以放荡的言语和行为猎取宝贵灵魂的妇人对比，而她则善于持家，勇敢地履行妇女的职责，手拿纺锤，为丈夫准备两件外衣，按时购买田地，妥善供应仆人的食物，以丰盛的餐桌款待朋友；以及其他一切细节，在那里他歌颂端庄勤劳的妇人。如今，若在这些方面赞美我的姊妹，就如同赞美一尊雕像的影子，或一头狮子的爪子，而未提及她最大的完美之处。谁比她更值得赞誉，而又如此避讳赞誉，使自己不被世人看见？谁更懂得端庄与欢愉的适当比例，使她的端庄不显冷酷，她的温柔不显放荡，而是在谨慎中结合了同情与尊严？你们这些贪图安逸、喜爱炫耀的妇人，听啊！你们轻视羞愧的面纱。有谁像她那样控制双眼？有谁如此轻视笑声，以至于一丝微笑对她而言都是大事？有谁更坚定地封闭耳朵？又有谁更向天主的圣言敞开耳朵？或者更确切地说，有谁让理智做舌头的统治者，宣称天主的审判？有谁像她那样约束嘴唇？\par

10. 于此，若你愿意，尚有另一端的卓绝之处：此端不仅她本人，亦系一切真正端庄贤淑的女性所不以为意者；然我等因那些过于偏爱装饰与炫耀、且拒听此类训诲之人，遂对此端多加看重。她从不曾以精工巧制之美金为饰，亦不以亚麻发辫——或全露或半掩——为美，更无螺旋卷发，亦无在可敬之首上搭建构型之人的不敬设计，亦无昂贵飘逸透明袍服之褶襞，亦无灿烂宝石之媚态——那些宝石染了周遭空气，于形体之上投下辉光；亦无画师之技艺与蛊惑，亦无那与天主相抗的地狱造物主之廉价之美，以其诡谲颜料掩蔽天主的造化，以其尊荣令之蒙羞，将娼妓的仿冒品置于渴慕的眼前，以取代天主的肖像，好让这私生之美窃走那本应为天主及来世保存的形象。然而，虽则她知晓女性外在装饰的种种多端，却无一比自己的品性与内在珍藏的光明更为宝贵。一种红晕为她所珍——羞赧的红晕；一种白净——节制的标记；至于颜料与描画、活生生的画像、流动的线条之美，她则留给舞台与街头的女子，留给所有认为羞耻为可耻、为羞辱之人。\par

11. 此类话题足矣。论及她的明智与虔诚，无人能给予充分叙述，亦难觅诸多典范，除却她本性及超性的父母——他们是她惟一的模范，而她的德行毫厘不爽地追随了他们，仅有一项最易认可的例外，即她既知晓又承认，他们乃她善行之源，亦是她自身光照之根。有何人的智识比她更敏锐？她被认作共同的顾问，不仅为其家族、同族、同一牧场之人，亦为周围所有人，他们视她的劝告与建议为不可违背的律法。有何言语比她更精辟？有何沉默比她更明智？既提及沉默，我当进而谈论她最为典型、最宜于女子、亦最利于今世的那一端。谁对天主之事有更充足的认识——既来自神圣谕示，亦来自自身的悟性？然而谁又更不轻言，将自己局限于女子当守的限度之内？此外，作为一位习得真实虔诚的女子本分——即那唯一堪当无餍之欲的光荣目标——谁像她那样以奉献物装饰圣殿，包括其他圣殿及这一座——如今她离去后，此殿恐难再得同样装饰？或更确切地说，谁将自己作为活殿呈献给天主？再者，谁如此崇敬司铎——尤其那位与她同为主之战士、虔诚导师之人，其美善的种子和一对奉献于天主的子女属谁？\par

12. 谁以更优雅和慷慨的款待向按天主之道生活的人敞开家门？而较此更为伟大的，谁以如此谦逊和敬神的问候迎接他们？此外，谁在苦难中显出更不动摇的心志？谁的灵魂对困苦者更富同情？谁的手对匮乏者更慷慨？我毫不迟疑以约伯之言尊崇她：\Quote{她的门向一切来者敞开，客旅未曾露宿街头。她为盲者是眼，为跛者是脚，为孤儿是母。}（参见\BibleRef{约 29:15-16; 31:32}）论及她对寡妇的怜悯，何须多言？其果实便是她自身从未被称为寡妇。她的家是她家族中一切贫乏者的共同居所；她的财物对一切需要者而言，与各人自己的财物同样公共。她\Quote{施舍散财，赒济贫乏}（\BibleRef{咏 112:9}），且按福音确实无误的真理，在上方的酒榨中积蓄了许多财宝，并时常在她曾是施惠者的人身上款待了基督（参见\BibleRef{玛 25:35-40}）。而最佳者，她毫无虚假的表白，在暗中于察见隐秘事的天主面前实践虔诚。一切她从这世界统治者手中夺回，一切她转移至安全的仓房。她未给大地留下任何事物，只留下身体。她以一切换取上方的希望；她留给子女的唯一财富是效法她的榜样，追随她的功德。\par

13. 但在这些令人难以置信的宽宏大量标记中，她并未将身体交给奢华与食欲的无节制快乐——那狂吠撕咬的狗——仿佛因她仁爱之工而自负，如同多数人一般，以怜悯穷人赎回自己的奢华，并以恶报善，而非以善医恶。她也没有在藉斋戒制服自己的尘土时，将硬卧褥的苦功留给他人；也没有在发现此端有属灵益处时，比任何人更少节制睡眠；也没有在这一点上如同已脱离身体般规范生活，当他人彻夜直立而战时，她却卧在地上——最克苦的人竭力如此。相反，她在此方面不仅超越女性，也超越最奉献的男性——藉其明智地咏唱圣咏，与神圣谕示交谈、阐发及恰当地引述，双膝跪地以致坚硬几乎生根，以痛悔之心与谦卑之灵流泪洗净污点，祈祷升向高天，心神脱离游荡而神往；在这一切中，或其中任何一项，有男或女能夸称超越了她？此外，说句重要的话，却是真实的：当她热切追求某些卓越之点时，对其他点她则是典范；有些点她乃是发现者，有些点她又超越他人。若在某一点上她被匹敌，她的卓越在于对全体的把握。她在一切点上如此成功，无人能在任何一点上达到中等程度；她达到每一点的如此完美，以至于任何一点本身足以取代全体。\par

14. 啊，未得照管的身体，褴褛的衣裳，其唯一的鲜花是德行！啊，灵魂，紧附于身体，因缺乏食物而几乎被减至无质状态；或者更确切地说，当身体被强行致死，甚至在分解之前，以使灵魂得以获得自由，并逃脱感官的纠缠！啊，守夜的夜晚，圣咏歌唱，以及从一天到另一天的站立！啊，达味，你的歌声从未使虔诚的灵魂感到厌倦！啊，柔弱的肢体，被抛在地上，违反自然，变得坚硬！啊，眼泪的泉源，在苦难中播种，为能欢乐地收割。啊，夜间的呼号，穿透云层，达于那住在天上的主！啊，神的热忱，在祈祷的渴望中，对夜间的犬、霜冻、雨水、雷声、冰雹和黑暗变得刚毅！啊，女性之本性，在共同的救恩之争中胜过男性，证明男女的区别在于身体而非灵魂！啊，圣洗的纯洁，啊灵魂，在你身体的纯洁内室中，你是基督的新娘！啊，痛苦的饮食！啊，厄娃，我们种族和罪恶的母亲！啊，狡诈的蛇和死亡，因她的自律而被克服！啊，基督的空虚自己、仆人的形状与苦难，因她的克苦而受尊崇！\par

15. 啊！我该如何数算她的一切特质，或是不伤及不知者，略过其多数呢？然而，在此应当补充她虔诚的赏报，因为我确实认为，你们这些熟知她生活的人，早已急切渴望在我的言辞中不仅听到现世之事，或她在彼处的喜乐——超越人的想象、听闻和眼见——也听到公义的赏报者在此世赐予她的一切：此事常有助于不信者的建树，使他们由微小之事达到对伟大之事的信德，由可见之物达到不可见之物。因此，我将提及一些众所周知的事实，以及另一些多数人从未知晓的隐密；那是因为她的基督徒原则注重不炫耀其眷顾。你们知道她那发狂的骡子如何拖着她的车奔跑，不幸翻覆，她被可怕地拖拽，严重受伤，令不信者因正义者遭遇如此横祸而反感，以及他们的不信如何迅速被纠正：因为，她全身受到挤压和挫伤，骨头和肢体，无论暴露的还是隐蔽的，她拒绝任何医师，只接受那准许此事的那位；既因她回避人的检视和手，即使在受苦中也保持她的端庄，也因她从准许此事的那位等待她的成义，以致她的获救除了祂以外别无所凭：结果，人们对她那出乎意料的康复之惊讶，不亚于对她的不幸之惊讶，并得出结论：这惨剧的发生是为了她藉苦难得光荣，痛苦是人的，康复是超人的，教训了后来的人，在苦难中彰显高度的信德，在灾祸中彰显忍耐，更高度地彰显天主对如此之人的慈爱。因为，对义人的美好应许\Quote{他虽跌倒，却不致彻底破碎}，如今又添加了更近的应许：\Quote{他虽彻底破碎，却必迅速被举起并受光荣。}因为，若她的不幸是不合理的，她的康复就是非凡的，以致健康很快掩盖了伤害，治愈变得比打击更著名。\par

16. 啊，非凡而奇妙的灾难！啊，比平安更高贵的伤害！啊，预言：\Quote{祂击打了，祂将包扎我们，使我们存活，三日后祂将使我们复活}，\BibleRef{欧 6:1}，它固然预示了一个更伟大、更崇高的事件，但也同样适用于格尔贡尼亚的苦难！此事众所周知，甚至远播远方，因为奇迹传遍各处，教训存入所有人的口舌和耳中，与天主其他奇事和德能并列。但接下来的事件，因我所说的基督徒原则以及她对虚荣和炫耀的虔诚回避，至今隐晦不为人知，您，这神圣羊群最优良、最完善的牧者，是否要我述说，并进一步赞同它？因为这秘密只托付给了我们，我们彼此见证了这奇迹，还是我们仍要对已逝的她持守信德？然而我认为，正如当时应默默无声，如今则是显明它的时候，不仅为光荣天主，也为安慰在苦难中的人。\par

17. 她身体患病，因一种异常而恶性的疾病而生命垂危，全身持续发热，血液时而激动沸腾，时而凝滞昏迷，呈难以置信的苍白，以及心智和肢体的麻痹：这并非长时间间隔，有时十分频繁。其毒性似乎超乎人力；医师们精心检查，单独会诊，都无济于事；她父母的泪水（常有大能力）也无用，公开的恳求和代祷（全体民众热切参与，仿佛为自己得救一样）也是如此：因为她的安全即众人的安全，而她的痛苦和疾病反成了共同的灾难。\par

18. 那么，这位伟大的灵魂，最伟大者的配得后裔，她为她疾病的良药是什么呢？因为我们如今已触及伟大的奥秘。在绝望于一切其他助佑后，她将自己交托给万有的医者，等待夜深人静的时刻，趁着疾病稍有间歇，她怀着信德走近祭台，呼求那在祭台上受尊崇的一位，以大声呼喊和各种祈求，忆起他昔日的一切大能作为，她深知那些古代和晚近的事迹，最后她鼓起一种虔诚而辉煌的大胆行为：她效法了那位因基督衣服的穗头而血漏枯干了的妇人。\BibleRef{玛 9:20}她做了什么？她又呼喊着将头靠在祭台上，泪流满面，正如那位曾用眼泪沾湿基督双脚的妇人，\BibleRef{路 7:38}并宣称若不痊愈绝不松手，然后她将药物涂遍全身，即她手中所珍藏的那份尊贵圣体圣血之\OriginalTerm{预像}{antitypes}，将泪水与之混合，啊，何等奇妙，她便离去，立时感到获得痊愈，身体、灵魂和心智都感到健康的轻松，因她的希望而获得了她所期望的，作为希望的赏报，借着属灵的力量获得身体上的医治。这些事尽管伟大，却并非虚假。你们所有人，无论有病或无病，都当相信，好能保持或恢复健康。我所说的并非夸口，这从她生前对我所揭示之事保持沉默便显而易见的。请确信，若非我担心如此伟大的奇迹对那些忠信者和不信者，无论今世还是后世，都完全隐藏起来，我如今也不会把它公开。\par

19. 这便是她的生命。她生活中的多数细节我未予陈述，以免我的言论过分冗长，又似乎过于贪图她的美名；但若我们丝毫不提她圣洁而光荣的死亡中的某些卓越之处，恐怕是亏待了她；尤其因为她如此渴望并期盼这死亡。因此，我将尽可能简明地加以叙述。她渴望离世，因为她对那召唤她的主大有胆量，而且她宁愿与基督同在，超越世上的一切。\BibleRef{斐 1:23}没有哪个最热情奔放的人，像她那样热爱自己的身体，以至于要挣脱这些锁链，逃离我们赖以生存的泥沼，纯洁地与那至美者结合，完全拥抱她的挚爱——我甚至要说，她的爱人——我们虽借着从他的光芒（如今虽微弱）得到光照，与他分离却得以认识他。这个愿望她也未曾落空，尽管这愿望神圣崇高；而更伟大的是，她借着预知和不断警醒，预先尝到了他的美。她仅有的睡眠将她带入无比的喜乐，而一次神视便涵盖了她预定的离世时刻，她已得知这一日，以便按天主的决定，她可预备妥妥而不受惊扰。\par

20. 她最近获得了洁净与成全的恩赐，即我们所有人都从天主领受的，作为我们\Emph{新}生命的共同恩赐和根基。或者更恰当地说，她的一生便是洁净与成全；虽然她从圣神领受了重生，但它的稳固乃是基于她先前的生命。我敢说，几乎唯独在她身上，这\OriginalTerm{奥迹}{mystery}更是一个印记，而非恩宠的恩赐。当她丈夫的成全成为她唯一尚存的愿望时（若你希望我简略描述那人，我不知除说他是她的丈夫外还能说什么），为使她全身心奉献于天主，不至于半成全地离世，或留下任何属于她而未成全的事物，这个祈求她也未曾落空，因为那成就敬畏他者之愿望并成全他们请求的主，必会应允。\par

21. 如今，当她一切顺心，愿望无一欠缺，预定的时刻临近时，她已为死亡和离世做好准备，便履行了此类情形中通行的法则，卧病在床。在向她的丈夫、子女和朋友们作了许多叮嘱之后（一如人们对一个充满夫妻之爱、母爱和手足之爱的人所期待的），并借着对天上之事的精彩讲论，使她的最后一日成为庄严的庆节之后，她便安眠了。她所充满的并非人的日子（因她不渴求这些日子，知道它们为她是邪恶的，且主要与我们的尘土和漂泊有关），而是极其充满天主的日子，远超我所想象任何一位在丰富白发中离世、历经多年的人。这样，她便得了释放，或更好说被接到天主那里，或飞升而去，或改变了居所，或略略提前了身体的离去。\par

22. 然而我几乎忽略了什么？但也许你，她的神师，不会允许我这样做，并且谨慎地隐瞒了这奇事，却让我知道了。这对她的卓越、我们对她的德行的记忆以及对她离去的遗憾而言，是一件大事。但我一想起这奇事，就战栗流泪。她正要离去，奄奄一息，周围是一群亲戚和朋友在行最后的善举，而她年迈的母亲俯身在她身上，灵魂因嫉妒她的离去而激动，痛苦和爱在所有人心灵中交织。有些人渴望听到一些炽热的话语，好铭记在心；另一些人急于发言，却无人敢开口；因为泪水无声，悲伤的剧痛无人安慰，因为似乎是一种亵渎，认为哀悼能成为对如此逝去者的荣耀。于是庄严肃静，仿佛她的死是一场宗教仪式。她躺在那里，看似无息、不动、无语；她身体的静止仿佛瘫痪，言语的器官似乎已死，因为那能驱动它们的东西已经离去。但她的牧者，在这场奇景中仔细注视着她，发现她的嘴唇在轻轻蠕动，于是将耳朵贴近它们，是他的性情和同情心鼓舞他这样做——但请你解释这神秘平静的含义，因为没有人会怀疑你的话！她低声重复着一篇圣咏——一篇圣咏的末句——说实话，这是她离去之勇敢的见证，有福的人能带着这些话安眠：\BibleQuote{在平安中我一躺下，即刻入睡。} 最美的女子，你就是这样歌唱的，这事就这样应验于你；圣歌成了现实，陪伴着你的离去，作为你的纪念，你已在受苦后进入甜蜜的安宁，并（超乎众人所得的安息）领受了那份属于爱人的安眠，正如一位在虔诚言语中生活并死去的人所应得的。\par

23. 我深知，你现在的命运，胜过眼所能见的，更为美好和珍贵：欢庆节日的歌咏、天使的群集、天上的万军、荣耀的景象，以及那至高的天主圣三所发出的纯洁无瑕、超越一切的光辉，不再被囚禁的心灵所限，被感官所分散，而是被完整的心灵所完全默观和拥有，并以全部神性的光明照射我们的灵魂：愿你充分享受这一切——那些你在地上时，通过你对其倾向的真实，只曾拥有其碎屑的事物。如果你仍关心我们的事务，且圣洁的灵魂从天主那里获得此特权，那么请接受这些话，以代替并优先于许多赞词——我曾在你们之前为凯撒里乌斯，并在你们之后为你献上这些赞词，因为我被保存下来为我的兄弟们致悼词。在你之后，若有人愿意给我同样的荣誉，我不得而知。然而，我唯一的荣誉是在天主内的荣誉，我的旅居和我的家是在我们的主耶稣基督内，愿光荣归于他，及圣父及圣神，直到永远。阿们。\par

\SourceRecord{Translated by Charles Gordon Browne and James Edward Swallow.；Nicene and Post-Nicene Fathers, Second Series；Vol. 7.；Edited by Philip Schaff and Henry Wace.；Buffalo, NY: Christian Literature Publishing Co.,；1894.；<http://www.newadvent.org/fathers/310208.htm>.}

% structure: p0001 source=h1 target=section reason=page-title

\section{第十二篇演讲}

\Emph{致父亲，当他把纳齐盎教会托付给我照管时。}\par

\EditorialNote{此篇演讲发表于公元372年。两年前，瓦伦斯将卡帕多细亚分为两省。提亚纳主教安提穆声称，教省应由帝国省份决定，要求对第二卡帕多细亚的教会拥有都主教权，这一主张与圣巴西略相悖，后者至今仍是未分裂省份的都主教。圣巴西略为维护其先前权利之持续性，在两省交界处的萨西玛设立一新主教区，并费尽周折劝服圣额我略接受祝圣为其首任主教。圣额我略虽俯首，却心有不甘，长时间不愿履任主教职责，最终由圣巴西略之弟、尼撒的圣额我略说服，尝试赴任。然而当他发现安提穆准备以武力阻止其进入时，便返回家中，毅然辞去主教之职，再次投身于他所深爱的独修生活。此后，他应父亲恳切请求，同意临时担任纳齐盎辅理主教之职，并在就职典礼上发表了这篇短讲。}\par

1. 我张开了口，汲取了圣神，我将我自己和我的一切献给圣神：我的行动和言语，我的不动和静默，惟愿祂扶持我、引导我，并移动手、心和舌头到正确之处，只要祂愿意；也按正确和有益的方式约束它们。我是天主的工具，一个理性的工具，一个由那精巧的工匠——圣神——调谐和弹奏的工具。昨日祂在我内的工作是静默。我默想不语。今日祂触动我的心吗？我的言语将被听见，我将默想发言。我既不如此多言，以致当祂倾向静默时仍欲说话；也不如此缄默无知，以致在当说话之时在我嘴唇前设防；而是按那同为一体的心智、圣言和圣神的旨意开启和关闭我的门。\par

2. 那么我要说话，既然我奉如此命令。我要向这里的善牧和他的神圣羊群说话，按我认为今日最适合我讲、你们听的方式。为什么你请求有人分担你的牧劳？我的讲话要从你开始，亲爱的、可敬的头，堪当亚郎的头，那属神和司铎的油膏沿它流下，流到胡须和衣袍上。为什么你虽有能力建立和引领许多人，且实际上以圣神的能力引领他们，却在属神工作中依赖杖和支柱？是因为你听说并知道，连显赫的亚郎也有亚郎的儿子厄肋阿匝尔和依塔玛尔受傅？(\\BibleRef{肋 8:2}) 我略过纳达布和阿彼胡，以免引喻不详；梅瑟在世时任命若苏厄接替他，作为向着许地前进者的立法者和将军？亚郎和胡尔在阿玛肋克被十字架(先已预表)打败的山上托住梅瑟双手的职务，(\\BibleRef{出 17:12}) 我情愿略过，认为不太适合或不适用于我们：因为梅瑟选择他们不是要分享立法工作，而是做祈祷的帮手和支持他疲倦双手的支柱。\par

3. 那么你有什么病？你的软弱是什么？是身体上的？我准备好支持你，是的，我已经支持，也被支持，如古时的雅各伯，藉着你父亲般的祝福。(\\BibleRef{创 27:28}) 是属灵上的？谁更强、更热切，尤其如今，当肉体的力量衰退消逝，如同许多妨碍和减弱光芒亮度的障碍？因为这些力量惯常彼此争战和对立，身体的健康以灵魂的疾病为代价，而当快乐随肉体一同平息消逝时，灵魂则繁茂向上看。但你之前的淳朴和高尚于我看来奇妙，你怎么无所畏惧，尤其是在这样的时代，不担心你的精神会被视为借口，大多数人会认为，尽管我们属灵的职业，我们这样做是出于肉体的动机？因为大多数人已使这职务被视为伟大、王侯般的，并伴有相当的享乐，即使一个人的管辖和统治比这更小的羊群，且其烦恼多于快乐。至此讲你的淳朴，或父母偏爱，如果是这样，它使你不承认自己，也不轻易怀疑他人任何可耻之事；因为难以被邪恶激动的心，也迟迟不会怀疑邪恶。我的第二个职责是简短地向你的子民，或现在也是我的子民，讲话。\par

4. 我被征服了，我的朋友们和弟兄们，因为我现在虽然当时没有，请求你们的帮助。我被父亲的老年，以及——用温和的话说——朋友的善意征服了。所以，请帮助我，你们谁能帮助，就向我这个被悔恨和热忱压碎和撕裂的人伸出手。悔恨建议逃遁、山岭和旷野，以及灵魂和身体的平静，并且心智应退回到自身，从可感觉的事物中召回其能力，以便与天主保持纯洁的共融，并被圣神闪射的光芒清楚地照亮，没有任何属地的或云雾之物掺杂或扰乱神圣之光，直到我们来到在此享受的光辉之源，悔恨和欲望都平息，当我们的镜子(\\BibleRef{格前 13:12})在真理之光中消逝时。另一种意愿是：我应该走上前来，为共同利益结出果实，并通过帮助他人而得到帮助；传播神圣之光，将一群子民献给天主作为祂的产业，圣洁的国民，王家的司祭，(\\BibleRef{伯前 2:9})并在许多灵魂中洁净祂的肖像。这是因为，正如一座公园比一棵树更好、更可取，整个天堂及其装饰比单独一颗星更好，身体比肢体更好，同样，在天主眼中，整个教会的革新比单个灵魂的进步更可取；因此，我不应只顾及自己的利益，也应顾及他人的利益。(\\BibleRef{斐 2:4}) 因为基督也是如此，当祂可能安享自己的尊荣和神性时，祂不仅如此空虚自己，取了奴仆的形状，而且还忍受了十字架，轻视了侮辱，(\\BibleRef{希 12:2})以便借自己的苦难毁灭罪恶，借死亡杀死死亡。前者是欲望的想象，后者是圣神的教训。而我，站在欲望与圣神之间，不知道应向哪一方让步，我将把我认为最好和最稳妥的道路交给你们，好使你们与我一同检验并参与我的计划。\par

5. 在我看来，取介于渴望与恐惧之间的中道，既部分满足渴望，又部分顺服圣神，乃是最佳且最稳妥之举。若我全然回避职分，从而拒绝恩宠，此为危险；若承担超出我能力的重担，又因它过于沉重。前者适合他人的身份，后者适合他人的能力，同时承担两者实属疯狂。然而，虔敬与安全都会建议我，按照能力来承担职分，如同对待食物，接受力所能及的，拒绝力所不及的，如此身体得健康，灵魂得安宁。因此，我现在同意分担我卓越父亲的重任，如同一只雏鹰，并非徒然地紧随着一只强大而高飞的老鹰。但此后，我将把我的翅膀交托给圣神，任由祂引领，无论何处、如何，都随祂的旨意：任何人不得违背祂的旨意，强迫或拖曳我往任何方向。因为承继父亲的辛劳是甘甜的，这群羊比起陌生异地的羊更为熟悉；我甚至敢说，在天主眼中更为宝贵，除非亲情令我迷醉，习惯使我失去觉察。再没有比自愿的牧者管理自愿的臣民更有益或更稳妥的道路了：因为我们的惯例不是靠武力或强迫带领，而是靠善意。否则，即便另一种政体也无法维系，因为被强迫维系的事物，一有机会就倾向于争取自由；而自由意志最维系我们的——我不称之为统治，而是——教导。因为虔敬的奥秘 \BibleRef{弟前 3:16} 属于那些自愿的人，而非被强迫的人。\par

6. 我的诸位好人，这就是我对你们所说的话，出于纯朴与善意，这也是我内心的秘密。愿那为你们和我双方都有益处的决定得胜，在圣神引导我们事务之下（因为我们的言论又回到同一点上），我们已将自己交付给祂，并用完美的油膏了头，归于全能圣父、独生圣子及圣神，祂就是天主。因为我们要\BibleQuote{将灯藏在斗底下}\BibleRef{玛 5:15}，向他人隐瞒对天主性的圆满认识，直到何时呢？现在理应将灯放在灯台上，光照所有的教会和灵魂，以及世界全部的圆满整体，不再借助比喻或理智的草图，而是明确的宣认。这确实是对那些在基督耶稣、我们的主内堪当此恩宠的人，最完善的天主论阐述。愿光荣、尊威与权能归于祂，至于永远。阿们。\par

\SourceRecord{Translated by Charles Gordon Browne and James Edward Swallow.；Nicene and Post-Nicene Fathers, Second Series；Vol. 7.；Edited by Philip Schaff and Henry Wace.；Buffalo, NY: Christian Literature Publishing Co.,；1894.；<http://www.newadvent.org/fathers/310212.htm>.}

% structure: p0001 source=h1 target=section reason=page-title

\section{第十六篇讲道}

\Emph{关于父亲因冰雹灾害而沉默}\par

\EditorialNote{本讲道属于公元373年。一系列灾难降于\Place{纳齐盎}的人民。一场致命的牛瘟摧毁了它们的畜群，随后是长期干旱，而今它们即将成熟的庄稼又被一场暴雨和冰雹毁坏。人民涌向教堂，发现\Person{圣额我略（长）}被这些可怕灾祸之感压得无法对他们说话，便恳求其助理登上讲坛。此场合无时间准备，因此圣额我略倾泻其情感于一席话中，从最充分的意义上讲，那是即兴之作。然而，如\Person{伯努瓦}所建议的，其现存形式可能源于对当时速记记录的后期润色。}\par

为何破坏公认的秩序？为何对那受法律约束的口舌施以暴力？为何挑战那服从圣神的言论？为何豁免了头，却急于奔向脚？为何越过\Person{亚郎}，而催促\Person{厄肋阿匝尔}？我不能容许源泉被堵，而溪流继续流淌；太阳被隐藏，而星星照耀；白发退隐，而青年制定法律；智慧沉默，而无经验贸然发言。倾盆大雨并非总比细雨更有益。不，若过于猛烈，它冲走泥土，增加农夫的损失；而轻柔的雨滴深深渗入，滋润土壤，有益农人，使庄稼长成丰盛。因此，流利的言辞并不比智慧更有益。因为前者虽或许带来短暂愉悦，却随即消散，如它所击打的空气一般毫无效果，尽管它以口才取悦贪婪的耳朵。但后者深入心灵，大大张开它，充满圣神，并显出其比起源更高贵，以寥寥数语产生丰盛收获。\par

我尚未提及那真正的和首要的智慧，为此我们杰出的农夫和牧者著称。首要的智慧是值得赞美的生活，为天主保持纯洁，或为那至纯至明者而净化，祂向我们要求祂唯一的祭献——净化，即痛悔的心和赞颂的祭献，以及在基督内的新受造物（\BibleRef{格后 5:17}），与新人（\BibleRef{弗 4:24}），以及如圣经所喜称的诸如此类。首要的智慧是轻视那由语言、修辞和虚假不必要修饰构成的智慧。我宁愿在教会中用我的理智说五句话，而不愿用舌音说一万句话（\BibleRef{格前 14:19}），以及像那无意义的号角声，不能唤醒我的士兵投入属灵争战。这就是我所赞美、所欢迎的智慧。借此，卑贱者得了名声，被轻视者得了最高荣誉。借此，一伙渔夫用福音网的网眼捕获了全世界，并以一句成就而缩短的话战胜了那归于无有的智慧（\BibleRef{格前 2:6}）。我不认为那巧言辞的人是智慧的，也不认为那口齿伶俐但灵魂不稳和未受纪律的人，它们如同坟墓，外面看起来华美，里面却充满尸骨（\BibleRef{玛 23:27}）和各种恶臭；而是那极少谈论德行，却在实践中给出众多例子，并且以生活证明其言语可信的人。\par

在我看来，我们可以凝视的美比用文字描绘的美更美；我们双手能持的财富比梦中想象的财富更有价值；行为所证实的智慧比华丽言辞所陈述的智慧更真实。因为\Quote{明智}，他说，\Quote{凡实行的人都有}，而非那宣告的人。时间是这种智慧的最佳试金石，\BibleQuote{白发是荣耀的冠冕}（\BibleRef{箴 16:31}）。因为，如我和\Person{撒罗满}所想，我们\BibleQuote{在死前不可称人有福}（\BibleRef{训 11:28}），而且\BibleQuote{今天会发生什么事，不可预知}（\BibleRef{箴 27:1}），因为我们在此世的生命有许多转折，我们卑微的身体（\BibleRef{斐 3:21}）不断上升、下降和变化；那么，那个几乎没有过失、几乎饮尽生命之杯、几乎抵达共同存在之海港口的人，必然比那仍有漫长航程的人更安全，也因此更令人羡慕。\par

不要约束那有许多高贵言论和果实的舌头，它生了许多义子——举目四望（\EditorialNote{依 49:18}），看有多少儿子，何等财宝！就是这全体人民，你们藉着福音在基督内所生（\BibleRef{格前 4:15}）。请勿吝惜我们那些精炼而非繁多的言语，且不要让我们提前尝到即将失去的滋味。请以那些虽少却对我亲爱且最甘甜的话发言，那些虽几乎听不见却因其精神呼声而被感知的话，正如天主听见\Person{梅瑟}的沉默，对心中祈祷的他说：\BibleQuote{你为什么向我呼号？}（\BibleRef{出 14:15}）安慰这人民，我恳求你，我，曾受你养育，后成为牧者，如今甚至是大牧者。请给我——在牧者技艺上，给这人民——在服从上，上一课。请就我们目前的沉重打击，就天主公义的审判，谈一谈，无论我们理解其意义，还是对其深渊无知。又\BibleQuote{慈悲被置于天平}（\BibleRef{依 28:17}），如\Person{圣依撒意亚}所言，因为良善不是没有辨别的，正如葡萄园中的首批工人（\BibleRef{玛 20:12}）所想象的，因为它们看不出同酬的人之间有什么区别；以及那愤怒——被称为\Quote{主手中的杯}和\Quote{使人跌倒、被饮尽的杯}——是与过犯成比例的，即使祂向所有人稍减其应付的，并以慈悲稀释祂愤怒的纯酒。因为他从严厉转向宽容，向着那些以敬畏接受管教的、那些在轻微磨难后怀孕并痛苦于悔改、最终产生完美的救恩精神的人；但他仍然保留愤怒的渣滓，即最后的点滴，要将其完整倾注在那些非但未被他的仁慈治愈，反而变得顽梗的人身上，如同那铁石心肠的\Person{法郎}，那残酷的监工，他被树立为天主在恶人身上权能的榜样（\BibleRef{罗 9:17}）。\par

5. 请告诉我们，这样的打击和鞭笞从何而来，我们又能对此作何解释。是因为宇宙中某种紊乱无序的运动，或某种无向导的激流，某种宇宙的失理性，仿佛宇宙没有统治者，因而被偶然性带着走——正如那些愚蠢地自诩智慧之人的学说，他们自己也被黑暗的混乱之灵随意牵引？还是说，宇宙的扰动与变化（这宇宙最初被建立、调和、结合、并按唯有赋予它运动的那位所知的和诸而运动）是在眷顾之缰绳引导下，被理性和秩序所支配？饥荒、龙卷风、冰雹——我们现时受到的警告打击——从何而来？瘟疫、疾病、地震、海啸、以及天上可怖之事又从何而来？那曾为人类享乐而安排、为他们共同而平等的喜悦的受造界，如今如何变成了惩罚不敬神之人的工具，为的是我们通过那在我们受尊崇时未曾感恩的事物而受惩戒，并在我们的苦难中认识到那在我们受惠时未曾认识的权能？为何有些人在主手中领受了加倍的罚，为他们的罪恶，\EditorialNote{依 40:2}而他们邪恶的量度被加倍填满，如同在以色列的惩戒中；而另一些人的罪却因七倍的报应归入他们怀中而被消除？尚未满盈的阿摩黎人的恶贯如何（\BibleRef{创 15:16}）？罪人是被放过，还是再次受罚？被放过，或许是因被保留到另一世界；受罚，是由于在此世借此得到医治。义人在不幸时可能受到考验，在顺遂时可能被观察，看他是否心灵贫穷，或未能远超越可见之物——正如我们的良心，我们内在且无误的审判官，所告知我们的。我们的灾祸是什么，它的原因又是什么？这是对德性的试炼，还是对邪恶的试金石？是将其视为惩罚而谦卑地顺服于天主大能的手下更好（\BibleRef{伯前 5:6}），还是将其视为试炼而超越它更好？在这些事上，求您教导并警戒我们，免得我们因眼前的灾祸而过于沮丧，或坠入邪恶的深渊而轻视它（因为这种情绪相当普遍），而是使我们能安静地承受训诫，不要因为对这训诫麻木而招致更严厉的。\par

6. 歉收的季节和农作物的损失是可怕的。当人们已在希望中欢欣，并几乎算作已收获的储藏时，更是如此。不时的收获又是可怕的，那时农夫带着沉重的心情劳作，仿佛坐在他们庄稼的坟墓旁——这些庄稼曾被温和的雨水滋养，却被狂暴的风雨连根拔起，割禾者不满手，捆禾者不满怀，也未曾得到路人赐予农夫的祝福。看到土地荒废、被清除、被剥夺装饰的景象的确悲惨，关于此，有福的岳厄尔在他最悲壮的关于土地荒芜和饥荒之鞭的描绘中哀叹（\BibleRef{岳 1:10}）；而另一位先知哀叹着，将土地昔日的美丽与最后的混乱对比，如此论及主击打大地时的愤怒：在祂面前是伊甸乐园，在祂背后是荒凉的旷野（\BibleRef{岳 2:3}）。这些事诚然可怕，甚至比可怕更甚，当我们只为眼前的事忧伤，尚未因更重打击的感觉而痛苦时；因为如同疾病，不断折磨我们的痛苦比未临到的痛苦更痛苦。但更可怕的是那些天主义怒的宝库中所含的，愿天主阻止你们去尝试；你们也不会尝试，如果你们逃往天主的仁慈，用眼泪赢得那愿意施仁慈者的心（\BibleRef{欧 6:6}），并通过悔改避免祂余怒的发作。至今，这还是温和、慈爱和温和的斥责，以及训练我们年轻岁月的鞭笞最初元素。至今，只是祂怒气的烟，祂折磨的前奏；尚未降下火焰，祂发怒的高潮；尚未燃起炭火，最终的鞭笞，祂曾威胁要降下一部分，当祂将另一部分举在我们之上时，祂强行忍住了一部分，当祂将另一部分降在我们身上时；祂用威胁和打击同等来教导我们，为祂的义怒开辟道路，出于祂良善的丰盈；从轻微的惩处开始，以免需要更严厉的；但若被迫，祂也准备好用更重的来教导我们。\par

7. 我知道那闪耀的剑（\BibleRef{则 21:9}），那在天上被浸醉的刀锋，被命令杀戮、毁灭、使无子女，且不怜惜皮肉、骨髓和骨头。我知道那一位，祂虽不受激情所动，却像失去幼崽的母熊，又像在亚述人路上的豹子，来迎接我们（\BibleRef{欧 13:7}）——不仅那时的亚述人，若有人现今在邪恶上成为亚述人，亦然；没有人能逃脱祂义怒的力量和迅速，当祂注视我们的不敬，以及祂的嫉妒——这嫉妒知道怎样吞噬祂的敌人，追赶祂的仇敌至死（\BibleRef{欧 8:3}）。我知道倾空、荒废、使成为荒凉、心肠消融、双膝相碰（\BibleRef{鸿 2:10}）——如此是不敬神之人的惩罚。我不详述将来的审判，因为此世的放纵使我们免于将来的审判，因为在此世被惩罚和洁净，比被移交到将来的苦刑更好——那时是惩罚的时候，不是洁净的时候。因为正如在此世记念天主的人战胜死亡（正如达味绝妙地歌唱过），同样，去世的人在地里没有宣认和恢复；因为天主将生命和行为限制在此世，而将来是对已做之事的审查。\par

8. 在眷顾之日，\EditorialNote{依撒意亚 10:3} 我们该做什么？这日子是一位先知用来恐吓我的，无论是天主向我们施行公义审判的日子，或是我们听说的那临到山岭的日子，或无论是什么、无论何时，当祂与我们理论、反对我们，将那些苦毒的控告者——我们的罪恶——摆在我们面前，将我们的过犯与祂的恩惠相比较，以思想对思想，以行为对行为，并要我们为那被邪恶所模糊损坏的肖像\BibleRef{创 1:26}交账，直到最后祂领我们离去，我们被自我定罪、自我谴责，再也不能说我们受了不公正的待遇——这个念头甚至在这里有时能安慰那些在谴责中受苦的人。\par

9. 但那时我们有什么辩护者？有什么托词？有什么虚假的借口？有什么听起来合理的诡计？有什么违背真理的伎俩能欺骗法庭，剥夺它公正的判决？这法庭为我们所有人将我们整个生命、行为、言语、思想放在天平上，将善与恶相称，直到那较重的得胜，判决按主要倾向作出；此后没有上诉，没有更高法庭，没有以随后行为来辩护的余地，没有从明智的童女或从卖油者那里为即将熄灭的灯获得油\BibleRef{玛 25:8}，没有那在火焰中受苦的富翁的悔改\BibleRef{路 16:24}，也没有为他的朋友恳求悔改，没有诉讼时效；只有那最终可怕的审判宝座，比可怕更公义；或者更确切地说，更可怕因为也是公义的；当宝座设立，万古常存者就座\BibleRef{达 7:9}，书卷展开，火河流出，祂面前的光和预备的黑暗；行善的进入复活得生命\BibleRef{若 5:29}，如今在基督内隐藏\BibleRef{哥 3:3}，将来与祂一同显现；行恶的进入复活受审判\BibleRef{若 5:29}，不信的人已因那审判他们的圣言被定罪。有些人将受到不可言说的光明和神圣君王的天主圣三的目睹的欢迎，这光明现在以更大的光辉和纯洁照耀着他们，将自己完全与整个灵魂结合——我认为天堂的王国唯独在于此，超越一切；其他人则要忍受被天主抛弃和那无限制的良心羞耻，此外还有别的痛苦。但这些容后再谈。\par

10. 如今我们该怎么办，我的弟兄们？当我们被压碎、被击倒、酩酊大醉，但并非因烈酒或葡萄酒\EditorialNote{依撒意亚 29:9}——这些只是暂时令人兴奋和迷惑——而是因主加给我们的打击。祂说：\BibleQuote{你这心啊，要激动战栗}\BibleRef{哈 2:16}，又将忧愁和沉睡的灵给轻视者喝。祂也对这些人说：\Quote{看，你们这些轻蔑的人，观看、惊奇而灭亡吧！}我们怎能承受祂的定罪？或当祂责备我们时，我们能回答什么？祂不仅责备我们忘恩负义于许多恩惠，还责备我们拒绝接受祂的惩戒和医治。祂称我们为祂的儿女\BibleRef{申 32:5}，却是不配的儿女；祂的儿子们，却是乖僻的儿子，在他们无路崎岖之地跛行跌倒。\Quote{我如何能教导你们呢？我不是已经做了吗？用温和的方法？我已经施行了。\BibleQuote{我越过了在埃及的第一灾，从井、河和一切蓄水池中喝血}\BibleRef{出 7:19}；我越过了青蛙、虱子、苍蝇等下一灾。我从第五灾开始，打击羊群和牛群，暂时放过理性受造物，我击打牲畜。你们轻视这打击，比那些被击打的牲畜更不理智、更不留意我。我停住雨水不下给你们；\BibleQuote{一块地下了雨，另一块未得雨而枯干}\BibleRef{亚 4:7}，你们却说：\InnerQuote{我们要硬挺过去。}}\par

11. 或许祂要对我说——我这甚至不被打击所改造的人——\Quote{我知道你是顽梗的，你的颈是铁筋\EditorialNote{依撒意亚 48:4}，疏忽的人仍然疏忽，不法的人仍然行不法，天上的管教毫无作用，鞭打也毫无作用。\BibleRef{耶 6:29}正如我从前藉耶肋米亚的口责备你们：工匠徒然炼银，你们的邪恶没有炼掉。\BibleRef{耶 6:29}你们能忍受我的愤怒吗？上主说。难道我的手没有能力将别的灾祸降在你们身上吗？在我的命令下还有那从炉灰中爆发的脓疮\BibleRef{出 9:10}，若将之向天撒去，梅瑟或任何天主行动的执行者，就能用疾病惩罚埃及。还有蝗虫、可触摸的黑暗，以及那按顺序最后、在痛苦和权能上却为首的一灾——长子被毁灭和死亡。}为躲避这灾，为使毁灭者转离，更妥当的是用伟大的拯救记号，即新盟约的血，洒在我们心思、默观和行动的门框上，藉着与基督一同被钉死，一同死亡，好使我们能与他一同复活、一同得荣耀、一同为王，无论现在还是在他最后的显现时；而不致被打破、压碎、哀叹，当那可悲的毁灭者在黑暗的此生中太迟地击打我们，摧毁我们的长子——我们生活中献给天主的产物和结果。\par

12. 但愿我永远不曾在其他惩罚中，被那良善者这样斥责——祂因我的乖戾而在愤怒中与我作对\BibleRef{肋 26:27}——\BibleQuote{我用热风和霉病打击你们，还有枯萎；却无效果。外面的刀剑\BibleRef{申 32:25}使你们丧子，但你们仍不归向我，上主说。}愿我成为不了那钟爱的葡萄树，它曾被种植、围护，并有篱笆和塔楼及一切可能的手段加以稳固，但当它野化并长出荆棘时，就因此被鄙视，它的塔楼被拆毁，篱笆被撤除，不再被修剪或挖掘，反被所有人吞噬、荒废和践踏！\EditorialNote{依 5:1}这是我感到必须说的话，关于我的恐惧；我为这一打击如此痛苦，我还要告诉你们，这是我的祈祷。\BibleQuote{我们犯了罪，行了不义，做了恶事}\BibleRef{达 9:5}，因为我们忘记了你的诫命，随从了我们自己的恶念\EditorialNote{依 65:2}，因为我们行事与你的基督的召叫和福音，以及祂为我们所受的苦难和屈辱不相称；我们成了你所钟爱的司铎和百姓的耻辱，我们一同犯错，我们都偏离了道路，一同成了无用的，没有一人行审判和正义，连一个也没有。我们以我们的邪恶和行为的乖张，缩短了你仁慈、善良、怜悯和同情，我们在其中背离了你。你是善的，但我们行了不义；你含忍，但我们应受鞭打；我们承认你的良善，虽然我们无知，我们只为少数过错受了鞭笞；你是可畏的，谁能抵抗你？群山在你面前颤抖；谁能与你臂膀的力量抗衡？你若关闭天空，谁能开启？你若放开你的洪流，谁能阻挡？在你眼中，使人贫穷和富有，使人活和死，打击和医治，都是小事，你的旨意是完美的行动。\BibleQuote{你发怒，我们犯了罪}\EditorialNote{依 64:5}，一位古人如此忏悔；现在是我说相反话的时候了：\Quote{我们犯了罪，你发怒；}因此我们成了邻人的耻辱。你转面不顾我们，我们就充满了羞辱。但是，求主止住，求主停止，求主宽恕，不要因我们的过犯永远抛弃我们，不要让我们的惩罚成为他人的警戒，而我们却可以从他人的考验中学习智慧。从谁那里？从那些不认识你、没有服从你权柄的列国和邦国。但是，上主，我们是你的子民，你产业的杖；因此求你纠正我们，但出于良善，而非你的愤怒，免得你使我们归于无有\BibleRef{耶 10:24}，并成为地上所有居民的轻蔑。\par

13. 我说这些话是求取怜悯；若能用全燔祭或牺牲平息祂的愤怒，我连这些也不会吝惜。你们也要效法你们战栗的司铎，我钟爱的孩子们，与我同享天主的管教和慈爱。以眼泪保守你们的灵魂，通过修正生活方式来止息祂的愤怒。祝圣一个斋期，召开圣会\BibleRef{岳 2:15}，如同有福的岳厄尔与我们一同嘱咐你们：召集长老和吃奶的婴孩，他们柔弱的年龄赢得我们的怜悯，尤其堪得天主的慈爱。我也知道，他嘱咐我，天主的仆人，以及你们，蒙受同样尊荣的人，要我们披上麻衣进入祂的殿，在门廊和祭台之间日夜哀哭，带着可悲的装束和更可悲的声音，不停地为我们自己和百姓大声呼求，不吝惜任何辛劳或言语，以平息天主：说\BibleQuote{上主，求你宽恕你的百姓，不要将你的产业交给侮辱}\BibleRef{岳 2:17}，以及其余的祷词；我们在对苦难的感受上，如同在地位上一样超越百姓，我们亲身教导他们痛悔和改正邪恶，以及随之而来的天主的含忍和鞭打的止息。\par

14. 来吧，我所有的弟兄，让我们俯伏敬拜，在主我们的造主面前哭泣；让我们在不同年龄和家族中设立公开的哀悼，让我们高声恳求；让这声音，而不是祂所憎恨的呼喊，进入万军的上主耳中。让我们通过认罪来预先面对祂的愤怒；让我们渴望看到祂在发怒后得到平息。他说：\BibleQuote{谁知道，祂是否会转意后悔，在身后留下祝福？}\BibleRef{岳 2:14}我确知这事，我是天主慈爱的保证人。当祂放下那不属于祂本性的愤怒时，祂将转向那属于祂本性的慈悲。前者是祂被我们强迫的，后者是祂所乐意的。如果祂被迫打击，祂当然会按本性克制。只求我们怜悯自己，为我们父亲的公义情感打开道路。让我们含泪播种，以便在喜乐中收割，让我们显明自己是尼尼微人，而非索多玛人。让我们改正我们的邪恶，免得我们被它消灭；让我们听从约纳的宣讲，免得被火和硫磺吞没；如果我们离开了索多玛，就让我们逃到山上，逃到左哈尔，在日出时进入它；让我们不要停留在整个平原上，不要回头看，免得我们被冻结成一根盐柱，一根真正不朽的柱子，控告那回到邪恶的灵魂。\par

15. 我们确知，不作恶实属超人性之事，只属于天主。我不提天使，免得为恶感留余地，给有害的争执以机会。我们的未痊愈状态源于我们邪恶未驯的本性，及其能力的运用。犯罪时的悔改，是人的行为，却表明善人属于那在救恩之道上的一份子。因为，即使我们的尘土沾染了些许邪恶，地上的帐幕压迫灵魂向上飞升——这灵魂至少被造得向上飞翔——\BibleRef{智 9:15}但要让肖像从污秽中洁净，并将与它同轭的肉身高举，以理智的翅膀托起；更好的是，我们既不需要此洁净，也无需被洁净，借着保持我们原始的尊严——我们在此世的训练正加速趋向这尊严——并且不因罪的苦味而被逐出生命树：虽然，犯错后回转，比跌倒时不受责罚更好。因为\BibleQuote{主所爱的，祂必惩戒}\BibleRef{箴 3:12}，而责备是父亲般的行动；每个不受惩戒的灵魂，都是未痊愈的。那么，不受惩戒不是一件难事吗？但受惩戒而不改正则更难。一位先知论到以色列，其心刚硬未受割礼，说：\BibleQuote{主啊，祢击打他们，他们却不忧伤；祢消灭他们，他们却拒绝接受管教}\BibleRef{耶 5:3}；又说：\BibleQuote{人民不转向击打他们的那一位}\EditorialNote{依撒意亚 9:13}；\BibleQuote{我的人民为何以永久的退后而退后}\BibleRef{耶 8:5}，为此他们将被彻底击碎毁灭？\par

16. 兄弟们，\BibleQuote{落入永生天主的手中}\BibleRef{希 10:31}是可怕的；主的面向作恶的人也是可怕的，祂以彻底的毁灭消除邪恶。天主的耳朵是可怕的，它甚至倾听亚伯尔通过无声之血所说的话。祂的脚是可怕的，追逐恶行。祂充满宇宙也是可怕的，以致无法在任何地方逃脱天主的行动\BibleRef{耶 23:24}，即使飞上天空，或进入阴府，或逃到极东之地，或藏身于海洋的深处和尽头。厄耳科士人纳鸿在我面前战栗，当他宣告尼尼微的神谕：\BibleQuote{天主是忌邪者，主在愤怒中向祂的仇敌复仇}\BibleRef{鸿 1:1}，并施用如此严酷的惩罚，以致对恶人再无施罚的余地。因为每当我听到依撒意亚威胁索多玛的人民和哈摩辣的统治者，说：\BibleQuote{你们为何还要被打，在罪上加罪？}\EditorialNote{依撒意亚 1:10}我几乎充满恐惧，泪流满面。他说：\BibleQuote{无法找到任何打击加到过去的打击之上，因为你们新添的罪过；你们彻底经历了全部，穷尽了各种形式的惩戒，永远以你们的邪恶招致新的惩戒。没有伤口、没有青肿、没有溃烂的疮；瘟疫感染全身，无法医治：因为无法涂抹药膏、油膏或绷带。}我略过其余的恐吓，免得加重你们目前的灾祸。\par

17. 唯愿我们认清这恶的目的。为何五谷枯萎，粮仓空虚，羊群草场枯竭，地上的果实被扣留，平原充满了羞耻而非肥美：为何山谷哀叹，不丰产谷物，群山不滴下甘甜——正如日后对义人所行的——反而被剥夺和羞辱，并且反受到基耳波亚的诅咒\BibleRef{撒下 1:21}？整个大地变得如起初尚未被美丽装饰之时。祢眷顾大地，使大地喝足——但这眷顾为的是恶，这饮水是毁灭性的。哀哉！何等景象！我们丰饶的庄稼化为残茬，我们所撒的种子从稀疏的遗存中辨认出来，我们的收成——其来临我们根据月份数计算，而非根据成熟谷物——几乎不能为主带来初熟之果。不敬虔者的财富就是如此，粗心播种者的收成也是如此；正如古时的诅咒所说：\BibleQuote{期望多多，收获少少}\BibleRef{盖 1:9}，\BibleQuote{播种而不收割，栽种而不压榨}\BibleRef{申 28:39}，\BibleQuote{十亩葡萄园只出一罢特}\EditorialNote{依撒意亚 5:10}，并听到他乡丰收，而我们自己被饥荒所迫。这是为什么？裂口的起因是什么？我们不要等待别人来指证，让我们作自己的监察员。对邪恶的一剂重要良药就是承认，并小心避免跌倒。我愿先这样做，因为我已经从高处向我的子民作了报告，尽了一位守望者的职责。因为我未曾隐瞒刀剑的来临，为能拯救自己的灵魂\BibleRef{则 33:3}和听众的灵魂。因此我现在要宣布我子民的不顺服，把他们的罪当作我自己的，或许这样我能得到一些温柔和缓解。\par

18. 我们中有人压迫穷人，夺取他的土地份，用欺诈或暴力非法侵占他的地界，将房屋连接房屋，田地连接田地，以掠夺邻人的某物，并渴望没有邻人，好独自住在地上。\EditorialNote{依撒意亚 5:8} 有人以高利贷和利息玷污土地，既在不曾播种的地方收取，又在不曾分散的地方收割，\BibleRef{玛 25:26} 耕种的不是土地，而是贫苦者的需要。有人夺取了属于天主——万有的赐予者——的初熟谷物和初榨葡萄酒，\BibleRef{拉 3:8} 显得既忘恩又愚昧，既不感谢已得到的，至少也不明智地为将来储备。有人不怜悯寡妇和孤儿，没有将自己的面包和微薄的食物分给穷人，或者更确切地说，分给基督，因为基督在那些即使只得到少许滋养的人身上接受了滋养；一个人或许因意外获得了大量财产——因为这是最不义的——他发现他的许多粮仓不够用了，填满一些又倒空另一些，为了为将来的收成建造更大的粮仓，\BibleRef{路 12:18} 不知道他正被带走，带着未实现的希望，要为他的财富和幻想交帐，被证明是他人财物的坏管家。有人使温良人的道路偏斜，\BibleRef{亚 2:7} 在恶人中使义人偏斜；有人憎恨在城门口责备人的人，\EditorialNote{依撒意亚 29:21} 厌恶说话正直的人；\BibleRef{亚 5:10} 有人向捕获甚多的网献祭，\BibleRef{哈 1:16} 将穷人的掠物藏在家中，\EditorialNote{依撒意亚 3:14} 要么不记念天主，要么记念得不好——说\Quote{愿主受赞美，因为我们富足了}，\BibleRef{匝 11:5} 邪恶地以为他从那位将要惩罚他的那一位那里得到了这些东西。因为因这些事，天主的义怒临到悖逆之子身上。\BibleRef{弗 5:6} 因这些事，天关闭了，或是为了惩罚我们而敞开；更要紧的是，如果我们不悔改，即使受打击，也不亲近祂——祂藉自然的力量亲近我们。\par

19. 对此，我们中那些买卖谷物、观察季节的艰难以图发财、在别人的不幸中纵情享乐、非法获取同胞财产——不像若瑟那样获取埃及人的财产，\BibleRef{创 41:39} 作为宏大政策的一部分（因为他既能适时收集和供应谷物，也能预知饥荒并远先防范）——的人，该说什么呢？因为他们说：\Quote{月朔何时过去，好让我们卖粮？安息日何时过去，好让我们打开粮仓？}\BibleRef{亚 8:5} 他们以各种量器和天平败坏正义，\BibleRef{箴 20:10} 并招致铅的厄法临到自己身上。\BibleRef{匝 5:8} 我们这些贪得无厌、崇拜金银如同古人崇拜巴耳、阿斯塔特和可憎的革摩士的人，\BibleRef{列上 11:33} 该说什么呢？我们看重贵宝石的光辉和柔软飘逸的衣裳——那是飞蛾的猎物、强盗和暴君与窃贼的战利品——我们以众多的奴隶和牲畜为荣，带着自己的所有、赢利和计谋，遍布平原和山地，像撒罗满的蚂蟥一样永不满足，\BibleRef{箴 30:15} 坟墓、大地、火和水也一样；我们为自己的占有寻求另一个世界，抱怨天主的界限太小，不足以容纳我们贪得无厌的欲望。那些坐在高台之上、抬高统治舞台、额头比戏台还高、不顾统御万有的天主和那不可接近的真国度的崇高、不以同样需要仁慈的态度统治如同仆人的属民的人，又如何呢？请你也看看那些躺在象牙床上的人——神圣的亚毛斯恰当地斥责他们——他们用上等香膏抹身，随着乐器歌唱，把易逝之物当作稳固的来依附，却没有为若瑟的苦难忧伤或怜悯；\BibleRef{亚 6:4} 尽管他们本应善待那些先遭灾祸的人，并因仁慈而获得仁慈；正如松树应当哀号，因为香柏树倒下了，\BibleRef{匝 11:2} 并从邻人的惩罚中受教，因别人的灾祸而调整自己的生活，有幸因前人的命运而得救，而不是自己成为他人的警戒。\par

20. 请你与我们一同思考这些问题，你这神圣庄严之人，以你所积累的丰富经验，即那智慧的源泉，这经验是你漫长一生中所收集的。请以此教导你的子民。教训他们：将你的饼分给饥饿的人；召聚无家可归的穷人；遮身裸露，不忽视骨肉之亲。\EditorialNote{依撒意亚 58:7} 如今更特别的是，我们可因匮乏而非丰裕得益，这结果比丰厚的祭品和大量礼物更中悦天主。此后——不，在此之前——我恳求你，今天作一个\Person{梅瑟}（\BibleRef{出 32:11}）或\Person{丕乃哈斯}。站在我们一边，赎罪，使瘟疫止息，无论是藉着属神的祭献（\BibleRef{伯前 2:5}），还是藉着祈祷和合理的转求（\BibleRef{罗 12:1}）。藉着你的调解，抑制上主的忿怒；避免鞭笞后续的击打。他知道尊重一位为子女代求的白发父亲。恳求我们过往的邪恶；为我们的未来作保。呈献一个因痛苦和恐惧而净化的子民。祈求身体的食粮，但更要祈求从天降下的天使之粮。如此行，你必使天主成为我们的天主，必使天心慰藉，必恢复早雨和晚雨（\BibleRef{岳 2:23}）。上主必将显示慈爱，我们的土地必出产果实；我们属地的土地出产暂时的果实，我们这归于尘土的身躯必结出永恒的果实，我们将藉你的手贮存在天上的酒醡中，你将我们和属于我们的在基督耶稣我们主内呈献，愿光荣归于祂，至于永远。阿们。\par

\SourceRecord{Translated by Charles Gordon Browne and James Edward Swallow.；Nicene and Post-Nicene Fathers, Second Series；Vol. 7.；Edited by Philip Schaff and Henry Wace.；Buffalo, NY: Christian Literature Publishing Co.,；1894.；<http://www.newadvent.org/fathers/310216.htm>.}

% structure: p0001 source=h1 target=section reason=page-title

\section{讲道集第十八篇}

\Emph{论其父之逝世}\par

\Emph{本篇讲道发表于公元374年。大圣额我略于该年年初逝世，据《希腊圣教历》记载，为1月1日，尽管克莱芒塞及其他学者将其逝世时间推后数月。其妻圣诺娜比他长寿，并出席聆听了本篇讲道，圣巴西略亦在场，他愿向这位祝圣他为主教的人致敬。这位年近百岁而逝的老圣人生前原属于一个名为\OriginalTerm{海普西斯塔里派}{Hypsistarii}的教派。我们对这一教派的存在及其教义的认识，皆来自本篇讲道及尼撒的圣额我略的寥寥数语（《驳欧诺米》卷一，1615年版，第12页），后者称其为\OriginalTerm{海普西斯塔派}{Hypsistians}。他在其妻圣诺娜的祈祷、影响和榜样下归化，受洗后不久，即被祝圣为纳齐盎主教。他是一位卓越的管理者、虔诚的基督徒、正统的导师、坚定的信仰宣认者、富有同情心的牧者、慈爱的父亲。其妻协助他的生活与工作，其三子格列高利、果尔果尼亚和凯撒留追随其后，他们的名字皆列于圣人名册。}\par

\Emph{在其父葬礼上的讲道辞——圣巴西略在场}\par

1. 天主的人啊，\BibleRef{苏 14:6}忠信的仆人，\BibleRef{户 12:7}天主奥秘的管家，\BibleRef{格前 4:1}有圣神渴望的人——因为圣经以此称呼那些卓越高超、超越可见事物的人。我也要称你为法郎的天主\BibleRef{出 7:1}，以及一切埃及和敌对势力的天主，并称你为教会的柱石和基础\BibleRef{弟前 7:15}，天主的旨意\EditorialNote{依撒意亚 62:4}，世上的光，把握着生命的话语，\BibleRef{斐 2:16}信德的支柱和圣神的安息之所。然而，我为何要一一列举那些你多方面的德行为你赢得并归于你本身的称号呢？\par

2. 但请告诉我，你从何处来，有何贵干，又给我们带来何种恩惠？因我知道你完全因天主并同天主而动，为的是有益于接待你的人。你是来视察我们，还是来寻找牧者，或是来照管羊群？你发现我们已不复存在，大多数已与他一同逝去，无法忍受我们苦难之地，尤其如今我们失去了那熟练的舵手、生命之光。我们仰望他，如同高悬于我们之上的救恩明灯，引导我们的航程。他已带着他所有的杰出才能和牧灵组织的一切力量离去，那是他长期积累的，充满时日与智慧，并如撒罗满所言，以荣耀的白发为冠冕。\BibleRef{箴 16:31}他的羊群荒凉沮丧，如你所见，充满灰心与消沉，不再安息在青翠的牧场上，在安慰的水边被养育，反而寻找悬崖、旷野和坑洞，在那里它将被分散而灭亡。\BibleRef{则 34:14}它对于再获一位智慧牧者已不抱希望，确信再也找不到像他这样的牧者，若能遇到一位相差不远者便心满意足了。\par

3. 如我所说，有三个同等的理由迫使你前来：我们自己、牧者和羊群。那么，请按照你心中服务的精神，给每个理由应得的分，并以公义引导你的言语，使我们比以往更加惊叹你的智慧。你将以何种方式引导它们？首先，以恰当的赞词颂扬他的德行，这不仅是对纯洁者纯粹的墓前祭文，也向他人展示他的言行和榜样，作为真虔诚的标记。然后，赐予我们一些关于生与死、灵魂与身体的结合与分离，以及两个世界——一个是现世而短暂的，另一个是精神感知而永存的——的简要劝言；并说服我们轻视那欺骗性、混乱、不平坦、如同波浪一般将我们托起又抛下的世界，却要紧紧抓住那坚固、稳定、神圣、恒常、远离一切动荡与困惑的世界。因为这会减轻我们对先离世友人的痛苦，不，我们甚应喜乐，如果你的言语将我们带离此地，置于高处，以未来隐藏现在的忧伤，并说服我们，我们也正在走向一位美善的主宰，我们的家乡胜过我们的流徙；并且，移居至那里，对于我们这些在此饱受风浪的人，犹如平静的港湾之于航海者；或者，如同在漫长旅程结束时，安逸与解脱临到那些摆脱了行旅艰辛的人一样，对于达到那边客栈的人，其生存比那些仍在行走这弯曲险峻的人生之路的人更为美好和可忍受。\par

4. 你或许也这样安慰我们；但羊群如何？你肯首先许诺自己如何照管并引导羊群吗？我们是何等乐意安息在你的翼下，渴望你的言语胜过口渴之人渴望最纯净的泉水！其次，请说服我们：那位为羊舍命的善牧\BibleRef{若 10:11}，如今并没有离开我们；他仍在，他牧养、引导，认识自己的羊，也被自己的羊所认识；虽然身体不可见，却为灵性所辨认；他护卫羊群不受豺狼侵害，不许任何人如盗贼般翻墙而入、出卖羊群，用陌生人的声音扭曲并偷窃那本在真理美好引导下的灵魂。是的，我深信，他现在代祷的效力更胜于昔日的教诲，因为既已挣脱肉体的锁链，使灵魂摆脱蒙蔽它的泥土，以赤裸的心智与那最为原始纯粹的赤裸心智相交，他就更亲近天主了；若不算轻率，不妨说，他被提升至天使的地位与信任。凭你的口才与圣神，你必能表述并阐发这些，远胜我所能勾勒。但为免因对亡者之卓越的无知，你的言辞与其功德相去太远，我将就我所知的亡者，简要勾画一幅轮廓，作为颂词的前期规划献给你——你是此类题材的杰出艺匠，由你来补充描绘他德行之美的细节，使之传于众人的耳目与心灵。\par

5. 我将颂词的法则留给他人去描述其乡土、家族、仪表的高贵、外表的辉煌，以及其余人类引以为傲的事物；我从最重要、最切近我们自身的事开始。他出身微贱，原非适宜于虔敬的土壤——我不以他的起源为耻，因我深信其生命的终局；这生命本非栽植于天主的殿宇，而是远离且生疏的，由两大对立之物——希腊的谬误与律法的虚伪——混合而成，其中各有一部分得以逃脱，又有一部分掺杂。一方面，他们拒绝偶像与祭献，却尊崇火与光；另一方面，他们遵守安息日与某些食物的细微条例，却轻视割损。这些卑微的人自称为\OriginalTerm{至尊崇拜者}{Hypsistarii}，宣称全能者是他们唯一崇拜的对象。这种双重不虔的倾向带来什么结果？我不知道更该赞美那召叫他的恩宠，还是他自身的意愿。然而，他如此洁净了蒙蔽心灵之眼的情欲，向着真理飞驰，以至于为了天上的父与真正的产业，甘愿一时失去母亲与财产；他更容易接受这种耻辱，胜于他人接受最大的荣耀。尽管此事实在奇妙，我却不太惊讶。为何？因为这荣耀是他与许多人共有的，所有人都得进入天主的大网，被渔夫的言语所捕获，虽然有人早有人晚，被福音所围住。但在他生命中特别令我惊叹的事，我有必要提及。\par

6. 甚至在他属于我们的羊栈之前，他已属于我们。他的品格使他成为我们中的一员。因为，正如我们自己的许多人并不与我们同在，他们的生活使他们远离共同的身体；同样，许多外面的人却站在我们这边，他们的品格先于他们的信仰，只是缺少那确实拥有的名号而已。我的父亲便是其中之一，一株异族的枝条，却以其生活倾向我们。他在自我节制上如此进步，以至于同时成为最受爱戴和最谦逊的人，这两种品质原难并存。他的正义有何更大更辉煌的见证，甚于他位居国家显要，却未以分文自肥，尽管他眼见众人都伸出百臂巨人\OriginalTerm{布里阿柔斯}{Briareus}之手攫取公帑，因不义之财而肿胀？我称之为不义之财。他的明智与此不无关系，但更详尽的细节将在我的讲辞中给出。我想，正是因为这样的行为，他才获得了信仰。此事如何发生，至关重要，不可忽略，我现在就要陈述。\par

7. 我曾听圣经说：\BibleQuote{谁能找到贤能的妇人？}并宣称这是天主的恩赐，美好的婚姻由主成就。连外邦人也持同样看法；他们说，人所能赢得的最好奖品莫过于贤妻，最坏的莫过于恶妇。但据我所知，没有人比他在这事上更为幸运。因为我认为，若有人从地极和万族中试图促成最好的婚姻，也找不到比这更好、更和谐的了。最卓越的男女如此结合，以至于他们的婚姻是德性的结合，而非肉体的结合；因为他们虽超越所有人，却不能彼此超越，因为在德行上他们完全匹配。\par

8. 那赐给亚当作为相称助手的女人，因人不独居为好，\newline{}\BibleRef{创 2:18} 非但没有成为助手，反而成了仇敌，没有成为同轭者，反而成了对手，借快感诱惑了男人，通过知善恶树使他远离了生命树。但我父亲由天主所得的妻子，不仅（这已属不凡）成了他的助手，甚至成了他的领袖，以言行引导他趋向至高的美德；在其余一切事上，她认为最好按婚姻之律顺从丈夫，但在虔敬之事上，她却不羞于自任教师。她的行为固然可佩，而他欣然顺从则更可佩。这是一位妇女——当他人因自然与人工之美而受赞誉时——她只承认一种美，即灵魂之美，以及尽可能保存或恢复天主的肖像。她摒弃颜料与妆饰，视之为舞台女子的行径。她认为唯一真正的高贵是虔敬，以及认识我们从何而来、往何处去。她认为唯一可靠而不可侵犯的财富，是为天主和穷人，尤其是为不幸的亲属，舍弃自己的财富；而仅提供必要的帮助，她视为不过是提醒而非解除他们的困苦；更为慷慨的施与则带来稳固的荣誉和最圆满的安慰。有些妇女在勤俭持家上出众，另一些在虔敬上出众，而她——虽将二德合一颇为困难——却在两件事上都超越众人，既因在每件事上卓越，又因唯独她将二者结合。她以细心和技巧，按\BibleBook{撒罗满}{箴言}论贤妇的训诫和法则，确保了家业的昌盛，仿佛她不知虔敬为何物；同时，她又如此专务天主和神性之事，仿佛完全摆脱了家务，不让任何一方面妨碍另一方面，反而使二者互相扶持。\par

9. 何时何地的祈祷曾逃过她？她每日以祈祷为先于一切之事；或者更确切地说，谁曾如她那样满怀希望，立即得到所求的回应？谁曾如此敬重司铎的手和面容？谁曾如此尊崇各种修道生活？谁曾以更持久的斋戒和守夜克苦肉身？谁曾在彻夜和每日的圣咏中站立如柱？谁对贞洁有更大的爱慕——虽然她自己容忍了婚姻的束缚？谁曾是孤儿和寡妇更好的保护人？谁曾如此帮助减轻哀悼者的不幸？这些事虽微不足道，也许在有些人看来不屑一顾，因为大多数人难以企及（凡难以企及之事，由于嫉妒，甚至被认为不可信），但在我的眼中却最为可敬，因为它们是她信德的发现和心灵热忱的业绩。同样，在圣会或圣地，除了礼仪中必要的应答之外，从未听到她的声音。\par

10. 既然祭台上从未举起铁器，\newline{}\BibleRef{申 27:5} 也看不见听不见任何凿子，乃是大事，那么，既然一切奉献给天主的都应当自然无华，她以缄默敬重圣所，也必是大事；她从未背对可敬的祭台，从未向神圣的地面吐痰；她从未握任何外邦女子的手，或亲吻她们的嘴唇，尽管她们在其他方面如何可敬，或与她血缘多近；她绝不——不说是自愿，即令被迫——与那来自亵渎不洁筵席的人同席共餐；她不能忍受违反良心的律法，经过或观看受玷污的房屋；她不能容许曾领受并宣讲神圣之事的耳朵或舌头，被希腊故事或戏剧歌曲玷污，因为不洁之事不合于神圣之事；更可奇者，她虽然因陌生人的不幸也深受感动，却从未屈服于外在的哀恸表现，在圣体圣事之前发出悲声，或从神秘被封印的眼睛滴下眼泪，或在庆节上留下任何哀悼的痕迹，尽管她自己的痛苦何等频繁；因为爱天主之灵魂应将每个人的经验置于天主之事以下。\par

11. 我缄默不言那更为不可言说的事，天主为此作证，那些忠信的女仆，母亲曾向她们吐露了这些事。或许关于我自身的事不值得提及，因为我已证明自己不配承受那对我的希望；但母亲在我出生前就把我许诺给天主，毫不畏惧未来，且在我出生后立即把我奉献，这在她方面是一桩伟大的事业。因天主的仁慈，她的祈祷并未完全落空，那吉祥的祭献也未被拒绝。其中一些事已经存在，另一些则在未来，通过逐渐的添加而成长。正如那最令人愉悦地洒下晨光的太阳，到正午时变得更为炽热和灿烂，同样，她从起初就表现出不小的虔诚证据，最终以更圆满的光辉照耀。那时，那将她安置在自己家中的人，在家中就有了一个不小虔诚的激励；她因自己的出身和世系，拥有对天主和基督的爱，并以德性为遗产；不像他（丈夫）那样，是从野橄榄上砍下而接在好橄榄上，却因信德的丰盛，不能忍受不相称的轭；因为，尽管她在忍耐和刚毅上超越众人，但她不能容忍这一点：即因那与她成为一体的人（丈夫）的疏离，而半心半意地结合于天主，以及在身体的结合之外，未能加上精神上的亲密联系；为此，她日夜俯伏在天主前，以许多禁食和眼泪恳求她头部（丈夫）的得救，且勤勉地致力于她的丈夫，通过责备、劝诫、关注、疏远，以及首要的是她自身对虔诚的热忱，以此多方影响他；灵魂尤其被此说服而软化，并甘愿顺从德性的压力。不断滴落在岩石上的水滴终将穿透它，并最终达成其渴望，正如后续所示。\par

12. 这些是她祈祷和希望的对象，出于信德的热忱而非青春的热情。事实上，没有人像她那样，因经验到天主的慷慨，而对所希望之事如此确信，如同对现世之事一样。为了我父亲的得救，他理性的逐渐说服与天主常赐予配得救恩之灵魂的梦境一同发生。那梦境是什么？这是这个故事中最令我愉悦的部分。他以为自己在歌唱，如同他从未做过的那样，尽管他的妻子频繁祈祷，这首歌正是出自圣王大卫的\WorkTitle{圣咏}中的句子：\BibleQuote{我真喜欢，因为有人向我说：我们要进入上主的圣殿。} 这首圣咏对他来说是陌生的，但随着话语，渴望临到了他。她一听到这梦，知道自己已获允所求，便抓住机会，回答说，若伴有实现，这梦将带来最大的喜乐，并以她的喜乐显示恩惠的伟大，她催促推动他的得救，以免有任何事介入阻碍这召唤，消散她渴望的对象。恰在此时，许多主教正赶往\Place{尼西亚}，以对抗亚略的疯狂，因为分裂神性的邪恶刚刚兴起；于是我的父亲将自己交付于天主和真理的宣讲者，承认了他的愿望，并向他们请求共同的救恩，其中一位是著名的\Person{莱昂提乌斯}，当时是我们自己的都主教。若我缄默不言那时由恩宠所赐予他的奇事，那将是恩宠的大不敬。这奇事的见证人不少。精确的教师们灵性上犯了错误，而恩宠是对未来的预兆，司铎职的仪式与望教者的接纳混在一起。啊，非自愿的入门礼！他屈膝跪下，接受了望教者地位的接纳形式，以至于许多人，不仅是最有智慧的，甚至是最低智力的，都预言了未来，由毫不模糊的迹象确知将要发生的事。\par

13. 短短间隔后，奇事接踵而至。我将这描述托付予信友的耳朵，因为对世俗之心而言，没有好事是可信的。他正接近那藉水和圣神而生的重生，藉此我们向天主承认基督模样之人的形成与完成，以及从属土的到属神的转变与重塑。他带着热烈的渴望和明亮的希望接近洗礼池，在尽一切可能的净化之后，且灵魂和身体的净化远超那些将要接受\Person{梅瑟}法版的人。他们的净化仅限于衣服、轻微的节食和暂时的节制。他过去的一生都为这光照作了准备，是一种确保恩赐的预先净化，为使完美能交付于纯洁，并使祝福在一颗确信拥有恩宠的灵魂中不冒风险。当他从水中上来时，一道光芒和荣耀在他周围闪耀，配得上他接近信德恩赐的态度；这甚至显明给一些其他人，他们当时隐藏了这奇事，因为害怕说出一个每个人都以为只有自己看见的景象，但不久后便彼此沟通了。然而，对施洗者和入教引导者来说，这如此清晰可见，以至于他甚至不能隐藏这奥秘，而是公开呼喊说，他正在以圣神傅油自己的继任者。\par

14. 确实，凡听见并知道梅瑟的人，都不会不信这事：梅瑟在世人的眼中微小，且尚微不足道时，被那焚烧而不熄灭的荆棘，或者更确切地说，被那在荆棘中显现的那一位所召唤，\newline{}\BibleRef{出 3:4} 并由那第一个奇事所鼓励。我说的是梅瑟，对谁而言，海水分开，玛纳降下，磐石涌出泉水，火柱和云柱轮流引导，他举手的姿势赢得了胜利，十字架的象征战胜了万千。\newline{}此外，依撒意亚曾看见色辣芬的荣耀，其后是耶肋米亚，受托拥有对抗万民和君王的伟大权能；\newline{}\BibleRef{耶 1:10} 前者听见了神圣的声音，并用一块火炭被洁净，以担任先知职；后者在他成形之前已被认识，在他出生之前已被圣化。\newline{}还有保禄，曾是一名迫害者，后来却成了真理的伟大宣讲者和外邦人在信德中的导师，被一道光环绕，\newline{}\BibleRef{宗 9:3} 并认出了他所迫害的那一位，受托伟大的职务，用福音充满了每一个耳朵和心灵。\par

15. 何以需要一一列举所有那些被天主召叫到自己身边、并与如此奇事相联系，从而坚定其虔诚的人呢？而且，在经历了如此难以置信、令人惊愕的开始之后，他后来的行为并没有使先前的事蒙羞——不像那些很快就对善事感到厌倦的人，他们忽视了一切进一步的长进，甚至完全堕落到恶习中。对他而言，不能这样说，因为他与自己以及早期的日子最为一致，使他在司铎职之前的生活与其卓越性相协调，而司铎职之后的生活又与之前的相协调，因为以一种方式开始而以另一种方式结束，或者朝向一个与他最初所追求不同的目标前进，是不相称的。\newline{}随后，他受托担任司铎职，并不像今日那样轻易和混乱，而是在一段短暂的间隔之后，以便在他自己的洁净之外，获得洁净他人的技巧和能力；因为这是属神秩序的法则。当他受托此职时，恩宠更得光荣，因为它实在是天主的恩宠，而非人的恩宠，并且，正如宣讲者所说，并非一种独立的冲动和心志。\par

16. 他接受了一座林间而质朴的教会，其牧职和照管并非来自远方，而是由他的一位前任——一位可钦慕的、具有天使般品格的、比我们现今的百姓统治者更纯朴的人——照管过；但在他被迅速接回天主身边后，由于失去了领袖，这教会大部分时间变得疏忽而荒废；因此，他起初以不严厉的方式，用牧灵知识的言语，并让自己作为榜样，如同一个属神的雕像，被磨光以彰显一切优秀行为之美，来软化百姓的习性。\newline{}随后，通过不断默想圣言（尽管他在这方面起步较晚），他在短时间内积累了如此多的智慧，以至于丝毫不逊于那些为此付出极大辛劳的人，并获得了天主的特殊恩宠，成为正统信仰的父亲和导师——不像我们现代的智者，顺应时代的精神，或用模糊而诡辩的语言为我们的信仰辩护，仿佛他们没有稳固的信仰，或者是在掺杂真理；但他在虔诚方面超越了那些拥有修辞能力的人，在修辞技巧上又超越了那些心地正直的人；或者更确切地说，作为演说家他位居第二，却在虔诚方面超越众人。\newline{}他承认一位在天主圣三中受敬拜的天主，以及三位在一位天主性中联合；既不因一体而陷入撒伯里乌异端，也不因三位而陷入亚略异端；既不通过收缩而无神论地消灭天主性，也不通过不平等本性或大小的区分将其撕裂。因为，既然其一切属性都是不可理解、超越我们理智能力的，我们又怎能在这类主题上感知或用定义表达那超越我们认知的事物呢？不可测量的怎能被测量，天主性怎能被归结为有限事物的状态，并根据大小程度来衡量呢？\par

17. 关于这位伟大的天主之人、真正的神学家，在这些主题上受圣神感动，我们还能说什么呢？只能说，通过他对这些要点的领悟，他就像伟大的诺厄——这第二个世界的父亲——使这教会被称为新耶路撒冷，成为在水上承载的第二只方舟；因为它既超越了灵魂的洪流和异端者的侮辱，又在声誉上超越所有其他教会，正如在人数上不及它们；并且它有着与神圣的白冷相同的命运——白冷无疑可以被说成是一座小城和世界的首都，因为它是基督的乳母和母亲，基督既创造了世界，又战胜了世界。\par

18. 为证明我所说的话：当教会中过分热心的那一方对我们激起骚乱，我们被一份文件和狡猾的用语诱骗而参与了邪恶时，唯独他被认为拥有未受损伤的心智和未被墨水玷污的灵魂，即使他因纯朴而受骗，并因其灵魂的坦诚而未能提防诡诈。\newline{}唯独是他，或者更确切地说，他率先以其对虔诚的热忱，使反对我们的派系——那最后离开我们、却最先回归的——与他自己和教会的其余部分和好；这是由于他们对他的尊敬以及他教义的纯洁，以至于教会中严重的风暴得以平息，飓风因他的祈祷和规劝而减弱为微风；而我，若可自夸一句，则是他在虔诚和行动上的同伴，在一切为善的努力上协助他，陪伴他、与他并肩奔跑，并在此次获准分担很大一部分辛劳。\newline{}关于这些事的叙述，在这里必须结束，虽然有些过早。\par

19. 谁能数尽他美德的全部，或若想跳过大部分，不难发现哪些可以省略？因为每一优点，当它浮现于脑海时，似乎都胜过先前的；它占据了我，我更不知所措，不知该跳过什么，而其他颂词作者却不知该说什么。因此，材料的丰富在某种程度上成了我的障碍，我的心灵本身在试图检验他的品质时也受到了考验，而在所有品质都相当时，无法找出一个胜过其余。因此，正如我们看到一块石子落入静水中，它成为一圈圈涟漪的中心和起点，每一圈连续扰动打破了外圈；这正是我的情况。因为当一件事进入我的脑海，另一件随之而来并取代它；我在选择中疲惫不堪，因为我已经抓住的总是让位于后续的。\par

20. 谁比他更关心公益？谁更善于理家？因为天主以巧妙变化安排万物，赐给了他一个家庭和相称的财富。谁对穷人——这一最被轻视却应享有同等尊荣的本性部分——更有同情之心，更慷慨之手？因为他实际上把自己的财产当作别人的，自己不过是管家，尽力救济贫困，不仅花费多余的，甚至花费必需的——这是爱穷人的明证，不仅按照撒罗满的训诫分给\BibleQuote{七分}（见：\BibleRef{训 11:2}），而且若有第八个前来，也不吝啬，而是比我们所知他人获取财富更乐于施舍；他除去轭和拣选（我想这指施舍时测试受者是否值得的一切卑劣行为）以及抱怨之言（\EditorialNote{依撒意亚 58:9}）于善行中。这是大多数人所做的：他们确实给予，但没有那种乐意，而乐意比单纯的给予更伟大、更完美。因为他认为，为了配得之人的缘故而慷慨对待不配之人，远胜于因害怕不配之人而剥夺配得之人的机会。这似乎是\BibleQuote{将我们的粮食撒在水面}（见：\BibleRef{训 11:1}）的职责，因为它在公义审判者的眼中不会被冲走或消失，而是会到达那里，我们一切所有都储存在那里，并在适当的时候与我们相遇，即使我们并不这样想。\par

21. 但其中最美好、最伟大的是，他的宽宏大量伴随着不求名利。其程度与特征我将说明。在视财富为共有、慷慨施与上，他与妻子竞相追求美德；但他将大部分施舍托付于她之手，认为她是此类事务最优越、最可靠的管家。她是怎样的女子啊！即使是大西洋，或更大的海洋，也无法满足她的索求。她的慷慨之爱如此博大无尽。相反，她与撒罗满所说的\OriginalTerm{蚂蟥}{horse-leech}（见：\BibleRef{箴 30:15}）相竞争，因对进步有不可遏止的渴望，克服了倒退的倾向，无法满足其行善的热忱。她不仅认为他们原有和后来所得的一切财产都不足以满足她的渴望，而且如我常听她所言，若有可能，她会欣然出卖自己和子女为奴，将所得用于穷人。她如此完全放纵自己的慷慨。我想，这比任何实例都更具说服力。在金钱上的宽宏大量在他人身上不难找到，无论是在国家的公共竞争中浪费，还是通过穷人借贷给天主（这是为施舍者存积的唯一方式）；但很难发现有人放弃了随之而来的名声。因为正是对名声的渴望促使大多数人乐于花费。而在施舍必须隐秘之处，施舍的倾向就不那么强烈了。\par

22. 他的手如此慷慨——其余细节我留给认识他的人，因此若在我身上也见证了类似行为，那也出自同一泉源，是那水流的一部分。谁更受天主引导而接纳人进入圣所，或愤慨于对圣所的侮辱，或怀着敬畏清除祭台上的不洁之物？谁以如此不偏不倚的审判和公义的天平，或裁决诉讼，或憎恨恶行，或尊崇美德，或擢升最优秀者？谁对罪人如此怜悯，对奔跑得好的人如此同情？谁更懂得使用棍杖与牧杖的时机，却最依赖牧杖？谁的眼睛更关注国内的信徒，尤其是那些通过隐修和独身生活而藐视尘世及尘世之物的人？\par

23. 有谁更能斥责骄傲、培养谦卑呢？而且不是以虚伪或外表的方式，如同现今许多装作有德之人，他们外表华丽如同最无知的女子，因为缺乏自身的美貌便求助于脂粉，可谓装扮得光彩夺目，其实美貌却不美，比原来更丑。因为他的谦卑不在乎衣着，而在乎心灵的修养；也不是靠弯下脖子、压低声音、垂目、长须、剃光的头或矜持的步伐来表现，这些可以一时装出，但很快就会被揭穿，因为一切做作之物都不能持久。不！他在生活中总是最高尚，心中却最谦卑；在德行上不可接近，在交往中最平易近人。他的衣着毫无奇特之处，既避开了华丽也避开了褴褛，而内在的光华却超越一切。对于肚腹的贪病，他比任何人都加以节制，却不张扬，以免因追求名声而滋生新的恶习，而又能保持谦卑而不自傲。因为他认为，为赢得外人的赞誉而说和做一切，是政治家的特点，其最高幸福在于现世生活；而属灵的人和基督徒只应关注一件事，就是自己的救恩，凡有助于此的便加以重视，凡无助于此的便以无价值而唾弃；因此轻视可见之物，只专注于内在的成全，并且最重视那些能促进自己进步、并通过自己吸引他人归向至善的事物。\par

24. 但他最杰出、最独特、却最不为人所知的是他的纯朴、无伪与无怨。因为无论在古时或当今的人中，每个人都据认为各有某种特殊的成就，正如每人恰从天主领受了某种特定的德行：约伯在不幸中有着不屈的忍耐（\BibleRef{约 1:21}），梅瑟（\BibleRef{户 12:3}）与达味有着温良，撒慕尔有预言和预见未来的能力（\BibleRef{撒上 9:9}），丕乃哈斯有热忱（\BibleRef{户 35:7}），为此他得享盛名，伯多禄和保禄有传道的热忱（\BibleRef{迦 2:7}），载伯德的儿子们有豪言壮语，因此他们也被称为\Quote{雷霆之子}（\BibleRef{谷 3:17}）。但我为何要在向了解此事的人说话时一一列举呢？斯德望和我父亲特别突出的标记是无恶意。因为斯德望即使在危险中也不恨他的攻击者，而是被石头砸死时还为砸他的人祈祷（\BibleRef{宗 7:59}），作为基督的门徒，他堪为基督受苦，因此，在他的忍耐中，为天主结出了比死亡更美的果实：而我父亲则在受攻击与宽恕之间不留任何间隔，以至于由于宽恕的迅速，他几乎连痛苦本身都被剥夺了。\par

25. 我们既相信也听说过天主愤怒的渣滓，即祂对待罪有应得之人的剩余：\Quote{因为主是报应的天主。}虽然祂本性仁慈，倾向于温和而非严厉，但祂并不完全赦免罪人，以免他们因祂的恩慈而变得更坏。然而我父亲对激怒他的人从不怀恨，他甚至完全不受愤怒的影响，尽管在属灵的事上他极富热忱；除非他已准备好、武装起来，以敌对阵势对抗那些前来伤害他的事物。因此，他的这种温和性情，正如俗语所说，即使成千上万的人也不会激怒他。因为他所有的愤怒不像蛇的毒液，在内部潜伏，随时自卫，急于爆发，一受惊扰就渴望立刻反击；而是像蜜蜂的螫刺，一击并不致命；他的仁慈超越人性。轮子和鞭笞常被威胁，而施用它们的人就在近旁；但危险最终不过是被人揪耳朵、拍脸颊、或打太阳穴：他就是这样缓和了威胁。他的衣服和凉鞋被扯掉，恶棍被打倒在地：那时他的愤怒不是针对攻击他的人，而是针对那个急于救助他的人，视之为邪恶的帮凶。还有什么能更有力地证明他良善，配得向天主献上祭品？因为常常，他自己非但不被激怒，反而为攻击他的人辩护，为他们的过犯而脸红，仿佛那些过犯是他自己的。\par

26. 露水更容易阻挡清晨的日光，也不能在他内残留一丝怒气；他一开口，愤怒便随言语而去，只留下对善的爱，且从不过夜；他也不怀怒——这怒甚至摧毁明智之人——也不在体内显出任何恶的迹象；反而，即使被激怒，他也保持平静。这结果极不寻常：不单是他独施斥责，而是他成为唯一一个既被受责备者爱戴又钦佩的人，因他的良善战胜了情感的激动；人们感到，被义人惩罚比被恶人玷污更有益，因为前者的严厉因其效用而变得愉快，后者的和善却因品行之恶而令人怀疑。但他的灵魂与品性如此单纯神圣，他的虔诚仍使傲慢者敬畏；或者说，他们敬畏的原因正是他们所蔑视的单纯。因为在他，无论祈愿或诅咒，都不可能不立即带来持久的祝福或短暂的痛苦。前者发自他灵魂深处，后者仅停留在唇边，作为慈父般的训诫。许多曾伤害他的人，并未遭受迟来的报应，或如诗人所言，\Quote{紧随人后的报复}；而是在他们情绪发作的当下，便被击打而悔改，前来跪在他面前，获得宽恕，荣耀地败退，并因惩罚与宽恕而改过。的确，宽恕之心常有巨大的拯救之力，以羞愧感制止作恶者，使他从恐惧转回爱——一种更安稳的心态。在惩罚中，有些人被负轭的牛掀翻——这些牛突然攻击他们，虽从未有类似行为；另一些人被最温顺驯良的马抛下踩踏；还有些人被难以忍受的热病及其恶行之幻象所攫；其余人则以不同方式受罚，并从所受之苦学会顺从。\par

27. 他的温和如此非凡，那么他在勤劳与践行美德方面是否逊于他人？绝非如此。他虽温和，却有着与温和相应的活力——若有人有的话。因为通常，仁慈与严厉这两种美德互相矛盾对立：前者温和却缺乏实践能力，后者有实践能力却不带同情；但在他身上，两者奇妙地结合：他的行动如严厉之人一样有力，却伴以温和；他的宽容看似不切实际，却带有活力——在他的庇护、言谈自由及各种职务中皆然。他结合了蛇的明智（关乎恶）与鸽的纯朴（关乎善），既不使明智沦为诡诈，也不使纯朴转为愚钝；而是尽己所能，将两者合为一种完美的美德形式。既然他有如此出生，如此司铎职的实践，如此在所有人心中的名声，那么他被视为堪当天主藉以确立真宗教的奇迹，又何足为奇呢？\par

28. 关于他的奇事之一，是他遭受疾病与身体的痛苦。但圣人也受困扰，这有何奇？或为净化他们虽轻微却属泥土的部分，或为磨炼美德并考验哲学，或为教育软弱者——他们从榜样学会忍耐而不屈服于厄运。他病了，时值神圣而显赫的逾越节，诸日之皇后，驱散罪恶黑暗的光明之夜——我们在这充足的光明中庆祝我们的救恩，与曾为我们被处死的光明同死，并与复活者一同复活。这正是他受苦的时刻。痛苦如何，我将简要说明。他全身因极高烧而灼热，力量衰竭，无法进食，睡眠消失，极度痛苦，心悸不安。口中上颚及整个上表面严重溃疡，疼痛异常，连吞咽水都困难危险。医术、朋友们最恳切的祈祷、以及一切可能的照顾，皆无效用。他自己身处此绝望境地，呼吸短促急促，对现世毫无知觉，完全超脱，沉浸于他长久渴望、如今已备妥的对象之中。我们则在圣殿中，将恳求与神圣礼仪相结合——因对其他一切绝望，我们投奔大医生、那夜的能力，以及最后的援助——我想说，我们是在守节，还是举哀？是庆贺，还是为一位已不在者行丧礼？啊，那些眼泪！那时全民众所流的眼泪。啊，声音、呼喊、与圣咏交织的赞美诗！从圣殿他们寻找司铎，从神圣礼仪寻找主礼，从天主那里寻找他们应得的牧者——由我的米黎盎带领他们，击鼓\BibleRef{出 15:20}，不是胜利之鼓，而是恳求之鼓；那时她才初次学会因不幸而羞愧，同时呼唤民众与天主：呼吁前者同情她的痛苦，慷慨流泪；呼吁后者垂听她的祈求——她以受苦的发明才能，在祂面前陈述祂古时的一切奇事。\par

29. 那么，那夜和病人的天主如何回应？我叙述此事时不禁颤栗。听者啊，你们尽管颤栗，但不要不信：因为这是不虔敬的，当我作为讲述者，且涉及他时。奥迹的时刻已到，庄严的站姿和秩序来临，此时为举行圣礼而保持静默；然后他被那使死者复活的天主和圣夜扶起。起初他微微动弹，然后更明显；接着他以微弱而含糊的声音呼唤了一名侍候他的仆人，叫他过来，带上他的衣服，用手扶他。仆人惊恐地过来，欣然服侍他；他则靠着仆人的手如同倚杖，效法山上的梅瑟，安排他虚弱的手祈祷，并与他的子民一起，或代表他们，热切地举行奥迹，用他体力允许的简短话语，但在我看来，怀着最完美的意向。何等奇迹！在没有圣所的神圣处所，没有祭台而献祭，远离圣礼的司祭：然而这一切都借圣神的力量临在于他，为他所认识，虽为在场者所不见。然后，他加上惯常的感恩经文，祝福百姓后，又回到床上，略进饮食，睡了片刻，恢复了精神；身体逐渐康复后，在节庆的新日，即我们称为复活节后第一个主日，他进入圣殿，在全体圣职人员伴同下，开始了保存下来的生命，奉献了感恩祭。在我看来，这与希则克雅的事迹相比毫不逊色，他在病中和祈祷中获得天主借延长寿命而荣耀他，这由日晷影子后退的标记所预示，\EditorialNote{依撒意亚 38:8}，这是按那位康复君王的请求，天主同时以恩宠和标记尊荣他，通过延长白日来保证他日子的延长。\par

30. 不久之后，同样的奇迹发生在我母亲身上。我认为也不宜忽略此事：因为我们将对她给予应得的荣誉——若对任何人而言有此必要——并且因她与他一同被纳入我们的叙述而使他喜悦。她一向强壮有力，终生无病，如今自己却受疾病袭击。由于极大的痛苦——为免冗长——尤其是无法进食，她的生命多日垂危，且找不到任何药物可治。天主如何扶助她？不是像为古以色列人那样降下玛纳，也不是劈开岩石给渴慕的人民饮水，或像厄里亚那样借乌鸦喂养她，\BibleRef{列上 17:6}，或像为饥饿的达尼尔在狮穴中遣先知腾空运送食物。而是如何？她梦见自己所钟爱的我——即使在梦中她也没有偏爱我们中的任何一人——在夜里突然来到她跟前，带着一篮纯白的饼，我照常祝福画十字，然后喂她，使她强壮，她便更强健了。夜间的异象是真实的行为。因为，结果她变得更有活力，更有盼望，由清晰明确的标记显明。次日清晨，我早早探望她，立刻看见她气色较好；我照常问她夜里过得如何，是否需要什么，她回答说：\Quote{我的孩子，你最乐意且慈祥地喂了我，现在却问我怎样。我很好，很安适。}她的婢女也向我示意不要抗拒，立即接受她的回答，以免真相暴露使她重陷沮丧。我再补充一件他们两人共享的事例。\par

31. 我曾从亚历山大港启航前往希腊，穿越帕忒尼安海。这趟航行实在不合季节，乘的是埃吉纳船，出于热切渴望的冲动；因为特别促使我启程的是，我遇到了一班熟悉的船员。航行了一段路程后，一场可怕的暴风雨降临，据同船者说，这样大的风暴他们很少见过。我们都害怕共同的死亡，但我最怕的是灵魂的死亡；因为我面临在悲惨中离世的危险，尚未受洗，我在死亡之水中渴望圣神之水。为此我呼喊、恳求、祈求稍得宽限。我的同船者，即使处于共同的危险中，也加入我的呼求，甚至我的亲属也不会如此，他们真是善良的灵魂，从危险中学到了同情。在这种状况下，我的父母为我感到心痛，我的危险通过夜间的异象传给了他们，他们从陆地上帮助我，以祈祷平息波涛——这是我后来登陆后推算时间得知的。这也在一次健康的睡眠中显示给我，当时在风暴稍稍平息时我经历了这睡眠。我似乎抓着一个复仇女神，面目可怖，预示着危险；因为黑夜清晰地把她展现在我眼前。我另一位同船者，一个极其友善且亲爱的青年，非常担忧我当时的景况，他梦见我的母亲在海面上行走，毫不费力地抓住船并拖向陆地。我们对这异象满怀信心，因为大海开始平静，不久我们就在没有太大艰辛的情况下抵达了罗得岛。我们自身因那次危险而成了祭品；因为我们向天主许诺，若得救，就将自己献给祂，得救后，我们果然把自己献给了祂。\par

32. 这就是他们共同的经历。但我以为，那些对他生平有确切了解的人，必已长久奇怪我们为何详述这些事，仿佛我们认为这是他成名的唯一理由，而推迟提及他那个时代的困难，他曾卓然抗击这些困难，好像我们对此一无所知，或认为无关紧要。现在，我们就来谈谈这个话题。我们时代的第一件、也是我认为最后一件坏事，就是那位背弃天主、背弃理智的皇帝，他认为征服波斯是小事，而使基督徒臣服于己才是大事；因此，他与引导并掌控他的魔鬼一起，穷尽各种不敬神的方式，通过劝说、恐吓和诡辩，竭力吸引人们归附他，甚至在其种种诡计之外还使用了武力。然而，他的阴谋暴露了，无论他试图用诡辩的伎俩掩盖迫害，还是公然动用权威——也就是这样或那样，藉欺骗或暴力，将我们置于其控制之下。谁能找到更彻底地藐视或击败他的人呢？在他对圣堂的轻蔑中，有许多迹象，其中之一就是派遣警察及其长官前去，意图自愿或强行占有它们：他曾攻击许多其他地方，同样怀着这样的意图来到这里，要求根据皇帝的谕令交出圣殿，但他远未达成任何愿望，要不是他很快就在我的父亲面前退让，或出于他自己的明智，或听了别人的劝告，他就不得不在脚部残废的情况下撤退，这位司铎为了保护他的圣地，是那样愤怒和热忱。那么，谁在导致他的灭亡中显然有更大的份呢？无论是在公开场合，通过他针对那可诅咒者所发的祈祷和联合恳求，不顾当时的[危险]；还是在私下，调集他夜间的装备：睡在地上，使他自己年老柔弱的身体日益衰弱；还有泪水，他几乎整整一年用泪泉浇灌土地，将这些实践只指向那洞察人心者，同时他试图避开我们的注意，怀着我所说过的那种隐藏的热心。他本会完全不为人知，要不是我有一次突然冲进他的房间，注意到他躺卧在地上的痕迹，便询问他的侍从这意味着什么，从而得知了夜晚的奥秘。\par

33. 这是同一时期、同一勇气的另一件事。\Place{凯撒勒雅}城因选举主教而骚动；因为一位刚刚离世，必须另选一位，而激烈的党派之争不易平息。因为这座城市自然容易产生党派精神，由于其信德的热情，而教座的显赫地位更增加了竞争。情况如此：几位主教已到来祝圣主教；民众分裂为几派，各有自己的候选人，正如这类情况中常见的那样，出于私人友谊或对天主的虔敬的影响；但最后全体民众达成一致，并借助当时驻扎在那里的一队士兵，抓住了他们的一位主要公民，一位生活卓越、但尚未以圣洗盖印的人，违背他的意愿将他带到圣所，把他放在主教们面前，恳求——恳求中夹杂着暴力——使他被祝圣并被宣布为主教，不是以最好的秩序，而是以全部真诚和热忱。而且，时间指出了比他更显赫和更虔诚的人是谁，这不可能说。那么，这场骚动导致了什么结果？他们的抵抗被克服了，他们给他付了洗，他们宣布了他，他们使他就座，这都是外在行动，而非属灵的判断和意向，正如后续所显示的那样。他们很高兴退下并恢复判断的自由，并商定了一个计划——我不知道这是否出自圣神的感动——认为所行的一切无效，任命无效，理由是暴力出自那位同样遭受暴力的人，并且抓住了一些当时所说的话，这些话过于激烈而非明智。但是，那位大司祭、行为的公正审查者，没有被他们的计划所动摇，不赞同他们的判断，而是保持不受影响、不动摇，仿佛完全没有受到压力一样。因为他看到，既然暴力是共同的，如果他们对他提出指控，他们自己就面临反指控；或者，如果他们宣告他无罪，他们自己或许也可被宣告无罪；但更公正地说，他们即使宣告他无罪，也无法确保自己的无罪：因为如果他们值得原谅，那么他当然也值得；如果他不值得，那他们就更不值得了：因为那时冒险抵抗到底，远比事后策划反对他更好，尤其是在这样的时刻，消除现有敌意比制造新敌意更有利。情况如下。\par

34. 皇帝曾来势汹汹地反对基督徒；他对选举感到愤怒，并威胁被选者，而城市正处于迫在眉睫的危险之中，不知在此日之后它是否应当不复存在，还是得以逃脱并受到某种程度的宽待。关于选举的革新是一个新的激怒原因，除了在繁荣时期摧毁了命运女神庙之外，这被视为对其权利的侵犯。省长也急于将这个机会据为己用，并对新任主教怀有敌意，由于他们政治观点不同，他从未与主教建立友好关系。因此，他致信召集祝圣者，要他们宣布选举无效，并且言辞并不温和，因为他们受到如同皇帝命令般的威胁。于是，当信函到达他手中时，他毫无畏惧或迟疑地回复——请看他回信的勇气和精神——\Quote{最优秀的省长，我们的一切行动都有一位监察者，一位皇帝，他的敌人正在武装反对他。祂将审查当前的祝圣，这是我们按照祂的旨意合法举行的。至于任何其他事情，你若愿意，可以极其轻易地对我们使用暴力。但没有人能阻止我们在此事上维护我们行动的合法性和正义性；除非你在这方面制定一条法律，而你却无权干预我们的事务。}这封信令收信人赞赏，尽管他一时为之恼怒，正如许多深知事实的人告诉我们的那样。它也阻止了皇帝的行动，使城市脱离了危险，并且使我们——不妨加上——免于耻辱。这是由一位微不足道且属于辅助教区的主教所成就的。难道这种领导权不是远胜于来自上级教区的头衔，以及一种基于行动而非名义的权力吗？\par

35. 有谁如此远离我们这个世界，以至于不知道那在顺序上最后、却是他能力最初和最伟大的证明？同一城市又为同一原因而骚动，因为那位以如此光荣的暴力选出的主教突然离世，他已归赴天主——为他曾在迫害中勇敢而高贵地奋战过的天主。骚动的热度与它的不合理成正比。那位杰出的人物并非不为人知，而是比众星间的太阳更加显眼，不仅在众人眼中，特别是在那批精选且最纯洁的子民眼中——那些在圣所事奉的人，以及我们中的纳齐尔人，这些任命应当，即使不完全是，也应尽可能地属于他们，这样教会就不会受到伤害，而不是属于最富有、最有权势的，或最暴力、最不合理的民众，尤其是其中最腐败的人。事实上，我几乎倾向于相信世俗政府比我们这被赋予神圣恩宠的教会更有秩序，而且这类事情因恐惧而比因理性管理得更好。因为哪个有理智的人，会忽视你那神圣而被主之手美化了的、无婚配的、无财产的、几乎无血肉的、言语上仅次于圣言本身的人，而去找另一个人？你是哲学家中的智者，世俗者中超越世界的人，我的同伴和同工，并且更大胆地说，与我共有同一灵魂的人，分享我生命和教养的人。但愿我能自由地向你说话，并在别人面前描述你，而不必因你的在场而不得不略过大部分话题，以免被怀疑为阿谀奉承。但正如我开始所说，圣神必然认识祂自己的人；然而他却成为嫉妒的目标，出自那些我羞于提及的人，我希望别人不会听到他们的名字，那些人刻意嘲讽我们的事务。让我们像河流中流的一块岩石那样掠过它，并以恭敬的沉默对待那应当被遗忘的话题，继续我们其余的议题。\par

36. 属圣神的人确切知晓圣神之事，他认为自己不可采取顺从的立场，也不可偏袒那些倚靠人的恩惠而非天主的派系与偏见，而必须以教会的利益和共同的救恩为唯一目标。因此，他写信、建言、努力团结百姓与圣职人员（无论是否在圣所供职），即使不在场也作证、裁决并投票，且因白发之威，对陌生人如同对己羊群一般行使权柄。最后，由于祝圣须合乎教规，且欠缺法定数量的主教来宣布，他便拖着年迈多病的身躯从卧榻上起身，毅然前往城中——更确切地说，是被抬去的，身体虽已奄奄一息，却深信若有何不测，这份热忱将是一块高贵的裹尸布。此时，又一次出现了并非不可信的奇事：他从劳苦中获得力量，从热忱中获得新生，主持礼仪，置身于争战之中，为那主教加冕，然后被护送回家——不再是躺在棺架上，而是在一艘神圣的方舟中。他的忍耐，我已对其赞美流连，在此事上更得以彰显。因为同僚们因战败的羞辱和这老人在公众中的影响力而恼怒，并以谩骂发泄；但他的坚韧如此强大，甚至超越于此，以谦逊为最强盟友，拒绝与他们互骂。因为他觉得，真正获胜之后反被口舌击败，实为可怕之事。结果，他的忍耐如此赢得了他们，以至于当时间助长了他们的判断，他们便将恼怒转为钦佩，跪在他面前请求宽恕，为先前行为感到羞愧，抛开仇恨，服从他为他们的族长、立法者与法官。\par

37. 同样的热忱也引发了他对异端者的反对，当时他们仗着皇帝的邪恶，发动进攻，妄想也压倒我们，将我们加入他们几乎已成功奴役的其他人之中。因在这事上，他藉着自己，也藉着虔敬的训练驱策我们如同良犬去攻击这些最凶残的野兽，给予我们不小的帮助。有一点我要责怪你们二位，求你们不要因我直言而见怪，若我表达痛苦之缘由令你们不悦。当我厌倦了世间的邪恶，渴望（若今世有人如此渴望）独处，并急于尽快逃离公共生活的喧嚣与尘土，逃入某个安全港湾时，是你们以司铎职那高贵的名号，抓住我并将我出卖给这个卑劣而险恶的灵魂市场。结果，祸患已降临于我，其余的尚可预见。因为过去的经验使人对将来多少有些怀疑，尽管理性有更好的建议。\par

38. 他的另一项卓越之处，我不可忽略。总的来说，他是一位极其坚忍的人，超越了他的血肉之躯；但在临终疾病的痛苦中，衰老的风险与重担又雪上加霜，他的软弱与常人无异；但这与其他奇事相称的结局，非但不平常，反而是他独有的。他从未脱离痛苦的折磨，但每日每时，有时每时辰，他唯一的缓解就是礼仪，痛苦如同被驱逐令一般向礼仪屈服。最后，在活了将近一百年之后，超过了\Person{达味}所定的我们寿数的界限，其中有四十五年——常人一生的年限——是在司铎职中度过的，他在美好的晚年结束了生命。以何种方式呢？以祈祷的言语与形式，没有留下任何恶迹，却留下了许多美德的回忆。人们对他所怀的敬意因此超越常人，无论在口舌上还是在心中。也不容易找到任何记得他而不\BibleQuote{用手掩口} \BibleRef{约 40:4}并纪念他的人。他的一生便是如此，他的成就与圆满也是如此。\par

39. 既然应当留下一些他慷慨的活的纪念，那么，除了他所建造的、献给天主和我们的这座圣殿之外，还需要什么呢？这座圣殿几乎完全由他私人财富支出，百姓的贡献很少。这是一项不应沉默地埋没的功业，因为论规模，它超过大多数；论美丽，它绝对胜过所有。它由八个等边四边形环绕，并凭借两层柱廊与门廊的美丽而高耸入云，其上放置的雕像栩栩如生；拱顶从上方向我们闪射光芒，从四面八方以丰富的光源使我们眼花缭乱，它确实是光的居所。它由最辉煌的材料构成的等角回廊环绕，中间有宽阔的区域，而它的门和门厅散发出优雅的光辉，从远处向走近的人表示欢迎。我尚未提及外部装饰，那些方形且榫接的石材的美与尺寸，无论是用于分隔转角的大理石柱基和柱头，还是来自我们自己的、绝不逊于外国的采石场；也没有提及那些从地基到屋脊突出的或镶嵌的多形多彩的带饰，它通过限制视野而夺走观赏者的注意力。谁能足够简洁地描述一件耗费如此多时间、辛劳和技艺的作品呢？或者说，这样说是否足够：在所有装饰其他城市的公私作品中，这座圣殿本身已能使我们在大众中享有盛誉？当这样一座圣殿需要一位司铎时，他也自费提供了一位，是否配得上这座圣殿，我无需置评。当需要祭献时，他也提供了，即在他儿子的不幸和他所受的试炼中，好让天主从他手中接受合理全燔祭和精神司祭，荣耀地被焚毁，以替代法律的祭献。\par

40. 我父，你怎么说？这样够了吗？你是否在这篇告别或埋葬的颂词中，为你为我学业所经受的一切辛劳找到了充分的补偿？你是否像往常一样，使我的舌头沉默，命令我适时停止，以免过度？还是你需要增加一些？我知道你叫我停止，因为我已经说得够多了。然而，请允许我再加一句。请告知我们你在荣耀中所处的位置，以及环绕你的光明；并请将很快随你而去的你的伴侣，以及你之前安息的孩子们，也接纳进同一居所，还有我，在这今生的苦难之后不再或只稍加增加时；在到达那居所之前，请在这甜美的石头中接纳我，这石头是你为我们俩竖立的，甚至在此地也是你献给你的同名的圣者的荣耀；请免除我对我已辞去的百姓的照顾，以及后来因你的缘故我所接受的百姓的照顾；愿你引导并保护整个羊群和全体圣职人员（他们称你为父）免于危险，我热切恳求，尤其是那位被你父性和精神的强制所压倒的人，使他不会完全认为那强制行为应受谴责。\par

41. 我言行的审判者，你对我们怎么看？如果我们说得足够，并合你的心意，请用你的决定确认它，我们接受；因为你的决定完全是天主的决定。但如果它远远不及他的荣耀和你的希望，我的后援并不远。请发出你的声音，他的功绩像及时雨一样等待着。的确，他对你有着最高的要求，即牧人对牧人、恩宠中的父亲对儿子的要求。如果那位曾通过你的声音向全世界发出雷鸣的人，自己也享受一些那声音，这有什么奇怪呢？还需要什么？只需将我们伟大的父亲\Person{亚巴郎}的生活伴侣和同路人，即属灵的\Person{撒辣}，结合在最后的基督徒礼中。\par

42. 我的母亲，天主的本性不同于人；实际上，一般而言，神性事物的本性不同于尘世事物的本性。它们拥有不变性和不朽性，以及绝对的存在及其后果，因为确定之物的属性是确定的。但我们的事物是怎样的呢？它们处于流变和朽坏中，不断经历某种新的变化。生命和死亡，正如它们被称呼的，表面上如此不同，却在某种意义上相互消解和相继。因为一个始于我们的母亲——朽坏，经由朽坏——即一切现有之物的更替——而运行，并在朽坏——即这生命的解体——中结束；而另一个，能够使我们脱离今生的苦难，并常常将我们提升到更高生命，在我看来并不准确地被称为死亡，其名称比现实更可怕；因此我们有危险不合理地惧怕不可怕之事，并把真正应当惧怕的事当作更可取的来追求。只有一种生命，就是仰望生命。只有一种死亡，就是罪，因为它是灵魂的毁灭。但其他一切，有些人引以为傲的，都是梦幻，嘲弄现实，是一系列引导灵魂误入歧途的幻象。母亲，如果情况是这样，我们就不会以生命为傲，也不会因死亡而大为受伤。从此世被转移到真实生命，有什么可抱怨的呢？从今生的变化、动荡、厌恶和所有我们必须付出的卑贱代价中解脱出来，进入稳定、无流变的事物中，作为较小的光体环绕大光旋转，这有什么可抱怨的呢？\BibleRef{创 1:16}\par

43. 分离的感觉使你痛苦吗？让望德来鼓励你。丧偶对你来说是沉重的吗？但对他却并非如此。如果爱德只给自己容易的事，而把更困难的事推给邻人，那么爱德有何益处呢？为什么它对于即将逝去的人会是沉重的呢？指定的日子临近了，痛苦不会长久。让我们不要通过卑劣的推理，使本来轻松的事情变成负担。我们遭受了巨大的损失，因为我们所享有的特恩是伟大的。损失对众人是共同的，这样的特恩却只有少数人享有。让我们借着后者的安慰来超越前者的想法。因为更合理的是，那更好的应得胜。你已以最勇敢的基督徒精神，承受了失去仍在壮年且有资格生存的子女的痛苦；也承受一位厌倦生活的人脱下他年老的身体，尽管他精神的活力仍为他保全了完好的感官。你想要有人照顾你吗？你的\Person{依撒格}在哪里？他为你留下他，好让他全面替代他的位置？向他要求小事，扶持他的手和服务，并以更大的事回报他：母亲的祝福和祈祷，以及随之而来的自由。你因被劝诫而烦恼吗？我为此赞美你。因为你已劝诫了许多在你长久生命中接触到的人。我所说的对你这样真正智慧的人毫无影响；但让它成为哀悼者的一般安慰良药，使他们知道自己是凡人跟随着凡人进入坟墓。\par

\SourceRecord{Translated by Charles Gordon Browne and James Edward Swallow.；Nicene and Post-Nicene Fathers, Second Series；Vol. 7.；Edited by Philip Schaff and Henry Wace.；Buffalo, NY: Christian Literature Publishing Co.,；1894.；<http://www.newadvent.org/fathers/310218.htm>.}

% structure: p0001 source=h1 target=section reason=page-title

\section{第二十一篇讲道}

\Emph{论伟大的\Person{亚大纳削}，\Place{亚历山大}的主教}\par

\EditorialNote{第22节提到最初在\Place{塞琉西亚}举行、后来在这座伟大城市聚会的会议，这毫无疑问地表明这篇讲道是在\Place{君士坦丁堡}发表的。在第5节的暗示中还可以找到更多的地方色彩。从对\Person{圣西彼廉}的颂词（\WorkTitle{讲道集}24,1）我们可以确知，\Place{君士坦丁堡}教会已有在每年圣人的节日举行庆典的惯例：目前东方教会守两个日子，即1月18日——\Person{圣亚大纳削}实际去世的日子，以及5月2日——纪念他的遗骸迁移到\Place{君士坦丁堡}的\Place{圣索菲亚}教堂。因此，很可能这篇讲道是在前一个日子发表的，\Person{阿塞马尼}认为\Person{圣亚大纳削}是在那一天去世的。\Person{帕佩布罗克}和\Person{布赖特博士}有些犹豫地赞成5月2日。\Person{蒂勒蒙}假定发表年份是公元379年；如果是这样，那么它一定是在\Person{圣国瑞}抵达该城之后不久发表的。然而，由于文中完全没有提到这件事，总体看来，它更可能应归为公元380年。这篇讲道在\Person{圣国瑞}的演说中占有很高地位，被视为教会颂词的典范。不过，它缺乏那种与家人朋友的讲道所特有的个人情感和对内在生活的亲密了解的魅力。}\par

1. 赞美\Person{亚大纳削}，就是赞美德行。谈论他与赞美德行是一回事，因为他曾经——或者更准确地说，现在依然——完全拥有了德行。因为所有按天主而生活的人，虽然离开此世，却仍为天主而活。为此，天主被称为\Person{亚巴郎}、\Person{依撒格}和\Person{雅各伯}的天主，因为祂不是死人的天主，而是活人的天主。\BibleRef{玛 22:32} 再者，赞美德行，就是赞美天主，祂将德行赐予人，并藉着与祂相似的启示，把人提升——或再次提升——到祂自己面前。\BibleRef{若一 1:5} 因为我们从天主所领受并将要领受的恩惠，无论多大、多少，无人能尽述，其中最好、最仁慈的，莫过于我们向祂的归向以及与祂的关系。因为天主之于理智之物，如同太阳之于可感之物。一个照亮可见的世界，另一个照亮不可见的世界。一个使我们肉体的眼睛看见太阳，另一个使我们理智的本性看见天主。同样，那赋予观看与被观看之物以观看与被观看之能力的，它本身就是可见之物中最美的；同样，天主为思维者与所思之物创造思维与被思维的能力，祂本身就是思维对象的最高者，一切欲望在祂那里找到终点，越过祂便再不能前行。因为即使是最有智慧、最锐利、最好奇的理智，也没有——也永不能有——一个更高超的对象。因为这是可欲之物的极致，达到它的人便在思辨中找到完全的安息。\par

2. 凡得以藉着理性与默观，脱离物质和这肉体的云雾或面纱（无论怎样称呼），并与天主相通，且按人性所能达成的程度，与最纯净的光明结合的人，他是有福的，既因他从这里上升，也因他在那里的神化——那是藉着真正的哲学和对物质二元论的超越，通过在天主圣三内所感知的合一而赐予的。凡因与肉身结合而败坏，且被泥土如此压制，以至不能注视真理之光，也不能超越下界事物的人——虽然他本是从上而生、被召向上——我认为他在盲目中是可怜的，即使他可能在此世的事物中富足；尤其是因为他成为他富足的玩物，并被它说服，相信别的东西是美的，而非那真正美的，收获他那可怜见解的可怜果实：黑暗的判决，或见到那不被承认为光的火。\par

3. 这样的哲学，不论现今还是古时，都只有少数人拥有——因为天主的人很少，虽然所有人都是祂的化工——在立法者、将领、司祭、先知、圣史、宗徒、牧者、教师，以及所有属灵的军团和队伍中——而在他们中，也有我们现在所赞美的人。我指这些人中的哪些呢？有如\Person{哈诺客}、\Person{诺厄}、\Person{亚巴郎}、\Person{依撒格}、\Person{雅各伯}、十二圣祖、\Person{梅瑟}、\Person{亚郎}、\Person{若苏厄}、民长、\Person{撒慕尔}、\Person{达味}、在某种程度上还有\Person{撒罗满}、\Person{厄里亚}、\Person{厄里叟}、被掳前的先知、被掳后的先知，以及——虽然在顺序上最后，但在真理上首位的——那些与基督降生成人、或祂摄取我们本性有关的人：那在光明之前的灯，那在圣言之前的声音，那在\OriginalTerm{中保}{Mediator}之前的中间者，那在新旧盟约之间的中间者——著名的\Person{若翰}；还有基督的门徒，以及基督之后、治理百姓的人，或在言语上卓著的人，或以奇迹引人注目的人，或藉着流血而成全的人。\par

4. 在与这些人的比较中，\Person{亚大纳削}与一些人竞争，被另一些人稍胜一筹，而对另一些人——若不怕大胆地说——他超过了他们：有些人他以他们的心智能力为榜样，有些以他们的行动，有些以他们的温和，有些以他们的热诚，有些以他们的冒险，有些在大多数方面，有些在全部方面；他从一个又一个那里收集各种形式的美（就像画家画理想中卓越的人物一样），并将它们融合在自己单一的灵魂中，从所有美中塑造了一个完美的德行形象。他在行动上胜过那些有才智的人，在才智上胜过那些有行动的人；或者，如果你愿意，他在才智上胜过那些以才智闻名的人，在行动上胜过那些有最大行动力的人；他超越了那些在两方面都有中等声誉的人，因他在每一方面的杰出，也超越了那些在一方面或另一方面站得最高的人，因他在两方面的能力；如果对于那些接受了榜样的人来说，能如此运用它以使自己依附于德行是一件大事，那么，他为了我们的益处而为后来人树立了榜样，他的名声也绝不逊色。\par

5. 要详尽地谈论和赞叹他，对于我目前演讲的目的来说，或许是一项过于冗长的任务，并且会采取历史的形式而非颂词的形式：我一直渴望把这样一部历史书写下来，以供后世愉悦和教训，正如他本人写了圣\Person{安当}的生平，并以叙述的形式阐述了隐修生活的规律。因此，在进入他的历史诸多细节中的几项——根据此刻记忆所及，认为最值得注意的——以满足我自己的渴望并履行符合这个节日的职责之后，其余许多事我们留给知道它们的人。因为，当不敬虔者的生活因被记念而受到尊崇时，沉默地经过那些敬虔生活的人，实在既不是敬虔的，也不是安全的，尤其是在一座城市里——这座城市几乎不能因许多德行的榜样而得救，因为它对神圣的事物，如同对赛马和戏剧一样，加以嘲弄。\par

6. 他自幼年起，就在宗教习惯与敬虔操练中成长，于此之前曾略学文哲，以免对此类知识全然生疏，或对他已决意轻蔑之事一无所知。因为他那高尚而热切的灵魂，不能容忍耽于虚妄，如同不谙技巧的运动员，对着空气挥舞而不击对手，以致失去奖赏。由于对\WorkTitle{旧约}与\WorkTitle{新约}每一卷书的默想，其深度无人能及，即使仅对其中一卷亦无人能及，他遂在默观中富有，在生活的光辉中富有，以那少有人能编织的金色纽带奇妙地将二者结合，以生活引导默观，以默观印证生活。因为敬畏主是智慧的开端，可谓其最初的襁褓；但当智慧挣破畏惧的束缚，上升至爱，它便使我们成为天主的朋友，不再是奴仆，而是儿子。\par

7. 他就是这样被培养和训练的，正如现今那些将要治理人民、依照天主的旨意和预知（这预知早早奠定了伟大功业的基础）掌管基督奥体的人所应当的那样。他被赋予这一重要职务，成为那些亲近那亲近我们的天主的人中的一员，并被认为配得神圣的职位和品级。在经历了整个阶序之后，简而言之，他被托付了治理人民的首要权柄——换言之，就是整个普世的责权。我无法说他是以司铎职作为德行的奖赏，还是作为教会之泉源和生命。因为教会，如同\Person{依市玛耳}（\BibleRef{创 21:19}）因渴慕真理而昏厥，需要被给予水喝，又如\Person{厄里亚}（\BibleRef{列上 17:4}）在旱灾时得溪水滋润；当她气息微弱时，被恢复生命并留存为\Place{以色列}的遗民（\EditorialNote{依 1:9}），使我们不至像\Place{索多玛}和\Place{哈摩辣}（\BibleRef{创 19:24}），其被火与硫磺毁灭的恶名昭彰甚至超过其邪恶。因此，当我们被贬抑时，为我们兴起了救恩的角（\BibleRef{路 1:69}），并在适当时机安置了角石（\EditorialNote{依 28:16}），将我们与它及彼此联结；又如火（\BibleRef{拉 3:2}）以炼净我们卑劣邪恶的材质，又如农夫的簸箕（\BibleRef{玛 3:12}）以扬去教义中轻浮者，分离出稳重的，又如剑以斩断邪恶的根源。于是圣言找到他作自己的盟友，圣神占有了那将为其发言之人。\par

8. 因此，基于这些理由，他并非像后来盛行的邪恶方式那样，也非通过流血与压迫，而是以宗徒式的、属神的方式，被全体人民的选举领上了\Person{圣马尔谷}的宝座，继承他的虔敬，不亚于继承其职位；就后者而言，他确实与他相距甚远；但就前者——这是真正继承的权利——而言，他却紧紧跟随。因为教义上的合一理应带来职位上的合一；而一位对立教师会设立对立的宝座；前者是真正意义上的继承者，后者只是名义上的。因为继承者并非那闯入者，而是其权利被侵犯者；不是违法者，而是合法受任命者；不是持相反意见者，而是持同一信仰者。如果这不是我们对继承者的理解，那么他继承的方式就如同疾病继承健康、黑暗继承光明、风暴继承平静、疯狂继承健全理智。\par

9. 他履行职务的精神与他被选任时的精神相同。因为他在登上宝座后，并未像那些意外攫取主权或遗产的人那样，因陶醉而傲慢。这是那不合法的、闯入的司铎的行径，他们不配自己的圣召；他们的司铎准备不费他们分文，他们未为德行忍受任何不便，只在被任命教导时才学习宗教，且在自己被洁净之前就着手洁净他人；昨日亵圣，今日司祭；昨日被摒于圣所之外，今日供职其中；精于恶习，疏于虔敬；是人的恩宠的产物，而非圣神的恩宠；他们使尽各种暴力，最终甚至以暴力对待虔敬；他们不以品格为职位增光，反而需要职位的声誉来支撑他们的品格，从而颠倒了二者间应有的关系；他们应当为自己献祭多于为民众的愚昧献祭；他们不可避免地落入两种错误之一：或因自身需要宽恕而过度宽恕，甚至教导而非遏制恶习，或以严厉的统治掩盖自己的罪过。这两种极端他都避免了：他在行动上崇高，在思想上谦卑；在德性上难以企及，在交往上极易亲近；温和，不轻易发怒，富有同情心，言语甘甜，性情更甘甜；外貌如天使，内心更胜天使；责备时平静，赞美时有说服力，且不因过度而破坏两者的良好效果，而是以父亲的温柔责备，以统治者的尊严赞美；他的温柔并未消散，他的严厉并未变得酸涩；因为前者是合理的，后者是明智的，两者皆真正智慧；他的性情足以培养属灵的子女，几乎不需言语；他的言语几乎不需棍棒（\BibleRef{格前 4:21}），而他节制地使用棍棒，就更无需用刀了。\par

10. 但我为何要为你们描画此人的肖像？\Person{圣保禄}早已预先勾勒了他。当他颂扬那位穿越诸天的伟大司祭时，便如此行（我甚至敢这样说，因为圣经能将那些按基督生活的人称为基督）；同样，当他按弟茂德前书中的规则，为未来的主教提供典范时，也是如此。因为若你将此法作为标准来衡量这位当受赞美者，你必将清楚察觉他的完美无缺。那么，请来助我完成这篇赞辞吧；因我在讲论中倍感沉重，虽想一点一点地略过，它们却接二连三地抓住我，我在这各方面匀称优美的形体中找不到任何卓越之处；因为每当我想到一点，它似乎都比其余的更美，从而强夺了我的言词。因此，我恳请你们这些曾仰慕并见证他的人，请彼此分享他的美德，勇敢地相互较量，无论男女老幼、青年与少女、老人与孩童、司铎与民众、独修者与共修者、生活简朴或严谨者、默观者或实践者。让一人赞美他的禁食与祈祷，仿佛他无肉身、非物质；另一人赞美他的不倦与对守夜和圣咏的热忱；另一人赞美他对贫者的庇护；另一人赞美他对权贵的不畏，或对卑微者的屈尊。让贞女们赞美新郎之友；让在婚姻羁绊中的人赞美他们的节制者，隐修士赞美他使他们展翅高飞，共修者赞美他们的立法者，纯朴者赞美他们的向导，默观者赞美那位神圣者，喜乐者赞美他们的缰绳，不幸者赞美他们的安慰，白发者赞美他们的支柱，青年赞美他们的导师，穷人赞美他们的恩人，富人赞美他们的管家。我想，连寡妇也会赞美她们的护佑者，孤儿赞美他们的父亲，贫者赞美他们的施恩者，异乡人赞美他们的款待者，弟兄赞美那富于手足情谊者，病人赞美他们的医生——无论何种疾病或疗法，健康者赞美健康的守护者，是的，所有人赞美那位对众人成为一切，为要赢得几乎——若非完全——所有人的人。\BibleRef{希 4:14}\BibleRef{若 3:29}\par

11. 基于这些理由，正如我所说，我将让其他有空闲欣赏他性格中细微之处的人去赞美和颂扬他。我称它们为细微之处，只是将他和他的性格与他自己的标准相比；因为那已得光荣的，即使极其辉煌，也因那超越的光荣而不再显赫，正如经上所载；的确，他卓越中的细微之处也足以使他人成名。但若我为了服侍次要细节而离弃圣言，那是不可容忍的；因此我必须转向他最杰出的特征；只有天主——我为之发言的那一位——才能使我说出配得上这样一位在圣言中如此崇高而强大的灵魂的话语。\BibleRef{格后 3:10}\BibleRef{宗 6:2}\par

12. 在教会全盛之时，一切安好，当时那种精心雕琢、牵强附会、矫揉造作的神学论说尚未进入神学学派；而用快速变幻的卵石欺骗眼睛的戏法，或在观众面前以多样而柔媚的扭曲姿态跳舞，与以不寻常或轻浮的方式谈论或聆听天主被视为一体。但自从\Person{塞克斯图斯}和\Person{皮洛}，以及那种反题风格，如同一种可怕的恶性疾病，感染了我们的教会，喋喋不休被视为文化，正如\BibleBook{宗徒}{大事录}论雅典人所说，我们除了谈论或听闻一些新奇之事外，别无所事。啊，哪一位\Person{耶肋米亚}能哀叹我们的混乱与盲目疯狂？只有他能发出与我们的不幸相称的哀歌。\BibleRef{宗 17:21}\BibleRef{哀 1:1}\par

13. 这种疯狂的起源是\Person{亚略}（他的名字源于疯狂），他因放肆的口舌而在一个亵渎之地死去，这是祈祷所致而非疾病，他像\Person{犹大}一样，因对圣言相似的背叛而肚腹崩裂。然后其他人感染了这疾病，组织了一种不虔敬的艺术，将神性局限于\OriginalTerm{非受生者}{Unbegotten}，不仅将\OriginalTerm{受生者}{Begotten}，也将\OriginalTerm{那出于者}{the Proceeding one}排除于神性之外，仅以名义上的共融来尊崇天主圣三，甚至拒绝为其保留这一共融。但那位有福者却不是这样，他真是天主的人、真理的强有力的号角。他意识到将三个位格紧缩为数字上的合一乃是异端，是\Person{撒柏略}的革新——他最先发明了神性的紧缩；而按本性的区别割裂三个位格，则是神性的一种违背自然的肢解。因此，他既幸福地保存了属于天主性的合一，又虔诚地教导了关于位格的圣三，既不在合一中混淆三个位格，也不在三个位格中分割实体，而是保持在虔敬的界限内，避免过分偏向或反对任何一方。\BibleRef{宗 1:18}\par

14. 因此，他首先在神圣的尼西亚大公会议中——那三百一十八位特选者的集会，因圣神而结合，尽己所能——遏止了疫病。虽然他尚未被列于主教之中，却在大公会议的成员中居于首位，因为品德与职位同等受重视。之后，火焰被恶者的吹动所助长，蔓延极广（由此产生了几乎整个大地和海洋都充满的悲剧），战斗激烈地围绕着他，那圣言的尊贵斗士。因为攻击最猛烈的地方就是抵抗点，而各种危险从四面包围；不虔敬善于设计恶事，并极为大胆地付诸实施：他们连神性都不宽恕，又怎能宽恕人呢？然而，其中一次攻击是至为危险的：我自己也对这一幕有所贡献；且容我为亲爱的祖国辩护无罪，因为这邪恶并非源于生养他们的土地，而是源于从事其事的人。因为那地的确是圣洁的，并因虔敬而广为所知；但这些人不配生养他们的教会；你们曾听说荆棘生在葡萄树中；叛徒\BibleRef{路 6:16} 就是犹达斯，门徒之一。\par

15. 有些人甚至不原谅我的同名者；他当时为求学而住在亚历山大，虽然曾受到他最仁慈的对待，如同最亲爱的儿女，并得到他特别的信任，却参与了反叛父亲与恩主的阴谋：因为虽然别人采取了积极行动，但正如俗语所说，阿贝沙隆之手与他们同在。如果你们中有人听过那以欺诈反对圣人的手、活人的尸体以及不义的放逐，他就明白我的意思。但我乐意忘记此事。因为在可疑之处，我倾向于认为我们应当偏向仁慈，宁可开释而非定罪。因为坏人会急忙定好人的罪，而好人不轻易定坏人的罪。因为不准备作恶的人，也不倾向于怀疑恶。现在我来谈事实而非传言，是确证为真理而非未经证实的猜疑。\par

16. 有一个怪物来自卡帕多细亚，生在我们最远的边界，出身卑微，心智更低，血统并不完全自由，而是混种，如同我们知道骡子那样；起初，依靠他人的餐桌为生，价格就是一块大麦饼，学会了为了肚子而说和做一切事；最后，在潜入公共生活之后，担任最低的职务，例如猪排承包商、军队口粮承包商，然后在这公共委托中因贪婪显出自己是恶棍，被剥光，设法逃脱；像流亡者一样从一个国家到另一个国家，一个城市到另一个城市；最后，对于基督徒团体而言在不祥的时刻，如同埃及的灾祸之一，到达了亚历山大。在那里，他的流浪停止，开始了他的恶行。在其他方面一无所长，没有文化，言谈不流利，甚至连恭敬的形式和伪装都没有，但他制造恶行和混乱的技巧无人能及。\par

17. 你们都知道他对圣人的无礼行为的一切细节。义人常被交在恶人手中\BibleRef{约 9:24}，并非为使恶人得尊荣，而是为使义人受考验；虽然恶人正如经上所记，要遭可怕的死亡，然而目前虔敬之人却成为笑柄，而天主的良善和为每个人所预备的丰富宝藏仍被隐藏。那时，言语、行为和思想都将在天主公义的天平上衡量，当祂起来审判大地，聚集计谋和作为，揭示祂所封存的事\BibleRef{达 12:9}。愿约伯的言行和受苦使你们信服此事：他是一位真实、无可指摘、公义、敬畏天主的人，并具有为他所证实的其他一切品质，然而却被那向他求取权柄者接连以如此显著的灾祸击打，以致虽然许多人在整个历史中常常受苦，有些人甚至可能遭受了严重的打击，但没有人能在不幸中与他相比。因为他不仅受苦——没有时间哀悼接踵而来的损失：钱财、财产、众多而美好的家庭，人人珍视的福分——而且最后被一种难以治愈、外观可怕的疾病击打；更糟的是，有一个妻子，她的唯一安慰是邪恶的劝告。因为他最大的痛苦是灵魂的折磨加于身体的折磨。他还有朋友中真正的苦恼安慰者，如他所说，他们不能帮助他。因为他们看到他的苦难，却不知其隐藏的意义，便以为他的灾祸是恶行的惩罚，而非德性的试金石。他们不仅这样想，甚至不羞于以他的命运责备他；而那时，即使他是因恶行受苦，他们也该以安慰的话对待他的悲伤。\par

18. 这就是约伯的命运：初看其历史如此。实际上，这是一场德行与嫉妒之间的较量：一方竭尽全力要战胜善，另一方则忍受一切，以求不被制服；一方试图通过惩罚义人来为恶铺路，另一方则紧握善，即使他们在不幸中超过别人。那么，那从旋风和云彩中回答约伯的（\BibleRef{约 38:1}），那迟于责罚、急于救助，那不让恶人的杖完全落在义人的份上、免得义人学习不义的，祂又做了什么？在较量结束时，祂以荣耀的宣告宣布了这位斗士的胜利，并揭示了他苦难的奥秘，说：\Quote{你以为我待你是为了别的目的，而不是要显明你的义吗？}这是医治他伤口的香膏，这是竞赛的冠冕，这是他忍耐的赏报。因为后来的繁荣也许是微不足道的，尽管在有些人看来似乎很大，并且是为小器的人而定的，即使他重新得到了比他失去的多一倍。\par

19. 既然这样，若乔治占了阿塔纳修的便宜，并不奇怪；倒不如说，义人不被凌辱之火试炼，才更奇怪；这也不奇怪，因为火焰若不止于此，才更奇怪。那时他退隐，极好地安排了流亡生活，因为他前往埃及那些神圣而神圣的默观之所，在那里，人们远离世界，迎接旷野，比所有有形躯壳中的人更活于天主。有些人在完全隐修和独居的生活中奋斗，只对自己和天主说话（\BibleRef{格前 14:28}），他们所知的世界仅是在旷野中所见。另一些人则在团体中珍视爱德的法律，既是独修者又是共修者，对其他人和公共事务的漩涡——这漩涡卷着我们，自己也在旋转，在突变中戏弄我们——是死的，他们彼此成为世界，在竞争中磨砺爱的锋刃。在与他们交往期间，伟大的阿塔纳修——他始终是所有人的中保和修和者，如同那藉着自己的血（\BibleRef{哥 1:20}）使对立者和好的那一位——将独修与团体生活调和：表明司祭职能够默观，而默观需要一位灵性导师。\par

20. 这样，他将两者结合，使安静行动和行动安静的拥护者如此合一，以致他们相信，隐修生活的特征在于性情的坚定，而非身体的退隐。因此，伟大的达味同时是最活跃和最独居的人，如果有人认为\Quote{我独处，直到我离去}这句诗在阐释这一主题时有价值和权威。因此，虽然他们在德行上超越众人，但他们在他的心思前，比众人在他们自己的心思前更加不足；他们对他司祭职的完美贡献甚少，却从默观中获得更大的助益。他所想的，就是为他们立的法律；他所不赞成的，他们就弃绝；他的判决对他们来说就是梅瑟的版，他们向他表示敬意，超过人应给予圣人们的敬意。是的，甚至当人们像猎捕野兽一样来搜捕这位圣人，到处搜寻却找不到时，他们对这些使者不发一言，而把头颈伸向刀剑，为基督的缘故冒死，认为为阿塔纳修遭受最残酷的苦难是迈向默观的重要一步，比他们常沉浸其中的长期禁食、硬卧和苦修更为神圣和崇高。\par

21. 当他赞同撒罗满的明智箴言——事事有时节（\BibleRef{训 3:1}）——时，这就是他的处境；于是他隐藏了一段时间，在战争时期逃避，为在后来和平到来时显现自己。与此同时，乔治因无人抵抗，以不敬神之权蹂躏埃及，摧毁叙利亚。他尽力夺取东方，像洪流卷走沿途之物一样，不断吸引弱者，攻击动摇或胆怯的人。他也赢得了皇帝（因我必须称他的不稳定为单纯，虽然我尊重他虔诚的动机）的单纯。因为说实话，他有热心，但不按知识（\BibleRef{罗 10:2}）。他收买了那些爱财胜过爱基督的掌权者——因为他手头有穷人的基金，他侵吞了——尤其是那些软弱无男子气、性别可疑但显然不敬神的人；我不知道为何或如何，罗马皇帝竟将统治人的权威委托给他们，尽管他们本应管理妇女。这就是那恶者的仆人、那撒莠子的人、那假基督先驱的力量；他是当时主教中首屈一指的演说家；如果有人愿意称他为演说家，他与其说是一个不敬神的人，不如说是一个敌对和好辩的辩论者——他的名字我乐于不提：他是他党派的手，用为敬神用途募集的黄金来歪曲真理，而恶人把它当作不敬神的工具。\par

22. 这一派系的登峰造极之举是首先在塞琉西亚——神圣而显赫的童贞女圣特克拉的城市——举行的会议，随后在这座大城举行，从而将他们的名字不再与高尚的关联相连，而是与最深的耻辱相连；无论我们应将这颠覆和扰乱一切的会议称为加兰塔，\BibleRef{创 11:4}——它理应混乱口舌，但愿他们的口舌因作恶的和谐而被混乱！——或称为盖法的公议会，在那里基督被定罪，或其他类似的名字。那捍卫天主圣三的古老而虔敬的教义被废除，通过设立栅栏和摧毁同质者（\OriginalTerm{同质者}{\GreekText{ὁμοούσιον}}）；藉着所写的内容为不敬神敞开大门，以尊重圣经和使用被认可的术语为借口，但实际上引入了非圣经的亚略主义。因为\Quote{按圣经，相似}这短语是对纯朴者的诱饵，隐藏着不敬神的钩子，一个看似向所有过路人注视的形像，一只适合作祟的靴子，随着每一阵风簸扬，\BibleRef{训 5:9}从新写的邪恶和对抗真理的诡计中获得权威。因为他们作恶有智慧，行善却无知。\BibleRef{耶 4:22}\par

23. 由此而来的是他们对异端者假装谴责，在言语上弃绝他们，以便为他们的努力赢得可信性，但实际上却助长了他们；不是指控他们无限的邪恶，而是指控他们夸张的言辞。由此而来的是亵渎圣人的审判官，新的组合，对神秘问题的公开审视和讨论，对生活行为的非法调查，雇佣的告密者，购买来的判决。有些人被不公正地从他们的主教座位上废黜，另一些人被强加进去，而在其他必要条件中，他们被迫签署不义的契约：墨水已备好，告密者在场。这是我们中大多数未屈服的人必须忍受的，虽然内心没有倒下，却被说服签署，从而与在两方面都邪恶的人联合，将自己卷入烟雾中，即使不是火焰中。对此我曾多次哭泣，当我默想当时不敬神的混乱，以及如今从圣言的庇护者那里兴起的对正统教义的迫害。\par

24. 因为事实上，正如圣经所说，\BibleQuote{牧人变成了畜类，许多牧人毁坏了我的葡萄园，玷污了我可喜悦的产业}，\BibleRef{耶 10:21}我指的是天主的教会，它是由许多在基督前后劳苦和流血的人以及受害者，甚至天主为我们所受的伟大苦难所聚集的。因为除了极少数例外，这些人或是因卑微而被忽视，或是因德行而英勇抵抗，作为遗民\EditorialNote{依撒意亚 1:9}被留给以色列，正如所预定的，作为种子和根，要在圣神的溪流中开花并复活，每个人都屈服于时代的影响，区别仅在于有些人较早，有些人较晚，有些人成为不敬神的冠军和领袖，而另一些人则被分配较低的等级，他们或因恐惧而动摇，或因需要而被奴役，或因奉承而被迷惑，或因无知而被欺骗；后者罪过最小，如果确实我们允许这甚至成为被托付领导人民之人的有效借口。因为正如狮子和其他动物的力量，或男人和女人的力量，或老年人和青年人的力量并不相同，而是由于年龄或种类有相当大的差异——统治者和他们的臣民也是如此。因为虽然我们在这种情况下可以原谅平信徒，而且他们常常免受惩罚，因为没有经受考验，但如何能原谅一位教师，他的职责是纠正他人的无知，除非他徒有其名？因为，难道不是荒谬吗——无论一个人多么粗鲁和缺乏教育，都不允许他对罗马法无知，也没有法律支持无知之罪——救恩奥秘的教师竟对救恩的基本原则无知，尽管他们在其他问题上头脑可能简单浅薄？但是，即使对无知而犯错的人给予宽容，对于其余的人，他们自诩头脑敏锐，却因我提到的原因屈服于宫廷派，并在长时间扮演虔敬角色后，在考验的时刻失败，又该说什么呢？\par

25. \Quote{再一次，}我听到圣经说，天地将要被震动，正如这事以前已经发生过一样，依我看，这预示着万物的明显更新。我们必须相信圣保禄所说的，\BibleRef{希 12:27}这最后的震动正是基督的第二次来临，以及宇宙向一种不可动摇的稳定状态的转化和改变。而且我想，这次当前的震动，那些默观者和爱天主者——他们提前操练天上的国民身份——从我们中间被震动，其重要性不亚于往日的任何一次。因为，尽管这些人在其他方面温和适度，但他们不能容忍为了安宁而将其理性推延到背叛天主事业的地步：相反，在这一点上，他们极其好战且战斗坚强；他们的热忱如此炽烈，以至于他们宁愿在动荡中走向极端，也不愿忽视任何错误之事。民众中有相当大一部分人正与他们一起分离，像一群鸟跟随领头的鸟一样飞走，甚至现在也不停止与他们一起飞翔。\par

26. \Person{亚大纳削}对于我们，当他临在时，就是教会之柱；甚至当他因恶人的凌辱而隐退时，也是如此。因为那些图谋攻取某坚固堡垒的人，当他们看不到其他容易接近或夺取的方法时，便诉诸计谋，然后用金钱或诡计诱惑守将，不费吹灰之力便占据了据点；或者，如果你愿意，如同那些图谋对付\Person{三松}的人，先剃去他的头发\BibleRef{民 16:19}——他的力量所在——然后抓住这民长，任意戏弄他，以报复他昔日的威力：同样，我们外来的敌人，在除掉了我们力量的源泉，剪除了教会的荣耀之后，也如此放纵于不敬的言论和行为。然后，敌对牧者的支持者和庇护者去世了，他用一个邪恶的结局为他本非邪恶的统治加冕，并在临终时无益地悔改（如人们所说），那时每个人面对更高的审判台，都是自己行为的审慎审判者。因为他承认，他对这三件有损他统治的恶事负有责任：杀害亲属、宣布背教者、以及革新信仰；据说他就这样离开了。这样，讲授真理之言的权利再次确立，那些曾遭受暴力的人如今有了不受干扰的言论自由，而嫉妒正在磨砺其怒火的武器。亚历山大里亚的人民也是如此，他们以一贯对傲慢者的不耐，无法容忍此人的放纵，因此用一种不寻常的死亡来标记他的邪恶，用不寻常的耻辱来标记他的死亡。因为你们知道那只骆驼、它奇特的负担、新的高举形式，以及第一次（我想也是唯一一次）的游行，至今傲慢者仍受此威胁。\par

27. 但是，当这不义的暴风、这虔诚的败坏者、这恶者的前驱受到了这样的报应——以一种我不能赞扬的方式，因为我们必须考虑的不是他该受什么，而是我们该做什么——但无论如何，作为公众愤怒和激动的结果，报应降下了：于是，我们的斗士从他那光辉的流放中复归（我这样称呼他为了天主圣三并以天主圣三的祝福而放逐的经历），城市人民和几乎全埃及都如此欢欣，以至于他们从四面八方、从国家最远的边界跑来，只为听到\Person{亚大纳削}的声音，或饱览他的面容，甚至（如同我们听说宗徒们的事）被他的身体的影子和不真实的形象所圣化\BibleRef{宗 5:15}：因此，尽管在时间的长河中，有许多荣誉和欢迎时常赐予不同的人，不仅是统治者和主教，还有最杰出的公民，但没有一次记录有比这次更众多或更辉煌的。而只有一次荣誉能与它相比，那是\Person{亚大纳削}自己曾领受过的，在他之前进入这座城市时，当时他从同一放逐中为相同原因返回。\par

28. 关于这个荣誉，当时还流传着如下一个传闻；我将冒昧提及它，尽管它或许是多余的，作为我演讲的一点调味，或为纪念他的进城而散落的花朵。在那次进城之后，一位曾任两任执政官的官员正骑马进入城市；他是我们中的一员，是\Place{卡帕多细亚}人中最著名的一位。我相信你们知道我说的是\Person{菲拉格里乌斯}，他赢得我们的爱戴远超任何人，受尊敬与受爱戴相当，如果我可以这样简略地列举他所有的荣誉：他曾被皇帝决定，应市民请求，第二次委托管理该城。然后，在场的其中一位平民，看到人群巨大得如同无边的大海，据说对他的一位同伴和朋友说——这种场合下常有之事——\Quote{告诉我，我的好伙计，你以前见过人们如此众多、如此热情地涌出以向任何人致敬吗？}\Quote{没有！}那年轻人说，\Quote{而且我想，即使是\Person{君士坦提乌斯}本人也不会受到这样的待遇。}他通过提及皇帝来指明可能荣誉的顶峰。\Quote{你说的是那个，}另一人带着甜美欢快的笑声说，\Quote{我几乎不敢相信，即便是伟大的\Person{亚大纳削}也会受到这样的欢迎。}同时加了一句我们本地的誓言来证实他的话。现在，他所说的话的要点（我想你们也清楚看到）是，他将我们颂扬的对象置于皇帝本人之上。\par

29. 众人对圣人的敬意如此之大，甚至现在我所描述的欢迎场面仍令人惊叹。因为若按出身、年龄和职业来划分（而城市在给予某人公开荣誉时，通常是这样安排的），我怎能用言语来描绘那宏伟的景象？他们汇成一条河流，这确实是一位诗人的任务，去描述那尼罗河，真有黄金般的河流，五谷丰登，从城市倒流回克厄鲁姆，我想，有一天的路程，甚至更远。请允许我再沉醉一会儿我的描述：因为我正要前往那里，连我的言语也难从那典礼中带回。他骑在一头小驴上，几乎，请别怪我痴迷，如同我的耶稣骑在另一头小驴上，\BibleRef{路 19:35}——无论那是外邦人，祂仁慈地骑上他们，使他们脱离无知的束缚，还是别的什么，圣经所启示的。人们用树枝、五彩斑斓的衣裳欢迎他，衣裳被撕下，撒在他面前和脚下：只有在那一刻，所有光荣、昂贵、无与伦比的东西都被这样轻慢地对待。再次，如同基督的荣进，那些在前面欢呼和后面跳舞的人；只是歌颂他的群众不仅是儿童，而是每个人口舌和谐，争相超越对方。我略过普遍的欢呼、香膏的倾倒、彻夜的宴乐、整座城闪耀的光芒、公宴和家宴，以及所有表达城市欢乐的方式，那时人们都慷慨而难以置信地给予他。就这样，这位非凡的人在如此盛大的迎接下来到他的城池。\par

30. 他随后过着适合如此子民统治者的生活，但他没有按自己所活的来教导吗？他的争战与他的教导不和谐吗？他的危险比任何据理力争的人少吗？他的荣誉低于他所奋斗的目标吗？他在受迎接后有任何辜负那迎接的行为吗？绝不。一切都是和谐的，如同一把竖琴上的旋律，同一调子；他的生活、教导、奋斗、危险、回归，以及回归后的行为，都是如此。因为一回到他的教会，他并不像那些被无节制的激情蒙蔽的人，在愤怒支配下，立即推开或击打任何迎面而来的事物，即使它是可以宽恕的。但他认为这是他维护名誉的特别时刻，因为受虐待的人通常克制，而有能力报复的人则不受控制。他对待那些伤害过他的人如此温和、仁慈，以至于他们自己，若我可这样说，也对他的回归不感到厌恶。\par

31. 他洁净了那使天主成为商品、贩卖基督之事的圣殿，在此也效法基督\BibleRef{若 2:15}；只是这事是用劝服的话语，而非用编成的鞭子完成的。他也使那些不和的人彼此和好，也与他和好，不需任何助手。他使受委屈者摆脱压迫，不分他们是自己一方还是对立一方。他还恢复了被推翻的教导：天主圣三再次被勇敢地宣讲，被放在灯台上，以独一天主性的光辉照射进所有人的灵魂。他为整个世界立法，并通过信件、邀请、教导那些不请自来的人，使所有人的心思受他影响，将唯一的法律——\OriginalTerm{善意}{Good will}——颁布给所有人。因为只有这善意能引导他们达到真正的结局。简言之，他体现了两种著名石头的德行——对攻击他的人，他是金刚石；对不和的人，他是磁铁，借着某种隐秘的自然力量吸引铁到自己，并影响最坚硬的物质。\par

32. 然而，嫉妒不可能容忍这一切，或看到教会因分裂创伤迅速愈合（如身体一样）而恢复到昔日的光荣和健康。因此，它激起了皇帝反对圣亚大纳削，一个反叛者，和他一样卑鄙，只是因时间不足而稍逊一筹，他是第一位狂怒反对基督的基督徒皇帝，在他得到机会时，突然生出了他长期怀胎的不虔敬之蛇怪，并在他被宣告为皇帝时，显明自己是背叛托付给他帝国的皇帝，更是双倍背叛拯救他的天主。他通过将诡诈与残忍混合，设计了最不人道的迫害，嫉妒殉道者在奋斗中赢得的荣誉；因着猜疑他们勇敢的名声，他把文字诡辩作为他的特征，或者说实话，他因本性热切地投身于此，模仿住在他内心的恶者的各种诡计。他以为征服整个基督徒种族是简单的任务；但发现战胜圣亚大纳削和他对我们教导的力量是巨大的。因为他看到，由于此人的抵抗和反对，反对我们的计谋无法成功；被剪除的基督徒地位迅速被填补，虽然看似惊人，却因外邦人的加入和圣亚大纳削的谨慎。因此，面对这一切，这位狡猾的败坏者和迫害者，不再隐藏他鄙俗的诡辩外衣，暴露了他的邪恶，公开将主教从城中放逐。因为这位杰出的战士必须通过三次斗争而得胜，从而成就他荣耀的完满称号。\par

33. 间隔短暂，公义便宣告判决，将冒犯者交与波斯人：差遣他是一位野心勃勃的君主——带回的却是一具无人怜悯的尸体；我听说，这尸体未得安葬于坟墓，反被它所震撼的大地摇出、抛起——我以为，这是他未来惩罚的预兆。\newline{}随后另一位王兴起，\BibleRef{出 1:8}不像前者那样面容无耻，也不以苛役和督工压迫以色列，而是最为虔诚、温和。他为帝国奠定最佳根基，以正义之举作为开端，合乎义理：他召回所有被放逐的主教，首当其冲的是那位德性卓绝、显著捍卫虔诚事业的人。此外，他查问我们那曾被撕裂、混乱、被多人分割为各种意见和部分的信仰真相；意图若有可能，藉圣神的合作使普世归于和谐与合一；若不能做到，则依附于最佳的一方，以便互相扶持，由此显示他在重大问题上思想的极度崇高与宏伟。在此，也极其显著地展现了\Person{亚大纳削}的纯朴，以及他对基督信仰的坚定。因为，当所有与我们同心的人分裂为三派，许多人在对圣子的概念上摇摆不定，更多人则在对圣神的概念上摇摆不定（在此点上稍有偏差即算正统），而在两点上都正确的人极少时，他率先、独自或仅与少数人共识，敢于以完全清晰明确的方式书写宣认三个位格的神性与本质的统一，并因此在后来，在圣神的感动下，达到了如同早期多数教父在圣子方面所得的、对圣神的同一信仰。这一宣认，一份真正崇高、宏伟的礼物，他呈献给了皇帝，以正统信仰的书面记述对抗不成文的革新，使一个皇帝被另一个皇帝战胜，理性被理性征服，论著被论著征服。\par

34. 这一宣认，看来受到西方和东方所有有识之士的尊重；有些人将虔诚埋藏于心中——如果我们相信他们所言——却不再前进，如同胎死腹中的婴儿；另一些人则以某种方式点燃了火花，足以躲避当时的困境，这火花或来自正统教徒中更热忱者，或来自民众的虔敬；还有些人大胆宣讲真理，我愿站在他们一边，因为我不敢再夸耀；我既不再顾及自己的怯懦——换言之，即比我更不健全之人的观点（在这方面我们已做得过多，既未从他人处获益，也未保护好自己的，如同恶劣的管家一般），而是将我的后代公诸于世，急切地滋养它，并随着它不断成长，将它展现于众人眼前。\par

35. 然而，这还不及他的行为令人钦佩。他既已为真理实际冒险，又何必以书面形式宣认？尽管如此，我仍将这一点补充到已说过的内容中，因为它在我看来特别奇妙，在当前这般纷争频发的时代，不可不受惩罚地忽略。因为他的行动，若我们留意，甚至会对今天的人有所教诲。如同在同一量度的水中，不仅有手汲取时残留在后的余沥，还有曾握于手中、从指缝间滴落的水滴；同样，在我们与那些不仅因不虔诚而疏离者之间，也与那些最虔诚者之间，存在着分别，既关乎那些无关紧要的教义（次要之事），也关乎意在表达同一含义的措辞。我们以正统意义使用\Quote{一个本质}与\Quote{三个位格}二词，前者指天主性的本性，后者指三位的特性；意大利人意思相同，但因词汇贫乏，无法区分本质与位格，故引入\Quote{位格}一词，以免被视为主张三个本质。结果，若非可悲，实为可笑。这细微的语言差异竟被视为信仰差异。于是，在\Quote{三位}的教义中怀疑撒伯里乌主义，在\Quote{三个位格}的教义中怀疑亚略主义，两者皆出于争辩精神。而后，随着争执（争辩的必然结果）逐渐而持续地加剧，整个世界有因音节之争而分裂的危险。我们这位蒙福之人，作为真正的天主之人、灵魂的伟大管家，看到并听到这一切，感到有责任不能忽略这种对话语如此荒谬无理的撕裂，便对疾病施以药物。他如何做呢？他以温和同情的态度与双方交谈，仔细权衡他们措辞的含义，发现其意义相同，教义毫无差异，便允许各方使用自己的术语，从而将他们团结在行动的统一之中。\par

36. 这本身比所有作者详述的漫长劳作与教诲更为有益，因为其中掺杂了些许野心，因而表达中也许带有些许新奇。这又比他的多次守夜与苦行更有价值，因为那些仅对实行者有益。这配得上我们英豪著名的流放与逃亡；因为他选择忍受这般苦难所追求的目标，在苦难过后仍继续追寻。他也没有停止在他人心中培育同样的热忱，赞扬一些人，温和地责备另一些人；激发慵懒者，约束激情者；有些情形急于防止跌倒，另一些情形则筹划跌倒后的恢复办法；性情淳朴，治理手法多样；辩论机敏，心智更机敏；俯就卑微者，超越高贵者；好客，保护投靠者，驱除灾祸，真正独自一人兼有希腊人分配给诸神的所有特质。此外，他既是已婚身分也是童贞身分的保护者，既和平又缔造和平，并陪伴那些由此世离去的人。啊，若我详述其多方面的卓越，他的德行为我提供了多少称号！\par

37. 在如此作为学徒与教师的历程之后，他的生活与习惯成为主教职的典范，他的教导成为正统信仰的法律，他为虔诚赢得了什么报酬？确实不应略过此事。他在美好的高龄结束了生命，被归到他的祖先——先祖、先知、宗徒和殉道者那里，他们曾为真理而争战。简而言之，在我为之撰写的墓志铭中，他离世时的荣耀甚至超过了他从流放归来时的荣耀；令多人流泪，他的荣光，储存在众人心中，远超一切可见的标记。然而，亲爱的圣人啊，你自身以你美好的声誉，特别展示了言语与沉默的恰当比例，请在此止住我的言语，它们虽已耗尽我所有能力，但远不及你应得的赞美。愿你从上面向我们投以仁慈的目光，引领这百姓完美地朝拜完美的天主圣三，我们默观并朝拜祂——圣父、圣子、圣神。愿你在我的命运若为平安时，拥有并助佑我的牧职；若经过挣扎，则扶持我，或提携我到你身旁，将我与你自己及类似你的人（虽然我求的是一件大事）安放在基督自己、我们的主内，愿一切光荣、尊威和权能归于祂，直到永远。阿们。\par

\SourceRecord{Translated by Charles Gordon Browne and James Edward Swallow.；Nicene and Post-Nicene Fathers, Second Series；Vol. 7.；Edited by Philip Schaff and Henry Wace.；Buffalo, NY: Christian Literature Publishing Co.,；1894.；<http://www.newadvent.org/fathers/310221.htm>.}

% structure: p0001 source=h1 target=section reason=page-title

\section{第一神学演讲（演讲第27篇）}

\Emph{第一神学演讲。}\par

\Emph{反驳欧诺米派之预备演说。}\par

一、 我要向那些以口才自夸的人发言；所以，先引用一句经文：\Quote{看，骄傲的人哪，我与你为敌，} \BibleRef{耶 50:31} 这不仅是针对你们的教导体系，也是针对你们的听闻和思想态度。因为有些人不仅耳朵 \BibleRef{弟后 4:3} 和舌头，甚至如我所见，连他们的手，也对我们的言语发痒；他们喜爱世俗的虚谈和所谓知识的错谬对立，以及无益的言辞争辩；因为\Person{保禄}，这位宣讲和建立\Quote{简短的言语} \BibleRef{罗 9:28} 的传道者，也是渔夫的弟子和老师，称一切过分或多余的言论为虚浮。至于我们提及的那些人，但愿他们那如此流畅而机巧地运用高尚和优美言辞的舌头，能同样关注实际行动。因为这样，也许过不多久，他们就会少一些诡辩，少一些愚蠢而奇怪的言辞杂耍——若允许我用一个可笑的说法来描述一件可笑的事。\par

二、 既然他们忽略一切正义之路，只关注这一点：即他们应当解开或绑住所提交的哪一个命题（如同那些在剧院里公开表演摔跤比赛的人，但那并非按照运动规则取胜的摔跤，而是一种欺骗无知者眼目、博取掌声的摔跤），每一个市场都必须因他们的谈论而嗡嗡作响；每一场宴席都被愚蠢的谈话和无聊折磨至死；每一个节日都变得不快乐且充满沮丧，每一个哀悼的场合都被更大的灾祸——他们的疑问——所慰藉；所有习惯于简朴的女眷住所都陷入混乱，并因他们言语的洪流而失去谦逊的花朵……既然，我说，情况如此，这邪恶无法容忍、不可承受，我们的伟大奥迹有被轻视的危险。那么，让这些窥探者容忍我们吧，我们怀着慈父般的同情，如圣\Person{耶肋米亚}所说，心肠撕裂； \BibleRef{耶 4:19} 让他们容忍我们，只要不野蛮地拒绝我们关于这个主题的演说；并且，如果他们能够，让他们暂时抑制舌头，倾听我们。无论如何，你们不会遭受任何损失。因为我们要么说给那些愿意听的人 \BibleRef{德 25:9}，我们的言语将结出某种果实，即对你们的益处（因为播种者将圣言撒在各种心田里；而好而肥沃的心田结出果实），要么你们会离开，像鄙弃其他讲论一样鄙弃我们的这次演说，并从中获取更多反对和辱骂我们的材料，从而更加自娱自乐。\par

如果我说出一种对你们而言陌生且违背你们习惯的语言，你们不必惊讶——你们自称无所不知，并以一种过于急躁和慷慨的方式教导一切……以免使你们痛苦，我说是无知和鲁莽。\par

三、 并非人人，我的朋友们，都适合谈论天主；并非人人；这主题并非如此廉价和低贱；我还要补充，并非在任何听众面前，也非在任何时候，也非在所有问题上；而是要在特定的场合、在特定的人面前、在特定的限度内。\par

并非人人，因为只有那些经过考验、在默想中纯熟、并且在灵魂和身体上事先被净化的人，或者至少正在被净化的人，才被允许。因为不洁者接触纯洁者，我们完全可以肯定，是不安全的，正如让虚弱的眼睛注视阳光是不安全的。那么，允许的场合是什么？就是我们摆脱了一切外在的污秽或扰乱，并且我们内在的统治者不被烦恼或错误的形象所混淆时；就像那些把好字与坏字混在一起，或将污秽混入香膏的芬芳中的人。因为真正闲下来才能认识天主；当我们能找到一个合适时机时，才能辨明神圣之事的直路。那么，哪些人是被允许的呢？是那些真正关心这主题的人，而不是那些像对待其他任何事物一样，在赛马、戏剧、音乐会、宴会，甚至更低下的活动之后，将其作为愉快闲聊的人。对这样的人来说，关于这些主题的轻浮玩笑和漂亮矛盾是他们娱乐的一部分。\par

四、 其次，我们可以在什么主题上、以及到什么程度进行哲学探讨？在我们可以触及的问题上，并且达到听众的智力和领悟力所能延伸的程度。不可超越，否则，正如过大的声音损害听觉，或过量的食物损害身体，或者，如果你愿意，正如超出承受能力的过重负担损害背负它们的人，或过量的雨水损害大地；同样，这些论点，被僵硬——如果我可以这样表达——所压迫和超载，甚至会在他们原本拥有的力量上遭受损失。\par

五、我并非说不必时常记念天主；……不可误解我，否则那些敏捷灵巧的人又要攻击我了。因为我们应该思念天主，甚至比呼吸还要频繁；若可这样说，我们不应做别的。是的，我完全赞同那命令我们昼夜默想，在早晚和中午述说，并时时赞美主的圣言；或用梅瑟的话说，无论人躺卧、起来、行路或做任何事——\BibleRef{申 6:7}——藉此记念，我们得以塑造为纯洁。所以，我不阻止不断地记念天主，只阻止谈论天主；也不是说谈论本身是错的，只是在不合时宜时；也不是全部教导，只是不知节制。因为即使是蜜，过饱过腻也会令人呕吐；\BibleRef{箴 25:16}而且，如撒罗满所说，我也认为凡事都有定时，\BibleRef{训 3:1}那好的若不以好的方式去做，就不再是好；正如冬季的花完全不当时令，男人的衣服不适于女人，女人的衣服不适于男人；又如哀悼时不适合谈论几何，宴乐时也不适合流泪；难道我们唯独在这件事上不顾及合宜之时，而这件事最应当尊重时机吗？决不可，我的朋友们和弟兄们（我仍要称你们为弟兄，尽管你们的行为不像弟兄）。让我们不要这样想，也不要像性急的烈马，摆脱骑手理性，抛弃使我们守节的敬畏，远远地跑离转折点，而要按自己的界限从事哲学探讨，不要被带到埃及，也不要被冲到亚述，\BibleRef{达 3:12}也不要在异地唱主的歌，我指的是在任何听众面前，无论陌生或亲属、敌对或友善、善意或相反，他们过分留意我们的行为，喜欢将我们的过错火花燃成火焰，暗暗点燃、煽动，用呼吸把它升到天上，比那吞没四周万物的巴比伦火焰更高。因为他们的力量不在于自己的教义，而是从我们的弱点中寻找。因此他们专注于我们的——我该说\Quote{不幸}还是\Quote{过失}？——如同苍蝇叮伤口。但至少我们不要再不认识自己，或对这类事情的适当秩序过于忽视。若不能结束现有的敌意，至少让我们达成一致：我们对奥秘要低声细语，以神圣的方式说神圣之事，不将不可说的投向亵渎的耳朵，也不证明我们不如崇拜魔鬼、事奉耻辱神话和事迹的人庄重；因为他们宁可把自己的血给未入教者，也不愿给某些话语。但要承认，正如衣着、饮食、欢笑和举止有某种礼节，言语和静默也有礼节；因为在天主的诸多称号和权能中，我们最尊崇圣言。因此，让我们的争论也保持适度。\par

六、为什么一个怀着敌意听这些话的人可以被允许听到关于天主的受生、或他的创造、或天主如何从虚无中造成、或分剖、分析、分割？为什么我们使控告者成为法官？为什么把刀剑交给敌人？你想，一个赞成奸淫、娈童、崇拜情欲、不能想象高于身体之事的人……直到最近还为自己设立神明，且是些以最卑鄙行为著称的神明——他会以怎样的心态和情绪接受这些论题？难道不是首先从物质的立场，可耻地、无知地，按他惯常的意义来理解吗？他不会把你的神学当作他自己神明和情欲的辩护吗？因为我们若自己轻率滥用这些词语，要说服他们接受我们的哲学便需要很长时间。而他们若自身是邪恶事物的发明者，当这些被提供给他们时，他们怎能不抓住呢？这样的结果来自我们彼此的争竞。这样的结果临到那些为圣言争斗而超越圣言所许的人；他们如同疯子，放火烧自己的房子，撕裂自己的儿女，或否认自己的父母，把他们当作陌生人。\par

七、当我们把这些不相干的人从谈话中赶走，并将那大军逐往深渊，赶入猪群（\BibleRef{路 8:31}）之后，接下来就要审视自己，像雕像一样美化我们的神学自我。首先要考虑的是：这滔滔不绝的争论和无休止的言谈究竟是什么？这贪得无厌的新毛病是什么？为什么我们捆住了双手，武装了舌头？我们不赞美好客，也不赞美兄弟之爱，也不赞美夫妻之情，也不赞美守贞；我们也不钦佩对穷人的慷慨，或吟咏圣咏，或彻夜守夜，或流泪。我们不以斋戒克制身体，也不以祈祷走向天主；我们也不使低劣的服从于更好的——我指尘土服从于精神——正如那些对我们复合本性作出公正判断的人所该做的那样；我们不使生活成为死亡的准备；我们也不做情欲的主人，不忘我们天上的尊贵；也不驯服膨胀和狂怒的愤怒，也不驯服导致跌倒的骄傲，也不驯服无理忧愁，也不驯服放纵的快乐，也不驯服轻浮的欢笑，也不驯服无节制的眼睛，也不驯服贪婪的耳朵，也不驯服多言，也不驯服荒谬的思想，也不驯服任何那恶者从我们内在根源中用来攻击我们的事物；给我们带来从窗户进来的死亡（\BibleRef{耶 9:21}），正如圣经所说，即通过感官。恰恰相反，我们做了完全相反的事，并给予他人情欲自由，就像君王在庆祝胜利时发布服役豁免一样，只要求他们倾向我们这边，并且以更大胆或更不虔敬的方式攻击天主。我们给予他们邪恶的报酬，为了一个不善之事——以口舌的自由换取他们的不虔敬。\par

八、然而，啊，多言的辩证者，我要问你一个小小的疑问（\BibleRef{约 38:3}），你要回答我，正如祂藉着旋风与云雾发出神圣训诫时对约伯所说的（\BibleRef{约 38:1}）。在天主的家里有许多住处，正如你所听到的，还是只有一个？你当然会承认有许多，而非只有一个。那么，它们都要被填满，还是只有一些被填满，而其他不被填满，以致有些住处将空着，并且被白白准备？当然所有住处都将被填满，因为天主所做的一切都不会是徒然的。你能告诉我你认为这住处是什么？是那里为有福者存留的安息与荣耀，还是别的东西？——不，不是别的。既然我们在这一点上一致，让我们进一步考察另一点。有没有什么东西能获得这些住处，如我所想，还是没有？——当然有。——那是什么？岂不是有各种行为方式和各种目标，一个引向一条路，另一个引向另一条路，按照信德的比例，这些我们称为\Quote{道路}？那么，我们必须走所有道路，还是其中一些……同一个人走所有这些道路，如果可能的话；如果不可能，则走尽可能多的路；或者只走其中一些？即使这也不可能，那么，至少在我看来，出色地走完一条路也是一件大事。——\Quote{你的想法是正确的。}——那么，当你听到只有一条路，而且是狭窄的路（\BibleRef{玛 7:14}）时，这个词向你表明了什么？是因着它的卓越而只有一条。因为它是一条，即使被分成了许多部分。它狭窄是因为它的艰难，也因为走这条路的人少，与众多敌手和行恶道的人相比。——\Quote{我也这样想。}——那么，我的好朋友，既然是这样，你为什么好像因某种贫乏而指责我们的教义，却一头扎进那条你所谓论证和思辨，而我所谓轻浮之谈和骗术的路呢？让保禄用那些严厉的责备来责备你吧，他在列举恩宠的恩赐之后说：\Quote{难道都是宗徒么？难道都是先知么？等等}（\BibleRef{格前 12:29}）。\par

IX. 但是，就算如此。你高不可及，甚至超出高远，超出云霄，如果你愿意，成为无形之物的观察者，不可言说之事的聆听者；你已像\Person{厄里亚}一样上升，像\Person{梅瑟}一样蒙受能见天主的光荣，像\Person{保禄}一样被提到天上；你为什么要在一天之内就把你的同伴们塑造成圣人，任命他们为神学家，仿佛向他们吹入教导，使他们成为许多无知神谕的会议？你为什么用你的蛛网缠绕那些软弱的人，如果这被视为伟大而智慧的事？你为什么激起黄蜂的巢穴来攻击信仰？你为什么突然向我们释放辩论的洪流，如同古老寓言中的巨人？你为什么收集人中间一切轻浮和懦弱的人，像乌合之众一样汇成一股洪流，又用奉承使他们更加柔弱，打造了一个新的作坊，巧妙地利用他们的无知为自己收获？你否认这样的事吗？其他事情对你无关紧要吗？你的舌头必须不惜代价地统治，你不能抑制你言语的阵痛吗？你可以找到许多其他值得讨论的高尚题目。把你的这病态转向这些，更有益。攻击\Person{毕达哥拉斯}的沉默，以及\OriginalTerm{俄耳甫斯的豆子}{Orphic beans}，以及关于\Quote{夫子说}的新奇夸耀。攻击\OriginalTerm{柏拉图的理念}{ideas of Plato}，以及我们灵魂的\OriginalTerm{转世}{transmigrations}和历程，以及\OriginalTerm{回忆}{reminiscences}，以及灵魂对可爱身体的可憎之爱。攻击\Person{伊壁鸠鲁}的无神论和他的\OriginalTerm{原子}{atoms}，以及他的非哲学的\OriginalTerm{快乐}{pleasure}；或者\Person{亚里士多德}的\OriginalTerm{渺小眷顾}{petty Providence}，他的牵强体系，以及关于灵魂可朽的论述，以及他人道主义的教义。攻击\OriginalTerm{斯多亚}{Stoa}的傲慢，或\OriginalTerm{犬儒}{Cynic}的贪婪和粗俗。攻击\Quote{\OriginalTerm{虚空与充实}{Void and Full}}（多么无稽），以及关于诸神、祭献、偶像和魔鬼的一切细节，无论是善意的还是恶意的，以及人们用占卜、召唤神或灵魂、星辰之力所做的所有诡计。如果这些在你看来不值得讨论，因为琐碎且已多次被驳斥，而你将坚持你的路线，并在其中寻求你野心的满足；那么在这里我也将为你提供广阔的道路。哲学讨论关于世界或诸世界；关于物质；关于灵魂；关于有理性的本性，善或恶；关于复活，关于审判，关于赏报，或基督的苦难。因为这些主题，击中目标并非无用，偏离也非危险。但我们将与天主交谈，今生仅有一点；但不稍久后，或许更完全地，在同一个我们的主耶稣基督内，愿荣耀归于祂，直到永远。阿们。\par

\SourceRecord{Translated by Charles Gordon Browne and James Edward Swallow.；Nicene and Post-Nicene Fathers, Second Series；Vol. 7.；Edited by Philip Schaff and Henry Wace.；Buffalo, NY: Christian Literature Publishing Co.,；1894.；<http://www.newadvent.org/fathers/310227.htm>.}

% structure: p0001 source=h1 target=section reason=page-title

\section{第二神学演讲（演讲28）}

\Emph{第二神学演讲。}\par

一、在上一篇演讲中，我们清楚地规定了神学家应有的品格、他应从事哲学思考的对象、时机与限度。我们指出：他应当尽可能纯洁，好让光明被光明所领悟；他应当与庄重之人交往，以免他的话因落于不毛之地而徒劳；合适的时机是当我们内心从外界纷扰中得享平静时，免得像疯子一样气喘吁吁；我们可达到的程度，乃是我们自己已经达到或正在前进的程度。既然这些事如此，我们已为自己破开了神性的荒地（\BibleRef{耶 4:3}），以免将种子撒在荆棘中（\BibleRef{玛 13:7}），并且平整了地面（\BibleRef{依 28:25}），借着圣经塑造自己并塑造他人……现在让我们进入神学问题，将我们的论述的基础——圣父、圣子、圣神——置于首位，因为我们所要论述的正是他们；愿圣父喜悦，圣子帮助我们，圣神启迪我们；更确切地说，愿那从唯一的天主而来的独一光明，在多样性中统一，在统一中多样，其中包含着奇妙的奥秘，降临于我们身上。\par

二、当我急切地登上那山（\BibleRef{出 24:1}）——用更真切的话说，当我既急切渴望，同时又害怕（一出于盼望，一出于软弱）进入云中，与天主交谈，因为天主如此命令；若有谁是亚郎，让他与我一同上去，让他站在近处，若有必要，就留在云外。但若有谁是纳达布或阿彼胡，或是长老品级的人，他固然可以上去，但要按他净化的程度远远站着。但若有谁属于群众，不配此默观高度，他若全然不洁，就完全不可靠近，因这对他有危险；但他若至少暂时净化了，就让他留在下面，只聆听那声音、那号角，即虔敬的单纯言语，并看见那山冒烟发光，令不能上去的人既恐惧又惊奇。但若有谁是邪恶凶猛的野兽，全然不能领会默观与神学的内容，让他不要恶毒地潜伏在树林中的洞穴里，猛然扑向某个教义或言论，用歪曲来撕裂健全的道理，而要让他远远站着，离开那山，否则他会被石头砸死、压碎，且在他邪恶中悲惨灭亡。因为对像野兽的人来说，真实而健全的言论就是石头。他若是一只豹子，就让他带着斑点死去（\BibleRef{耶 13:23}）；若是一只吼叫的狮子，寻找可吞噬我们灵魂或言语的（\BibleRef{伯前 5:8}）；或是一只野猪，践踏真理的珍贵透明珍珠（\BibleRef{玛 7:6}）；或是一只外来的阿拉伯豺狼，或在辩论诡计上比这些更锋利的；或是一只狐狸，即背信弃义的灵魂，根据环境或需要改变形态，以死尸或腐肉为食，或在逃过大葡萄园后偷食小葡萄园；或是任何其他被法律视为不洁、不可用于食物或娱乐的肉食野兽；我们的言论必须远离这些，而被刻在坚固的石版上，且是双面的，因为法律一部分可见，一部分隐藏；一部分属于留在下面的大众，另一部分属于少数向上攀登那山的人。\par

三、朋友们、同入教者、真理的共同爱好者，我遭遇了何事？我奔跑去把握天主，就这样上了山，拉开云幕，远离物质与物质事物，尽我所能退入自身。然后当我仰望时，我勉强看见了天主的背面（\BibleRef{出 33:23}），尽管我被那磐石——为我们而成为肉身的圣言——所遮蔽。当我稍加细看时，我看到的不是那第一且纯粹的本性，它认识自身——我是指圣三；不是那停留在第一层帐幔内、被革鲁宾隐藏的本性；而只是那最终甚至延伸到我们的本性。就我所知，那是威严，或如圣达味所称的，彰显在受造物中的荣耀，即祂所产生并治理的。因为这些都是天主的背面，祂留下的犹如祂自身的影子，如同太阳在水中的倒影与映像，因我们无法直视太阳本身，因祂的纯粹光明超越了我们的知觉能力。这样，你就当谈论天主；即使你是梅瑟，是法郎的天主（\BibleRef{出 4:2}）；即使你像保禄一样被提到第三层天（\BibleRef{格后 12:2}），听见了不可言说的话语；即使你超越他们二者，被提升到天使或总领天使的地位与尊严。因为纵然一件事物是完全属天的，或超越诸天，其本性更高，且比我们更接近天主，但它距离天主及对其本性的完全领悟，仍比它高出我们这复合、卑下、下沉于尘世的构造的程度更为遥远。\par

四．因此我们必须重新开始。如同一位希腊神学教师——依我看来并非不熟练——所教导的：构思天主是困难的，但用言语定义祂却是不可能的；他这样说，意图使人以为他已领悟了天主，同时又可因无法表达领悟而免于被指为无知。但依我之见，表达祂是不可能的，而构思祂则更为不可能。因为可以构思的，或许能借着语言，对并非完全耳聋或理解迟钝的人，说得相当清楚，至少不完美地表明出来。但要完全领悟如此伟大的对象，是完全不可能、不可行的，不仅对完全疏忽和愚昧的人，就是对那些高度提升、热爱天主的人，以及一切受造的本性，也是如此；因为此世的黑暗和肉体的厚障阻碍了对真理的圆满理解。我不知道对于更高的本性和更纯洁的理智来说是否相同，他们因亲近天主并蒙受祂全部光照，可能看到，即使不是全部，至少比我们更完美、更清晰；有些或许更多，有些更少，按他们的等级而定。\par

五．关于这一点，说得够多了。关于我们的事，不仅天主的平安\BibleRef{斐 4:7}超越一切理解和知识，也不仅天主在许诺中为义人储备的、\Quote{眼睛未曾看见，耳朵未曾听见，人心也未曾想到}的事（除了一点点），也不仅是对受造物的准确认识。因为即使关于这一点，我愿你明白，当你听到这些话：\Quote{我要观看你手指所造的天，并你所陈设的月亮星宿}及其中的固定秩序时，你只得到一个影子；好像他现在不是在观看，而是注定将来要观看。但远在这些之上的是那超越它们的本性，万有由此而出，那不可理解、不可限量的——我指的不是祂的存在这一事实，而是祂的本性。因为我们的宣讲不是空洞的，我们的信德也不是虚妄的\BibleRef{格前 15:19}，这也不是我们所宣扬的道理；因为不愿你拿我们坦诚的陈述作为开端，狡猾地否认天主，或因承认无知而傲慢。因为确信一件事存在，与知道它是什么，是两回事。\par

六．如今，我们的眼睛和自然律教导我们：天主存在，祂是万有的创造者和维持者。眼睛，因为它们落在可见之物上，看见它们在美善的稳定和进程中，仿佛不动地运动、旋转着；自然律，因为通过这些可见之物及其秩序，它追溯至它们的作者。因为宇宙怎能生成或组织起来，除非天主召它存在并维持它？因为每一个看见一把制作精美的琴，并思考它如何巧妙地装配和安排，或听见它的旋律的人，只会想到制琴者或弹琴者，并在心中转向他，即使未曾亲眼见过他。同样，向我们也显示出那创造、推动并保存一切受造物的那一位，即使他不被理智领悟。非常缺乏理性的人，是那些不愿自愿在这点上追随自然证明的人；但即使是我们想象、塑造或理性为我们描绘的，也不能证明天主的存在。但如果某人甚至在一定程度上领悟了这一点，天主的存有又如何能证明？谁曾达到这智慧的极限？谁曾配得如此伟大的恩赐？谁曾张开他心智的口，吸入了圣神，以便借着那洞察一切，连天主的深奥也洞察的圣神\BibleRef{格前 2:10}来接纳天主，不再需要进步，因为他已经拥有了欲望的终极对象，以及一切社会生活和最优秀之人的一切智慧所奔赴的目标？\par

七．因为如果你依赖理性的所有近似，你将把神体构思成什么？理性又将把你带往何处，最富哲思的人，最好的神学家，以熟悉无限者自夸的人？祂是身体吗？那么祂如何是无限、无量、无形、不可触摸、不可见的？这些是身体的特征吗？多么傲慢！因为这不是身体的本性！或者你会说祂有身体，但没有这些特征？何等愚蠢，神体竟不拥有比我们更多的东西！因为如果祂被限制，祂如何是崇拜的对象？或者祂如何能免于由元素组成，因而可以分解为元素，甚至完全消散？因为每一复合物都是争端的起点，争端导致分离，分离导致解体。但解体完全与天主和第一本性无关。因此不能有分离，以免有解体；不能有争端，以免有分离；不能有复合，以免有争端。因而也必须没有身体，以免有复合。如此，论证从后往前建立起来了。\par

VIII. 然而，我们又当如何保留天主渗透万有并充满万有的真理——如经上所载：\Quote{我岂不充满天地？——上主说，}\BibleRef{耶 23:24}以及\Quote{上主的神充满了世界，}\BibleRef{智 1:7}——如果天主部分包含、部分被包含呢？因为祂要么占据一个空无的宇宙，那样万有对我们而言便消失了，结果是我们因使天主成为一身体而侮辱了祂，并夺走了祂所造的一切；否则祂将是一个被包含在其他身体内的身体——这是不可能的；或者祂被包裹在它们之中，或与它们相混合，如同液体混合一样，一个分开另一个、又被另一个分开——这种观点比\Person{伊璧鸠鲁}的原子论还要荒谬和愚妄；因此这个关于身体的论证将落空，毫无实质基础。但若我们要主张祂是非物质的（例如某些人所想象的\OriginalTerm{第五元素}{Fifth Element}），并且祂在圆周运动中回旋……我们就假设祂是非物质的，且祂是第五元素；如果他们愿意，也让祂依其论证的独立倾向与安排而无形体；因为目前我在这一点上不与他们争辩；那么，祂在何种意义上会是那些运动和骚动的事物之一？更不必说使造物主与受造物遭受同样的运动，以及那位承载万有的（如果他们连这一点也允许）与祂所承载者成为一体，其中所含的侮辱。再者，什么是推动你的第五元素的力量？什么是推动万有的力量？什么推动那力量？那推动的力量又是什么？如此这般，以至\Emph{无穷}。如果祂服从于运动，祂怎能不完全被空间所包含？但如果他们主张祂是不同于此第五元素的某物；假设他们将一种天使性体归于祂，他们将如何证明天使是有形体的，或他们有何种身体？在这种情况下，天使所服务的天主又能在多大程度上超越天使？而如果祂高于他们，就又引入了一堆不合理的身体和深度的胡言，毫无立足之地。\par

IX. 因此，我们看到天主并不是一个身体。因为没有一位受默感的教师曾主张或承认这种观念，我们法庭的判决也不允许。因此，剩下的只有将祂设想为无形体的。但\Quote{无形体}这一术语，即使被许可，也不能向我们呈明或包含祂的本质，正如\Quote{无始}、\Quote{无源}、\Quote{不变}、\Quote{不朽}或任何其他用于天主或与天主相关的谓词一样。因为祂无始、不变、不受限制，这对祂的存在或实体产生了什么效果？不，祂存在的全部问题仍需留待那真正拥有天主心思并在默观中精进者进一步思考和阐述。因为正如说\Quote{它是一个身体}或\Quote{它是受生的}，并不足以清楚地向心灵呈现这些谓词所指的各种对象，你还必须表达你所使用它们的主体，如果你要清晰而充分地呈现你的思想对象（因为这些谓词——有形体的、受生的、有死的——每一个都可用于人、牛或马）。同样，那热切追求自有者本性的人，不会止步于说祂\Emph{不是}什么，而必须超越祂不是的，说出祂\Emph{是}什么；因为把握某一点，比无休止地逐一否定、以无穷细节否认更易，从而通过排除否定和肯定肯定来达到对这一主题的理解。\par

但一个人只陈述天主不是什么，而不进而说祂是什么，他的行为与以下情况大致相同：当被问及二乘以五是多少时，回答：\Quote{不是二、不是三、不是四、不是五、不是二十、不是三十，简言之，不是十以下的任何数，也不是十的倍数；}却不肯回答\Quote{十}，也不愿将提问者的心思稳定在答案的坚实基础上。因为通过事物是什么来显示它不是什么，远比通过剥除它不是什么来证明它是什么要容易和简洁得多。这一点，想必对每个人都是显而易见的。\par

X. 既然我们已经确知天主是无形体的，让我们再深入考察一下。祂是无处还是处于某处？如果祂无处，那么某个非常好问的人可能会问：祂如何甚至能够存在？因为不存在的东西是无处的，那么无处的东西或许也是不存在的。但如果祂处于某处，祂必然要么在宇宙之内，要么在宇宙之上。如果祂\Emph{在}宇宙之内，那么祂必然要么在部分之中，要么在整体之中。如果在部分之中，那么祂将被那小于祂的部分所限制；但如果在整体之中（即无处不在），则被那更远更大的——我指的是包含特殊者在内的宇宙整体——所限制；如果宇宙要被宇宙所包含，就没有任何地方能免于限制。如果祂被包含在宇宙之内，就会得出这一点。此外，在宇宙被创造之前祂在哪里？因为这是一个不小的难题。但如果祂在宇宙\Emph{之上}，那么有什么能区分这与宇宙呢？这\Quote{之上}又位于何处？而这种超越与被超越者如何在思想中被区分，假设没有界限来划分和界定它们？难道不需要某个中间物来标记宇宙与宇宙之上者的界限吗？这若不是空间又是什么？而我们已经否定了空间。因为我还没有提出这一点：如果天主甚至在思想中可被理解，祂就会完全被限制；因为理解是限制的一种形式。\par

十一、现在，我为什么讲了这一切？或许对于大多数人来说，这样过于细致，也符合当今轻视高贵简朴、引入曲折复杂风格的演说方式。这是为了使树由果实可以认出（见：\BibleRef{路 6:44}），我是说，为使在这样的教导中运行的黑暗，由论点的隐晦被认出。我这样做的目的，并非为了借解结和解难（这是\Person{达尼尔}伟大的奇妙恩赐）而为自己博得惊人言论或过度智慧的声誉，而是为了阐明我的论点从一开始就瞄准的要点。这要点是什么呢？那就是：天主性不能为人的理性所领悟，我们甚至不能完全想象其伟大。这并非出于嫉妒，因为嫉妒远离天主性——它是无情的、唯一的善和万有的主；尤其不会嫉妒祂所有受造物中最尊贵的。因为圣言更喜欢什么胜于有理性、会言语的受造物呢？为什么？甚至他们的存在本身就是祂至高美善的证明。这种不可理解性也不是为了祂自己的光荣和尊威——祂是充满的（见：\BibleRef{依 1:11}）——仿佛祂拥有光荣和威严取决于不能接近祂。因为这是极其诡辩且与品格不相称的，我不说与天主，即使与任何一个有适度善意、对自己有正确认识的人，为寻求自己的至高地位而给别人设置障碍。\par

十二、但是否还有其它原因，让那些更亲近天主、亲眼目睹并注视祂不可探测的判断的人来看吧（见：\BibleRef{罗 11:33}），如果有人是如此杰出于美德，并行走于无限者的道路中，按俗话说。然而，就我们这些以小小尺度衡量难解之事的人所达到的而言，也许一个原因是防止我们因为太容易获得而轻易丢掉所有物。因为人对于辛苦得来的东西抓得很紧；而容易得来的，由于容易重新获得，很快就丢掉了。因此，对所有明智的人来说，若这个祝福不太容易，反而成了祝福。或者，也许是为了使我们不致遭受\Person{路济弗尔}的命运，他堕落了，并且因为接受了全部光明而对全能的主硬着颈项，遭受了最可怜的坠落，从所达到的高处坠落。或者，也许是为了给那些在这里已经洁净、并对他们所渴望的长期忍耐的人，将来对他们的辛劳和光荣生活给予更大的赏报。\par

因此，这身体的黑暗置于我们与天主之间，如同古时云彩置于埃及人和希伯来人之间（见：\BibleRef{出 14:20}）；这也许就是所谓\BibleQuote{祂以黑暗为藏身之处}所意指的，即我们的迟钝，借着它，很少人能看见哪怕一丁点。但关于这一点，让那些以此为职责的人去讨论吧；让他们在探究中尽量上升。对于我们这些（正如\Person{耶肋米亚}所说的）\BibleQuote{地上的囚徒}（见：\BibleRef{哀 3:34}）且被血肉的厚重覆盖的人，至少这一点是知道的：正如一个人不能跨过自己的影子，无论他跑得多快（因为影子总是跟随着，追赶的速度一样快），又如眼睛不能离开中间的气和光靠近可见物体，鱼不能在水外游动；同样，那些在身体内的人，完全脱离身体对象而与纯粹思想对象交往是完全不可能的。因为当我们心智最完全地脱离可见事物，收敛自身，试图接近那些与它相类而不可见的事物时，我们环境中的某种东西总是渗入。\par

十三、以下将为你们阐明此事：——\Quote{圣神、火、光、爱、智慧、正义、理智、理性}等，不正是第一性体的名称吗？那么，你们能设想圣神脱离运动与扩散吗？或火脱离其燃料、向上的运动及其固有的颜色与形式吗？或光不与空气混合、脱离那如同其父与源的事物吗？你们又如何设想理智？岂不是那内在于另一主体之中，而它的运动是思想——或沉默或表达？至于理性……你们还能想到别的吗，除了那在我们内静默的，或那涌流出来的（因为我避免说\Quote{被释放}）？若你们设想智慧，它岂不是你们所知的智性习惯，涉及神圣或属人的默观？正义与爱，岂不是可称赞的性情，一者与不义相对，一者与仇恨相对，时而加强自身，时而放松，时而占据我们，时而离弃我们，总之，造就我们如其所是，如同颜色改变身体？抑或我们应撇开这一切，尽我们所能，从它的肖像中收集片面的感知，绝对地注视神性？那么，这微妙之物是什么，它属于这些，却又不是这些？那在本性上不可复合、不可比较的合一，如何能同时是这一切，且完美地是每一个？因此，我们的心智在超越有形之物、与无形者交往时，会晕眩，只要它带着固有的软弱，注视超出其力量的事物。因为每一理性本性都渴望天主和第一原因，但如前所述，它无法把握他。因此，它们因渴望而虚弱，仿佛因无力而不耐烦，便尝试第二条路：或注视可见之物，并从中造出一个神……（拙劣的发明，因为可见之物在何种方面、多大程度上比观看者更高、更像神，竟要受其崇拜？）或透过可见之物的美和秩序，达到那超乎视觉者；但不要因可见之物的壮丽而丧失天主。\par

十四、因此，一些人将太阳奉为神，另一些人将月亮，或星辰之众，或连同其所有天体的天本身奉为神，并根据它们运动的性质或量，将宇宙的治理归于它们。另一些人则崇拜元素——土、气、水、火——因其有益的本性，因为无人能离之而生。另一些人则崇拜任何偶见的可见之物，将他们所见最美者立为神。还有那些崇拜图画和雕像的人，最初确实是他们祖先的雕像——至少，在更具情感和感官性的人中是如此——并以纪念物敬拜逝者；后来，连陌生人的雕像也被后世之人崇拜，因他们与初时相隔甚远；由于对第一性体的无知，他们视传统敬礼为合法而必要；因为习俗经时间巩固，便被奉为法律。我认为，有些谄媚专制权力、颂扬身体力量、赞美美貌的人，最终将他们所尊崇的人奉为神，或许还借用某种传说来助长其欺诈。\par

十五、他们中最易受激情支配的人，将自己的情欲神化，或在其诸神中尊崇它们：愤怒、嗜血、贪欲、酗酒，以及一切类似的恶行；并以此为自己罪恶的卑劣不公的借口。有些神祇他们置于地上，有些藏于地下（这是他们仅有的一点智慧），有些则升于天上。荒谬的继承分配！然后，他们凭其错误的权威和私意，将神或魔的名字赋予这些概念中的每一个，并设立以奢华为陷阱的雕像，以为用血和祭品的烟气，有时甚至以最可耻的行为、疯狂和杀人来尊崇它们，是恰当的。因为这样的尊荣正适合这样的神。人们甚至曾以崇拜怪物、四足兽、爬虫，以及极度卑劣可笑的受造物来侮辱自己，\BibleRef{罗 1:23}并将天主的荣耀献给他们；以至于我们难以决定，是更应鄙视崇拜者，还是他们所崇拜的对象。或许崇拜者更可鄙，因为他们虽有理性本性，并蒙天主恩宠，却以较劣者取代较优者。这是恶者的诡计，他滥用善以达到恶的目的，如同他在大多数恶行中所做的。因为他抓住人那在寻求天主时游移的渴望，为将能力扭曲归向自己，并偷窃那渴望，像盲人问路一样牵引着它；他将人抛掷并分散在各方，投入同一个死亡与毁灭的深渊。\par

XVI. 这就是他们的做法。然而，理性在我们对天主的渴望以及在我们无法没有领袖和向导的感觉中接纳了我们，然后引导我们关注可见之物，并遇见自太初以来存在的事物，但它并不停留于此。因为将主权赋予那些观察告诉我们同等地位的事物，并非智慧的本分。因此，理性通过这些事物引导我们走向超越它们的那一位，正是通过祂，这些事物才得以存在。因为是谁安排了天上和地上的事物，以及那些穿越空气的和生活在水中的事物？或者更确切地说，是那些先于这些的事物：天、地、空气和水？谁调和了它们，谁又区分了它们？每一事物与其他事物共同拥有什么？以及它们的相互依赖与和谐？因为我称赞那个人，尽管他是个外邦人，他说：\Quote{是什么给予它们运动，并驱使它们不停息、无障碍地运行？难道不是那创造者吗？祂将理性置于它们所有之内，宇宙便依照此理性运行并受管理。难道不是祂创造了它们并使它们存在？}因为我们不能将此能力归于\OriginalTerm{偶然}{Accidental}。因为，假设它的存在是偶然的，那么你将让我们把它的秩序归于什么？如果你愿意，我们可以把这让给你：那么你又将把它按其最初被造的条件得以保存和守护归于什么？这属于\OriginalTerm{偶然}{Accidental}，还是其他某物？当然不属于\OriginalTerm{偶然}{Accidental}。而这其他某物不是天主又能是什么？因此，从天主而来的理性，即从一开始就植入我们内并成为我们首要法则的理性，且贯串于一切之中，通过可见的事物将我们提升到天主面前。让我们重新开始，并推理出这一点。\par

XVII. 至于天主在本性和本质是什么，从来没有人发现过，也不可能发现。将来是否会被人发现，这是一个可以由任何人探讨和决定的问题。依我之见，当在我们内那神性而神圣的部分，即我们的心智和理性，与它的同侪相融，并且肖像上升到它现在所渴望的原型时，它就会被发现。而我认为这就是那棘手问题的解答，即\Quote{我们将认识，如同我们被认识一样}。但在我们现世生活中，临到我们的一切不过是一束微弱的光流，仿佛来自巨大光明的微小光辉。因此，如果有人认识天主，或有圣经为他认识天主作证，我们应明白此人拥有某种程度的知识，这使他显得比另一位没有同样光照的人更受启迪；而这种相对的优势被说成是绝对的知识，并非因为实际上如此，而是通过与他人的能力相比较。\par

XVIII. 于是，\Person{厄诺士}\Quote{希望呼求主的名}。希望是他受称赞的原因；这不是因为他\Emph{认识}天主，而是因为他\Emph{呼求}祂。\Person{哈诺客}被提升，但尚不清楚是因为他已领悟了神性，还是为了使他领悟。\Person{诺厄}\BibleRef{创 6:8}的荣耀是他蒙天主喜悦；他受托从洪水中拯救全世界，或者更确切地说，拯救世界的种子，藉着一只小方舟逃脱了洪水。\Person{亚巴郎}，虽然是伟大的圣祖，却因信德而成义，并献上奇异的牺牲，这牺牲是伟大祭献的预像。然而他并未看见天主，如同天主，而是像人一样给祂食物。他因尽其所能地朝拜而被认可。\BibleRef{创 32:28}\Person{雅各伯}梦见一个高耸的天梯和天使的梯子，并神秘地用油膏抹了一根柱子——也许是为象征那为我们受膏的磐石——并将那地方命名为\Place{天主之家}，以尊崇他所看见的那一位；并且与天主以人形搏斗；无论这人与天主的搏斗意味着什么……可能是指人的德行与天主的比较；他身上带着搏斗的标记，显示出受造本性的挫败；作为他虔敬的报酬，他的名字被改变：不再叫\Person{雅各伯}，而是叫\Person{以色列}——那个伟大而光荣的名字。然而，直到今日，他和他的子孙，即十二支派中的任何一个，都不能夸耀自己领悟了天主的全部本性或纯净的视见。\par

十九、至于厄里亚，正如你从故事中所学到的，不是强风、不是火、也不是地震，而是一丝微风隐约显现了天主的临在，但即便如此也并非祂的本性。这位厄里亚是谁？他是被火马车提到天上的人，象征着义人超人的卓越。你难道不惊讶于古代的民长玛诺亚和后来的门徒伯多禄吗？前者连一个带有天主肖像的显现都无法忍受，说道：\BibleQuote{我们完了，妻子，我们见了天主。}（\BibleRef{民 13:22}）就如同说，连一个天主的异象都不能被人把握，更不用说天主的本性了；后者无法忍受基督在他的船上，因此请求祂离开（\BibleRef{路 5:8}），尽管伯多禄比其他人更热切于认识基督，并因此受到了祝福（\BibleRef{玛 16:16}），并被托付了最大的恩赐。关于依撒意亚或厄则克尔——他们亲眼目睹了极大的奥秘——以及其他的先知，你能说什么呢？因为其中一位看见了万军的上主坐在光荣的宝座上，被六翼的色辣芬环绕、赞颂和遮掩，他自己被活炭洁净，并为先知职务作了装备。另一位描述了天主的革鲁宾战车（\BibleRef{则 1:4}），以及其上的宝座，和宝座之上的穹苍，以及在穹苍中显现的那一位，还有声音、能力和作为。至于这是白天只对圣人显现的景象，还是夜间无误的异象，或是心灵将未来如同现在一样交谈的印象，还是某种其他不可言喻的预言形式，我不能说；先知的天主知道，那些受启示的人也知道。但我所说的这些人和他们的同伴，谁也没有像经上记载的那样，站在天主的议会和本质面前，或看见或宣告了天主的本性。\par

二十、如果保禄被允许说出第三层天（\BibleRef{格后 12:2}）所涵盖的事，以及他进入、上升或被提的事，那么也许我们会更多知道一些关于天主的本性，如果这就是被提的奥秘的话。但既然那事不可言传，我们也就以沉默尊重它。我们听保禄这样说到它：\BibleQuote{我们现在所知道的，只是局部的，我们做先知所讲的，也只是局部的}（\BibleRef{格前 13:9}）。这是那在知识上并不粗鲁（\BibleRef{格后 11:6}）、并威胁要证明基督在他内说话的那位伟大博士和真理捍卫者的承认。因此，他估计地上的一切知识都如同藉着镜子观看，模糊不清（\BibleRef{格前 13:12}），立足于真理的微小形象。现在，除非有人觉得我过于谨慎，过于焦虑于这件事的考察，也许正是这件事，圣言自己暗示了有些事现在不能承受，但将来要承受并阐明（\BibleRef{若 16:12}），而圣言的前驱洗者若翰、伟大的真理之声，曾宣称甚至整个世界也容不下。\par

二十一、真理，以及全部的圣言，充满困难和晦暗；我们仿佛用微小的工具进行伟大的工作，当我们用仅仅属人的智慧来追求对自有者的认识时，并与感官为伴（或不离感官），被它们带到这里那里，引入错误，我们就致力于寻求那只能靠理智把握的事，而无法以赤裸的理智相遇赤裸的现实，好更接近真理，并以概念塑造心智。\par

如今，关于天主的话题更难达到，因为它比其他一切更完美，并且面临更多的反对意见，解决这些意见也更加费力。因为每一个反对意见，无论多么微小，都会阻止和妨碍我们论证的进程，并切断其进一步的进展，就像有人突然勒住全速前进的马的缰绳，使它们因意外的冲击而转身。因此，撒罗满——他是有史以来最有智慧的人（\BibleRef{列上 3:12}），无论在他之前还是他同时代的人，天主赐予他宽阔的心胸和如沙般丰富的默观——他越深入深渊，就越发眩晕，并宣告智慧的终点在于发现智慧离他有多远（\BibleRef{训 7:23}）。保禄也试图达到——我不敢说天主的本性，因为他知道这完全不可能——而只是天主的判断；但他既找不到出路，也找不到上升途中的停留之处，而且他心灵对知识的热切探寻没有以任何确定的结论告终，因为一些新的未达之境持续向他显露（啊，奇妙，我也有类似的经历），他以惊叹结束论述，称此为天主的丰富（\BibleRef{罗 11:33}）和深渊，并承认天主的判断不可测度，几乎是用达味的话——达味时而称天主的判断为大渊，其根基无法度量或感觉触及；时而又说天主对他的认识和他自身的构造是奇妙的，超出了他自己的力量或把握。\par

第二十二。因为，他说，若我暂且撇开一切其他事物，只思考我自己、人的整个本性与构造，我们如何被混合，我们的运动是什么，有死的如何与不朽的复合，我如何向下流却又被向上托起，灵魂如何被限制；它如何赋予生命并分担情感；理智如何同时被限制又不受限制，常驻于我们内却又以迅速的流动漫游宇宙；它如何被言语接纳并传递，穿过空气，进入一切事物；它如何参与感觉，却又远离感觉而隐藏。甚至在这些问题之前——我们在自然的工坊中最初是如何被塑造和组合的，我们最后的形成和完成是什么？对营养的欲望和给予是什么，谁使我们不由自主地到达那些最初的生命泉源和源头？身体如何靠食物滋养，灵魂如何靠理性滋养？自然的吸引是什么，父母与子女之间的相互关系是什么，竟能因爱的魔力而结合？物种如何是永久的，特性各不相同，尽管数量多到无法描述个体特征？同一动物如何既有死的又有不朽的——一是由于死亡，一是由于生成？因为一个离去，另一个接替，如同河流的水流，永不停歇却又始终恒常。你还可以讨论更多关于人的肢体和部分，它们为用途和美丽而相互适应，有些相连有些分离，有些更优越有些较不雅观，有些结合有些分开，有些包含有些被包含，按自然的规律和理性。关于声音和耳朵也有许多可说。声音如何由发声器官发出，由耳朵接收，两者通过空气介质的冲击和回响而连接？关于眼睛也有许多可说，它们与可见之物有难以言喻的交流，仅凭意志而移动，且同时移动，其受影响的方式与理智完全相同。因为理智与思想对象结合的速度，与眼睛与可见对象结合的速度相等。关于其他感官也有许多，不是理性研究所能探究的。关于我们在睡眠中的安息、梦中的幻象、记忆和回忆也有许多；关于计算、愤怒和欲望；简言之，这一切支配着这个被称为\Quote{人}的小世界。\par

第二十三。我要为你列举其他动物与我们之间以及它们彼此之间的差异吗？——本性、产生、营养、区域、性情以及社交生活等方面的差异。如何有一些是群居的，另一些是独居的；一些食草，另一些食肉；一些凶猛，另一些温顺；一些亲近人且被驯养，另一些不可驯服且自由？有些我们可称之为近乎有理性和学习能力，而另一些则完全缺乏理性，无法教导。有些感官更丰富，另一些较少；一些不动，另一些能行走，有些非常迅速，有些非常缓慢；有些在大小或美丽上超越，或在这两方面之一，有些非常小或非常丑，或两者兼有；有些强壮，有些软弱；有些善于自卫，有些胆小狡猾，另一些则毫无防备。有些勤劳节俭，另一些完全懒惰且无远见。在谈到这些问题之前，如何有一些是爬行的，另一些是直立的；一些固定在一处，一些两栖；一些喜爱美丽，另一些不加装饰；一些结婚，一些单身；一些节制，一些不节制；一些多产，一些不产；一些长寿，一些短命？详述所有细节将会是一篇令人厌倦的讲辞。\par

第二十四。也要看看游鱼一族，它们滑过水域，仿佛在液体元素中飞翔，呼吸着自己的空气，但与我们的空气接触时却处于危险之中，就像我们在水中一样。留意它们的习性和脾性，它们的交往和生育，它们的大小和美丽，它们对地点的喜爱和漫游，它们的聚集和离去，以及它们与陆栖动物如此相似的特质；在某些情况下，特质的相似与对比，无论是名称还是形状。再考虑鸟类的族群，以及它们形态和颜色的多样性，包括无声的和鸣禽。它们旋律的原因是什么？是谁给了它们？谁给了蚱蜢胸前的琴弦，以及树枝上的歌唱和啁啾，当它们被太阳感动，在正午奏响音乐，在树丛中歌唱，用声音护送旅人？谁为天鹅编织了歌声，当它展翅迎风，让翅膀的沙沙声奏出旋律？因为我不愿谈论被迫的声音，以及所有其余人为的矫饰反对真理。孔雀，那只美地亚的虚荣之鸟，从哪里获得对美丽和赞美之爱（因为它完全意识到自己的美丽），以致当它看见有人靠近，或如人们所说，要在雌孔雀面前炫耀时，它昂起脖子，展开尾巴环绕自己，金光闪闪，缀满星星，以高傲的步伐向爱慕者展示它的美丽？如今，\WorkTitle{圣经}赞赏妇女纺织的巧妙，说：\BibleQuote{谁给了妇女纺织的技巧和刺绣艺术的巧妙？}这属于有理性的活物，超越智慧，甚至一直到达天上的事物。\par

XXV. 但我愿你也惊叹于无理性受造物的自然知识；若能，请解释其原因。鸟儿如何以岩石、树木和屋顶为巢，并使之既安全又美观，且适宜哺育幼雏？蜜蜂和蜘蛛从何处获得勤劳与技艺？前者用以设计蜂巢，以六角形且相协调的管道将其连接，并通过隔板与角度的交替以直线构筑基底——即使是在如此昏暗的蜂房与巢脾中；后者则用以织造精巧的网，以如此轻盈、几乎如空气般的丝线向各方伸展，几乎始于不可见的起点，迅即成为佳美的居所，并为捕捉弱小而供其果腹之猎物设下陷阱。哪一位\Person{欧几里得}能模仿这些？他虽然以不存在之线探求哲学，耗神于证明。哪一位\Person{帕拉墨得斯}曾发明鹤类的战术、运动和阵形——如人们所说——以及它们列队移动与复杂飞行的体系？谁是它们的\Person{斐狄亚斯}与\Person{宙克西斯}？谁是懂得描绘与塑造至美之物的\Person{帕拉西亚斯}与\Person{阿格拉欧丰}？什么样的\Person{代达罗斯}的和谐\Place{克诺索斯}歌舞，为少女臻至最高美境？什么样的\Place{克里特}迷宫，难入难解，如诗人们所言，因构造之巧而不断自相交错？我不谈蚂蚁的仓库和仓储官，以及它们根据所需时间而储存的木材量，也不谈它们行军、领袖和工作中良好秩序的一切已知细节。\par

XXVI. 倘若这知识已为你所及，且你熟悉这些学科，那么也请观察植物之间的差异，直至叶子富有艺术的形态：它既为眼睛提供至高的愉悦，亦对果实最为有利。再看果实的多样与丰饶，尤其是那些最必需之物的美妙。思考根、汁液、花朵和气味的能力——不仅极其甘美，而且有医药之效；色彩的优雅与特质；宝石的珍贵价值与璀璨透明。因为大自然将一切置于你面前，如同丰盛的宴席，人人可取，既有生活必需品，亦有奢侈品——为的是，即使别无他由，你至少能因祂的恩惠而知天主，并通过自身的匮乏而比从前更智慧。接下来，我恳求你，遍历大地的长阔高深——这万物之母；以及彼此相连并与陆地相连的海湾；美丽的森林；丰沛永流的江河与泉源——不仅有地表上冰凉宜饮之水，亦有在地下奔流、流经洞穴，然后被猛烈之气挤压、回推，因这争斗与抗拒之力而变热，继而一旦有隙便逐渐喷涌而出，遂在多地提供我们所需之温泉，与冷泉一起给予我们免费且自发的医治。告诉我，这些如何而来，从何而来？这未经技艺编织的宏大网络是什么？这些事物，无论就其相互关系还是单独观察，其可惊叹之处绝不稍减。\par

大地如何稳固坚定、毫不动摇？它被支撑在什么之上？什么支撑着它，而它又立足于何物？因为甚至理性本身也无从倚靠，唯有天主的旨意。为何一部分隆起为山峰，一部分下陷为平原，且以各种不同方式分布？由于每一微小变化本身虽小，却既更丰盛地供应我们的需要，又因多样而更显美丽；一部分分配为居所，另一部分则荒无人烟——即所有高耸的山脉及从大陆割裂的各种海岸峭壁。这难道不是天主威能的最清晰明证吗？\par

XXVII. 关于海洋，即使我不惊叹其浩瀚，也当惊叹其温和：它虽流动不息，却安止于界内；若不惊叹其温和，也必惊叹其浩瀚；但我既惊叹两者，就当赞美那同时存在于两者中的权能。什么汇集了它？什么为其划界？它如何被兴起，又如何平息，如同尊重邻地一般？再者，它如何容纳所有江河，却仍保持自身不变——因其广袤无垠之极度丰富（若可如此说）？如此巨大元素的边界，怎会只是沙粒？你们那些精通无用琐事的自然哲学家——我指的是那些实际上试图用酒杯度量大海、用自己概念度量如此伟业的人——能告诉我们什么吗？还是让我从圣经中更简洁、却比最冗长的论证更令人满意且真实地给出真正科学的解释：\BibleQuote{祂以命令围住了水面的边界}。这就是流动本性之锁链。祂又如何使那居于旱地的航海者（即人）乘小舟、借微风航行其上（你难道不惊叹此景——你的心智不为之震撼吗？），使陆地与海洋因需要与贸易而相连，使本性如此分离的事物为人而合为一体？泉源的最初源头是什么？人啊，请寻求，若你能追溯或找到这些事中之任何一件。是谁为江河劈开平原与山脉，赐予它们无障碍的河道？另一边，海洋永不泛滥，江河永不枯竭，这奇事又如何而来？水的滋养能力是什么？其中有何差异？有些从上灌溉，有些从根汲取——若我可在谈论天主丰盛恩赐时稍纵情于言辞的话。\par

XXVIII. 如今，且让我们离开大地与地上之物，藉思想的翅膀翱翔于空中，使我们的论述能循正途前进；然后我愿领你升向天上的事物，直至天上本身，以及超越天上者；因为我的论述犹豫不敢升到那超越者，但仍要尽力而为。谁将空气这一伟大而丰饶的财富倾泻而出？它不按人的等级或运气来衡量，不受疆界限制，不按年龄分配，却像\OriginalTerm{玛纳}{Manna}的分配一样（\BibleRef{出 16:18}），每人都得到足够的份量，其公平分配为人所珍惜；它是飞鸟的战车，是风的宝座，是季节的调节者，是活物的激发者，或更确切地说，是肉体中自然生命的保存者；身体在其中存在，藉它我们说话；光与它所照耀的一切，以及流经它的视力，都在其中。并请留意随后之事。我不能将人们认为属于空气的一切权柄都归于空气。风的仓库是什么？（\BibleRef{约 37:9}）雪库又是什么？如圣经所说，谁孕育了露珠？冰出自谁的胎？谁将水束缚在云中，并将一部分固定于云中（奇妙啊！），以祂的\OriginalTerm{圣言}{Word}维系那本性流动之物，又将其余部分倾注于整个地面，按时适量地散布，既不使全部湿气不受约束地自由流出（因为在诺厄的日子，净化是足够的；祂不能撒谎，不会忘记自己的\OriginalTerm{盟约}{covenant}），……也不完全加以约束，以致我们需要一位\Person{厄里亚}（\BibleRef{列上 18:44}）来终结干旱。它说：若祂关闭天空，谁能开启？若祂开放水闸，谁能关闭？（\BibleRef{约 12:14}）谁能使雨水过多或不足，除非祂以自己的尺度和天平治理宇宙？请教，你这位从地上打雷，却连一星真理之光也发不出来的人，能对雷和闪电制定什么科学规律？你将以地上的何种蒸汽归咎于云的创造？或是由于空气的某种稠化，或是过于稀薄的云的挤压或碰撞，使你以为挤压是闪电的原因，碰撞是雷声的成因？或是风被压缩而无出路，你会以它的压缩解释闪电，以它的爆发解释雷声？\par

如今，如果你在思想中已穿越了空气和空气的一切，请随我到达天上和天上的事物。且让\OriginalTerm{信德}{faith}引导我们，而非理性——至少你已学到了后者在与你相近之事上的软弱，并藉认识那超越理性之物而认识了理性，以免你完全属于大地或出于大地，因为你连自己的无知也不知道。\par

XXIX. 谁将天空铺展在我们周围，并排列星辰？或更恰当地说，首先，凭你自身对天上事物的知识，能告诉我天与星辰是\Emph{是}什么？你连自己脚下是什么都不知道，甚至不能测量自己，却必须忙于你本性之上的事，凝视那无限者？因为，即或你明白轨道与周期、盈与缺、升与落、某些度与分，以及其他所有使你以奇妙知识而自豪的事；你并未达到对现实本身的理解，而只观察了一些运动，当经过长期实践证实，并将许多人的观察概括为一，从而推导出定律，便获得了\Quote{科学}之名（正如月相现象已普遍为我们的视觉所知），成为这一知识的基础。但如果你在这一主题上非常科学，并有正当理由值得称赞，请告诉我这秩序与运动的起因是什么。太阳何以成为全世界的灯塔之火，在众目之中如同合唱团的领唱，以其明亮遮蔽所有其他星辰，比它们彼此遮蔽更为彻底？证据是它们在他面前发光，但他胜过了它们，甚至不让人察觉它们与他同时升起；美丽如新郎，迅速而伟大如巨人；因为我不允许他的赞美从我自己的圣经以外的任何源头唱出——他力量如此强大，从世界的一端到另一端，他以热能拥抱万物，没有什么能躲过其感觉，但它以光充满每一只眼，以热充满每个有形的受造物；以其温和的性情和运动的秩序温暖而不灼烧，临在于一切，同等地拥抱一切。\par

XXX. 你曾考虑过这一事实的重要性吗？一位外教作家论及太阳在物质对象中的地位如同天主在思想对象中的地位。因为前者给眼睛带来光，如同后者给心智带来光；并且是可见物中最美的，如同天主是思想对象中最美的。但最初是谁给了它运动？是什么持续推动它在轨道中运行，尽管它本性稳定不动，真正不知疲倦，是生命的赐予者和维持者，以及诗人公正歌颂它的所有其他称号，且永不停止它的行程与恩惠？它如何成为白天的创造者（当它在大地之上时），以及黑夜的创造者（当它在大地之下时）？或者，当人默观太阳时，什么才是正确的表达？昼夜之间的相互进攻与退让是什么？以及它们有规律的失常——用一个有些奇怪的表达？它如何成为季节的制造者和划分者？这些季节按规律秩序来去，如同舞蹈般相互交织，或依爱之律与秩序之律分别，并逐渐混合，悄然接近相邻者，正如昼夜所为，以免因骤变而使我们痛苦。关于太阳，这就够了。\par

你可知月亮的本性与现象，光的尺度与行程，何以太阳主管白昼，月亮统辖黑夜？月赋予野兽胆气，日则唤醒人劳作，或升或降，各随其宜。你可知道昴宿的结系，或参宿的栅栏，正如那核点星辰数目、一一呼名者？\BibleRef{约 38:31}你可知道每一荣耀的差异\BibleRef{格前 15:41}，及其运行的秩序，好叫我信任你，因你藉它们编织人间事务，举起受造物反对造物主？\par

三十一。你说什么？我们在此止步，除了物质与可见之物外不再讨论，抑或——既然圣言知道\Person{梅瑟}的会幕是全受造界（我意指整个可见与不可见之物体系）的预象——我们要穿过第一层帐幔，越过感觉领域，进入圣所，即理智与天上的受造界？然而，即便是这（理智受造界），我们也不能以无形的方式看见，尽管它是无形的，因为它被称为——或者说是——\OriginalTerm{火与神体}{Fire and Spirit}。因为祂曾说：\Quote{祂使自己的天使为神体，祂的仆役为火焰。}……然而，或许这个\Quote{制造}，意指藉着那使他们存在的话语而保存他们。天使于是被称为神体与火：神体，因他是理智领域的受造物；火，因他有净化之性；因为我知道同样的名称也属于第一本性。但是，至少相对于我们而言，我们必须认为天使的本性是无形的，或者至少尽可能接近无形。你可看见我们如何对此主题头晕目眩，无法前进任何一点，除非到此地步：我们知道有\OriginalTerm{天使}{Angels}、\OriginalTerm{总领天使}{Archangels}、\OriginalTerm{上座者}{Thrones}、\OriginalTerm{宰制者}{Dominions}、\OriginalTerm{率领者}{Princedoms}、\OriginalTerm{异力者}{Powers}、\OriginalTerm{光辉者}{Splendours}、\OriginalTerm{升天者}{Ascents}、\OriginalTerm{理智异力或理智体}{Intelligent Powers or Intelligencies}，纯洁无杂的本性，对恶毫不动摇，或几乎不动摇；永远环绕第一因合唱（不然我们该如何赞美他们？），自那里被最纯净的光照所光照，或这天使照一一度，那天使照另一度，按各自的本性与等级比例……如此美丽相称、被塑造，以至成为第二等的光明，能藉第一光的洋溢与慷慨光照他人？他们是天主旨意的仆役，兼具天赋与分授的力量，穿行万有，因服事的热忱与本性的敏捷，随时随处临在……他们中不同的分子拥抱世界不同的部分，或被委派管理宇宙的不同区域，正如那安排并分配一切的祂所知晓。将万有合而为一，单单着眼于造物主万有的旨意；他们是神性威严的赞美者，永恒地默观永恒的光荣——并非为天主借此获得荣耀的增加，因为圆满者无可加添——祂向自身以外的一切供应美善，而是为使这些在神之后的首批本性永不停止领受福佑。如果我们按这些事物之所值讲述，那是因着天主圣三与三个位格中的唯一神性的恩宠；但若不如我们所愿那般完美，即便如此，我们的论述也已达成其目的。因为这正是我们努力要表明的：连第二等的本性也超越我们理智的能力，更何况那第一的（我惧怕于仅仅说那在万有之上的）、唯一的本性呢？\par

\SourceRecord{Translated by Charles Gordon Browne and James Edward Swallow.；Nicene and Post-Nicene Fathers, Second Series；Vol. 7.；Edited by Philip Schaff and Henry Wace.；Buffalo, NY: Christian Literature Publishing Co.,；1894.；<http://www.newadvent.org/fathers/310228.htm>.}

% structure: p0001 source=h1 target=section reason=page-title

\section{第三篇神学演讲（演讲29）}

\Emph{第三篇神学演讲。}\par

\Emph{论圣子。}\par

以上就是为遏制对手好辩的性情及其轻率作风（这使他们在一切事上，尤其在有关天主的讨论中缺乏安全感）可说的简要的话。然而，责备他人毫无困难，是非常容易的事，谁喜欢都能做到；但以自身的信仰取代他们的信仰，却是有德性、有智慧之人的本分。因此，让我们依靠在对手那里受辱、在我们这里受敬拜的圣神，将我们自己对天主性的看法——无论这些看法是什么——像某种高贵而适时的产儿一样带到光中。这并不是说我以前一直沉默；因为唯独在这件事上我充满青年般的活力与勇气；但事实上，在目前的情况下，我更加勇于宣讲真理，以免（用圣经的话说）因退缩而落入不讨天主喜悦的定罪中。既然每一篇讲论都有双重性质——一部分是确立自己的立场，另一部分是推翻对手的立场——让我们首先陈述自己的立场，然后努力反驳对手的立场——两者都尽可能简明扼要，使我们的论点一目了然（就像他们为欺骗愚昧或单纯之人所杜撰的基础论著那样），并且使我们的思想不致因讲论冗长而分散，如同水不被引导在沟渠中，而是漫无目的地流散在开阔地上。\par

关于天主，有三种最古老的观点：\OriginalTerm{无统治}{Anarchia}、\OriginalTerm{多人统治}{Polyarchia}和\OriginalTerm{独一统治}{Monarchia}。前两种是希腊人的子孙的儿戏，愿他们继续如此。因为无统治是混乱无序的；而多人统治是分党结派的，因而也是无政府的，因而是混乱的。因为两者都趋向同一件事，即混乱；而混乱导致解体，因为混乱是解体的第一步。\par

但我们所尊崇的是独一统治。然而，这独一统治并不仅限于一位，因为统一若与自己不睦，就可能陷入多元的状态；而是由性体平等、意志合一、行动同一、以及各要素汇聚于一体所构成的——这是受造性体所不能及的——因此，虽然在数目上各别，但本质并无分割。如此，统一从永恒出发，通过运动到达二元，在三位中得享安息。这就是我们所说的圣父、圣子和圣神。圣父是生育者与发出者——当然是无痛苦、无时间、非物质的方式。圣子是受生者，圣神是发出者；因为我不知如何用完全排除有形事物的词汇来表达此事。我们不敢像一位希腊哲学家胆敢说的那样，说什么\Quote{美善的溢出}，仿佛它是一个满溢的碗，而且他在关于第一因和第二因的讲论中竟直言不讳。我们绝不可视此受生为不自主的，如同某种自然的、难以遏制的溢出，完全不符合我们对天主性的理解。因此，让我们安守本分，按天主圣言亲自所说，讲论无始者、受生者及由圣父所发者。\par

这些何时存在？它们超越一切\Quote{何时}。但如果我要更大胆地说——当圣父存在之时。圣父何时存在？从未有过祂不存在的时刻。圣子和圣神也是如此。你再问我，我再回答你：圣子何时受生？当圣父不受生之时。圣神何时发出？当圣子不是发出而是受生之时——超越时间，超乎理性；尽管我们无法表达那超越时间的事，如果我们像所愿的那样避免任何暗示时间的表达。因为诸如\Quote{何时}、\Quote{以前}、\Quote{以后}和\Quote{起初}之类的说法，无论我们如何勉强，都不是无时间的；除非我们采用\OriginalTerm{永世}{Æon}（即与永恒事物同延的间隔），它不被任何运动或太阳的绕行所分割或度量，如时间被度量那样。\par

那么，既然它们同永，为何不都像无始？因为它们来自祂，尽管不是在祂之后。因为无始者是永恒的，但永恒者不一定是无始的，只要它可以追溯到圣父为其根源。因此，就根源而言，它们不是无始的；但显而易见，根源不一定先于其效果，因为太阳并不先于它的光。然而，就时间而言，它们在某种意义上又是无始的，即使你以诡辩吓唬头脑简单的人，因为时间的根源不在时间之内。\par

然而，这受生如何是无痛苦的？因为它不是物质的。因为如果物质受生包含痛苦，非物质受生则排除痛苦。我要反过来问你：如果祂是受造的，祂怎能是天主？因为受造的不是天主。我不提醒你，如果从物质意义上理解创造，这里也有痛苦，如时间、欲望、想象、思想、希望、痛苦、风险、失败、成功等——所有这一切，以及更多的，都存在于受造物中，这是人人都清楚的事。不，我惊讶的是，你竟不敢做到如此地步，以至于构想婚姻、怀孕期和流产的危险，好像圣父若不如此受生就根本不能受生；或者，你又不敢列举飞禽、走兽和鱼类的受生方式，并将神圣而不可言喻的受生归入其中之一，甚或从你的新假设中消除圣子。你甚至看不到这一点：正如祂按肉身的受生与一切其他受生不同（因为在你所知的人中，哪里有一位童贞母亲？），祂在属神的受生上也是如此；或者更确切地说，祂的存在与我们的不同，因而在受生上与我们也不同。\par

五、那么，那位没有开始的圣父是谁？是一位其存在本身没有开始者；因为存在有开始者，也必然开始为圣父。祂并非在开始存在之后才成为圣父，因为祂的存在没有开始。祂是绝对意义上的圣父，因为祂也不是圣子；正如圣子是绝对意义上的圣子，因为祂也不是圣父。这些称谓不属于我们绝对的意义，因为我们两者都是，并不此多彼少；我们出于两者，而非仅出于其一；因此我们是被分割的，逐渐成为人，甚至可能不是人，成为我们并不愿成为的，离开与被离开，以至于只留下关系，而无实在的基础。\par

但是，反对者说，\Quote{祂生了}和\Quote{祂受生}这一表达形式本身就引入了受生的开始观念。然而，如果你不使用这种表达，而说\Quote{祂从起初就受生}，以便轻易回避你那些牵强附会和喜好时间的反对意见，那又怎样呢？你会引用圣经来反对我们，好像我们捏造了某些与圣经和真理相悖的东西吗？其实，众所周知，在实践中我们经常发现，在谈论时间时，时态是互换的；尤其是圣经的习惯，不仅在过去时和现在时中如此，甚至在未来时中也是如此，例如\BibleQuote{外邦为什么嚣张？}\BibleRef{咏 2:1}，那时他们尚未嚣张；以及\BibleQuote{他们将要步行过河}\BibleRef{咏 66:6}，其意思是他们确实过了河。要列举出学者们注意到的所有这类表达，将是一项漫长的任务。\par

六、关于这一点就谈到这里。他们下一个反对意见是什么？充满了好斗和厚颜无耻？他们说，祂要么自愿生圣子，要么非自愿。接着，他们以为从两方面用绳索捆绑了我们；但这些绳索并不结实，而是非常脆弱。因为，他们说，若是非自愿，祂就受制于某一位，而这位是谁？若祂受其支配，祂又何以是天主？但若是自愿，那么圣子就是意志之子；那么祂如何是出自圣父？——他们就这样为祂发明了一种新的母亲——意志——来代替圣父。关于他们的这个论点，有一点可以算是好的，那就是他们抛弃了\Quote{受生}，而躲进了\Quote{意志}。因为意志不是受生。\par

其次，让我们看看他们论点的力量。最好是先与他们进行短兵相接的较量。你本人，那个如此轻率断言任何你之所想的人，你是被你父亲自愿生的还是非自愿生的？若是非自愿，那么他受制于某个暴君（哦！可怕的暴力！），而暴君是谁？你很难说是本性——因为本性是容忍贞洁的。若是自愿，那么几个音节便消除了你的父亲，因为你被证明是意志之子，而不是你父亲之子。但我转向天主与受造物之间的关系，并将你自己的问题向你的智慧提出：天主是自愿创造了万物，还是被迫的？若是被迫，这里也有暴政和施行暴政者；若是自愿，那么受造物也被剥夺了他们的天主，而你首当其冲，因为你发明了这样的论证和诡辩。因为在创造者与受造物之间设置了意志作为障碍。然而，我认为愿意的主体与愿意的行为是不同的；生育者与生育的行为是不同的；说话者与言语是不同的，否则我们都非常愚蠢。一方面我们有动因，另一方面有所谓的运动。因此，所愿意的并不是意志的产物，因为它并非总是由此产生；所生的也不是生育的产物，所听到的也不是言语的产物，而是出于那愿意、生育或说话的主体。但天主的事超越这一切，因为对祂来说，愿意生育也许就是生育，没有中间的行为（如果我们完全接受这一点，而不是认为生育高于意志）。\par

七、那么，你愿意让我在这个\Quote{圣父}一词上稍微戏谑一番吗？因为你的例子鼓励我如此大胆。圣父是天主，或是愿意的，或是不愿意的；你如何逃脱你自己过度的机敏？若是愿意的，祂从何时开始愿意？不可能是祂开始存在之前，因为祂之前没有任何事物。或者祂的一部分是意志，另一部分是意志的对象？若是如此，祂就是可分的。因此，根据你的论证，问题就产生了：祂自己是不是意志之子？若是不愿意的，是什么迫使祂存在？如果祂是被迫的——而且恰恰是被迫成为天主——祂又何以是天主？那么，我的对手说，祂是如何受生的？祂是如何被造的，如果如你所说，祂是被造的？因为这也是同样困难的一部分。也许你会说，凭借意志和圣言。你还没有解决全部的困难；因为你还需要说明意志和圣言是如何获得行动能力的。因为人并不是这样被造的。\par

VIII. 那么，祂是如何被生的呢？这个生本不是什么大事，倘若你能理解它的话——你连自己的生都未曾真正认识，或至少理解得极少，而且连那一点点也羞于谈论，而你还以为自己知道全部吗？你必须费很大功夫才能发现组合、形成、显现的法则，以及灵魂与身体——理智与灵魂、理性与理智——结合的纽带；还有运动、增长、食物的同化、感觉、记忆、回忆，以及组成你的其他一切部分；其中哪些属于灵魂与身体共同，哪些各自独立，哪些彼此接受。因为那些成熟较晚的部分，在受孕时便已接受了它们的法则。告诉我这些法则是什么？然后也不要冒险去猜测天主的生；因为那是不安全的。即使你完全知道自己的事，也绝不知道天主的事。如果你不理解自己的事，又如何能知道天主的事？因为天主比人更难追踪，同样，天上的生也比你的生更难理解。但如果你因为不能理解它，就断言祂不能被生，那么你就有时间否定许多你无法理解的现存事物；首先就是天主本身。因为你无法说出祂是什么，即使你非常鲁莽，过于自负于你的智力。首先，抛开你那些关于流淌、分裂和切割的观念，以及你像对待物质生育那样对待非物质生育的概念，然后你或许才能恰当地领会天主的生。祂是如何被生的？——我愤慨地重复这个问题。天主的生当以沉默为尊。你得知祂是被生的，这已是大事。但祂生的方式，我们连天使都无法领会，何况是你呢？要我告诉你它是如何的吗？它是按那生了子的圣父和那被生的圣子所知道的方式。除此之外，都被乌云遮蔽，逃出你昏暗的视线。\par

IX. 好吧，但圣父生的圣子要么是先前存在的，要么是不存在的。这是多么荒谬！这个问题适用于你或我：我们一方面已经存在，例如肋未在亚巴郎的腰中（\BibleRef{希 7:10}），另一方面又进入存在；因此从某种意义上说，我们部分来自已存在的，部分来自不存在的。而原始物质的情况则相反，它肯定是从不存在之物中被造的，尽管有些人假装它是无始的。但在这种情况下，\Quote{被生}从一开始就与\Quote{存在}同时发生。那么，你将这个狡猾的问题建立在什么基础上呢？因为有什么比从元始而来的更古老，如果我们把圣子的先前存在或不存在放在那里？无论哪种情况，我们都破坏了它作为元始的主张。或者，也许你会说，如果我们问你圣父是由存在的还是不存在的质料生成的，祂是双重的，部分先存，部分存在；或者祂的情况与圣子相同，即祂是从不存在的质料中被造的——因为你那些可笑的问题和沙土之屋，连最轻微的涟漪都经不住。\par

我不接受任何一种解决，并宣称你的问题包含荒谬，而不是难以回答。然而，如果你按照你的辩证假设，认为这些选择之一必定每一情况都为真，那么让我问你一个小问题：时间是在时间之内，还是不在时间之内？如果它在时间之内，那么在什么时间之内，它除了那个时间又是什么，以及它如何包含它？但如果它不在时间之内，那么能够设想一个无时间的时间，那是什么样的超越智慧？现在，关于这个表述：\Quote{我现在正在说谎}，请承认其中一个选择：要么它是真的，要么它是假的，不加限定（因为我们不能承认它既是真又是假）。但这不可能。因为必然，他要么是在说谎，因而是在说真话，要么是在说真话，因而是在说谎。那么，既然如此，有些情况下相反的事物为真，那么有什么奇怪，在另一种情况下它们可能都为假，因而你的巧妙谜题只是愚蠢？再为我解一个谜。你是在你的生之时在场，并且现在你也在自己面前，还是两种情况都不是？如果你曾在场且仍在场，你是谁，你和谁在一起？你单独的自我如何同时成为主体和客体？但如果以上都不是，你如何与自己分离，这种分离的原因是什么？但你会说，对一个人是否与自己同在的问题大惊小怪是愚蠢的；因为这种表述不是用于自己，而是用于他人。那么，你可以肯定，讨论那从元始就被生的东西在它生之前是否存在，是更加愚蠢的。因为这样的问题只出现在按时间分割的质料中。\par

X. 但他们说：\Quote{\OriginalTerm{无始者}{Unbegotten}与\OriginalTerm{受生者}{Begotten}并不相同；若是这样，圣子也与圣父不相同。} 不用明说也知道，这种论证路线显然是将圣子或圣父排除在天主性之外。因为如果\OriginalTerm{无始}{Unbegotten}是天主的本质，那么受生就不是那本质；如果相反的情况成立，则无始者就被排除。有什么论证能反驳这一点呢？那么，我的新神学的好发明家，如果你无论如何都执意要拥抱一种亵渎，就请选择你更喜欢的亵渎吧。其次，你断言无始者与受生者不相同，是在什么意义上说的？如果你是指非受造者与受造者不相同，我同意你的看法；因为\OriginalTerm{无源者}{Unoriginate}与受造者确实不是同一本性。但如果你说生者与受生者不相同，这种说法是不准确的。因为实际上它们相同是一个必然真理。因为父与子的关系的本性就是：子嗣与父母具有同一本性。或者我们这样论证：你说无始与受生是什么意思？如果你是指无始或受生的单纯事实，它们并不相同；但如果你是指这些名称所指向的那一位，它们如何不相同呢？例如，智慧与不智慧本身并不相同，但两者都是同一个人——人的属性；它们标记的不是本质的差异，而是本质之外的差异。不朽、无罪、不变也是天主的本质吗？如果是这样，天主就有许多本质而非一个；或者天主性由这些复合而成。因为如果它们是本质，天主就不可能没有复合而同时拥有这一切。\par

XI. 他们并不这样断言，因为这些属性也通用于其他存有。但天主的本质是唯独属于天主、且是祂所特有的。然而，那些认为物质和形式是\OriginalTerm{无始}{unbegotten}的人，不会承认无始是天主独有的属性（因为我们必须进一步摒弃摩尼教人的黑暗）。但假设它是天主独有的属性。那亚当呢？他不是唯独天主直接创造的吗？你会说：是的。那么，他就是唯一的人了吗？决不是。为什么？因为人性并不在于直接创造？因为受生的也是人。同样，\OriginalTerm{无始者}{Unbegotten}也不是唯一的天主，虽然祂是唯一的父。但请承认，\OriginalTerm{受生者}{Begotten}也是天主；因为祂出自天主，即使你固守你的无始，你也必须承认这一点。那么，你如何描述天主的本质呢？不是通过宣告它是什么，而是通过否定它不是什么。因为你的词语表示祂不是受生的；它并没有向你呈现那没有受生者的真实本性或状态。那么天主的本质\Emph{是}什么呢？既然你对祂的受生也如此焦虑，就该由你的愚昧来定义它；但对我们来说，如果将来有一天，当这黑暗和愚钝为我们除去时，我们能够知道它，正如那位不说谎者所应许的，那将是一件非常伟大的事。对于那些为达此目的而净化自己的人，这也许是他们的思想和希望。至于我们，我们将敢于这样说：如果圣父\OriginalTerm{无源}{Unoriginate}是一件伟大的事，那么圣子由这样一位圣父受生也丝毫不逊色。因为祂不仅分享了\OriginalTerm{无源者}{Unoriginate}的光荣——因为祂出自无源者——而且还有祂受生的额外光荣，这在所有那些不是完全粗俗和物质思维的人眼中，是如此伟大和崇高。\par

XII. 但是，他们说：\Quote{如果圣子在本质上与圣父相同，那么如果圣父是\OriginalTerm{无始的}{unbegotten}，圣子也必须是无始的。} 完全如此——如果天主的本质在于\OriginalTerm{无始}{Unbegotten}；那么祂就会是一个奇怪的混合体：\OriginalTerm{受生方式的无始}{begottenly unbegotten}。然而，如果差异在本质之外，你怎么敢如此肯定地说论它呢？难道你也是你父亲的父亲，以至于在本质相同的前提下，在任何一个方面都不欠缺你的父亲吗？显然，我们若探究天主本质的本性，将绝对不影响位格。但\OriginalTerm{无始}{Unbegotten}并非天主的同义词，证明如下：如果是这样，既然天主是一个关系词，那么无始也必须是关系词；或者，既然无始是一个绝对词，那么天主也必须是绝对词……即谁的天主都不是。因为完全相同的词有相同的用法。但\Quote{无始}并非用作关系词。因为它是相对于什么？天主是哪些事物的天主？为什么是所有事物的。那么天主和无始怎么可能是同义词呢？再者，既然受生与无始是矛盾的，就像拥有与缺乏一样，那么矛盾的本质就会共存，这是不可能的。或者，既然拥有先于缺乏，而缺乏是对拥有的否定，那么按照你的假设，不仅圣子的本质必须先于圣父的本质，而且圣子的本质必须被圣父所毁灭。\par

XIII. 如今他们那不可抗拒的论点还剩下什么呢？也许他们最后会逃遁于此：如果天主从未停止生育，那么生育就是不完善的；祂何时才会停止呢？但如果祂停止了，那么祂必然曾经开始过。这样，这些属血肉的人再次提出属血肉的论点。祂是否永远受生，我暂时不说，直到我更精确地审视这句话：\BibleQuote{在一切山丘之先，祂生了我，}\BibleRef{箴 8:25}。但我看不出他们结论的必然性。因为，如他们所说，凡有终结者必有开始，那么当然，凡无终结者必无开始。那么，他们对于灵魂或天使本性将如何决定？如果它有开始，它也会有终结；如果它没有终结，那么显然按照他们的说法，它就没有开始。但事实是它有开始，却永不会有终结。因此，他们的断言——那将有终结的也曾经有开始——是不真实的。然而我们的立场是：正如对于一匹马、一头牛或一个人，同一物种的所有个体适用同一个定义，凡分享这定义者也有权分享这名称；同样，天主有一个同一的本质、一个同一的本性、一个同一的名称，尽管按照我们思想中的区别，我们使用不同的名称；并且凡真正用这名称称呼的，实在是天主；祂在本性上是怎样的，祂就真正被怎样称呼——至少如果我们坚持真理不在于名号而在于实在的话。但我们的对手，仿佛害怕留下任何一块石头不去翻转以推翻真理，尽管当他们被论据和证据所迫时，承认圣子是天主，但他们只意味着祂是在模棱两可的意义上是天主，并且祂只是分享这名称而已。\par

XIV. 当我们向他们提出这个异议时说：\Quote{那么你们想要说什么呢？圣子不是真正的天主，就像一幅动物的图画不是真正的动物一样吗？如果不是真正的天主，祂究竟在什么意义上是天主呢？}他们回答说：\Quote{为什么这些术语不能是模棱两可的，并且在两种情况下以恰当的意义使用呢？}他们会给我们一些例子，诸如陆犬和鲨鱼（dogfish）；其中\Quote{狗}这个词是模棱两可的，但在两种情况下都是恰当使用的，因为在模棱两可的名称中有这样的物种，或者任何其他同样的名称用于两种不同本性的事物的情况。但是，我的好朋友，在现在的情况下，当你把两个本性包括在同一个名称之下时，你并没有断言其中一个优于另一个，或者一个在先而另一个在后，或者一个程度更大而另一个程度更小地拥有那共同谓述的东西，因为没有任何联系迫使这种必然性加于它们。一个并不比另一个更是狗，一个也不更少；无论是鲨鱼比陆犬，或是陆犬比鲨鱼，都不更甚。为什么它们会是，或根据什么原则？但这里名称的共同性存在于同等价值、尽管本性不同的事物之间。但在我们所谈的场合，你把天主的名称与可敬的威严相连，使它超越一切本质和本性（这唯独属于天主的属性），然后你将这名称归于圣父，却剥夺圣子对它的权利，使圣子服从圣父，只给祂次等的尊荣和敬拜；即使你在言辞中赐予祂同等的尊荣，实际上你却割断了祂的神性，并且恶意地从使用同一个名称（意指完全的平等）转到联系不等同的事物。因此，在你的口中，图画中的人与活人比我所举的狗更适用于说明神性关系的比喻。否则，你既承认两者有相同的名称，也必须承认两者有同等的本性尊严——即使你为了论证引入了这些不同的本性；这样你就破坏了你那狗的类比，即你发明来作为不平等例子的类比。如果你所区别的那些荣誉并不相等，那么你那模棱两可的例子的力量何在？因为你以狗为庇护，不是为了证明平等，而是为了证明不平等。还有谁能更清楚地表明他既与自己的论点作战，又与神性作战呢？\par

XV. 倘若我们承认，就起因而言，圣父大于圣子，他们却假定祂是因本性而为起因，然后推导出结论说祂也是因本性而更大，那就很难说他们误导自己更多还是误导与他们辩论的人更多。因为并非绝对地，凡是对一个种类可谓述的，也可以对组成该种类的所有个体谓述；因为不同的个别属性可能属于不同的个体。因为如果我自己采取同样的前提，即圣父因本性更大，然后再加上这一点：\Quote{然而并非在各方面都因本性更大，也非在各方面是圣父}——从而得出结论：\Quote{因此更大的并非在各方面都更大，圣父也非在各方面都是圣父}——有什么妨碍我呢？或者，若你愿意，让我们这样说：天主是一个本质；但并非每一个本质都是天主；然后请你自己下结论——因此天主并非在每一个场合都是天主。我认为这里的谬误是从有条件的用法推理到无条件的用法，用逻辑学家的术语来说。当我们把这\Quote{更大}一词归于祂作为起因的本性时，他们却推断它关乎祂自有的本性。这就像我们说某人是一个死人（dead man），他们却推断他只是一个人（Man）。\par

XVI. 我们怎能略过以下一点呢？这一点与其他部分同样令人惊叹。他们说，\Quote{圣父}是一个本质或一个行动的名称，企图从两方面束缚我们。如果我们说它是本质的名称，他们就会说我们同意他们的观点，即圣子出自另一个本质，因为天主的本质只有一个，而按他们的说法，这本质已被圣父所占据。另一方面，如果我们说它是行动的名称，我们就会被认为公开承认圣子是受造而非受生的。因为哪里有主动者，哪里就必有效果。他们会说，他们奇怪受造物怎能与造物者相同。我自己若是必须接受二者之一，而不将两者都搁置，提出第三个更真实的说法，即最聪明的先生们，\Quote{圣父}既非本质的名称，也非行动的名称，而是圣父对圣子、圣子对圣父的关系的名称。因为如同在我们中间，这些名字表明一种真实而亲密的关系，同样，在眼前的情形中，它们也表明受生者与生者之间的本性同一。但让我们向你们让步，假定\Quote{圣父}是本质的名称，它仍会引入圣子的概念，而不会使其具有不同的本性，这是根据通常的概念和这些名称的力量。如果你们高兴，就假定它是行动的名称；你们也无法以此击败我们。因为这种行动的结果确实会是\OriginalTerm{同体}{Homoousion}，否则在此事上行动的概念就是荒谬的。你们看，即使你们试图不公平地争斗，我们仍避开了你们的诡辩。但现在，既然我们已知道你们在论点和诡辩方面是多么不可战胜，让我们看看你们在天主圣言中的力量吧，或许你们可能选择从圣经中说服我们。\par

XVII. 因为我们从那些伟大而崇高的言词学会了相信并教导圣子的天主性。这些言词是什么呢？这些：天主——圣言——在起初就与元始同在的祂，以及元始。\Quote{在起初已有圣言，圣言与天主同在，圣言就是天主，}\BibleRef{若 1:1} 以及\Quote{在你那里有元始，}以及\Quote{他从起初就称呼她为元始。}\BibleRef{依 41:4} 然后圣子是\OriginalTerm{独生子}{Only-begotten}：\Quote{那在父怀里的独生子，祂给我们详述了。}\BibleRef{若 1:18} 道路、真理、生命、光明。\Quote{我是道路、真理、生命}；以及\Quote{我是世界的光。}智慧与能力，\Quote{基督是天主的智慧和天主的德能。}\BibleRef{格前 1:24} \OriginalTerm{光辉}{Effulgence}、\OriginalTerm{印迹}{Impress}、\OriginalTerm{肖像}{Image}、\OriginalTerm{印记}{Seal}；\Quote{祂是祂荣耀的光辉和祂本体的印迹，}以及\Quote{祂良善的肖像，}\BibleRef{智 7:26} 以及\Quote{天主父以印记印证了祂。}\BibleRef{若 6:27} 主、君王、自有者、全能者。\Quote{主从天主那里降下火雨；}\BibleRef{创 19:24} 以及\Quote{正义的权杖是你王国的权杖；}以及\Quote{那今在、昔在、将来永在的全能者}\BibleRef{默 1:8}——所有这些都明显地是指圣子，连同所有其他同样含义的经文，其中没有一个是后来才想到的，或是后来才加给圣子或圣神的，正如也不是加给圣父自己一样。因为他们的完美不受添加的影响。从来没有一个时候，祂没有圣言，或者祂不是父，或者祂不是真实的，或者不是智慧的，或者不是德能的，或者缺乏生命、光辉或良善。\par

但是，与所有这些相反，请你们为我数算那些助长你们无知傲慢的言辞，例如\Quote{我的天主和你们的天主}，或更大的、受造的、被造的、被圣化的；如果你们愿意，再加上\Quote{仆人}\BibleRef{斐 2:7}和\Quote{服从}\BibleRef{斐 2:8}和\Quote{赐予}\BibleRef{若 1:12}和\Quote{学习}\BibleRef{希 5:8}，以及被命令、被派遣、自己什么也不能做、或说、或判断、或赐予、或意愿。还有这些——\Quote{祂的无知}\BibleRef{谷 13:32}、\Quote{服从}\BibleRef{格前 15:28}、\Quote{祈祷}\BibleRef{路 6:12}、\Quote{请求}\BibleRef{若 14:16}、\Quote{成长}\BibleRef{路 2:52}、\Quote{被成全}。如果你们愿意，甚至比这些更卑微的，例如提到祂的睡觉、饥饿、\Quote{处于痛苦中}\BibleRef{路 22:44}和\Quote{害怕}\BibleRef{希 5:7}；或许你们甚至将祂的十字架和死亡也作为对祂的羞辱。我想你们会把祂的复活和升天留给我，因为在这些事中可以找到支持\Emph{我们的}立场的东西。如果你们想要拼凑那个模棱两可、硬塞给你们的假神——祂对我们而言是真天主，与父同等——你们可能还会收集许多其他事情。因为这些点中的每一个，如果我们逐一探讨，都可以很容易地用最敬畏的意义向你们解释，经文字面的绊脚石也可以被清除——前提是你们绊倒是因为诚实，而非恶意。用一句话给你们解释：崇高的事，你们应归于天主性，归于祂那个超越苦难、无形体的本性；而卑微的事，则归于祂为了你们而自谦、降生成人的复合状态——是的，说祂成了人，之后又被高举，这并没有什么不好。结果是你们将抛弃这些属血肉的、卑下的学说，学习更高超的事，与祂的天主性一同上升，不再永久停留在可见的事物中，而是与祂一同上升到思想的世界，并认识哪些经文指涉祂的本性，哪些指涉祂的取了人性。\par

十九. 因为那如今被你们藐视的那一位，曾是在你们之上的。那如今是人的那一位，曾是\OriginalTerm{非复合的}{Uncompounded}。祂所是的，祂继续是；祂所不是的，祂取为自己。起初，祂是无原因的；因为什么是天主的原因？但后来为了一个原因，祂被生了。那原因就是你们能得救——你们侮辱祂、蔑视祂的天主性，因为祂取了你们更密实的本性，借着理智与肉身相交。祂较低的本性，即人性，成了天主，因为它与天主结合，并成为一位，因为较高的本性得胜，为使我也能成为神，正如祂成为人。祂被生——但祂是受生的；祂由女人所生——但她是童贞女。前者是人性的，后者是天主性的。在祂的人性内，祂没有父亲，但也在祂的天主性内，祂没有母亲。这两点都属于天主性。祂居住在母胎中——但祂被那仍在自己母胎中的先知认出了，他为\OriginalTerm{圣言}{Word}的缘故而跳跃，因圣言他得以存在。祂被包裹在襁褓中\BibleRef{路 2:41}——但祂的复活脱去了坟墓的缠带。祂被放在马槽中——但祂被天使光荣，被星宿宣告，被\OriginalTerm{贤士}{Magi}朝拜。你们为什么因呈现在你们眼前的事物而绊倒，因为你们不愿看那呈现在你们心智中的事物？祂被驱逐到埃及——但祂驱逐了埃及的偶像。在犹太人眼中，祂没有威仪，也没有俊美\BibleRef{依 53:2}——但在达味看来，祂比世人更俊美。在山顶上，祂明亮如闪电，比太阳更明亮\BibleRef{玛 17:2}，向我们启示未来的奥秘。\par

二十. 祂作为人受洗——但祂作为天主赦罪——并非祂自己需要净化礼节，而是为圣化水元素。祂作为人受试探，但祂作为天主得胜；是的，祂令我们放心，因为祂已战胜了世界。\BibleRef{若 16:33} 祂饿了——但祂喂饱了数千人；是的，祂是生命之粮，且是从天降下的。祂渴了——但祂呼喊说：\BibleQuote{凡渴的，请到我这里来喝}。是的，祂应许从信者身上要流出活水的江河。祂疲倦了，但祂是疲惫和负重者之安息。\BibleRef{玛 11:28} 祂沉睡，但祂在海上轻松行走。祂斥责风，使伯多禄在开始下沉时变轻。祂纳税，但取自鱼腹；是的，祂是那些要求者之王。\BibleRef{若 19:19} 祂被称为\OriginalTerm{撒玛黎雅人}{Samaritan}和\OriginalTerm{附魔者}{demoniac}；——但祂拯救了那从耶路撒冷下来、落在强盗手中的人；恶魔承认祂，祂驱逐恶魔，将成群污鬼沉入海中，\BibleRef{路 8:28} 并看见\OriginalTerm{魔君}{the Prince of the demons}像闪电般坠落。祂被石头砸，但未被捉住。祂祈祷，但祂垂听祈祷。祂哭泣，但祂止息眼泪。祂问拉匝禄安葬何处，因为祂是人；但祂复活拉匝禄，因为祂是天主。\BibleRef{若 11:43} 祂被出卖，且极廉价，仅值三十块银钱；\BibleRef{玛 26:15} 但祂赎回世界，且付上高价，因为代价是祂自己的血。\BibleRef{伯前 1:19} 祂如羊被牵到宰杀之地，\BibleRef{依 53:7} 但祂是\OriginalTerm{以色列的牧者}{the Shepherd of Israel}，如今也是全世界的牧者。祂如羔羊沉默不言，但祂是\OriginalTerm{圣言}{the Word}，由那在旷野中呼喊者的声音所宣告。\BibleRef{若 1:23} 祂被击打和受伤，但祂医治各种疾病和灾殃。\BibleRef{依 53:23} 祂被举起钉在木架上，但借着\OriginalTerm{生命树}{the Tree of Life}祂恢复我们；是的，祂甚至拯救了与祂同钉的强盗；\BibleRef{路 23:43} 是的，祂使可见的世界陷入黑暗。他们给祂喝搀着苦胆的醋。谁？那曾将水变成酒的那一位，\BibleRef{若 2:1} 祂是苦味的摧毁者，祂是甘甜和全然可爱者。\BibleRef{歌 5:16} 祂舍弃生命，但祂有权再取回；\BibleRef{若 10:18} 圣所帐幔裂开，因为天上的奥秘之门打开了；磐石崩裂，死者复活。\BibleRef{玛 27:51} 祂死了，但祂赐予生命，且以死亡毁灭死亡。祂被埋葬，但祂复活了；祂下到阴府，但祂带出灵魂；祂升天，并将再来审判生者死者，试验你们这样的话语。若前一部分给你们错误之起点，愿后一部分终止它。\par

二十一. 这就是我们对那些想要困惑我们的人的答复；并非心甘情愿地给出（因为轻率的言谈和词语的矛盾并非忠信者所喜，有一个仇敌对我们已足够），而是出于必要，为了我们的攻击者（因为药物因疾病而存在），好使他们能被引导看到，他们并非全智，也非在那些使福音无效的冗长论证中不可战胜。因为当我们停止相信，只靠论证的力量保护自己，并以质疑摧毁\OriginalTerm{圣神}{the Spirit}在我们信德上的要求，然后我们的论证对于主题的重要性又不够强大（这必然发生，因为它是通过这样一个微小的器官即我们的心智来发动的），结果会是什么？论证的软弱似乎属于奥秘本身，于是口才的优雅使十字架无效，如保禄也认为的。\BibleRef{格前 1:17} 因为信德是完成我们论证的。但愿那宣布结合、释放被捆绑者，并给我们心智能力解开他们非自然教义之结的那一位，若可能，改变这些人，使他们成为有信德而非修辞学家，成为基督徒而非他们如今被称为的那样。我们确实为此恳求，因基督的缘故。\BibleQuote{你们要与天主和好}，\BibleRef{格后 5:20} 且不要\BibleQuote{熄灭圣神}；\BibleRef{得前 5:19} 或者，更确切地说，愿基督与你们和好，愿圣神光照你们，虽然如此迟延。但若你们过于喜爱争执，我们无论如何将持守\OriginalTerm{天主圣三}{the Trinity}，愿我们因天主圣三而得救，保持纯洁无玷，直到我们渴望的那更为完美的显现，在祂内，我们的主基督，愿光荣归于祂，直到永远。阿们。\par

\SourceRecord{Translated by Charles Gordon Browne and James Edward Swallow.；Nicene and Post-Nicene Fathers, Second Series；Vol. 7.；Edited by Philip Schaff and Henry Wace.；Buffalo, NY: Christian Literature Publishing Co.,；1894.；<http://www.newadvent.org/fathers/310229.htm>.}

% structure: p0001 source=h1 target=section reason=page-title

\section{第四篇神学演讲（演讲三十）}

\Emph{第二篇：论圣子}\par

一、我既已藉圣神之能，充分推翻了他们论证的微妙与错综之处，并已从总体上解决了从圣经引来的异议与反驳——这些亵渎的窃夺圣经者与偷窃经义者，正是借此将群众拉拢到他们一边，并混淆真理之路的；而且我做得并非不清楚，我相信所有公正的人都会这样说：将较高和较神圣的表述归于天主性，将较低和较人性的表述归于那为我们人类而成为第二亚当、且为对抗罪恶而成为能受苦的天主的祂；但由于我们论证的匆促，尚未逐条细论这些经文。既然你们要求我们对每一条给出简要的解释，以免被他们论证的表面合理所迷惑，那么我们便摘要说明，将其分条编号以便记忆。\par

二、在他们看来，以下经文是太现成了：\BibleQuote{主在祂道路的起始创造了我，为祂的作为之先}\BibleRef{箴 8:22}。我们如何应对？我们要控告\Person{撒罗满}吗？还是因他后来堕落而拒绝他先前的话？我们要说，这话是智慧本身所说，仿佛知识以及创造者圣言所说——万物都是按照这话而造的？因为圣经常将众多甚至无生命的东西拟人化；例如：\Quote{海说了如此这般的话}；\Quote{深渊说，不在我内}；\Quote{诸天宣扬天主的荣耀}；又向刀剑发出命令；且问群山与丘陵为何跳跃。我们不援引任何这类例证，尽管我们的一些前辈曾用它们作为有力的论证。但让我们承认，这话是用在我们救主、真智慧身上。我们一起思想一个微小的问题。在一切存在物中，哪一个是无始的？天主性。因为无人能说出天主的起源，否则那起源就比天主更古老了。但天主为我们而取的人性，其缘由是什么？无疑是我们得救。还能是什么？既然我们在此清晰地同时发现\OriginalTerm{受造}{Created}与\OriginalTerm{受生}{Begetteth Me}，那么论证便简单了：凡我们发现与一个原因相连的，我们就归于人性；而一切绝对和无始的，我们就归于祂的天主性。那么，这\Quote{受造}岂非是与一个原因相连而说的？\Quote{祂在祂道路的起始创造了我，为祂的作为之先}，正是如此说。祂双手的作为是真理与审判；为它们的缘故，祂被天主性所傅油；因为此傅油是属人性的；但\Quote{祂生我}并不与一个原因相连；否则你们要指出它的附属成分。那么，还有什么论证能推翻以下观点：智慧被称为受造物，是就较低层次的出生而言；但就第一和更难理解的出生而言，祂是受生的？\par

三、接下来是祂被称为仆人并服事许多人的事实，以及被称为天主的儿子是一件大事。因为事实上，祂为使我们和所有在罪恶捆绑下得祂拯救的人得自由，曾服事过肉身、出生和我们生活的种种条件。人性能与天主相融合，并藉此融合而神化，我们能被来自高天的旭日如此眷顾，以至于那要诞生的圣者被称为至高者的儿子，并赐给祂一个超越万名的名字——除了天主这个名字还能是什么？——并且万膝都向那为我们而被贬抑的、兼有天主性体与奴仆形体者跪拜，且所有以色列家族都应知道，天主已立祂为主、为基督。这一切都是藉受生者的行动和生祂者的美意而成就的。\par

四、那么，他们那第二个不可抗拒的伟大经文是什么？\BibleQuote{祂必须为王}\BibleRef{格前 15:25}，直到某一时候……\BibleQuote{被天收留直到万物复兴的时候}\BibleRef{宗 3:21}，\BibleQuote{坐在右边直到仇敌被推翻}\BibleRef{咏 110:1}。但此后呢？祂必须停止为王，或被移出天堂吗？谁会、或为何要使祂停止？你真是个大胆而无政府主义的解释者；然而你曾听到，\BibleQuote{祂的王权没有终结}\BibleRef{路 1:33}。你的错误在于不理解\Quote{直到}一词并非总是排除之后的事，而是肯定\Quote{直到}那时刻，而不否认之后的事。举一个例子——你如何理解\BibleQuote{看，我天天与你们在一起，直到今世的终结}\BibleRef{玛 28:20}？难道意味着祂之后不再如此？理由何在？但这并非你错误的唯一原因；你也未能区分所指之事。祂被称为王，一方面作为全能之君，统御愿意和不愿意者；另一方面则在使我们服从，使我们甘愿承认祂的主权而置于祂的王权之下。从第一种意义上说，祂的王权没有终结。但就第二种意义而言，有何终结？当我们进入得救的状态后，祂接受我们为仆人。因为当我们已经顺服时，何必再使我们顺服？此后祂起来审判大地，将得救者与丧亡者分开。此后，祂要像神一样站在群神之中，即站在得救者之中，分别并决定每一个人配得何等尊荣和居所。\par

五. 其次，请考虑你所提出的\OriginalTerm{屈从}{subjection}，藉此你将子屈从于父。你问，现在祂不是屈从吗？或者如果祂是天主，祂必须屈从天主吗？你正在构思你的论点，仿佛它涉及某个强盗或某个敌对的神祇。但请这样看待它：因我的缘故，祂被称为\OriginalTerm{诅咒}{curse}，祂摧毁了我的诅咒；以及\OriginalTerm{罪}{sin}，祂除免了世罪；并成为\OriginalTerm{新亚当}{new Adam}以代替旧亚当；同样，祂将我的\OriginalTerm{悖逆}{disobedience}视为自己的，作为全身之首。因此，当我悖逆和反叛时，既因否认天主又因我的私欲偏情，基督也因我的缘故被称为悖逆。但当万有一方面因承认祂，另一方面因改造而屈从于祂时，祂自己也将完成祂的屈从，带领祂所拯救的我到天主面前。因为在我看来，这就是基督的屈从：即承行父的旨意。但如子使万有屈从于父，父也如此使万有屈从于子；一位藉其工作，另一位藉其美意，如我们已说过的。因此，那使万有屈从者便将所屈从的呈献给天主，将我们的状况视作自己的状况。在我看来，同类的表达是：\BibleQuote{我的天主，我的天主，你为什么舍弃了我？} 并非祂被父舍弃，也非被祂自己的天主性舍弃，如某些人所想，好像天主性畏惧苦难，因此在祂受难时抽离自身（因为谁迫使祂降生世上或被举在十字架上呢？）但如我所说，祂亲自代表我们。因为我们从前是被舍弃和被轻视的，但现在藉着那不能受苦者的苦难，我们被接纳并得救。同样，祂将我们的愚昧和过犯视为自己的，并说出圣咏接下来说的话，因为显然\BibleBook{圣咏}{第二十一篇}是指基督而言。\par

六. 同样的考虑适用于另一处经文：\BibleQuote{由所受的苦难，学习了服从}，以及祂的\Quote{大声哀号和眼泪}、祂的\Quote{恳求}、祂的\Quote{蒙应允}、祂的\Quote{虔敬}，所有这些祂都奇妙地演绎出来，如同一出为我们而构思的戏剧。因为作为圣言，祂既非服从也非不服从。因为这样的表达属于仆人、属下，前者适用于较好的仆人，后者属于应受惩罚的仆人。但作为\OriginalTerm{奴仆的形体}{form of a servant}，祂\OriginalTerm{屈尊就卑}{condescends}与同仆，甚至与仆人同等，并接受了陌生的形态，在自身内担负着我和我的一切，为在祂内耗尽恶，如火消融蜡，或如太阳驱散地上的雾气，并使我藉着\OriginalTerm{结合}{blending}分享祂的本性。因此，祂用行动尊崇\OriginalTerm{服从}{obedience}，并藉苦难经验地证明了它。因为拥有意愿是不够的，正如我们若不藉行为证明它也是不够的；因为行为是意愿的证明。\par

而且，也许这样假设也不错：藉着祂爱人的技巧，祂衡量我们的\OriginalTerm{服从}{obedience}，并以自己的苦难为比较标准来衡量一切，以便祂藉自己的经历认识我们的状况，知道向我们要求多少，我们付出多少，并将我们的软弱与环境一起考虑。因为若透过面纱照耀黑暗，即此世的光，被另一个黑暗（我指恶者和诱惑者）迫害，那么黑暗作为更软弱的，岂不更受迫害？这有什么可惊奇的呢？尽管祂完全逃脱了，我们却至少部分被追上了？因为祂被迫逐比我们被捕获更奇妙——至少对所有对此事深思熟虑的人是这样。我还要在所提及的经文之外另加一段，因为我认为它明显指向同一意思。我说的是：\BibleQuote{祂既然亲自受试探而受苦，就能扶助那些受试探的人}。但天主在万物\OriginalTerm{复兴}{restitution}的时候将成为\OriginalTerm{万有中的万有}{all in all}；这不是指只有父存在，而子完全消融于父，如同火炬投入大火堆，从那里被夺走片刻又放回（因为我不愿甚至\OriginalTerm{撒伯里乌派}{Sabellians}因这种说法而受害）；但整个天主性，当我们不再被分割（如我们现今被运动和私欲偏情所分割），并且丝毫不含有天主，或极少含有，而是完全相似时。\par

七. 作为你的第三个论点，你算到圣言是\OriginalTerm{更大的}{Greater}；至于第四个，到我的天主和你们的天主。的确，如果祂被称为更大的，而\Quote{相等}一词不曾出现，这或许可能成为有利于他们的论据。但如果我们发现两个词都明确使用了，这些先生们会怎么说呢？这如何能加强他们的论点？他们如何调和不可调和之事？因为同一事物同时大于且等于同一事物是不可能的；显然的解答是：\Quote{更大}指\OriginalTerm{起源}{origination}，而\Quote{相等}属于\OriginalTerm{本性}{Nature}；我们很乐意承认这一点。但也许有人会支持我们对你论点的攻击，并断言：由这样的\OriginalTerm{原因}{Cause}而来的那一位，并不低于那无\OriginalTerm{原因}{Cause}的一位；因为祂将分享\OriginalTerm{无始者}{Unoriginate}的光荣，因为祂源于无始者。此外，还有\OriginalTerm{受生}{Generation}，这对所有人来说都是如此奇妙而庄严的事。因为说祂大于被视为人的子，确实是真的，但并非什么大事。如果天主大于人，有什么稀奇呢？这确实足以回应他们关于\Quote{更大}的谈论。\par

八. 至于其他经文，\Quote{我的天主}一词所用，不是指圣言，而是指可见的圣言。因为那真正是天主者的天主，怎能存在呢？同样，他不是可见者的父，而是圣言的父；因为我们的主有两个性体；以致在这两种情形中，每一情况中一个用语是恰当使用，另一个是不恰当使用；但其方式与我们在我们身上使用它们的方式恰好相反。因为关于我们，天主恰当地是我们的天主，却不恰当地是我们的父。这就是异端者错误的起因，即这两个名号的合并，它们因性体的结合而互换。这点的迹象可见于：在我们的思想中，每当性体彼此区分时，名号也就区分；正如你听到保禄的话：\BibleQuote{我们主耶稣基督的天主，荣耀的父。}基督的天主，却是荣耀的父。因为虽然这两个用语表达同一个位格，但这并非由于性体的合一，而是由于两者的结合。还有什么比这更清楚的呢？\par

九. 第五，有人可能主张，论及他说他领受生命、审判、外邦人的产业、或掌管一切血肉的权能、或荣耀、或门徒、或任何其他提到的事物。这同样属于人性；然而，你若将之归于天主性，也并无荒谬。因为你不会如此归于，好像这是新获得的，而是由于本性从一开始就属于他，并非出于恩赐。\par

十. 第六，有人可能断言，经上记载：\BibleQuote{子凭自己不能作什么，只有看见父作什么，他才能作。}对此的解答如下：能与不能并非只有一种含义，而是有多种含义。一方面，它们有时用于指力量不足，有时用于指时间方面，有时相对于某个对象；例如，一个孩子不能成为运动员，或一只小狗不能看见，或与某某打斗。也许有一天，孩子会成为运动员，小狗会看见，会与那个打斗，尽管它可能仍然不能与别的打斗。或者，它们可用于指普遍真实的事情；例如，\BibleQuote{建在山上的城是不能隐藏的；}然而，它却可能被另一座更高的山与之成一直线而隐藏。或在另一种意义上，它们用于不合理的事；例如，\BibleQuote{当新郎与伴郎同在时，伴郎们岂能禁食？}无论他是被视为以身体形可见的（因为他在我们中间居住的时期不是哀伤而是喜乐的），还是作为圣言。因为那些被圣言洁净的人为何要守身体的禁食呢？或者，它们又用于指违背意愿的事；例如，\BibleQuote{他因他们的不信，在那里不能行大能的事，}即那些应接受奇迹的人。因为为了治愈，既需要病人的信德，也需要医治者的能力，当两者之一缺乏时，另一个就不可能。不过，这个意义大概也应归入不合理之类别。因为对于那些后来会因不信而受损的人，治愈是不合理的。这话：\BibleQuote{世界不能恨你们，}也属于同一类别，同样地：\BibleQuote{你们既是恶的，怎能说出善来？}因为就哪一种意义而言，两者是不可能的，除非是违背意愿？在表明本不可能的事，若天主愿意就可能的话语中，也有某种类似的意义；例如，\BibleQuote{人不能重生，}或\BibleQuote{针眼不能让骆驼穿过。}因为如果天主愿意，有什么事能阻止这两件事发生呢？\par

十一、除此之外，还有我们现在所研究的完全不可能且不可接受之事。因为我们主张，天主不可能为恶，也不可能不存在——这表示天主的软弱而非刚强——亦不可能使不存在者存在，或使二加二既等于四又等于十；同样，子行父所不行之事，也是不可能且不可思议的。因为凡父所有的，都是子的；反之，凡属于子的，也都是父的。因此没有什么是独有的，因为一切都是共有的。因为他们的存有本身是共有且平等的，即使子是由父领受了它。正是为此，经上说：\BibleQuote{我因父而生活}；这不是说他的生命和存有是靠父维系，而是因为他从父那里获得他的存有，超越一切时间和原因。但是，他如何看见父做什么，他也照样做呢？是否如同那些临摹图画和文字的人，因为他们若不注视原本并受其引导，便无法达到真实？然而智慧岂需要教师，或未经教导便不能行事？父在现在或过去\Quote{做}什么？他是否在现今这个世界之前创造了另一个世界，或将要创造一个未来的世界？子是否注视了那个世界而创造了这个世界？或者他望去另一个世界，而做一个相像的？按此论证，必然有四个世界：两个由父创造，两个由子创造。何等荒谬！他洁净癞病人，从魔鬼和疾病中解救人们，复活死者，在海面上行走，并行了其他一切工程；但是，父在哪一种情形下或何时在他之前行了这些事？难道不是父赋予了这些同样行动的理念，而圣言使之实现，却不是以奴隶或无技巧的方式，而是以全知和卓越的方式，或者更恰当地说，如同父一样？因为在这意义上，我理解那些话：\BibleQuote{凡父所做的，子也照样做}；不是因为所做之事的相似，而是出于权柄。这可能也是那段经文的意义，即：\BibleQuote{父做事直到现在，子也做事}；不仅如此，它亦指他对他所造之物的治理和保存；正如经上所说：\BibleQuote{他使他的天使成为神风}，\BibleQuote{大地坚定在基础上}（尽管这些都一次性地确立并造成了），以及\BibleQuote{建立雷霆，创造风}。对于这一切，圣言一次性地被赐予，但行动至今仍在持续。\par

十二、让他们在第七位引用：子从天降下，不是为了行他自己的旨意，而是为了行那派遣他者的旨意。如果这话不是由那降下者亲口所说，我们会说这措辞是由自人性发出的，而非出自那被视为救主者，因为他人性意志不可能与天主相悖，因为它已完全被纳入天主；但若单纯地设想在我们本性中，人意志不完全随从天主，而是时常与之斗争和抵抗。因为我们以同样的方式理解这些话：\BibleQuote{父啊！若是可能，就让这杯离开我吧！但不要照我的意愿，而要照你的意愿}。因为他不像不知道这是否可能，或以意志反对意志。但由于这是那取了我们本性的那一位（因为正是他降下）的语言，而非他所取本性的语言，我们必须以这种方式回应这种反对：这段经文并非说子有他自己的特殊意志，有别于父的意志，而是说他根本没有；因此其含义是：\Quote{不是要行我自己的旨意，因为除了那属于我和你的共同旨意外，没有我自己的旨意；因为我们既有一个天主性，故我们也有一个旨意。}因为许多这样的表达是与这共有性相关，并且不是以肯定方式，而是以否定方式表达的；例如：\BibleQuote{天主不赐予圣神，因为没有限度}（事实上，他没有把圣神赐给子，也没有\Emph{度量}它，因为天主不被天主度量）；又如：\BibleQuote{不是我的过犯，也不是我的罪恶}。这些话不是因为他有这些事，而是因为他没有这些事。又如：\BibleQuote{不是因我们所行的义德}，因为我们并未行过任何义德。在后面的经文中也可清楚看到这一含义。因为他说：\BibleQuote{我父的旨意是什么呢？就是要凡相信子的人，获得拯救，并得到最终的复活。}那么，这是父的旨意，而不是子的吗？难道他违背自己的意志宣讲福音、接受人的信仰？谁能相信呢？此外，那段说所听的圣言不是子的，而是父的经文，也具有同样的力量。因为我不能明白如何两者所共有之物，能说成仅属于其中一位，无论我怎样思考，我也不认为有谁能明白。因此，你若对旨意持这种看法，那么你的看法便是正确且虔敬的，如我所想，也如所有正直之人所想。\par

XIII. 第八段经文是：\BibleQuote{使他们认识祢，唯一的真天主，和祢所派遣的耶稣基督} \BibleRef{若 17:3} 以及\BibleQuote{除了天主一个外，没有谁是善的} \BibleRef{谷 10:18}。在我看来，解答十分容易。因为如果你将这句话仅归于圣父，那么\Quote{\OriginalTerm{真理本身}{Very Truth}}将置于何处？如果你如此理解\Quote{唯一的智者天主} \BibleRef{罗 16:27}、\Quote{独享不死不灭，住在不可接近的光中} \BibleRef{弟前 6:16}、或\Quote{永世的君王，不朽不灭，不可见的唯一智者天主} \BibleRef{弟前 1:17}的含义，那么圣子就处在死亡或黑暗的判决之下，至少被判定为既不智、也不为王、不可见、甚至根本不是天主，而这些点概括了所有。你又如何能阻止祂的\Quote{善}——这尤其唯独属于天主——与其他一切一同消失？然而，我认为经文\BibleQuote{使他们认识祢，唯一的真天主}是为了推翻那些冒称为神的假神，因为祂不会加上\BibleQuote{和祢所派遣的耶稣基督}，如果\BibleQuote{唯一的真天主}是与祂相对比，而这句话不是在共同天主性的基础上来说的。\BibleQuote{没有谁是善的}是针对那试探的法律学士，他视为人子的基督之善而作见证。因为祂说，圆满的善唯独属于天主，即使一个人也被称为完全的善。例如：\BibleQuote{善人从心里所存的善就发出善来} \BibleRef{玛 12:35}；以及\BibleQuote{我要将王国赐给一个比你更好的人} \BibleRef{撒上 15:28}……这是天主论及撒乌耳时所说话。又如：\BibleQuote{上主，求祢善待善人} \BibleRef{咏 125:4}，以及所有类似对我们这些受称赞之人的表达，他们是从至善者涌流而出的一种溢出，并以次要方式降临于他们。最好的做法是，如果我们能说服你相信这一点。但若不能，对这一相反的建议你又能说什么：按照你的假设，圣子也被称为唯一的天主。在哪里？请看这句：\BibleQuote{这就是你的天主，没有别的可比拟祂} \BibleRef{巴 3:35}，稍后又言：\BibleQuote{此后祂在地上显现，与人往来} \BibleRef{巴 3:37}。这个附加证明清楚表明，这些话不是论及圣父，而是论及圣子；因为正是祂以肉体形式与我们交往，并在这下界之中。如果我们决定将这些话理解为与圣父对比，而非与假神对比，那么我们就失去了圣父，所用的正是我们曾用以攻击圣子的那些措辞。还有什么比这种胜利更灾难性的呢？\par

XIV. 第九点，他们引述\BibleQuote{祂常活着，为我们转求} \BibleRef{希 7:25}。哦，何其美好、奥秘且仁慈！因为\Quote{转求}并不意味着寻求报复，如大多数人的方式（因为那会带有卑微），而是因祂的中介身份 \OriginalTerm{中介身份}{Mediatorship} 为我们恳求，正如圣神也被说成为我们代求 \BibleRef{罗 8:26}。因为\BibleQuote{只有一个天主，也只有一个天主与人之间的中保，就是基督耶稣，祂是人} \BibleRef{弟前 2:5}。因为祂如今仍以人身份为我的得救而恳求；祂继续穿着祂所摄取的肉体，直到祂以降生成人的德能使我成为天主；虽然祂不再按肉体被认识——我是指如同我们的一样（除了罪之外）的肉体情欲。因此，我们有\BibleQuote{\OriginalTerm{辩护者}{Advocate}耶稣基督} \BibleRef{若一 2:1}，祂并不是在圣父面前为我们卑躬屈膝、奴性地俯伏……摒弃如此奴性且不配圣神的猜疑！因为这种猜疑既不适合圣父要求，也不适合圣子接受；对天主如此设想是不义的。但祂作为人所受的苦难，祂作为圣言和\OriginalTerm{策士}{Counsellor}劝勉圣父忍耐。我认为这就是祂转求 \OriginalTerm{转求}{intercession} 的含义。\par

XV. 第十个异议是祂的无知 \OriginalTerm{无知}{ignorance}，以及这句话：\BibleQuote{至于那日子和那时刻，没有人知道，连天上的天使也不知道，连子也不知道，只有父知道} \BibleRef{玛 24:36}。然而，智慧 \OriginalTerm{智慧}{Wisdom} 怎么可能不知道任何事？——智慧是创造诸世界的、完善它们的、改造它们的、是一切受造物的界限 \OriginalTerm{界限}{Limit}、知道天主的事如同人的灵知道人的事？因为还有什么比这种知识更完美的呢？那么，你怎能说祂准确知道那时刻之前的一切，以及关于末期将要发生的一切，却不知道那个时刻本身？因为这样的事就像谜语一般；好比有人说他准确知道墙前的一切，却不知道墙本身；或者，知道白天的尽头，却不知道夜晚的开始——而对一者的认识必然带来对另一者的认识。因此，每个人都必须看见：祂作为天主知道，作为人却不知道——如果我们可以将可见的与仅由思维所辨识的分开的话。在这段经文中，\Quote{子}这个名称的绝对且无条件的用法（没有加上\Quote{谁的子}）给了我们这个想法，即我们应当以最虔敬的方式理解这无知，将其归于人性，而非天主性。\par

十六、若此论证已足，我们就此止步，不再深入追问。但若不足，则我们第二个论证如下：——正如我们在其他一切事上所做的那样，我们也要将祂对最大事件的知悉，为尊崇圣父而归於那原因。我认为，任何人，即便不像我们一位学生那样解读，也会很快看出，就连圣子知道那日子和时辰，也不异于圣父。因为我们由此得出什么结论呢？即既然圣父知道，那么圣子也知道，因为显然只有第一性才能认知和领会这事。我们还有必要解释关于祂领受诫命、遵守诫命、常做祂喜悦的事的经文；以及关于祂被成全、被高举、藉所受苦难学习服从；还有祂的大司祭职、祂的祭献、祂的被出卖、祂向那能救祂免死的祈祷、祂的痛苦、血汗及祷告，等等；除非对每个人而言，这些话语显然不是关乎那不变且超乎一切受苦能力的性体，而是关乎那能受苦的人性。所以，关于这些反对意见的论证，至此作为对那些更善于将探讨推至更完满解决者的基础与备忘。然而，也许值得，并且符合已说的话，不略过圣子的实际称号（它们很多，涉及祂众多属性），而是将每个称号的含义摆在你面前，并指出这些名称的奥义。\par

十七、我们如此开始：天主性无法以言语表达。这不仅由论证证明给我们，也由最智慧、最古老的希伯来人，就他们给予我们推测的根据而言。因为他们将某些字符专属於尊崇天主性，甚至不允许任何低于天主之物的名字与天主的名字用相同的字母书写，因为他们认为，天主性甚至在那程度上让任何受造物分享自身是不合适的。那么他们又怎能允许那不可见且分离的性体由可分的言语解释呢？因为从未有人呼出整个空气，也没有任何心智完全理解或言词穷尽地包含天主的存有。但我们藉其属性描绘祂，由此获得关于祂的某种模糊、微弱且部分的观念；我们最好的神学家是那一位，并非发现了全部（因为我们现时的锁链不允许我们看到全部），而是比他人更多领悟了祂，并在自身内收集了更多相似或真理的影像，或随我们如何称呼它。\par

十八、就我们所能达到的而言，\InnerQuote{祂是}和\InnerQuote{天主}是祂本质的专属名字；其中特别\InnerQuote{祂是}，不仅因为当祂在山上传话给梅瑟、梅瑟问祂的名字时，祂如此自称，命他告诉人民\Quote{那\InnerQuote{我是}（\BibleQuote{我是自有者}）派遣了我}(\BibleRef{出 3:14}），还因为我们发现这个名字更为严格切当。因为\OriginalTerm{神}{Qeo/j}这个名字，即使如那些精通此事者所说，源于\OriginalTerm{奔跑}{Qe/ein}或\OriginalTerm{燃烧}{Ai!qein}，源自不断的运动，并因祂吞噬恶的事物（由此祂也被称为吞噬之火），它仍是一个相对的名字，而非绝对的名字；同样，主也是一个名字，它也被称为天主的名。祂说：\BibleQuote{我是上主，你的天主}，这是我的名；\BibleQuote{上主是他的名}。但我们所探究的是一个其存有是绝对的性体，而非与别的东西相连的存有。但\InnerQuote{存有}在其本义上专属天主，完全属于祂，不受任何之前或之后所限制或截断，因为於祂没有过去或未来。\par

十九、在其它称号中，有些显然属于祂权威的名字，有些属于祂眷顾世界的名字，而且此眷顾从双重角度看待，一在降生成人之前，一在之后。例如，全能者、光荣的君王、或诸时代的君王、或万军的君王、或受爱者的君王、或诸王的君王。又如万军的上主，即万军的、或万能的、或诸主的上主；这些显然属于祂权威的称号。但拯救的天主、或报应的天主、或和平的天主、或正义的天主；或亚巴郎、依撒格、雅各伯的天主，以及所有看见天主的属神以色列——这些属于祂的眷顾。因为我们被这三件事所治理：对惩罚的恐惧、对拯救及随之而来的荣耀的希望、以及获得这些的美德的实践，所以报应的天主的名管辖恐惧，拯救的天主的名管辖我们的希望，美德的天主的名管辖我们的实践；以便任何达到其中任何一项的人，因携带天主在自己内，更能趋向成全，并趋向由美德而来的亲密关系。现在这些是天主性共有的名字，但无始者（\OriginalTerm{无始者}{Unoriginate}）的专名是父，无始而生者（\OriginalTerm{无始而生者}{unoriginately Begotten}）的专名是子，非受生而发者（\OriginalTerm{非受生而发者}{unbegottenly Proceeding}）的专名是圣神。现在让我们继续谈到圣子的名字，这是我们这部分论证的起点。\par

XX. 按我的意见，祂之所以被称为子，是因为祂在本质上与圣父同一；且不仅为此，还因为祂出自圣父。祂被称为独生子，不仅因为祂是唯一的子、惟独出自圣父、且只是子，也因为祂的子职方式为祂所独有，不为身体所分享。祂被称为圣言，因为祂与圣父的关系如同言语与心智；这不仅出于祂无受生的生成，也出于结合及宣示的功能。也许这种关系可比作定义与被定义者之间的关系，因为后者也被称为\OriginalTerm{圣言}{Lo/goj}。因为经上说，那在理智上领悟了子的（因为\Quote{看见了}意即如此），也领悟了圣父；而子是圣父本性的简明证明和容易的陈述。因为凡受生者，皆是生者之无声言语。若有人说这名号是因祂存在于一切万有之中而赐予祂的，\par

他也不算错。因为有什么事物不是藉着圣言而存在的呢？祂也被称为智慧，作为对神事与人事的知识。因为那创造万有的，如何能不知道祂所造之物的道理呢？又被称为能力，作为万有受造物的维持者，并赐予它们保持自身的力量。又被称为真理，因其本性为一而非多（因为真理是一，虚假是多），并作为圣父纯洁的印记及最无谬误的印鉴。又被称为肖像，因其与祂同体，且因祂出自圣父，而非圣父出自祂。因为肖像的本性正是如此：是其原型的复制品，并属于那以其名称之者；只是在此处有更多含义。因为在日常用语中，肖像是活动之物的静止表现；但在此处，却是永生者的活复制品，比塞特之于亚当，或任何儿子之于父亲，更为精确地相似。因为单纯存有的本性正是如此：不能说它们在某方面相似、另一方面不相似；而是完全相似，更应称为同一而非相似。再者，祂被称为光，作为那些藉着言语和生命得以洁净的灵魂之光明。因为若无知和罪是黑暗，那么知识和虔敬的生活便是光。……祂被称为生命，因为祂是光，是一切理性灵魂的构成和造化能力。因为我们在祂内生活、行动、存在，这是依据那吹入我们的双重能力；因为我们都由祂的气息所感，而凡能领受的，并按照我们开启理智之口的程度，也由天主圣神所感。祂是正义，因为祂按我们所应得的分配，并是公义的仲裁者，既为在法律之下的人，也为在恩宠之下的人，为灵魂和身体，好使前者管理，后者服从，使更高的统御更低的；免得较差的背叛较好的。祂是圣化，作为纯洁，使纯洁的能为纯洁所容纳。祂是救赎，因为祂释放我们这些被罪所俘虏的，将祂自己作为赎价给了我们，作为赎世祭献。祂是复活，因为祂从此世提升我们，并使我们这些被罪所杀的获得新生。\par

XXI. 然而，这些名号仍是那超越我们者与那为我们而来的那一位所共有的。但还有一些却是我们特有的，属于祂所取的本性。因此祂被称为人，不仅为使藉着身体，有形体的受造物能领悟祂——否则因祂不可理解的本性，这是不可能的——也为使祂藉着自己圣化人性，如同酵母对整个面团一般；并藉着将自身与被定罪者结合，释放它脱离一切定罪，为我们成了所是的一切，除了罪——身体、灵魂、心智，以及一切藉以死亡者——如此祂成了人，即这一切的结合；有形可见的天主，因为祂保留了那惟独理智所察觉的。祂被称为人子，既因亚当，也因祂所出的童贞女；一位是先祖，另一位是母亲，既符合生育之律，又超越它。祂是基督，因祂的天主性。因为这是祂人性的受傅，并不像其他所有受傅者那样，以行动圣化，而是以受傅者的圆满临在；其效果是，那傅油的被称为人，并使那受傅的成为天主。祂是道路，因为祂藉着自己引导我们；是门，因祂让我们进入；是善牧，因祂使我们安卧在青草地上，领我们到憩息之水旁，引导我们，保护我们免于野兽，回转迷途者，领回失落的，包扎破碎的，守护强壮的，并在将来的羊栈中将他们聚集，以牧灵的知识。是羊，作为牺牲；是羔羊，作为完美的；是大司祭，作为奉献者；是默基瑟德，在超越我们的本性中无母，在我们的人性中无父；且没有族谱（因为，经上说，\BibleQuote{谁能述说祂的世代？}）；此外，是撒冷王，即和平之王，正义之王，并从圣祖们收取十分之一，当他们战胜恶势力时。这些是子的名号。请行经它们：那些崇高的，以神性的方式；那些属于身体的，以适合它们的方式；或更好说，全然以神性的方式，使你成为神，从下而上，为了那从高处降下为我们之故的那一位。在一切中并超越一切，持守这一点，你便永不会错，无论在高超的名号或卑微的名号中；耶稣基督昨天、今天，在降生中，以及在圣神内，直到永远，始终如一。阿们。\par

\SourceRecord{Translated by Charles Gordon Browne and James Edward Swallow.；Nicene and Post-Nicene Fathers, Second Series；Vol. 7.；Edited by Philip Schaff and Henry Wace.；Buffalo, NY: Christian Literature Publishing Co.,；1894.；<http://www.newadvent.org/fathers/310230.htm>.}

% structure: p0001 source=h1 target=section reason=page-title

\section{\WorkTitle{第五篇神学演讲（演讲31）}}

第五篇神学演讲。论圣神。\par

如此便是关于\OriginalTerm{圣子}{Son}的记述，他以此方式逃脱了那些要石击他的人，从他们中间经过。因为\OriginalTerm{圣言}{Word}不被石击，却在他乐意时投石；并用机弦攻击野兽——那就是言语——以不圣洁的方式接近那山。但是，他们继续说：\Quote{关于\OriginalTerm{圣神}{Holy Ghost}，你有什么可说的？从何处给我们引进这位陌生的天主，圣经对他默不作声？}甚至连那些对\OriginalTerm{圣子}{Son}保持分寸的人也这样说。正如我们在道路和河流的情况中所见，它们彼此分开又汇合，同样，在这件事上，由于不敬之心的过度，各方面都不同的人也在某些点上达成一致，以至你永远无法肯定他们究竟是同心一致，还是在相互冲突。\par

现在圣神的主题提出了特别的困难，不仅因为当这些人在关于圣子的辩论中疲惫时，他们以更炽烈的热情攻击圣神（因为对他们来说，似乎绝对必须有一个对象来表达他们的不敬之心，否则生命对他们而言似乎不再值得活下去），而且还因为我们自己也被他们众多的问题搞得精疲力竭，处境与那些失去胃口的人有些相似；他们对某种食物产生了厌恶，便回避一切食物；同样，我们也厌恶一切讨论。然而，愿圣神赐予我们恩宠，那样话语将进行，天主将受光荣。那么，我们将把仔细考查和区分\Quote{神}或\Quote{圣}这两个词在圣经中有多少种用法和理解的任务，连同适合这种研究的证据，以及说明除了这些之外，这两个词——我是指\Quote{圣神}——的组合如何以一种独特的意义使用，留给其他为我们也为他们自己在这方面工作过的人，正如我们也为他们工作过一样；但我们自己将致力于主题的其余部分。\par

那些因为我们引进一位陌生的或外加的天主——即圣神——而向我们发怒，且如此激烈地为字句争辩的人，应该知道他们是无故害怕；我要让他们清楚明白，他们对字句的热爱不过是他们不敬之心的掩饰，正如稍后我们在尽最大力量反驳他们的反对意见时将要表明的那样。但我们如此信赖我们所钦崇的圣神的天主性，以致我们要通过将属于天主圣三的名号归于他来开始关于他天主性的教导，即使有些人可能认为我们过于大胆。圣父是那照亮进入世界每个人的真光。圣子是那照亮进入世界每个人的真光。另一位护慰者是那照亮进入世界每个人的真光。曾是、曾是、曾是，但只是一体。光重复了三次；但只有一个光和一位天主。这就是大卫很久以前心中所表示的，当他说：\BibleQuote{在你的光明中，我们必得见光明}。现在我们既看见又简明地宣告天主圣三的道理，从光（圣父）、光（圣子）、在光（圣神）中领悟。拒绝的人，让他拒绝吧；行不义的人，让他行不义吧；我们宣讲我们所领悟的。我们将登上一座高山，呼喊，即使我们在下面不被听见；我们将高举圣神；我们将不害怕；如果我们害怕，那将是害怕沉默，而不是害怕宣讲。\par

如果曾有一个时期圣父不存在，那么也曾有一个时期圣子不存在。如果曾有一个时期圣子不存在，那么也曾有一个时期圣神不存在。如果那一位从起初就有，那么三位也是如此。如果你推倒了一位，我敢断言，你也没有建立另外两位。因为不完美的天主性有什么益处？或者，如果天主性不完美，又怎会是天主性？缺失了完美中某些东西的，如何能是完美的？如果它没有圣神，肯定缺少些什么；因为它若没有圣神，怎会有圣神？因为要么圣洁是与圣神不同的东西，如果是这样，请某人告诉我它被认为是什么；要么是一样的，那么它为何不是从起初就有的，仿佛天主在一段时间内不完美且没有圣神更好似的？如果圣神不是从起初就有的，他就与我自己处于同一等级，即使在我之前一点；因为我们都因时间而与天主性分离。如果他与我处于同一等级，他如何能使我成为神，或使我与天主性结合？\par

五、或者，让我从一个稍早的起点与你们论述祂，因为我们已经讨论了天主圣三。撒杜塞人完全否认圣神的存在，正如他们否认天使和复活一样；我不晓得依据什么理由，他们拒绝了旧约中关于祂的重要见证。在希腊人中，那些更倾向于谈论天主、并最接近我们的人，在我看来，对祂形成了一些概念，尽管他们对祂的名称有分歧，称祂为世界的心灵，或外在心灵，以及诸如此类。但在我们中间有智慧的人中，有些人将祂视为一个行动，有些人视为一个受造物，有些人视为天主；有些人则不确定该如何称呼祂，说是出于对圣经的敬畏，似乎圣经没有把这事说得清楚。因此他们既不崇拜祂，也不侮辱祂，而是采取一种中立的立场——或者不如说，对祂采取一种十分可怜的立场。在那些认为祂是天主的人中，有些只在理智上是正统的，而另一些则敢于在口头上也如此。我听说过一些更聪明的人，他们衡量神性；这些人同意我们有三个概念；但他们将这些概念彼此完全分离，以至于使其中一个在本质和能力上都是无限的，第二个在能力上无限但本质上有限，第三个在两者上都有限；这样以另一种方式仿效那些称之为创造者、合作者、服务者的人，并认为这些名称所附带的秩序和尊严也符合事实的次序。\par

六、但我们不能与那些甚至不信祂存在的人、也不与希腊的饶舌者进行任何讨论（因为我们不愿以罪人的油来使我们的争论丰富）。至于其他人，我们将这样辩论。圣神必定要么被归入自有者的范畴，要么被归入在另一者中被沉思的事物的范畴；熟悉这些事的人称前者为实体，后者为属性。现在，如果祂是一个属性，祂将是天主的一个行动，因为祂还能是什么，或是谁的，因为这确实最避免复合？如果祂是一个行动，祂将被产生，但不会产生，而且一旦被产生就停止存在，因为这是行动的本性。那么，祂如何行事、说话、定义、忧伤、发怒，并且拥有所有明确属于一个主动者而非运动的特性呢？但如果祂是一个实体，而不是实体的一个属性，祂将被视为要么是天主的受造物，要么是天主。因为任何介于两者之间的事物，无论是与两者毫无共同之处，还是两者的复合，连发明山羊鹿的人也无法想象。现在，如果祂是受造物，我们如何相信祂，如何在祂内得以成全？因为相信\Emph{在}某事物和相信\Emph{关于}它并非同一回事。前者属于神性，后者属于——任何事物。但如果祂是天主，那么祂既不是受造物，也不是制造物，也不是同仆，也不是这些低贱的称呼中的任何一个。\par

七、话已交给你们。让投石器发射吧；让三段论编织起来。要么祂完全是非受生者，要么祂是受生者。如果祂是非受生者，就有两位无始者。如果祂是受生者，你们必须进一步划分。祂要么由圣父所生，要么由圣子所生。如果由圣父所生，就有两位圣子，他们是兄弟。你们可以乐意地称他们为孪生子，或一个年长、一个年幼，因为你们如此喜爱肉体上的概念。但如果由圣子所生，那么这样的一个人会说，我们瞥见了一位孙子神，没有什么比这更荒谬的了。然而就我而言，如果我看到这种区分的必要性，我会毫不畏惧那些名称而承认事实。因为，圣子之所以在某种更高的关系中为圣子（因为我们只能以此指出祂出自天主并与父同性同体），并不意味着所有下界和我们的亲属关系的名称都必须转移到神性上。或者，按照同样的论证，你们可能会认为我们的天主是男性，因为祂被称为天主和父亲，而神性是阴性（由于该词的性），圣神是中性，因为祂与生育无关。但如果你们愚蠢到像古老的神话和寓言那样说，天主通过与自己的意志结婚而生下了圣子，我们就会引介给\Person{马西昂}和\Person{瓦伦提努斯}的雌雄同体的神，他们想象了这些新奇的永世。\par

八、但由于我们不接受你们的第一个划分——即断言在受生者和非受生者之间没有中间项——那么，连同你们的宏伟划分一起，你们的兄弟和孙子也立刻消失了，正如一条复杂链条的第一个环节断裂时，它们随之断裂，并从你们的神学体系中消失。因为请告诉我，你们将把那个处于你们划分的两项之间、并由一位比你们更好的神学家——我们的救主祂自己——引入的发出者置于什么位置？或者，也许你们为了你们的第三遗诏而从福音中取出了那个词：圣神，祂由圣父所发出；祂，既然是从那源头发出，就不是受造物；既然不是受生者，就不是圣子；既然处于非受生者和受生者之间，就是天主。这样，祂逃脱了你们三段论的罗网，彰显自己为天主，比你们的划分更强有力。那么，什么是发出？你们告诉我圣父的非受生性是什么，我就向你们解释圣子受生和圣神发出的生理学，而我们将因窥探天主的奥迹而都变得疯狂。我们是谁，竟能做这些事？我们连脚下的东西都看不见，不能数清海沙、雨滴或永恒的日子，更何况进入天主的深渊，为那不可思议、超越一切言语的本性提供说明？\par

IX. 那么，他们说，圣神有什么欠缺，以至于祂不能成为圣子？因为如果祂没有欠缺，祂就会是圣子了。我们断言，祂没有任何欠缺——因为天主没有缺陷。但是，表现的差异，如果可以这样说的话，更确切地说，是它们彼此之间相互关系的差异，引起了它们名称的差异。因为，事实上，圣子不是圣父并非由于某些欠缺（因为圣子身份并非一种欠缺），然而祂不是圣父。按照这一论证，圣父也必然有某种欠缺，因为祂不是圣子。因为圣父不是圣子，而这既不是由于本质的欠缺，也不是由于本质的从属；而是\OriginalTerm{不受生}{Unbegotten}、\OriginalTerm{受生}{Begotten}或\OriginalTerm{发出}{Proceeding}这一事实，赋予了第一位以圣父之名，第二位以圣子之名，第三位，即我们所说的那一位，以圣神之名，以便三位之间的区别得以保存在天主性的同一本质与尊荣中。因为圣子不是圣父，因为圣父是独一的，但祂是圣父所是的；圣神也不是圣子，因为祂出于天主，因为独生子是独一的，但祂是圣子所是的。三位在天主性内是一，而一位在属性上是三；这样，合一就不是\OriginalTerm{撒伯里乌主义}{Sabellian}的，圣三也不赞同现今邪恶的分裂。\par

X. 那么，圣神是天主吗？当然。那么，祂是\OriginalTerm{同体}{Consubstantial}的吗？是，如果祂是天主。我的对手说，请向我承认，从同一源泉产生一位圣子和一位非圣子，而这两者与源泉同体，那么我就承认天主和天主。不，如果你向我承认有另一位天主和另一天主性，我就给你同样的圣三，具有同样的名称和事实。但由于天主是独一的，至高的本性是独一的，我如何向你呈现相似性？或者你将再次在较低的区域，在你周围的事物中寻找它？这是非常可耻的，不仅可耻，而且非常愚蠢，从下面的事物猜测上面的事物，从变幻的本性猜测不变的事物，并且如\Person{依撒意亚}所说，\Quote{在死人中寻找活人}。但我仍会为你尝试，甚至从那源头给你一些论证的帮助。我想我将略过其他点，尽管我可能从动物史中提出许多点，有些广为人知，有些只有少数人知道，那是自然以奇妙的设计在动物繁衍中所创造的。因为不仅相似者生相似者，相异者生相异者，而且相似者也由相异者所生，相异者由相似者所生。如果我们相信传说，还有另一种繁衍方式，当一个动物自我消耗并自我生成。还有一些生物在某种程度上偏离其真实本性，通过自然的壮丽，经历从一种生物到另一种生物的变化和转化。事实上，有时在同一物种中，部分可能是生成的，部分不是；然而全部是同体的；这更接近我们当前的主题。我只提一个我们本性的、人人都知道的事实，然后我将继续下一部分。\par

XI. \Person{亚当}是什么？天主的受造物。那么\Person{厄娃}是什么？受造物的一部分。而\Person{舍特}是什么？两者的生子。那么在你看来，受造物、部分和生子是同一回事吗？当然不是。但这些人的本质不是相同的吗？当然是。那么，这里是一个公认的事实：不同的人可能有相同的本质。我说这些，不是要把创造、区分或任何身体的属性归于天主性（愿你们中没有字句的争辩者再次攻击我），而是为了在它们之中，如同在舞台上，默观那些纯理智的对象。因为不可能精确地描绘出任何达到全部真理的图像。但是，他们说，这一切是什么意思？难道一个不是生于那一位的后裔，另一个是那一位的其他东西吗？\Person{厄娃}和\Person{舍特}不是都来自同一位\Person{亚当}吗？他们两个都是由他生的吗？不；一个是他的部分，另一个是由他生的。然而两者是同一回事；都是人；没有人会否认。那么，你会放弃你反对圣神的争论，即祂要么完全是受生的，否则就不能是同体或天主；并从人的例子中承认我们立场的可能性吗？我认为这对你是有益的，除非你执意要非常好斗，并与已被证明的事实抗争。\par

XII. 但他说，古往今来有谁敬拜过圣神？有谁向祂祈祷过？哪里记载我们应当敬拜祂，或向祂祈祷，而你从何处得来这信理？日后我们讨论未成文之事时，将给出更圆满的理由；目前只消说：我们是在圣神内敬拜，在圣神内祈祷。因为圣经说：\Quote{天主是神，朝拜祂的人，应当以心神以真理去朝拜。}（\BibleRef{若 4:24}）又说：\Quote{我们不知道我们如何祈求才对，但圣神却亲自以无可言喻的叹息，代我们转求。}（\BibleRef{罗 8:26}）又说：\Quote{我要以神魂祈祷，也要以理智祈祷。}（\BibleRef{格前 14:15}）——即是在理智和圣神内。因此，在我看来，朝拜或向圣神祈祷，就是祂亲自向自己献上祈祷或朝拜。哪个敬虔或有学识的人会不赞同这点，因为事实上，由于三者之间尊荣与神性的平等，对一位的朝拜就是对三位的朝拜？所以我不会被\Quote{万物都是藉着圣子造成的}这一论点吓倒，仿佛圣神也是其中之一。因为经上说\Quote{一切受造的}，而非\Quote{一切}。因为圣父不是受造的，任何非受造之物也不是受造的。你若证明祂是受造的，就把祂归于圣子，把祂算在受造物中；但若你无法证明，你从这全称断言中为你的不敬得不到任何好处。因为若祂是受造的，那必是藉着基督；我自己不否认。但若祂不是受造的，祂怎能是\Quote{万有}之一，或是藉着基督的呢？那么，停止你在反对独生者时羞辱圣父（因为给祂呈现一个受造物来剥夺祂更宝贵的圣子，并非真正的尊荣），并在反对圣神时羞辱圣子。因为祂不是同仆的创造者，而是与同等尊荣的一位同受光荣。不可将圣三的任何一位与你自身同列，免得你从圣三中坠落；不可将同一可敬的本性从任何一位切割，因为若你推翻三者中的任何一位，你就推翻了全体。宁可对合一持有不足的看法，也不可冒险陷入彻底的不敬。\par

XIII. 我们的论述现在到了关键点；我悲哀的是，一个早已沉寂并让位于信德的问题如今又被重新挑起；然而我们必须对抗这些饶舌者，不让判决因缺席而成立，因为我们有圣言站在我们这边，为圣神辩护。他们说，若有一位天主、一位天主、一位天主，为何没有三位神，或者为何被光荣的不是多个本原？谁说这些话？那些已达到更彻底不敬的人，还是那些在某种程度上对圣子持温和态度、占了次要部分的人？因为我的论证部分是对两者共同的，部分特别针对后者。我对这些人的回答如下：你们这些崇拜圣子却背弃圣神的人，有什么权利称我们为\OriginalTerm{三神论者}{Tritheists}？你们岂不是\OriginalTerm{二神论者}{Ditheists}？因为若你们也否认对独生者的崇拜，你们显然已站在我们对手一边。我们何必对你们客气，仿佛你们不是\Emph{完全}死了？但若你们崇拜祂，并因此走在救恩之路上，我们要问你们，你们有什么理由来解释你们的二神论，若有人指控你们的二神论？若你们内有智慧之言，请回答，也为我们打开一条回答之路。因为用于反驳二神论指控的同一理由，也足以反驳三神论的指控。这样，我们将藉着利用你们这些控告者作为我们的辩护者而获胜，再没有比这更慷慨的了。\par

XIV. 我们与两者的争吵和争论是什么？对我们来说，只有一位天主，因为天主性是唯一的，凡从祂而出的一切都归于一位，尽管我们相信三位。因为一位并不比另一位更是天主，也没有一位在前、另一位在后；他们不分意志，能力也不分离；你在它们中找不到任何可分事物的性质；简言之，天主性在分开的位格中是不可分的；如同三个太阳彼此结合，有一道光混合在一起。因此，当我们注视天主性、第一因或\OriginalTerm{独一统治}{Monarchia}时，我们所构思的是唯一；但当我们注视那些天主性居于其中的位格，以及那些无时间地、同等荣耀地从第一因获得存有者时，我们崇拜的就是三位。\par

XV. 那又怎样，他们或许会说，难道希腊人不也相信一位天主性，正如他们更先进的哲学家所宣称的？在我们看来，人性是唯一的，即整个人类；但他们却有许多神，而不是一位，正如有许多人一样。但在这种情况下，共同的本性具有一种仅在思想上可想象的一致性；而个体彼此之间非常遥远，无论是时间、性情还是能力上。因为我们不仅是复合的存在，而且也是矛盾的存在，既与他人也与自己矛盾；我们甚至一天之内也不完全保持同一，更不用说一生了，而是身体和灵魂都处于不断的流动和变化中。天使以及那仅次于圣三的整个更高本性，或许也可作同样说明；尽管他们在某种程度上是单纯的，并因接近至善而更稳固地持守于善中。\par

XVI. 那些希腊人所敬拜的鬼神（他们自己称之为\OriginalTerm{鬼神}{dæmons}），无需我们作任何控告；他们的神学家自己的证言已定了他们的罪，有些受情欲支配，有些好结党，充满无数的邪恶与变化，彼此敌对，不仅互相敌对，甚至敌对他们的第一因——他们称之为\OriginalTerm{Oceani}{Oceani}、\OriginalTerm{Tethyes}{Tethyes}、\OriginalTerm{Phanetes}{Phanetes}以及其他许多名字；最后，有一位神因贪权而憎恨自己的子女，又因贪婪吞下所有其余，为要成为所有人神之父——他悲惨地吞下他们，然后又吐出来。如果这些不过是神话和寓言，像他们为了逃避故事的可耻而说的那样，那么对于\Quote{万物分为三份，每位神治理宇宙的不同部分，各有其疆界和品级}这一论断，他们又作何解释？但我们的信仰并非如此，我的神学家说，雅各伯的产业也不是这样。然而，这些位格中的每一位，与其所联合者之间的合一，不亚于其自身的合一，这是由于本质与能力的同一。这就是我们所理解的关于合一的道理。如果这道理是真的，就让我们感谢天主赐予我们这一瞥；如果不是，就让我们寻求更好的。\par

XVII. 至于你用来推翻我们所主张的合一的那些论据，我不知道该说你是开玩笑还是认真的。因为这论据是什么？你说，同一本质的事物会\OriginalTerm{共同计数}{connumeration}，而所谓共同计数，你指的是它们被归入一个数字。但不同本质的事物则不会这样计数……因此，按照这论据，你无法避免说到三个神，而我们则完全没有这种危险，因为我们断言它们不是\OriginalTerm{同质}{consubstantial}的。所以，你用一句话就摆脱了麻烦，并赢得了一场有害的胜利；实际上，你做的事就好像那些因怕死而自缢的人一样。为了在捍卫\OriginalTerm{单一统治}{Monarchia}时省事，你否定了\OriginalTerm{神性}{Godhead}，把问题留给了对手。但对我来说，即使需要劳苦，我也不会放弃我朝拜的对象。然而在这一点上，我看不出困难在哪里。\par

XVIII. 你说：同一本质的事物是\OriginalTerm{共同计数}{connumeration}的，但非同质的则一个一个地计数。你从哪里得来这个？从什么教义或神话的老师那里？难道你不知道，每一个数字表达的是其所包含事物的数量，而非事物的本性吗？但我如此守旧，或许该说如此无知，以至于即使事物本性不同，我也用\Quote{三}来指那数量，并用\Quote{一、一、一}以不同的方式指那些单位，即使它们在本质上合而为一——我关注的不是事物本身，而是所枚举的事物在数量上的多少。既然你如此紧咬字面（尽管你是在反驳字面），那就请从这源头拿你的论证吧。在\BibleBook{}{箴言}中有三样行走得好：狮子、公山羊和公鸡；再加上第四样：在人民面前演说的君王——这是略过那里列举的其他四样，尽管它们本性各异。我在梅瑟那里也发现两个\OriginalTerm{革鲁宾}{Cherubim}被单独计数。但是，按照你的术语，前者虽本性大不相同，难道不能称为三吗？后者虽如此紧密相连、同一本性，难道不能作为单位处理吗？因为如果我谈论天主和\OriginalTerm{玛门}{Mammon}，作为两个主人，归在同一名下——它们彼此如此不同——我大概会因这种共同计数而更被人嘲笑。\par

XIX. 但在我看来，他说，那些被称为共同计数且同一本质的事物，其名称也对应，例如\Quote{三个人}或\Quote{三个神}，而不是\Quote{三个这个和那个}。这让步意味着什么？这是为名称立法的人应当做的，而非为真理辩护的人。因为我也要断言，伯多禄、雅各伯和若望不是三个或同质，只要我不能说\Quote{三个伯多禄}、\Quote{三个雅各伯}或\Quote{三个若望}；因为你们为普通名称所保留的，我们也要求为专有名称保留，按照你们的安排；否则你们若不将你们自己所假设的也给予别人，就是不公。那么，若望在他的公函中说\BibleQuote{作证的有三个：圣神、水及血}，你们认为他在胡说吗？首先，因为他竟敢将不同本质的事物归在同一个数字下——而你们说这只应在同质的事物中做，谁能断言这些东西是同质的呢？其次，因为他所用术语的方式不一致：他先用阳性\Quote{三}，然后加上三个中性词，违背了你们和你们的文法家所下的定义和规律。因为先写阳性\Quote{三}，然后加上中性的\Quote{一、一、一}，与先写阳性的\Quote{一、一、一}然后不用阳性而用中性的\Quote{三}——你们自己在神性上就否认后者——这有什么区别？至于\OriginalTerm{蟹}{Crab}，它可以指动物、器械或星座，你们又怎么说？还有\OriginalTerm{狗}{Dog}，它时而指陆地、时而指水中、时而指天上的，你们又怎么说？难道你们不说\Quote{三只蟹}或\Quote{三条狗}吗？当然是的。那么它们因此是同一实体吗？只有傻子才会这么说。所以你们看，从共同计数而来的论证完全站不住脚，被所有这些实例反驳了。因为如果同一实体的事物不总是归在同一个数字下，而非同一实体的事物却这样计数，并且在两种情况下都用一次称呼，那你们的教义又得到了什么好处？\par

二十. 我还要看看这进一步的一点，它并非与主题无关。一加一等于二；二分又成为一和一，这是非常明显的。然而，如果如你们理论所要求的，相加的元素必须是\OriginalTerm{同体}{consubstantial}，而分离的则是\OriginalTerm{异质}{heterogeneous}，那么同一事物必须既是同体又是异质。不！我嘲笑你们引以为豪的\OriginalTerm{先后的计数}{Counting Before and Counting After}，好像事实本身依赖于名称的顺序。如果是这样，根据同一法则，由于同一事物因其本性的平等，在圣经中被计数时有时在先有时在后，什么阻止它们同时比自己更尊贵和不尊贵呢？我同样说\Quote{天主}和\Quote{主}这些名称，以及介词\OriginalTerm{从祂、藉祂、在祂内}{Of Whom, and By Whom, and In Whom}，你们以此按照艺术规则为我们描述神性，将第一个归于圣父，第二个归于圣子，第三个归于圣神。因为如果每个这种表达都恒常地分配予每一位，你们会怎么做呢？事实上它们被用于所有位格，正如研究过此问题的人所明了的，你们却使它们成为这种本性及尊荣不平等的基础。这对所有并非完全缺乏理智的人已经足够。但是，既然你们一旦攻击圣神之后，就难以控制自己的冲动，反而像一头狂暴的野猪，将争论进行到底，并扑向刀口，直到在自己胸前接受全部创伤；那么让我们继续看看还有什么论据留给你们。\par

二十一. 反复多次，你们用\OriginalTerm{圣经的沉默}{silence of Scripture}来攻击我们。但这并非陌生教义，亦非事后之想，而是被古人及我们同时代的许多人承认并明确陈述，这已由许多论述此主题的人所证明，他们处理圣经，并非漠不关心或当作消遣，而是深入到\OriginalTerm{字面与内在意义}{the letter and the inner meaning}，并被堪当看到隐藏之美，且被知识之光照耀。然而我们反过来将尽可能简要地证明它，以免显得过分好奇或不恰当地好胜，在别人的根基上建造。但由于圣经并不十分清楚或经常以明确的词句写祂为天主（如同它首先写圣父，然后写圣子），这事实成为你们亵渎及这种过度多言和不敬的机缘，我们将通过简短讨论事物和名称，特别是它们在圣经中的用法，使你们从这一不便中解脱。\par

二十二. 有些事物不存在，却被谈论；另一些存在，却不被谈论；有些既不存在也不被谈论；有些既存在又被谈论。你们问我证明吗？我准备给出。根据圣经，天主睡眠乃醒，发怒，行走，以\OriginalTerm{革鲁宾}{Cherubim}为祂的宝座。然而祂何时变得易于受情感影响？你们可曾听说天主有身体？那么，这尽管不是事实，却是一种修辞表达。因为我们按照自己的理解，从我们自己的属性中将名称给予了天主的属性。祂远离我们而沉默，好像不关心我们，原因祂自己知道，我们称之为祂睡眠；因为我们自己的睡眠便是这样一种不活动的状态。同样，祂突然转向为我们行善是苏醒；因为苏醒是睡眠的解散，正如眷顾是转离的解散。祂惩罚时，我们说祂发怒；因为对我们而言，惩罚是愤怒的结果。祂时此时彼的行动，我们称之为行走；因为行走是从一处到另一处的变化。祂在神圣天军中安息，好像喜爱居住其中，是祂的坐于宝座；这同样来自我们自己，因为天主没有像在圣人们身上那样安息。祂行动的迅捷被称为飞翔，祂警醒的眷顾被称为祂的面容，祂的给予和赏赐是祂的手；总之，天主的所有其他能力或活动都为我们描画了某种相应的肉身能力。\par

二十三. 此外，你们从哪里得到你们的\OriginalTerm{无始}{Unbegotten}和\OriginalTerm{无起源}{Unoriginate}，你们立场的这两座堡垒，或者我们从何处得到我们的\OriginalTerm{不朽}{Immortal}？请用同样多的词句向我展示这些，否则我们就要把它们搁置一边，或当作圣经未包含而擦除；你们就被你们自己的原则所杀，你们所依赖的名称被推翻，连同你们所信赖的避难之墙。难道不明显它们源自暗示它们的经文，尽管词句并未实际出现？这些经文是什么？——\BibleQuote{我是元始，我是终末，在我以前没有别的神，以后也不会有}\BibleRef{依 44:6}。因为所有依赖于那\InnerQuote{自有}的都支持我，因为它既无开始也无结束。当你们接受这一点，即祂之前没有事物，祂也没有更古老的原因，你们就隐含地给了祂\InnerQuote{无始}和\InnerQuote{无起源}的称号。而说祂的存在没有终结，就是称祂为不朽和不可毁灭。那么，我提到的第一对就这样解释了。但哪些事物既不存在也不被说？天主是恶的；球是方的；过去是现在；人不是复合的存在。你们可曾见过如此愚蠢的人敢于思考或断言任何这样的事情？还剩下要表明那些既存在也被言说的事物：天主、人、天使、审判、\OriginalTerm{虚妄}{Vanity}（即像你们这样的论证），以及\OriginalTerm{信仰的颠覆}{subversion of faith}和\OriginalTerm{奥秘的掏空}{emptying of the mystery}。\par

XXIV. 既然在术语和实体上有如此大的差异，为何你如此拘泥于字句，偏袒犹太人的智慧，为音节而牺牲事实？然而，如果你说了两次五或两次七，我从你的话中得出结论你指的是十或十四；或者如果你说了一个理性的、必死的动物，你指的是人，你会认为我在胡说吗？当然不会，因为我只是重复了你的意思；因为词语并不更多地属于说话者，而是属于唤起它们的人。那么，在这种情况下，我关注的与其说是所用的术语，不如说是它们所要传达的思想；同样，如果我在圣经言语中发现某种意思中包含了未表达或未清楚表达的东西，我也不应因畏惧你对术语的诡辩技巧而避免表达它。这样，我们将坚持我们自己的立场，对抗半正统者——我可不把你算在其中。因为既然你否认圣子的许多清晰称号，显然即使你学到更多更清晰的称号，你也不会被感动而敬畏。但现在我要稍微回溯一下论证，向你展示——尽管你很聪明——这整个神秘体系的原因。\par

XXV. 在世界存续的整个时期中，人类生活有过两次显著的变化，也被称为两个盟约，或由于此事广为人知而称为两次地震；一次是从偶像到法律，另一次是从法律到福音。我们在福音书中得知第三次地震，即从这大地到那不可动摇或移动的国度。这两个盟约在这一点上是相似的：变化并非突然发生，也不是在努力之初就发生。为什么（因为这是我们必须得知的一点）？以免我们遭受暴力，而是通过说服被感动。因为非自愿的东西不能持久，如同被强力遏制的溪流或树木。而自愿的则更持久、更安全。前者出于使用暴力者，后者属于我们自己；前者出于天主的温和，后者出于暴虐的权威。因此，天主认为他不应恩待不情愿者，而应施恩于愿意者。所以，如同一位导师或医生，他部分移除、部分容忍祖先的习惯，允许一些倾向于快乐的东西，正如医生对待病人一样，巧妙地将药物与美味混合，以便药物被服用。因为改变那些被习俗和习惯尊崇的东西并非易事。例如，第一次砍掉了偶像，但保留了祭献；第二次摧毁了祭献，但没有禁止割损。然后，一旦人们接受了削减，他们也交出了先前被允许的东西：第一次是祭献，第二次是割损；他们从外邦人变成了犹太人，从犹太人变成了基督徒，被逐步的变化引入福音。\Person{保禄}是一个证明；因为他曾一度施行割损，并服从法律上的净化，最后进步到可以说：\BibleQuote{弟兄们，如果我仍宣讲割损，为什么还受迫害呢？}他先前的行为属于暂时的安排，后来的属于成熟。\par

XXVI. 我可以将神学与此比较，只是它以相反的方式进展。因为在我用来说明的情况下，变化是通过逐次削减而实现的；而在这里，完美是通过添加而达到的。情形如下：旧约公开宣告了圣父，而圣子则较为隐晦。新约彰显了圣子，并暗示了圣神的天主性。如今圣神亲自居住在我们中间，并赐予我们关于他自己的更清晰的证明。因为当圣父的天主性尚未被承认时，公开宣告圣子是不安全的；而当圣子的天主性尚未被接受时，再加重（请允许我如此大胆地表达）加上圣神也不安全；否则，人们可能像负重过度的人，眼睛尚未能承受太阳的光芒，反而会连他们力所能及的东西也失去；而通过逐步添加，以及如\Person{达味}所说的\BibleQuote{上行}，进步，从光荣到光荣，使天主圣三的光辉照耀那些更蒙光照的人。因此，我认为，圣神是\Emph{逐步}来居住于门徒中，根据他们接受他的能力而分配自己：在福音开始时、在受难后、在升天后，使他们的能力得以成全，被吹气在他们身上，并以火舌显现。而且，他也逐渐被耶稣所宣告，如果你更仔细地阅读，你就会自己发现。他说：\BibleQuote{我要请求父，他必会赐给你们另一位护慰者，就是真理之神。}他说这话是为了不显得是一位竞争的天主，或借另一权威向他们讲论。又说：\BibleQuote{他必会派遣他，但是在我的名内。}他略去了\BibleQuote{我要请求}，但保留了\BibleQuote{必会派遣}；然后又说：\BibleQuote{我要派遣}——他自己的尊威。然后：\BibleQuote{他要来}——圣神的权威。\par

二十七. 你看见光照逐渐照临我们；神学的次序，我们最好遵守，既不太过突然地宣布，也不一直隐藏到底。前者不合科学，后者近乎无神；前者会惊动外人，后者会使我们自己的人疏离。我要在我所说的之上再加一点；这点或许很容易进入某些人的思想，但我认为是我个人思索的果实。我们的救主有一些事，祂说当时门徒们还不能承受（尽管他们已充满许多教诲），或许就是出于我提到的理由；因此它们被隐藏了。祂又说，当圣神来住在我们中间时，祂会教导我们一切事。在这些事中，我认为其中之一就是圣神本身的天主性，后来当这样的知识变得合时宜、并在我们救主恢复之后能够被接受时，才被显明，这时它不再因其奇妙性而被怀疑不信。因为无论是祂所应许的，还是圣神所教导的，有什么比这更伟大呢？如果真有被视为伟大且堪当天主尊威的事被应许或教导。\par

二十八. 这就是我对此事的立场，我希望这永远是我的立场，以及所有亲爱我者的立场：朝拜天主圣父、天主圣子、天主圣神，三个位格，一个天主性，在尊荣、光荣、本体和国度上不可分割，正如我们当中一位不久前离世的神感哲学家所表明的。愿那持不同心思、或随从时代风尚、时而这样想时而那样想、在最高事物上思想不健全的人，不得看见晨星的升起，如经上所说，也不得见其光耀的荣光。因为如果祂不该受朝拜，祂怎能藉着圣洗使我成为神呢？但如果该受朝拜，那么祂必定是钦崇的对象，而若是钦崇的对象，祂必定是天主；二者相连，一条真正金链和救恩之链。实际上，来自圣神的是我们的新生，来自新生的是我们的新创造，来自新创造的是我们对祂的尊严的更深刻认识，我们正是从祂那里领受了这一切。\par

二十九. 那么，这就是一个承认圣经沉默的人可以说的。但现在，一大群见证将要涌向你，从中圣神的天主性将向所有不是极其愚蠢、或全然敌视圣神的人显明，在圣经中最清楚地被认出。请看这些事实：基督诞生，圣神是祂的前驱。祂受洗，圣神作证。祂受试探，圣神带领祂上去。祂行奇迹，圣神伴随。祂升天，圣神取代祂的位置。在天主的概念中，有什么大事是祂没有能力做的呢？有什么属于天主的称号没有应用于祂，除了无生者和受生者？因为必须保持圣父和圣子各自的特性，免得在将一切——甚至混乱本身——安排得井然有序的天主性中产生混乱。实际上，当我想到那丰富的称号，以及那些攻击圣神的人亵渎了多少名号时，我不禁战栗。祂被称为天主之神、基督之神、基督的理智、主之神，而祂本身是主，义子之神、真理之神、自由之神；智慧之神、聪敏之神、超见之神、刚毅之神、明达之神、孝爱之神、敬畏之神。因为祂是这一切的创造者，以自己的本体充满一切，包含万有，以本体充满世界，却不能被世界以其能力所理解；良善、正直、尊贵，出于本性而非出于义子；圣化而非被圣化；衡量而非被衡量；分施而非分受；充满而非被充满；包含而非被包含；被继承、被光荣、与圣父和圣子同列；被作为威胁提出；天主的手指；像天主的火；据我理解，这是显示祂的\OriginalTerm{同体性}{consubstantiality}；创造者圣神，祂藉着圣洗和复活创造新生命；那知晓万有、教导、随己意吹拂、引导、说话、派遣、分离、发怒或受试探的圣神；那启示、光照、赋予生命，或者更确切地说，祂本身就是光明和生命；那建造圣殿；那使人神化；那得以成全甚至预先行洗，却在洗礼之后作为一个分别的恩赐被寻求；那行天主所做的一切事；分成火舌；分施恩赐；使成为宗徒、先知、传福音者、牧者和教师；理解力多样、明澈、锐利、无玷、无阻碍，这与祂的行动极其智慧和变化多端是相同的；使一切事物清晰明了；具有独立的能力、不变、全能、全见、穿透一切有理智的、纯洁的、极其精微的神体（我想是天使的军旅）；同样地，也穿透一切先知和宗徒的神体，以同样的方式而非同样的地方；因为他们生活在不同的地方；从而显示祂是不受限制的。\par

三十. 那些讲论并教导这些事的人，此外还称祂为另一位\OriginalTerm{辩护者}{Paraclete}，意思是一位另外的天主，他们知道唯独亵渎祂的罪不得赦免，并用如此可怕的恶名标记了\Person{阿纳尼亚}和\Person{撒斐辣}，因为他们向圣神说谎，你们认为这些人怎样？他们宣称圣神是天主，还是别的什么？实际上，如果你对此有任何怀疑，需要别人教导你，那么你必定是极其迟钝，远离圣神。所以，祂的名号如此重要，如此鲜明。为何还需要向你们摆出那些言词本身所包含的见证呢？此外，凡在此处用更卑微的方式所说的，例如祂被赐予、被派遣、被分开；祂是恩赐、赏报、默感、应许、为我们代祷，以及不再细说的其他任何这类表达，都应当归因于第一因，以表明祂是从谁而来的，并使人们不致以异教方式承认三个本原。因为混淆位格正如\Person{撒伯里乌}派，或分离本性正如\Person{亚略}派，同样是亵渎的。\par

三十一、我在自己的心中非常仔细地思考了这个问题，从每一个角度审视它，以期为此极其重要的主题找到某种比喻，但我未能在地上发现任何可与\OriginalTerm{天主性}{Godhead}相比拟的事物。因为即使我偶然找到一点微小的相似之处，它也多半从我手中溜走，只留下我连同我的例子一起落在地上。我像前人一样，在心中构想出一眼泉、一道水流、一条江河，想看看是否第一个可类比于圣父，第二个可类比于圣子，第三个可类比于圣神。因为在它们之间没有时间上的区别，也没有被切断与彼此的联系，尽管它们似乎被三个\OriginalTerm{位格}{personalities}分开。但我首先害怕会呈现出天主性中一种不停滞的流动；其次，通过这个形象，会引入数字上的单一。因为泉、水流和江河在数字上是同一个，尽管形式不同。\par

三十二、我又想到了太阳、光线和光。但这里再次有一种担忧，惟恐人们会在那\OriginalTerm{单纯的本性}{Uncompounded Nature}中引入组成的概念，正如在太阳和太阳中的事物中那样。其次，恐怕我们会将\OriginalTerm{本体}{Essence}赋予圣父，却否认其他位的位格，使它们只成为天主的\OriginalTerm{能力}{Powers}，存在于祂内而非位格性的。因为光线和光并非另一个太阳，它们只是来自太阳的光辉，是祂本体的属性。这样，随着比喻，我们就会将\OriginalTerm{存在和非存在}{Being and Not-being}都归于天主，这更加荒谬。我也听说有人提出过如下比喻：一束阳光照射在墙上，随着光束在半空中所吸收的湿气的运动而颤动，然后被坚硬的物体阻挡，从而引发一种奇特的振动。因为它以许多快速运动颤动着，它既不是一个多于多个，也不是多个多于一个；因为由于它合一和分离的迅捷，它未等眼睛看见就溜走了。\par

三十三、但就连这个比喻，\Emph{我}也无法使用；因为很明显是什么给了光线以运动；但在天主之前却没有任何事物能推动祂；因为祂自己就是万物的原因，祂没有先在的原因。其次，在这里也暗示了诸如组合、扩散和不稳定易变的本性之类的东西……这些我们都不能设想存在于天主性中。总之，在这些比喻中，没有任何东西能为我的思想提供一个立足点，使我能够观察我试图为自己呈现的对象，除非有人可以宽厚地接受形象中的某一点而拒绝其余。最后，在我看来，最好让这些形象和阴影离去，因为它们具有欺骗性，远非真理；而我自己则依附于更虔诚的概念，安于少量的言辞，藉着圣神的引导，将我从祂所领受的光照作为我的真诚伴侣和同伴保持到底，并度过此世，尽我所能说服所有人一同敬拜圣父、圣子、圣神，同一个天主性和\OriginalTerm{权能}{Power}。愿光荣、尊威和权能归于祂，直到永远。阿们。\par

\SourceRecord{Translated by Charles Gordon Browne and James Edward Swallow.；Nicene and Post-Nicene Fathers, Second Series；Vol. 7.；Edited by Philip Schaff and Henry Wace.；Buffalo, NY: Christian Literature Publishing Co.,；1894.；<http://www.newadvent.org/fathers/310231.htm>.}

% structure: p0001 source=h1 target=section reason=page-title

\section{第三十三篇讲道}

\Emph{反亚略派，并论及自身。}\par

\EditorialNote{于公元380年中叶在君士坦丁堡宣讲。}\par

责难我们贫穷而自夸富裕、以人数多少来定义教会、轻视小小羊群、以天平衡量天主性、称量人民、看重沙粒而轻忽天上星体、珍藏卵石而忽视珍珠的人在哪里？他们不知道沙粒并不比星体更多，卵石不若璀璨宝石——前者比后者更纯净、更珍贵。你们又愤怒了吗？又武装自己了吗？又侮辱我们了吗？这是新信仰吗？且稍止你们的威吓，容我发言。我们不会侮辱你们，但要驳斥你们；不会威胁，但要指责你们；不会打击，但要治愈你们。这也算是侮辱！何等骄傲！你们在此也视同等者为奴隶吗？若非如此，容我直言；因为即使是兄弟，若被欺骗，也会责备兄弟。\par

你们愿意我对你们说出天主对那顽梗悖逆的以色列所说的话吗？\BibleQuote{我的百姓啊，我向你们做了什么？或在何事上伤害了你们，或在何事上使你们厌烦？}\BibleRef{米 6:3}这话从我口中对辱骂我的你们说出，更为合适。可悲的是，我们彼此伺机攻击，因意见分歧而摧毁了共融的精神，变得几乎比现在与我们交战的蛮族更不近人情、更野蛮——这些蛮族被我们所分离的\OriginalTerm{天主圣三}{Trinity}联合起来，一同攻击我们；不同之处在于，我们并非外国人对外国人的侵扰劫掠，也非不同语言之民族间的战争（在灾祸中这算是些许安慰），而是彼此为敌，几乎是同室操戈；或者，若你们愿意，我们是同一身体上的肢体，互相吞噬、彼此消耗。灾难还不止于此，尽管这已够糟，因为我们甚至视自己的削弱为收获。既然我们处于这种状况，并按时代来调节信仰，让我们彼此比较时代：你们以你们的皇帝，我以我的主权者；你们以阿哈布，我以约史雅。请告诉我你们的节制，我必宣扬我的暴力。但你们的暴行已被许多书籍和口舌宣扬，我想后世将视之为你们行为的永世耻辱柱；我则要宣告我的所为。\par

我曾率领何种骚动的暴民攻击你们？曾武装何种士兵？曾怂恿何种暴怒的将军——比其雇主更野蛮，甚至连基督徒也不是，反将他对我们的不敬作为私人崇拜献给他自己的神明？我曾围攻那些正在祈祷、向天主举手的人吗？我曾用号角打断圣咏吗？或将圣体之血与屠杀之血混在一起？我曾用死亡的嚎叫止息心灵的叹息，或用悲剧的眼泪止息忏悔的泪水？我曾将祈祷之所变为葬地？谁曾将不可触碰的礼仪器皿交给邪恶之人的手，交给\Person{乃步匝辣当}（\BibleRef{列下 25:11}）——厨役长，或交给\Person{贝耳沙匝}——他邪恶地用圣器宴饮（\BibleRef{达 5:3}），后来为他的疯狂付出应有代价？\BibleQuote{可爱的祭台}，如圣经所说，但\BibleQuote{如今已被亵渎。}哪个放荡的青年曾因我们的缘故，以可耻的扭捏作态侮辱你们？尊贵的宝座，配得尊贵之人的居所和安息，从古至今一直由敬虔的司铎继承，他们教导神圣奥理；哪个异教的通俗演说家或邪恶的舌头曾登上你，攻击基督徒的信仰？贞女的谦逊与庄严，连德行高尚之人的目光都难以承受，我们中谁曾羞辱你们，因暴露不可见之物而冒犯你们，并将可怜的景象（配得上索多玛之火）展示给不敬之人的眼睛？我绝口不提那些死亡，它们比这羞辱更可忍受。\par

我们曾放出何种野兽来攻击圣人的身体——如同有些人出卖人性一般——只凭一个罪名：即不认同他们的不敬；或者通过与他们共融而玷污自己？我们避之如蛇毒，并非因为它伤害身体，而是因为它玷污灵魂的深处。我们曾指控谁埋葬死人（连走兽也敬重死者）为犯罪行为？何等罪状，配得上另一个剧场和另一群野兽！哪个主教的苍老肌肤被钩子刮过——在门徒面前，他们无力相助，只能以泪洗面，与基督一同被悬挂，在苦难中得胜，以宝贵的血洒向百姓，最后被带走处死，与基督一同被钉、埋葬并受光荣——与那位以如此牺牲者征服世界的基督一同受光荣？哪些司铎被那些敌对的元素——火与水——分开，在海上升起奇异的烽火，与所乘坐的船只一同被焚？谁（为掩盖我们苦难的多数部分而保持沉默）被施予如此恩惠的官员指控为残忍？因为即使他们顺从那些人的私欲，至少他们憎恶其意图的残酷。前者是机会主义，后者是算计；前者来自皇帝的无法律，后者来自他们必须据以审判的法律意识。\par

五、再谈一些更久远的事，因为那些也属于同一团体。我曾砍下过谁的——无论是生者还是死者——手，为了诬蔑圣人，并靠夸夸其谈而凯旋于信德？我曾把谁的流放当成恩惠来计数，竟不尊敬神圣的哲学家们的圣学院，甚至向那些学院求助过他们恳求的人？绝非如此；正相反，我把那些为真理而招致愤怒的人都算作了殉道者。你们指责我言语放荡，但我在谁身上引入了那些几乎无血肉的妓女？我把哪个信徒逐出本国，交给无法无天的人之手，使他们像野兽一样被关在没有光线的房间里，并且（这是悲剧中最可悲的部分）彼此分离，忍受饥渴的困苦，食物经量配给，他们必须通过狭窄的开口接收，甚至不允许他们见到难友？而这些受苦的人是谁？是世界不配有的人。\BibleRef{希 11:38} 你们就是这样尊崇信德吗？这就是你们的善待吗？你们不知道这些事的大部分，这也很合理，因为这些事实众多，且行为本身令人愉快。但受苦的人记得更清楚。甚至有一些人比时代本身更残酷，像野猪冲向篱笆。我要求你们交出昨天的受害者——那位老人，亚巴郎般的神父，当他从流放归来时，你们在光天化日之下和城市中央用石头迎接他。而我们，若如此说不算冒犯，甚至连我们的凶手也免于危险。天主在某处圣经中说：\BibleQuote{我怎能为此赦免你呢？}\BibleRef{耶 5:7} 这些事中哪一件我该称赞？或者，我该为哪一件给你戴上花冠？\par

六、既然你们的前科如此，我也希望你们来告诉我我的罪行，以便我或能改过自新，或蒙受羞辱。我最希望的是能发现自己全然无罪；但如果不能，至少能从罪行中悔改；因为这是审慎者的次优选择。因为即使我不像义人那样首先自我控告，\BibleRef{箴 18:17}但我仍然乐意接受他人的医治。你们对我说：\Quote{你们的城市很小，或者根本算不上城市，只是一个村庄，干旱、不美、居民稀少。}但是，我的好朋友，这是我的不幸，而非我的过错——如果这确实是不幸的话；并且如果这是违反我意愿的，那么我为我的厄运而可怜（如果我可以这么说的话）；但如果是自愿的，那么我就是一个哲学家。其中哪一样是罪行？谁会责备海豚不是陆地动物，或责备牛不是水生动物，或责备七鳃鳗是两栖动物？但是你们继续说，我们有城墙、剧院、赛马场、宫殿、美丽宏伟的柱廊，以及那奇妙的地下和地上河流，还有那宏伟而令人钦佩的圆柱，熙攘的市场和不安分的人民，以及由高尚之人组成的著名元老院。\par

七、你们为什么不再提地点的便利，以及我所称的陆地与海洋之间的竞争，看谁拥有该城，谁用一切美物装饰我们的皇家之城？那么这就是我们的罪行：你们伟大而辉煌，而我们微小且来自小地方？许多人这样冤枉你们，实际上所有被你们超越的人都这样；我们必须死，因为我们没有建造一座城，没有在它周围筑墙，不能夸耀我们的赛马场、竞技场和猎犬群，以及与此相关的所有蠢事；也不能夸耀我们浴池的美丽和辉煌，以及它们的大理石、绘画和各种金绣的昂贵，几乎胜过大自然？我们也还没有为自己围海造地，或混合季节，就像你们，新的造物主，所做的，以便我们可以生活在最愉快又最安全的方式中。如果你们愿意，可以添加其他指控，你们说：\BibleQuote{银子是我的，金子是我的}，\BibleRef{盖 2:8} 这些是天主的话。我们既不看重财富——如果财富增加，我们的法律禁止我们用心思——也不计算年收入和日收入；也不彼此争胜，用诱惑物堆满我们的桌子，以满足我们无知的肚腹。因为我们也不高度珍视那些被我们吞下后都变成同等价值——或者更该说无价值——并被排出的东西。但我们生活如此简朴，仅够糊口，与那些饮食不用器具、不假人工的野兽几乎没有什么区别。\par

八、你们是否还批评我衣服的破烂，以及我面容的缺乏修饰？因为我看到有些非常微不足道的人以此为傲。你们能不能不要管我的头，不要嘲笑它，就像孩子们嘲笑\Person{厄里叟}一样？接下来发生的事我就不提了。你们是否还要指责我的缺乏教育，以及你们所认为的我的言辞粗鲁和乡土气息？你们又该把我没有满嘴闲聊、不是受欢迎的小丑、不热衷集市、不与随便什么人闲聊说各种话题、甚至使谈话变得沉重、不常去\Place{策宇西普斯}（那个新耶路撒冷）、也不挨家挨户奉承和填饱肚子、而是大部分时间待在家里、精神低落、面容忧郁、安静地与自己相处（我是自己行为的真正审判官）、也许因无用而该被监禁——这些事放在哪里？你们怎么能宽恕我这一切，而不指责我？你们是多么甜蜜和善良！\par

IX. 但我如此守旧且富有哲思，竟相信诸天为众人所共有；太阳与月亮的运行、星辰的秩序与排列也是如此；众人均等地分享白昼与黑夜、季节的更替、雨水、果实以及大气的孕育之力；奔流不息的江河是众人共同的丰沛财富；大地是同一的，她是母亲也是坟墓，我们从中被取出，也将归于她，无人拥有比他人更多的份额。此外，我们共享理性、法律、先知，以及基督的苦难——藉此，我们所有毫无例外地分享同一亚当、被蛇引入歧途、被罪所杀、由天上的亚当拯救、藉羞耻之树被带回我们曾堕落的生命之树的人，都被重新创造。\par

X. 我也曾被撒慕尔的辣玛所欺骗——那是一位伟人的小小父家；这并非对先知的羞辱，因为它的荣耀与其说出自自身，不如说出自他；他并未因此受阻于在出生前就被献给天主，或发出神谕并预见未来；不仅如此，他还傅油君王和司祭，并审判著名城邑的人民。我也听说过撒乌耳：他在寻找父亲的驴子时找到了王位。连达味自己，也从羊圈中被选为以色列的牧者。亚毛斯又如何？他岂不是一边牧羊、刮取桑葚，一边被托付了预言的恩赐？我怎能略过若瑟？他既是奴隶，又是给埃及供应粮食的人，还是那预许给亚巴郎的万千子孙之父。是的，我还被厄里亚的加尔默耳山所欺骗——他接受了火战车；被厄里叟的羊皮所欺骗——它比丝绸织物或金线织成的衣服更有力量。我被若翰的旷野所欺骗——那里拥有妇人所生者中最大的那位，以及我们所知的那衣着、那食物、那腰带。我甚至敢于超越这些，发现天主亲自眷顾我的质朴。我将与白冷为伍，分享马槽的羞辱；因为你们以此为理由拒绝尊重天主，那么你们以此为理由轻视祂的传报者也不足为奇。我要向你们提出渔夫们，以及那些贫穷的、福音所传报的对象，他们被选在众多富人之前。你们何时才会停止以你们的城市自夸？你们何时才会尊重那你们所憎恶轻视的旷野？我尚未说金子产于沙中；也未说透明宝石是岩石的产物和赠礼；因为若以这些来反驳城市中所有可耻之事，我使用言论自由或许并非善终。\par

XI. 但也许某个心胸狭隘、属血肉的人会说：\Quote{但我们的传报者是个陌生人和外邦人。}宗徒们呢？他们岂不也是许多民族和城市中的陌生人——他们被分散到那些地方，为的是福音能在各处自由传播，使万物不失去三重光明的照耀，或不被真理所光照；而是为那些坐在黑暗和死影中的人消除无知的黑暗？你们听过保禄的话：\Quote{叫我们往外邦人那里去，他们往受割礼的人那里去。}\BibleRef{迦 2:9} 就算犹太是伯多禄的家；保禄与外邦人有何相干？路加与阿哈雅有何相干？安德肋与厄毗鲁有何相干？若望与厄弗所有何相干？多默与印度有何相干？马尔谷与意大利有何相干？至于其余的人，不一而足，他们与所去之处又有何相干？因此，你们要么谴责他们，要么原谅我，要么证明你们这些真福音的使者正因小事而受侮辱。但既然我已就这些小事与你们争辩，现在我将采取更宏大、更哲思的视角来继续。\par

XII. 我的朋友，每个胸怀大志的人都有一个祖国，就是天上的耶路撒冷，我们在那里储存我们的公民身份。所有人都有一家——若你看这尘世，那是泥土；若你看更高处，那是我们所分有、并受命保持的灵气，我将带着它站在我的审判者面前，交代我的天上高贵和天主肖像。因此，凡通过德行和对其原型的认同而守护了这肖像的人都是高贵的。反之，凡与邪恶混杂，给自己穿上另一种形式——蛇的形式——的人都是卑贱的。这些地上的国家和家庭是我们这暂时生命和演出的玩物。因为我们的祖国就是每个人最初占据的，或许是作为暴君，或许是不幸所得；在这方面，我们全都是陌生人和朝圣者，无论我们怎样玩弄名号。家族被称为高贵，或因自古富有，或因新近升迁；出身卑微的是那些父母贫苦的，或因不幸，或因缺乏雄心。因为从上而来的高贵如何能时而开始、时而结束；且不赐给某些人，却通过特许状赐给另一些人呢？这就是我在这件事上的看法。因此，我把以坟墓或神话为荣耀之事留给你们，我自己则尽力清除欺骗，以便若可能，保持我的高贵，或重新获得它。\par

XIII. 于是，因这些缘故，我这个出身卑微、来自无闻之地的人，来到了你们中间，而且并非出于自愿，也不是自我差遣，像许多现在夺取高位的人那样；而是因为我受邀请、被迫，并随从良心的顾虑和圣神的召唤。若不然，愿我继续在此空争斗，不使任何人脱离谬误，而愿那些寻求我灵魂荒芜的人得偿所愿，若我说谎。但我既然来了，或许带着并非可轻视的能力（若我可稍为愚妄自夸），我仿效了那些贪得无厌者中的哪一个？我效法了什么投机取巧？虽然我有这许多榜样，甚至离开它们也很难罕见不坏？我曾与你们争论过什么教会或财产？虽然你们两者都绰绰有余，而其他人则太少？我们拒绝并效法了什么御令？我们曾向什么统治者奉承以反对你们？我们谴责了谁的狂妄？而对方对我做了什么？\BibleQuote{主，不要将这罪归给他们，}那时我也这样说，因为我适时记起了斯德望的话，\BibleRef{宗 7:59}，现在我如此祈祷。\BibleQuote{被人辱骂，我们就祝福；被人亵渎，我们就退让。}\BibleRef{格前 4:12}\par

XIV. 若我在此事上做错了，即当受暴虐时我忍受了，求你宽免我这错；我也曾受过他人的暴虐；我感谢我的温和竟招来了愚妄的指控。因为我这样计算，用完全高于你们任何人的考量；这些试炼，与基督所忍受的唾污和鞭打相比，算得什么？祂为了谁并靠谁帮助我们承受这些危险，祂忍受了。我根本不把它们加在一起，看作与那茨冠等值，那茨冠夺去了我们胜利者的冠冕，为了祂的缘故，我也学得因生活的艰苦而蒙受冠冕。我不认为它们值得那一根苇子，那腐朽的帝国就是被这苇子摧毁的；仅以苦胆，仅以酸醋，我们就因此治愈了苦涩的味道；仅以祂在受难中所显示的温良。祂曾被以亲吻出卖吗？祂以亲吻责备，但不击打。祂突然被逮捕吗？祂确实斥责，但随从；若你因热忱用剑砍下玛耳曷的耳朵，祂会发怒，并复原它。若有一人披着麻布逃走，\BibleRef{谷 14:51}，祂会保护他。若你为祂的逮捕者求索多玛之火，祂不会倾降；若祂接纳一个因罪悬在十字架上的强盗，祂会因祂的良善将他带入乐园。愿所有爱人之人的行为都是爱的，正如基督的一切苦难那样，对此我们所能加添的，没有比当天主甚至为我们死了，我们却拒绝宽恕邻人最微小的过犯更大的了。\par

XV. 此外，这点我也曾思考，仍继续思考；你们看是否完全正确。我以前常与你们讨论这点。这些人拥有房屋，但我们拥有住在屋内的那位；他们拥有圣殿，但我们拥有天主；此外，我们自己是永生天主的活殿，活的祭献，合理的全燔祭，完美的祭献，是的，因着钦崇天主圣三，我们成了神。他们拥有人民，但我们拥有天使；他们鲁莽的胆量，但我们有信德；他们恐吓，但我们祈祷；他们击打，我们坚忍；他们金银，但我们有纯洁的圣言。\BibleQuote{你为自己建造了宽大的房屋和宽敞的楼房，一间有天花板、窗户通透的房屋}（你认出圣经的话）。\BibleRef{耶 22:14}但这还不比我信德更高，也不比我正被带往的天堂更高。我的羊群是不是小羊群？但它没有被引向悬崖。我的羊栈是否狭窄？但它是狼不能靠近的；强盗不能进入，小偷和陌生人也不能攀越。我必定会看到它，我知道，更加宽广。而且许多现在为狼的，我必定算作我的羊，甚至可能算作牧者。这是善牧带给我的喜讯，为了祂的缘故我为羊舍命。我不怕小羊群；因为它一目了然。我认识我的羊，也为我所认识。他们就是这样认识天主并为天主所认识的人。我的羊听我的声音，这声音我从天主的圣言中听见，我从圣教父们所学得，我在各种场合同样宣讲，不随波逐流，并且我永不止息地教导；我生于其中，也将离世于其中。\par

XVI. 我按名字呼唤他们（因为他们并非像那些被数算且有名字的星辰一样无名），他们跟随我，因为我将他们牧养在休息的水边；他们也跟随每个这样的牧人，因为他们的声音他们乐意听，正如你们所见；但陌生人他们不跟随，反而逃避他，因为他们习惯了辨别自己牧人的声音与陌生人的声音。他们要逃避华伦提努及其将一分为二的学说，拒绝相信创造者不同于善者。他们要逃避\OriginalTerm{深渊与静默}{Depth and Silence}，以及神话般的\OriginalTerm{永世}{Æons}，这些确实配得深渊与静默。他们要逃避马尔西翁的由元素和数字合成的神；逃避孟他努的邪恶而阴性的灵；逃避摩尼的物质与黑暗；逃避诺瓦替安的夸耀和空谈的纯洁自许；逃避撒伯里乌的分析与混淆，以及若我可以这样说，他的吸收，将三位压缩为一位，而非将一位定义为三位格；逃避亚略及其追随者所教导的性体差异，以及他们的新犹太主义，将神性局限于\OriginalTerm{无始者}{Unbegotten}；逃避福提努的属地的基督，他始于玛利亚。但他们钦崇圣父、圣子和圣神，同一神性；天主圣父、天主圣子，以及（不要动怒）天主圣神，同一性体在三位格中，理智的、完美的、自存的，数目上分开，但在神性上不可分。\par

XVII. 让今日所有威胁我的人接受这些话；其余的人，谁愿拿去就拿去吧。圣父不能忍受被剥夺圣子，圣子也不能忍受被剥夺\OriginalTerm{圣神}{Holy Ghost}。但如果他们受限于时间，且是被造之物，那么必然会发生那样的事……因为被造的并非天主。我也不能忍受被剥夺我的\OriginalTerm{祝圣}{consecration}；一个主，一个信德，一个圣洗。如果这一点被取消，我从谁那里得到第二次呢？你们这些毁灭圣洗或重复圣洗的人，说什么呢？人没有圣神能成为属神的人吗？那没有尊敬圣神的人，与圣神有分吗？那受洗归入一个受造物和同仆的人，能尊敬他吗？不是这样，不是这样，尽管你满口言辞。我不会背叛你，\OriginalTerm{无始之父}{Unoriginate Father}啊，也不会背叛你，\OriginalTerm{独生圣言}{Only-begotten Word}啊，也不会背叛你，圣神啊。我知道我宣认了谁，弃绝了谁，又投靠了谁。我决不允许自己，在学会了信者的言语之后，再去学习不信者的言语；宣认真理，然后又与谎言为伍；下来领受祝圣，却带着更少的圣洁回去；受洗是为得以生活，却被水杀死，就像在分娩剧痛中死去的婴孩，出生同时即死亡。为什么使我同时有福和可怜，新受光照和未受光照，神圣和无神，以致连\OriginalTerm{重生}{regeneration}的希望也遭破坏呢？几句话就够了。请记住你的宣认。你受洗归入什么？归入圣父？好，但仍属犹太。归入圣子？……好……但还不圆满。归入圣神？……非常好……这才是圆满。那么，是单单归入他们，还是归入他们的一个共同名称？后者。那共同名称是什么？就是天主。相信这个共同名称，并顺利前行、为王，并从此进入\OriginalTerm{天堂之福}{Bliss of Heaven}。我想，这就是对他们更清晰的认识；愿我们都能达到这认识，在同一基督、我们的天主内，愿光荣和权能归于他，与无始之父及赋予生命的圣神，从今世直到永远，世世无尽。阿们。\par

\SourceRecord{Translated by Charles Gordon Browne and James Edward Swallow.；Nicene and Post-Nicene Fathers, Second Series；Vol. 7.；Edited by Philip Schaff and Henry Wace.；Buffalo, NY: Christian Literature Publishing Co.,；1894.；<http://www.newadvent.org/fathers/310233.htm>.}

% structure: p0001 source=h1 target=section reason=page-title

\section{讲道辞 34}

\Emph{论埃及人的到来。}\par

\Emph{这篇讲道辞于三八〇年在\Place{君士坦丁堡}宣讲，背景如下：\Place{亚历山大}宗主教\Person{伯多禄}派遣了五位他的附属主教前去祝圣冒充者\Person{玛克西穆}，以取代\Person{额我略}所占据的宝座。这引发了诸多麻烦，但最终侵入者被驱逐并流放。不久后，一支埃及船队（很可能是常规运粮船）抵达\Place{君士坦丁堡}，似乎是在某个庆节的前一天。船上的船员们次日登陆前往教堂，途经亚略派信徒占据的众多教堂，却直奔\Place{小阿纳斯塔西亚堂}。圣\Person{额我略}深感打动，鉴于最近发生的事，特别就这一行为向他们表示祝贺，于是发表了以下这篇讲道。}\par

一、我将合宜地向那些从埃及来的人讲话；因为他们热切地来到了这里，以热忱战胜了恶意——来自那条河流所滋养的埃及，它从地下涌出如雨，在其季节中如同大海——如果我也能追随那些在此事上如此雄辩之人，尽管只是微薄之量；这埃及也为我主基督所滋养：祂曾逃往埃及，现在又由埃及供应；前者是当祂逃避\Person{黑落德}屠杀婴儿时\BibleRef{玛 2:13}；现在则因父辈对子女的爱，藉着基督——饥求善者的新食粮\BibleRef{若 6:33}——是历史所载、人所相信的最大的粮食施舍；即那从天上降下、赐生命给世界的食粮，是那不可朽坏、不可分解的生命；关于祂，我现在似乎听到圣父在说：\BibleQuote{我从埃及召叫了我的儿子}\BibleRef{欧 11:1}。\par

二、因为圣言从你们那里传向众人，健康地被人相信和宣讲；而你们是所有结果实者中最优秀的，尤其是那些现在持守正确信仰的人——就我所知，我不但是这种食粮的爱好者，也是它的分施者，不仅在家，也在外。因为你们确实以慈爱供给人民和城市肉身的食物；你们也供给精神的食粮，不是给某个民族，也不是给这个或那个被狭隘界限所限制的城市——尽管其人民可能认为它非常辉煌——而是几乎给整个世界。你们带来的不是面包的饥荒或口渴的干旱\BibleRef{亚 8:11}——这不是很可怕的饥荒，避免它也容易——而是听主圣言的饥荒，现在遭受这种饥荒是最悲惨的，医治它也是最费力的，因为罪恶增多\BibleRef{玛 24:12}，而我几乎在任何地方都找不到它真正的治疗者。\par

三、你们的粮仓总管\Person{若瑟}就是如此——我也可以称他为我们的；他凭借超凡的智慧既能预见饥荒，又能通过政府的法令加以医治，用丰腴肥美的牲畜治愈了丑陋饥饿的牲畜。事实上，你们可以按自己的意愿理解若瑟是指哪位：或是那位伟大的爱慕者、创造者、永生之同名者，或是他在宝座、圣言和白发上的继任者——我们的新\Person{伯多禄}，在德行和名声上不逊于那位摧毁并粉碎中间路线的人（尽管它仍像蛇被切断后的尾巴那样微弱地扭动）；前者在经历许多争战和搏斗后在善终中离世，我深知他从上注视我们，并向为正义而劳苦的人伸出援手——这在他脱离束缚后更甚；而后者正奔向同样的结局或生命的终结，已接近天上的居民，但仍留在肉身中，所需的是为圣言提供最后的援助，并带着更丰富的供应踏上旅程。\par

四、你们是这些伟人、圣师、真理的战士和得胜者的乳儿和后裔；对于他们，无论是时间还是暴君，理智还是嫉妒，恐惧还是控告者，诽谤者——无论是公开与他们为敌，还是暗中密谋——没有谁看似我们一方的人，没有陌生人，没有金子（那隐藏的暴君，如今几乎一切都因此被颠覆并听凭掷骰子的运气），没有谄媚或威胁，没有长久远方的流放（因为只有他们不受充公影响，因他们巨大的财富即是一无所有），也没有任何其他——无论是缺席、在场或预见到的东西——能诱使他们采取较差的一方，以任何方式背叛天主圣三，或在神性上遭受损失。相反，他们在危险中变得坚强，对真宗教更加热忱。因为为基督如此受苦增进爱，对高尚灵魂而言，它仿佛是未来斗争的定金。埃及啊，这就是你们现今的事迹和奇事。\par

五、从前你们赞美我的门德斯的山羊、孟斐斯的阿匹斯（一头肥硕多肉的牛犊）、伊西丝的礼仪、奥西里斯的肢解，以及你们可敬的塞拉皮斯——一根受神话、时代和崇拜者疯狂尊崇的木料，仿佛某种未知而神圣之物，尽管它可能得到谎言的助长；还有比这些更可耻的事物，各种怪异野兽和爬虫的形象——所有这些都被基督和基督的使者战胜了，包括其他各时代显赫的人，以及我刚才提到的教父们；因着他们，令人赞叹的国度啊，你今天比所有其他——无论是古代还是现代历史中的——加起来更加著名。\par

六、因此，我拥抱并问候你们，最高贵的子民，最虔诚的基督徒，满怀热心，堪当你们的领袖；因为我觉得没有比这更伟大的话可以对你们说，也没有更好的方式来欢迎你们。我以唇舌稍许问候你们，却以我情感的悸动由衷地致意。\BibleRef{迦 2:9}啊，我的人民，因为我称你们为我的子民，同一心意，同一信仰，受教于同一的教父们，朝拜同一的天主圣三。我的人民，因为你们是我的，尽管在那些嫉妒我的人眼中似乎并非如此。为使这些心加伤痛，看，我在众目睽睽和众目未睹的证人面前，伸出右手相交之礼。我以此新善举消除旧毁谤。啊，我的人民，因为你们是我的，尽管我这样说是最卑微的人在要求最伟大的事物，因为圣神的恩宠使同心之人同享尊荣。啊，我的人民，因为你们是我的，尽管相隔遥远，因为我们被神性地结合在一起，\BibleRef{依 62:4}且与肉性之人的结合全然不同，因为身体在空间上合一，灵魂却由圣神凝聚。啊，我的人民，你们从前学习为基督受苦，但现在若听从我，就当学习不再行事，而视行事的能力为足够的收益，并认为无论在过去你们忍耐的日子，还是在现今你们温良的日子里，都是在向基督献上祭献。啊，人民，主已预备好向你们行善，正如向你们的仇敌行恶。啊，人民，主从万民中拣选了你们归祂自己；啊，你们是刻在主手上的子民，主对你们说：\Quote{你是我的意愿}；又说：\Quote{你的城门是雕刻的}，以及其他一切论及得救者的话。啊，人民——不要因我贪得无厌多次重复你们的名字而惊奇，因为我喜悦这样不停地呼唤你们，如同那些对某些景象或声音乐此不疲的人。\par

七、然而，天主的人民，也是我的人民，你们昨日在海上的集会也十分美好，若曾有景象赏心悦目，那便是如此：当时我见海面如森林，被一座人手造的云遮盖，你们的船只美丽而迅疾，如同列队而行，顺风微拂，仿佛特意护送你们，将你们的海上之城吹向京城。但我们如今所见的这次集会更为美丽、更为壮观。因为你们并未急于混入更大的人群，也不以人数衡量宗教，也不满足于做一盘散沙，而非由天主圣言净化的人民；而是正当地将凯撒的归凯撒，又将天主的归天主；前者出于习惯，后者出于敬畏；你们以货物养民之后，自己前来由我们喂养。因为我们也分施谷物，我们的分施也许不比你们的逊色。\BibleQuote{你们来，吃我的饼，喝我给你们调和的酒吧。}\BibleRef{箴 9:5}我与智慧一同邀请你们赴我的筵席。因为我对你们的美意称赞不已，并急于回应你们的热切心意，因为你们来到我们这里，如同进入自己的港湾，奔向自己的同类；你们看重同信之谊，并认为这是怪事：那些侮辱更高事物的人彼此和谐、思想一致，且认为可以通过共同的阴谋来弥补各人虚假的错谬，如同绳子由绞合而更加坚固；而你们却与那些心意相同、与你们更有理由联合的人不相往来、不相汇聚，因为我们在天主性中也合而为一。为使你们看出，你们来到我们这里并非徒然，也没有停泊在陌生人和外邦人的港口，而是停泊在你们自己人当中，并且由圣善的圣神引导得当——我们将向你们略述有关天主的事；你们要认出自己的东西，如同人们用臂章识别亲属一样。\par

八、我发现在存在之物中有两种最高的区别，即：统治与服务；这不是像我们当中由暴政切割或由贫穷分割的那样，而是由本性所区分的——如果有人喜欢用这个词的话。因为那首要的也是超乎本性的。前者是创造性的、本原性的、不变的；而后者是被创造的、从属的、变化的；或者说得更清楚些，一个是超时间的，另一个是在时间之下的。前者被称为天主，存在于三个最伟大的位格中：原因、创造者与成全者；我指的是圣父、圣子与圣神，他们既非彼此分离以致本性分裂，也非如此收缩以致被一个位格所限定；前者是亚略异端的疯狂，后者是撒伯流派异端的主张；但他们一方面比完全分裂者更为合一，另一方面比完全单一者更为丰富。另一个层面是在我们中的，被称为受造界，虽然按其接近天主的程度，一个可能比另一个更为尊崇。\par

九、既然如此，谁若站在主的一边，就让他与我们同行，\BibleRef{出 32:26}让我们朝拜三位中唯一的天主性；不将任何屈辱之名归于那不可接近的光荣，口中常有三一之主的赞美。因为既然我们甚至不能恰当地描述祂本性的伟大，因其无限与不可界定，我们又怎能主张祂的屈辱呢？但若有人与天主疏远，因而将唯一至高的本体分裂为不平等的本性，倘若这样的人不被刀剑所割，与不信者同分，\BibleRef{路 12:46}现在和将来都收获他那恶念的恶果，那倒是奇事了。\par

十、关于圣父，我们该说什么？凡被自然观念所占据的人，都一致承认祂，虽然祂已遭受了不敬的开端，被古代的革新首先分割为\Quote{善者}与\Quote{创造者}。至于圣子和圣神，请看我们将如何简洁明了地论述。若有人能说二者中任何一位是可变的或会改变的；或说祂能在时间、地方、能力或行动上被衡量；或说祂不是本性善的，或不是自动的，或不是自由的，或是仆役，或是颂唱者；或说祂恐惧，或是得到自由的接受者，或不被算为与天主同在；——让他证明这一点，我们便默认，并将因我们同仆的尊威而受荣耀，即使我们失去了我们的天主。但若圣父所有的一切也属于圣子，除了父的主动成因；圣子所有的一切也属于圣神，除了祂的子身份，以及凡关于祂为\Quote{我}，为人，为我的救恩而\Quote{降生成人}所说的话——祂取了我的，以便藉这新的混合而赐予祂的——那么，请停止你们的空谈，尽管已经太迟了，你们这些虚言成空、立刻坠地的诡辩家！因为\Quote{以色列家啊，你们为何要死呢？}（\BibleRef{则 18:31}）——如果我能用圣经的话为你们哀悼。\par

十一、至于我，我也尊敬圣言的许多名号，它们如此众多、崇高而伟大，连魔鬼也尊重它们。我也尊敬圣神的同等身分，并畏惧那针对亵渎祂者所发出的威胁。亵渎不是不称祂为天主，而是将祂与天主性分离。在此你该注意：被亵渎的是主，被报复的是圣神，显然祂是主。我不能在接受光照之后仍处于未光照的状态，即在我受洗所归入的三位中，以不同的印记标记任何一位；这样我实在是被埋在水里，所领受的不是重生，而是死亡。\par

十二、我敢说一些什么，啊，天主圣三，愿我的愚昧得到宽恕，因为冒险的是我的灵魂。我也是天主的肖像，属天的光荣的肖像，虽然我居于地上。我不能相信我是被与我同等者所拯救。如果圣神不是天主，就让祂先成为天主，然后让祂——与我同等者——使我成为神。但现在，恩宠，或更确切地说，赐予恩宠者所行的欺瞒何其大——相信天主，却空无天主地离去；由一套问题和宣认导向另一套结论。唉，这美好的名声啊！如果洗后我变黑了，如果我要见到那些尚未洁净者比我更光明，如果我被施洗者的异端所欺骗，如果我寻求更强的圣神却找不着。在你对第一次洗礼存有恶念之前，请给我第二次泉源。你为什么吝惜给我完全的再生？你为什么使我——作为圣神的殿如同天主的殿——成为受造物的居所？你为什么尊崇我的一部分，而羞辱另一部分，对天主性作出错误的判断，因而将我与恩赐隔绝，或更确切地说，藉着恩赐将我分裂？要么尊崇全部，要么羞辱全部，新的神学家啊！这样，如果你是邪恶的，至少你与自己一致，不偏袒地判断同等的本性。\par

十三、总结我的言论：——与革鲁宾一同光荣祂，他们联合三个\Quote{圣}为一个\Quote{主}（\EditorialNote{依撒意亚 6:3}），并以此藉翅膀向勤勉者暗示原始本体。与达味一同接受光照，他对光说：\Quote{在你的光中，我们必得见光}，即是在圣神中我们将看见圣子；还有什么比这光线更远呢？与若望一同发出雷鸣，论及天主，他发出不高也不属地的声音，而是崇高和属天的：\Quote{在起初已有圣言，圣言与天主同在，圣言就是天主}（\BibleRef{若 1:1}），\Quote{真天主}出于\Quote{真父}，而非一位仅被尊为\Quote{子}名的善仆；以及\Quote{另一位护慰者}（与说话者——天主的圣言——相\Quote{另}）。当你读到\Quote{我与父原是一体}时，你当注视本体的合一；但当你看见\Quote{我们要到他那里去，并要在他那里作我们的住所}（\BibleRef{若 14:23}）时，记住位格的区分；当你看见\Quote{父、子、圣神}这些名字时，思想三个位格。\par

十四、当你研读《宗徒大事录》时，与路加一同受感动。你为什么将自己与阿纳尼雅和撒斐辣放在一起？他们是徒然的盗用者（如果盗用自己的财产算是徒然的事），他们侵吞的并非银子或其他廉价无价值之物，像一块金子（\BibleRef{苏 7:21}）或古时一个贪心的兵士所行的\Quote{双德拉克玛}，而是偷窃了天主性本身，并撒谎——不是对人，而是对天主，正如你所听到的。什么？你竟不尊重圣神的权能吗？祂随意向谁、何时、如何吹气。祂在洗礼之前来到科尔乃略及其同伴身上，在洗礼之后藉宗徒之手来到其他人身上；这样，从两方面——既因祂以主人的姿态而非仆人的姿态来临，又因祂被寻求以成全人——圣神的天主性得到证明。\par

十五、与保禄一同谈论天主，他曾被提到第三层天（\BibleRef{格后 12:2}），他有时列举三个位格，且以变化的顺序，不保持同一顺序，而是将同一位格有时列为第一，有时第二，有时第三；为什么？为显示本性的平等。有时他提到三个，有时两个或一个，因为未提到的也被包含在内。有时他将天主的行动归于圣神，仿佛与祂毫无区别；有时他以基督代替圣神；有时他区分位格说：\Quote{只有一个天主，万物都出于祂，而我们也归于祂；也只有一个主耶稣基督，万物都藉着祂而有，而我们也藉着祂而有}（\BibleRef{格前 8:6}）；有时他又将独一天主性集合起来：\Quote{因为万物都出于祂，藉着祂，在祂内}（\BibleRef{罗 11:36}）——即藉着圣神，正如圣经多处所示。愿光荣归于祂，直到永远。阿们。\par

\SourceRecord{Translated by Charles Gordon Browne and James Edward Swallow.；Nicene and Post-Nicene Fathers, Second Series；Vol. 7.；Edited by Philip Schaff and Henry Wace.；Buffalo, NY: Christian Literature Publishing Co.,；1894.；<http://www.newadvent.org/fathers/310234.htm>.}

% structure: p0001 source=h1 target=section reason=page-title

\section{讲道词第三十七}

论福音的话，\Quote{当耶稣讲完了这些话}，等等——\BibleRef{玛 19:1}\par

I. 拣选渔夫的耶稣，祂自己也使用网，且更换地方。为什么？不仅为祂能藉着祂的临在赢得更多爱慕天主的人；而且，依我看来，也为祂能祝圣更多的地方。对犹太人，祂成了犹太人，为能赢得犹太人；对法律之下的人，祂成了法律之下的，为能赎出那些在法律之下的；对软弱的人，祂成了软弱的，为能拯救软弱的人。祂成了众人中的一切，为能赢得一切。为什么我说\Quote{众人中的一切}？因为我发现，保禄不能忍受说自己的事，救主却忍受了。因为祂不仅成了犹太人，不仅将一切怪诞可鄙的名称加于己身，甚至其中最为怪诞的，即罪与咒诅本身；并非祂是这样，而是如此称呼。因为使人脱离罪的那位如何能成为罪，且从法律的咒诅中救赎我们的那位如何能成为咒诅？但这是为使祂将谦逊的表现推延至此，并塑造我们，使我们进入那产生高举的谦逊。如我当时所说，祂成了渔夫；祂屈尊就卑，降及一切；祂撒网；祂忍受一切，为从深渊中拉上鱼来，也就是在生活不安而苦涩的波浪中游荡的人。\par

II. 因此如今，当祂讲完了这些话，祂离开加里肋亚，来到约但河对岸的犹太境内；祂适宜地居住在加里肋亚，为使坐在黑暗中的百姓看见大光。\BibleRef{依 9:1}祂迁往犹太，为劝人从文字兴起，跟随圣神。祂时而山上教导，时而平原讲论；时而渡船，时而斥责风浪。也许祂入睡，为祝福睡眠；也许祂疲倦，为祝圣辛劳；也许祂哭泣，为使眼泪蒙福。祂从一处移往另一处，那不受任何地方所局限的；那不受时间限制、无形体、无界限的，同一的这位过去与现在都是；祂既超越时间，又降生于时间；曾是不可见的，如今被看见。\BibleQuote{祂在起初，与天主同在，祂就是天主。}\BibleRef{若 1:1}\Quote{是}字出现三次，为以数目确认。祂放下了祂所是的，接受了祂所不是的；并非祂成为两个，而是祂屈尊成为由两者所成的一个。因为两者都是天主：那接受的与那被接受的；两个本体会合于一个，不是两个儿子（我们不要错解这结合）。这位如此伟大的——但我遭遇了什么？我落入了人类的言语。因为那绝对者如何能被称为\Quote{如此伟大}，那无量的如何能被称为\Quote{如此}？但请宽恕这话，因为我用有限的工具谈论最伟大的事物。那伟大、宽容、无形、无体的本性将容忍这一点，即我的言语如同有身体，且弱于真理。因为如果祂曾屈尊于肉体，祂也将容忍这样的言语。\par

III. 有许多群众跟随祂，祂就在那里，在那群众更多的地方医治了他们。如果祂停留在自己的崇高上，如果祂没有屈尊于软弱，如果祂保持祂所是的，让自己不可接近、不可理解，也许只有少数人会跟随祂——也许连少数也没有，或许只有梅瑟——且祂仅能勉强看见天主的背后。因为祂穿透了云彩，要么是脱离了身体的重量，要么是脱离了感官；因为祂，作为有形体、使用物质眼睛的，怎能凝视天主的精妙、无形体，或我不知道如何称呼的东西呢？但因祂为我们脱下自己，因祂降下（并论及一种\OriginalTerm{空虚自己}{exinanition}，可以说是一种脱去和祂荣耀的减损），祂由此成为可理解的。\par

IV. 同时，请宽恕我，我又一次体验人的情感。我为我的基督感到愤慨和悲痛（但愿你们同情我），当我看见我的基督因那最应受尊敬的事而受辱。请告诉我，祂因为了你们而谦卑，就该受辱吗？祂因为眷顾受造物，就成为受造物吗？祂因为看顾那些受时间约束的人，就成为时间约束的吗？不，祂承受一切，忍耐一切。\BibleRef{格前 13:7}这有什么惊奇？祂忍受了鞭打，忍受了唾污，为我尝了苦胆。甚至如今，祂仍忍受被石头砸，不仅被那些侮辱祂的人，也被我们这些似乎尊敬祂的人。因为用有形体的名称谈论无形体者，也许正是那些侮辱祂、用石头砸祂之人的作为；但请宽恕，我再说一次我们的软弱，因为我并非故意用石头砸祂；只是没有其他言语可用，我们就用我们所有的。祢被称为圣言，祢却超越圣言；祢超越光，却被称为光；祢被称为火，不是可感觉的，而是因为祢炼净轻浮无价值的事物；一把剑，因为祢将坏的从好的中分开；一把簸箕，因为祢簸扬禾场，吹走所有轻浮如风之物，将沉重饱满的收在天上的仓里；一把斧头，因为祢砍下不结果的无花果树，在长久忍耐之后，因为祢斩断邪恶的根；门，因为祢引入；道路，因为我们直行；羔羊，因为祢是祭品；大司祭，因为祢献上身体，即圣子；因为祢出自圣父。我又一次激动人的舌头；又一次有人妄议基督，或者不如说妄议我，我蒙恩作了圣言的使者。我像若翰——\BibleQuote{在旷野里有呼喊者的声音}\BibleRef{玛 3:3}——一片曾经干燥的旷野，现在却太过拥挤了。\par

V. 但是，如我所说，回到我的论点；为此，一大群人跟随了祂，因为祂俯就我们的软弱。接下来呢？经上说，法利塞人也来见祂，试探祂，对祂说：\BibleQuote{人无论什么缘故都可以休妻么？}法利塞人再次试探祂；读法律的人再次不认识法律；解释法律的人再次需要别人教导他们。撒杜塞人关于复活试探祂，经师关于成全质问祂，黑落德党人关于人头税质问祂，其他人关于权柄质问祂，这还不够；还有人必须就婚姻问题询问那不可试探者、婚姻的创造者、那从第一因造出整个人类种族的那位。祂回答他们说：\BibleQuote{你们没有念过，那起初造人的，是造男造女么？}祂知道如何解答他们的一些问题，并遏制另一些问题。当有人问：\BibleQuote{你凭什么权柄做这些事？}祂自己，因为问话者的极端无知，用另一个问题回答：\BibleQuote{若翰的洗礼是从天上来的，还是从人间来的？}祂从两边使问话者陷入困境，以便我们也能效法基督的榜样，有时阻止那些过于多事与我们对论的人，用更加荒谬的问题来化解他们问题的荒谬。因为我们在虚浮之事上有时也聪明，如果我夸耀愚妄之事的话。但当祂看到一个需要推理的问题时，祂并不认为问话者不配得到明智的回答。\par

VI. 你所提出的问题，在我看来，是对贞洁的尊重，并需要一个友善的回答。贞洁，关于这一点，我看到大多数人的倾向不良，他们的法律不平等且不规则。因为，为何他们约束女人，却放纵男人；女人对丈夫的床行恶便是淫妇，为此法律的惩罚非常严厉；但如果丈夫对妻子行淫，他无需交代？我不接受这种立法；我不赞同这种习俗。制定法律的是男人，因此他们的立法对女人严厉，因为他们也将子女置于父亲的权下，却忽略了较弱的性别。天主并非如此；而是说：\BibleQuote{应孝敬你的父亲和你的母亲，这是第一条带应许的诫命；为使你获得幸福，并在地上延年益寿。}又说：\BibleQuote{咒骂父亲或母亲的，应处死刑。}同样地，祂对善给予尊荣，对恶给予惩罚。并且，\BibleQuote{父亲的祝福，坚固子女的房屋；母亲的诅咒，拔除子女的根基。}\BibleRef{德 3:11}看立法的平等。男人和女人有一位造主；子女对双亲欠下同一份债。\par

VII. 那么，你如何要求贞洁，而你自己却不遵守？你如何要求你所不给予的？你同样是身体，为何不平等地立法？如果你查考更糟糕的——女人犯了罪，亚当也犯了罪。\BibleRef{创 3:6}蛇欺骗了他们两个；没有发现一个更强、一个更弱。但你要看更好的吗？基督藉祂的受难拯救了二者。祂为男人成了血肉吗？也为女人。祂为男人死了吗？女人也因祂的死而得救。祂被称为达味的后裔，\BibleRef{罗 1:3}因此你也许以为男人受尊荣；但祂由童贞女所生，这站在女人一边。祂说：\BibleQuote{二人成为一体}；所以，同一肉身应有同等的尊荣。保禄也以基督的榜样为贞洁立法。如何，以什么方式？他说：\BibleQuote{这奥秘是伟大的，但我是指基督和教会说的。}\BibleRef{弗 5:32}妻子藉着丈夫尊敬基督是好的；丈夫不藉着妻子羞辱教会也是好的。他说：\BibleQuote{妻子应当敬重丈夫}，因为如此她尊敬基督；但祂也吩咐丈夫爱妻子，因为如此基督爱教会。那么，让我们进一步思考这句话。\par

VIII. \BibleQuote{搅动牛奶，必出奶油；}\BibleRef{箴 30:33}审视这一点，也许你会在其中发现一些更具营养的东西。因为我认为圣言在此似乎不赞同第二次婚姻。因为，如果有两个基督，就可以有两个丈夫或两个妻子；但基督只有一个，教会只有一个头，那么也只有一个肉身，第二个应被拒绝；如果第二个被阻止，第三者还有什么可说的？第一次是法律，第二次是宽容，第三次是过犯，再往后便是猪猡行径，这样的邪恶甚至没有许多例子。现在，法律准许因任何缘故离婚；但基督不是因任何缘故；祂只允许因淫乱而分离；在其他一切事上，祂命令忍耐。祂允许休弃淫妇，因为她败坏了后代；但在其他一切事上，让我们忍耐和承受；或者，凡是接受了婚姻之轭的，要坚忍和忍耐。如果你在她身上看见纹路或斑点，就除去她的装饰；如果是急躁的舌头，就约束它；如果是卖笑，就使其端庄；如果是无度的花费或饮酒，就减少它；如果不是时宜的出门，就束缚它；如果是高傲的眼，就惩戒它。分离者与被分离者，谁处于危险之中，是不确定的。经上说：\BibleQuote{你的泉源只应属于你，不可让陌生人分享；}\BibleRef{箴 5:17}又说：\BibleQuote{愿你的恩宠之驹和爱情之鹿与你相伴；}\BibleRef{箴 5:19}那么你要小心，不要成为别人的溪流，也不要讨别人喜欢超过你自己的妻子。但如果你被带到别处，那么你也为你的伴侣制定了淫荡的法律。救主如此说。\par

九、但法利塞人呢？在他们听来，这话似乎刺耳。是的，因为他们也厌恶其他高尚的言辞——无论是古代的法利塞人，还是今日的法利塞人。因为使人成为法利塞人的，不仅是种族，还有性情。因此，我也把那按品行归于这类人的人，算作亚述人或埃及人。那么，法利塞人说什么呢？\BibleQuote{若是人与妻子的情形是这样，倒不如不娶。}（\BibleRef{玛 19:10}）哦，法利塞人，你现在才明白\Quote{不娶更好}吗？你先前不曾看见寡居、孤苦、夭亡、喜乐后的哀悼、婚礼后的葬礼、无子女，以及与此相关的一切喜剧或悲剧吗？两者都是最恰当的说法。结婚是好的，我也承认，因为\BibleQuote{婚姻在各方面都应当受人尊重，床笫也应是无玷污的。}（\BibleRef{希 13:4}）对节制的人是好，对贪得无厌、想给肉体过分尊荣的人则不然。当婚姻只是婚姻、结合及渴望延续子嗣时，婚姻是光荣的，因为它给世界带来更多中悦天主的人。但当婚姻激起情欲、用荆棘围困我们、并像开辟出罪恶之路时，那么我也说：\Quote{结婚是不好的。}\par

十、婚姻是光荣的，但我不能说它比童贞更高超；因为童贞若不比一件好事更好，就算不得什么大事。然而，你们这些负轭的妇女，不要发怒。我们必须服从天主，而非服从人。但你们，童贞女和已婚者，要结合为一，在主内合而为一，彼此装点。若没有婚姻，就不会有独身者。因为童贞女从何处进入此生呢？婚姻若非为天主和生命孕育了童贞的果实，便不会受人尊敬。你也要尊敬你的母亲，你从她而生。你也要尊敬那既出自母亲又身为母亲的人。她不是母亲，却是基督的净配。可见的美并不隐藏，但那不可见的美却是天主可见的。君王女儿的一切荣耀都在内部，以金线绣成，或用行动，或用默观绣成。那负轭的，也要在某种程度上属于基督；童贞女则完全属于基督。愿那一个不完全被世界缠住（\BibleRef{路 8:14}），愿那一个完全不属于世界。因为对负轭者来说，那只是部分；对童贞女来说，却是全部。你选择了天使的生活吗？你被列于无轭者之中吗？不要沉沦于肉体；不要沉沦于物质；否则你虽未嫁，却与物质结合。淫荡的眼睛不守护童贞；淫词荡语与恶者混合；行走放荡的脚表明疾病或危险。愿心灵也是童贞的；不要游荡；不要徘徊；不要在自身中携带邪恶事物的形象（因为形象是淫乱的一部分）；不要在灵魂中制造可憎之物的偶像。\par

十一、但祂对他们说：\BibleQuote{这话不是人人都能领受的，唯独那些得了恩赐的人才能领受。}你看见这事的高超吗？它几乎是不可理解的。因为由肉体所生的不再为肉体而生，这确实超乎肉体。那与肉体结合却按肉体生活，反超越其本性，确实是天使般的。肉体把她束缚于世界，但理性把她领向天主。肉体压她向下，理性给她翅膀；肉体捆绑她，但欲望释放了她。童贞女啊，你要全心专注于天主（我也同样吩咐男人和女人）；在其他方面，不要像众人那样看待光荣的事：家族、财富、王座、王朝，以及那表现在肤色与肢体组合中的美——那是时间与疾病的玩物。你若已将全部爱情倾注于天主；若你没有两个欲望对象——暂时的与永恒的、可见与不可见的——那么你已被选择的箭刺透，并如此学会了新郎的美，以致你也能用婚歌说：\Quote{你是甘饴的，全然可爱。}\par

十二、你看见水在铅管中受约束，因被高度压缩并导向一点，常常如此脱离水的本性，以致后面的水常向上涌流。同样，你若约束你的欲望，完全依附于天主，你就不会堕落；你不会分散；你将完全属于基督，直到你看见你的新郎基督。你要保持自己不可接近，无论在言语、行为、生活、思想和行动上。恶者从四面八方干扰你；他到处窥探你，看从何处打击，何处伤害你；愿他找不到任何裸露而易于他打击的地方。他见你越纯洁，就越努力玷污你，因为闪亮衣服上的污渍更显眼。愿眼不吸引眼，欢笑、亲昵、黑夜也不吸引眼，免得黑夜带来毁灭。因为那渐渐被诱惑并偷走的，当时不易察觉，却终究达到邪恶的完成。\par

十三、祂说：\Quote{众人不能都领受这话，惟独那些得恩赐的人，才能领受。}你们听见\Quote{得恩赐}这话，不要以异端方式理解，而引进本性的差异：属地的、属灵的和混合的。因为有人心术不正，以为有些人的本性是完全败坏的，有些人的本性是得救的，另一些人的性情则随其意志所趋，或好或坏。因为我认为人可能有某种禀赋，有人多，有人少，但这禀赋本身并不足以使人达于成全，而是需要理性将其发挥出来，使本性付诸行动，犹如用铁击燧石便生出火来。当你们听见\Quote{得恩赐}时，要加上：\Quote{也是赐给那些蒙召并愿意如此倾向的人。}因为当你们听见\BibleQuote{不在乎那愿意的，也不在乎那奔跑的，只在乎那显示慈悲的天主}时，\BibleRef{罗 9:16}我劝你们也这样想。因为有些人过于自负其成就，将一切归于自己，而不归于造他们、赐他们智慧、供应他们美善的天主；这些人便由此话学到：连愿意行善也需要天主的助佑，甚至选择正确之事也是神圣的，是天主慈悲的恩赐。因为我们必须自主，同时我们的救恩也必须来自天主。故此祂说\Quote{不在乎那愿意的}——即不只在乎那愿意的，也不只在乎那奔跑的，而在乎那显示慈悲的天主。其次，既然愿意也是出于天主，祂有理由将一切归于天主。无论你们怎样奔跑、怎样角力，仍需要一位来颁发桂冠。除非主建造房屋，建造的人就枉然劳力；除非主看守城池，看守的人就枉然警醒。祂说：\BibleQuote{我知道，跑得不快，力战不胜，胜利不属于勇士，良港不属于好舵手，因为得胜和使船安全靠港都是天主的作为}\BibleRef{训 9:11}\par

十四、在另一处也有类似的话和理解，或许我应该将它附加在已说过的内容之后，好把我的财富也分施给你们。\Person{载伯德}儿子的母亲，出于母爱之情，无知地请求了一件她不知轻重的事，但因她过分的爱和对自己孩子应有的慈爱，也是可原谅的。因为没有什么比母亲更富爱心的了——我说这话是为订一条尊崇母亲的法律。于是，他们的母亲求耶稣让她的两个儿子，一个坐在祂右边，一个坐在祂左边。但救主怎么说？祂先问他们能否饮祂将要饮的杯；当他们表示能，而救主接受了这表示（因为祂知道他们正因此而被成全，或者说将要因此而被成全）时，祂说什么？\Quote{他们必要饮这杯；但坐在我的右边和左边，不是我可以赐的，而是为谁预备的，就赐给谁。}那么，理性就毫无作用吗？劳作毫无作用？推理毫无作用？哲学毫无作用？禁食毫无作用？守夜、睡在地上、流如泉的眼泪，都毫无作用吗？难道不是因着这些，而是因着某种抽签的方式，使\Person{耶肋米亚}成圣，而另一些人自母胎就被疏远了吗？\par

十五、我担心会有某种怪异的推论出现，例如灵魂在此之前曾生活在别处，然后被束缚于这个身体，并从那另一生命中，有些人领受了预言的恩赐，另一些人则被判罪，即那些生活不善的人。但这观念太荒谬，且违反教会的传统（别人若喜欢，可以玩味这些学说，但对我们而言，玩味它们是不安全的）；因此，我们在此处也必须在\Quote{为谁预备的}这话上加上\Quote{就是那些相称的人}，这些人不仅从父那里领受了这特征，也自己赋予了它。\par

十六、因为有些阉人是从母胎生来就是阉人，等等。我很想就阉人讲一些大胆的话。你们这些生来的阉人，不要骄傲。因为从节制方面说，这或许是不情愿的。因为它未经考验，你们的节制也未经试验证明。因为天生的善不算功劳，出于意愿的善才值得称赞。火有何功劳，因为它燃烧是出于本性？水有何功劳，因为它下流是造物主赋予的特性？雪因寒冷、太阳因照耀，又当得起什么感谢？——太阳即使不愿意，也要照耀。你们若愿意，可因选择更好的事而自夸。你们若在血肉之中使自己成为属神的，若被铅重的肉躯往下拖，却因理性获得翅膀，若出身卑贱，却被发现是属天的，若被束缚于肉身，却显示出超越肉身，你们便可自夸。\par

十七、既然天生的贞洁不算功劳，我就向阉人要求另外的事。不要在有关天主性上行邪淫。既然已与基督成婚，就不要羞辱基督。既然由圣神所成全，就不要使圣神与你平等。\Person{保禄}说：\BibleQuote{若我还求人的欢心，我就不是基督的仆人了}\BibleRef{迦 1:10}。若我崇拜受造物，我就不应被称为基督徒。因为基督徒为何宝贵？岂不是因为基督是天主吗？除非我与祂因爱结合只是出于人的情欲。然而，我尊敬\Person{伯多禄}，却不被称为伯多禄派；尊敬\Person{保禄}，也从未被称为保禄派。我不能容许自己因一个生于天主的人的名字而被称呼。所以，如果你们因相信祂是天主而被称为基督徒，愿你们永远如此称呼，并保持在名与实中；但如果你们只因对基督有感情而被称为基督徒，那么你们所归给祂的，并不比那些因某种实践或事实而来的其他称号更多。\par

XVIII. 试想那些热衷于赛马的人。他们按自己所支持的颜色和派别而被称呼。无需我指明，你们也知道这些名称。若你们以同样方式获得基督徒之名，那么单凭这个称号，即使你们引以为傲，也是微不足道的。但若你们是因相信祂是天主，就要以行为显明你们的信德。如果圣子是受造物，那么你们直到如今仍在敬拜受造物而非造物主。如果圣神是受造物，你们的受洗就是徒然，只在两边正确，甚至两边也不正确；但在一方面你们全然危险。想像天主圣三是一颗珍珠，四面都一样，同样闪耀。若珍珠的任何部分受损，整颗宝石的美就消失了。同样，当你们为光荣圣父而侮辱圣子时，圣父并不接受你们的荣耀。圣父不以圣子的侮辱为光荣。因为\BibleQuote{智慧之子使父亲喜乐}，\BibleRef{箴 10:1}，何况圣子的荣耀不更成为圣父的荣耀吗？若你们也接受这句话：\BibleQuote{我儿，不要因你父亲的侮辱而自夸}，\BibleRef{德 3:10}，同样，圣父也不以圣子的侮辱为光荣。若你们侮辱圣神，圣子也不接受你们的荣耀。因为圣神虽不似圣子那样来自圣父，却同出一源。要么光荣全体，要么侮辱全体，以保持心意一致。我不能接受你们一半的虔诚。我要你们全然虔诚，但必须按我所愿的方式。请原谅我的挚爱：我甚至为那些恨我的人悲伤。你们曾是我的肢体，即使现在被割断了；也许你们还会再成为肢体，因此我慈爱地说话。以上是为阉者而说的，好使他们于天主性方面保持贞洁。\par

XIX. 因为不仅是身体的罪被称为奸淫和通奸，而且你犯的任何罪，尤其是违反神圣的事，也是如此。也许你问我们如何证明这一点：——经上说：他们随着自己的发明行淫。你看见无耻的奸淫行为吗？还有，他们在树林中行通奸。你们看见一种奸淫的宗教吗？那么就不要犯属灵的奸淫，却保持身体贞洁。不要因你未能在你\Emph{能}犯奸淫之处保持贞洁，而显出你身体的贞洁是不情愿的。你为何行不虔敬？为何你急于行恶，以致称一个人为阉者或恶棍都一样？使你们自己站在男子的一边，即使如此晚，也要有些男子汉的思想。避开闺房；不要让宣告的耻辱加于名字的耻辱之上。你们希望我们在这篇讲论中再坚持一会儿，还是已经对我们所说的感到厌倦？不，接下来让阉者也受尊敬。因为这话是赞美的话。\par

XX. 祂说：有些阉者是从母腹生来如此的；有些阉者是被人阉割的；有些阉者为天国而自阉。\BibleQuote{能领受的，就领受吧。}我想这篇讲论会从身体分离，用身体的形象代表更高的事；因为将意义停留于身体上的阉者是渺小且非常软弱的，不配那圣言；我们必须额外理解一些配得圣神的事。于是，有些人似乎本性倾向于善。当我说本性时，我并非轻视自由意志，而是假定两者——向善的倾向，以及使自然倾向产生效果的东西。还有另一些人，理性将他们割断情欲而洁净。我想这些就是那些被人阉割的人所指的，当教导之言区分善恶，拒绝其一而命令另一（如经文：\BibleQuote{离恶行善}），便成就了属灵的贞洁。我赞成这种阉割；而且我极赞美教师和受教者，因为前者高贵地施予了割除，后者更高贵地忍受了割除。\par

XXI. 还有那些为天国而自阉的阉者。另外，那些没有遇到老师的人，他们自己成了值得称赞的老师。没有父亲或母亲，没有司铎或主教，也没有任何受委托教导的人，教导你你的本分；但通过在你内发动理性，并通过你的自由意志点燃善的火花，你使自己成为阉者，并获得了如此的美德习惯，以至于向恶的冲动对你几乎成为不可能。因此，我也赞美这种阉割，或许甚至超过其他。\BibleQuote{能领受的，就领受吧。}选择你愿意的部分；要么跟随老师，要么作自己的老师。唯有一事是可耻的——就是情欲未被根除。它们如何被根除并不重要。老师是天主的受造物；你也有同样的起源；无论是老师抓住了这恩宠，还是这善是你自己的——同样都是好的。\par

XXII. 只要我们割断自己脱离情欲，\BibleQuote{免得任何苦根长出来烦扰我们}，\BibleRef{希 12:15}；只要我们跟随那肖像；只要我们敬畏我们的原形。割除身体的情欲；也割除属灵的情欲。因为灵魂比身体贵重多少，洁净灵魂就比洁净身体贵重多少。如果洁净身体是值得称赞的行为，请看，我请求你们，洁净灵魂是何等更大更高的事。割除亚略的不虔敬；割除撒伯流的错误意见；不要超过当有的联合，或错误地分开；不要将三个位格混淆为一个，或制造三个本性的差异。那一位如果正确理解就是值得赞美的；三位如果正确分开，当分开的是位格，而非天主性时。\par

二十三、我亦为平信徒制定此律，也吩咐所有司铎和受命治理者。凡天主赐予能力相助的，都来助圣言一臂之力。制止谋杀、惩罚奸淫、惩戒盗窃，固然是大事；但以法律确立虔敬、赐予健全教义，更为重要。我的言语在捍卫天主圣三上，远不及你们的诏令有效——只要你们能约束心怀恶意者，帮助受迫害者，阻止杀人者，并防止人被杀。我所说的不仅是肉体的屠杀，更是属灵的屠杀。因为一切罪都是灵魂的死亡。我的演讲到此结束。\par

二十四、但我还需要为聚集在此的人献上祈祷。丈夫与妻子、统治者与被统治者、老年人、少年人、少女，无论何种年龄，都当忍受钱财或身体的一切损失，但惟有一件事不可容忍——失去天主性。我钦崇圣父，我钦崇圣子，我钦崇圣神；更确切地说，我们钦崇祂们；我——正在发言者——在万有之前，在万有之后，与万有同在，都在同一位基督我们的主内，愿光荣与德能归于祂，直到永远。阿们。\par

\SourceRecord{Translated by Charles Gordon Browne and James Edward Swallow.；Nicene and Post-Nicene Fathers, Second Series；Vol. 7.；Edited by Philip Schaff and Henry Wace.；Buffalo, NY: Christian Literature Publishing Co.,；1894.；<http://www.newadvent.org/fathers/310237.htm>.}

% structure: p0001 source=h1 target=section reason=page-title

\section{演讲集第三十八篇}

\Emph{论主显节，或称基督诞辰。}\par

\EditorialNote{本演讲的标题引发了一个疑问：它是在380年12月25日宣讲的，还是在381年1月6日宣讲的？\Quote{主显}一词是主显节的名称，但根据沙夫的说法，原本同时庆祝主的诞生与受洗。这两个词似乎都被用在最简单意义上的\Quote{天主的显现}，而且确实被用于圣诞节。例如，苏伊达斯说：\Quote{主显节是救主的降生成人}；埃皮法尼乌斯（《异端》53）说：\Quote{主显节的日子就是基督按肉体诞生的日子}。但圣热罗尼莫将这个词用于基督的受洗：\Quote{主显节的日子仍然可敬，并非如一些人所以为的，是因为他在肉体内的诞生；因为那时他是隐藏的，而非显现的；而是与\InnerQuote{这是我的爱子}这话所说的时刻相符}（《厄则克耳注释》卷一）。还有一篇归属于圣金口若望的讲道《论基督受洗》，其中明确否认\Quote{主显}一词适用于圣诞节。然而，本演讲本身有证据表明主的诞生节是在更早的日期举行的；因为在第16章中，讲道者说：\Quote{稍后你们会看见耶稣为了我的净化工于约旦河中接受净化}。另一项证据见于《论圣光》演讲第14章：\Quote{在他诞生时，我们适当地举行了庆节，我是庆典的带领者，你们也是。现在我们来到基督的另一个行动和另一个奥迹。}}\par

一、基督诞生了，请光荣他。基督从天而来，请前去迎接他。基督在地上，请你们高升。\Quote{普世大地，请向上主歌唱}；为使我把两者合而为一：\Quote{愿诸天欢乐，愿大地踊跃}，为了那来自天上、然后来自地上的那一位。基督在肉身中，请以颤抖和喜乐欢庆；颤抖是因为你们的罪过，喜乐是因为你们的希望。基督出自童贞女；哦，已婚的妇女们，请你们像贞女一样生活，好成为基督的母亲。谁不朝拜那从起初就存在的？谁不光荣那在末后的？\par

二、黑暗再次过去；光明再次被造；埃及再次被黑暗惩罚；以色列再次被云柱光照。\BibleRef{出 14:20}\BibleQuote{那坐在无知黑暗中的百姓，要看见全知的大光明。}\BibleRef{依 9:6}\BibleQuote{旧的已成过去，请看，一切都成了新的。}\BibleRef{格前 5:17}文字让位，圣神来到前面。阴影消逝，真理进入它们当中。默基瑟德被终结。那无母的成为无父的（在他的先前状态中无母，在他的第二次状态中无父）。自然的法则被颠覆；上界必须被充满。基督命令它，我们不要反对他。\Quote{万民哪，你们要鼓掌欢呼}，因为\BibleQuote{有一个婴孩为我们诞生了，有一个儿子赐给了我们；他肩负着统治权（因为借着十字架它被高举），他的名字被称为\BibleQuote{伟大学者的天使}。}\BibleRef{依 9:6}让若翰呼喊：\BibleQuote{预备上主的道路；}\BibleRef{玛 3:3}我也要呼喊这一天的大能。那非血肉的成了血肉；天主子成为人子，耶稣基督\BibleQuote{昨天、今天、直到永远，常是一样。}\BibleRef{希 13:8}\BibleQuote{让犹太人恼怒，让希腊人嗤笑；}\BibleRef{格前 1:23}让异端者说到舌头痛。那时他们才会相信，当他们看见他升天时；若不是那时，那么当他们看见他从天而降、坐着审判时。\par

三、这些将在以后谈论；当前的庆节是\OriginalTerm{主显节}{Theophany}或\OriginalTerm{诞辰}{Birthday}，因为二者都指同一件事，有两个称谓。因为天主借着诞生向人显现。一方面，他是自有永有的，出自永恒者，超越原因和言语，因为在\OriginalTerm{圣言}{Word}之前没有言语；另一方面，为了我们他也成了受造的，为使我们获得存在的那一位也使我们获得美好的存在，或者更确切地说，当我们因邪恶而从美好存在堕落时，他借着\OriginalTerm{降生成人}{Incarnation}恢复我们。\Quote{主显节}这个名称是就显现而言，\Quote{诞辰}是就他的诞生而言。\par

四、这是我们当前的庆节；我们今天庆祝的正是这一件事：天主降临到人那里，为使我们能出去，或者更准确地说（因为这是更恰当的表达），为使我们能回归天主——使我们脱去旧人，穿上新人；并且\BibleQuote{如同我们在亚当内死了，也照样在基督内活着，}\BibleRef{格前 15:22}与基督同生、同钉、同葬、同复活。\BibleRef{哥 2:11}因为我必须经历这美好的转变，正如痛苦跟随更幸福的事，同样更幸福的事也必须从痛苦中产生。因为\BibleQuote{罪恶在哪里越多，恩宠在那里也格外丰富；}\BibleRef{罗 5:20}如果一个禁果定了我们的罪，那么基督的苦难岂不更要使我们成义？因此，让我们守节，不是按照外邦节期的方式，而是按照虔诚的方式；不是按照世俗的方式，而是以超越世俗的方式；不是作为我们自己的，而是作为属于他那属于我们的，或者更确切地说，作为我们主人的；不是出于软弱，而是出于医治；不是出于创造，而是出于再造。\par

五、这如何能成就？让我们不要装饰门廊，不要安排舞蹈，不要装饰街道；让我们不要宴请眼目，不要用音乐迷惑耳朵，不要用香水使鼻孔倦怠，不要败坏口味，不要放纵触觉——这些是易于犯罪和进入罪恶的道路；让我们不要在柔软飘逸的衣着上柔靡（其美在于无用），也不要用宝石的闪烁或金子的光泽\BibleRef{罗 13:13}或颜色的诡计（这些诡计歪曲自然之美，为侮辱天主的形象而发明）；\BibleQuote{不可在宴乐和醉酒中}，我深知，与这些混杂的是床笫和放荡，因为恶师所教的是恶；或者不如说，无用种子的收成是无用的。让我们不要设立高高的树叶床，为肚腹搭建属于荒淫的帐棚。让我们不要品评酒的芳香、厨师的精致、香膏的巨大花费。不要让海和地带给我们贵重的粪土（这是我学会评价奢侈的方式）；让我们不要在放纵上互相争胜（因为在我看来，一切多余都是放纵，一切超出绝对需要的）——而同时别人却在饥饿和匮乏中，他们是用同样的泥土、以同样的方式造成的。\par

六、让我们把这一切留给希腊人以及希腊人的排场和节期吧，他们以神的名号称呼那些喜爱祭献烟气、并且始终用肚腹来朝拜的存有；他们是恶的发明者和恶神的朝拜者。但我们，我们朝拜的对象是\OriginalTerm{圣言}{the Word}，如果必须要在某种程度上享受奢侈，就让我们在言语中、在神圣法律中、在历史中寻求奢侈；尤其是那些属于这庆节起源的历史；好使我们的奢侈与那召集我们的那一位相近，而不远离。或者你们愿意（因为今天我款待你们），让我尽可能丰富而高贵地将这些事的故事摆在我善良的客人面前，好使你们知道一个外邦人如何能款待本地人，一个乡下人如何能款待城里人，一个不关心奢侈的人如何能款待喜欢奢侈的人，一个贫穷无家之人如何能款待那些因财富而显赫的人？\par

我们开始于此点；我请求你们这些喜爱此类事的人洁净你们的心思、耳朵和思念，因为我们的言论将关乎天主和神圣之事；当你们离去时，你们可以享受那永不消逝的喜乐。而这同一篇言论将同时非常完备又非常简洁，使你们既不因它的不足而不悦，也不因饱足而感到不快。\par

七. 天主总是曾在，总是现在在，总是将在。或者不如说，天主总是\Quote{是}。因为\Quote{曾是}和\Quote{将是}是我们时间的片段，属于变化的本性，而祂是永恒的存在。这是祂在山上向\Person{梅瑟}颁布训谕时，给自己所起的名字。因为祂在自身内总括并包含一切存在，既无过去的开始，也无未来的终结；犹如一片无限无际的伟大存在之海，超越一切时间和本性的概念，仅被心灵隐约而微弱地摹绘……不是以祂的本质，而是以祂的环境；一个形象由这一源头得来，另一个由那一源头得来，组合成某种真理的呈现，但我们在抓住它之前就已失去它，在构思它之前就已逃逸，当我们的主导部分即便被洁净时，它如闪电般闪耀其上，那闪电不停留其行程，之于我们的视线也是如此……按我的理解，为的是通过我们可以理解的那部分吸引我们归向祂（因为完全不可理解的事物超出了希望的范围，不在努力的能力之内），并通过我们不能理解的那部分激发我们的惊奇，因惊奇而更加渴望，因渴望而得以净化，因净化而使我们相似天主；\BibleRef{若 10:15} 如此，当我们变得相似祂时，天主——容我直言——将与我们交往如神，与我们合一，或许正如祂已经认识那些认识祂的人一样程度。神圣的本性既是无限而难以理解的；我们对祂所能理解的只是祂的无限性；即使有人可能认为，因为祂是单纯的本性，所以祂要么完全不可理解，要么完全可以理解。因为让我们进一步探究\Quote{是单纯的本性}的含义。因为这单纯本身肯定不是它的本性，正如复合本身不是复合物的本质一样。\par

VIII. 当无限从两个观点被考虑时，即开始和终结（因为超越它们且不被它们限制的就是无限），当心灵向上看深处，没有立足之地，并依靠现象来形成对天主的概念时，它将那里所发现的无限不可接近者称为\Quote{无起源者}。当它向下看深处，看未来时，它称祂为\Quote{不死不朽者}。当它从整体得出结论时，它称祂为\OriginalTerm{永恒者}{\GreekText{αἴωνιος}}。因为\OriginalTerm{永恒}{\GreekText{αἵων}}既非时间，也非时间的部分；因为它不能被度量。但时间（由太阳的行程所度量）之于我们，正如永恒之于永生者，即一种类似时间的运动与间隔，与其存在同延。然而，我关于天主只能说到这里；因为现在不是适当的时候，我现在的主题不是关于天主的道理，而是关于降生成人的道理。但当我提到天主时，我指的是圣父、圣子和圣神。因为天主性既未扩散超越这些，而引入一群众神；也未局限在比这些更小的范围内，以致谴责我们对天主性有贫乏的概念；要么为保存\OriginalTerm{单权}{Monarchia}而犹太化，要么因众神之多而堕入异教。因为两边的邪恶相同，虽然方向相反。这便是至圣所，连色辣芬也隐藏，并以三重\Quote{圣}颂来光荣，归于主和天主的一个称号，正如我们的一位前人以最美妙、最高超的方式所指出的。\par

IX. 但由于这种自我默观的运动本身不能满足善，而善必须被倾注并超越自身以增多其恩惠的对象，因为这对最高的善是必不可少的，祂首先构思了天上的和天使的能力。这个构思由祂的圣言完成，并由祂的圣神成全。于是，次要的光辉出现了，作为首要光辉的仆役；我们是将它们构思为有灵的灵体，还是非物质的、不朽的火，或某种尽可能接近这个的其他本性。我本想说，它们不能向邪恶移动，而只能向善移动，因为它们在主身边，接受来自主的第一光线——因为地上的存有只有第二光照；但我不能那样说，而只能构思并谈论它们只是难以移动，因为那一位，因光辉被称为\Person{路济弗尔}，但因其骄傲变成了并被称为黑暗；以及服从他的背叛的众军，通过他们对善的背叛而成为邪恶的创造者和我们的挑唆者。\par

X. 如此，基于这些缘由，祂便赋予理智世界以存在，就我所能推究的这些事，并以我粗浅的言辞估量伟大事物而言。当祂的初次创造井然有序时，祂构思了第二个世界，物质而可见的；这是天地及其间万物的体系与组合——当我们观看每个部分优美的形式时，确是令人赞叹的创造，但当我们思量整体的和谐与一致，以及每个部分如何以美好的秩序与其它部分相配，且全部趋向于作为统一体的世界的完美完成时，则更值得赞叹。这为要表明：祂不仅能召唤一种与自己相似的本性进入存在，也能召唤一种完全相异于自己的本性。因为与神性相近的是那些理智的、唯有心智才能领悟的本性；而感官所能察觉的一切则完全相异于祂；其中最远隔的是那些完全缺乏灵魂和运动能力者。但也许有些过于欢庆和冲动的人会说：\Quote{这一切与我们何干？赶你的马到终点吧。向我们讲讲这庆节，以及我们今天聚集在此的理由。}是的，这正是我将要做的，尽管我从稍远的地方开始，被爱和我的论证所需所迫。\par

XI. 如此，心智和感官，彼此这样区分，已停留在各自的界限内，并在自身中携带创造者圣言的宏伟，作为祂大能之工的无声赞美者和激动人心的传报者。还没有两者的混合，也没有这些对立物的混合——那是更伟大的智慧和慷慨在创造本性中的标记；善的丰富也尚未完全显明。如今，创造者圣言决意显明这一点，并从两者——即可见与不可见的受造物——中产生一个单一的活物，便塑造了人；从已有的物质中取得一个身体，并在其中放置一个从祂自己取来的气息\BibleRef{创 2:7}，这气息，圣言知道，是理智灵魂和天主的肖像，犹如一个第二世界。祂将他安置于地上，伟大而渺小；一个新的天使，一个混合的敬拜者，完全被引入可见的创造，但仅部分进入理智的；地上万物的君王，却服从于天上的君王；属地的又属天的；暂时的又不朽的；可见的又理智的；介于伟大与卑贱之间；在一个位格中结合了精神和血肉：精神，因所受的宠遇；血肉，因被提升的高度；前者使他能继续存活并赞美他的恩主，后者使他能受苦，并因受苦而被提醒，若因伟大而骄傲则受纠正。一个在此受训练、然后迁往别处的活物；并且，为完成这奥迹，因倾向天主而被神化。因为，我想，我们在此有限拥有的真理之光，正是趋向于此：使我们既能看见也能体验天主的辉煌，这辉煌与那造了我们并要以更高方式重塑我们的祂相称。\par

XII. 祂将这存在置于乐园中——无论那乐园可能为何——以自由意志的恩赐尊荣了他（为使天主不仅因祂播下善的种子而属于他，也因他自己的选择而属于他），使他耕种不朽的植物，这或许意指神圣的意念，既有较简单的也有更完美的；赤身露体，生活简朴而毫无技巧，没有任何遮盖或遮蔽；因为由起初就如此的人应当是这般。又赐给他一条法律，作为他的自由意志可以行动的材料。这法律是一条诫命，关于他可吃哪些植物，以及哪一棵他不可触摸。后者是知善恶树；然而，并非因为它从起初被种植时就是邪恶的；也不是因为天主嫉妒我们而禁止它……愿天主的仇敌不要在那里摇动舌头，或仿效那蛇……但若在适当的时候食用，它本是好的，因为按照我的理论，这树乃是默观，只有那些达到了习惯的成熟的人才能安全地进入；但对于那些在习惯上仍有些单纯和贪婪的人，则是不好的；正如固体食物对于尚柔弱、需要奶水的人是不好的。\BibleRef{希 5:12} 但是，因着魔鬼的恶意和女人的任性——她作为较柔弱者屈服于此，并将它加于男人，因她更善于说服——唉，我的软弱！（因为我初祖的软弱就是我的软弱），他忘记了赐给他的诫命；\BibleRef{创 3:5} 他屈服于有害的果实；因他的罪，他立即被逐出生命树、乐园和天主；并穿上了皮衣……那或许是指更粗糙的血肉，既有死又有矛盾。他首先学到的是自己的羞耻；\BibleRef{罗 1:22} 并躲避天主。然而，在这里他也获得了益处，即死亡和罪的消除，为使邪恶不至于不朽。这样，他的惩罚变成了仁慈的；因为我深信，天主施加惩罚是出于仁慈。\par

十三. 既然他（人类）首先藉多种方式被惩戒（因为他的罪恶甚多，罪恶的根源通过不同缘由、在不同时期滋生），即藉言辞、法律、先知、恩惠、威吓、灾祸、洪水、烈火、战争、胜利、失败、天上、空中、地上和海中的征兆，藉人、城、邦国出乎意料的变迁（其目的在于摧毁邪恶），最后他需要一种更强的治剂，因为他的疾病日趋恶化：彼此残杀、奸淫、背誓、逆性之罪，以及首要和最终的万恶——偶像崇拜和对受造物的崇拜取代对造物主的崇拜。既然这些需要更大的帮助，它们也就得到了更大的帮助。那就是：\OriginalTerm{天主的圣言}{Word of God}自身——祂在万世之前，无形无象，不可理解，无体，\OriginalTerm{始之始}{Beginning of Beginning}，\OriginalTerm{光明之光}{Light of Light}，\OriginalTerm{生命与不死之源}{Source of Life and Immortality}，\OriginalTerm{原始美的肖像}{Image of the Archetypal Beauty}，\OriginalTerm{不可移动的印记}{immovable Seal}，\OriginalTerm{不变的形象}{unchangeable Image}，\OriginalTerm{圣父的界限与圣言}{Father's Definition and Word}——降临到祂自己的肖像，并因我们的肉体而取了肉体，为我的灵魂而与理性灵魂结合，以同类净化同类；除了无罪以外，完全成了人。由童贞女怀孕（\BibleRef{路 1:35}），她首先在灵魂和肉体上被圣神圣化（因为必须尊崇生育，并使童贞获得更高的尊荣），然后祂以天主的身分，带着所取的人性而出，\OriginalTerm{一位二性}{One Person in two Natures}，即肉体和精神，后者使前者神化。哦，新的混合；哦，奇特的结合；\OriginalTerm{自有者}{Self-Existent}成为受造，\OriginalTerm{非受造者}{Uncreate}被造，不可容纳者被容纳，藉着理性灵魂的中介，这灵魂在神性与肉体的有形体质之间调和。那使人富足的，成为贫穷，因为祂承袭了我肉身的贫穷，为使我能承袭祂天主性的富足。那充满者倒空自己，因为祂暂时倒空了自己的荣耀，为使我能有分于祂的圆满。祂良善的富足是什么？这环绕着我的奥秘是什么？我曾有分于肖像，却没有保全它；祂有分于我的肉身，为能既拯救肖像，又使肉身不死。祂进行了第二次交往，远比第一次奇妙，因为第一次祂分施了更优越的本性，如今祂亲自承受了更劣的本性。这比前一个行动更神圣，在一切明智者眼中更为崇高。\par

十四. 对这些，那些吹毛求疵者——那些对神性尖酸推理者、一切可赞美之事的诽谤者、光明的黑暗者、在智慧上未受教养者、基督为他们白白死去之人、那些忘恩负义者、恶者的作品——有什么可说？你是否将这恩惠视为对天主的侮辱？你竟因此认为祂渺小，因为祂为你而自卑，因为\OriginalTerm{善牧}{Good Shepherd}（\BibleRef{若 10:11}），祂为羊舍命，来寻找那曾在山上和山丘间迷失的（你在那里献祭），寻得了迷途者；既寻得（\BibleRef{路 15:4}），就扛在肩上——祂也曾以此肩负\OriginalTerm{十字架的木头}{Wood of the Cross}——然后带它回到更高的生命；既带回，就将其列在从未迷失者之中。因为祂点了一盏灯——祂自己的肉身——并打扫房屋，从罪恶中洁净世界；寻找那枚钱币，即被情欲掩盖的\OriginalTerm{王权肖像}{Royal Image}。祂在找到钱币时，召集天使朋友们，使他们分享祂的喜乐——祂曾使他们也分享降生成人的奥秘。因为\OriginalTerm{先导者的蜡烛}{candle of the Forerunner}之后，跟随那更明亮的光；话语之后，圣言接续；\OriginalTerm{新郎的朋友}{Bridegroom's friend}之后，\OriginalTerm{新郎}{Bridegroom}接续；那为天主预备一个特选的子民、以水在圣神面前洁净他们的人？你难道为这一切而责备天主？你因此认为祂被贬低，因为祂束上毛巾，洗门徒的脚，表明谦卑是走向高举的最佳道路？因为对于那俯伏于地的灵魂，祂谦卑自己，为能同自己一起举起那因罪恶重负而摇摇欲坠的灵魂？你为何不也指控祂与税吏同席（\BibleRef{路 5:29}），并收税吏为徒，为使自己也有所收获……那收获是什么？——罪人的救恩。若如此，我们必须责备医生俯身于病痛，忍受恶臭以治愈病人；或责备那按法律命令弯腰入沟以拯救坠入其中的牲畜的人。\par

十五、祂被派遣，但作为人，因为祂具有双重本性；因为祂按可朽之体的本性，感到疲倦、饥饿、口渴、忧苦，并流下眼泪。即使这个说法也用于天主，意思是圣父的善意应被视为一种派遣，因为祂将自己一切所是都归于此；既为尊崇无始之源，也为避免被视为敌对的天主。经上记载，祂被出卖，也自己交出自己；祂被圣父举起，被接到天上；另一方面，祂自己复活并\Emph{升天}；每一对中的前者指圣父的善意，后者指祂自己的权能。那么，你是否被允许只强调祂屈辱的事，而忽略所有高举祂的事，并把你那一方算作祂所受的苦，却不考虑这是出于祂的旨意？请看，如今圣言还要受什么苦：有人尊祂为天主，却将祂与圣父混为一谈；有人辱祂为纯血肉，将祂与天主性分离。祂会更恼怒哪一方？或者，祂会宽恕谁？是那些侮辱性地混淆祂的人，还是那些分裂祂的人？前者本应区分，后者本应联合；一者在于数目，一者在于天主性。你因祂的肉体而跌倒吗？犹太人也如此。你称祂为撒玛黎雅人，并且……我不说其余的话。你不信祂的天主性吗？连魔鬼也不这样，哦，你比魔鬼更不信，比犹太人更愚钝。那些（魔鬼）知道\Quote{圣子}的名号意味着同等尊荣；这些（犹太人）知道驱逐他们的那位是天主，因为他们由自身的经验而确信。但你既不承认同等尊荣，也不承认天主性。你倒不如是个犹太人或附魔的人（若我可以说一句荒谬的话），总好过未受割损、身体康健，却存有如此邪恶和不虔敬的心态。\par

十六、稍后，你将看见耶稣为我的净化而在约旦河受净洗，更确切地说，是藉祂的净化圣化诸水（因为祂无需净化，除去世罪的祂），以及天裂开，与祂同性同体的圣神为祂作证；你将看见祂受试探并得胜，有天使服侍祂，治愈各种疾病和各种灾殃，使死人复活（但愿祂也使你们因异端而死的人复活），驱逐魔鬼，有时亲自，有时藉祂的门徒；以少量饼饱食众多群众；在海面上干足行走；被出卖并被钉十字架，并将我的罪与祂一同钉在十字架上；作为羔羊被奉献，又作为司祭奉献；作为人埋葬于坟墓，又作为天主复活；然后升天，并将在自己的荣耀中再来。看，基督的每一奥迹中有多少隆重的庆节！所有这些只有一个圆满，就是我的成全和返回亚当最初的状态。\par

十七、现在，我请你接受祂的孕娠，并在祂面前踊跃；若不似若翰在母胎中\BibleRef{路 1:41}，至少也要似达味因约柜安息而踊跃\BibleRef{撒下 6:14}。敬畏那册籍，因你被记在天上；朝拜那诞生，因你从出生的锁链中被释放\BibleRef{路 2:1}；尊敬小城白冷，它引你返回乐园；崇拜马槽，你虽无理性，却藉此被圣言喂养。如同依撒意亚所嘱咐的，像牛认识你的主人，像驴认识你主人的槽；如果你属于那洁净可作合法食物、并反刍圣言之食、堪当祭献的人。如果你属于尚不洁净、不可食用、不堪祭献的人，属于外邦人的部分，就同明星奔跑，与贤士们同献礼物：黄金、乳香、没药\BibleRef{依 1:3}，如同献给君王、天主、为你而死的那一位。与牧童一同光荣祂\EditorialNote{玛窦福音第二章}；与天使一同欢唱；与总领天使一同歌颂。愿这庆节成为天上和地上万军共有的\BibleRef{路 2:14}。因为我深信今天天上的众军与我们一同喜乐，同庆这大节……因为他们爱人，也爱天主，正如达味在祂受难后所介绍的，那些与基督一同升天、前来迎接祂、彼此呼唤要举起城门的人。\par

十八. 关于基督的诞生，有一件事我要你们憎恨...黑落德对婴儿的屠杀。\BibleRef{玛 2:16} 或者不如说，你们也应崇敬这事，即与基督同年龄的牺牲，在新祭献奉献之前被杀。若祂逃往埃及，就欣然作祂流亡的伴侣。与受迫害的基督同流亡，是件大事。若祂在埃及久留，就以在那里对祂的虔敬敬拜，将祂从埃及召出。无过地经历基督生命的每一阶段和官能。被净化；受割损；除去从出生就遮盖你的面纱。之后在圣殿里教导，驱逐亵渎的商贩。\BibleRef{若 2:15} 若有必要，忍受被石头砸，因我很清楚你将从投石者那里被隐藏；你甚至要从他们当中逃脱，如同天主。若你被带到黑落德面前，大部分时间不要回答。\BibleRef{路 23:9} 他会尊重你的沉默，胜过多数人的长篇大论。若你受鞭打，求他们留下的部分。\BibleRef{若 19:1} 为尝味道而尝苦胆；\BibleRef{玛 27:34} 喝醋；\BibleRef{若 19:29} 寻求唾污；接受击打，戴上茨冠，即虔诚生活的坚硬；穿上紫袍，手持芦苇，接受那些嘲弄真理者的戏弄崇拜；最后，与祂同钉十字架，并欣然分担祂的死亡与埋葬，为使你与祂一同复活、一同受光荣、一同为王。注视并受那伟大的天主注视，祂在天主圣三中被敬拜和光荣，并且我们宣告，祂如今在我们血肉的束缚所允许的范围内，清楚地呈现在你们面前，在我们的主耶稣基督内，愿光荣归于祂，直到永远。阿们。\par

\SourceRecord{Translated by Charles Gordon Browne and James Edward Swallow.；Nicene and Post-Nicene Fathers, Second Series；Vol. 7.；Edited by Philip Schaff and Henry Wace.；Buffalo, NY: Christian Literature Publishing Co.,；1894.；<http://www.newadvent.org/fathers/310238.htm>.}

% structure: p0001 source=h1 target=section reason=page-title

\section{第三十九讲道辞}

\Emph{\Quote{论圣光讲道辞}.}\par

\Emph{\Quote{论圣光讲道辞}于381年主显节宣讲，次日继之以论圣洗讲道辞。在东方教会中，此节日更特别视为纪念我主基督圣洗之日，因此是隆重施行圣事的重要日子之一。此节日通常称为Theophania，礼仪中的福音是\BibleRef{玛 3:13-17}。八日庆期内的主日称为\OriginalTerm{圣光后}{\GreekText{μετὰ} \GreekText{τὰ} \GreekText{φῶτα}}，指明当时此庆日被称为\Quote{圣光}，圣额我略时代似乎即是如此。此名源于圣洗圣事，在古代常称为\Quote{光照}，新教友手持点燃的火炬或蜡烛，即与此名（源于圣事的神性恩宠）有关。此庆日的庆典似乎持续两天，第二天用于隆重授予圣事。因此，我们发现有属于此庆日的两篇讲道辞。第一篇在当天宣讲，他特别论述了其中所纪念的庆日和我主基督圣洗的奥迹；然后他谈到不同种类的圣洗，列举了五种，即：—}\par

\begin{itemize}
\item 以色列藉梅瑟在云中和海中领受的预像性圣洗。
\end{itemize}

\begin{itemize}
\item 由圣若翰洗者施行的预备性悔改圣洗。
\end{itemize}

\begin{itemize}
\item 由我主赐予我们的水和圣神的精神性圣洗。
\end{itemize}

\begin{itemize}
\item 光荣的殉道圣洗。
\end{itemize}

\begin{itemize}
\item 痛苦的补赎圣洗。
\end{itemize}

在谈到最后一种时，他借机驳斥了诺瓦提安追随者的极端严格主义，他们否认对某种领洗后所犯的罪可予赦免。\par

在次日宣讲的第二篇讲道辞中，他论述了圣洗圣事及其精神效果；并借机斥责当时仍然盛行的将圣洗推迟至濒死之际的做法。他还论述了圣事的有效性及精神效果完全独立于施行圣事的司铎的品级或功德；最后，他概述了领受圣事所涉及的责任，并附有对信经以及施行圣事时所伴随礼仪的极为宝贵的阐述。\par

\OriginalTerm{一、}{I.}再次，我的耶稣，再次，一个奥迹；不是欺骗性的，也不是混乱的，不属于希腊的错误或醉态（因为我如此称呼他们的庆典，我想每个有健全理智的人也如此认为）；而是一个崇高而神圣的奥迹，与上天的荣耀相连。因为我们所来到并正在庆祝的圣光节，其根源在于我的基督——那真光，来到世界的每个人所被照亮的——的圣洗，\BibleRef{若 1:9}并实现了我的净化，并助益那我们从起初从上主那里领受的光，但此光被我们的罪所昏暗和混淆。\par

\OriginalTerm{二、}{II.}因此，请听天主的声音，这声音对我——既是这些奥迹的门徒又是导师——显得如此清晰，唯愿天主也如此向你们发声；\BibleQuote{我是世界的光}。\BibleRef{若 8:12}因此，你们当接近祂并受光照，愿你们的脸面不蒙羞，被真光所印记。这是新生的时节，\BibleRef{若 3:3}让我们重生。这是更新的时节，让我们再次接纳第一个亚当。让我们不要停留在现状，而是成为我们曾经的模样。光在黑暗中照耀，在这生命和肉身中，被黑暗追逐，却不被它追上——我指的是那敌对力量，无耻地扑向可见的亚当，却遇见了天主而被击败——这样，我们弃绝黑暗，接近那光，然后成为完美的光，完美之光的子女。看这日子的恩宠，看这奥迹的能力。你们不是被提升离地吗？你们不是明显被高举，因我们的声音和默想而超升吗？当圣言兴盛我言语的进程时，你们将被举得更高。\par

\OriginalTerm{三、}{III.}律法的影儿般的洁净中有任何这样的事吗？它以暂时的洒水和母牛的灰烬洒在不洁者身上；或者外邦人在他们的奥迹中庆祝这样的事吗？他们的所有仪式和奥迹，在我看来都是无稽之谈，是恶魔的黑暗发明，是不幸心灵的虚构，借时间助长，被寓言隐藏。因为他们当作真理崇拜的，却作为神话掩盖。但如果这些事是真的，就不该被称为神话，而应证明并非可耻；\BibleRef{希 7:13}而如果是假的，就不该被当作奇事；人不该如此轻率地对同一件事持有最相反的意见，如同在市场与孩童或真正邪恶之人游戏，而非与有理智的人、崇拜圣言的人讨论，尽管他们藐视这种做作的似是而非。\par

IV. 在这些奥迹中，我们所关心的并非\Person{宙斯}的诞生和克里特暴君的偷窃（尽管希腊人不喜欢这样的称号），也不是\Person{库瑞忒斯}之名和武装的舞蹈，那是用来掩盖一位哭泣神祇的哀号，好让他逃脱父亲的仇恨。因为被当作石头吞下的人竟然像孩童一样哭泣，这实在奇怪。我们也不涉及弗利吉亚的阉割、笛子和\Person{科律班忒斯}，以及人们关于\Person{瑞亚}的一切狂乱，他们将人奉献给众神之母，并接受适合这样一位神祇之母的仪式。我们也没有少女的被掳，\Person{得墨忒耳}的漂泊，她与\Person{刻勒俄斯}、\Person{特里普托勒摩斯}和龙的亲密关系，以及她的所作所为和所受的苦……因为我羞于将那黑夜的仪式暴露在光天化日之下，并将污秽之事奉为神圣奥秘。\Place{厄琉息斯}知道这些，那些亲眼目睹那里被沉默保守、并且理应如此之事的人也知道。我们的纪念也不是关于\Person{狄俄尼索斯}，以及那因不完全生育而受苦的大腿，正如之前另一个头颅曾因生育而受苦；也不是关于阴阳神，或一群醉醺醺、萎靡不振的合唱队，或崇敬他的底比斯人的愚行，或他们所崇拜的\Person{塞墨勒}的霹雳。也不是\Person{阿佛洛狄忒}的娼妓奥秘，据他们承认，她出身卑微且受辱；这里没有\OriginalTerm{法罗斯}{Phalli}和\OriginalTerm{伊提法罗斯}{Ithyphalli}，在形状和行动上都是可耻的；没有陶里斯对外乡人的屠杀；也没有拉科尼亚青年的血洒在祭坛上，他们用鞭子鞭打自己；唯独在这种情况下，他们错误地运用了勇气，竟崇拜一位女神，且是一位贞女。因为这些人既崇尚柔弱，又崇拜鲁莽。\par

V. 你将把\Person{珀罗普斯}的屠杀置于何地？那屠杀款待了饥饿的众神，是痛苦且不人道的款待。\Person{赫卡忒}可怕阴暗的幽灵、地下幼稚的把戏和\Person{特洛福尼俄斯}的巫术，或\Place{多多那橡树}的喃喃低语，或\Place{德尔斐}三脚架的诡计，或\Place{卡斯塔利亚}的预言之饮——它能预言一切，唯独不能预言自己将归于沉寂。还有\Person{玛吉}的祭献之术和他们的内脏预兆，或\Person{迦勒底}的天文学和星象，将我们的生命与天体的运行相比——而天体连自己是什么或将是什么都不知道。也不是那些色雷斯秘仪，据说\Quote{\OriginalTerm{敬拜}{\GreekText{θρησκεία}}}一词源于此；也不是\Person{俄耳甫斯}的仪式和奥秘——希腊人如此钦佩他的智慧，以至于为他设想了一把能以音乐吸引万物的竖琴。也不是\Person{密特拉}的折磨——那些能忍受接受这种秘仪的人理应遭受的折磨；也不是\Person{俄西里斯}的肢解，那是埃及人敬重的另一场灾难；也不是\Person{伊西斯}的不幸，以及比\Person{门德斯人}更可敬的山羊，和\Person{阿庇斯}的畜栏——那头在\Person{孟斐斯人}的愚昧中自我放纵的牛犊；也不是他们用以凌辱\Place{尼罗河}的一切尊荣，同时又颂扬它为果实和谷物的赐予者，以及以它的腕尺来衡量幸福的尺度。\par

VI. 我略过他们对爬行动物的尊崇，以及对卑贱之物的崇拜——每一种都有其独特的祭仪和节庆，且都共同分享一种相同的魔性。因此，如果他们绝对必须不虔敬、背离尊崇天主、被引向偶像和工艺作品及人手所造之物，有理智的人也无法为自己祈愿更糟的事，只能祈愿他们正好崇拜这类东西、并以这种方式尊崇它们。这样，就如\Person{保禄}所说，他们就在自己身上受到所应得的\Quote{错谬的报应}(\BibleRef{罗 1:27})，这报应就在他们崇拜的对象本身；他们并非尊崇它们，而是因它们受辱；因错谬而可憎，又因他们所敬拜和尊崇之物的卑贱而更为可憎；以至于他们甚至比他们所崇拜的对象更无理智，因其愚蠢的程度恰如后者的卑劣。\par

VII. 好吧，让这些事成为希腊人及其神祇——他们的愚昧所归因者——的消遣吧；这些神祇将天主的光荣转向自己，以可耻的思想和幻想从各方面分裂人类，自从他们藉着那棵不合时宜、不恰当传递给我们的知识之树，将我们赶离生命之树以后，便趁我们现在比以前更软弱而攻击我们；彻底夺去了我们内中掌权的理智，为情欲打开了门。因为他们天性嫉妒且憎恨人类，或者不如说因自己的邪恶而变成这样，他们既不能容忍我们这些在下者上升到那在上者那里——他们自己已从上坠落到了地上；也不容忍他们的光荣和最初本性的这种转变已经发生。这就是他们迫害受造物的意义。为此，天主的肖像受了凌辱；因为我们不喜欢遵守诫命，我们就任凭自己错谬的独立自主。当我们错谬时，我们就因所崇拜的对象而蒙羞。因为这不仅是一场灾祸——我们被造为善工，以光荣和赞美我们的造物主，并尽可能效法天主，却成了一切情欲的巢穴，残忍地吞噬和消耗那内在的人——而且还有进一步的邪恶：人竟然将自己的情欲奉为神祇的辩护者，这样，罪就不仅可以被认为无须负责，甚至被视为神圣的，以所崇拜的对象作为托辞。\par

八、然而，既然我们蒙受恩宠，得以逃离迷信的错误，与真理结合，事奉永生而真实的天主，超越了受造界，越过一切属于时间和原动的事物；就让我们以相称于所赐恩宠的方式，来注视并推论天主和神圣的事物。让我们从最合适的起点开始讨论。最合适的起点，正如撒罗满为我们所指定的：他说，\BibleQuote{智慧的开端，是获取智慧}（\BibleRef{箴 4:7}）。他又告诉我们，智慧的开端是敬畏。因为我们必须一开始就以默观为始，却以敬畏为终（因为不受约束的默观也许会推我们坠入深渊），而应以敬畏为基础，得到净化，可说是因敬畏而变轻，从而被提升到高处。因为哪里有敬畏，哪里就有遵守诫命；哪里有遵守诫命，哪里就有肉身的净化——那层遮蔽灵魂、不让它看见神光的云雾。哪里有净化，哪里就有光照；光照乃是对于那些渴慕最伟大事物者，或那最伟大者，或那超越一切伟大的事物者，其渴望的满足。\par

九、因此，我们必须先净化自己，然后才接近与那纯洁者的交往；否则，我们就会遭遇与以色列同样的经历（\BibleRef{出 34:30}），他们不能忍受梅瑟脸上的荣耀，因而要求一个帕子（\BibleRef{格后 3:7}）；或者，与玛诺亚一同感受并说：\Quote{我们必定要死，妻子啊，因为我们看见了天主}（\BibleRef{民 13:23}），虽然那只是他幻想中的天主；或者，如同伯多禄将耶稣请出船（\BibleRef{路 5:8}），因为我们自己不配这样的造访；我说伯多禄，是指那曾步行海面的人（\BibleRef{玛 14:29}）；或者，如同保禄眼睛失明（\BibleRef{宗 9:3}），在他还未从他迫害的罪责中得洁净时，与那他所迫害者交谈——更确切地说，只是那大光的一闪；或者，如同百夫长（\BibleRef{玛 8:8}）寻求医治，却因可称赞的敬畏而不肯接待医治者进入他的家。我们每一个人，只要仍未洁净，仍是百夫长，指挥许多恶事，在凯撒的军中服役——那世界统治者，领导那些被拖下的人——也该如此说：\Quote{我不堪当你到我的檐下来}。但当他仰望了耶稣，虽然身材矮小如同古时的匝凯（\BibleRef{路 19:3}），爬上一棵野桑树的顶上，藉着治死他在地上的肢体（\BibleRef{哥 3:5}），并已超越那卑微的身体，那时他必要接待圣言，且有话对他说：\Quote{今天救恩临到了这一家}（\BibleRef{路 19:9}）。然后，让他抓住救恩，结出更完美的果实，把那作为税吏错误聚敛的，正确地进行分散和倾注。\par

十、因为同一圣言，一方面，就其本质而言，对不配的人是可怕的；另一方面，因着它的慈爱，那些如此预备好的人——他们已从灵魂中驱除了不洁和世俗的灵，藉着自我反省扫除并装饰了自己的灵魂，没有让灵魂空闲或无所事事，以免再次被七邪灵以更大的武装占据……七的数目与德行的数目相同（因为最难对抗的敌人要求最严厉的努力）……但除了逃避邪恶，还要实践德行，让基督完全地，或至少尽可能多地，住在他们内——可以接受它，好使恶的力量找不到任何空处再以自己充满它，并使那人的末后状况比先前更坏，因着它攻击的更大力量以及那堡垒的更大强度和不可攻破。但是，当我们以一切谨慎守护了灵魂，在我们心中设立了上升的阶梯，开垦了我们的荒地（\BibleRef{耶 4:3}），播种了正义（\BibleRef{箴 11:18}），正如达味、撒罗满和耶肋米亚所命令我们的，让我们以知识之光光照自己，然后谈论那隐藏在奥秘中（\BibleRef{格后 2:6}）的天主的智慧，并光照他人。同时，让我们净化自己，领受圣言的基本入门，好使我们自己获得最大的益处，使我们成为像天主一样的，并在圣言来临时接待祂；不仅如此，还要紧紧抓住祂并向他人彰显。\par

十一、现在，藉着以上所述净化了剧场，让我们略略谈论这庆节，以喜庆而虔诚的灵魂一同庆祝这庆典。既然庆节的主要点是纪念天主，就让我们追念天主。因为我认为，那些在\Emph{那里}——那是一切有福者居住之处——庆节的人的声音，无非就是所有堪当那城的人所歌唱的天主的赞美诗和颂词。如果我要说的包含一些先前说过的话，没有人会感到惊讶；因为我不仅会重复同样的话语，而且当我为你们谈论天主时，也会同样战战兢兢，无论是舌头、心智还是思想，好使你们也能分享这值得赞美和蒙福的情感。当我谈论天主时，你们必须立即被一次闪光和三次闪光所光照。三个在于个体性或位格——如果有人更愿意这样称呼，或称为\Quote{位}——因为我们不会在名称上争论，只要音节指向同一意义；但就本体而言，是一——即天主性。因为，如果可以这样说，它们是分开而不分离的；它们在分开中联合。因为天主性在三中是一，而三是在一之中，天主性在它们内，或更准确地说，它们就是天主性。我们将略过过度和不足，既不使统一性成为混淆，也不使区分成为分离。我们要同样远离撒伯里乌的混淆和亚略的分割，它们是截然相反的恶，但在邪恶上相等。因为有什么必要异端地将天主融合在一起，或将祂切割成不平等呢？\par

十二、因为对我们来说，只有一位天主，就是圣父，万物都出于祂；也只有一位主耶稣基督，万物都藉着祂；并且只有一位圣神，万物都在祂内；\BibleRef{格后 8:6}然而，这些\Quote{出于}、\Quote{藉着}、\Quote{在……内}的词并不表示本性的差异（因为如果是这样，这三个介词或三个名字的顺序就绝不会改变），而是标示同一不混乱本性的位格。这由以下事实所证明：如果你不草率地阅读同一宗徒的另一段话，它们又被收归为一：\Quote{万物都出于祂、藉着祂、归于祂；愿光荣归于祂，直到永远。阿们。}\BibleRef{罗 11:36}圣父是父，祂是\OriginalTerm{无源}{Unoriginate}，因为祂不从任何人而来；圣子是子，祂不是无源的，因为祂来自圣父。但如果你从时间意义上理解\Quote{源}，那么祂也是无源的，因为祂是时间的创造者，并不受时间约束。圣神确实是神，祂发自圣父，但并非以子的方式，因为祂不是藉着\OriginalTerm{生}{Generation}，而是藉着\OriginalTerm{发}{Procession}（为了清晰起见，我不得不创造这个词）；因为圣父并没有因生了什么而不再是\OriginalTerm{无生者}{Unbegotten}，圣子也没有因出自无生者而不再是\OriginalTerm{受生者}{Begotten}（这怎么可能呢？），圣神也没有因祂发或因祂是天主（尽管不敬神的人不相信）而变成圣父或圣子。因为\OriginalTerm{位格}{Personality}是不变的；否则，位格如何能保持不变，如果它可改变并能从一个转移到另一个？但那些把\Quote{无生者}和\Quote{受生者}当作模棱两可的神的本性的人，或许会认为\Person{亚当}和\Person{塞特}在本性上不同，因为前者不是由血肉所生（他是被造的），而后者是由\Person{亚当}和\Person{厄娃}所生。所以，正如我们所说，在三位中只有一位天主，而这三位是一体。\par

十三、既然这些事是如此，或者说，既然祂是如此；而且对祂的朝拜不应只由天上的存在者献上，也应有地上的朝拜者，为使万物充满天主的荣耀（因为它们已充满天主本身）；因此，人受造并被尊以天主的双手和肖像。但因魔鬼的嫉妒和罪恶的苦味，人可悲地与造物主天主分离——轻视人，这并非天主的本性。那么，做了什么，那关乎我们的伟大奥迹是什么？本性被革新，天主成了人。\Quote{祂乘驾于诸天之上，在东方的}自己的荣耀与威严，在我们卑微与低下的西方被显扬。天主之子屈尊成为并被称为人子；并非改变祂所是的（因为祂是不变的），而是取了祂所不是的（因为祂满爱世人），为使不可领悟者得以被领悟，藉着肉体作为帷幕与我们交谈；因为那有生死朽坏的本性无法承受祂无遮蔽的天主性。因此，\OriginalTerm{不混杂者}{Unmingled}被\OriginalTerm{混合}{mingled}；不仅天主与出生、圣神与血肉、永恒与时间、无限与度量相混合；而且生育与童贞、羞辱与远超一切尊荣者、不受苦者与苦难、不朽者与可朽者相混合。因为那欺骗者自以为他的恶毒不可战胜，在我们被他以成为神的希望欺骗之后，他自己却被天主取了我们本性的计谋所欺骗；这样，在他以为攻击\Person{亚当}的时候，实际上却遇到了天主，因而新亚当拯救了旧亚当，肉身的定罪被废除，死亡被肉身所杀。\par

十四、在祂诞生时，我们适当地举行了庆节，包括我——这庆节的领导者，以及你们，以及世上和天上的一切。我们与星一同奔跑，与贤士一同朝拜，与牧羊人一同被光照，与天使一同光荣祂，与\Person{西默盎}一同将祂抱在怀中，与年老而贞洁的\Person{亚纳}一同作出应答的宣认。感谢那位以陌生者的形象来到自己处所者，因为祂光荣了那陌生者。现在，我们来到基督的另一行动和另一奥迹。我无法抑制我的喜乐；我被提升到天主那里。几乎像\Person{若翰}一样，我宣布喜讯；我虽不是前驱，却来自于旷野。基督受光照，让我们与祂一同发光。基督受洗，让我们与祂一同降下，好能与祂一同上升。耶稣受洗；但我们必须留心注意的不仅这一点，还有其他几点。祂是谁？谁为祂施洗？在什么时候？祂是全洁者；由\Person{若翰}施洗；时间是在祂奇迹的开端。我们要学习和教导什么？首先洁净自己；其次谦卑自下；第三，只在灵性和身体成熟时才宣讲。第一点特别针对那些仓促领洗、未作适当准备，或未以向善的心境确保圣洗恩宠之稳定的人。因为既然恩宠包含了对过去的赦免（因为它是\Emph{恩宠}），就更应敬畏，以免我们又回到那呕吐物中。第二点针对那些反抗这奥迹的管家的人，如果他们的地位更高。第三点针对那些自信于青春，认为任何时候都适合教导或主持的人。耶稣被洁净，而你还鄙视洁净吗？……且由\Person{若翰}洁净，而你还反抗你的前驱吗？……且在三十岁时，而你在未长胡须之前就敢教训长者，或认为你教训他们，尽管你的年龄并不令人尊敬，甚至或许你的品行也是如此？但有人会说，\Person{达尼尔}和某某人在年轻时作判官，这些例子常在你口中；因为每个作恶者都准备为自己辩护。但我回答说，罕见之事并非教会的法律。因为一燕不成夏，一线不成几何学家，一次航行不成水手。\par

十五、但若翰施洗，耶稣来到他跟前（见：\BibleRef{玛 3:14}）……也许是为圣化洗者本人，但更是为了将整个旧亚当埋葬在水里；而且在此之先并为此之故，来圣化约但河；因为祂既是圣神又是血肉，同样以圣神和水来圣化我们（见：\BibleRef{若 5:35}）。若翰不愿接受祂；耶稣坚持。\Quote{我需要祢给我施洗}（见：\BibleRef{玛 3:17}），这是声音对圣言、朋友对新郎说的（见：\BibleRef{若 3:39}）；在妇人所生之中最大的那位（见：\BibleRef{玛 11:11}）对一切受造的首生者说的（见：\BibleRef{哥 1:5}）；在母腹中跳跃的那位（见：\BibleRef{路 1:41}）对在母腹中被朝拜的那位说的；那位是、也必将为前驱者，对那位是、也必将被彰显的那位说的。\Quote{我需要祢给我施洗}；再加上\Quote{也为祢}；因为他知道他将要殉道而受洗，或者，像伯多禄那样，不仅要洗脚（见：\BibleRef{若 13:9}）。\Quote{祢反而来就我吗？}这也是预言性的；因为他知道在黑落德之后，将发生比拉多的疯狂，所以当他先行之后，基督也会随他而去。但耶稣说什么？\Quote{你暂且容许吧}，因为这是祂降生成人的时期；因为他知道再过不久，祂也要给洗者若翰施洗。那么什么是\Quote{簸箕？}就是净化。什么是\Quote{火？}就是烧灭糠秕的火焰，以及圣神的热力。什么是\Quote{斧头？}就是那灵魂的砍除，即使施加了粪肥也无药可救（见：\BibleRef{路 13:8}）。什么是\Quote{剑？}就是圣言的切割，它将较劣的与较优的分开（见：\BibleRef{希 4:12}），在信者与不信者之间作划分（见：\BibleRef{玛 10:35}）；激起儿子、女儿、新娘反对父亲、母亲、婆婆（见：\BibleRef{米 7:6}），年幼新鲜的反抗年长暗淡的。什么是鞋带，那是你——给耶稣施洗的若翰——所不能解的？（见：\BibleRef{若 1:27}）你，这旷野之人，不食人间烟火，新的厄里亚（见：\BibleRef{路 7:26}），比先知更大，因为你看见了你所预言的祂；你是旧约与新约的中保。这是什么？也许是那临在的消息与降生成人，其中没有最微小的一点可被解开，我所说的不是那些仍在血肉中、在基督里为婴孩的人，甚至也不是那些像若翰在圣神里的人。\par

十六、然而——耶稣从水中上来……因为祂将自己与整个世界一同提升……看见天门敞开，那是亚当因自己及其所有后裔所关闭的（见：\BibleRef{创 3:24}），如同乐园的大门因火剑而关闭。圣神为祂的天主性作证，因为祂降临在一位与自己相似者身上，正如来自天上的声音（因为被见证的那位是从那里来的），且如鸽子一般，因为祂尊重那身体（因为这身体也是天主，因它与天主的结合）而以有形的形象显现；此外，鸽子自古以来惯于宣告洪水的终结。但如果你以体积和重量来判断天主性，认为圣神是微小之物，因为祂以鸽子的形像降临，你这个人，对于最伟大的事物怀着轻贱渺小的思想，那么为了保持一致，你也必须蔑视天国，因为天国被比作一粒芥菜子（见：\BibleRef{玛 13:31}）；并且你必须高举那仇敌在耶稣的威严之上，因为它被称为一座大山（见：\BibleRef{匝 4:7}）、利维坦和水生之物的王，而基督被称为羔羊（见：\BibleRef{依 53:7}）、珍珠（见：\BibleRef{玛 13:46}）、水滴及类似的名字。\par

十七、既然我们的庆节是圣洗的，我们必须忍受一点艰难，与那位为了我们而取了形体、受洗并被钉十字架者一起受些苦；让我们谈论不同种类的洗礼，使我们出来后得到净化。梅瑟也曾施洗（\EditorialNote{肋未纪十一}），但那是在水里，并且在此之前还有云中和海中的洗礼（见：\BibleRef{格前 10:2}）。保禄说，这是预像；那海是水的预像，那云是圣神的预像；玛纳是生命之粮的预像；那饮料是神圣饮料的预像。若翰也曾施洗，但这不像犹太人的洗礼，因为它不仅在水里，而且也是\Quote{为悔改}。然而它并非完全属神，因为他没有加上\Quote{并在圣神内}。耶稣也曾施洗，但这是在圣神内。这是完美的洗礼。如果我可以稍作离题，那使你也被神化的一位，怎能不是天主呢？我也知道第四种洗礼——就是殉道与血的洗礼，基督自己也曾经历；这洗礼远比所有其他洗礼更为崇高，因为它不能被后来的污点所玷污。是的，我还知道第五种，即眼泪的洗礼，这要艰辛得多，由那每夜用眼泪洗涤床铺、浸湿床榻的人所领受；他的伤痕因他的邪恶而发臭；他哀恸行走，面容忧伤；他效法默纳舍的悔改（见：\BibleRef{编下 38:12}）和尼尼微人的谦卑（见：\BibleRef{纳 3:7}），天主曾怜悯了他们；他在圣殿里说出税吏的话，而被算为正义，胜过那顽固的法利塞人（见：\BibleRef{路 18:13}）；他像客纳罕妇人那样俯首乞求怜悯和碎屑，就是那非常饥饿的狗的食物（见：\BibleRef{玛 15:27}）。\par

十八. 然而我——因为我承认自己是一个人，即一种善变且易变的动物——我既热切领受此圣洗，也敬拜赐予我圣洗的那位，并将它传授给他人；且藉施行仁慈而为仁慈作准备。因为我知道自己也为软弱所困\BibleQuote{也为软弱所困}\BibleRef{希 5:2}，并且\BibleQuote{你们用什么尺度量给人，也要用什么尺度量给你们}\BibleRef{玛 7:2}。但你说什么，哦，新法利塞人，名号洁净而心不洁，你将\Person{诺瓦提安}的论点倾泻在我们身上，尽管你同样有软弱？你不给哭泣留任何余地吗？你不流一滴眼泪吗？愿你不遇到像你自己一样的审判者！你难道不因\Person{耶稣}的仁慈而羞愧吗？\BibleQuote{祂承担了我们的脆弱，背负了我们的疾病}\BibleRef{玛 8:17}，祂来不是召义人，而是召罪人悔改；祂喜爱仁慈胜过祭献；祂宽恕罪过直到七十个七次。如果真的是纯洁而非骄傲，使你那超然地位何等有福，制定超出人性能力的法律，并以绝望摧毁改进！因为两者同样邪恶：不受智慧约束的纵容，与永不宽恕的谴责；前者放松一切缰绳，后者以严厉扼杀一切。向我展示你的纯洁，我便会赞同你的大胆。然而事实上，我担心你满身疮痍，却会使它们不治。你难道不接纳\Person{达味}的悔改吗？他的痛悔甚至为他保全了先知之恩。还有伟大的\Person{伯多禄}本人，他在我们救主受难时陷入了人性的软弱？然而\Person{耶稣}接纳了他，并以三重提问与宣认治好了他的三重否认。或者你甚至拒绝承认他藉血得以成全（因为你的愚昧甚至到了那地步）？或者\Place{格林多}的犯过者？但\Person{保禄}看到他的改过，便肯定了对他之爱，并说明了理由：\BibleQuote{免得这样的人被过度的忧苦吞没}\BibleRef{格后 2:7}，被惩罚过重所压垮。而你却拒绝准许年轻寡妇因她们年龄易跌倒而结婚？\Person{保禄}冒然做了；当然你可以教导他，因为你已被提到第四层天，到了另一个乐园，听到了不可言传的话，并理解了你的福音更广的范围。\par

十九. 但你说这些罪不是在圣洗之后。你的证据在哪里？要么证明它，要么停止定罪；若有任何疑问，让爱德得胜。但你说\Person{诺瓦提安}不肯接纳那些在迫害中失足的人。你这是什么意思？如果他们不肯悔改，他是对的；我也拒绝接纳那些完全不屈服或不够充分屈服，且不肯以改过来抵消其罪的人；当我接纳他们时，我会给他们分配合适的地位；但如果他拒绝那些以痛哭消磨自己的人，我不会效仿他。为何\Person{诺瓦提安}缺乏爱德要成为我的标准？他从未惩罚贪婪，那是第二偶像崇拜；但他谴责奸淫，好像他自己不是血肉之躯。你说什么？这些话是否说服了你？来站在我们这边，即人道这边。让我们一起赞美上主。你们中任何人，即使对自己很有把握，也不敢说：\Quote{不要摸我，因为我是圣洁的，谁像我这样圣洁？}也让我们分享你的光亮。但或许我们没有说服你？那么我们将为你哭泣。若他们愿意，让他们跟随我们的路，即基督的路；若他们不愿，让他们走自己的路。或许他们将在火中受洗，在最后那个更痛苦更长的洗礼中，那火\BibleQuote{吞噬木柴如同草芥}\BibleRef{格前 3:12}，并烧尽一切邪恶的碎秸。\par

二十. 但让我们今天敬拜基督的圣洗；让我们善度庆节，不是放纵口腹，而是在心神中喜乐。我们如何奢侈呢？\BibleQuote{你们要洗濯，自洁。}\BibleRef{依 1:17} 若你因罪成为朱红，且血色较浅，就变成雪白；若你成为红色，且人浸浴在血中，也要被带到羊毛般的洁白。无论如何要净化，你便会洁净（因为天主最喜欢人的改过和救恩，为此而有每一篇讲论和每一件圣事），使你们好似世上的光，成为赋予其他一切人生命的力量；使你们作为完善的光，站立在那大光旁边，并学习天上光照的奥秘，被天主圣三更纯洁、更清晰地光照，而你们现在正按度量接受来自同一神性的同一光辉，经由我们的主耶稣基督；愿光荣与权能归于祂，至于无穷之世。阿们。\par

\SourceRecord{Translated by Charles Gordon Browne and James Edward Swallow.；Nicene and Post-Nicene Fathers, Second Series；Vol. 7.；Edited by Philip Schaff and Henry Wace.；Buffalo, NY: Christian Literature Publishing Co.,；1894.；<http://www.newadvent.org/fathers/310239.htm>.}

% structure: p0001 source=h1 target=section reason=page-title

\section{讲道第四十篇}

\Emph{论圣洗圣事。}\par

\EditorialNote{公元381年1月6日讲于君士坦丁堡，即《论神圣之光》讲道次日。}\par

一、昨天我们隆重庆祝了光辉的神圣之光节日；因为理应为我们得救恩而欢庆，且远胜于婚礼、生日、命名日、乔迁、孩子登记、周年纪念等人类为属世朋友所行的一切庆祝。而今天，让我们简要谈论圣洗及其带给我们的益处，虽然昨天的讲道已简略提及；部分因为时间紧迫，部分因为讲道必须避免冗长。因为讲道过长对人的耳朵是敌人，正如食物过多对身体一样……值得你们用心听取我们所说，并以准备好的心态领受我们关于如此重要主题的讲道，因为认识这圣事的能力，本身就是\OriginalTerm{光照}{Enlightenment}。\par

二、圣言为我们承认三种诞生：即自然诞生、圣洗的诞生和复活的诞生。其中，第一种在夜间，属于奴役，带有情欲；第二种在白昼，毁灭情欲，割断所有源于诞生的帷幕，引领向更高的生命；第三种更为可畏且短促，瞬间聚集全人类，站在其创造者面前，为在此世的服役和交往交账：看它是随从了肉身，还是随从了圣神上升，并崇拜了其新造的恩宠。我主耶稣基督亲自以自身显示了祂尊重所有这些诞生：第一种，通过那首次赋予生命的嘘气（\BibleRef{创 2:7}）；第二种，通过祂的降生成人和祂自己所受的圣洗；第三种，通过祂作为初果的复活；祂屈尊成为众弟兄中的首生者（\BibleRef{罗 8:29}），也同样成为从死者中首生者（\BibleRef{哥 1:18}）。\par

三、关于其中两种诞生，即第一种和最后一种，当前无须谈论。让我们讨论第二种，它现在为我们所必需，并使光明节得名。光照是灵魂的光辉，生命的转变，对天主的良心所提出的问题。它是对我们软弱的援助，对肉身的弃绝，对圣神的追随，与圣言的共融，受造物的改善，对罪的淹没，对光的参与，对黑暗的消散。它是通向天主的马车，与基督同死，心智的成全，信德之壁垒，天国钥匙，生命的改变，奴役的移除，锁链的解开，全人的重塑。我何必详细论述？光照是天主恩赐中最伟大、最辉煌的。因为正如我们称至圣所、雅歌要比其他更为全面和卓越，同样，这被称为光照，因其比我们所拥有的任何其他光照更为神圣。\par

四、正如赐予此恩的基督有众多不同称呼，这恩赐本身也一样，或因其本性的极大喜乐（如非常喜爱某物的人乐于使用其名），或因其恩惠的多样性反映于其名称上。我们称它为：\OriginalTerm{恩赐}{Gift}、\OriginalTerm{恩宠}{Grace}、\OriginalTerm{圣洗}{Baptism}、\OriginalTerm{傅油}{Unction}、\OriginalTerm{光照}{Illumination}、\OriginalTerm{不朽的外衣}{Clothing of Immortality}、\OriginalTerm{重生的洗濯}{Laver of Regeneration}、\OriginalTerm{印记}{Seal}，以及一切尊贵的称呼。我们称它为恩赐，因为它是白白赐予我们的；恩宠，因为它甚至赐予负债者；圣洗，因为罪与它一同被埋在水中；傅油，如同司祭和君王的傅油，因为曾被傅油者正是如此；光照，因其光辉；外衣，因其遮盖我们的羞耻；洗濯，因其洗涤我们；印记，因其保存我们，而且也是主权的标志。天空因之喜乐；天使因其光辉荣耀它。它是天上福乐的形象。我们确实渴望歌颂它的赞美，却不能相称地做到。\par

五、天主是光（\BibleRef{若一 1:5}）：至高的、不可接近的、不可言喻的，既不能用理智理解，也不能用嘴唇说出（\BibleRef{弟前 6:16}），祂赐生命给每一个理性的受造物（\BibleRef{若 1:9}）。祂在思想世界，如同太阳在感官世界；按我们被净化的程度向我们呈现；按祂呈现于我们理智的程度被爱；又按我们爱祂的程度被理解；祂自己默观并理解自己，并把自己倾注于外在。我所指的光，是在圣父、圣子和圣神内被默观的光，其财富是祂们本性的合一，以及祂们光辉的同一涌出。第二种光是天使，是那第一种光的一种流出或交通，其光芒来自对其的倾向和服从；我不知道它的光芒是按照其状态秩序分布，还是其秩序源于其各自光芒的程度。第三种光是人，是一种对外物可见的光。因为人称人为光，是由于我们内言语的能力。这名称又用于那些更相似天主、比他人更接近天主的人。我还承认另一种光，原初的黑暗被它驱散或刺穿。它是有形受造物中首先被召叫进入存在的；它照耀整个宇宙、星辰的环绕轨道和天上所有灯塔之火。\par

第六、光明也是那赐给首生之人的首条诫命（因为法律的诫命是灯，是光；\newline{}\BibleRef{箴 6:23}；又说：因为祢的判决是地上的光）；虽然嫉妒的黑暗潜入，行了邪恶。\newline{}那适合其臣民的、典型的光明，就是成文的法律，它预表真理和那大光明的圣事，因为梅瑟的脸曾因它而荣耀。\newline{}\BibleRef{出 34:30}再者，为了提及更多的光明——\newline{}那从火中向梅瑟显现的光明，它焚烧了荆棘，却没有烧毁它，以显示其本性并宣告其中所含的能力。\newline{}那在火柱中的光明，曾引导以色列并驯服了旷野。\newline{}\BibleRef{出 13:21}那是光明，在火车上厄里亚提升上去，\newline{}\BibleRef{列下 2:11}却没有在提升时烧毁他。\newline{}那是光明，当永恒的光明与暂时相混时，环绕在牧人们周围。\newline{}\BibleRef{路 2:9}那是光明，那星辰的美丽，走在前面前往白冷，引导贤士们的道路，\newline{}\BibleRef{玛 2:9}并护卫那在我们之上的光明，当祂来到我们中间时。\newline{}光明是那在山上向门徒们显示的天主性——但对他们的眼睛来说稍微太强了。\newline{}光明是那向保禄显现的异象，\newline{}\BibleRef{宗 9:3}并通过伤害他的眼睛，医治了他灵魂的黑暗。\newline{}光明也是天堂的光辉，给那些在此地已被净化的人，那时义人将如同太阳一样发光，\newline{}\BibleRef{玛 13:43}天主将站在他们中间，\newline{}\BibleRef{智 3:7}作为神和君王，决定和区分天堂真福的等级。\newline{}此外，还有一种特殊的光明，就是我们正在谈论的圣洗圣事的光照；因为它包含了一个伟大而奇妙的我们救恩的圣事。\par

第七、因为完全无罪只属于天主，属于那第一和纯一的性体（因为纯一是和平的，不受纷争影响），我敢说也属于天使的性体；至少，我主张那性体几乎是无罪的，因为它接近天主；但犯罪是人性的，属于地上的复合体（因为复合是分离的开始）；因此，主认为不宜让祂的受造物得不到帮助，或忽略祂与自己分离的危险；相反，正如祂赋予不存在者存在，同样祂赋予存在者新的创造，一种比第一次更神圣、更高的创造，这创造对于刚开始生活的人是一个印记，对于年龄较大者是一份恩赐和那因罪而堕落的肖像的恢复，好使我们不致因绝望而变得更坏，永远被带向更邪恶的事，因气馁而完全远离善与德行；并且，如经上所说，坠入邪恶的深渊而轻视祂；而应像那些在长途旅行中在客栈稍作休息的人那样，我们得以精力充沛、充满勇气地完成剩余的道路。这就是圣洗的恩宠和能力；不是像古时那样淹没世界，而是净化每个人的罪，彻底清除罪的一切瘀伤和污点。\par

第八、既然我们是双重构成的，即身体与灵魂，一部分可见，另一部分不可见，那么洁净也是双重的：通过水和圣神；一个可见地在身体中领受，另一个无形地与它结合，在身体之外；一个是象征的，另一个是真实的，洗净深处。而这洁净帮助我们第一次的出生，使我们变得更加而非旧，更像天主而非我们现在这样；它不通过火重塑我们，不破坏我们而重新创造我们。因为，一言以蔽之，圣洗的德能应被理解为与天主订立的盟约，为获得第二次生命和更纯洁的品行。事实上，人人都需要非常敬畏这一点，并各自以全部谨慎守护自己的灵魂，以免在这宣认上成为说谎者。因为如果天主被呼求作为中保来批准人的宣认，那么如果我们被发现违背了与天主自己所立的盟约，那危险是多么大；并且如果我们在那谎言面前，在真理本身面前被判有罪，除了我们的其他过犯之外……而那时没有第二次重生，或再创造，或恢复到我们以前的状态，即使我们竭尽全力，带着许多叹息和眼泪，藉此它才得以结痂（依我看，非常困难，尽管我们都相信它可以结痂）。然而，如果我们能甚至擦去伤疤，我会很高兴，因为我也需要怜悯。但更好的是不需要第二次洁净，而停留在第一次，我知道这第一次是众人共享的，无需劳作，对奴隶、主人、穷人、富人、卑微者、尊贵者、温良者、纯朴者、负债者、无债者价值相等；就像呼吸空气、倾泻光明、四季更替、瞻仰受造物（我们共同分享的巨大喜悦）以及信仰的平等分施一样。\par

IX. 因为以更痛苦的疗法取代无痛的治疗是奇怪的事；抛弃慈悲的恩宠，而欠下惩罚的债；以我们的改过来衡量罪。因为我们必须流出多少眼泪才能使它们相等于圣洗之泉？谁能为我们担保死亡会等待我们的痊愈，而审判台不会在我们仍为负债者、尚需要来世的火时召唤我们？你也许，像一位善良而可怜的农夫，恳求主人仍宽恕那棵无花果树，\BibleRef{路 13:8} 不要因为它不结果而砍下它；而是允许你用眼泪、叹息、祈祷、睡地、守夜、克苦灵魂与身体，以及藉告解与谦卑生活的改正来给它培土。但主人是否会宽恕它并不确定，因为它使另一位求怜悯的人的地土变得荒芜，并因对这棵树的忍耐而变坏。那么，让我们藉圣洗与基督同葬，为的是也与祂同复活；让我们与祂同降下，为的是也与祂同升高；让我们与祂同上升，为的是也同受光荣。\par

X. 如果圣洗之后，光明的迫害者和诱惑者攻击你（因为他甚至藉着帐幔攻击了我的天主圣言，藉着那显现出来的攻击了隐藏的光），你有战胜他的方法。不要惧怕争战；用水保护自己；用圣神保护自己，藉着祂，恶者的一切火箭将被熄灭。\BibleRef{弗 6:16} 那是圣神，却是撕裂山岭的圣神。\BibleRef{列上 19:11} 那是水，却是熄灭火焰的水。如果他因你的缺乏攻击你（正如他胆敢攻击基督），并要求石头变成饼，不要不知道他的诡计。\BibleRef{格后 2:11} 教导他所未学到的。用生命之圣言保护自己，祂是从天上降下的食粮，并赐生命给世界。\BibleRef{若 6:33} 如果他以虚荣谋害你（正如他对基督所做，当他带祂到圣殿顶端，对祂说：跳下去\BibleRef{玛 4:6}，作为你天主性的证明），不要被得意压倒。如果你被这所胜，他不会停在这里。因为他贪得无厌，他抓住一切。他以美好的借口谄媚你，但结局是恶；这是他战斗的方式。是的，这强盗精通圣经。一方面是关于饼的《经上记载》，另一方面是关于天使的《经上记载》。他说：\Quote{经上记载：\BibleQuote{祂必为你吩咐祂的天使，他们要用手托着你。}}卑鄙的诡辩家！你为何隐藏了后面的话？因为我深知，即使你默默跳过。我要使你践踏毒蛇和蜥蜴，我要践踏蛇和蝎子，被天主圣三所保护。如果他藉着贪婪向你摔跤，在瞬间、眨眼之间将万国展示给你，如同属于他，并要求你崇拜，视他为乞丐而轻视他。依靠印记对他说：\Quote{我本人是天主的肖像；我尚未因你的骄傲从天上荣耀中被抛下；我穿上了基督；我藉圣洗被转变为基督；你崇拜我吧。} 我深知，他将因此失败受辱而离开；正如他从基督——原初之光——那里离开，他也会从那些被基督光照的人那里离开。洗池将如此福分赐予领悟它的人；它为正确饥饿的人提供如此丰盛的宴席。\par

XI. 那么，让我们受洗以赢得胜利；让我们领受那洁净之水，比牛膝草更洁净，比律法之血更纯净，比洒在不洁者身上的小母牛灰更神圣，\BibleRef{希 10:4} 那水提供身体的暂时洁净，却不能完全除去罪；因为如果一次得了洁净，为何还需要进一步的净化？让我们今天受洗，免得明天受苦；不要推迟福分如同它是伤害，也不要等待变得更邪恶以便更多得赦免；不要成为贩卖基督的人，免得我们担负过重，全军覆没，使恩赐失事，因期望过多而丧失一切。当你仍是自己思想的主人时，奔向恩赐。当你身体或精神尚未患病，在旁人看来也不如此（其实你心智健全）；当你的好处尚未在他人权下，你自己仍是它的主人；当你的舌头尚未结巴或干涸，或（不再多说）失去说出圣事言语的能力；当你仍能被明认为信徒，而非猜测；并且仍能领受祝贺而非怜悯；当恩赐对你仍清晰，且无疑问；当恩宠能达到你灵魂的深处，而不仅仅是身体被洗为埋葬；在眼泪环绕你宣告你的死亡之前——甚至这些眼泪或许因你而被克制——你的妻子儿女会延迟你的离去，并倾听你的临终遗言；在医生无力帮助你，只给你数小时可活——这数小时并非他所能给——并以点头权衡你的救恩，在你死后博学地谈论你的疾病，或通过撤走来加重收费，或暗示绝望之前；在要为你施洗的人与谋取你钱财的人之间有争夺之前，一人争取使你领受你的天路行粮，另一人争取被写入你的遗嘱为继承人——而两者没有时间兼顾。\par

XII. 为何要等待热病将这恩赐带给你，却拒绝从天主手中领受它？为何你要藉着时间的流逝得到它，而不是藉着理性？为何你要将它归功于一个图谋不轨的朋友，而不是一种有益的渴望？为何你要被迫接受它，而不是出于自由意志；出于必要性，而非自由？为何你必得从别人口中听到你的死亡，而不是现在就想它如同已经来临？为何你寻求那些毫无益处的药物，或是危机时的汗液，而死亡的汗或许已经临到你身上？在极端之前医治你自己；怜悯你自己，你是你疾病唯一真正的医治者；将那真正有益的医药施于你自己；当你还在顺风中航行时，要惧怕船难，如果你利用恐惧作为助手，那么你遭遇船难的危险就会更小。让你有机会以欢宴庆祝这恩赐，而不是以哀悼；让那塔冷通得以培植，而不是埋在地里；让恩宠与死亡之间有一段间隔，这样不仅罪过的账目被涂销，而且有更好的东西写在它的位置上；这样你不仅拥有恩赐，也拥有赏报；你不仅逃脱火刑，也能继承荣耀，这荣耀是由培植恩赐而来的。因为对于小灵魂的人而言，逃脱折磨已是大事；但对于大灵魂的人而言，他们的目标也在于获得赏报。\par

XIII. 我知道得救者中有三类：奴隶、雇工和儿子。如果你是个奴隶，要惧怕鞭子；如果你是个雇工，只指望得到工钱；如果你超越这些，是个儿子，就要敬畏祂如父，并因为服从父是好的而去行善；即使对你而言没有赏报，这本身便是赏报，即你使父喜悦。那么，我们就当小心，不可轻视这些事。追求金钱而抛弃健康，这是多么荒谬；对身体洁净慷慨，却对灵魂洁净吝啬；寻求脱离地上的奴役，却不关心天上的自由；竭尽全力使房屋和衣着华丽，却从不思考如何使自己真正变得宝贵；热心向他人行善，却毫无行善于己的意愿，这又是多么荒谬。如果善可以用钱买到，你会不惜代价；但如果怜悯白白地在你脚下，你却因它廉价而蔑视它。每次洗礼时机都适合你，因为任何时候都可能是你的死期。我同保禄一起大声向你们呼喊：\Quote{看，如今正是悦纳的时候；看，如今正是救恩的日子；} \BibleRef{格后 6:2} 而那个\Quote{如今}并不指向某一个时间，而是每一个当下。同样，\Quote{你这睡眠的，醒来吧，基督必光照你，} \BibleRef{弗 5:14} 驱散罪的黑暗。因为正如依撒意亚所说，在夜里希望是邪恶的，在清晨被接纳更为有益。\par

XIV. 按时播种，按时收聚，该开门仓时就开门；按时栽植，让葡萄串成熟时被摘下，在春天大胆启航，在冬初海上开始翻腾时再将船拖上岸。也当有战争之时与和平之时；有结婚之时与禁婚之时；有友谊之时与不和之时，若有需要；简言之，万事皆有定时，若你听从撒罗满的劝告。这样做最好，因为这劝告有益。但你得救的工作是你应当时时从事的；且让每个时刻都成为你领受圣洗的确定时刻。若你总是推延今天而等待明天，你会在不知不觉中被那恶者以其惯用的伎俩，因你小小的拖延而欺骗。他对你说：把现在给我，把未来给天主；把你的青春给我，把老年给天主；把你的快乐给我，把你的无用给祂。环绕你的危险何等巨大！多少意外的不幸！战争消耗了你；地震压倒了你；大海吞没了你；野兽带走了你；疾病杀死了你；一块面包屑走错了路（一件最微不足道的事，但有什么比人死更容易呢？尽管你如此以天主的肖像自豪）；一次过度放纵的饮酒；一阵风把你吹倒；一匹马把你甩下；一剂居心叵测的药，或原本想有益健康却发现有害；一个不人道的法官；一个无情的刽子手；或任何使变化最迅速、非人力所能及的事物。\par

XV. 但若你事先以上印来巩固自己，并用最好最有力的帮助为自己确保将来，在身体和灵魂上都以傅油留下印记，如同古时以色列在夜间以那血和傅油保护长子一样，\BibleRef{出 12:22} 那么有什么事能临到你呢？又为你成就了什么？请听箴言。\Quote{你若坐下，他说，必无所惧；你若躺卧，你的睡眠必甘甜。} \BibleRef{箴 3:24} 再听达味向你报喜讯：\Quote{你不必害怕黑夜的惊骇，也不必害怕白日飞行的箭矢，黑暗中的瘟疫，正午的毒害。} 这在你活着时，将大大有助于你的安全感（因为一只被印封的羊不易被诱捕，而未标记的羊则是盗贼的易得之物），而在你死时，则是一件幸运的殓衣，比黄金更宝贵，比坟墓更华丽，比无果的奠祭更虔敬，比成熟的初熟之果更及时，这是死人献给死人的，以习俗制定法律。不，即使一切事物都离弃你，\BibleRef{路 9:60} 或被强行夺走——金钱、财产、王座、尊荣，以及这早期动荡所归属的一切——你仍能安全地躺下生命，在天主赐予你救恩的帮助上毫无损失。\par

XVI. 但你难道害怕毁掉这恩赐，因此拖延你的净化，因为你不能再次获得它？怎么？难道你不害怕在迫害时期失去你拥有的最宝贵的东西——基督？你因此避免成为基督徒？绝无此念。这样的恐惧不属于理智正常的人；这样的论证表明疯狂。哦，粗心的谨慎，如果我可以这么说。哦，恶者的诡计！他实在是黑暗，却装作光明；当他不能再在公开战争中得胜时，就暗中设下陷阱，给出表面良善实则邪恶的建议，为要至少通过某种诡计得逞，使我们无法逃脱他的阴谋。这显然就是他在此例中所图谋的。因为他无法说服你轻视圣洗，就通过一种虚假的安全感使你遭受损失；结果因你的恐惧，你可能无意识地遭受你所害怕的事；因为你害怕毁掉恩赐，反而因此完全失去恩赐。这是他的本性；只要他看到我们向上攀登那他从其中坠落的天堂，他就绝不会停止他的诡诈。因此，属神的人，你要认清你对手的阴谋；因为战斗是针对那已经拥有的，并且关乎最重要的利益。不要以你的敌人为顾问；不要轻视成为并被称为信徒。当你还是慕道者时，你只在信仰的门廊；你必须进入里面，穿过庭院，观察圣事，瞻仰至圣所，并与天主圣三为伴。你所为之战斗的利益是伟大的，你所需要的稳固也是伟大的。以信德的盾牌保护自己。你若用这武器武装作战，他便怕你，因此他想要剥去你的恩赐，好使他更容易战胜手无寸铁、毫无防备的你。他攻击每个年龄、每种生活方式；必须被所有人击退。\par

XVII. 你是年轻人吗？抵制你的情欲；加入天主的军队联盟；勇敢地对抗哥肋雅。\BibleRef{撒上 17：32} 夺取你的成千上万；这样享受你的青年时期；但不要让你的青春枯萎，被不完全的信德所杀。你是老人，且临近命定的终结吗？助益你为数不多的剩余日子。将净化托付给你的老年。你为什么在深浓的老年和临终时害怕青春的情欲？或者你要等到死后才被洗洁，那样不是可怜而是可厌？你是在痛悔快乐的渣滓，而你自己也处于生命的渣滓中？可耻的是，你确实过了青春的盛年，却没有过了你的邪恶；要么仍然陷在其中，至少因拖延净化而显得如此。你有婴孩吗？不要让罪恶有任何机会，而是让他从童年就被圣化；从他最柔嫩的年龄就被圣神祝圣。你因本性的软弱而害怕印号吗？哦，多么心胸狭隘、信德薄弱的母亲！为什么？亚纳甚至在撒慕尔出生前\BibleRef{撒上 1：10}就把他许诺给天主，并在出生后立即将他祝圣，以司铎的装扮养育他，不惧怕人性中的任何事，而是信赖天主。你不需要护身符或咒语，这些邪魔也会介入，在更虚荣的人心中将敬拜从天主那里偷窃归给自己。将天主圣三赐给你的孩子，那是伟大而尊贵的守护者。\par

XVIII. 更进一步？你生活在贞洁中吗？被这个净化所封印；让它成为你生命的分享者和伴侣。让它指导你的生命、你的言语、每个肢体、每个行动、每个感官。尊崇它，好使它尊崇你；它给你的头戴上恩宠的花冠，并用喜乐的花冠保护你。\BibleRef{德 32：3} 你受婚姻约束吗？也受印号约束；让它与你同住，作为你节制的守护者，比无数太监或看门人更安全。你尚未与肉体结合吗？不要害怕这祝圣；你在婚姻后仍是纯洁的。我承担那个风险。我为你主持婚礼。我为新娘装饰。我们并未因给予贞洁更高的尊荣而羞辱婚姻。我将效法基督，纯洁的傧相和新郎，祂既在婚宴上行奇迹，又以祂的临在尊崇婚姻。\BibleRef{若 2：1} 只愿婚姻纯洁，不与污秽的情欲混杂。我只要求这一点：从恩赐中接受安全，并将贞洁的祭献在适当时辰——当固定的祈祷时刻到来，以及那比俗务更宝贵之事时——献给恩赐。你们要彼此同意和认可而行这事。因为我们不是命令，而是劝勉；我们要从你们那里收受一些东西，为你们的益处和你们二人的共同安全。总之，没有一种生活状态或职业，圣洗不是有益的。你为自由人，受它约束；你为奴隶，得与平等；你为悲伤者，得安慰；喜乐者受管教；穷人获得不能夺去的财富；富人能够成为自己财产的好管家。不要耍诡计或阴谋反对你自己的救恩。因为即使我们能欺骗他人，也不能欺骗自己。因此，玩弄自己是非常危险和愚蠢的。\par

XIX. 但你必须在公务中生活，并被它们玷污；浪费这恩宠将是可怕的事。回答很简单：若能如此，连法庭也要逃避，与良善的同伴一起，为自己装上鹰的翅膀，或更恰当地说，鸽子的翅膀……因为凯撒的事或凯撒的物与你何干？……直到你能安息在没有罪、没有玷污、没有咬人的蛇阻拦你敬虔脚步的地方。将你的灵魂从世界夺回；逃离索多玛；逃离焚烧；前行不要回头，免得你被凝固成一根盐柱。\BibleRef{创 19:26} 逃往山上，免得你与平原一同毁灭。但若你已被必然的锁链捆绑和约束，就这样对自己推理；或更恰当地说，容我这样对你推理。达到善并保持净化两者都更好。但若不能两全，那么，为了公务而稍受玷污，总比完全失去恩宠要好；正如我认为，受父亲或主人一点惩罚，总比被赶出门外要好；受一点光照，总比全然黑暗要好。明智人的选择，在好事上选择更大更完美的，在坏事上选择更小更轻的。因此，不要过分害怕净化。因为我们的成功总是由公正而仁慈的审判者根据我们在生活中的处境来评判；往往一个在公共生活中取得微小成功的人，比一个享受自由却没有完全成功的人得到更大的赏报；正如我认为，一个戴着镣铐的人稍有进步，比一个没有负重的人奔跑更令人惊叹；或者在泥泞中行走只溅上一点泥点，比在干净路上完全干净更了不起。为证明我所说的：——妓女辣哈布仅因一事成义，即她的款待，尽管她其余的行为未得称赞；税吏因一事被高举，即他的谦卑，\BibleRef{路 18:14} 尽管他其他方面未得见证；这样你就能学会不要轻易对自己绝望。\par

XX. 但有人会说：\Quote{如果我被圣洗所占据，因我自己的仓促而斩断了生活的享乐，而当时我本可以放纵取乐，然后再获得恩宠，那我将得到什么？因为葡萄园中工作时间最长的工人并未因此获得更多，因为给最后来的人也付了同样的工资。}你使我免去了些麻烦，无论你是谁，因为最终你艰难地说出了你迟延的秘密；虽然我不能称赞你的诡诈，但我称赞你的告白。但请过来听这比喻的解释，免得你因缺乏知识而被圣经伤害。首先，这里不是关于圣洗的问题，而是关于在不同时间相信并进入教会这座好葡萄园的人。因为从每个人相信的那天和那一刻起，他便被要求开始工作。而且，虽然先进入的人在劳动的度量上贡献更多，但他们在意志的度量上并未贡献更多；不，也许在这一点上，最后来的人还应当得更多，尽管这话听起来可能矛盾。因为他们后来进入的原因，是他们后来被召叫到葡萄园的工作。在其他所有方面，让我们看看他们有多么不同。先来的人是在商定了工钱之后才相信并进入的；但其他人没有协议就上前来工作，这是更大信德的证明。并且先来的人被发现具有嫉妒和抱怨的本性，但没有人对其他人提出这样的指控。对于先来的人，所给的是工资，尽管他们是无用的家伙；对于最后来的人，那是恩赐。因此，先来的人被证明为愚妄，并有理被剥夺了更大的赏报。让我们看看如果他们是迟来的会发生什么。为什么，显然会得到同样的工资。那么他们怎能因平等而指责主人不义呢？因为这一切都从先来的人那里夺去了他们劳动的功绩，尽管他们先开始工作；因此，如果你用善意来衡量劳动，平等的工资分配是公正的。\par

XXI. 但假设按照你的解释，这比喻的确描绘了圣洗池的能力，那么，如果你先进来并忍受了炎热，从而避免嫉妒最后来的人，藉此仁爱本身你或许能获得更多，并领受赏报，不凭恩宠，而凭债务，是什么阻止你呢？其次，领受工资的工人是那些进入了葡萄园的，而不是那些错过了的；后者似乎就是你的情况。因此，如果你确定会获得这恩赐，尽管你有这样的心态，并恶毒地扣留了一些劳动，你或许可以被原谅，因为你躲在这些论据中，并想从主人的仁慈中非法获利；虽然我可以向你保证，对于任何不完全是唯利是图的人来说，能够劳动本身就是更大的赏报。但由于你通过讨价还价有被完全关在葡萄园外的危险，并因停下拾取小利而损失资本，但愿你们被我的话说服，放弃虚假的解释和矛盾，不作争辩地前来领受这恩赐，免得你们在实现希望之前被夺去，并发现你们炮制这些诡辩是自招损失。\par

XXII. 但你说，难道天主不是仁慈的吗？既然祂知道我们的思想，洞察我们的欲望，祂岂不会以圣洗的愿望来代替圣洗吗？你是在说谜语，如果你的意思是，因天主的仁慈，未受光照者在祂眼中成为光照；而仅仅渴望进入天国、却不做属于天国之事的人，就在天国里。然而，我要大胆说出我对这些事的看法；我认为所有其他有见识的人都会站在我这一边。在领受恩赐的人中，有些人与天主和救恩完全疏离，他们沉溺于各种罪恶，并且渴望作恶；另一些人是半恶的，处于善恶之间的中间状态；还有些人虽行恶，却不赞同自己的行为，犹如发热的人不喜欢自己的疾病。还有一些人甚至在受光照之前就值得称赞；部分出于天性，部分出于他们为圣洗而准备的用心。这些人在领受圣事后变得明显更好，更不易跌倒；前者是为获取善，后者是为保存善。在这些人中，那些屈服于\Emph{某些}恶的，比那些完全败坏的人要好；而那些更加热心、在圣洗前就开垦荒地的人，比那些稍有屈服的人更好；他们拥有已经劳苦的优势，因为圣洗泉并不像消除罪恶那样消除善行。但比这些更好的，是那些也培育恩赐、并尽最大可能美化自己的人。\par

XXIII. 同样，在那些未能领受恩赐的人中，有些人是完全属血气的或兽性的，根据他们是愚蠢或邪恶；并且我认为，这要加在他们的其他罪恶上，即他们对这恩赐毫无敬畏，只视之为一份礼物——若给了他们，他们便接受；若不给，他们便忽视。另一些人知道并尊崇这恩赐，却推迟领受；有的出于懒惰，有的出于贪婪。还有一些人没有条件领受它，或许是由于年幼，或某种完全非自愿的情况，使他们即使愿意也无法领受。正如我们在前一种情况中发现许多差异，在此也如此。完全藐视它的人，比因贪婪或疏忽而忽视它的人更坏。这些比那些因无知或暴政而失去恩赐的人更坏，因为暴政不过是一种非自愿的错误。我认为，第一种人将受惩罚，既为所有罪恶，也为藐视圣洗；第二种人也将受惩罚，但较轻，因为他们的失败更多出于愚昧而非邪恶；第三种人既不被公义判官荣耀，也不受惩罚，因为他们未被封印，却也不是邪恶之人，而是受难者而非作恶者。因为并非每个不够坏到受惩罚的人都好到配受荣耀；同样，并非每个不够好到配受荣耀的人都坏到配受惩罚。我也从另一个角度看待此事。如果你单凭意愿就判断一个有意杀人者，而不考虑杀人行为，那么你也可以将渴望圣洗却未领受的人算作受洗者。但如果你不能做前者，又怎能做后者？我不明白。或者，如果你愿意，我们这样说：如果你认为愿望与实际圣洗有同等效力，那么对待荣耀也采取同样标准，你可以满足于渴望荣耀，好像渴望本身就是荣耀。只要你渴望荣耀，得不到实际的荣耀又有什么损害呢？\par

XXIV. 因此，既然你听到了这些话，就上前来领受光照，你的脸就不会因错过恩宠而羞愧。那么，按时领受光照，免得黑暗追赶你，抓住你，将你与光明隔绝。黑夜来临，那时没有人能工作（\BibleRef{若 12:35}），在我们离开之后。前者是大卫的声音，后者是真光的声音，那光光照每一个来到世界的人。请看撒罗满如何责备你们这些过于懒惰或懈怠的人，他说：懒惰人，你要睡到几时？你何时睡醒呢？（\BibleRef{箴 6:9}）你依靠这个或那个，并在罪恶中\Quote{编造借口；}我等待主显节；我更喜欢复活节；我要等待五旬节。与基督一同受洗更好，在祂复活之日与基督一同复活（\BibleRef{玛 24:50}），以尊崇圣神的显现。然后呢？结局会在你所不期待的日子、你所不知道的时刻突然来临；那时你将缺少恩宠为伴；你将在那些丰富的良善中饥饿，尽管你本应从相反的道路收获相反的果实，即由勤勉获得的丰收，和从圣洗泉得来的清凉，如同渴慕的牡鹿急奔泉源，用水消除奔跑的辛劳；而不致处于依市玛耳的情形，因缺水而干渴，或如寓言所说，在泉源中间受渴受罚。让集市日过去再找工作是可悲的。让玛纳过去再渴求食物是可悲的。迟来的劝告，只有在无法弥补时才意识到损失，也是可悲的；那就是在我们离开之后，每个人生命行为的痛苦终结，罪人的惩罚，被净化者的荣耀。因此，不要迟延来领受恩宠，要赶快，免得强盗超过你，免得奸淫者越过你，免得贪得无厌者在你之前得满足，免得凶手先夺走福分，或税吏、或奸淫者，或这些以暴力夺取天国的人中的任何一个（\BibleRef{玛 11:12}）。因为天国甘愿受暴力，并且因良善而受压迫。\par

XXV. 请接受我的劝告，我的朋友，作恶要迟缓，但奔向你的救恩却要迅速；因为对邪恶的急迫和对善的迟延同样不良。若你受邀参加狂欢，不要急于前往；若受邀背教，则速速跳开；若一伙作恶者对你说：\BibleQuote{来与我们同流合污，让我们藏匿一个义人，不义地将他埋在地下，}\BibleRef{箴 1:11}，连你的耳朵也不要借给他们。这样，你将获得两个极大的益处：你既使对方知道他的罪，又使自己脱离恶伴。但若伟大的达味对你说：\Quote{来，让我们在主内欢欣}；或另一位先知说：\Quote{来，让我们登上主的山}\BibleRef{米 4:2}；或我们的救主亲自说：\Quote{凡劳苦和负重担的，你们都到我这里来，我要使你们安息}\BibleRef{玛 11:28}；或说：\Quote{起来，我们离开这里，闪耀发光，洁白胜过雪，白于奶，亮过蓝宝石}；我们不要抗拒或迟延。让我们像伯多禄和若望一样，赶快前去；\BibleRef{若 20:3} 如同他们奔向坟墓与复活那样，我们也奔向洗礼池；一同奔跑，彼此竞赛，争先获得这一恩赐。你不要说：\BibleQuote{去吧，改天再来，明天我受洗}\BibleRef{箴 3:28}，而你今天便可获得这恩赐。\Quote{我要与父亲、母亲、兄弟、妻子、儿女、朋友和我所珍视的一切人同在，然后我再得救；但使我成为光明的适宜时刻尚未到来。}因为若你这样说，恐怕那些你希望分享你喜乐的人，倒成了与你同忧的人。若他们与你同在，固然好——但不要等候他们。因为说以下的话是可耻的：\Quote{可是，我为受洗所备的献仪在哪里？使我成为光明的洗礼袍在哪里？款待为我施洗者所需的物品在哪里？好叫在这些事上，我也值得被注意？因为你看，这一切都是必需的，为此恩宠会减少。}不要如此轻慢大事，也不要容许自己如此卑劣地思想。圣事比可见的环境更为伟大。献上\Emph{你自己}；穿上基督，以你的品行宴请我；我乐于受到如此款待，而赐予这些大恩的天主也乐于如此。在天主眼中，没有什么大事是穷人所不能献的，这样穷人在此也不致落后，因为他们不能与富人竞争。在其他事上，穷人与富人有区别，但在此处，更情愿者更为富有。\par

XXVI. 不要让任何事阻碍你前进，也不要使你从你的急切中分心。趁着你的渴望仍炽烈，抓住你所渴望的。趁着铁还热，让它被冷水淬炼，免得期间有事情发生，消灭你的渴望。我是斐理伯，你当是甘达刻的太监。\BibleRef{宗 8:36} 你也说：\Quote{看，这里有水，什么能阻止我受洗呢？}抓住机会；为这恩赐而大大欢喜；说了就受洗；受了洗就得以得救；即使你是埃塞俄比亚人的身体，也要使灵魂变白。不要说：\Quote{一位主教将给我施洗——而且他是一位都主教——而且是耶路撒冷的主教（因为恩宠不是来自地方，而是来自圣神）——而且他出身高贵，因为我的高贵若被一个出身低微的人施洗，那可真是可悲的事。}不要说：\Quote{我不介意一位普通的司铎，只要他是独身的、虔诚的、度天使般生活的；因为我在洁净的时刻反被玷污，那是可悲的。}不要询问讲道者或施洗者的资格凭证。因为判断他的另有其人，\BibleRef{撒上 16:7} 并且他要审查你所不能看见的事。原来人看外貌，上主却看人心。但对你而言，任何人都可以是洁净你的可信赖者，只要他是已被认可的人之一，不是公开被定罪的人，也不是与教会不相干的人。你这位需要医治的人，不要判断你的法官；也不要在那些洁净你的人的地位上挑剔，或对你的灵性父亲们吹毛求疵。一个可能比另一个高或低，但所有的人都高于你。你且这样看：一个可能是金的，另一个是铁的，但两者都是戒指，并且刻着同一个君王的肖像；因此当他们印在蜡上时，一个的印痕与另一个的印痕有什么区别呢？没有。你若如此聪明，就在蜡上分辨材料吧。告诉我哪个是铁戒指的印痕，哪个是金的。它们如何成为同一个？区别在于材料，而非印痕。因此任何人都可以为你施洗；因为虽然一个人可能在生活上优于另一个人，但洗礼的恩宠是相同的，任何一位具有同样信仰的人都可以圣化你。\par

XXVII. 不要轻视与穷人一同受洗，即使你是富人；也不要轻视与出身低贱的人一同受洗，即使你是贵族；也不要轻视与直到如今是你奴隶的人一同受洗，即使你是主人。即使这样，你也不会像基督那样自卑，你今天就是受洗归于祂，祂为了你的缘故甚至取了奴隶的形状。从你重生的那一天起，所有旧的标记都被抹去，基督穿在了所有人身上，形状相同。不要轻视告明你的罪，要知道若翰是如何施洗的，好藉着现在的羞愧逃脱将来的羞愧（因为这也是将来惩罚的一部分）；并通过公开展示你的罪，将其视为可鄙之事而战胜它，来证明你真正憎恨罪恶。不要拒绝驱魔的药剂，也不要因其冗长而拒绝它。这也是检验你是否对恩宠有正确态度的试金石。你所做的苦工，与埃塞俄比亚女王\BibleRef{列上 10:1}相比算什么呢？她起身从地极来到，要听撒罗满的智慧。看，这里有一位大于撒罗满的\BibleRef{玛 12:42}，按那些成熟推理之人的判断。不要因旅途遥远、海路遥远、火焰（如果这也摆在你面前）、或任何大小障碍而犹豫，使你可能获得这恩赐。但如果你可以毫不费力、毫无麻烦地得到你所渴望的，那么拖延恩赐是多么愚蠢呢：\BibleQuote{凡口渴的，请到水边来！}\BibleRef{依 55:1}，依撒意亚邀请你，\BibleQuote{没有钱的，也来买酒和奶，不用金钱，不用代价。}哦，祂仁慈的快速！哦，盟约的容易！这祝福只需你愿意就可以买到；祂将渴望本身视为巨大的代价；祂渴望被人渴望；祂给所有渴望喝的人喝；祂将被人求恩惠视为一种恩惠；祂乐意且慷慨地给予；祂给予的喜悦比他人领受的更大。\BibleRef{宗 20:35}只是不要让我们因祈求微小的、与施予者不相称的东西而被定为轻浮。耶稣从那位撒玛黎雅妇人那里求水喝，并赐给她涌到永生的水泉，这样的人是有福的。\BibleRef{若 4:7}那播种在一切水边\BibleRef{依 32:20}的人是有福的，在每一灵魂上，今天用牛和驴践踏、干涸无水、被非理性压迫的土地，明天将被犁耕和浇灌。那虽然是\Quote{芦苇谷}，却从上主圣殿得到浇灌的人是有福的；因为他结出了果实而非芦苇，产出了人的食物，而非粗糙无益的东西。为此，我们必须非常小心，不要错过这恩宠。\par

XXVIII. 诚然，有人会说，对于请求圣洗的人，情况如此；但你对那些仍是孩童、既不知损失也不知恩宠的人，有什么可说的呢？我们也要给他们施洗吗？当然，如果有危险迫近。因为他们在不知不觉中被圣化，比他们未受印、未入门而离去更好。\par

对此的一个证明是在第八天行割损礼，这是一种预像之印，在孩童有理智之前就赋予他们。同样，涂门框\BibleRef{出 12:22}也保存了长子，尽管涂在没有知觉的东西上。但关于其他人，我建议等到三岁末，或多一点或少一点，那时他们能够听并回答一些关于圣事的问题；这样，即使他们不完全理解，至少也知道了轮廓；然后以我们奉献的大圣事圣化他们的灵魂和身体。因为情况是这样：那时他们开始为自己的生命负责，当理智成熟，他们学到生命的奥秘（因为由于年幼，他们对无知的罪没有责任），而藉着圣洗得到坚固，在所有方面都是更有益的，因为突然的危险袭击我们，比我们的帮助者更强。\par

XXIX. 但有人说，基督受洗时三十岁\BibleRef{路 3:23}，尽管祂是天主；而你却要我们催促我们的圣洗吗？——当你说祂是天主时，你就解决了困难。因为祂是绝对的洁净；祂不需要洁净；但祂是为了你而被净化，正如为了你，祂虽然没有肉体，却穿上了肉体。推迟受洗对祂没有任何危险，因为祂安排了自己的苦难，如同安排了自己的出生一样。但在你的情况中，危险非同小可，如果你生于腐朽，没有穿上不朽就离去。还有一点我要考虑，那个特殊的受洗时间对祂是必要的，但你的情况不同。祂在出生后第三十年显明自己，而不是更早；首先，为使祂不显得炫耀，这是庸俗之心的品性；其次，因为那个年龄充分考验美德，是教导的正确时间。既然祂必须经受拯救世界的苦难，那么属于苦难的一切也必须配合苦难：显现、圣洗、来自天上的见证、宣告、群众的聚集、奇迹；它们应该像是一个身体，不被撕开，也不被间隔打断。因为从圣洗和宣告中产生了人们聚集的地震——圣经这样称呼那个时间\BibleRef{玛 21:10}——从群众中产生了征兆和奇迹的显现，这些引领到福音。从这些产生了嫉妒，从嫉妒产生了仇恨，从仇恨产生了谋害祂的阴谋和背叛；从这些产生了十字架，以及我们得救所藉的其他事件。关于基督的理由，就我们所能达到的而言，就是这样。也许还可以找到另一个更秘密的理由。\par

XXX. 但是，对于你而言，有什么必要去效法那些远高于你的榜样，而作出这样对自身不利的事呢？因为福音史中有许多细节与今天的情况大不相同，而且时间也不吻合。例如，基督在祂受试探前守斋，我们在复活节前守斋。就守斋的日子而言是相同的，但时间的差异却非同小可。祂用守斋来武装自己以对抗试探，而对我们来说，这斋期象征着与基督同死，是为庆节作准备的净化。祂完全守了四十天的斋，因为祂是天主；而我们按自己的能力来衡量守斋，尽管有些人因热心而超出自己的力量。再者，祂在楼房里，晚餐后，受难前一天，将逾越节圣事赐给门徒；但我们是在祈祷所中，在食物之前，在祂复活后举行。祂第三天复活；我们的复活则在很长时间之后。但与祂有关的事，既不完全与我们分离，在时间上也未与我们的事紧密结合；它们被传授给我们，只是作为我们行为的榜样，并刻意避免完全和精确的相似。\par

XXXI. 因此，如果你愿意听从我，你就要永远告别所有这些争论，欣然接受这祝福，并开始双重斗争：首先，通过净化自己来为圣洗预备；其次，保存圣洗的恩赐——因为获得我们尚未拥有的祝福与保持已获得的祝福同样困难。因为热心所获得的东西，往往被懒惰所毁灭；而犹豫所失去的，勤勉又能够重新获得。为实现你所渴望的，守夜、斋戒、卧地、祈祷、眼泪、对需要者的怜悯与施舍都是极大的帮助。让这些成为你对所领受的感恩，同时也是它们的保障。你领受的恩惠提醒你许多诫命，所以不要违背它们。一个穷人走近你吗？记住你曾是多么贫穷，又变得多么富有。一个缺食少饮的人，或许是另一个拉匝禄，被放在你的门口；尊重你所走近的圣体桌，你所分享的饼，你所共饮的杯，它们因基督的苦难而祝圣。如果一个陌生人、无家可归的外方人扑倒在你脚下，你要在他内欢迎那位因你而成为陌生人，且是在自己人中，\BibleRef{若 1:11}，并藉祂的恩宠来住在你内、引领你走向天上居所的那一位。要成为匝凯，他昨天还是税吏，今天却成了慷慨的人；将一切奉献给基督的来临，这样即使身材矮小，你也能因高尚地默观基督而显得伟大。一个患病或受伤的人躺在你面前；尊重你自己的健康，以及基督从你身上解除的创伤。若你看见赤身露体的人，给他衣服穿，以尊崇你不朽的衣裳，那就是基督，因为凡受洗归于基督的，就是穿上了基督。\BibleRef{迦 3:27} 若你发现一个欠债者扑倒在你脚下，撕毁每一张字据，无论合宜与否。要记得基督曾免除你的一万塔冷通，不要苛刻地追讨一笔小债——而且是从谁那里？从你的同伴身上，你尚且从主人那里得到了更大的宽免。否则，你必须向祂的仁慈交代，因为你没有效法祂的仁慈。\par

XXXII. 让圣洗之水不仅作用于你的身体，也作用于你内的天主的肖像；不仅洗去你的罪，也纠正你的脾气；不仅洗去旧的污秽，也净化源头。让它不仅促使你正当地获得，也教导你体面地失去拥有；或者，更容易的是，归还你非法所得。因为你的罪已被赦免，但你所造成的损失若不补偿给受损害者，又有什么益处呢？你的良心上有两项罪：一是你得了不义之财，二是你保留了这些财物。你得了第一项的宽恕，但就第二项而言，你仍在罪中，因为你仍占有别人的东西；你的罪并未终止，只是被时间分隔。一部分罪在你受洗前犯下，另一部分在受洗后仍存留；因为圣洗带来对过去、而非现在罪的宽恕；它的净化不可被戏弄，而必须真实地印刻在你身上；你必须完全洁净，而不只是被涂抹；你领受的恩赐不是仅仅遮盖罪，而是彻底洗除罪。\BibleQuote{罪恶得赦免的人是有福的}……这是彻底净化所成就的……\BibleQuote{罪过被掩盖的人是有福的}……这属于那些内心深处尚未痊愈的人。\BibleQuote{上主不归咎罪过的人是有福的}……这是第三类罪人，他们的行为虽不可称道，但意向清白。\par

三十三、那么我说什么，我的论点是什么？昨天你是一个被罪弯曲的客纳罕灵魂，今天你被圣言拉直了。不要再被弯曲，被定罪于地，好像被魔鬼用木轭压着，也不要患上不治的弯曲。昨天你因大量的血漏而干涸\BibleRef{玛 9:20}，因为你流出了朱红的罪；今天你被止住并重新繁茂，因为你触摸了基督的衣边，你的漏症停止了。我求你，守护这洁净，免得你又血漏，并且不能抓住基督偷取救恩；因为基督不喜欢经常被偷窃，虽然祂非常仁慈。昨天你被抛在床上，精疲力竭、瘫痪，当水被搅动时，没有人把你放入水池。今天你有了祂，祂既是人又是天主，或者更准确地说，既是天主又是人。你从床上被抬起，或者更准确地说，你拿起你的床，公开承认了这恩惠。不要再因犯罪而被抛在床上，在肉体被其享乐麻痹的邪恶安息中。但正如你现在这样，行走，记住那命令：\Quote{看，你已痊愈了；不要再犯罪，免得你遭受更糟的事}，如果你在领受祝福后证明自己是恶的。你曾听见那大声呼喊：\Quote{拉匝禄，出来！}\BibleRef{若 11:43}，当你躺在坟墓里时；然而，不是四天之后，而是许多天之后；你从殓布的束缚中被释放。不要再次死去，也不要与住在坟墓里的人同住\BibleRef{谷 5:3}；也不要被你自己罪的绳索捆绑；因为直到最后的普遍复活，你能否从坟墓中复活是不确定的，那复活将使每一件工作受审判\BibleRef{训 12:14}，不是为得医治，而是为受审判，并为它所积存的一切善恶交账。\par

三十四、如果你满身麻风病，那无形的恶，但你刮去了恶物，重新获得完整的肖像。向我这司祭显示你的洁净，让我认识到它比法律的洁净更珍贵多少。不要把自己列在九个忘恩的人中，而要效法第十个。因为虽然他是撒玛黎雅人，但他比其他人更有见识。确保你不会再长出恶疮，发现身体的不适难以治愈。昨天吝啬和贪婪枯萎了你的手；今天让慷慨和仁慈伸展它。对于一只软弱的手，一个高贵的治疗是分散施舍，赒济穷人，倾泻我们所拥有的丰富事物，直到我们触及最底层；也许这会给您涌出食物，如同给匝尔法特的妇人那样，尤其是如果你恰好供养一个厄里亚，认识到为了基督的缘故而匮乏是一种好的富足，祂为了我们而成为贫穷。如果你又聋又哑，让圣言在你的耳中发声\BibleRef{谷 7:32}，或者更确切地说，让那发声者住在那里。不要向主的训诲和祂的劝告闭耳，如同毒蛇不听符咒。如果你眼瞎未蒙光照，照亮你的眼睛，免得你沉睡至死。在天主的光中看见光，并在天主之神中由圣子得光照，那是三重且不可分的光。如果你接受全部的圣言，你将在自己的灵魂上带来基督所有的医治能力，这些能力分别医治了这些人。只是不可不知恩宠的尺度；只是不可让仇敌在你睡觉时恶意地撒下莠子\BibleRef{玛 13:25}。只是要小心，既然你因洁净而成为恶者的仇敌，不要再因犯罪使自己成为可怜的对象。只是要谨慎，免得因祝福而喜乐并过分高抬时，又因骄傲而跌倒。只是要勤于你的洁净，\Quote{在你心中设置上升阶梯}，并竭尽全力保持你所领受的赦免之恩，以便赦免来自天主，而保存它也来自你自己。\par

三十五、这事如何成就？要常想起这比喻\BibleRef{路 11:24}，这样你将最好、最完美地帮助自己。那污秽而恶毒的灵从你身上出去了，被圣洗所驱逐。他不服从这驱逐，不愿无家可归：他走过无水之地，干涸于天主之流之外，并在那里渴望停留。他游荡，寻找安息；却找不着。他落在已受洗的灵魂上，他们的罪已被圣洗池洗去。他惧怕水；他被这洁净窒息，正如那军旅在海中一样\BibleRef{谷 5:13}。他又回到他出来的房屋。他无耻，他好争，他重新攻击，重新尝试。如果他发现基督已在那里居住，并充满了他所空出的地方，他就又被赶回，悻悻而去，成为他流浪状态中可怜的对象。但若他在你身上发现一个地方，确实打扫干净、装饰整齐，却空虚无主，随时准备接纳首先占据它的任何事物，他就一跃而入，带着更大的队伍住下；那最后的状态比先前更坏，因为先前尚有改过与得救的希望，但现在邪恶猖獗，因逃离善而拖入罪，因此那居留者对占据更加稳固。\par

XXXVI. 我要再次提醒你们关于光照之事，且要经常提醒，并从圣经中列举它们。因为我本人因记起它们而感到更幸福（对那些尝过光的人而言，有什么比光更甜蜜呢？），并且我要用我的言语使你们眩目。义人有了兴起的光，它的伴侣是喜乐欢喜。而且，\BibleQuote{义人的光恒久不变}；\BibleRef{箴 13:9}\Quote{你从永恒的山上奇妙地闪耀}，这是对天主说的，我想是指那些帮助我们向善努力的天使的能力。你们也听到了达味的话：\Quote{主是我的光明，我的救恩，我还怕谁呢？}现在他求光明和真理为他发出，现在感谢他分享光明，因为天主的光已印在他身上；也就是说，所赐光照的标记已印在他身上并被认出。只有一个光让我们逃避——那就是那有害之火所生的光；\BibleQuote{让我们不要行走在我们自己火的光中，也不要在我们所点燃的火焰中行走}\BibleRef{依 50:11}因为我知道一种洁净之火，是基督来要投到地上的，\BibleRef{路 12:49}祂自己也被\OriginalTerm{寓意地}{anagogically}称为火。这火除去一切物质的和恶习的东西；祂渴望尽快点燃这火，因为祂渴望快快地赐福我们，甚至赐给我们炭火来帮助我们。我也知道一种火不是洁净的，而是复仇的；或是索多玛的火，\BibleRef{创 19:24}祂倾倒在所有罪人身上，夹杂着硫磺和暴风；或是为魔鬼及其使者所预备的火；\BibleRef{玛 25:41}或是从主的面容而出的火，要烧灭祂四围的仇敌；还有比这些更可怕的：那与不死之虫并列的不灭之火，为恶人是永恒的。因为所有这些都属于毁灭之力；虽然有些人甚至在这里可能对火采取更仁慈的看法，配得上那惩治者。\par

XXXVII. 正如我知道两种火，我也知道两种光。一种是我们统治之力的光，指引我们的脚步符合天主的旨意；另一种是欺骗性的、多管闲事的光，与真光截然相反，却假装是那光，以便用外表欺骗我们。这实际上是黑暗，却有午正的外表，光的高峰。因此我读到那段论及那些不断在午正黑暗里逃遁的人的话；\BibleRef{依 16:3}因为那真是黑夜，却被那些被奢华毁坏的人认为是明亮的光。因为达味说什么？\Quote{黑夜围绕着我，我却不自知，因为我以为我的奢华就是光照。}但这样的人就是如此，处于这种状态；但让我们为自己点燃认识的光。这将通过播种正义和收获生命的果实而实现，因为行动是默观的保护者，这样我们也能学习什么是真光，什么是假光，并避免不慎陷入邪恶伪装成善良的陷阱。让我们成为光，正如大光对门徒们所说的：\BibleQuote{你们是世界的光}\BibleRef{玛 5:14}让我们成为世界上的光，\BibleQuote{持守生命之言}\BibleRef{斐 2:15}也就是说，让我们成为他人的赋予生命的力量。让我们抓住天主性；让我们抓住那第一的、最明亮的光。让我们向着祂行走、闪耀，\BibleQuote{免得我们的脚绊倒在黑暗与敌意的山上}\BibleRef{耶 42:16}趁着白日，\BibleQuote{让我们端正地行走如同在白昼，不在宴乐和醉酒中，不在好色和放荡中}\BibleRef{罗 13:13}这些都是黑夜的羞耻。\par

XXXVIII. 弟兄们，让我们洁净每一个肢体，让我们净化每一种感官；让我们身上没有一样是不完美的，或是属于初生的；让我们不留下任何未受光照的部分。\BibleQuote{让我们照明我们的眼睛}\BibleRef{箴 4:25}使我们能正视前方，不因好奇而忙碌的目光而在自己内携带任何娼妓偶像；因为即使我们不崇拜情欲，我们的灵魂也会被玷污。\BibleQuote{如果有梁木或木屑}\BibleRef{玛 7:2}让我们除去它，使我们也能看见他人的。让我们的耳朵被光照；让我们的舌头被光照，使我们能听天主上主要说什么，并使祂在早晨使我们听到祂的慈爱，使我们听到喜乐和快乐的声音，传入虔诚的耳朵，使我们不是一把利剑，也不是磨快的剃刀，也不在我们的舌头下转动劳苦和辛劳，而是\BibleQuote{我们在奥秘中讲论天主的智慧，那隐藏的智慧}\BibleRef{格前 2:7}敬畏那火舌。\BibleRef{宗 2:3}也让我们的嗅觉得到治愈，使不成为柔弱的人；并且被撒上灰尘代替香膏，\BibleRef{依 3:34}却能\BibleQuote{闻到为我们倾倒的香膏}\BibleRef{歌 1:3}以属灵的方式接受它；并如此被它塑造和转化，以致从我们身上也能闻到芬芳。让我们洁净我们的触觉、味觉、喉咙，不要过于温柔地触摸它们，也不以光滑的东西为乐，而是以配得上祂——那为我们而降生成人的圣言——的方式来对待它们；并效法多默的榜样，\BibleRef{若 20:28}不以美味和调味料来纵容它们，那些是更为有害纵容的弟兄，而是品尝和学习主是美善的，以那更美好且持久的味觉；不要短暂地滋润那有害且忘恩的尘土，它让给它的一切过去，不能持守；而是以比蜜更甜蜜的话语来喜悦它。\par

XXXIX. 此外，在所说的一切之外，我们应以洁净的头（如同洁净感官之工坊的头）紧握基督的头，\BibleRef{弗 4:16} 全身都靠祂联络得合式，紧密结合；并要打倒那自高自大的罪，当它企图将我们高举于我们更好的部分之上时。肩头也当被圣化、被洁净，好能背起基督的十字架，这事并非人人都能轻易做到。手与脚也当被祝圣：手，为能在各处举起圣洁的手，\BibleRef{弟前 2:8}并把握基督的训诲，免得主有时发怒；且使圣言藉行为得着信服，如同先知手中所交付的言语一般；\BibleRef{盖 1:1}脚，为不使它们急速流人血，也不奔向邪恶，而能敏捷地奔向福音和那从上头召叫的奖赏，\BibleRef{斐 3:14}并迎接那洗涤洁净它们的基督。若能也洁净那收纳并消化圣言之食的肚腹，那也是好的：不使它因奢侈和那必朽坏的食物成为神，\BibleRef{若 6:27}反而要尽可能洁净它，使它更为瘦削，好能在最深处领受天主圣言，并为\Place{以色列}的罪过而体面地忧伤。\BibleRef{耶 4:19}我也发现心和内腑被尊为有荣耀。\Person{达味}使我相信这一点，当他祈求在他内造一颗洁净的心，在他的内腑更新正直的灵时，我想，他指的是心智及其运动或思念。\par

XL. 至于腰与肾，我们不可忽略这些。让洁净也抓住这些。我们的腰当用节制束紧并约束，如同古时法律吩咐\Place{以色列}吃逾越节羔羊时所行的。\BibleRef{出 12:11}因为没有人能纯洁地走出埃及，或逃脱那毁灭者，除非他约束了这些。肾也当藉那美好的皈依而改变，好使它们将所有情感转向天主，以致能说：\Quote{主，我的一切愿望都在祢面前，我不渴求人的日子}，\BibleRef{约 17:16}因为你必须成为一个有愿望的人，\BibleRef{达 10:11}但必须是属神的愿望。因为这样你就能摧毁那将大部分力量放在肚脐和腰上的龙，\BibleRef{约 39:16}杀死来自这些的力量。不要惊讶我更加尊荣那不体面的部分，\BibleRef{格前 12:23}藉我的言语抑制它们并使它们贞洁，反对肉身。让我们献给天主那在我们地上的每个肢体；\BibleRef{哥 3:5}让我们全部祝圣它们：不仅是肝叶或肾与脂油，也不是我们身体的部分，时而这，时而那（为何轻视其余呢？），而是让我们献上整个自己，成为合理的全燔祭，\BibleRef{罗 12:1}完美的祭献；我们也不只拿肩膀或胸脯作为司铎取去的份，\BibleRef{肋 7:34}因为那是小事，而是让我们献上整个自己，好能完整地收回自己；因为当我们将自己献给天主、以我们自己的救恩为祭献时，这就是完全地领受。\par

XLI. 除此之外，并超越一切，我恳求你持守那美好的寄托——我赖以生活和工作、并愿它成为我离世之伴侣的寄托；因着它我忍受一切痛苦，轻看一切快乐——那就是对圣父、圣子及圣神的宣认。今天我把它交托给你；以此我将为你授洗并使你成长。我给你分享这个，并愿你终生捍卫它：那在合一中的三位、又分别包含三位的唯一天主性及德能，在实体或本性上既无不等，也不因优越或卑下而增减；各方面平等，各方面同一，正如诸天的美丽与伟大是一般；三位无限者的无限结合，每一位独自被视为天主：圣父如何，圣子也如何，圣子如何，圣神也如何；三者同观则是一天主：每一位天主，因为同体；一位天主，因为那\OriginalTerm{独一统御}{Monarchia}。我刚想到那一位，就被三位的光辉照亮；我刚区分他们，就被带回那一位。当我想到三位中任何一位时，我就将其视为全部，我的眼睛充满，而我思考的大部分却逃离了我。我无法抓住那一位的伟大，以致能归给其余更大的伟大。当我同观三位时，我只看到一个火炬，不能分割或衡量那不可分的光。\par

XLII. 你害怕谈论\Quote{生}，恐怕你会归给那不能受苦的天主任何受动性吗？相反，我害怕谈论\Quote{受造}，恐怕我会因那侮辱和不真实的区分而毁灭天主：要么将圣子从圣父割裂，要么将圣神的实体从圣子割裂。因为这包含一个悖论：不仅一个受造的生命被那些拙劣衡量天主性的人塞进天主性里，而且这个受造的生命自身也彼此分裂。因为正如这些卑下世俗的思想使圣子服从圣父，同样，圣神的品位也被贬低至圣子之下，直到天主和受造的生命都被这新神学侮辱。不，我的朋友，在\OriginalTerm{天主圣三}{Trinity}内没有什么是奴仆，没有什么是受造的，没有什么是偶然的，正如我听见一位智者所说的。宗徒说：\Quote{如果我仍讨人的喜欢，我就不是基督的仆人了}；\BibleRef{迦 1:10}如果我仍敬拜一个受造物，或受洗归于一个受造物，我就不能成为神，也没有改变我最初的出生。对于那些敬拜\Person{阿斯塔特}或\Person{革摩士}——\Place{漆冬}人的可憎之物，或星辰的形像\BibleRef{亚 5:26}——一个对于这些拜偶像者略高于他们的神，却仍是受造物和手工作品的人，我能说什么呢？当我自己要么不敬拜那两位我受洗归于其联合之名者，要么敬拜我的同仆——因为他们是同仆，即使等级略高；因为在同仆之间必然存在差异。\par

XLIII. 我愿称圣父为更大，因为从他流出平等与平等者的存有（这众所公认），但我畏惧使用\Quote{根源}一词，唯恐我将他造成卑下者的根源，从而以尊荣的先后侮辱了他。因为出自他的那几位被贬低，并非根源的光荣。再者，我怀疑你贪得无厌，恐怕你抓住\Quote{更大}这个词，分割本性，在\Emph{一切}意义上使用\Quote{更大}，而它并不适用于本性，只适用于根源。因为在同体的诸位中，就实体而言，没有更大或更小。我愿在圣神之前尊圣子为子，但藉着圣神圣化我的圣洗，却不允许这样。但你害怕被指责为三神论吗？你当占有这善事，即三位中的合一，而让我去作战。让我作造船者，你使用船；或若别人造船，你当我是房屋的设计师，你安全地住在里面，尽管你没有在它上面劳苦。你的航行不会比我造船者更不顺利，你的居所也不会比我建造者更不安全，因为你没有在上面劳苦。看我的宽容多么大；看圣神的良善；战争属于我，成就属于你；我将置身火线，你将在平安中生活；但请与你的保护者一起祈祷，并藉着信德把你的手给我。我有三块石头，要向培肋舍特人投掷；\BibleRef{撒上 17:49} 我有三次吹气，对付撒勒法妇人的儿子；\BibleRef{列上 17:21} 我要以之复活那被杀者；我有三股洪水，对付那木柴，我要以水祝圣祭献，引发最意外的火；我要藉着圣事的能力，推倒那些可耻的先知。\par

XLIV. 我何须再多言？现在是教导的时候，不是争辩的时候。我在天主和蒙选的天使面前作证，\BibleRef{弟前 5:21} 你要在这信德中受洗。如果你的心以不同于我教导的方式被书写，来，让我更改这书写；我不是这些真理的拙劣书写者。我写那刻在我自己心上的；我教导那从我起初直到这白发之年所受教并持守的。风险是我的；愿奖赏也是我的，作为你灵魂的引导者，藉着圣洗圣化你。但如果你已经正确准备，并印有那美好的铭文，你要保存所写的，在变化的时间中，对于不变的事保持不变。在更好的意义上效法比拉多的榜样；你正确被书写的人，效法那错误书写的人。对那些想以不同方式说服你的人，说：我所写的，我已经写了。\BibleRef{若 19:22} 因为我实在羞愧，如果那错误的仍坚持不变，而正确的却如此容易被弯曲；然而我们应从较坏的容易转向更好的，但从更好的转向更坏的则不可动摇。如果这样，依照这教导你来领受圣洗，看，我不会闭口不言，看，我将我的手借给圣神；让我们催促你的救恩。圣神热切，祝圣者准备，恩赐预备。但如果你仍在迟疑，不接受天主性的圆满，去找别人为你施洗——或者更确切地说，去溺死你：我没有时间切割天主性，使你在一瞬间因重生而死亡，使你既无恩赐也无恩宠的希望，却在如此短的时间内使你的救恩沉没。因为你从三位天主性中减去任何一点，就会推翻全部，并毁灭你自己的被造完美。\par

四十五、但也许在你的灵魂上还没有形成任何或善或恶的文字；而你们今天愿意被书写，并被我们塑造至完美。让我们进入云中。把你的心版给我；我将作你的梅瑟，尽管这样说是一句大胆的话；我要以天主的手指在其上写下一个新的十诫。\BibleRef{出 38:28} 我要在其上写下一种更短的救恩方法。如果有任何异端或无理性的野兽，就让它留在下面，否则它将有被真理之言砸死的危险。我要因父及子及圣神之名给你授洗，使你成为门徒；这三者有一个共同的名字，即天主性。你们将藉着显现\BibleRef{玛 28:19}和言语知道你们弃绝一切不敬虔，并与整个天主性合一。要相信世界上的一切，无论有形无形，都是由天主从无中创造，并受造物主的天主眷顾所治理，并将改变进入更美好的状态。要相信恶没有本体或国度，无论是无始的、自有存在的、或是天主所造的；它是我们和那恶者的工作，因我们的疏忽而临于我们，却不是来自我们的造物主。要相信圣子，永恒的圣言，由圣父在万世之前无体而生，在末世为了你而成为人子，由童贞玛利亚以不可言喻、无玷的方式诞生（因为哪里有天主，哪里就不会玷污，救恩由此而来），祂亲自同时是完整的人和完美的天主，为了整个受苦者，好把救恩赐予你的整个存有，摧毁了你所有罪的定罪：按其天主性是不受苦难的，按其所取的（人性）是可受苦难的；祂为你成为了人，正如你为祂而成为了神。要相信祂为了我们罪人被领到死亡；被钉十字架并埋葬，以至于尝了死味；并且祂第三天复活了，升了天，为要带你这个躺在低处的人与祂一同上升；并且祂要带着祂荣耀的显现再来，审判生者死者；不再有血肉，也不是没有身体，而是按照祂独自知道的一种更似神的身体的法则，为叫那些刺透祂的人看见祂\BibleRef{默 1:7}，另一方面作为天主而没有肉体。此外，要接受复活、审判和按照天主公义的天平所给予的赏报；并相信这光将赐给那些心灵洁净的人（即见到并认识天主），与他们的洁净程度相称，我们称之为天国；但对于那些统治官能患有盲目的人，则是黑暗，即远离天主，与他们在世上的盲目相称。然后，第十条，要在教义的基础上行善；因为没有行为的信德是死的，\BibleRef{雅 2:17}正如没有信德的行为也是死的。这就是可以公开的关于圣事的全部内容，且是不禁止众人听闻的。其余的，你要在教会内藉着天主圣三的恩宠学习；那些事你要密封稳妥地隐藏在自己内。\par

四十六、但我再宣讲一件事给你们。你们在领洗后即将站在大堂前的站礼，是未来光荣的预像。迎接你们的圣咏，是天上圣咏的前奏；你们将点燃的灯，是那里光照的圣事，我们将以此迎接新郎，闪耀着纯洁的灵魂，以我们信德的灯照耀，不可因疏忽而沉睡，免得祂若忽然来到，我们错过所等候的那一位；也不可无油、无灯、缺乏善工，免得被逐出新房。因为我看见这样的情况多么可怜。当喊声要求迎接时，祂要来到；那些明智的将迎接祂，他们的灯闪耀，油料充足；但其他人则太迟地向有油的人寻求油。祂要急速来到，前者与祂一同进去，后者却被关在外面，因为他们在准备中浪费了进门的时间；他们将痛哭，因为太迟才学到自己懒惰的惩罚，那时他们无论如何恳求也不能再进入新房，因为他们已因自己的罪将门关闭，以另一种方式效法那些错过婚宴之人的榜样：这婚宴是良善的圣父为良善的新郎所设的；一人因新婚的妻子，另一人因新买的田地，第三人因一对牛；他和他所得到的都是不幸，因为为了一点小事失去了大事。因为在那里没有傲慢的人，没有懒惰的人，也没有穿着污秽的破衣而不穿婚宴礼服的人，即使他们在这里可能以为自己配得在那里穿明亮的袍子，并暗自闯入，以虚妄的希望欺骗自己。然后呢？当我们进入后，新郎知道祂要教导我们什么，以及祂要如何与同祂一起进来的灵魂交谈。我想，祂会以更完善、更纯洁的教导与他们交谈。愿我们众人，无论是教师还是受教的，同享这些，在同一的基督我们的主内，愿光荣与王权归于祂，直到永远。阿们。\par

\SourceRecord{Translated by Charles Gordon Browne and James Edward Swallow.；Nicene and Post-Nicene Fathers, Second Series；Vol. 7.；Edited by Philip Schaff and Henry Wace.；Buffalo, NY: Christian Literature Publishing Co.,；1894.；<http://www.newadvent.org/fathers/310240.htm>.}

% structure: p0001 source=h1 target=section reason=page-title

\section{\WorkTitle{讲道 41}}

\Emph{论五旬节}\par

\Emph{不确定下列讲道属于哪一年，但它确实是在君士坦丁堡宣讲的；本笃会编辑认为是在381年，若是如此，则该日为5月16日。一个有助于确定此日期的迹象见于第14章，其中表达了对自身因勇敢阐述真信仰而面临人身危险的忧虑。事实上，在该年年初，德奥多西皇帝将亚略派德莫斐鲁斯逐出并废黜后，使作者就任宗主教宝座，他险些在亚略派手中遇刺。}\par

本篇讲道再次处理第五篇神学讲道的主题，即圣神的天主性问题，但采用与前一篇演讲完全不同的论据来确立这一点，前一篇演讲的要点在此处不再重复。\par

讲道者首先评论犹太人、外邦人和基督徒庆祝节日的不同方式。然后，他论述数字七的神秘意义，并以若干实例加以说明；接着进入其主要主题。\par

天主的圣神，他说，完成了基督的工作。那些视祂为受造物的人，如同马其顿派的追随者，因此犯了亵渎和不敬之罪。真信德承认祂是天主；这一信仰对于救恩是必要的；但在用此名号称呼祂时需有所保留。我们确实必须坚持承认祂拥有天主性的全部属性；并且我们必须至少容忍那些像讲道者本人一样也给予祂天主之名的人，他希望所有听众都能从圣神获得恩宠如此行。然后，他继续从圣经证明，事实上天主性的所有属性都确实属于圣神；而祂独特的位格标记是：祂既不像圣父那样非受生，也不像圣子那样受生。他没有触及双重发出的问题。\par

从第8章中的一些表述来看，这篇讲道似乎不是向他一贯的听众，而是向一个\Quote{修道者}大会宣讲的。\par

这篇讲道的标题在不同手稿中有所不同。因此有些手稿标有\Quote{同一位论五旬节}，另一份手稿增加了\Quote{并论圣神}；还有一份手稿写为\Quote{同一位，五旬节讲道}。本笃会之前的印刷版为\Quote{论圣五旬节}。\par

一、让我们稍微推理一下这庆节，好使我们属神地庆祝它。因为不同的人有不同的庆节方式；但对于圣言的敬拜者而言，讲道似乎是最好的；而在讲道中，最合时宜的讲道是最好的。在一切美的事物中，没有比使庆节的爱好者属神地庆祝它更让美之爱慕者喜乐的了。让我们如此思考。犹太人和我们一样庆祝庆节，但只是在字面上。因为他们虽追求肉体的法律，却没有达到属神的法律。希腊人也庆祝庆节，但只是在肉体上，并且是为了尊崇他们自己的神灵和魔鬼，其中有些据他们自己承认是激情的创造者，另一些则是出于激情而被尊崇的。因此，他们庆祝庆节的方式也是充满激情的，仿佛他们的罪本身就是对天主的尊崇，而他们的激情在天主那里找到避难所，引以为荣。我们也庆祝庆节，但我们是按圣神所喜悦的方式庆祝。圣神喜悦我们通过履行某种义务，无论是行动还是言语，来庆祝庆节。因此，我们庆祝庆节的方式是，在我们的灵魂中珍藏那些永久且将与之紧密结合的事物，而不是那些将离弃我们并被毁灭、只暂时刺激我们感官的事物；而这些事物，按我的判断，至少大多是有害和毁灭性的。因为身体的恶已经足够了。那火还需要更多的燃料吗？或者那野兽还需要更丰富的食物，使它更难以控制，对理性过于猛烈吗？\par

二、因此，我们必须在精神上庆祝佳节。这是我言论的开端；因为即使我的言语似乎有点过于散漫，也必须说出来；为了爱好学问的人，我们必须勤勉，好像在我们的庄严庆节中掺入一些调味品。希伯来人按照梅瑟的律法尊崇七数（后来毕达哥拉斯派也尊崇四数，他们甚至习惯以四数起誓，就如西蒙派与马西昂派以八数和三十数起誓一样，因为他们以这些数的永世体系命名并加以敬拜）；我说不出这数是根据什么类比、或由于什么力量；反正他们尊崇它。有一点是明显的：天主在六天内创造物质，给它形式，并以各种形体和混合安排它，造了这可见的世界，第七天就停止了所有工作，正如安息日本身名称所示，在希伯来文中意为休息。然而，若有比这更高深的理由，就让别人去讨论吧。但他们尊崇七数不仅限于日子，也延伸到年份。属于日子的，由安息日证明，因为他们持续遵守它；与此相应，除酵也持续那么多天。\BibleRef{出 12:15} 属于年份的，由第七年，即豁免年表明；它不仅有七周期，还有七周期的七周期，同样用于日子和年份。日子的七周期产生五旬节，在他们中被称为圣日；年份的七周期产生所谓的喜年，其中也有土地的豁免、奴隶的解放和所购财产的归还。因为这民族不仅将后代和首生子的初果奉献给天主，也将日子和年份的初果奉献。如此，对七数的敬礼也产生了对五旬节的敬礼。因为七乘七产生五十除去一天，我们从将来世界借用这一天，它既是第八日又是第一日，或者说是唯一且不可毁灭的。因为我们灵魂现行的安息在那里才能找到终止，好将一份给予七，也给予八（我们的一些前辈就是这样解释撒罗满的这段话的）。\par

三、关于七数的尊荣有许多证据，但我们只从众多中取少许。例如，提到七种宝贵的神恩（\OriginalTerm{神恩}{spirits}）；因为我认为依撒意亚\BibleRef{依 11:2}喜欢称圣神的活动为神恩；主的训言，按达味所说，被纯净七次；义人从六次患难中被救出，第七次不受打击。\BibleRef{约 5:19} 但罪人得赦免，不是七次，而是七十个七次。\BibleRef{玛 18:22} 我们也可从反面看到（因为邪恶的惩罚当受赞美）：加音受报复七倍，即因杀弟受罚；拉默客七十个七倍，\BibleRef{创 4:24} 因为他在律法和谴责之后仍行凶杀。邪恶的邻人将七倍收回怀中；智慧之家立于七根柱石上，\EditorialNote{箴 9:1}则鲁巴贝耳的石头镶有七眼；\BibleRef{匝 3:9} 天主一日受赞美七次。再者，不生育的生了七个，\BibleRef{撒上 2:5}这是圆满的数，她与那子女不完美的相反。\par

IV. 若我们也必须考察古代历史，我见\Person{哈诺客}\BibleRef{犹 1:14}——我们祖先中的第七位——因被提升而受尊荣。我也见第二十一位\Person{亚巴郎}\BibleRef{创 5:22}因一个更大奥秘的加入，而获赐宗主教职的光荣。因为三次重复的七日便得出此数。一个极大胆的人甚至会冒险论及新亚当——我的天主和主\Person{耶稣基督}，祂按\BibleBook{路加}{福音}的倒叙家谱，被算作从坠入罪恶的旧亚当起的第七十七代。\BibleRef{路 3:34} 我又想起\Person{若苏厄}（\Person{农}的儿子）的七支号角，以及相同次数的绕行、日子和司铎，耶里哥城墙便因此倒塌。同样，绕城七次；正如\Person{厄里亚}先知的三次吹气——他借此使匝尔法特寡妇的儿子复活\BibleRef{列上 17:21}——以及他七次浇水于木柴，藉由天主降火烧尽祭献，并在他的挑战下定了那不能行同样之事的羞耻先知们的罪一样，都含有奥秘。还有那七次观望云彩的命令，加于幼仆身上；\Person{厄里叟}七次伸展自身于叔能女子的小孩身上，藉此伸展，生命的气息得以恢复。我认为，同一种教义也属于（若我略过圣殿的七枝七盏灯台）：司铎祝圣礼仪持续七天；\BibleRef{肋 8:33} 癞病人的洁净礼仪为七天，圣殿奉献礼\BibleRef{列上 8:6}也为同样的数目；第七十年人民从充军之地归来；\BibleRef{编下 36:32} 好让单元中的一切也出现在十数中，并使七日的奥秘在一个更完美的数目中受到敬重。但我为何谈论遥远的过去呢？\Person{耶稣}本身是纯全完善的，祂能在旷野中用五个饼喂饱五千人，又用七个饼喂饱四千人。他们吃饱后剩下的，前者是十二筐，后者是七筐；我想，这并非毫无理由或与圣神不相称。你若亲自阅读，可注意到许多数字含有的深意超越表面。但论及一个在此时对我们最有益的实例，并非因这些原因或其他非常相似、甚至更神圣的原因，希伯来人尊崇五旬节，我们也尊崇它；正如希伯来人的其他礼仪，我们遵守……它们曾被他们象征性地遵守，而为我们所圣事性地恢复。现在，关于这日子说了这么多作为前言，让我们继续我们要进一步说的。\par

V. 我们正在庆祝五旬节和圣神降临的节期，以及应许的指定时间和我们盼望的实现。这奥秘是多么伟大，多么庄严！基督身体的各项神恩计划已结束；或者更确切地说，属于祂肉身来临的那些（因为我犹豫是否称其为祂身体的神恩计划，只要没有论述说服我脱去身体是更好的），而圣神的计划正开始。那些属于基督的事是什么？\Person{童贞玛利亚}、诞生、马槽、襁褓、天使光荣祂、牧人奔向祂、星宿的运行、贤士朝拜祂并献礼物、\Person{黑落德}屠杀婴儿、\Person{耶稣}逃往埃及、从埃及返回、割损礼、圣洗、从天而来的见证、受试探、为我们被石头砸（因为祂必须被赐予我们作为忍受苦难为圣言的榜样）、被出卖、被钉、被埋葬、复活、升天；这些事中，至今祂仍在基督的敌人手中遭受许多侮辱；祂忍受着，因为祂是宽仁的。但从爱祂的人那里，祂接受一切光荣。祂延迟，正如前一种情形中的祂的愤怒，也如后一种情形中的祂的仁慈；在他们身上或许是为赐予他们悔改的恩宠，在我们身上是为考验我们的爱德：看我们是否在磨难\BibleRef{弗 3:13}和为真宗教的斗争中不灰心，正如从起初祂神圣救恩计划的秩序和祂不可测度的判断，祂以这些明智地安排一切关于我们的事。这些便是基督的奥秘。我们将看见后续的事更为光荣；愿我们也被看见。关于圣神的事，愿圣神与我同在，并赐我言语，如同我所愿；若非如此，也至少与这季节相称。无论如何，祂将作为我的主与我同在；并非以奴役的样式，也不等候命令，如有些人所想。因为祂随意向谁吹，向谁，到什么程度。\BibleRef{若 3:8} 这样我们得蒙灵感，思想并谈论圣神。\par

VI. 那些将圣神贬低到受造物等级的人是亵渎者和邪恶的仆人，且是邪恶者中最恶劣的。因为邪恶仆人的部分在于藐视主权、反叛统治，并使那自由者成为他们的同仆。而那些视祂为天主的人是受天主感召，并在心灵上卓著的；那些更进一步\Emph{称}祂为天主的，若是对善意的听者，则被高举；若是对低级者，则不够审慎，因为他们把珍珠投给泥土，把雷声给软弱的耳朵，把太阳给昏花的眼睛，把固体食物给仍用奶的人，\BibleRef{希 5:12}而他们本应逐步引导他们超乎自身，并把他们带到更高的真理；光照上再加光，真理上加真理。因此，我们将保留那更成熟的论述——时机未到——并与他们如下说话。\par

七、诸位朋友，若你们不肯承认圣神是非受造、永恒的；那么，这种心态显然是出于相反的神——恕我因热忱而说得稍显大胆。然而，若你们足够健全，能避开这明显的亵渎，并将赐予你们自由的那一位置于奴役之外，那么就请你们自己藉着圣神和我们，看看将有何后果。因为我深信你们多少已分享于祂，所以我将视你们为同道，与你们探讨这个问题。要么向我指出主与仆之间的某种中间状态，使我能将圣神的地位置于其中；要么，若你们羞于将奴役归于祂，那么你们所要寻求的对象应居何地位便毋庸置疑了。然而你们对字眼不满，对词语踟蹰，它成了你们的绊脚石、使人跌倒的磐石；因为对某些人而言，基督也是如此。这终究是人之常情。让我们以属灵的方式彼此相遇；让我们充满兄弟之爱，而非自爱。请将天主性的权能赐予我们，我们便将名称的使用权让给你们。用你们更为敬畏的其它词语来承认这本性，而我们将如医治病人般治愈你们，悄悄拿走你们所喜爱的某些东西。因为可耻，实在可耻且完全不合逻辑，当你们灵魂健全时，却对声音纠缠不休，并隐藏宝藏，仿佛嫉妒别人，或害怕连自己的舌头也被圣化。但更可耻的是我们处于我们所指责你们的状态，一边谴责你们对词语的吹毛求疵，一边自己却在字母上吹毛求疵。\par

八、诸位朋友，请承认天主圣三具有同一的天主性；或若你们愿意，承认同一的本性；我们将祈求圣神赐予你们\Quote{神}这个字。我深知祂会赐予你们，因为祂已经赐予了你们头一份和第二份；尤其如果我们所争论的只是某种属灵的胆怯，而非魔鬼的异议。让我说得更清楚、更简洁：你们不要因我们更崇高的词语而向我们问责（因为这种上升与嫉妒无关），而我们也不会批评你们所达到的，直到你们经由另一条路被带到同一个安息所。因为我们寻求的不是胜利，而是赢得弟兄，他们因与我们分离而令我们痛心。我们将这点让与你们，因我们在你们身上确实见到某种生命的真理，你们在圣子方面是健全的。我们钦佩你们的生活，但不完全赞同你们的教义。你们既拥有圣神的事物，就再接纳祂本身，使你们不仅竞赛，而且正规地竞赛\BibleRef{弟后 2:5}，这正是你们得冠冕的条件。愿你们蒙恩获得这交谈的赏报，使你们完善地承认圣神，并与我们，且先于我们，宣告祂所应得的一切。是的，我甚至要大胆为你们这样做；我要说出宗徒的愿望。我是如此依恋你们，如此敬畏你们的行列、你们节制的色彩、那些神圣的集会、可敬的童贞、净化、彻夜的咏唱，以及你们对穷人、弟兄和客旅的爱，以致我甘愿从基督那里被受诅咒，甚至如同被定罪者一样受苦，只要你们能站在我们身边，我们一起光荣天主圣三。因为其余的人，我何必多言？他们显然已死（唯有基督能复活他们，祂以自己的权能复活死者），并因他们的教义而相连，却在地方上可悲地分离；他们彼此争吵，如同斜视的双眼看同一物件，差异不在视线而在位置——若我们只指控他们斜视，而非完全盲目。现在我已在一定程度上阐述了你们的立场，来，让我们回到圣神的话题上，我想你们现在会跟随我。\par

九、因此，圣神始终存在，现在存在，且永远存在。祂既无开端，亦无终结；而是与父和子一同永远地被排列和计数。因为任何时候父缺少子，或子缺少圣神，都是不合适的。否则，天主性将在其最伟大的方面被剥夺荣耀，因为它似乎像是事后才达到圆满的完成。因此，祂总是被分享，却不分享；成全，却不被成全；圣化，却不被圣化；使成神，却不被成神；祂自身永远与自身相同，与祂所排列的那几位相同；不可见，永恒，不可理解，不变，无性质，无量度，无形体，不可触摸，自我推动，永恒运动，拥有自由意志，自我权能，全能（尽管所有属于圣神的一切都归于第一因，就如所有属于独生子的一切一样）；生命与赋予生命者；光明与赋予光明者；绝对的善，善的泉源；正义的、君王的圣神；主，派遣者，分离者；自己圣殿的建造者；引领，照祂的旨意运行；分施祂自己的恩赐；义子的灵、真理的灵、智慧的灵、明达的灵、知识的灵、虔诚的灵、谋略的灵、敬畏的灵（这些都被归于祂）；藉着祂，父被认识，子被光荣；并且藉着祂，\Emph{唯独}祂被认识；同一等次，同一服务，同一敬拜，同一权能，同一成全，同一圣化。何须赘言？凡父所有的，子也都有，除了非受生；凡子所有的，圣神也都有，除了受生。而这两件事并不分开实体，就我所理解的，而是实体内的区分。\par

十、你是在努力提出反驳吗？好，我也在努力继续我的讲论。尊敬圣神之日；若你能稍加收敛你的舌头。现在是谈论其他舌音的时候——敬畏它们或惧怕它们，当你看见它们如同火焰。今天让我们以信理教导；明天我们可以讨论。今天让我们庆祝庆节；明天将有足够时间行为不端——前者奥秘地，后者戏剧性地；前者在教会内，后者在市场中；前者在清醒者中，后者在醉酒者中；前者如同强烈渴慕者所当行的，后者如同嘲弄圣神者所行的。既已结束了那外来的元素，现在让我们彻底装备我们自己的朋友。\par

十一、祂首先在天使和天上的能力中工作，就是在天主之后、围绕天主的那等。因为它们的完美和光明，以及它们难以或不可能被引向罪恶，并非来自别的源泉，而是来自圣神。其次，在圣祖和先知中，圣祖见过天主的异象，或认识祂，先知也预知未来，他们的主要部分由圣神塑造，并与未来之事相联如同现在，因为这就是圣神的能力。再其次，在基督的门徒中（因为我不提基督本身，圣神住在祂内，不是作为运作，而是作为平等者陪伴着），并且以三种方式，按照他们能接受祂的程度，在三个场合：在基督受苦难被光荣之前，在祂复活的被光荣之后，以及在祂升天（或说是恢复，或我们当称为上天）之后。现在第一个场合显示祂——医治病人、驱逐魔鬼，这些都是不能离开圣神的；复活后向他们呼吸（这显然是神圣的感动）也是如此；同样，我们现在纪念的火舌的分施也是。但第一个场合隐约地显示了祂，第二个更明确，这个现在更完全，因为祂不再仅以能量临在，而是我们可以说，实体性地，与我们交往，并住在我们内。因为正如圣子曾以肉体形态与我们生活——同样，圣神也应以肉体形态显现；并且在基督返回祂的地方之后，祂应降临到我们这里——\Emph{来临}，因为祂是主；\Emph{被派遣}，因为祂不是敌对的天主。因为这些话既显示了一致，也标明了单独的位格。\par

十二、因此，祂在基督之后来临，为使我们不缺一位护慰者；而是\Emph{另一位}护慰者，好让你承认祂的同等尊荣。因为这\Quote{另一位}一词标记了一个\OriginalTerm{另一个我}{Alter Ego}，一个同等主权的名字，而非不平等。因为我知道，\Quote{另一位}不是说不同种类，而是说同质的事物。祂以舌头的形态来临，因为祂与圣言有密切的关系。这些舌头是火的，也许因为祂的净化能力（因为我们的圣经知道一种净化之火，任何愿意的人都能找到），或者因为祂的本体。因为我们的天主是吞噬之火，是\BibleRef{希 12:20}烧尽恶人的火；\BibleRef{申 4:24}尽管你可能因\OriginalTerm{同质}{Consubstantiality}而陷入困难，再次对这些话吹毛求疵。舌头是分开的，因为神恩的多样性；它们停留（坐下）以表示祂在王权中在圣者间的安息，并且因为革鲁宾是天主的宝座。这事发生在楼房里（我希望我不显得过于冗长），因为那些应当领受的人要上升并被提升到地上；因为某些楼房也被神圣之水覆盖，天主的赞歌由此唱出。耶稣本人也在楼房中将圣体圣事的共融给了那些正在被引入更高奥迹的人，由此一方面显示天主必须降到我们这里，正如我知道祂古时对梅瑟所做的那样；另一方面显示我们必须上升到祂那里，并且这样天主与人之间会因尊荣的汇聚而发生共融。因为只要任何一方停留在自己的立足点上——一位在祂的荣耀中，另一位在他的卑微中——天主的慈善就不能与我们混合，祂的慈爱就无法沟通，并且有一道不可逾越的巨大深渊分开他们；这深渊不仅分开了富翁与拉匝禄及他所渴望的亚巴郎的怀抱，也分开了受造而变化的本性与那永恒不变的本性。\par

十三、先知们在以下段落中宣告了这一点：——\BibleQuote{上主的神临于我}；\BibleRef{依 61:1} 又有，七神将安息在他身上；并上主的神降下引导他们；以及知识的神充满贝匝肋耳，\BibleRef{出 31:3}他会幕的建造者；还有那激怒以民的神，\BibleRef{依 63:10}；并以火车火马接去厄里亚的神，\BibleRef{列下 2:11}且厄里叟以双倍份量寻求；以及达味被善良而君王般的圣神引导和坚固。这应许首先借约耳的口预言，他说：\BibleQuote{在末日，我要将我的神倾注在一切血肉上（就是一切信者），并倾注在你们的儿女上}，\BibleRef{岳 2:28}及其他；然后后来借耶稣，被祂光荣，并将光荣归给祂，如同祂被父光荣并光荣了父。\BibleRef{若 14:16}这应许是多么丰富！祂将永远住下，并与你们同在，无论是现在与那些在时间范围内堪当的人，还是将来与那些被算为堪当那世界的人，那时我们完全藉着此地的生活留住祂，并因我们犯罪而没有拒绝祂。\par

第十四。圣神与圣子一同施行创造与复活，你可由以下圣经得知：\Quote{因上主的一句话，诸天造成，一切万象因祂口中的气而成；}以及：\BibleQuote{天主的神造了我，全能者的气息教训了我。}\BibleRef{约 33:4}又说：\Quote{你将发出你的神，它们便被造成，你也将更新地面。}祂是神性重生的创始者。以下是你的证据：——\BibleQuote{除非由圣神重生，没有人能看见或进入天国，}\BibleRef{若 3:3}并从第一次出生（那是黑夜的奥秘）中被洗净，借着白昼与光的重塑，每个人都被重新创造。圣神是至智慧、至慈爱的，\BibleRef{智 1:6}若祂占据一个牧人，便使他成为圣咏作者，藉他的歌制服邪灵，\BibleRef{撒上 16:23}并宣告他为君王；若祂占据一个牧放山羊、采集野果的人，\BibleRef{亚 7:14}便使他成为先知。请回想\Person{达味}和\Person{亚毛斯}。若祂占据一个俊美少年，便使他成为长者的审判者，甚至超越他的年龄，正如\Person{达尼尔}所见证的，他在狮子的洞穴中征服了狮子。\BibleRef{达 6:22}若祂占据渔夫，便使他们以基督的网捕捉整个世界，将他们在圣言的网眼中提起。请看\Person{伯多禄}、\Person{安德肋}和雷霆之子，他们发出圣神之事如雷。若占据税吏，祂便使他们为门徒之道获利，使他们成为灵魂的商人；以\Person{玛窦}为证，昨日是税吏，今日是福音作者。若占据热心的迫害者，祂便扭转他们热心的方向，使他们成为\Person{保禄}而非\Person{扫禄}，正如祂发现他们满满作恶时一样，如今也满满虔诚。祂是溫良的圣神，却为犯罪者所激怒。因此，让我们借着承认祂的尊位，体验祂的温柔，而非祂的震怒；不要希望看见祂不妥协的怒火。也是祂使我今日成为向你们大胆宣告的传报者——若使我不得安宁，愿天主受感谢；但若有风险，同样感谢祂；一方面，好使祂宽恕那些仇恨我们的人；另一方面，好使祂圣化我们，在领受这宣讲福音的赏报时，得以借血而完美。\par

第十五。他们说起异语，不是本地的方言；这奇事甚大，未学语言的人竟能说出语言。这征兆是为不信的人，\BibleRef{格前 14:22}而不是为信的人，好成为不信者的控告，正如经上记载：\BibleQuote{我要藉其他舌音和其他口唇向这人民说话，即使如此，他们仍不听从我}\BibleRef{依 28:11}——\Person{主}说。但他们听见了。在此稍停，提出一个问题：你如何划分这话？因为这句话有歧义，要靠标点来决定。是他们各人听见自己的方言——容我这样说——一个声音发出，却有多方听到；空气如此被打动，并且可以这样说，发出的声音比原声更清晰？还是我们应在\Quote{他们听见}后停顿，然后将\Quote{他们用自己的语言说话}加在后面，这样的话，对于听者而言是说着他们自己的语言，但对于说者却是外语？我倾向于后一种解释；因为另一种理解，奇迹更在于听者而非说者；而在此，奇迹应在于说者；正是他们被责骂醉酒，显然是因为他们藉圣神行了舌音的奇迹。\par

第十六。但正如古代的舌音混淆是可赞美的，当人们用同一种语言行恶与不敬——甚至像今日有些人胆敢如此——建造巴贝耳塔时；\BibleRef{创 11:7}因为藉着混淆他们的语言，他们意图的合一被打破，他们的工程被摧毁；同样，如今的奇迹更值得赞美。因为从同一圣神倾注在许多人身上，使他们重新和谐。并且有神恩的多样性，需要另一种神恩来辨别哪个最好，虽然所有都值得赞美。那分裂也可称为高贵，达味说：\Quote{主啊，淹没他们，分裂他们的舌头。}为什么？因为他们喜爱一切淹没的话语，诡诈的舌头。他在这里几乎明言责斥当今分裂神性的舌头。关于这一点，就说到这里。\par

第十七。其次，既然舌音是对\Place{耶路撒冷}的居民、最虔诚的犹太人、帕提亚人、玛代人、厄蓝人、埃及人、利比亚人、克里特人、阿拉伯人、美索不达米亚人以及我自己的卡帕多细亚人说的，也是对犹太人（若有人更喜欢这样理解）——从天下各国聚集而来的人；值得看一看这些人是谁，以及他们属于何种俘虏。因为埃及和巴比伦的俘虏是有限的，而且早已因回归而结束；至于罗马人统治下的俘虏，那是因他们对我救主的妄为而施加的，尚未发生，虽然已在近期。那么，只能理解为\Person{安提约古}时期的俘虏，此事发生在不久之前。但若有人不接受这种解释，认为过于繁琐，因这俘虏既非远古，也未遍布天下，而希望找到一个更可靠的——或许最好的方式如下：民族曾被多次迁移，如\Person{厄斯德拉}所记载；有些支派被收回，有些则遗留下来；其中很可能（当他们散居在各民族中）有些人当时在场并分享了这奇迹。\par

十八、这些疑问，好学之士先前已经探究过，或许也并非没有理由；而现今任何人若再有所补充，他都将与我们同道。但现在我们有责任解散这场聚会，因为所言已足。但这节日永不会结束；如今我们确实以肉身守节，但稍后将在那里全然地以精神守节，在那里我们将更纯洁、更清晰地看到这些事的缘由，就在圣言本身、天主、我们的主\Person{耶稣基督}内，他是得救者的真节日与喜乐——愿光荣与朝拜归于他，偕同\Person{圣父}及\Person{圣神}，从今世直到永远。阿们。\par

\SourceRecord{Translated by Charles Gordon Browne and James Edward Swallow.；Nicene and Post-Nicene Fathers, Second Series；Vol. 7.；Edited by Philip Schaff and Henry Wace.；Buffalo, NY: Christian Literature Publishing Co.,；1894.；<http://www.newadvent.org/fathers/310241.htm>.}

% structure: p0001 source=h1 target=section reason=page-title

\section{第42篇演讲}

\Emph{最后告别辞——在150位主教面前}\par

\Emph{此篇演讲于公元381年在君士坦丁堡举行的第二届大公会议期间发表。历史与个人的动机使这一场合具有极深的意义。听众包括参加大公会议的东方教会150位主教，以及讲者自己的羊群，即君士坦丁堡的正统基督徒。他亲自努力聚集了这群羊，当时他们已被异端教师蹂躏。他以勇敢捍卫信仰、清晰教导和慈父般关怀属灵需要，赢得了成员的钦佩和爱戴。他虽非所愿，却在欢呼中被拥立为东方教会最高的教会职务，并被召主持聚集主教的公会议。当他发现自己无法就一个极其重要的问题引导公会议的审议，并察觉到一些主教将他本人及其地位当作新的纷争缘由时，他觉得必须辞去高位，并试图通过这一个人牺牲恢复教会的和平。他的言辞与这一场合相称。他必须处理引起纷争的话题，却以温和而明晰的分寸感加以处理；他在自我辩护中极其克制；他以最温柔的感情讲述自己与羊群的共同经历；他庄严清晰地作了最后的公开信仰阐述；最后，以优美的语言，用一位老者颤抖的声调，向他的羊群、他的主教座堂和他的宝座，以及所有令人动情的回忆，温柔告别。这一场合的悲怆在历史上无与伦比。演讲者和听众都深受感动，这种情感在所有读过他言辞并想象当时场景的人心中重新燃起。}\par

1. 亲爱的牧人们和同牧者，你们如何看我们的事？你们的脚是何等美丽，因为你们带来和平的喜讯，以及你们所带来的福事；在你们眼中，你们的脚又是美丽的——你们适时来到我们这里，不是为了寻找迷途的羊（\BibleRef{玛 18:12}），而是为了与一位朝圣的牧人交谈。你们如何看待我们这次朝圣？以及它的果实，毋宁说是那在我們内的圣神（\BibleRef{迦 5:22}）的果实（\BibleRef{弟后 1:14}）？我们一直被祂所感动（\BibleRef{宗 17:28}），尤其是现在，渴望拥有，或许确实拥有，没有什么是属于自己的。你们是自己明白并觉察，并且是我们行动的公道评判者吗？还是我们必须像那些被要求就其军事指挥、民政管理或国库行政作出交代的人一样，公开并亲自向你们呈报我们管理的账目？当然，并非我们羞于被审判，因为我们自己也同样是审判者，两者都基于同样的爱德。但这法律是古老的：因为甚至连保禄也将自己的福音传达给了宗徒们（\BibleRef{迦 2:2}）——不是出于炫耀，因为圣神远避一切炫耀——而是为了确立他的成功并纠正他的失败，如果他的言行中确实有任何失败的话，正如他在论及自己时所宣称的那样。因为就连先知的神魂也服从先知（\BibleRef{格前 14:32}），按照那妥善治理和分配万有的圣神的秩序。你们不要惊奇，他私下向一些人交账，而我却公开向所有人交账。因为我的需要比他更大，需要借助批评者的自由，如果我在责任上有所缺失，以免我白白奔跑，或已经白跑了（\BibleRef{迦 2:2}）。而唯一可能的自辩方式，就是在了解事实的人面前发言。\par

2. 那么，我的辩护是什么？（\BibleRef{格前 9:3}）如果它是虚假的，你们必须定我的罪；但如果是真实的，你们，我为之并当着你们说话的人，必须为此作证。因为你们就是我的辩护、我的见证，以及我喜乐的冠冕（\BibleRef{得前 2:19}），如果我也敢稍微借用宗徒的话夸口的话。这群羊，当它又小又穷，就外表而言，甚至不成其为羊群，而是羊群的一丝痕迹和遗物，没有秩序、没有牧人、没有边界，既无权拥有牧场，也没有羊圈的保护，流落在山上和地上洞穴与岩窟中（\BibleRef{希 11:38}），四处分散飘零，各自寻找栖身之所或牧场，或感恩地保全自己的性命；就像那被狮子骚扰、被暴风驱散、或在黑暗中流散的羊群，先知们哀叹它，将其比作以色列的厄运（\EditorialNote{厄则克耳 31:ii}），被交在外邦人手中；我们也曾为此哀叹，只要我们的命运仍值得哀叹。因为事实上，我们同样被驱逐、被抛弃，散落在每座山上和每个丘陵上，因缺乏牧人；一场可怕的暴风雨降临在教会，凶猛的野兽袭击她，即使在风平浪静之后，现在它们也不放过我们，毫不羞愧，行使着比时间应许更大的权势；而一片阴暗的黑暗，远比埃及的第九灾——那可触摸的黑暗（\BibleRef{出 10:21}）——更加沉重，笼罩并隐藏了一切，使我们几乎彼此都看不见。\par

3. 为了更动情地说话，我信赖那曾离弃我的祂，如同信赖一位父亲，\Quote{亚巴郎不认识我们，以色列不承认我们，但祢是我们的父亲，我们仰望祢；\EditorialNote{依撒意亚 63:16}在祢以外我们不认识别的，我们提说祢的名。}因此，耶肋米亚说：我要与祢争辩，我要与祢理论。\BibleRef{耶 12:1}我们变得如同起初，那时祢没有治理我们，\EditorialNote{依撒意亚 63:19}祢忘记了祢的圣盟约，向我们将祢的仁慈关闭。因此我们，崇拜天主圣三，完美神性的完美恳求者，成了祢所爱者的羞辱，既不敢将任何超乎我们的事物降低到我们的水平，也不敢如此起来反抗那些与天主作战的亵渎之舌，以至使祂的尊威与我们同列；但显然，我们因其他罪恶而被交出，因为我们的行为不配祢的诫命，我们随从自己邪恶的心思而行。因为，还有什么其他理由使我们被交与地上所有居民中最不义、最邪恶的人呢？先是拿步高折磨了我们，\BibleRef{耶 51:34}他在基督徒时代怀着反基督的愤怒，只因借着基督获得救恩就恨基督，并将圣书换为对非神的祭献。他吞噬了我，撕裂了我，一片轻微的黑暗笼罩了我，即便在哀悼中我也愿保持圣经的语言。若不是主帮助了我，并公义地将祂交给了无法无天之人的手，弃绝了祂（这些都是天主的审判）给波斯人，他们的血公义地流了，因为祂不圣洁的血，唯有在此情况下，公义甚至不能容忍迟延，我的灵魂早已居于坟墓。第二个并不更仁慈，甚至更令人痛苦，因为他虽带着基督之名，却是假基督，同时是基督徒的负担和羞辱，因为顺从他是不敬神的，在他手下受苦是不光荣的，因为他们甚至不显得受冤，也未因苦难获得殉道者的荣耀称号，因为真理在此被歪曲，他们虽作为基督徒受苦，却被当作异端者受罚。唉！我们在不幸中是何等富有，因为烈火吞噬了世间的美物。\BibleRef{岳 1:19}蝻子所剩的，蝗虫吃了；蝗虫所剩的，蚂蚱吃了；然后来了蠹虫，然后我不知道还有什么，一个灾祸接踵而来。但我为何要悲怆地描述这时代的灾祸，以及对我们的惩罚，或者更恰当地说，我们所受的考验和净化？无论如何，我们经过了火与水，并因天主我们救主的恩惠，达到了一个安息之处。\par

4. 回到我最初的起点。这是我的田地，它当初又小又贫瘠，不仅不配天主，祂一直并用美丽的种子和虔敬的教理培育整个世界，而且显然也不配任何贫乏困苦、资财微薄的人。不，它甚至不配称为田地，不需要谷仓和打谷场，甚至不配镰刀；没有禾捆或禾束，只有小而不合时的禾束，像屋顶上的禾束，不充满收割者的手，也不叫过路的人给予祝福。这就是我的田地，我的收获；在观看隐密者的眼中，它广大、饱满、肥美，成为这样一位农夫应有的样子，它的丰裕来自那用圣言善耕的灵魂谷地；然而在公开场合不被认识，没有被聚集，而是零散地收集，像在残茬中拾起的麦穗，像在葡萄采收中拾起的葡萄，那里没有留下整串。我想我可以恰当地加上一句：我找到以色列如同旷野中的无花果树，又像未熟葡萄串中的一两颗熟葡萄，被保存为主所赐的祝福，\EditorialNote{依撒意亚 65:8}和一份奉献的初果，虽然尚小且稀少，不满吃者的口；又像山上的旗帜，山巅的火炬，或其他只被少数人看见的孤寂之物。这就是它先前的贫乏和卑微。\par

5. 但是，那使人穷也使人富，杀死也使人复生，制造并改变一切，将黑夜变为白昼，\BibleRef{亚 5:8}冬天变为春天，风暴变为平静，干旱变为大雨，且常因\BibleRef{列上 18:42}一位备受迫害的义人\BibleRef{雅 5:16}的祈祷而行此事的天主；那将谦卑者高举，将恶人打倒在地的天主；天主对自己说：我确实看见了以色列的困苦；\BibleRef{出 3:7}他们不再因泥土和砖块受折磨；祂一说话就来眷顾，在眷顾中拯救，用大能的手和伸展的臂，借着祂所拣选的梅瑟和亚郎，领出了祂的子民——结果如何，行了何等奇事？那些经书和纪念碑所记载的事。因为，除了路上的所有奇迹和那巨大的吼声（简言之），若瑟独自进入埃及，\BibleRef{创 37:28}不久后六十万人离开埃及。\BibleRef{出 12:37}还有什么比这更奇妙？还有什么更能证明天主的慷慨，当祂愿意从无资财的人那里为公共事务提供资财？应许之地通过一位被恨恶的人分配，那被出卖的\BibleRef{创 49:22}赶出了众民族，自己成为一个大民族，那小小的嫩枝变为茂盛的葡萄树，\BibleRef{欧 10:1}如此巨大，以至伸展到河流，延伸到海洋，从边界到边界蔓延，用其荣耀的高度隐藏了群山，高耸于香柏之上，甚至天主的香柏之上，不论这些山和香柏意指什么。\par

6. 这就是从前那羊群，现在也是如此：如此健康、茁壮，即使尚未达于完美，它也因不断增长而向前迈进，我预言它将前进。这是圣神告知我的，若我有任何先知性的直觉和对未来的洞察。藉着以前的事，我能充满信心，并藉推理认出这一点，因为我是理性的养育者。因为从那种状态达到如今的发展，比从现在达到荣耀的顶峰，更不可能得多。因为自从它开始被那使死者复活的祂，\BibleRef{罗 4:17}骨与骨相连，关节与关节相接，生命与重生的神在它们的枯干中赐给了它们，\BibleRef{则 37:7}它的整个复活，我很清楚，必定要应验：好叫悖逆者不得自高，那些抓住影子的、或醒时之梦的、或散去的微风的、或水中船迹的人，不以为他们拥有什么。叹息吧，柏树，因为香柏树倒了！\BibleRef{匝 11:2}让他们从别人的不幸受教，学习穷人不会永远被遗忘，天主性不会，如哈巴谷所说，在愤怒中击穿有权者的头颅\BibleRef{哈 3:13}——这被击穿并被不虔敬地分为统治者与被统治者的天主性，为以贬低它而极大侮辱天主性，并藉与天主性平等来压迫受造物。\par

7. 我确实似乎听到那声音，来自那聚集破碎者、欢迎受压迫者的祂：拓展你的绳索，向右向左扩展，打入你的橛子，不要吝惜你的帷幔。\EditorialNote{依撒意亚 54:2}我放弃了你们，但我将帮助你们。在微小的怒气中我击打了你们，但以永远的仁慈我必荣耀你们。祂的慈爱之量超过祂惩戒之量。前者是因我们的邪恶，后者是因当受钦崇的天主圣三：前者是为我们的洁净，后者是为我的荣耀，我要荣耀那些荣耀我的人，\BibleRef{撒上 2:30}并激动那些激动我的人。看，这事与我一同密封，这是不可解除的报应之律。但你们以墙壁、石版、丰富镶嵌的石头、长长的柱廊和走廊包围自己，闪耀发光如金子，部分如水倾泻，部分如沙堆积；不知晓有信德、除了天别无屋顶遮盖的，强于在财富中翻滚的不虔敬，而三位因主名聚在一起\BibleRef{玛 18:20}在天主眼中比数以万计否认天主性的人更宝贵。你们更愿要整个客纳罕人而不要\Person{亚巴郎}一人？或索多玛人而不要\Person{罗特}？或米德杨人而不要\Person{梅瑟}，\BibleRef{出 2:15}那时他们每一人都是旅客和异乡人？\Person{基德红}那三百勇敢舔水的勇士，\BibleRef{民 7:5}与被打败的数千人相比如何？或者\Person{亚巴郎}的仆役，人数略多于他们，与许多君王和数万大军相比，他们人数虽少，却追上并击败了他们？\BibleRef{创 14:14}或者你们如何理解这段经文：即使以色列子民的数目如同海沙，但余民将得救？再者，我为自己留下七千人，未曾向巴耳屈膝。情形不是这样吗？不是这样吗？天主不以数目为乐。\par

8. 你们数算数以万计，天主数算那些处于救恩状态的人；你们数算无法计数的尘土，我数算蒙拣选的器皿。因为在天主眼中，没有比纯正教义、以及完全拥有真理一切教义的灵魂更尊贵的了。——因为没有任何事物堪当那创造万有、藉祂万有而有、并为祂万有而有的祂，\BibleRef{格前 8:6}可以给予或献给天主：不仅任何个人的手工或工具，即便我们愿意藉联合全人类的所有财产和手工来尊崇祂。我岂不充满天地？\BibleRef{耶 23:24}主说！你们要为我建造什么房屋？或我安息之所是什么呢？\EditorialNote{依撒意亚 66:1}但是，既然人难免不及那堪当的，我向你们要求那最接近的，就是孝爱虔诚，这财富是众人共有的、在我眼中同等，其中最贫穷者若胸怀高尚，可超越最显赫者。因为这种荣耀取决于意图，而非财富。你们要确信，我必从你们手中接纳这些。你们不得进入我的庭院，惟有温良者的脚要踏入，他们恰当地、真诚地承认了我、我的独生圣言和圣神。你们要继承我的圣山到几时？我的约柜在外邦人中要到几时？\BibleRef{撒上 6:1}如今你们稍多耽溺于属于他人的事物，满足你们的欲望。因为正如你们设计拒绝我，我也要拒绝你们，\BibleRef{欧 4:6}全能的上主说。\par

9. 我似乎听到祂这样说，看到祂这样做，并且还听到祂向祂的子民呼喊——那子民从前稀少、分散、可怜，如今却变得众多、相当团结且令人羡慕——\Quote{穿过我的门，你们要宽广。}\EditorialNote{依撒意亚 62:10}你们必须永远处于困境、住在帐棚里，而烦扰你们的人却大大欢乐吗？又向那些管理的天使——因为我相信，正如\Person{若望}在他的\WorkTitle{默示录}中所教导我的，每一个教会都有它的守护天使，\BibleRef{默 2:1}——\Quote{为我的子民预备道路，将石头从路上除掉，}\EditorialNote{依撒意亚 62:10}使人民在神圣的道路和入口上没有绊脚石或障碍，如今是通往人手所造的圣殿，\BibleRef{宗 7:48}但不久之后，通往天上的耶路撒冷，\BibleRef{迦 4:26}和那里的至圣所——我知道，这将是那些在此勇敢行走道路之人的痛苦和奋斗的终结。在他们之中，你们也被召为圣人，\BibleRef{罗 1:6}是特选的子民、王家的司祭，\BibleRef{伯前 2:9}是主最优越的一份，从一滴水而成的整条河流，从火星而成的天上明灯，从芥菜籽而成的大树，\BibleRef{玛 13:21}飞鸟都来栖息其上。\par

10. 这些我们呈献给你，亲爱的牧者们，这些我们献给你，用这些我们欢迎我们的朋友、宾客和同行的朝圣者。我们没有更美丽或更辉煌的东西可以献给你，因为我们已经选了所有我们财宝中最大的，好让你看见，虽然我们是外人，我们却一无所缺，虽然贫穷，却使许多人富足。\BibleRef{格后 6:10}如果这些事是微小而不值得注意的，我倒想知道什么是更大、更重要的事。因为，如果用健全的道理建立并巩固了一座城市——这城市是宇宙的眼睛，拥有极大的海陆力量，如同东西岸之间的纽带，世界的各个极端从四面八方在此汇合，又从这里，作为信仰的共同市场，而发源，一座被如此多语言的漩涡水流抛来抛去的城市——若这不是什么大事，那么任何事都很难被认为是伟大或可敬的。但如果这确实值得赞扬，那么就因这个缘故允许我们有一些荣耀，因为我们对你所看到的这些结果贡献了一部分。\par

11. 举目向四周观看，且看，\EditorialNote{依撒意亚 60:4}你，我话语的批评者！看那为\Place{厄弗辣因}的佣工所编织的冠冕，和狂傲的冠冕；看长老们的集会——他们因年岁和智慧而受尊敬，执事们的美好秩序——他们离同一个圣神不远，读经员的良好品行，人民对教导的热忱，无论男女，在德行上同样著名：男人们，无论是哲学家或普通百姓，都在神圣之事上同样明智；无论是统治者或被统治者，在这方面都适当地从属于规律；无论是军人还是贵族、学生还是文人，都是天主的士兵\BibleRef{弟后 2:3}，虽然在其他方面温和，却准备好为圣神争战，所有人都尊敬上面的集会——我们不是靠单纯的文字，而是靠赐生命的神进入那集会，所有人实在都是理性的人，是那真正是圣言的祂的崇拜者；女人们，如果已婚，则由神圣而非血肉的纽带结合；如果未婚而自由，则完全奉献给天主；无论年轻或年老，有些人光荣地迈向老年，有些人热切地努力保持不朽，被最美好的希望所更新。\par

12. 对那些编织这冠冕的人——我所说的，不是按主说的，\BibleRef{格后 11:17}但我仍要说——我也提供了帮助。他们中有些是我言语的结果，不是我们随意说出的言语，而是我们所爱的言语；也不是那些卖弄的言语——虽然娼妓的言语和态度曾被诽谤地归于我——而是那些最庄重的言语。他们中有些是我圣神的后代和果实，因为圣神能生养那些超越身体的人。对此，我毫不怀疑你们中那些仁慈的人，不，你们所有的人，都会作证，因为我一直是所有人的农夫：我唯一的报酬就是你们的宣认。因为我们现在没有，过去也没有任何其他目的。因为德行，为了保持其为德行，是没有报酬的，它的眼睛只注视那善。\par

13. 你愿意我说一些更大胆的话吗？你看见敌人的舌头变得温和，那些攻击天主性而反对我的人平静了吗？这也是我们圣神的成果，是我们耕耘的结果。因为我们在运用纪律时并非没有纪律，我们也不像许多人那样投掷侮辱——他们不攻击论点而攻击发言者，有时试图用谩骂来掩盖他们推理的弱点；正如乌贼被说成在它们面前喷出墨汁，为了逃脱追捕者，或者在不被察觉时自己捕猎他人。但我们表明，我们的争战是为了基督，我们像基督一样争战，祂是温良和善的，\BibleRef{玛 11:29}祂担当了我们的软弱，祂曾争战。虽然温和，我们却不因让步些许真理之言而损害它，以获得合理性的名声；因为我们不以恶行追求善：我们之所以温和，是因为我们争战的合法性，它局限于我们自己的范围和圣神的规则。关于这些要点，这是我的决定，我为所有灵魂的管家和圣言的分施者制定这条律则：既不要因严厉而激怒他人，也不要因顺从而让他们傲慢；而是要在讨论圣言时使用良善的言语，并且两方面都不要越过中庸之道。\par

14. 但你们也许渴望我尽可能阐述信仰。因为我将因回忆的努力而圣化，百姓也将因这类讨论的特殊喜悦而受益，你们将完全承认这一点——除非我们遭到无端的嫉妒，如同那些在真理的彰显上与我们所不及之人竞争的人。因为，如同深水，有些在深处完全隐藏，有些遇到阻碍泛起泡沫，在破碎前犹豫片刻（如它们向我们的耳朵所承诺的），有些则真的破碎了；同样，在神圣哲学的教导者中——撇开完全误入歧途者——有些将他们的虔诚完全秘密地隐藏在自己内，有些离产痛不远，避免不虔敬，却未说出他们的虔诚，或因教导上的谨慎保留，或因恐惧的压力，他们自己，如他们所说，心智健全，却未使他们的百姓健全，仿佛他们受托管理自己灵魂的治理，而非他人的；而另一些人公开他们的宝藏，无法克制自己不去生出他们的虔诚，并认为只拯救自己而不将福泽的溢出施予他人，那并非救恩。在这些之中，我将自己算上，以及所有与我一同高贵地敢于宣认真理的人。\par

15. 我们教导的一个简明宣告，一个众人可理解的题词，就是这个百姓，如此真诚地敬拜天主圣三，以至于他们宁愿将任何人从这生命中剪除，也不愿将三者中的一位从天主性中剪除：同心同德，同等热忱，并因教义的合一而彼此联合，与我们、与天主圣三联合。简略概述细节：那无始的、是元始的、与元始同在的，是一位天主。因为无始者的本性不在于无始或无生，因为任何事物的本性在于它是什么，而不在于它不是什么。它是肯定其所是，而非否定其所不是。而元始，因为它是元始，并不与那无始的分离。因为它的元始并非它的本性，正如无始并非另一位的本性。因为这些是本性的伴随物，而非本性本身。那与无始者和元始者同在的，也不是别的，就是他们所是的。现在，那无始者的名字是圣父，元始者是圣子，与元始同在者是圣神，三者有同一本性——天主。而合一在于圣父，\OriginalTerm{位格}{personae}的次序从他而出并归于他，不是混淆，而是拥有，没有时间、意愿或能力的分别。因为在我们中，这些造成了众多个体，因为每一个都与其他性质以及同一性质的其他个体拥有相分离。但对于那些本性单纯、本质相同者，\Quote{一}这个词以最高意义属于他们。\par

16. 那么，让我们告别一切争辩性的摇摆和平衡真理，既不似撒伯里乌派那样为了合一而攻击天主圣三，以邪恶的混淆破坏区别；也不似亚流派那样为了三而攻击合一，以不虔敬的区别推翻一体。因为我们的目标不是以恶易恶，而是确保我们获得善。这些是恶者的玩物，他不断将我们的命运导向邪恶。但我们行走于两端之间的王道，那是德行所在，如权威们所言，相信圣父、圣子和圣神，同一实体和光荣；在祂们内，圣洗也达到圆满，名义上和实际上（你知道你已受启蒙！）是否定无神论和承认天主性的；这样我们重生了，承认本质和不可分割敬拜中的合一，以及\OriginalTerm{位格}{\GreekText{ὑποστάσεις}}或\OriginalTerm{位格}{personae}中的三（有些人更喜欢这个词）。让那些在这些点上争辩的人不要发出可耻的嘲笑，仿佛我们的信仰依赖于词汇而非实在。因为你们主张三个位格的人是什么意思？你们用这个词暗示三个本质吗？我确信你们会大声反对这样做的人。因为你们教导三者的本质是同一的。你们主张三个位格的人是什么意思？你们想象一个单一的复合存在，有三个面，或完全人形？愿这思想消失！你们也会大声回答说，这样想的人将永远看不到天主的容貌，无论它是什么。一方所说的\OriginalTerm{位格}{\GreekText{ὑποστάσεις}}和另一方所说的\OriginalTerm{位格}{personae}是什么意思？进一步问：它们是三个，按性质区分而非按属性区分。好极了。人们怎能比你们更一致和谐，即使你们使用的音节有差异？你们看，我是一个多么好的调解人，使你们从文字回到意义，如同我们对旧约和新约所做的那样。\par

17. 然而，容我回过来：让我们谈论那 \OriginalTerm{不生者}{Unbegotten}、\OriginalTerm{受生者}{Begotten} 与 \OriginalTerm{发出者}{Proceeding}，如果有人喜欢造这些名称；因为我们不必害怕把形体观念加到那无形体者身上，如同那些诋毁 \OriginalTerm{天主性}{Godhead} 的人所想的那样。因为受造物必须被称为天主的，这对我们而言是件大事，但永不能称为天主。否则，我将承认天主是受造物，如果我成为天主——在严格意义上。因为这是真理。如果是天主，祂就不是受造物；因为受造物与不属于天主的我们同类。如果是受造物，祂就不是天主，因为祂在时间上有开始。并且那开始者曾有一个时间不存在。而那不存在是其先前状态的，就不是严格意义上的存有。那严格来说不存在的，如何能是天主？因此，三者中没有一个是受造物，更糟的是，也没有一个是为我而存在的；因为若是如此，祂不仅是受造物，而且在尊荣上低于我们。因为，如果我是为天主的荣耀，而祂是为我而存在，如同钳子是为车，锯子是为门，那我就在因果关系上优于祂。因为无论天主超越受造物多少，那为我而存在者，就同样低于那为天主而存在的我。\par

18. 此外，\Place{摩阿布人} 和 \Place{阿孟人} 甚至不可被允许进入 \BibleRef{申 23:3} 天主的教会——我是指那些诡辩的、有害的论证，它们好奇地探究天主那不可言说的 \OriginalTerm{生成}{generation} 和 \OriginalTerm{发出}{procession}，并鲁莽地列阵反对天主性；仿佛那些言语无法表达的事，要么必须只有他们能理解，要么就因为未被他们理解而不存在。然而，我们遵循圣经，并除去瞎子路上的绊脚石，将紧握救恩，敢于做任何事，只不敢傲慢反对天主。至于证据，我们留给别人，因为许多人已陈述过，我们自己也以不少心力陈述过。事实上，此刻我要去收集那些始终被相信的论据，是非常可耻的。因为，即使在小事和无关紧要的事上，先教后学都不是最好的次序，更何况在神圣而如此重要的事上。再者，现在也不适合解释和理清圣经的疑难，那需要比当前目的所允许的更充分、更仔细的考虑。总而言之，这就是我们的教导。我详述这些，并非要与敌手争辩：因为我已经常与他们辩论，即使不完美；而是为了向你们展示我教导的特质，使你们看到我是否在维护你们自己的辩护上有所参与，是否与你们站在同一立场，抵抗同一敌人。\par

19. 朋友们，你们现在已听到了我在此在场的辩护：如果值得称赞，那么为此感谢天主，并感谢你们——召叫我的人；如果低于你们的期望，我仍然为此感恩。因为我确信它并非全然该受责备，并且我相信你们也承认这一点。我曾 \BibleRef{格后 12:17} 占过这人民的便宜吗？我曾如我常见的那样，谋过自己的私利吗？我曾使教会烦恼吗？对其他人——他们以为因我们缺席而赢得了判决——我们在辩论中已经予以驳斥；但就我所知，对你们却毫不如此。伟大的 \Person{撒慕尔} 在反对以色列关于君王的事上说，\Quote{我没有取过你们的一头牛，\BibleRef{列上 12:2} 也没有取过你们灵魂的赎罪祭，主在你们中间作证}，没有这也没有那，他列举更多，以免我一一细数；但我保持了司祭职的纯洁无瑕。如果我曾贪爱权力、宝座的高度、或践踏君王的宫殿，愿我永远得不到任何尊荣，或者如果得到了，也愿我被夺去。\par

20. 那么我是什么意思呢？我不是无赏的德行之士，尚未达到如此高的德行。请赐予我劳苦的报酬。什么报酬？不是某些容易怀疑的人所想的那一种，而是我寻求而无害的一种。请让我从长久的辛劳中歇息；请给我的外方服务以尊荣；请另选一人代替我，就是那位为你们所热切寻求的，一位手洁的，一位声音不拙的，一位能在所有方面取悦你们、并与你们分担教会事务的人；因为这正是此类的时机。但请看，我恳求你们，这身体的状况——被时间、疾病、辛劳耗尽。你们需要一个胆怯而不刚强的老人吗？他可以说是日复一日地死去，不仅在身体上，甚至在智力上也如此，发现很难在你们面前详述这些？不要违抗你们导师的声音：因为你们确实从未违抗过。我厌倦了被指责为温和。我厌倦了被敌人和我们自己的人以言语和嫉妒攻击。有些人瞄准我的胸膛，但他们的努力不太成功，因为公开的敌人可以防备。另一些人埋伏攻击我的背部，造成更大的痛苦，因为未预料的打击更致命。如果我曾是一名舵手，我也算是最熟练的之一；大海在我们周围汹涌，在船旁沸腾，乘客们一片喧哗，他们总是为某事争吵，彼此咆哮，也向海浪咆哮。我坐在船舵旁，与大海和乘客同时搏斗，要在这双重风暴中将船安全地驶入陆地，这是多大的挣扎啊！如果他们事事支持我，安全也来之不易；而当他们反对我时，又如何可能避免沉船呢？\par

21. 还有什么需要说的？但我怎能承受这圣战？因为据说有一场圣战，正如波斯战争一样。我怎能将与之为敌的各教座、敌对的牧者，以及因之分裂并对立的子民联合起来——他们仿佛被地震造成的裂缝分隔在邻近的地区；又如同在瘟疫中，家仆和家庭成员相继染病；此外，地球上的各个区域都受到派系精神的影响，以致东方和西方对垒，似乎要在意见上和地理上同样分裂。我的和你的、老和新、更理性的和更属灵的、更高贵的和更卑贱的、更多的和更少的团体，还要持续多久？我为我的老年感到羞耻，因为我被基督拯救后，却被人以别人的名字称呼。\par

22. 我无法忍受你们的赛马场和戏院，以及这种在费用和派系精神上争奇斗艳的狂热。我们解下缰绳，又从另一边套上，彼此嘶鸣，几乎像他们一样拍打空气，向天空扬起尘土，如同那些兴奋的马匹；并在其他面具下满足我们自己的竞争心，成为不良的竞争裁判和糊涂的事务判断者。今天共享同一宝座和意见——如果我们的领袖如此带领我们；明天又因意见和立场的敌对——如果风向相反。在友谊和仇恨的变化中，我们的名字也不断变化：最可怕的是，我们毫不羞耻地向同一群听众提出相反的道理；我们也不能始终如一，因着好争论而不同时候有不同的表现。它们就像某个狭窄海峡的潮涨潮落。因为当孩子们在市场中间玩耍时，我们若丢下家务加入他们，那是最可耻、最不合适的；因为儿童的玩具不适合老年：同样，当别人争斗时，即使我比大多数人有更好的认识，我也不会让自己成为他们中的一员，而是如现在一样，享受隐没的自由。因为除此之外，我的感觉是：我在大多数问题上与多数人意见不一致，不能忍受走同一条路。虽然这可能是鲁莽和愚蠢的，但我就是这种感觉。别人喜欢的东西使我痛苦，别人痛苦的东西我却喜欢。因此，即使我被当作一个令人讨厌的人而囚禁，被大多数人认为失去理智，我也不感到惊讶——据说有一位希腊哲学家就是这样，他的节制使他被指控为疯狂，因为他嘲笑一切，因为他看到大多数人热切追求的对象是可笑的；或者甚至被认为充满新酒，正如后来基督的门徒那样，因为他们说方言，\BibleRef{宗 2:4}，因为人们不知道这是圣神的能力，而不是心智混乱。\par

23. 现在，请考虑对我们的指控。据说，你统治教会已经这么久，并受到时间的推移和君主的影响（这是最重要的事）。我们注意到了什么变化？从前有多少人暴虐地对待我们？我们经历了什么苦难？虐待？威胁？流放？抢劫？没收？将司铎烧死在海上？用圣人的血亵渎圣殿，以至于圣殿变成了埋骨所？公开杀害年长的主教，更准确地说，杀害宗主教？唯独对虔诚者禁止进入任何地方？事实上，任何能想到的苦难？我们对这些作恶者报以了什么？因为命运之轮给了我们权力，可以正当对待那些如此对待我们的人，我们的迫害者应该受到教训。除了所有其他事情，只谈我们的经历，不提你们自己的——我们不是被逼迫、虐待、赶出教堂、房屋，甚至最可怕的是，赶出旷野吗？我们不是忍受了愤怒的民众、傲慢的官员、皇帝及其法令的漠视吗？结果是什么？我们变得更坚强，我们的迫害者逃跑了。事实正是如此。在我看来，能够报复他们就是对那些欺负我们的人足够的报复。这些人却不这么想；因为他们对报复极为精确和公正：因此他们要求现状允许的惩罚。他们说，哪个官员被罚了款？哪里的民众受到了惩罚？民众的头目呢？我们为将来激起了他们对我们多大的畏惧？\par

24. 也许我们会受到责备，如同从前一样，因为我们的宴席精美、衣着华丽、前驱的侍从、对遇见者的傲慢。我不知道我们应当与执政官、总督、最显赫的将军们较量，他们没有机会挥霍他们的收入；也不知道我们的肚腹应当渴望享用穷人的财物，将他们的必需品花费在奢侈品上，并在祭台前打嗝。我不知道我们应当骑高头大马，乘华丽马车，有队伍前导，受鼓掌围绕，人人都为我们让路，如同我们是野兽一般，开辟一条通道，使我们的到来从远处就可以看见。如果这些苦难曾经受过，现在它们已经过去了：请原谅我这过错。\BibleRef{格后 12:13}另选一位能取悦多数人的人吧；把沙漠、乡间生活、我的天主赐给我，我只需取悦祂，并藉着简朴的生活取悦祂。失去演讲、会议、公众集会、以及像现在这样使我思绪飞扬的鼓掌、亲友、荣誉、城市的美丽与宏大、及其眩目（迷惑那些只观表面而不探究事物本质的人），固然是痛苦的；但还不如在公众骚动与动荡中（它们随风转舵）被喧嚷和玷污那么痛苦。因为他们寻求的不是司铎，而是演说家；不是灵魂的管家，而是金钱的司库；不是纯洁的祭献奉献者，而是强有力的庇护者。我要为他们说句辩护的话：\Quote{我们就是这样训练了他们，成为一切人中的一切\BibleRef{格前 9:22}，是为了拯救一切还是毁灭一切，我不知道。}\par

25. 你们怎么说？你们被我的话说服了吗？还是我必须用更强的言辞才能说服你们？是的，我指着你们和我共同敬拜的天主圣三本身、指着我们共同的希望、并为了这百姓的合一，请赐我这恩惠：用你们的祈祷送我离去；让我这争斗的宣告是这样；把退休证明给我，如同君王给他们的士兵一样；如果你们愿意，就带着好评，好让我享受这荣誉；如果不愿，随你们便；这对我无关紧要，直到天主看清我的实情。那么我们要选谁做继任者？天主必会亲自预备一位牧者来担任此职，正如祂从前预备了一只羔羊作燔祭一样。\BibleRef{创 22:8}我只再请求这一点——愿他成为被人嫉妒而非怜悯的对象；不是凡事顺从一切人，而是能在某些方面为最佳利益而抵抗的人。因为前者最令人愉快，后者最有益处。所以请你们为我准备告别致辞：我现在要向你们辞别了。\par

26. 再会，我的\Place{阿纳斯塔西娅}{Anastasia}，你的名字散发着虔诚的芬芳：因为你为我们振兴了曾被轻视的教义；再会，我们共同胜利的场所，现代的\Place{史罗}{Shiloh}，\BibleRef{苏 18:1}是会幕首次安设之处，此前在旷野中漂泊了四十年。同样再会，宏大而著名的圣殿，我们的新产业，其伟大如今是由于圣言，它曾是\Place{耶步斯}{Jebus}，\BibleRef{编上 11:4}而今被我们建成了耶路撒冷。再会，你们所有其他人，在美丽上仅次于此，散布在城中各处，如同许多链环，将各自的邻里联结起来，充满了那些我们曾绝望存在的敬拜者，不是藉着我这软弱者，而是藉着与我同在的恩宠。\BibleRef{格前 15:10}再会，你们宗徒们，这里的高贵定居者，我斗争中的主宰们；如果我并不常与你们一起过节，那可能是因为我像你们中的一位，\Person{圣保禄}{S. Paul}\BibleRef{格后 12:7}一样，身上带着\OriginalTerm{撒殚}{Satan}，为了我的益处，也正是我现在离开你们的原因。再会，我的宝座，令人羡慕又危险的高位；再会，大司祭的集会，因其司铎的尊严和年岁而受尊崇，以及你们所有在祭台旁亲近那亲近你们的天主的天主仆人。\BibleRef{雅 4:8}再会，\OriginalTerm{纳齐尔人}{Nazarites}的歌咏团，\OriginalTerm{圣咏集}{Psalter}的和谐，彻夜的守夜，可敬的贞女，端庄的已婚妇女，寡妇和孤儿的聚会，以及你们穷人的眼睛，转向天主也转向我。再会，好客且爱基督的居所，我软弱的帮助者。再会，我讲论的爱好者们，你们的热情和聚集，看不见的和看得见的笔，以及那被向前拥挤听道的人所紧靠的栏杆。再会，皇帝们、皇宫、大臣们和皇室，无论对皇帝是否忠诚——我不知道——但大部分对天主不忠诚。拍手吧，欢呼吧，将你们演说家捧上天。这有害而饶舌的舌头已停止对你们说话。虽然它不会完全停止说话：因为它将用手和墨水作战；但现在我们已停止说话。\par

27. 再会，伟大的爱基督之城。我将为真理作证，尽管你们的热心并不合乎真知。\BibleRef{罗 10:2} 我们的分离使彼此更加温良。请接近真理：在这迟暮之时悔改吧。要尊崇天主，胜过你们素常所为。改变并非耻辱，而执著于恶则是致命的。再会，东方和西方，我曾经为你们并反对你们而战；祂是证人，祂将赐给你们和平，只要有人效法我的退隐。因为放弃自己宝座的人，并不会失去天主，反而会拥有那更崇高、更安稳的高天之位。最后，也是最重要的，我要呼喊——再会，天使们，这圣殿的守护者，以及我在此处和朝圣之旅的守护者，因为我们的处境全在天主手中。再会，天主圣三，我的默想与我的荣耀。愿你们由在此的人们所保存，并保存他们，我的子民：因为他们是属于我的，即使我被派往别处；愿我得知你们在言语和行为上永远受颂扬、受光荣。我的孩子们，我恳求你们，守护所委托于你们的。\BibleRef{弟前 6:20} 请记得我的被石击。\BibleRef{哥 4:18} 愿我们主耶稣基督的恩宠与你们众人同在。阿们。\par

\SourceRecord{Translated by Charles Gordon Browne and James Edward Swallow.；Nicene and Post-Nicene Fathers, Second Series；Vol. 7.；Edited by Philip Schaff and Henry Wace.；Buffalo, NY: Christian Literature Publishing Co.,；1894.；<http://www.newadvent.org/fathers/310242.htm>.}

% structure: p0001 source=h1 target=section reason=page-title

\section{演说第四十三篇}

\Emph{关于加帕多家凯撒勒雅主教大巴西略的葬礼演说。}\par

\EditorialNote{圣巴西略卒于主历三七九年一月一日。一场重病及其他原因使圣额我略未能出席他的葬礼（书信79）。Benoît认为，圣额我略在《墓志铭》第119篇第38行中说他的嘴唇被缚，证明他当时仍在塞琉西亚退隐。这是一个无根据的推论。在本演说的第2节中，圣人轻描淡写地提到自己的疾病，却将他在君士坦丁堡的劳苦作为缺席的更迫切理由，并说他乃依照圣巴西略的判断承担那任务。这暗示圣额我略是在圣巴西略去世前就去了君士坦丁堡，或者那时他已受朋友劝告影响而正准备出发——更可能是前者，因为我们可以确信，如果圣额我略仍留在塞琉西亚，除非身体不支，没有理由不陪伴在朋友身边。他在君士坦丁堡的紧迫职务和长途旅行的困难是他致信尼撒的圣额我略的其他原因；而且我们知道他在君士坦丁堡患过重病（诗歌11.887；演说23.1）。圣额我略于主历三八一年六月离开君士坦丁堡，Tillemont将本演说的年代定在他返回纳齐盎后不久。Benoît认为它可能是在圣巴西略逝世周年纪念日发表的。所有评注家一致认为，这篇演说极具力量和美感。其长度（对开本62页）、发言者体力虚弱以及即使有兴趣的听众的耐力极限，使我们倾向于认为它并非以现存形式口述的。我们无法轻易排除那些明确指向实际宣讲的表述，但它后来可能被增补过。}\par

1. 因此，伟大的巴西略——他常常为我提供讲论的题材，对此他比其他任何人更为自豪——如今亲自为我提供了演说家所曾遇到的最伟大的题材。因为我认为，如果有人想试验自己的口才，以他所偏爱的所有题材中那一个为标准（如同画家对待划时代的画作），他会选择其他所有题材中居于首位的那个，但会把它搁置一边，视为人的口才所不能及的。赞美这样一个人，不仅对于我这个早已抛弃一切竞争念头的人，甚至对于那些为口才而活、唯一目标就是通过这类题材获得荣耀的人，都是一项如此艰巨的任务。这是我的看法，而且我自认为十分公正。但我不知道，如果我不处理这个题材，我还能处理什么题材能显出我的口才；或者我还能给自己、给美德的仰慕者、或给口才本身，带来比表达我们对这个人的钦佩更大的好处。对我来说，这是偿还一笔最神圣的债务。而我们的言语，对于那些特别蒙受言语恩赐的人，是超越其他一切的债务。对于美德的仰慕者，一篇讲论既是愉悦，也是对美德的激励。因为当我得知对人们的赞美时，我就清晰地知道他们的进步：如今，我们中没有任何人不在其能力范围内，可以在那进步中达到任何程度。至于口才本身，无论怎样，一切都必顺利进行。因为，如果讲论几乎配得上其主题——口才就展示了它的力量；如果它远远不及，如同在陈述巴西略的赞辞时必然发生的那样，那么通过实际证明其无能，它就宣告了其主题的卓越超乎一切言语的表达。\par

2. 这些理由促使我发言，并迎接这场竞赛。对于我这样迟到，在他已得到许多人——无论在公开还是私下尊崇他——的赞美之后才出现，愿无人感到惊讶。是的，愿那个神圣的灵魂、我恒常敬畏的对象宽恕我！而且，当他还在我们中间时，他常常根据朋友的权利和更高的法则在许多方面纠正我——我说这话并不羞耻，因为他是我们所有人的美德标准——所以如今，他从上俯视我，会宽容地对待我。我也请求在场的他的最热忱仰慕者中任何人的原谅——如果真有人能比别人更热忱，而且我们不是所有人都并驾齐驱地热忱于他的美名。因为使我未能达到可能预期的，并非轻视；我也并非如此不顾美德或友谊的要求；我也未曾认为赞美他更适合其他任何人而非我。不！我的第一个理由是，我畏避这项任务——因为我要说出真相，如同那些在声音和心灵洁净之前就接近神圣职分的司铎那样。其次，我要提醒你们——尽管你们很清楚——我所从事的关于真道的任务，那任务本是适当加给我的，并把我从家中带出来，按我猜想是出于天主的旨意，而且一定是按照我们那位卓越的真理捍卫者的判断——他一生的气息就是那促进全世界救恩的虔诚教义。至于我的身体健康，也许我不该斗胆提及，因为我的主题是一个如此刚毅地征服了身体的人——甚至在他离开此世之前——并坚持认为灵魂的任何能力都不应受这个锁链的妨碍。这就是我的辩护。我向那些如此清楚我事务的你们谈论他时，我认为无需再多加辩解。我现在必须继续我的颂词，将我交托给他的天主，以使我的赞美不致成为对这个人的侮辱，也不致让我远远落后于其他所有人；尽管我们所有人都同样达不到他应得的赞美，如同那些仰望天空或太阳光芒的人一样。\par

3. 如果我曾看到他为自己出身和出身权利而骄傲，或为那些眼光向下的人所无限微小的任何事物而骄傲，我们就得审视一份关于英雄的新名单。关于他的祖先，我本可以调动多少细节！甚至历史也不会比我更有优势，因为我声称有这个优势：他的声誉并非依赖于虚构或传说，而是依赖于许多见证人所证实的实际事实。在父亲一方，\Place{本都}为我提供许多细节，丝毫不亚于它古时的奇迹——那些奇迹充满了所有历史和诗歌；还有许多其他细节涉及这个我的故土\Place{加帕多家}的杰出人物，它因年轻的后代而出名，不亚于因马匹而出名。因此，我们用他母亲的家族来匹配他父亲的家族。还有哪个家族拥有更多或更杰出的将领和总督，或朝臣，或富人和崇高的宝座、公共荣誉、以及演说家的名声？如果允许我提及它们，我会不把\Person{珀罗普斯家族}、\Person{开克洛普斯家族}、\Person{阿尔克迈翁家族}、\Person{埃阿科斯家族}和\Person{赫拉克勒斯家族}以及其他最高贵的家族放在眼里：因为他们，由于家中没有公共功绩，就求助于不确定的领域，声称半神和神祇——纯属神话人物——作为他们祖先的荣耀，他们最吹嘘的细节不可信，而我们可以相信的那些则是耻辱。\par

4. 但既然我们的主题是一位坚持认为每个人的尊贵应根据其自身价值来判断的人，并且正如形态和颜色，以及我们最著名和最臭名昭著的马匹，由其本身的特性来检验，同样我们也不应借用他人的羽毛来描绘自己；在提及一两个特征后（这些特征虽来自祖先，但他通过自己的生活使之成为自己的，并且特别能令听众愉悦），我将进而讨论他本人。不同的家族和个人有不同的显著之处和兴趣，大小不一，如同长短不一的祖传遗产，传于后代：他家族双方的显著之处是虔敬，我现在将展示这一点。\par

5. 有一场迫害，最可怕和最严厉的；我是指，如你们所知，\Person{马克西米努斯}的迫害，它紧随其前任者之后，以其过度的胆大妄为和渴望赢得不虔敬之暴行的冠冕，使它们都显得温和。它被我们的许多勇士所克服，他们与之搏斗至死，或几乎至死，只留下足够活命的生命以存活于胜利之后，并不要在斗争中逝去；他们留在世上成为德行的训练者、活的见证、呼吸的战利品、无声的劝勉，在他们众多的行列中，发现了\Person{巴西略}的父系祖先，在他们实践各种虔敬形式时，那个时代赐予了他们许多美丽的花冠。他们如此预备和决意，愿意欣然承受这一切，因这些，基督便为那些效法祂为我们而争战的人戴上冠冕。\par

6. 但既然他们的争斗必须是合法的，而殉道律法同样禁止我们自愿前去迎向它（出于对迫害者和软弱者的考虑），或是在它临到时退缩；因为前者显示鲁莽，后者显示怯懦；在这方面他们向立法者给予了应有的尊敬；但他们的计策是什么，更确切地说，他们被那在一切事上引导他们的天主眷顾引领到了何处？他们前往\Place{本都}山中的一处丛林，那里有许多幽深广阔的丛林，带着极少的逃亡同伴或照料他们需要的仆从。让其他人惊叹于此时间之长，因为他们的逃亡极其漫长，约七年或稍多一点，而他们原本娇生惯养的生活方式变得困窘而不寻常，可想而知，要忍受霜冻、炎热和雨水的困扰；荒野不允许与朋友有团契或交谈：这对习惯于众多随从的侍候和尊荣的人来说是一个巨大的考验。但我将要谈论更加伟大和非凡的事：没有人会不相信它，除了那些在软弱而危险的判断中，轻视为基督所受的迫害和危险的人。\par

7. 这些高尚的人，因时间流逝而受苦，对普通食物感到厌恶，渴望更开胃的东西。他们确实不像以色列人那样说话，因为他们不像他们那样在出埃及后的旷野苦难中成为怨言者\BibleRef{格前 10:10}——那埃及对他们来说本比旷野更好，因那里丰盛供应的肉锅以及其他美味他们留在了那里，因为那时他们在愚昧中似乎对制砖和泥土毫不在意——而是以更虔敬和忠信的方式。因为他们说：\Quote{为何不可信呢？那行奇迹的天主，曾在旷野中慷慨养活祂无家可归、逃亡的百姓，降下玛纳\BibleRef{出 16:13}，并赐予丰盛的鹌鹑，不仅以必需品喂养他们，甚至以奢侈品；那一位分开大海、停止太阳、分开河流，以及行了其他一切事的天主；因为在这样的情况下，人心常会回顾历史，歌颂天主的诸多奇迹。他们接着说，那一位今天难道不能以美味喂养我们这些虔敬的勇士吗？许多从富人餐桌逃脱的动物，在这山中藏身，许多可食的鸟类飞过我们渴望的头顶，只要您一句话，它们都一定可以被捕捉！}话音刚落，他们的猎物便出现在面前，食物自动到来，一席盛筵不劳而备，雄鹿忽然从山间某处出现。它们多么壮丽！多么肥美！多么适于屠宰！几乎可以想象它们因未被更早召唤而生气。其中一些鹿示意引领其他鹿跟随，其余的便跟随它们。谁追赶并驱赶了它们？没有人。什么骑手？什么样的狗，什么样的吠叫或喊声，或按狩猎规则占据了出口的年轻人？它们是祈祷和正义恳求的俘虏。在当天或任何一天的人中，谁曾见过这样的狩猎？\par

8. 哦，何等奇妙！他们自己就是狩猎的掌管者；凡他们所愿的，仅凭意愿就被捕获；留下的，他们便放回丛林，供另一餐食用。厨师是即兴的，宴席精致，宾客们为这奇妙的希望预尝而心怀感激。由此，他们更加热切地奋斗，以报答所领受的恩惠。这就是我的故事。而你，我的迫害者，在你对传说的赞叹中，请讲述你的女猎手、\Person{俄里翁}、\Person{阿克泰翁}——那些命运多舛的猎人，以及那替代少女的母鹿——若此类事物使你奋起效仿，且我们承认这故事并非传说。这故事的结局过于不体面。因为若一个少女被救是为了学会谋杀她的宾客，并以不人道回报人道，这交换有何益处？仅此一例，权且如此，从众多中选出，代表其余，就我而言。我讲述它并非为增添他的声誉：因为大海并不需要流入其中的众多江河，无论它们多么浩大，而我当前赞美的主题也不需要任何对他美名的增补。不！我的目的是展现他先辈的品格，以及他眼前的榜样，他远远超越了这榜样。因为若其他人从祖先那里获得部分荣誉是一个巨大的额外优势，那么他则是对原本的资本作了如此增补，以致溪流似乎逆流而上。\par

9. 他父母的结合，既由美德共同体所巩固，也由同居生活所维系，因诸多原因而卓著，尤其是对穷人的慷慨、好客、因自律而来的灵魂纯洁、将部分财产奉献于天主——此事在当时尚未如现今这般为多数人所重视，它因先前那些使之彰显的榜样而增长——以及所有其他在\Place{本都}和\Place{卡帕多西亚}广为流传、令多人满意的优点。然而，依我之见，他们最大的特点是其子女的卓越。传说中确有子女众多且俊美的例子，但实际经验向我们呈现了这样一对父母：他们自身的品格，即使没有子女的品格，已足以令其美名远扬；而他们子女的品格，即使没有他们自身的卓越，也会因其子女的卓越而超越所有人。因为若有一两个子女获得卓越，可归因于天性；但若所有子女都卓越，这荣耀显然归于教养他们的人。这由受祝福的司铎与贞女之列，以及那些在婚姻中未允许任何事物阻碍他们获得同等声誉，从而使他们之间的区别在于身份而非生活方式的人们所证明。\par

10. 谁不认识我们的总主教之父\Person{巴西略}？他对每个人都是一个伟大的名字，他实现了父亲的祈祷——若曾有人如此，我敢说无人能及。因为他在美德上超越众人，只是被他的儿子阻止了获得头奖。谁不认识\Person{埃美利亚}？她的名字预示了她的所成，或者说她的生活是她名字的写照。因为她有资格享有这意指优雅的名字，并且简而言之，在妇女中占有她丈夫在男子中同样的地位。因此，当决定应将我们聚会所荣耀的那一位赐予人类，使其服从本性的束缚——如同古时任何一位由天主为共同益处而赐予者——他既不适合由其他父母所生，他们也不适合有另一个儿子：于是两件事恰当地一同发生。如今，我服从神圣律法——这律法命令我们向父母致以全部敬意——已将赞美的初果献给我所纪念的他们，然后继续论述\Person{巴西略}本人，并预先说明这一点：我认为所有认识他的人都将视为真实，即我们只需要他自己的声音来发表对他的颂词。因为他既是一个光辉的赞颂主题，又是唯一口才足以处理这主题的人。至于美貌、力量和身材——我见大多数人在这些方面喜悦——我让给任何愿意领取的人——并非在这些方面他年轻时、尚未以苦行削弱肉体时次于那些汲汲于身体的渺小人物——而是为了使我避免不熟练的竞技者的命运：他们在次要目标上徒耗力量，因而在决定性的搏斗中落败，而搏斗的结果是胜利和冠冕的荣耀。因此，我将为他争取的赞美基于以下理由：我认为无人会认为这些理由多余或超出我演讲的范围。\par

11. 我以为是明理之人所公认的，我们的首要益处乃是\OriginalTerm{教育}{education}；不仅是我们这种更高贵的形式，它不顾修辞的装饰与荣耀，而坚持救恩和\OriginalTerm{默观}{contemplation}对象的美；而且是那种外在的文化，许多基督徒不智地憎恶它，视其为奸诈危险，使我们远离天主。因为正如我们不应忽视天、地、空气及一切此类事物，只因有人错误地攫取它们，尊崇天主的工程而非天主；而应从它们中为我们的生活和享受获取所能得到的益处，同时避免其危险；不象愚昧人那样使受造物起来反抗造物主，而是从自然工程中领悟工程师，并如神圣宗徒所说，将一切思想掳去顺从基督：\BibleRef{格后 10:5}；又正如我们知道，火、食物、铁或任何其他元素，本身并非最为有用或最有害，而取决于使用者的意愿；并且我们曾从某些爬虫中配制出健康的药物；同样，我们从世俗文学中接受了探究与思辨的原则，同时拒绝了它们的偶像崇拜、恐怖与毁灭之坑。不，甚至这些也通过我们对更坏与更好的对比，以及从他们教义的软弱中获得我们教义的力量，而帮助了我们的宗教。因此，我们不可因有些人乐意这样做就羞辱教育，而应认为此等人是粗俗无教养的，他们希望所有人都像他们一样，以便隐藏于大众之中，逃避其缺乏文化的察觉。但来吧，在草述我们的主题并承认这些之后，让我们默观\Person{巴西略}的生平。\par

12. 在他最早的岁月里，他在那位伟大的父亲——当时被\Place{本都}公认为公共\OriginalTerm{德行}{virtue}导师——的指导下，被包裹并塑造，以神圣的\Person{达味}所说的那种最好、最纯洁的、与黑夜相对、日复一日地进行塑造。在他之下，随着生命与理性一同增长，我们杰出的朋友受到了教育：他并不夸耀\Place{色撒利}的山洞作为他德行的工坊，也不夸耀某个吹嘘的\OriginalTerm{半人半马}{Centaur}作为当时英雄们的导师；他也没有在这样的教导下学会射野兔、追小鹿、猎雄鹿、或在战争中出众、或驯服马驹，同时将同一个人用作教师和马匹；他也没有以神话中的鹿髓和狮髓为食，而是接受了通识教育，练习敬拜天主，简言之，通过初步的指导被引向未来的\OriginalTerm{成全}{perfection}。因为那些只在生活或文学上成功，而在另一方面有缺陷的人，在我看来与独眼人毫无差别，他们的损失很大，但他们的畸形更大，无论在他们自己眼中还是在他人眼中。而那些在两者上都同样卓越并左右逢源的人，既拥有成全，也在天堂的福乐中度过一生。这正是他所得着的，他在家中有一个在善行上的德行典范，仅仅看见这典范就使他从一开始就杰出。正如我们看到小马驹和牛犊从出生起就在母亲身边跳跃，他也如此，像小马驹一样顽皮地紧跟在父亲身旁，在他向德行的崇高冲动中并未被远远落下，或者，如果你愿意，他勾画并展示了他未来德行之美的痕迹，并在成全的时刻到来之前画出了成全的轮廓。\par

13. 当在家受过充分训练后，既然他不应在任何形式的卓越上有所欠缺，也不应被那从每朵花采集最有用之物的忙碌蜜蜂所超越，他便出发前往\Place{凯撒勒雅}城，以便在那里求学——我是指我们这座著名的城市，因为它是我学业的引导者和女主人，不亚于其他城市的文学之都，她超越并统治着这些城市：如果有人剥夺她的文学力量，他就夺去了她最美丽、最特殊的标志。其他城市以其他装饰为荣，无论是古老的还是近来的，以便有可以描述或观看的东西。文学是我们的标志，仿佛是战场或舞台上我们的徽章。他后来的生活，让那些训练过他或享受过他训练的人来详细描述，他对于他的老师们如何，他对于他的同学们如何：在每一种文化形式上，他与前者相等，超越后者；他在短时间内从所有普通人以及国家领袖那里赢得了怎样的声誉；他展现出超越年龄的文化和超越文化的性格坚定。甚至在修辞学家的教席之前，他就是演说家中的演说家；甚至在哲学家的教义之前，就是\OriginalTerm{哲学}{philosophy}家中的哲学家；最高的是，甚至在\OriginalTerm{司铎职}{priesthood}之前，他就是基督徒中的司铎。所有人在各方面都对他如此恭敬。雄辩是他的副业，他从中采集了足够多，使之成为他基督徒哲学的辅助，因为表达\OriginalTerm{默观}{contemplation}对象需要这种能力。因为无法表达自己的心智如同昏睡之人的动作。他的追求是哲学，是从世界脱离，与天主共融，在属世事物中关心属天事物，并在一切不稳定、变动之处赢取那稳定而持久的事物。\par

14. 然后他到\Place{拜占庭}，东方帝国都城，因其修辞与哲学教师的卓越而闻名；凭借其敏锐而有力的才智，他很快吸收了其中最宝贵的课程。随后，受天主派遣，并因他对文化的热切追求，他到了\Place{雅典}——文学之乡。\Place{雅典}对我来说，若对任何人而言，确是一座黄金之城，是一切美善的守护者。因为它使我更完美地认识了\Person{巴西略}，尽管此前我并非不认识他；在追求学问的过程中，我获得了幸福；并且以另一种方式，我经历了与\Person{撒乌耳}相同的事，\BibleRef{撒上 9:3}，他寻找父亲的驴，却找到了王国，无意中得到了比他原定目标更重要的东西。至此，我的叙述清晰明了，沿着平坦易行的王道引领着我的赞词；从此以后，我不知道该说什么，也不知往何处去，因为我的任务变得艰巨。此刻，我心中焦虑，并借此机会从我自身的经历中为我的演说增添一些内容，稍作停留，叙述我们友谊——或者更确切地说，我们生活与性情的合一——的起因与境况。因为正如我们的眼睛不容易从吸引人的对象上移开，若强行移开，它们又常常转回去；同样，我们在描述最甜蜜的事物时也会流连忘返。我害怕这项工作的艰难。然而，我将尽力保持节制。若我因怀念而有些失控，请宽恕这种最合理的情感，在有情感的人眼中，缺乏这种情感将是巨大的损失。\par

15. 我们在\Place{雅典}被容纳，如同河水的两条支流；因为离开故乡的共同源头后，我们在各自追求文化的过程中分离，如今又因天主的推动以及我们自己的同意而重新结合。我比他稍早一点，但他很快跟随着我，带着巨大而辉煌的期望而来。因为他到达之前就已精通多种语言，我们任何一方在学业上超越对方都是一件大事。我不妨再加一段简短的叙述作为演说的调味，这将成为了解者的回忆，不了解者的知识来源。\Place{雅典}的大多数年轻人在愚蠢中疯狂追求修辞技巧——不仅是出身卑贱无名之辈，甚至包括高尚显赫之人，在难以控制的年轻人群体中。他们就像那些沉迷于马匹和表演的人一样，正如我们在赛马场上所见：他们跳跃、呼喊、扬起尘土，在座位上策马，以手指为鞭抽打空气（而非马匹），套马、解马，尽管那些马并不属于他们；他们随时交换车夫、马匹、位置、领头者；而做这些事的是谁？往往是贫困潦倒、连一天生计都无法维持的人。这正是学生们对待自己导师和对手的态度，急于增加自己一方的人数从而丰富自己。这件事绝对荒谬愚蠢。城市、道路、港口、山顶、海岸线——总之，\Place{阿提卡}或\Place{希腊}其他地方的每一处，连同大部分居民，都被占领；因为连这些也被敌对派系瓜分了。\par

16. 每当有新人到来，落入那些抓取他的人们手中——无论是被迫还是自愿——他们便遵守这条半开玩笑半认真的\Place{阿提卡}法律。新人首先被带到最先接纳他的人之一的家中，或是他的朋友、亲戚、同乡，或是那些辩论能力卓越、提供论据、因此在他们中特别受尊敬的人；他们的报酬在于获得追随者。随后，他便遭受任何愿意之人的嘲弄，我想其意图是抑制新人的自负，并立即使他们屈服。嘲弄根据嘲弄者的粗野或文雅而更加无礼或更具辩论性；这一表演，对不知者而言似乎非常可怕和粗暴，但对经历过的人来说却非常愉快和人性化：因为其威胁更像是假装而非真实。接着，他被列队穿过市场前往浴场。行列由那些为年轻人负责的人组成，他们分成两列，间隔一段距离，在他前面走到浴场。但临近浴场时，他们如着魔般叫喊狂跳，高喊不能前进，必须停下，因为浴场不接纳他们；同时猛烈敲门吓唬年轻人。然后让他进入，此刻他们赋予他自由，浴后便以平等者的身份接纳他，视其为己出一员。他们认为这是仪式中最愉快的一部分，因为烦恼迅速得到了交换和解除。在这次场合，我不仅因尊重我伟大的朋友\Person{巴西略}品格的庄重和推理能力的成熟而拒绝使他蒙羞，还说服了其他所有恰好不认识他的学生同样对待他。因为他从一开始就受到大部分人尊敬，其名声早已先行。结果，他是唯一逃脱常规的人，获得了比新生地位更高的荣誉。\par

17. 这是我们友谊的前奏。这是我们结合的火花：我们就这样感受到了彼此的爱之伤。然后发生了这样一件事，因为我认为连这一点也不该省略。我发现亚美尼亚人不是一个单纯的种族，而是非常狡猾和诡诈的。当时，他的一些特别亲密的同伴和朋友，甚至在他父亲教导的早期就与他亲密无间——因为他们是他学派的成员——以友谊的伪装来到他面前，但内心怀着嫉妒而非善意，向他提出一些好辩而非理性的问题，试图在第一次交锋中就压倒他，因为他们早已知道他天生的禀赋，无法容忍他当时所获得的荣誉。他们认为，那些已经穿上长袍、练习过辩论的人，竟然无法战胜一个陌生人和新手，这真是奇怪的事。我也出于对雅典的虚妄之爱，相信他们的表白而没有察觉他们的嫉妒，当他们退让和转身时，我因雅典的声誉在他们身上被摧毁而感到愤慨，这么快就蒙羞，于是支持了那些年轻人，并恢复了辩论；借助我额外的份量——因为在这种情况下，一点小小的添加就能改变一切，如诗人所说，\Quote{使他们在战斗中势均力敌}。但当我察觉到争论的隐秘动机——它再也无法掩盖，最终清楚地暴露出来——我立刻后退，离开他们的行列，站到他一边，使胜利成为决定性的。他立刻对所发生的事感到高兴，因为他的洞察力非常敏锐，充满热忱——用荷马的话充分描述他——他用论据追击那些勇敢的年轻人，用三段论打击他们，直到他们彻底溃败，他才停止，并明确赢得了与其能力相称的荣誉。于是，友谊再次被点燃，不再是火花，而是显明可见的火焰。\par

18. 他们的努力就这样徒劳无功，他们严厉责备自己的鲁莽，对我怀有的恼怒竟发展成公开的敌意，并指控我背叛——不仅背叛他们，也背叛雅典本身：因为他们在第一次交锋中就被一个甚至还没来得及获得自信的学生驳倒并蒙羞。而他呢，根据人类那种一旦突然达到所怀有的崇高希望，就觉得结果低于期望的情感，他感到不悦和烦恼，无法从抵达中感到快乐。他在寻找他所期望的，并称雅典为空洞的幸福。我却试图通过论战和推理的魅力来消除他的烦恼；指出——正如事实一样——一个人的性情不可能在没有长时间和更持续交往的情况下立即被察觉，文化修养也同样不会在几次努力和短时间内就被那些尝试她的人所认识。这样我恢复了他的愉快，通过这种相互体验，他更紧密地与我结合。\par

19. 随着时间的推移，当我们承认彼此的爱慕，并以哲学为目标时，我们彼此就是一切：同屋、共餐、亲密无间，有一个共同的人生目标，彼此的爱慕日益增长、更加牢固。对肉体吸引的爱，因其对象是短暂的，就像春天的花朵一样短暂。因为当燃料耗尽时，火焰无法存留，并随着点燃它的东西一同离去；当刺激物消逝时，欲望也不会停留。但属天且受约束的爱，因其对象是稳定的，不仅更持久，而且越看清美善，就越把那些爱慕同一对象的心彼此紧密地联结在一起。这就是我们超然之爱的法则。我感到自己正被过分地卷走，不知如何进入这一点，但无法克制自己不描述它。因为如果我遗漏了什么，似乎紧接着就变得紧迫重要，比我所优先提及的更为重要。如果有人试图暴虐地推动我前进，我就变得像章鱼一样——当被从洞里拖出时，它们用吸盘紧贴岩石，直到最后一个吸盘被施加必要的力才能分开。那么，如果你允许我，我就如愿了；如果不允许，我也必须自己从中取出。\par

20. 当我们这样，如同品达所说，以\Quote{金柱支撑的坚固厅堂}相互支持时，我们就在天主的恩宠和我们彼此之爱的共同影响下前进。噢！我怎能不流泪地提及这些事。\par

我们被相同的希望推动，从事于一种特别容易招致嫉妒的追求——学问。然而，我们不认识嫉妒，竞争对我们有益。我们争夺的，不是为自己赢得首位，而是将首位让给对方；因为我们把彼此的声誉当作自己的。我们似乎有一个灵魂，居住在两个身体里。如果我们不能相信那些主张\Quote{万物在万物之中}的学说，那么在我们的情况下，这一说法是值得相信的，因为我们如此活在彼此之中并彼此相共。我们两人的唯一事业就是美德，并为将来的希望而生活，在我们实际离开这个世界之前就已隐退其中。为此，我们的整个生活和行动都指向它，在诫命的引导下，我们彼此磨砺美德之武器；如果我说这不算什么大事——作为彼此正确与不正确的规则和标准。我们的伙伴不是最放荡的，而是最稳重的同伴；不是最好斗的，而是最和平的，他们的亲密最有益：因为知道沾染恶习比传递美德更容易；正如我们更容易感染疾病，而不是给予健康。我们最珍爱的学习不是最愉快的，而是最优越的；这是塑造年轻心灵走向美德或恶习的一种方式。\par

21. 我们向来知道两条道路：第一条较为贵重，第二条较为轻微；一条通向我们的神圣殿堂及其中的导师，另一条通向世俗的教师。其余一切——宴乐、剧场、集会、筵席——我们都留给那些愿意追求它们的人。因为依我之见，除却那通向德行并使修德者臻于完善的事以外，没有任何事是有价值的。不同的人从父亲、家族、职业、功业获得不同的名称；我们却只有一件大事和一个名号：成为并被称为基督徒。我们对此的看重，远胜于\Person{盖吉兹}（若这并非传说）对其戒指的转动——他的吕底亚王权全赖于此；也胜于\Person{迈达斯}对其黄金的追求——他因祈求凡物皆化为金而丧命，那是另一个弗里吉亚传说。我又何必提起极北族\Person{阿巴里斯}的箭，或阿尔戈斯的\Person{珀伽索斯}？对他们而言，空中飞行尚不及我们凭着彼此相助、向着天主的上升那样重要。\Place{雅典}尽管在属灵之事上对他人有害——这对虔诚者并非小事，因为该城在那些邪恶的财富（即偶像）上比希腊其他地方更为富足，人很难不被其崇拜者和信徒所裹挟——然而我们，因心志紧闭、固守防御，未受任何损伤。相反，看似奇特的是，我们反而因此更坚定了信仰，因为我们看清了他们的诡计与虚妄，以致在众神受拜之地也藐视那些神明。若有一条淡水河流过海洋，或有一动物能在吞噬万物的火焰中舞蹈，那么在我们所有同伴中，我们便是如此。\par

22. 而最美好的是，我们身边环绕着一群绝非低贱的伙伴，在他的教导和引领下，因同一目标而喜悦，如同徒步追随那吕底亚战车——他的行径与性情；于是我们声名远播，不仅在本族导师和同伴中，甚至遍及全\Place{希腊}，尤其为那里最杰出的人物所瞩目。我们甚至超越了希腊的疆界，这由许多证据得以显明。因为凡认识\Place{雅典}的人，都认识我们的导师；凡认识导师的人，也认识我们，我们成为谈论的对象，或被亲眼目睹，或被传闻所知，犹如一对显赫的伴侣。在他们眼中，\Person{奥瑞斯忒斯}和\Person{皮拉得斯}与我们相比实在微不足道，或如\Person{莫利奥尼之子}——荷马卷轴中的奇观，因患难与共、共执缰绳和鞭策的辉煌车技而闻名。但我无意间竟陷入自我赞美，而这是我绝不允许他人所为的。然而这也无怪，因为从他的友谊中，我在此处得了些益处；正如生前他助我于德行，自他离世后，他又助我于声名。但我必须回归正题。\par

23. 谁在未及白发之前，便已拥有如此程度的暮年智慧？因为\BibleRef{智 4:8}是\Person{撒罗满}用以界定暮年的。谁对长者和幼者——不仅对我们同代的人，甚至对远在我们之前的人——如此恭敬？谁因自身品格而较少需要教育？然而谁又以其品格而如此饱学？哪一门学问他不曾遍历？且成就非凡，穿越所有学问，如同无人能穿越其中任何一个；他在每一门学问上都达到如此高峰，仿佛那是他唯一的研究。艺术和科学中能力的两大源泉——天赋与勤勉——在他身上等量齐观。因为他既因所付辛劳而无需太多天资，又因其天资而无需劳苦；二者如此合作与统一，以致难以看出他更以哪方面著称。谁在修辞学——那如烈火般喷涌的能力——上有如此功力？尽管他的性情与修辞家迥异。谁在文法——那完善我们的希腊语言、编纂历史、主管韵律、为诗歌立法——上有如此造诣？谁在哲学——那真正崇高而高远的学问，无论是实践的还是思辨的，抑或是其中涉及逻辑证明的对立与争辩的部分（即所谓辩证法）——上有如此造诣？在辩证法中，若有必要，要逃脱他言辞的罗网，比逃离迷宫更为困难。对于天文、几何和数之比，他掌握得如此透彻，以至那些擅长此等学问的人也无法难倒他；但他鄙视过度的钻研，认为这对渴慕虔诚的人毫无用处。因此，既可钦佩他之所选，更甚于他所弃；或钦佩他所弃，更甚于他所选。医学——哲学与勤勉的成果——因他体质纤弱和照顾病人而成为必要。由此开端，他达到了对这门技艺的精通，不仅在经验和实践的方面，也在其理论和原理上。然而，这些虽甚辉煌，与他的人格操守相比，又算得了什么呢？凡与他接触过的人，都视\Person{米诺斯}和\Person{拉达曼堤斯}为儿戏——希腊人认为他们配得上阿福花草地和乐土原野，那是他们按照\WorkTitle{梅瑟之书}（也是我们的书）中对乐园的描绘而设想的，尽管名称不同，但所指即为此，只是换了名目。\par

24. 情况正是如此，他的大帆船满载着人性所能达到的一切学问；因为过了\Place{加的斯}再无航路。剩下的唯一需要，就是上升至更完善的生命，紧握我们一致同意的那些希望。我们离别的日子临近了，伴随着告别辞、送行、邀约返回、哀叹、拥抱与泪水。因为对于任何人，没有什么比离开\Place{雅典}和彼此分离，对曾是那里同伴的人更痛苦的了。在那场合，人们见到了一个值得记载的凄惨景象。我们周围聚集着同窗、同班同学以及一些教师，他们恳求、强迫、劝说，无论如何不会放我们走；说着、做着悲苦之人能做的一切。在此我既要指控自己，也要（尽管大胆）指控那神圣无瑕的灵魂。因为他详述急于回家的理由，得以战胜他们挽留他的愿望，他们虽不情愿，却被迫同意他离开。而我则被留在\Place{雅典}，一部分（说实话）是因为我被说服了——一部分是因为他背叛了我，他被说服抛弃并交出一个拒绝抛弃他的人给那些俘获他的人。这种事在发生之前令人难以置信。因为这就像将一个身体劈成两半，毁掉每一部分，或像将共享同一马槽同一轭的两头公牛分离，它们可怜地相互吼叫，抗议分离。不过，我的损失并未持续太久，因为我不能长期忍受被人看到可怜的样子，也不能向每个人解释我们的分离：所以在\Place{雅典}短暂停留后，我的渴望使我像\Person{荷马}笔下的骏马，挣断束缚者的缰绳，越过平原，奔向我的同伴。\par

25. 我们回来后，对尘世和舞台稍作迁就，足以满足大众的期望，并非出于对戏剧表演的偏爱，我们很快就独立了，并且从无髭少年的行列晋升为成人，在哲学的道路上大胆迈进；因为尽管嫉妒使我们无法再同在一处，我们却因热切的渴望而结合。\Place{凯撒勒雅}城占据了他，视他为第二奠基人和护佑者；但随着时间的推移，他偶尔会缺席，这是由于我们分离的必要，也是为了我们既定的哲学行程。由于要侍奉年迈的父母，以及一连串的不幸，使我一直与他分离，也许这并不合情理，但事实如此。我倾向于将我生命中所有的不一致与困难，以及哲学道路上的障碍（这有负于我的渴望与目标）都归咎于此原因。至于我的命运，让它随意引向天主所愿之处，但愿我的道路能因他的代祷而更好。关于他本人，天主对人类的丰盛之爱，\BibleRef{铎 3:4}以及祂对我们种族的眷顾，在诸多境遇中越来越辉煌地彰显了他的功绩之后，将他立为教会中一道显著而辉煌的光，提升他至司铎的圣座，使他通过\Place{凯撒勒雅}一城向全世界放射光芒。这是以何种方式呢？并非仓促提升，也非即时洁净他、使他智慧（如同现今许多求晋升者之所为），而是按灵性成长的正当顺序赐予他殊荣。\par

26. 因为我不赞赏我们中间有时存在的混乱和无序，甚至在那些管理圣所的人中也是如此。我不敢，也不公正，去指责他们全体。我赞同航海的习俗：先给未来的舵手划桨，然后才把他引到船尾，托付指挥权，让他坐在舵边，但前提是在漫长的击浪和观察风向之后。军事方面也是如此：士兵、将领、统帅。这种先后顺序对下属最有益、最有利。如果我们这里也如此，那就会大有裨益。但事实上，最神圣的职位在我们中间有成为最可笑之事的危险。因为晋升并不取决于德行，而是取决于恶行；神圣的宝座不是落于最配得者，而是落于最有势力者。先见未来的\Person{撒慕尔}在先知中，但被弃的\Person{撒乌耳}也在其中。\Person{撒罗满}的儿子\Person{勒哈贝罕}在君王中，但奴隶和背教者\Person{雅洛贝罕}也在其中。没有一个医生或画家不是先研究疾病的性质、或调和多种颜色、或练习素描：但一位主教却很容易找到，无需艰苦训练，名声刚建立，一下子就播种发芽，就像巨人传说那样。我们一日之间制造圣者，并吩咐那些没有受过教训、对他们的尊严毫无贡献（只有意愿）的人成为智者。因此，一个人满足于低微的职位，安于卑下的地位，而他是配得高位、深思过默启之言、通过许多法律使肉体顺服精神的人；而另一个人傲慢地居首，在高明者面前扬起眉毛，不为自己的地位战栗，也不因看到在他之下受纪律约束的人而惊恐；并且错以为自己不仅在地位上，也在智慧上超越他人，因地位影响而失去理智。\par

27. 我们的伟大而杰出的\Person{巴西略}并非如此。在这恩宠上，如同在其他一切事上，他都是公众的榜样。因为他首先向人民诵读\WorkTitle{圣经}，当时他已能解释\WorkTitle{圣经}，且不认为自己配得上在圣所中的这一等级，于是便在长老的席位上赞美主，随后又在主教的席位上赞美主，他获得此职并非出于隐秘或暴力，而是不求荣誉，反被荣誉所求，且视之为来自天主而非人情的恩赐。关于他主教职位的叙述须稍缓：让我们在他下属的职务上稍作停留，因为在我讲述的过程中，这事几乎被我忽略了。\par

28. 在他与治理此教会的前任之间出现了一场分歧：其根源与性质最好默然不提，但它确实发生了。那位前任在其他方面绝非庸碌之辈，且以其虔诚令人钦佩，正如当时的迫害及对他的反对所证明的，但他对\Person{巴西略}的感情却是人所难免的。因为\Person{摩摩斯}不仅抓住普通群众，也抓住最优秀的人，所以惟有全然不受此类情感影响并抵抗它们是属于天主的。教会中所有更卓越、更明智的部分都起来反对他，如果那些远离世俗、将自己生命奉献给天主的人比多数人更明智的话。我指的是我们时代的纳齐尔人，以及那些献身于这类追求的人。他们因自己的领袖被忽视、侮辱和拒绝而不悦，并冒险采取了一项极其危险的行动。他们决定叛离并自教会身体中分裂出来——这教会是容不得派系的——同时带走不少百姓，包括下层人士和有权势者。由于三个强有力的原因，此事极为容易。首先，此人（指巴西略）的声望超越了我所认为的我们时代任何一位哲学家，且若他愿意，能够鼓舞同谋者的勇气。其次，他的对手因就职时的骚乱而被本城怀疑是以专断方式获得提升，不符合法律和教规。此外，还有一些西方的主教在场，将教会中所有正统派成员吸引到自己一边。\par

29. 那么，我们这位高贵的友人，和平之主的门徒做了什么呢？他既不习惯抵抗诬蔑者或党派分子，也不认为战斗、撕裂教会身体是他的本分——这教会本就因其他原因而受到攻击，且因异端者的强大势力而处境艰难。在我的建议和热切鼓励下，他与我一同离开此地前往\Place{本都}，在那里主持默观之所。他本人也建立了一座值得提及的隐修院，因为他与\Person{厄里亚}和\Person{若翰}这两位苦修大师一同迎接旷野；他认为这对于他而言比针对当前局势设计任何有损其哲学精神的计划更为有益，也比在风暴中毁掉他正平稳前进的航程更为明智——当时争议的浪潮已经平息。尽管如此，他那富于哲思的隐退是多么奇妙，我们将发现他的回归更为奇妙。因为事情是这样的。\par

30. 当我们正忙于这些事务时，突然升起了一团充满冰雹的乌云，带着毁灭性的咆哮，袭击并攫取了它所冲击的每一座教会：一位最贪爱黄金、最敌视基督的皇帝，感染了这两种严重的疾病——贪得无厌和亵渎神明；他接续了那位迫害者，作为迫害者，又接续了那位背教者，虽然不是背教者，但对于基督徒，或者更确切地说，对于更虔诚和纯洁的基督徒群体——他们崇拜我所称为唯一真正虔诚和救恩道理的天主圣三——他并不更好。因为我们并不将神性分割成部分，也不以不自然的分离将那唯一不可接近的本性从其自身中放逐；也不以一种邪恶来医治另一种邪恶，用更不敬虔的分割与分离来摧毁\Person{撒伯流}的亵渎性混淆；这就是\Person{亚略}的错误，其名字就表明了其疯狂，他是扰乱和毁灭教会一大部分的祸首。因为他并非通过不相等的神性等级羞辱圣子来尊敬圣父。但我们承认圣父的唯一光荣，即独生子的平等；圣子的唯一光荣，即圣神的光荣。我们认为，使三者中的任何一位从属于其他，就是摧毁整体。因为我们崇拜并承认它们按其特性为三位，但按其神性为一体。然而他（指亚略）却没有这样的观念，因为他不能向上仰望，反而被引导他的人所贬低，并胆敢连神性的本性也一同贬低，成为一个邪恶的受造物，将尊严置于奴役之下，并使那未受造、永恒的本性与受造物并列。\par

31. 这就是他的心思，他就以如此不敬向我们开战。因为我们不得不认为，这无异于一次蛮族入侵——不是摧毁城墙、城市和房屋，以及其他由人手建造、尚可修复、价值微薄之物，而是将破坏倾泻在人的灵魂上。一支可观的军队加入了他的攻击：教会的邪恶统治者，他那遍布世界的帝国的残暴总督。有些教会他们已掌控，有些正在进攻，另一些则指望通过皇帝已然施加的影响以及他所威胁的暴力来夺取。但在他们企图诱惑我们自己的计划中，他们的信心特别建立在那些我所提及之人的器量狭隘、我们主教之缺乏经验，以及我们中间盛行的软弱之上。这场斗争将十分激烈：众多部队的热忱绝非卑劣，但他们的阵列却虚弱无力，因为缺少一位领袖和战略家，以圣言和圣神的能力为他们奋战。那么，这个高贵、宽宏大量、真正热爱基督的灵魂做了什么呢？无需许多言辞来敦促他前来援助。他一看到我负此使命——因为为信德而战对我们两人是共同的——就立即答应了我的恳求；并以一个基于灵性理由的绝妙区分做出决定：挑剔讲究之时（如果我们果真会允许自己有这种情绪的话）是安全之时，而忍耐则在必要之时。他立刻随我从本都返回，像一个热心的志愿者投身于为濒危真理而战，并奉献自己为他母亲教会服务。\par

32. 那么，他实际的努力是否辜负了他起初的热忱？它们是否由勇气而非审慎引导，或者由技艺引导，同时他却畏惧危险？或者，尽管在这些方面达到了无与伦比的完美，他心中是否还残留着一丝恼怒？远非如此。他立即完全和解，并参与了每一个计划和努力。他清除了我们道路上的一切荆棘和绊脚石——敌人正是倚靠这些来攻击我们。他抓住一个，握住另一个，推开第三个。对某些人，他成了坚固的墙垒和壁垒，如\BibleRef{耶 1:18}所说；对另一些人，他成了击碎石头的斧头，或荆棘中的火焰，如神圣圣经所说，轻易烧毁了那些侮辱天主性的柴捆。如果他的巴尔纳伯（讲述并记录这些事的人）在斗争中曾有助于保禄，那么功劳当归于保禄，因为他选择并使他成为自己战斗中的同伴。\par

33. 于是敌人失败了，这些卑鄙之人平生第一次遭到卑劣的羞耻和挫败，学会了再也不敢轻视卡帕多细亚人——世上所有人中，他们的特质是在信德上的坚定和对天主圣三的忠诚奉献；归于祂的是他们的合一与力量，并从祂那里获得比他们所能给予的更大更强的援助。巴西略接下来的事务和目标是安抚主教，消除猜疑，说服所有人：所感受到的恼怒是由于恶者的诱惑和努力，他嫉妒德行上的和睦：他谨慎地服从于服从和灵性秩序的法则。于是他带着训导和建议去拜访主教。在顺从主教意愿的同时，他对主教来说就是一切：良好的顾问、熟练的助手、天主圣意的阐释者、行为的引导、他晚年的手杖、信德的支柱、内部最忠实的人、外部最务实的人——总之，他多么倾向于善意，就曾多么被认为敌对。于是，教会的权力几乎（如果不是完全）落入了他的手中，与教区主教相当。因为作为对他善意的回报，他得到了权威。他们的和谐与力量的结合是奇妙的：一人是民众的领袖，另一人则是领袖的领袖，如同驯狮者，巧妙地安抚着权力的拥有者。因为主教刚刚就任，仍多少受世俗影响，尚未具备圣神的事物，在围攻教会的敌人激流之中，他需要有人扶持和协助。因此他接受了这联盟，并想象自己是战胜了那战胜他之人的胜利者。\par

34. 他对教会的关心与保护，还有许多其他证据：他对城中总督及其他最有权势者的勇敢无惧；裁决争端，毫不迟疑地被接受，凭他一句话便生效，他的意向被视为决定性的；他对困苦者的支持，大多是精神上的，也有不少是肉体上的——因为这也常藉着善意影响灵魂，使其屈服；对贫苦人的支援，对客旅的款待，对少女的照顾；为隐修生活所立的成文与不成文的规诫；祈祷的安排，圣所的装饰，以及其他真正属神的人为天主工作而造福百姓的方式：其中一件尤为重大且值得注意。曾有一次饥荒，是前所未有最严重的。全城陷于困苦，毫无援助或解救的源头。因为沿海城市能轻易承受这样的匮乏时期，用本地产品交换进口物资；但像我们这样的内陆城市，既不能以盈余获利，也不能满足需要——无论是出售已有的，还是输入所缺的；然而所有这类困苦中最艰难的部分，是那些拥有存粮之人的麻木与贪得无厌。因为他们窥伺时机，利用苦难牟利，以不幸为生；毫不理会\BibleQuote{怜悯贫苦人的，就是借给上主}（\BibleRef{箴 19:17}），也不理会\BibleQuote{屯积粮食的，必为人民所咒骂}；以及其他对仁爱者的许诺、对残忍者的威胁。但他们过于贪得无厌，政策失当；因为他们向同胞关闭了自己的心肠，也就向自己关闭了天主的心肠，忘记了他们更需要天主，甚于别人需要他们。这些买粮食卖粮食的人便是如此，既不尊重同胞，也不感谢赐予他们所有的天主，而其他人却在困顿之中。\par

35. 他确实不能用祈祷使面包从天上降下（\BibleRef{出 16:15}），滋养在旷野中逃出的百姓；也不能从因倒空而满溢的器皿深处，无偿地供应食物之泉（\BibleRef{列上 17:14}），以惊人的回报来支持那位支持他的妇人；更不能用五个饼喂饱数千人，而饼的碎块竟成为许多桌宴席的额外供应（\BibleRef{玛 14:19}）。这些是梅瑟、厄里亚以及我的天主的事迹，他们也是从祂获得能力。或许这也是他们时代与环境的特点：因为\BibleQuote{奇迹是为不信的人，而不是为信的人}（\BibleRef{格前 14:22}）。但他却以同样的信德，筹划并执行了与之相应且倾向同一方向的事。因为他凭自己的言语和劝告，打开了拥有粮食之人的仓库，于是按圣经所载，将食物给饥饿的人（\EditorialNote{依撒意亚 58:7}），用面包满足穷人，在饥荒时喂养他们，以美物充满饥饿的灵魂。这是如何做的呢？因为这绝非对他赞词的微小增添。他聚集了饥荒的受害者，以及一些刚刚康复的人——男人、女人、婴儿、老人，每一个在困苦中的年龄层——并征集了各种能缓解饥荒的食物，在他们面前摆上汤盆和在我们中间保存下来的肉类，这些是穷人的食物。然后，他效法基督的服务，基督束上毛巾，不嫌卑下地洗门徒的脚；他借助自己的仆人以及同仆的协助，亲身照顾那些需要之人的身体和灵魂，将个人的尊重与满足他们的急需结合起来，从而给予了双重的缓解。\par

36. 这就是我们年轻的粮食供应者，第二个若瑟：尽管我们对他还能说更多。因为第一个若瑟从饥荒中获利，用他的仁爱买下了埃及，以丰年治理预备荒年，利用别人的梦境达到此目的。但这一位的服务却是无偿的，他对饥荒的援助未获分文利益，只有一个目的：透过善意的对待赢得善意，并藉他分发的粮食赢得天上的祝福。此外，他提供了圣言的滋养，以及那更完美的恩赐与分配，那真是来自天上的——如果圣言是天使的食粮，灵魂靠它得喂养、得畅饮，这些灵魂饥渴于天主，寻求一种不会消逝、不会短缺、永远存留的食物。这食物，他——我所认识的最贫穷、最匮乏的人——丰富地供应了，为解除的不是粮食的饥荒，也不是水的干渴，而是对那圣言的渴求（\BibleRef{亚 8:11}），这圣言真正赋予生命、滋养人，并使那适当进食的人成长至属灵的成年。\par

37. 在这些及类似行动之后——我何必一一细述？——当那位名字象征虔诚的主教，在\Person{巴西略}的怀中安然离世后，他被升上了崇高的主教宝座，这并非没有困难，也并非没有他本地主教们的嫉妒争斗，他们一方竟有城中最大的恶棍支持。但圣神必须获胜——而且果然决定性地获胜。因为圣神从远处带来了一些杰出而热切于虔诚的人为他傅油，与他们同来的还有新\Person{亚巴郎}，我们的宗主教，我指的是我的父亲，关于他发生了一件非凡的事。因为他虽因年岁衰迈，疾病缠身几至最后一息，仍冒险踏上旅途，凭圣神的助佑以他的投票给予支持。简言之，他被放入轿中，如同尸体被放入坟墓，回来时却焕发青春活力，昂首挺胸，因了覆手和傅油而获得力量——甚至可以说，因了那受傅者的头而获得力量。这必须添加到古时的实例中，证明劳苦带来健康，热忱的目标使死者复活，年老者在圣神的傅油下跳跃起来。\par

38. 他既堪当此首席之职——按常理，唯有活出如此生命并赢得如此恩宠与尊重者方宜担任——他的此后行为并未辱没自己的修行，也未辜负那信赖他之人的期望。反而他超越自己，正如迄今所示他超越他人一般；他在这点上的见解极为卓越且富有智慧。因为他认为：仅在私人生活中远避邪恶并达到某种程度的善，固然是德行；但身为领袖与治理者，尤其是在此职上，若不能远超众人，并以不断进步使自己的德性与其尊严和宝座相称，便是恶行。因为身居高位者难以保持中庸，而必须藉德行之卓越将子民提升至黄金中道。或者，为更圆满地处理此事，我以为其结果与我主救主的情形相同，也与我猜想每一位特具智慧之人的情形相同，当他以那超越我们却又属于我们的形态与我们同在时。因为福音记载，祂在智慧、身量及恩宠上日益增长（\BibleRef{路 2:52}），并非这些品质在祂身上有所增长——因为那起初便成全的，如何能变得更成全？——而是逐渐显明并彰显出来。所以我认为巴西略的德行虽未增加，却在此时获得了更广泛的施展，因其权能为他提供了更丰富的素材。\par

39. 他首先表明，他的权位并非来自人的恩惠，而是天主的恩赐。这一点也将由我的行为所证明。因为在那段时期，有哪一种哲学研究他不曾与我共同进行？故此，人人都以为我会在事后跑去迎接他，表示我的喜悦（这或许是任何其他人都可能做的），并依照我们友谊所推断的，要求分享他的权柄，而非在他身旁统治。但我因极度忧虑当时所可能引发的烦扰与嫉妒，尤其是他的处境依然痛苦而动荡，我便留在家中，强力克制我热切的愿望；巴西略虽责备我，却接受了我的借口。当我后来抵达时，出于同样的理由，我拒绝了此席位及在司铎中的尊贵地位，他便仁慈地不予责备，反倒赞扬我；他宁愿让一小撮不明我们策略的人指责他骄傲，也不愿做任何违反正当理性及其决断的事。的确，一个人能如何更好地显示其灵魂超乎一切谄媚与奉承，并且其唯一目标是义的法则，正是通过他以如此方式对待我——他视我为朋友与同工中的首要人物。\par

40. 他的下一步工作是平息并藉宽宏的对待消除针对他的反对；这毫无谄媚或卑躬屈膝的痕迹，而是以极为侠义和宽宏的方式；不仅着眼于当前的紧急情况，也旨在培养未来的顺服。因为他看到，温和导致松懈与懒散，严厉则引发固执与任性，他便能藉结合两者避免每种路线之危险，将他的纠正与体谅相混合，柔和与坚定相混合；他大多藉行为而非言词影响人：不是以技巧奴役他们，而是以善良赢得他们，并藉少用而非常用权柄来吸引他们。最为重要的是，他们被引领认识到他智力的优越及德性的不可企及，认为唯有站在他一边并受他统率才是安全的，唯一危险在于与他作对，并以为与他分歧即与天主疏离。于是他们甘心屈服投降，如同遭遇雷击般自我降伏，争相以道歉互抢先机，将当初的敌意强度变为同等的善意强度，并在修行中长进——他们发现这乃是唯一真正有效的防御。对于少数例外者，他则放过不加理会，因他们的恶意无可救药，他们消耗精力自损其身，正如锈蚀吞噬其所附之铁一般。\par

41. 家中的事务既已按他心意安定下来——这在不信他、不认识他的人看来是不可能的——他的谋划变得更为宏大，境界更为高远。因为其他人全都只盯着眼前的地面，注意力集中于自己即刻的关切，只要这些安妥，便不再烦扰自己，无力从事任何伟大而侠义的谋划或事业；而他，虽在其他方面节制，却在这件事上不能节制；他昂首挺胸，精神之眼环顾四周，涵盖了救恩圣言所传播的整个世界。当他看到天主的伟大产业——由祂的言语、法律和苦难所赎买的，那神圣的国度、王家的司祭（\BibleRef{伯前 2:9}）——陷入如此悲惨的境地，被分裂成万千见解与错误；又看到那从埃及（即不敬神与黑暗无知的埃及）被领出并移植的葡萄树，已长得如此美好而无边无际，整个大地都被它的荫影遮蔽，高过山岭和香柏，如今却遭受那邪恶的野猪——魔鬼——所蹂躏，他不能仅限于默默哀叹不幸，只向天主举手祈求祂驱散迫在眉睫的灾祸，而自己却沉睡；他觉得自己必须付上代价来援助她。\par

42. 有什么比这场灾祸更令人悲痛，更迫切地呼唤一位高瞻远瞩的人为公益而尽力呢？个人的成败对团体无关紧要，但团体的成败必然涉及个人的境遇。怀此信念与目的，那守护和庇护团体的人（正如撒罗满真实所说：敏锐的心如骨骼中的蛀虫，麻木不仁者则欣然自信，而同情心是痛苦的根源，不断思虑耗损心灵），他因此被许多创伤所折磨痛苦；像约纳和达味一样，他愿自己死去（\BibleRef{纳 4:8}），不给自己眼睛睡眠，也不让眼皮打盹，他将残存的血肉消耗在思虑中，直到发现那邪恶的疗方：并向天主和人寻求援助，以阻止那普遍的烈火，驱散那笼罩我们的阴霾。\par

43. 他的一个策略极其有用。在尽己所能的退省和私下神交之后，他考虑了所有人性论据，并潜入圣经的深奥之处，草拟了一幅虔诚教义的纲要，通过与他们敌对的交手和攻击，击退了异端的大胆进击：以口舌的短兵相接，推翻了那些逼近之人；以墨翎的飞箭，打击了那些远处之人，这墨箭并不逊于碑刻；他并非像那犹太人那样，仅为一个弱小民族关于饮食、暂时祭献和肉身洁净的指示（\BibleRef{希 9:10}），而是为天下万邦和世界各处，关于真理的圣言，我们救恩的源泉。再者，既然盲目的行动和不实际的推论都同样无效，他就以行动带来的援助补充了他的推论：他拜访、致函、接见、教导、申斥、谴责（\BibleRef{弟后 4:2}）、威胁、责备，为邦国、城邑和个人辩护，设想各种救助，从各处觅得疾病的良药：他成为第二个贝匝肋耳，神圣会幕的建筑师（\BibleRef{出 31:2}），将一切材料和艺术应用于工程，并将所有元素结合为和谐而超越的美。\par

44. 何必再多说细节？我们再次受到反基督皇帝的袭击，那位信仰的暴君，带着更丰盛的不敬和更猛烈的攻势，因为争端必须与更强的对手进行，正如那不洁的邪魔，当它离开人游荡之后，又返回带着更多的魔鬼住进那人里面，正如我们在福音中所听到的（\BibleRef{路 11:24}）。他效法了这个灵，既重开以往败北的争战，又增添了原本的努力。他认为这是一件奇怪而难以容忍的事：他，统治这么多邦国，赢得如此声望，并用不敬之权使周围所有人屈服，战胜了每一位对手，竟被一个人、一座城公开打败，因而不仅受到那些引导他的不虔诚者的嘲笑，而且按他所想，也受到所有人的嘲笑。\par

45. 据说波斯王在远征希腊时，不仅因他愤怒和骄傲中所率领的各族人眾數量而得意，从而被激发起过分的威胁；而且他还想通过使他们对他的恐惧来进一步恐吓他们，这是因为他对元素的新奇处理。人们听闻了一片奇怪的陆地与海洋，那是新创造者的杰作；一支军队航行于旱地，徒步于海洋；岛屿被移走，海洋被鞭打，以及那个军队和远征的所有其他疯狂行径；虽然这些事使卑贱者惊恐，但在勇敢坚毅的人眼中却是可笑的。在我们的远征中不需要任何这类东西，但更坏、更有害的是，据说皇帝说了并做了这件事。他向天张口，说亵渎至高者的话，他的舌头遍行世界。受默感的达味早在我们的日子之前就出色地描述了那位使天俯就地，并以造化衡量那超世俗的本性，而造化甚至无法容纳这本性，尽管出于对人的仁慈，它在一定程度上来到我们中间，以便将躺在地上的我们吸引到自己身边。\par

46. 他最初的狂妄行径确实狂怒，而他最后对我们的攻击更加狂怒。我先说什么？流放、驱逐、没收、公开和秘密的阴谋、在时间允许时的说服、在说服不可能时的暴力。那些像我们一样坚持正统信仰的人被逐出他们的教堂；另一些赞同皇帝那毁灭灵魂的教义的人被强行安插进来，他们乞求不敬的证明，并签署比这些更为苛刻的声明。司铎在海上的焚烧，不敬的将军——不是那些征服波斯人、制伏西徐亚人或降服任何其他野蛮民族的人，而是那些攻打教堂、在祭台上胜利跳舞、用人和牺牲的血玷污无血祭献、并侮辱贞女端庄的人。目的是什么？将圣祖雅各伯逐出，而代之以那在出生前就被憎恨的厄撒乌（\BibleRef{罗 9:11}）。这就是他最初狂妄行径的描述，即使现在，只要回忆和提起这些事，就会使我们大多数人流泪。\par

47. 因此，当他遍历各地之后，发动进攻，企图征服这座坚不可摧、令人敬畏的教会之母——真理尚存的最后一颗未熄的火星——时，他发现自己首次失了算。因为他像一枚击中坚硬物体的弹丸那样被弹回，像一根断了的缆绳那样退却。他所遇到的教会主教便是如此，这堵壁垒使他的一切努力都化为乌有。其他细节可以从那些亲身经历并讲述这些事的人口中听到——而讲述这些事的人无不全然亲历。但所有知晓当时那些斗争、袭击、许诺、威胁、以及他先期派来试图说服我们的那些使节——那些司法和军事官员、来自后宫的人（他们是在女人中的男人、在男人中的女人，他们的唯一男子气概在于不虔敬，并且因不能自然放纵，只能用舌头行淫）——以及首席厨师\Person{乃步匝辣当}，他用他手艺的武器威胁我们，却被自己的火焰所派遣——所有人必定充满钦佩。而特别令我惊叹、即使我想也不能略过的事，我将尽可能简洁地描述。\par

48. 谁没有听说过当时的总督？他本人是由另一方接纳（或许是交付）圣洗的，对我们表现得如此傲慢无礼；他试图通过超越指令的字面意思、在每一细节上迎合他的主人，来保证并维护自己的权力。尽管他对教会大发雷霆，摆出狮子般的面孔，像狮子一样咆哮，以致大多数人不敢靠近他，但我们高贵的主教被带进——毋宁说是他自己走进——他的法庭，仿佛是被邀请赴宴，而不是受审。我如何能完全描述总督的傲慢，以及圣人应对它的智慧？\Quote{巴西略先生，}他直呼其名，还不屑于称他为主教，\Quote{您胆敢像无人敢做的那样，抗拒并反对如此伟大的君主，这是什么意思？} \Quote{在哪方面？}我们高贵的勇士说，\Quote{我的鲁莽表现在哪里？因为我还不明白。} \Quote{在您不尊重您君主的宗教，而其他所有人都已屈服归顺时。} \Quote{因为，}他说，\Quote{这不是我真正君主的旨意；我既是天主的受造物，且被命定为天主，就不能敬拜任何受造物。} \Quote{那么我们呢，}总督说，\Quote{在您看来算什么呢？我们这些向您下令的人，难道什么都不是吗？您对此怎么说？与我们为伍，难道不是一件大事吗？} \Quote{我不否认，}他说，\Quote{您是一位总督，而且是一位显赫的总督，但终究不比天主更尊贵。而与您为伍确是一件大事——因为您是天主的受造物——但与我任何一个子民为伍也是如此。因为信德，而非个人地位，才是基督信仰的标志。}\par

49. 于是总督激动起来，从座位上站起，怒气冲天，言语也更加严厉。\Quote{什么？}他说，\Quote{您不怕我的权柄吗？} \Quote{怕什么？}巴西略说，\Quote{它怎能影响我？} \Quote{怕什么？怕我权力中的任何一种手段。} \Quote{有哪些？}巴西略说，\Quote{请告诉我。} \Quote{抄家、流放、酷刑、死亡。} \Quote{您没有别的威胁了吗？}他说，\Quote{因为其中没有一样能触及我。} \Quote{这怎么可能？}总督说。\Quote{因为，}他回答，\Quote{一个一无所有的人，是抄家所不能及的——除非您要我那些破烂衣衫和几本书，那是我仅有的财产。流放对我来说是不可能的，我不受任何地点的限制，既不认为我现在居住的土地是我的，也不认为我将被扔到的那片土地是我的；或者更确切地说，我把一切土地都看作天主的，我是祂的客人和寄居者。至于酷刑，对于一个身体已经不再存在的人，它们能起什么作用？除非您指的是第一击，因为这唯独在您手中。死亡是我的恩人，因为它将更快地把我送到天主那里，我为祂而活、而存在、几乎已死，并且我长久以来一直在奔向祂。}\par

50. 总督听了这番话，大为惊异，说：\Quote{从没有人如此大胆地对\Person{莫德斯图斯}说过这样的话。} \Quote{也许，}巴西略说，\Quote{您从未遇到过一位主教，或者在他为这些利益辩护时，他会使用完全相同的语言。因为按照我们律法的诫命，我们在通常的情况下是谦逊的，对每个人都顺从。我们不可傲慢地对待任何一个普通人，更不用说如此伟大的君主了。但是，当天主的利益处于危急时，我们就不顾其他一切，只以这些为唯一目标。火与剑、野兽与撕裂皮肉的铁耙，我们以此为乐，毫不畏惧。您可以进一步侮辱和威胁我们，尽您权力所能做任何事。皇帝本人也可以听到这一点——无论用暴力还是劝说，你们都不能使我们与不虔敬同流合污，即使你们的威胁变得更为可怕。}\par

51. 这场交谈结束时，总督被巴西略的态度所折服，确信他完全不为威胁和权势所动，便将他从法庭上释放，先前的威胁态度被某种尊敬和顺从取代。总督本人迅速觐见皇帝，说：\Quote{陛下，我们被这教会的主教打败了。他不怕威胁，辩才无碍，不为说服所动。我们必须找一个更软弱的人来试探；在这种情况下，要么诉诸公开的暴力，要么任我们的威胁失效。}于是皇帝被对巴西略的称赞所迫而谴责自己的行为（因为就连敌人也能钦佩一个人的美德），就不允许对他使用暴力；而且，如同被火软化却仍是铁的铁一样，虽从威胁转为赞赏，却因羞耻而无法改变自己的道路，不与巴西略共融，反而试图用他能找到的最合理的借口为自己的行为辩护，如下文所示。\par

52. 因为他带着全部随从进入了教堂；那时正是主显节，教堂里挤满了人，他通过置身于民众之中表达了合一。此事不可轻忽。他进入时，被圣咏的雷鸣般回响、会众的人山人海以及遍布圣所及其周围的天使般的而非人间的秩序所震撼：而巴西略则主持着他的子民，挺身站立，正如圣经论及撒慕尔所说，\BibleRef{撒上 19:20}身体、眼睛和心灵毫不动摇，仿佛没有发生任何新事，而是专注于天主和圣所，可以说，就像一尊雕像，而他的侍从则怀着敬畏侍立在他周围。见此景象——这确实是无与伦比的景象——他因人性的软弱而被打败，他的眼睛变得模糊眩晕，内心充满恐惧。这一点当时大多数人尚未察觉。但当他在天主的祭台上奉献礼品时——他必须亲自行这事，因为照常没有人会协助他，因为不确定巴西略是否接纳他——他的情感泄露了出来。因为他步履蹒跚，若不是圣所中有人伸出手扶住他摇晃的脚步，他就会跌倒在地上，令人惋惜。此情此景，如此而已。\par

53. 至于他与皇帝会谈的智慧——皇帝在我们所谓的共融中进入帐幔内，与他见面交谈，正如他长久以来所渴望的那样——我还能说什么呢，只能说是充满神感的话语，被朝臣和我们这些与他们一同进入的人所听到。这是皇帝对我们友好情感的开端和首次确立；这次接待所产生的印象终结了大部分如洪水般侵袭我们的迫害。\par

54. 另一件事与我提到的那些事同样重要。恶人得胜，将他放逐的法令签署了，这完全满足了那些促成此事的人。夜幕降临，马车备好，我们的仇敌欢欣鼓舞，虔诚者陷入绝望，我们围住这位热心的旅者，他荣誉的耻辱什么都不缺。接下来怎样？天主使之归于无有。因为那位击打埃及长子（\BibleRef{出 12:29}）以惩罚其严苛对待以色列的天主，也击打了皇帝的儿子，使其患病。何等迅速！放逐的判决与疾病的判决同时发生：恶书记的手被阻止，圣人得以保全，虔诚之人因那使皇帝的傲慢恢复理智的热病而出现在我们面前。还有什么比这更公正或更迅速呢？事情经过如下：皇帝的孩子患病，身体疼痛。父亲为之痛苦，因为父亲能做什么呢？他四处寻求帮助，召来最好的医生，极其热切地恳求，并伏在地上。苦难甚至使皇帝谦卑，这不奇怪，因为达味为他儿子所受的类似痛苦，圣经为我们记载了（\BibleRef{撒下 12:16}）。但无论何处都找不到医治这恶疾的方法，于是他求助于巴西略的信德，没有亲自召见他，因对最近的虐待感到羞愧，而是委托给他最亲近的朋友们。巴西略一到，没有表现出任何别人可能会有的延迟或勉强，疾病立刻缓解，父亲怀抱更好的希望；如果他没有将咸水与淡水混合，即同时求助于异端者又召请巴西略，孩子本应恢复健康并得以保留在父亲的怀抱中。这确实是当时在场并分担痛苦的人们的信念。\par

55. 据说同样的不幸也降临在总督身上。他也因病不得不俯伏在圣人的手下；实在，对有理智的人而言，天主的眷顾带来教训，而痛苦常比顺遂更好。他病倒了，流泪痛苦，派人去请巴西略，恳求他，呼喊说：\Quote{我承认你是对的；只求你救我！}他的请求得到了应允，正如他本人所承认的，并且说服了许多先前不知此事的人；因为他不断惊叹并描述这位主教的能力。这就是他在这些事例中的行事及其结果。那么，他是否以不同的方式对待他人，为琐事而争执不休，或者未能达到哲学的高度，做那些不值得称赞、最好略过不提的事呢？绝非如此。那曾兴起恶人哈达得攻击以色列的\BibleRef{列上 11:14}，也兴起本都省的总督来攻击他；名义上是因一个卑微妇人的烦扰，而实际上是邪恶对抗真理的斗争的一部分。我略过他对巴西略的其他侮辱，或者说，对天主的侮辱；因为战斗正是针对天主并为着天主而进行的。然而，其中有一件事特别给攻击者带来耻辱，并高举了对手——如果哲学和哲学上的卓越是一件伟大而崇高的事——我将详细叙述。\par

56. 一位法官的参事正试图强迫一位出身高贵、丈夫新丧的贵妇接受一桩她不情愿的婚姻。她无法逃脱这种高压手段，便采用了一个既谨慎又大胆的计策。她逃到圣洁的祭台前，将自己置于天主的保护之下，以免受暴行。我以天主圣三之名——如果我可以在我这篇颂词中引入一点法庭风格的话——请问，应该做什么？我不是说伟大的巴西略——他在这些事上为我们所有人制定了法则——而是任何一位远不如他的司铎，难道不应该允许她的诉求，保护她、照顾她，伸手捍卫天主的仁慈和尊崇祭台的法律吗？难道不应该甘愿做任何事、受任何苦，也不参与任何不人道的阴谋，同时亵渎祭台和她在那里寻求庇护的信德吗？不！被挫败的法官说，所有人都应服从我的权威，基督徒应背叛自己的法律。他所要求的女请求者被不惜一切代价地保护着。因此，他愤怒之下，最终派了几名官吏搜查圣人的卧室，目的是羞辱他，而非出于任何必要。什么！搜查一位如此无欲无求、天使敬畏、妇女甚至不敢直视的人的住所？不仅如此，他还传唤他，迫使他为自己辩护；而且态度毫不温和友善，如同对待罪犯。当巴西略出现时，像我的耶稣站在比拉多审判台前一样，他坐在审判席上，充满愤怒和骄傲。然而，雷霆并未降下，天主的剑仍闪烁发光，等待时机；而祂的弓虽然张开，却被克制。这确实是天主的习惯。\par

57. 再考虑另一场我们这位勇士与其迫害者之间的斗争。当他的破旧大披肩被命令扯去时，他说：\Quote{如果你愿意，我也脱去我的外衣。}他枯瘦的身体受到鞭打的威胁，他提出愿意被刺梳撕裂，并说：\Quote{这样的撕裂将治好我的肝病，如你所见，它正在消耗我。}这就是他们的辩论。但当全城的人目睹这暴行和共同的危险——因为每个人都认为这种傲慢是自身的危险——时，全城怒火中烧；如同一窝蜂被烟熏起，一个接一个地被搅动并起身，每个种族、每个年龄层，尤其是来自小武器工场和皇家织布厂的男人们。因为这些行业的工人由于享有自由而特别暴躁和勇敢。每人手持当时使用的工具，或手边任何东西。他们手持火炬，在飞石中，举起棍棒，一起奔跑呼喊，同心协力。愤怒造就可怕的士兵或将军。妇女们在这种情形下被激怒，也并非手无寸铁。她们的发针成了长矛，不再像妇女，而因热忱的力量获得了男子的勇气。简而言之：她们认为将分沾杀死他的敬虔，并视最先动手对付胆敢如此行事之人为最敬虔。那么，这位傲慢而大胆的法官作了什么？他陷入可怜的困境，乞求怜悯，前所未有地卑躬屈膝，直到那位不流血而成为殉道者、不挨打而赢得荣冠的人到来；现在他以个人影响力平息了民众，并救下了那位侮辱他而今寻求他保护的人。这是圣人天主的行为，祂运作并改变一切为至善，祂拒绝骄傲的人，却赐恩宠给谦卑的人。\BibleRef{雅 4:6} 那曾分开红海、阻断约旦河、掌管自然元素、并伸出援手为拯救祂流亡的子民而树立凯旋标记的祂，为何不能也拯救这人脱离危险呢？\par

58. 这就是在天主的眷顾下，与世界之战的终结与幸运的结局，一个与他的信德相称的结局。然而，与主教们及其盟友的战争却在此处立即开始，这既带来了巨大的耻辱，也对他们的属民造成了更大的伤害。因为，当他们的主教们如此行事时，谁能说服他人节制呢？长久以来，他们对他怀有恶意，出于三个理由。他们在信德的事上与他意见不合，除非绝对被迫服从多数信徒，否则绝不认同。他们也并未完全放下因他的当选而对他怀有的怨恨。而最令他们痛心、却又最羞于承认的是——他的声望远超过他们。此外，还有一个进一步的分歧重新点燃了其他争端。当我们的国家被划分为两个省和两个都主教区，且原属前者的很大一部分被划归新省时，这再次激起了他们好斗的情绪。一方认为，教会的边界应依世俗边界而定；因此，他主张那些新划归的地区属于他，并且与原先的都主教分离。另一方则坚持古老的传统以及从我们祖先传下来的划分。许多痛苦的后果随之而来，或者正在未来的母腹中挣扎。不义地召集了由新都主教主持的会议，并夺取了收益。一些教会中的司铎拒绝服从，另一些则被争取过去。结果，教会的事务陷入了可悲的纷争与分裂之中。确实，新奇的事物对人有一种魅力，他们轻易地将事件转向对自己有利的方向，而推翻已经建立的事物比恢复被推翻的事物更容易。然而，最激怒他的是，在他眼前流过的\OriginalTerm{金牛山}{Taurus}的收益竟归于他的对手，还有\Person{圣奥雷斯特}的奉献，他极想从中收获果实。他甚至有一次在巴西略骑马经过自己的道路时，抓住他的骡子的缰绳，并带着一群强盗挡住了去路。这是何等冠冕堂皇的借口——对他属灵子女和交托给他的灵魂的关怀，以及对信德的捍卫——这些借口掩盖了那最普遍的恶习，即贪得无厌；此外，还有向异端者缴纳费用的不正当性，而任何令他不满的人都被称为异端者。\par

59. 然而，身为天上真耶路撒冷的都主教，这位天主之人既没有被堕落之人的失败所冲昏，也不容许自己忽视这种行为，更没有渴求任何不足的补救措施。让我们看看这补救是多么伟大而奇妙，或者，我应该说，多么与他的灵魂相称。他将纷争作为教会增长的原因，而这场灾难，在他极其高明的管理下，导致了该国主教人数的增加。由此产生了三个最理想的结果：对灵魂更大的关怀，每个城市自主管理自己的事务，以及该地区战事的停止。恐怕我自己也被当作这个计划的一个附属品。我很难用其他词语来描述这个处境。尽管我钦佩他整个的行为，其程度甚至超出我的表达能力，但唯独这一点，我发现不可能赞同，因为我将承认我对它的感受，尽管这些感受从其他来源你们多数人已并非不知。我指的是他对我的改变与不信实，这是一种连时间也未能抹去的痛苦。因为我这一生所有矛盾和混乱的根源；它夺去了我实践哲学或至少享有哲学名声的机会——尽管后者微不足道。你们或许会允许我为他的辩护是：他的观念是超人的，在死前已变得超越世俗影响，他唯一的兴趣在于圣神的事；而他对友谊的重视，在他准备——也只在此时——忽视友谊的要求时，丝毫未减，当这些要求与他首要的对天主的责任相冲突，且他所追求的目标比他不得不搁置的利益更为重要时。\par

60. 我担心，为了避免那些想了解所有关于他能说之事的人的冷漠指责，我又会遭受那些以中庸为理想的人对我冗长的谴责。因为巴西略本人对后者最为尊重，特别致力于那句箴言\Quote{在一切事上，中庸是最好的}，并终生付诸实践。然而，我不顾那些追求过度简洁或过度冗长的人，就这样继续我的演讲。不同的人以不同的方式获得成功，有些人只致力于众多美德形式中的一种，但据我所知，迄今为止没有人能在所有方面都达到最高的卓越；在我看来，最好的人是在最广阔的领域里赢得桂冠，或在某个单一特定方面获得尽可能高的声誉。然而，巴西略的声望如此之高，以至于他成了人类的骄傲。让我们这样来思考问题。有人致力于贫穷、一无所有的生活、免于多余之物吗？除了他的身体和必要的遮体之物外，他还拥有什么？他的财富就是一无所有，他认为与他同在的十字架比巨大的财富更珍贵。因为无论多么渴望，没有人能获得万物，但任何人都可以学会轻视万物，从而证明自己超越万物。既然他具有这样的心思，过着这样的生活，他就不需要祭坛和虚荣，也不需要像\Quote{克拉特斯解放忒拜人克拉特斯}这样的公开宣告。因为他的目标始终是成为最卓越者，而非看似最卓越者。他也并非住在木桶里，在市场中间，通过沉溺于公众注目而将贫穷变为富有：而是贫穷且不修边幅，却毫无炫耀，并欣然将曾拥有的一切抛诸脑后，轻快地渡过人生之海。\par

61. 节制、少欲、不受快乐的统治、不受那残忍而屈辱的主母——肚腹——的束缚，是一件奇妙的事。有谁像他那样不依赖食物，并且毫不夸张地说，那样自由于肉身？因为他把一切的饱足和过度留给了无理性的受造物，它们的生命是奴役和屈辱的。他对那些仅次于食欲而同等重要的东西很少留意，而尽可能地只靠最必需的东西生活，他唯一的奢侈就是证明自己并不奢侈，并因此没有更多的需要；但他却观看百合花和飞鸟（\BibleRef{玛 6:26}），它们的美是无心的，它们的食物是偶然的，这是按照我基督的重要教诲，祂为了我们而甘愿在肉身中成为贫穷的（\BibleRef{格后 8:9}），好使我们能享受祂天主性的富足。因此，他穿一件单衣和磨损的外套，睡在光秃秃的地上，守夜，不洗澡（这就是他的装饰），他最甜美的食物和调味品是面包和盐，他的新美味，还有清冽丰沛的饮料，那是泉水供应给无忧无虑之人的。这些事情的成果或伴随，就是照顾病人和行医，这是我们共同的智力追求。因为，虽然我在其他方面都不如他，但在这苦难上我却不得不与他相等。\par

62. 童贞、独身、与天使和那单一的本性同列，是一件伟大的事；因为我不敢称之为属于基督的，祂虽然为了我们这些生者而愿意生于童贞女，从而制定了童贞的法律，为引导我们远离此生，切断世界的力量，或者不如说，将此世传至彼世，将现今传至未来。那么，谁更尊崇童贞，或更克制肉身，不仅在他个人的榜样上，而且在他所照顾的人身上？那些隐修院和成文规章是谁的？他藉此克制了每一种感官，规范了每一个肢体，并赢得了真正的童贞实践，将美的视线转向内在，从可见到不可见；又消磨外在，从火焰中撤去燃料，向天主揭示心中的秘密——祂是纯洁灵魂的唯一新郎，将警醒的灵魂接纳到祂自己内，如果她们带着燃烧的灯和充足的油去迎接祂？（\BibleRef{玛 25:2}）此外，他极美妙地调和并联合了独修生活和团体生活。这两者在许多方面曾彼此分歧和争论，而两者都没有绝对纯粹地拥有善或恶：一种较为宁静安定，倾向于与天主结合，却难免骄傲，因为其德行无法测试或比较；另一种更具实际服务性，却难免趋向骚动。他为苦修者和隐士建立了小室，但离他的共修团体不远，而不是像用一道墙那样区分和隔离彼此，他将它们聚集并联合起来，为使默观精神不致与社会隔绝，也使积极生活不致不受默观影响；而是像海洋和陆地那样，通过彼此恩赐的交换，可以联合促进同一个目标：天主的荣耀。\par

63. 还有什么呢？博爱是一件崇高的事，扶助穷人、援助人类的软弱也是。请出城稍微走远一点，看看那座新城——虔诚的宝库、富人的公共金库。由于他的劝导，他们财富中的剩余，甚至必需品，都储存在那里，免于蛀虫的侵蚀，\newline{}\BibleRef{玛 6:19}，不再让盗贼眼红，也避开了嫉妒的竞争和时间的腐蚀。在那里，疾病被以宗教的眼光看待，灾难被视为祝福，同情心受到考验。我怎能将这项事业与有可见城门的\Place{底比斯}、埃及的\Place{底比斯}、\Place{巴比伦}的城墙、卡里亚的\Person{摩索拉斯}陵墓、金字塔、无重的铜像巨像、或已不复存在的神龛的规模和美丽，以及所有其他令人类惊叹、载入史册的事物相提并论呢？那些事物的建造者除了些许名声之外，一无所获。我的主题是最奇妙的一切：通往救恩的捷径，升入天堂的最易攀爬之路。我们眼前不再有那可怕而可怜的景象：那些四肢大半坏死、如同活尸的人，被赶出城市、家园、公共场所和泉水，甚至被赶离他们最亲爱的人，只能靠名字而非相貌来辨认；他们不再出现在我们的集会、聚会、日常交往和联合中，不再因其疾病而成为憎恨（而非同情）的对象；他们谱写凄惨的歌曲——如果其中还有人留得住声音的话。我为何要用悲剧风格来表达我们所有的经历？任何语言都无法描述他们的困苦。然而，是他率先敦促那些身为男子汉的人：不应轻视自己的同胞，也不应以不人道的行为羞辱众人唯一的首领基督；而应利用他人的不幸来稳固自己的命运，并将自己所需的怜悯借给天主。因此，他虽然出身高贵、祖先显赫、声名卓著，却不屑于以嘴唇来尊崇这种疾病，而是称他们为弟兄——并非如某些人可能臆想的那样出于虚荣（有谁比他更远离这种感觉呢？），而是率先走近照料他们，作为其哲学的结果，给予不仅有言的教导，也有无声的教诲。其效果不仅见于城中，也见于乡村和远方；甚至社会的领袖们也竞相以博爱和宽宏大量对待他们。其他人有厨师、华丽的餐桌、糕点师的巧思和美味、精致的马车和柔软飘逸的衣袍；\Person{巴西略}的关怀却是针对病人，缓解他们的创伤，效法基督，不用言语，而以行动洁净癞病。\par

64. 对于这一切，那些指责他骄傲和傲慢的人将说什么呢？他们是此等行为的严厉批评者，将那些根本不是标准的人视为标准，来衡量这位以生命为标准的人。难道亲吻癞病人，并如此谦卑自己的人，会傲慢地对待健康的人吗？难道用斋戒消耗肉体，却用空洞的傲慢使灵魂膨胀吗？难道可以谴责法利塞人，阐述傲慢的贬抑作用，却认识那自甘为奴仆形像的基督，与税吏同席，为门徒洗脚，不轻看十字架——为将我的罪钉在其上；而更令人难以置信的是，看见天主被钉十字架——是的，与强盗一起，被路人嘲笑，而祂本是不受苦难、超越痛苦的；然而照诽谤者的想象，祂自己却高踞云层之上，认为没有什么是能与祂平等的吗？相反，他们所谓的骄傲，依我看，正是他品格的坚定、稳固和不动摇。在我看来，这些人会轻易将勇敢称作鲁莽，将谨慎者称作懦夫，将节制者称作厌世者，将正义者称作吝啬。因为这条哲学格言极好：恶习紧邻德性，从某种意义上说，它们是近邻；且对于那些在此类训练上有缺陷的人来说，最容易将人误认为非其本相。因为谁比他更尊崇德性、谴责恶习？谁对正直者更仁慈、对作恶者更严厉？他的微笑常是赞许，他的沉默是斥责，在隐秘的良心中折磨恶人。若一个人不是饶舌者、戏谑者、闲谈者，也不是受欢迎的人——因为他为了取悦众人而将自己变成众人\newline{}\BibleRef{格前 9:22}——那又怎样？在有理智的人眼中，他岂不更应受赞美而非指责？除非狮子因其可怖而威严、不像猿猴、其扑跃高贵且因其非凡而受珍视，这是一个错误；而舞台演员因其讨喜和博爱的性格应赢得我们的赞赏，因为他们取悦了大众，并用响亮的耳光引人发笑。然而，若这真是我们的目标，谁在遇见他时如此令人愉快——如我所知，我有最长的经验？谁的故事更和蔼，机智更文雅，责备更温柔？他的责备不引起傲慢，他的放松不导致放荡，而是避免两者过度，按着\Person{所罗门}的规则，适时适度地运用两者，因为\Person{所罗门}给每件事物分配了季节。\newline{}\BibleRef{训 3:1}\par

65. 但这些与他辩才的声名和教导的能力（这已赢得世界各地的喜爱）相比，又算什么呢？至今我们一直在环绕山脚，却忽略了山顶；至今我们一直在穿越海峡，却未注意到那浩瀚深邃的海洋。因为我想，如果有人曾经或能够成为号角，以其响彻远方的声音，或成为天主的声音，包容宇宙，或成为世界的地震，以某种闻所未闻的奇迹，那么他的声音和才智才配得上这些称号，因为他超越并凌驾于所有人之上，正如我们超越无理性的受造物一样。谁比他更藉圣神洁净自己，使自己配得陈述神圣之事？谁更被知识之光所照亮，更深入地洞察圣神的深渊，并藉天主的助佑看见天主之事？谁的语言能更好地表达理智的真理，而不像大多数人那样一瘸一拐，要么无法表达思想，要么让口才超越推理能力？在这两方面他都获得了同样的卓越，并显示出自身的平等和绝对完美。探索一切，甚至天主的深奥之事\BibleRef{格前 2:10}，按照圣保禄的见证，是圣神的职务，不是因为祂对它们无知，而是因为祂乐于默观它们。现在巴西略已经完全探究了圣神的一切，并从中为各种性格的人抽出他的教训，他的崇高教诲，以及劝勉我们舍弃现世，适应来世。\par

66. 太阳被达味赞美，因其美丽、伟大、迅捷的行程和能力，如新郎般灿烂，如巨人般威严；而从其运行的广度，它有如此力量，将光平等地从天这边照耀到天那边，其热力在距离中毫无减弱。巴西略的美德是德性，伟大是神学，行程是不断上升直到天主面前的永恒运动，能力是圣言的播种和分配。所以我不犹豫地说出这一点：他的话语传遍大地，他的言辞的能力传到地极：正如圣保禄论及宗徒所说\BibleRef{罗 10:18}，借用达味的话。今日任何聚会中还有什么其他魅力？宴席、法庭、教堂中还有什么快乐？在上者和在下者中还有什么喜悦？隐修士或团体修士中呢？悠闲阶级或事忙者中呢？外教哲学学派或我们自己的学派中呢？有一种东西贯穿一切，是最伟大的——他的著作和劳作。作家们除了他的教导或著作，不需要其他素材。旧日所有对神圣启示的辛勤研究都已沉寂，而新的研究挂在每个人嘴上，我们中间最好的教师就是最熟悉他的作品、并讲述和解释给我们听的人。因为仅他一人就足以代替所有其他人，尤其是对那些特别渴望教导的人。\par

67. 我只想这样说：每当我拿起他的\WorkTitle{六天创世论}，口诵其词时，我就被带到造物主面前，理解了创造的话语，并且比以前更加赞美造物主，将我的老师作为我唯一的视觉媒介。每当我拿起他的论战著作，我就看见索多玛之火\BibleRef{创 19:24}，邪恶叛逆的舌头被它烧成灰烬，或者那邪恶建造的加兰塔\BibleRef{创 11:4}被正义地摧毁。每当我读他关于圣神的著作，我找到我所拥有的天主，并因他的神学和默观的支持而勇敢地宣讲真理。他的其他论文，为那些目光短浅者提供解释，用刻在他心版上的三重铭文，引导我从单纯的字面或象征解释走向更广阔的视野，当我从一个深渊走向另一个深渊，呼唤深渊之后又深渊，找到光之后又光，直到我达到最高峰。当我研读他颂扬我们斗士的赞美辞时，我蔑视身体，享受他所赞扬者的团契，并激励自己投入战斗。他的道德和实践演讲净化灵魂和身体，使我成为适合天主的殿宇，被圣神敲击的乐器，以其音调颂扬天主的荣耀和德能。事实上，他使我和谐有序，并藉神圣的转化改变了我。\par

68. 既然我提到了神学和他在这门科学中最崇高的论文，我要对已说的话再补充一点。因为这对团体极有益处，以免他们因对他不公正的低估而受损。我的言论是针对那些心怀恶意的人，他们以自己的恶行来掩盖对他人的诽谤。他在维护正统教义以及天主圣三的合一与同等神性时——用我认为尽可能精确清晰的术语来说——他不仅会欣然欢迎被逐出他的主教座（他原本无意登坐此座）为收获而非危险，甚至流亡、死亡及其前的折磨也是如此。这从他的实际行为和所受苦难中显而易见。因为当他因真理而被判流放时，他唯一的反应是吩咐一名仆人拿起书写板跟随他。他认为，按照达味圣王的劝诫，必须谨慎地引导自己的言语，暂时忍受战争时期和异端者得势的时期，直到继之以自由和宁静的时期，那时才容许言论自由。敌人伺机等候那未经修饰的话\Quote{圣神是天主}——尽管这是真的，但他们和他们不敬的恶庇护者却以为是不敬——以便将他和他教导神学的权能逐出城外，而他们自己便能夺取教会，使之成为起点和堡垒，从那里用他们邪恶的教义侵袭世界其余部分。因此，借助其他术语、意义无异的陈述以及必然导向这一结论的论证，他如此压制了对手，以致他们哑口无言，陷入自己承认的网罗——这是辩证能力和技巧的最大证明。他关于这一主题的论文进一步证明了这一点，显然是由借自圣神宝库的笔所写成。他将使用确切术语之事暂时推迟，恳求圣神本身和祂热心的护卫者不要因他的\OriginalTerm{经济}{economy}而恼怒，也不要因固守一个措辞而导致整个事业受损，出于一种不妥协的性情，在信仰面临危机的时刻。他向他们保证，他们不会因稍改措辞或用其他术语教导同一真理而受损。因为我们的救恩与其说在于言辞，不如说在于行动；因为我们不会拒绝犹太人，如果他们渴望与我们合一，却暂时寻求使用\Quote{受傅者}一词而不是\Quote{基督}：而如果教会竟被异端夺取，团体将遭受极其严重的损害。\par

69. 他，不亚于任何人，承认圣神是天主，这从他在有机会时经常公开宣讲此真理，并在私下被问及时热忱地承认，即可明了。但他在与我交谈中使之更清楚，在我们就此主题的会谈中，他对我毫无隐瞒。不满足于单纯肯定这一点，他还进一步（此前他极少这样做）诅咒自己承受与圣神分离的最可怕命运，如果他不朝拜圣神为与圣父圣子\OriginalTerm{同质同等}{consubstantial and coequal}。如果有人愿意接受我为他在此事业中的同工，我将陈述一件至今大多数人不知道的事。在当时的困难压力下，他亲自承担了\OriginalTerm{经济}{economy}，同时让我自由发言——我不太可能被从隐没中拖出受审或流放——好让我们共同努力使我们的福音得以稳固建立。我提到这一点，不是要为他辩护——因为他比攻击者更强大，如果有任何攻击者的话——而是要防止人们以为他著作中使用的术语就是真理的极限，从而削弱信德，并认为自己的错误受到他神学的支持（他的神学是时代影响和圣神影响的共同结果），而不去考虑他著作的意义和写作目的，以便更接近真理，并能封住不敬党派之人的口。无论如何，愿他的神学成为我的神学，以及所有亲近我之人的神学！我对他在此方面的清白如此自信，以致我将他视为我在这事上以及一切事上的伙伴：愿我的归于他，他的归于我，无论是在天主面前，还是在最有智慧的人面前！因为我们不会说福音作者彼此矛盾，因为有些人更多关注基督的人性，另一些人关注祂的神性；有些人从我们经验之内开始叙述，另一些人从我们之上开始；并且通过这样分享他们信息的实质，他们使接受者得益，并追随内在圣神的感动。\par

70. 那么来吧，从古至今有许多因虔敬而闻名的人，如立法者、将领、先知、教师，以及流血牺牲的勇士。让我们将我们的主教与他们相比，从而认识他的功绩。亚当曾蒙天主的亲手所造\BibleRef{创 1:27}，享受乐园的欢愉，并获赐最初的律法；但是，除非我冤枉了我们原祖的名声，他没有遵守诫命。而巴西略却接受并遵守了诫命，未从知善恶树受到伤害，逃脱了那旋转的火剑，并且我深信，他已抵达乐园。厄诺士首先呼求主名。巴西略不但自己呼求主，而且更伟大的是，他向别人宣讲主。哈诺客被提升天，因着一点虔诚（因为那时信仰仍在影中）获得提升\OriginalTerm{提升}{translation}，逃离了余生之险，但巴西略的整个一生就是一种提升，他在完满的一生中完全经受了考验。诺厄受托管理方舟\BibleRef{创 6:13}，以及一个新世界的种子被托付给一个木制的小屋，以保存它们免受洪水之灾。巴西略逃脱了不敬之心的洪水，并使自己的城市成为一座安全的方舟，轻快地在异端者之上航行，后来恢复了整个世界。\par

71. 亚巴郎是位伟人，一位圣祖，他献上了新的祭献，将所应许的子孙作为甘愿献上的祭品，急切地献给了赐予者。但巴西略的奉献并不微小，他把自己献给了天主，没有任何替代之物（因为那怎么可能呢？）因此他的祭献得以圆满。依撒格在出生前就已被应许，巴西略则将自己应许给天主，并娶了黎贝加为妻，我说的是教会，不是由仆人从远处领来，而是由天主在近处赐予并托付给他；他也没有在子女的偏爱上受骗，而是按照圣神的判断，毫无欺骗地将各人应得的赐予他们。我赞扬雅各伯的梯子，以及他给天主傅油的石柱，还有他与天主的角力，无论那是什么；在我看来，那是人的高度与天主的至高之间的对比和对抗，导致了他种族失败的征兆。我也赞扬他在畜牧上的巧妙策略和成功，以及他的子女——十二位圣祖，还有他祝福的分赠，以及对未来的光荣预言。但我更赞扬巴西略，因为他所见的不仅是梯子，而是通过不断迈向卓越的阶梯登上它；他傅油的不是石柱，而是为天主树立了石柱，通过示众不敬者的教导而树立；他还进行了角力，不是与天主，而是为了天主，以推翻异端者；他的牧灵关怀使他富有，因为他为自己赢得了许多有标记的羊，超过没有标记的羊；他在属灵子女上卓有成效，他的祝福建立了许多人。\par

72. 若瑟是粮食的供应者\BibleRef{创 41:40}，但仅限于埃及，且次数不多，供应的是身体的食物。巴西略则造福众人，且无时无刻，供应的是属灵的食物，因此，在我看来，他的职责更为尊贵。像乌兹人约伯\BibleRef{约 1:1}一样，他受到试探，并战胜了试探，在奋斗结束时获得了光辉的荣耀，未被众多攻击者中的任何一个动摇，对试探者的努力取得了决定性的胜利，并堵住了那些不明白他痛苦奥秘之朋友们的无理之言。\Quote{在祂的司祭中，有梅瑟和亚郎。}梅瑟确是伟大的，他降下灾祸于埃及，在众多征兆和奇事中带领百姓，进入云中，并认可了双重法律：外在的属字母，内在的属圣神。亚郎是梅瑟的兄弟，无论是自然还是属灵上，他为百姓献祭和祈祷，作为那伟大而神圣会幕的大司祭\OriginalTerm{大司祭}{hierophant}，这幕是主所支搭的，不是人支搭的\BibleRef{希 8:2}。巴西略与他们二人都可匹敌，因为他折磨异端者的埃及种族时，用的不是身体上的灾祸，而是精神上和思想上的灾祸；他带领归主的人民，即那些热心行善的选民\BibleRef{铎 2:14}，进入应许之地；他颁布了法律，不再是隐晦的，而是全然属灵的，写在\BibleRef{格后 3:3}不会破碎而是保留下来的石板上；他进入了至圣所，不是一年一次，而是常常，我可以说每天，并从那里向我们启示了天主圣三；他洁净百姓，不是用临时的洒血，而是用永恒的净化。若苏厄\BibleRef{苏 1:2}的特殊长处是什么？是他的统帅才能，土地的分配，以及对圣地的占领。难道巴西略不是一位\OriginalTerm{统领}{Exarch}吗？难道他不是那些因信德得救\BibleRef{弗 2:8}之人的将领吗？难道他没有按照天主的旨意，为他的跟随者分配不同的产业和住所吗？因此，他也能用这句话：\Quote{我的产业落在佳美之处；} 以及：\Quote{我的产业在祢手中，} 这产业比我们在地上的产业更宝贵，且不能夺去。\par

73. 再者，要历数民长，或最卓越的民长，有\Quote{撒慕尔在那些呼求祂名的人中}，他在出生前就被奉献给天主\BibleRef{撒上 1:20}，出生后立即被祝圣，并以角膏抹君王和司祭。但巴西略不也是从母胎中就被奉献给天主，带着外衣\BibleRef{撒上 2:19}被献在祭台上吗？他不也是天上之事的先知，主的受傅者，以及那以圣神成全者的傅油者吗？在君王中，达味备受颂扬，他的胜利和来自敌人的战利品\BibleRef{撒下 5:1}都有记载，但他最显著的特征是他的温和，以及在他为王之前，他弹琴的能力，甚至能安抚恶灵。撒罗满向天主求得了宽阔的心胸\BibleRef{列上 4:29}，在智慧和默观上取得了最大的进步，成为他那个时代最著名的人。在我看来，巴西略在温和方面毫不逊色，或在智慧方面也毫不逊色，或者只是稍逊一筹，以致他安抚了暴怒君王的傲慢；他不仅使南方女王从地极前来拜访他，因他智慧的名声，而且使他的智慧传遍了地极。我略去撒罗满其余的生涯。即使我们放过它，也是对所有人显而易见的。\par

74. 你赞扬厄里亚在暴君面前的勇气\BibleRef{列下 1:1}和他乘火升天吗？或是厄里叟的美好遗产，即羊皮披肩，以及厄里亚的灵？你也必须赞扬巴西略的生活，他是在火中度过的，我是指在众多试炼中，以及他穿过火而逃脱，那火焚烧却没有烧毁\Quote{荆棘}\BibleRef{出 3:1}的奥秘，以及那来自高天的美丽皮衣，即他对肉身的无动于衷。我略过其余的事：三个青年在火中被滋润\BibleRef{达 3:5}，那位逃避的先知在鲸鱼腹中祈祷\BibleRef{纳 2:1}，并从这兽中出来，如同从内室出来；那位义人在狮坑中抑制了狮子的狂怒\BibleRef{达 6:22}，以及七位玛加伯\BibleRef{加下 7:1}的奋斗，他们与父母一起在血中和各种酷刑中得以成全。巴西略与他们一样坚忍，并赢得了他们的荣耀。\par

75. 现在转向新约，将他的一生与那些在这里卓越的人相比，我将在导师们身上找到他们的弟子的荣耀之源。谁是耶稣的前驱？\BibleRef{路 1:76}若翰，圣言的呼声，真光的明灯，在他之前甚至在母胎中就已欢跃\BibleRef{路 2:41}，并且他先于耶稣进入阴间，因黑落德的怒气而被派遣\BibleRef{玛 14:10}，在那里也宣告那将要来的那一位。如果我的言语在任何人看来是鲁莽的，让我事先向他保证，在做这个比较时，我既不偏爱巴西略，也不认为他与那超越一切妇人所生者相等\BibleRef{玛 11:11}，而只是表明他被激励去效法，并在某种程度上拥有了他那些引人注目的特征。因为对于热忱者来说，即使是在微小程度上效法最伟大的人，也不是小事。难道不明显吗？巴西略是若翰苦修生活的翻版。他也住在旷野，在夜间守望中穿着褴褛的衣服，在他退缩隐退时；他也喜爱相似的食物，通过禁欲为天主净化自己；他也被认为配作基督的宣告者，如果不是前驱的话，不仅全周围地区的人出来到他那里，连边界之外的人也来；他也站在两个盟约之间，通过施行精神废除一个盟约的字面意义，并藉着解除那显明的法律而实现那隐藏的法律。\par

76. 他效法伯多禄的热忱\BibleRef{宗 4:8}，保禄的殷切，这两位有名有姓之人的信德，载伯德儿子的崇高言论，以及所有门徒的节俭和纯朴。因此，他也受托掌管天国的钥匙\BibleRef{玛 16:1}，不仅从耶路撒冷直到依里黎苛\BibleRef{罗 15:1}，而且他在福音中拥抱更广的范围；他不被称为，而成为\Quote{雷霆之子}；并且躺在耶稣的胸膛上，从中汲取他言语的能力和他的思想深度。他本可能成为斯德望\BibleRef{宗 7:58}，他固然渴望，但敬畏之心中止了那些想要用石头砸死他之人的手。我能更简洁地总结，以避免就这些点详细论及每个人。在某些方面，他发现了善，在某些方面他效法了善，在其他方面他超越了善。他多方面的德行使他超越了当代所有的人。我只剩下一点要说，而且用几句话。\par

77. 他的德行如此伟大，名声如此卓越，以致他的许多次要特征，甚至他的身体缺陷，也被别人为了出名而效仿。例如他的苍白、胡须、步态、沉思以及通常冥想时的说话犹豫，这在许多人轻率而不加考虑的模仿中，变成了忧郁的样子。此外，他衣服的样式、床的形状和进食的方式，这些对他都不是重要的事，而只是偶然和巧合的结果。因此，你可能会看到许多在外貌上像巴西略的人，在这些轮廓的雕像中，因为称他们为他的遥远回声也太过分了。因为回声虽然是声音的逐渐消失，但无论如何也能非常清晰地再现它，而这些人则远远不及他，甚至不能满足他们接近他的愿望。能够偶然遇见他或为他做些什么，或者带走他开玩笑或认真时所说或所做的一些纪念品，这也不是小事，而是一件有充分理由被高度重视的事：我知道我自己也常常为此感到自豪；因为他的即兴发挥比别人的劳苦之作更珍贵、更出色。\par

78. 然而，当他跑完赛程，保持了信德\BibleRef{弟后 4:7}，渴慕离世，以及接受冠冕的时刻临近时\BibleRef{斐 1:23}，他并没有听见召唤：\BibleQuote{你上这山来，死在那里}\BibleRef{申 32:49}，而是\Quote{死，然后上到我们这里来。}在这里，他又行了一件不亚于前述的奇事。因为当他几乎死去、气息奄奄、丧失了大半能力时，他却在最后的言辞中变得更强，以致以虔诚的话语离世，并通过祝圣他最杰出的随从，将他的手和圣神都赐予了他们：这样，那些曾协助他司铎职务的门徒，就不会被剥夺司铎职。我余下的任务，我虽不情愿，却要着手，因为这些话由别人说出会比我自己说出更充分。因为我无法以哲学态度面对我的不幸，即使我多么渴望如此，当我回想起这损失是大家共同的，而且不幸降临到整个世界时。\par

79. 他躺着，呼出最后一口气，天上的歌咏团在等待着他，他早已将目光投向那里。全城的人涌向他周围，无法承受他的离去，谴责他的离世，仿佛这是一种压迫，紧紧抓住他的灵魂，好像它能被他们的手或祈祷所约束或强迫。他们的痛苦使他们疯狂，所有人都渴望，如果可能，将自己的生命分一部分给他。当他们失败时——因为他必须被证明是人——他用最后的话说：\BibleQuote{我把我的灵魂交在你手中}，然后欢欣地将灵魂交托给带他离开的天使照料；并非没有给在场的人留下一些虔诚的指示和训诫——随后发生了一件比先前任何事都更令人惊奇的奇迹。\par

80. 圣人被抬出，由圣人们的手高高举起，每个人都急切地：有人想抓住他衣服的边\BibleRef{路 8:44}，有人只想触摸他的影子\BibleRef{宗 5:15}，或担架——上面载着他神圣的遗体（因为有什么比那身体更神圣、更纯洁呢？），有人想靠近抬担架的人，有人只想享受这景象，仿佛连这也有益。市场、门廊、两三层的房子里挤满了护送、前导、跟随、陪伴他的人，人们互相践踏；成千上万的各种种族和年龄的人，前所未有。圣歌被哀伤压倒，哲学的顺从被不幸压垮。我们的人与陌生人、犹太人、希腊人、外邦人竞争，他们也与我们竞争，通过更多的哀悼来获得更大的益处。结束我的故事：这场灾难以危险告终；许多灵魂因推挤和混乱而随他一同离去，这些人在离世中被视为有福的，与他一同步入死亡，也许某个热情的演说家会称他们为\Quote{殡葬的牺牲}。遗体最终摆脱了那些想抢夺它的人，穿过前面的人群，被安放在他祖先的坟墓中，大司祭被添加到司铎中，在我耳中回响的洪亮声音被添加到传报者中，殉道者被添加到殉道者中。如今他在天堂，如果我没有弄错，他在为我们献祭，为人民祈祷，因为他虽然离开了我们，却没有完全离开我们。而我，\Person{额我略}，半死不活，被劈成两半，从我们伟大的结合中被撕开，拖着一份痛苦的生活——想必在与他分离后，这生活并不轻松——现在失去了指引，不知道我的结局将如何。即使现在，我仍在夜间的异象中得到劝告和教导，如果我在职责上有任何亏损。我现在的目的与其说是将哀叹与赞美混合，或是描绘这人的公共生活，或是发布一幅万世共通的德行图画，和对所有教会、所有灵魂都有益的榜样，我们可以将其视为一部活的法律，从而正确引导我们的生活，不如说是劝告你们——已经完全领受了他的教义——将目光定睛在他身上，如同一个看着你们也被你们所看的人，从而被圣神所完美。\par

81. 那么，到这里来，围绕着我，你们所有他团体中的成员，无论是圣职人员还是平信徒，无论是本国人还是外国人；请帮助我完成我的颂词，每人都提供或要求讲述他的某些优点。你们，法官席上的，请看立法者；你们，政治家，请看国务活动家；你们，平民，请看他的秩序；你们，文人，请看导师；你们，贞女，请看新娘的引导者；你们，已婚者，请看节制者；你们，隐修士，请看赐予你们翅膀者；你们，同修会者，请看审判者；你们，淳朴者，请看向导；你们，默观者，请看神视者；你们，快乐者，请看缰绳；你们，不幸者，请看安慰者，白发者的手杖，青年的引导者，贫困者的救助，丰裕的管理者。我想，寡妇也会赞扬她们的护佑者，孤儿赞扬他们的父亲，穷人赞扬他们的朋友，异乡人赞扬他们的款待者，弟兄赞扬那有兄弟之爱的人，病人赞扬他们的医生——无论他们患什么病，需要什么医治——健康者赞扬健康的守护者，所有人赞扬那成为一切人的一切，以赢得大多数——即使不是全部——的人。\par

82. 这是我奉献给你的，巴西略，由那曾是你的同伴、与你年龄和地位相仿者口中说出，这声音曾是你以为最悦耳的。若它接近了你的功绩，那应感谢你，因为是由于对你的信赖，我才敢谈论你。但若它远未达到你的期望，我们这些因年迈、疾病和对你思念而憔悴的人，又该作何感受？然而，当天主使我们能做我们所能做的事时，祂是喜悦的。但愿你能从高处垂视我们，你这神圣而可敬的圣者；或借着你的恳求，为我们平息那由天主为管教我们而赐予的肉体上的刺（\\BibleRef{格后 12:7}），或使我们能勇敢地承受它，并引导我们整个生命走向最有益于我们的事。若我们被接升，也请你在自己的帐幕里接纳我们，好使我们同住一起，更清晰、更圆满地瞻仰那圣而荣福的天主圣三——我们现在已在一定程度上领受了它的肖像——我们的渴望终得满足，以此作为我们所有战斗和所忍受的攻击的酬报。这些是我们为你所说的话；在我们离世之后，还有谁会赞美我们呢？即使我们能在我们主基督耶稣内提出任何堪受言语或赞美的主题，愿光荣归于祂，直到永远。阿们。\par

\SourceRecord{Translated by Charles Gordon Browne and James Edward Swallow.；Nicene and Post-Nicene Fathers, Second Series；Vol. 7.；Edited by Philip Schaff and Henry Wace.；Buffalo, NY: Christian Literature Publishing Co.,；1894.；<http://www.newadvent.org/fathers/310243.htm>.}

% structure: p0001 source=h1 target=section reason=page-title

\section{讲道集 45}

\Emph{第二篇复活节讲道。}\par

这篇讲道并非如其标题可能使我们以为的那样，是在第一篇之后立即发表的；事实上，两者之间相隔了许多年，彼此并无关联。按年代顺序，它们是圣额我略讲道中的第一篇和最后一篇。第二篇是在\Place{Arianzus}教堂内宣讲的，该村位于\Place{Nazianzus}附近，额我略在此继承了一些产业。他辞去\Place{君士坦丁堡}总主教之职后便退隐于此，随后，由于年事渐高、体力日衰，连小小的Nazianzus主教区的管理工作也感到力不从心，于是在公元383年底前后退居\Place{Arianzus}，并于389年或390年在那里去世。\Quote{这篇讲道的开场白完全是圣经的风格；讲道者在此描述并假借复活天使之口说话。他的目的是要显示这节日庆典的重要性，并以寓意的方式解释古代逾越节的一切细节，将其应用于基督和基督徒的生活。有两段文字是逐字借用自曾在君士坦丁堡宣讲的圣诞讲道的}（\Person{Benoît}）。\par

本笃会编辑者表示，他们无法确定这种重复是来自圣额我略本人，还是某些抄写员的疏忽。\par

一、\Quote{我要站在我的守望台上，}可敬的哈巴谷说，\BibleRef{哈 2:1}；今天我也要按着圣神赐给我的权柄和观察，站在他旁边；我要观看，要看有什么事对我说。\Quote{于是，我站立观看；看，有一个人骑着云彩，他极其崇高，他的容貌好似天使的容貌，} \BibleRef{民 13:6} \Quote{他的衣服犹如闪电的光辉；他向东伸出他的手，大声呼喊。他的声音好像号筒的声音；在他周围，犹如有一大队天军；他说：\InnerQuote{今天救恩临到了世界，临到那可见的与那不可见的。基督从死人中复活了，你们要与他一同复活。基督又回到他自己，你们也要回转。基督从坟墓中释放了，你们也要从罪恶的捆绑中释放。地狱的门已敞开，死亡已被消灭，旧亚当被废弃，新亚当得以成全；谁若在基督内，他就是一个新受造物。}} \BibleRef{格后 5:17} \InnerQuote{你们要更新。}他这样说着；其余的便歌唱起来，正如从前基督藉着祂的诞生在地上向我们显现时所唱的一样：\Quote{光荣归于至高之处的天主，平安归于地上，善意归于人。}我也在你们中间与他们一同说出同样的话。但愿我能得到一个与天使相称的声音，响彻大地的四极。\par

二、主的逾越节，逾越节，我再次说逾越节，为光荣天主圣三。这对我们来说，是庆节中的庆节，是盛典中的盛典，远超一切其他的（不仅那些纯粹属人的、匍匐于地面的，甚至那些属于基督本身、为尊崇祂而举行的），犹如太阳远超众星。昨天我们的光彩行列和灯火辉煌确是美丽的，无论是在公开场合还是私下里，我们各色人等、几乎各个阶层都参与其中，以我们密集的火光照亮黑夜，仿效那大光的形式，包括我们头顶上天体点燃的灯塔之光，以及那在天之上的、在天使中间之光（那第一性的光明本性，仅次于万有的第一性，因为它是直接从那第一性涌流出来的），以及那在圣三内的光，一切光皆由之而有，从这不可分的光中分有并受到尊崇。然而今天更为美丽、更为辉煌；因为昨天的光是那大光升起的前驱，并像是为节日作准备的一种喜庆；而今天，我们正在庆祝复活本身，不再是作为期待的对象，而是作为已经发生的事实，并召集整个世界归向它。让不同的人在这一时节带来不同的果实和不同的奉献，或大或小……即那些属神的、天主所喜悦的祭献，按各人的能力。因为连天使本身——那些首要的、理智的、纯洁的存有，也是天上光荣的目击者——恐怕也难以献上与它的尊位相称的礼物，即使他们能达到完全的赞美。至于我们，我们要献上一篇讲辞，这是我们所有物中最好、最宝贵的——尤其是当我们为圣言赐予理性受造物的恩惠而赞美祂时。我要从这一点开始。因为我不能忍受，当我为这伟大的祭献和这最伟大的日子献上嘴唇的祭献时，竟不回溯到天主，不以祂为我的开端。因此我求你们，一切喜爱此类主题的人，请洁净你们的心思、耳朵和思想，因为这篇讲辞关乎天主，并将是神圣的；这样你们离开时，必充满那不会消逝于虚无的喜乐。它将是既充实又简洁的，既不因欠缺而令你们苦恼，也不因过度而令你们不悦。\par

III. 天主始终存在，现在存在，将来也永远存在；更确切地说，天主永远是\Quote{存在}者。因为\Quote{过去是}和\Quote{将来是}是我们时间与可变化本性的片断。但祂是永恒之\Quote{有}；当祂在山上向\Person{梅瑟}颁布神谕时，这便是祂自称之名。因为在祂内，祂总结并包含一切存在，既无过去的起始，也无未来的终结……如同一个存在的汪洋大海，无垠无际，超越一切时间与本性之概念，仅被理性模糊地、依稀地影射……不是凭其本质，而是凭其环境；从一个来源获得一个形象，从另一来源获得另一个形象，组合成某种真理的呈现，而这真理在我们尚未抓住时便已逃逸，在我们尚未想及时便已飞逝，即使在我们净化的理智上闪耀，如同闪电不停止其行程照耀我们的眼睛……依我之见，这是为了通过我们所能理解的部分吸引我们归向祂（因为完全不可理解之物超出了希望的范围，也不在努力所及之内）；而通过我们无法理解的部分引发我们的惊叹；作为惊叹的对象，成为更令人渴望的对象；被渴望而净化；净化而使我们肖似天主；如此，当我们变得相似于祂时，天主——允许我大胆地说——便会如同天主与我们交往，与我们合一，为我们所知；并且这或许达到祂早已认识那些认识祂的人同样的程度。故此，天主本性是无限而难以理解的，我们对祂所能理解的一切就是祂的无限性；即使有人可能认为，因为祂是单纯的本性，所以祂要么完全不可理解，要么完全可理解。让我们进一步探讨\Quote{具有单纯的本性}意味着什么。因为十分确定的是，这种单纯本身并非其本性，正如复合本身并非复合体的本质一样。\par

IV. 当从两个角度——起始与终结——来考虑无限性时（因为超越这两者而不受其限制的便是无限性），当理性仰望高处的深渊，没有立足之处，而倚靠\Emph{现象}以形成关于天主的观念时，它将在那里发现的无限与不可接近之物称为\OriginalTerm{无始}{Unoriginate}。当它俯视低处的深渊与未来时，它称祂为不死与不朽。当它从整体得出推论时，它称祂为永恒。因为永恒既非时间，亦非时间的一部分，因它不可度量。但对我们而言，由太阳行程所度量的时间，对永恒者而言，永恒便是一种类似时间内运动与间隔，与其存在等长。不过，这些就是我目前必须关于天主所说的一切；因为当前不是合适的时机，我现在的主题不是论天主的道理，而是论降生成人的道理。当我说\Quote{天主}时，我意指圣父、圣子及圣神；因为天主性既未超越这些位格而扩散，以致引入一群神明，也未局限于比这些更小的范围，以致判我们为贫乏的天主观，要么犹太化以保全唯一统治，要么因我们众多神明而堕入异教。因为两边的邪恶相同，尽管方向相反。至圣之所即是如此，它甚至对\OriginalTerm{色辣芬}{Seraphim}也是隐藏的，并因三次重复的\Quote{圣}而在归于\Quote{主与天主}的称号中受到光荣，正如我们的一位先贤所最美妙、最高远地推论出的那样。\par

V. 然而，既然仅凭自我默观这一运动不能满足善，而是善必须倾注自己并超越自身而出，以增广其恩惠的对象（因为这对于最高的善是必要的），祂首先构想出天使与天上的权能。这一构思是由祂的圣言完成并由祂的圣神所成全的。于是次级光辉便出现了，作为首要光辉的仆役（不论我们视它们为理智之神，或是非物质、无形体的火，或是某种尽可能接近此形态的本性）。我愿说它们不能朝向恶运动，而只可接受善的运动，因为它们在临近天主并由天主的第一光芒照耀（因为地上的受造只有次等光照），但我不得不止于说它们是不可动摇的，而只构想并宣称它们是难以动摇的，因为那位因他的光辉被称为\OriginalTerm{路济弗尔}{Lucifer}者，却因骄傲而变为并被称为\Quote{黑暗}；以及隶属于他的叛变之军，他们因反叛善而成为邪恶的创造者，与我们的煽动者。\par

VI. 如此，鉴于这些理由，祂赋予思想世界以存在——就我所能推理这些问题，并以其鄙陋语言估算伟大事物而言。然后，当祂的第一个创造井然有序时，祂构想出第二个世界——物质且可见的世界；这是一个由地与天以及其间的万物所构成的体系。当我们观看每个部分的优美形态时，这确实是一个令人赞叹的创造；但当我们考虑整体的和谐与一致，以及每个部分如何以美好秩序与其他部分相配合，且所有部分与整体相配合，达成世界作为一个统一体的完美完成时，则更令人赞叹。这是要表明，祂不仅能召唤出一种与祂相似的本性，也能召唤出一种完全与祂相异的本性。因为与天主相似的，是那些理智的、仅由心灵可理解的本性；而所有感官可觉察的，都完全与祂相异；其中与祂最远的，是那些完全缺乏灵魂与运动能力的。\par

七、理智与感官，如此彼此区分，本各自安于其界，自身蕴藏着\OriginalTerm{造物主圣言}{Creator-Word}的宏伟，是祂大能的沉默颂扬者和激动人心的宣告者。两者尚未有任何混合，也未有这些相反者的交织，这是创造万有时更大智慧与慷慨的征兆；善的整个财富也尚未被显明。如今，造物主圣言决意显示这一点，并从两者（我指不可见的与可见的受造界）中造出一个单一的生命体——祂塑造了人；从已存在的物质中取了一个身体，并将从祂自己取来的一口\OriginalTerm{气息}{Breath}（圣言知道这是一有灵的灵魂，是天主的肖像）放在其中，如同一个第二世界，伟大的渺小，祂将人安置在地上，一个新天使，一位混成的敬拜者，完全被引入可见的受造界，却只部分被引入智性的受造界；地上万物的君王，却服从于天上的君王；属地的又属天的；暂时的又不朽的；可见的又智性的；处于伟大与卑微之间；在一个位格中结合了神魂与肉身；神魂是由于施予他的恩宠，肉身是由于他被提升的高度；前者使他能继续存活并荣耀其恩主，后者使他能受苦，并藉受苦被提醒，若他在伟大中变得骄傲则得纠正；一个生命体，在此受操练而后迁往他处；为完成这奥迹，因他的倾向天主而被神化……因为，我想，我们此处所获的、仅限一点的真理之光，其趋向正是如此：使我们既看见又经验天主的荣光，这荣光配得上那造我们、将解散我们又将以更高尚样式重造我们的祂。\par

八、祂将这个存在物安置在乐园里——无论那是怎样的乐园（祂以自由意志的恩赐尊荣了他，为使善能属于他，作为他选择的结果，同样也属于那播下善种的那一位）——为耕种那些不死的植物，这或许是指神圣的观念，包括较简朴的和较完美的；赤身露体，过着单纯且无技巧的生活；毫无遮盖或遮蔽；因为从起初就是如此是合宜的。祂又赐给他一条法律，作为他的自由意志所依凭的素材。这条法律是一条诫命，关于他可享用哪些植物，而哪一棵他不可触摸。后者就是\OriginalTerm{知善恶树}{Tree of Knowledge}；然而，并非因为它从被栽种之初就是恶的；也不是因为天主嫉妒人而禁止它——愿天主的仇敌不要在这方面鼓舌，或效仿那蛇。但若在适当的时候享用，它本是好的；因为按我的理论，这树是默观，只有那些达到习惯成熟的人才能安全地进入；但对于那些仍有些单纯和贪婪的人却不好；正如固体食物对于那些尚幼嫩、需要奶水的人也不好。但是，当因着魔鬼的恶意和女人的任性（\BibleRef{智 2:24}），她作为较柔弱者屈服于它，并把它带到男人身上，因她更善于说服——哀哉我的软弱，因为我初祖的软弱也成了我的；他忘记了所赐给他的诫命，屈从于那有害的果子；因他的罪，他立刻被逐出生命树、乐园和天主；并穿上了\OriginalTerm{皮衣}{coats of skins}，那或许是指更粗糙的肉身，既必死又矛盾的。这是他首先学到的事——自己的羞耻——他便躲避天主。然而，在此他也获得了一个益处，即死亡和罪的断绝，为使恶不致成为不朽的。如此，他的惩罚变成了仁慈，因为我深信，天主施行惩罚原是出于仁慈。\par

IX. 祂先以多种方式受到惩罚，因为祂的罪恶众多，这罪恶的根源通过不同的原因和不同的时代萌发——藉言语、法律、先知、恩惠、威胁、灾祸、洪水、烈火、战争、胜利、失败、天上的征兆、空中的征兆、地上的征兆和海洋的征兆；藉人、城邦、民族出乎意料的变迁（其目的在于摧毁邪恶）——最后，祂需要一种更强的治疗，因为祂的疾病日益恶化：相互残杀、通奸、伪证、违反本性的罪行，以及一切罪恶中首要和最终的邪恶——偶像崇拜，将敬拜从造物主转向受造物。既然这些需要更大的援助，它们也获得了更大的援助。那就是天主的圣言，祂在万世之前，不可见、不可理解、无形体，是万始之始，光明之光，生命与不朽之源，原型的肖像，不可移动的印鉴，不变的形象，圣父的定义和圣言——来到祂自己的肖像前，为了我们的肉体而取了肉体，并为了我的灵魂而将自己与一个理性的灵魂结合，以相似净化相似；除了罪以外，在各方面都成为人；由童贞女受孕，她首先在身体和灵魂上被圣神净化，因为必须既尊敬生育，又使童贞获得更高的荣誉。然后祂出来，作为天主，带着祂所取的那一位；一位在两性中，肉体和神性，后者使前者神化。啊，新的混合！啊，奇异的结合！自有者成为存在，非受造者被造，那不可包容者藉着介于神性和肉体的肉体性之间的理性灵魂的中介而被包容。那赐予财富者成为贫穷，因为祂取了我肉体的贫穷，好使我获得祂神性的财富。那充满者倒空自己，因为祂暂时倒空自己的荣耀，好使我分有祂的圆满。祂的良善的财富是什么？这环绕我的奥秘是什么？我曾分有肖像，却没有守住它；祂分有我的肉体，为要既拯救肖像，又使肉体不朽。祂传递了第二次共融，比第一次更奇妙，因为那时祂传递了更好的本性，但现在祂亲自取了较差的。这比先前的行动更似天主；这在所有有智慧的人眼中更为崇高。\par

X. 但或许有些过于急躁和欢乐的人会说：\Quote{这一切与我们何干？快马加鞭奔向目标；向我们谈论这庆节和我们今日聚集在此的原因。}是的，这正是我打算要做的，尽管我已从稍早之处开始，且因我论证的需要而被迫如此。在学者和爱美者眼中，若我们就逾越节一词本身说几句话，并无妨害，因为这样的补充在他们的耳中并非无用。这个伟大而可敬的逾越节，在希伯来语中被称为\OriginalTerm{逾越节}{Pascha}，其意为\Quote{越过}。从历史上看，源自他们从埃及逃往客纳罕地的迁移；从属灵上看，源自从低处向高处和应许之地的进步和上升。我们注意到，圣经中常见的一件事，即某些名词从不明确的意义变为更清晰的意义，或从粗俗变为更文雅，在此也发生了。因为有些人以为这是神圣苦难的名称，于是将Phi和Kappa改为Pi和Chi，用希腊语称这日子为\OriginalTerm{逾越节}{Pascha}。习俗采纳并肯定了这词，借助大多数人的耳朵，对他们而言，它听起来更为虔诚。\par

XI. 但在我们之前，圣宗徒已宣布法律不过是未来之事的阴影，\BibleRef{希 10:1}这些事由思想所领悟。而天主，在更古老的时代向梅瑟颁布神谕时，关于这些事立法说：\BibleQuote{你应谨慎，在山上按指示你的样式制造一切，}\BibleRef{出 25:40}当时祂向梅瑟显示可见之物，作为不可见之物的预表和图案。我深信这些事没有一件是徒然设立的，没有一件是没有理由的，没有一件是卑下的或不配天主的立法和梅瑟的职分，尽管在每一预像中都难以找到一种理论，能够细述最微妙的细节，涉及帐幕本身及其尺寸和材料，以及搬运这些的肋未人和司铎，以及所有关于祭献、洁净和供奉的律例；尽管这些只有那些在德行上与梅瑟等同，或最接近其学识的人才能理解。因为在这同一座山上，天主被人看见：一方面通过祂自崇高居所的降临，另一方面通过祂将我们从地上的卑贱中提升，好使不可理解者在一定程度上，且按安全范围，被必死的本性所理解。因为除此之外，物质身体的迟钝和被禁锢的理智不可能意识到天主，除非藉祂的助佑。因此，那时并非所有的人都似乎被视作配得同等的等级和地位；而是各在不同的位置，我想，各按自己净化的程度。有些人甚至被完全驱逐，只被允许听到来自高天的声音——即那些性情完全如野兽、不配神圣奥秘的人。\par

十二. 但我们站在那些完全迟钝的人与那些非常沉思和崇高的人之间，既不可完全闲散不动，也不可过于忙碌以致达不到目的而偏离目标——因为前者是犹太式的且非常低下，后者只适合梦占者，两者都该受谴责——让我们在这些问题上发言，只要在我们所能及的范围内，且不是非常荒谬或受大众嘲笑。我们的信仰是：由于我们这些因原罪而堕落、被快乐引诱甚至到偶像崇拜和非法流血的人，应当因天主父的仁慈而被召回并重新提升到原来的地位，祂不能容忍祂亲手所造的如此高贵作品——人——失去，我们再造的方式和所应做的事是这样：所有激烈的方法都不被认可，因为不太可能说服我们，而且很可能因我们根深蒂固的骄傲而加重灾祸；但天主以温和友善的治疗方法安排我们的恢复。因为弯曲的树苗不能承受突然向反方向弯曲或矫正之手施加的暴力，它更容易折断而非变直；而一匹烈性且超过一定年龄的马不会在没有任何抚慰和鼓励的情况下忍受马嚼子的控制。因此，法律被赐予我们作为帮助，好比天主与偶像之间的边界之墙，使我们远离一方而转向另一方。它起初先让一小步，以便获得更大的。它暂时容忍祭献，以便在我们心中建立天主，然后时机成熟时也废除祭献；如此通过逐步移除，智慧地改变我们的心意，并在我们已经受训能有即时服从后，将我们引向福音。\par

十三. 为此原因，成文法律介入，带领我们归向基督；这是我对祭献的解释，也是我如此解释的理由。为使你们不致不知道祂智慧的深奥和祂不可探测的判断的丰富（\BibleRef{罗 11:33}），祂甚至没有完全使这些祭献成为不洁或无用的，或只包含血的空壳。但那伟大的、如果可以如此说，就其最初本性而言不可被献为祭的\OriginalTerm{祭品}{Victim}，与法律的祭献混合，成为净化，不是为一部分世界，也不是为短时间，而是为全世界和所有时间。为此，一只羔羊因其纯洁而被选，它遮盖了原始的赤身。因为这样就是为我们奉献的祭品，祂在名字和事实上都是不朽的衣服。祂是一个完美的祭品，不仅因祂的\OriginalTerm{天主性}{Godhead}（没有什么比这更完美），而且因祂所取的人性已受\OriginalTerm{神性}{Deity}傅油，并与傅油者成为一体，我敢说，与天主平等。公的，因为是为\Person{亚当}而奉献；或者更确切地说，是为强壮者中的更强壮者，当第一个人陷入罪中时；主要因为在祂内没有女性的、不刚毅的成分；但祂大力冲破童贞母亲的子宫束缚，由女先知所生，一个男性，正如\Person{依撒意亚}宣告的喜讯（\BibleRef{依 13:3}）。一岁的，因为祂是正义的太阳（\BibleRef{拉 4:2}），从天堂出发，被祂可见的本性所限制，并回归自身。\Quote{美善的蒙福冠冕}——各方面都相似于自身；不仅如此，还赋予所有美德之环生命，彼此温柔地混合交织，按照爱与秩序的法律。\OriginalTerm{无玷}{Immaculate}而无诡诈，作为过犯、以及来自罪的缺陷和污点的治疗者。因为虽然祂承担了我们的罪，背负了我们的疾病（\BibleRef{依 53:4}），但祂自己并没有遭受任何需要治疗的事。因为祂在一切事上受过试探，像我们一样，却没有罪（\BibleRef{希 4:15}）。因为那迫害在黑暗中照耀的光的，不能赶上祂。\par

十四. 还有什么？正月被引入，或者更确切地说，月月的开始，无论这在希伯来人中起初就是如此，还是后来因这个缘故才如此，并因奥秘而成为第一个；初十日，因为这是最完全的数，是单位中的第一个完满单位，是完美的根源。它被保留到第五日，也许因为我所讲的祭品净化五种感官，人因之犯罪跌倒，战争围绕这些感官激烈进行，因为它们易于受到犯罪的诱惑。它不仅从羔羊中选出，也从较低等的种类中，即放在左边的（\BibleRef{玛 25:33}）——山羊羔中选出；因为祂不仅为义人，也为罪人被祭献；也许甚至更为这些，因为我们更需要祂的仁慈。一只羔羊为一个家庭被要求作为最好的方式，但如果不能，则可由一场（因贫困）合献一只用于一个家族的家庭；因为显然，最好每个人都足以使自己完美，并向召唤他的天主奉献自己活的圣祭，时时刻刻在各方面被祝圣。但如果不能，则可利用那些在德行上相似且性情相同的人作为帮助。因为我认为这个规定意味着，若有需要，应将祭品分给最近的人。\par

十五、然后神圣之夜到来，这是今世混乱黑暗的周年纪念，原始黑暗在其中消解，万物进入生命、秩序和形式，混沌归于有序。然后我们从埃及逃离，即从阴郁迫害的罪恶；从看不见的暴君法老和苦工头目那里逃离，迁往上面的世界；得救脱离泥土和制砖，脱离这肉身状况的糠秕和危险，这状况对大多数人来说，几乎不被单纯的糠秕般的思虑所压倒。然后羔羊被宰杀，行为与言语以宝血封印；即习惯与行动，我们门框的侧柱；我指的是心意和意见的运动，它们由默观正确地开启和关闭，因为连思想也有界限。然后最后的、最重的灾祸临到迫害者，真正配得上这夜晚；埃及哀悼她自己推理和行动的长子，这在圣经中也被称为\Quote{加色丁人的种子}\BibleRef{友 5:6}被除去，巴比伦的儿女被摔在磐石上毁灭；整个空气中充满了埃及人的哭喊和喧嚷；然后，他们的毁灭者将因尊重傅油而离开我们。然后除去旧酵；即除去陈腐酸涩的邪恶，而非那使人活泼、做面包的酵；为期七天，这是一个最神秘的数字，与今世相协调，使我们不储备任何埃及面团，或法利塞人或不敬虔教训的遗物。\par

十六、好吧，让他们哀悼；我们将在傍晚吃羔羊——因为基督的受难是在时日的圆满；也因为祂在傍晚以祂的圣事分施给祂的门徒，摧毁了罪恶的黑暗；不是水煮，而是火烤——好叫我们的话中没有任何草率、稀薄或容易消失的东西；而是完全一致、坚实，远离一切不洁和一切虚浮。让我们借助好炭火\BibleRef{依 6:6}，从那来地上投火的那一位\BibleRef{路 12:49}得到帮助，这火将毁灭一切恶习，并促使其燃烧。因此，圣言中一切坚实滋养的部分，将与内心的部位和思想的隐藏处一同被吃，被消化并归于属灵的消化；是的，从头到脚，即从对天主性的最初默观直到对降生成人的最后思想。我们也不可把它带出外面，也不可留到早晨；因为我们大部分的奥迹不可传于外人，在这夜之后也没有进一步的净化；拖延对于有份于圣言的人来说是不可取的。因为正如不让愤怒延续整日，而在日落前消除它\BibleRef{弗 4:26}是善而中悦天主的——无论你从时间意义还是神秘意义来理解，因为正义的太阳落在我们的愤怒上对我们是不安全的——同样，我们也不该让这样的食物过夜，也不该推迟到明天。但是，凡是骨质的、不适合食用、甚至难以理解的部分，不可折断；也就是说，不要不好地猜测和误解（我无需说，在历史上耶稣没有一根骨头被折断，尽管钉死他的人因安息日而加速了他的死亡）；也不可剥去扔掉，免得圣物被丢给狗\BibleRef{玛 7:6}，即给邪恶的圣言听众；正如圣言荣美的珍珠不可丢在猪前；但它要由那探查知道万事的圣神用火——也是全燔祭所用的火——焚烧，被提炼和保存，而不是被水毁灭，也不是被撒散，如以色列人仓促铸造的牛犊头被梅瑟撒散\BibleRef{出 32:20}，作为他们心硬的羞辱。\par

十七、我们也不可忽略这吃的方式，因为法律并没有这么做，而是在字面规定中将它神秘劳作贯彻到这一点。让我们急忙吃这祭品，用无酵饼、苦菜，腰束带，脚穿鞋，像老人一样拄杖；急忙，以免我们陷入罗特被禁止犯的过错\BibleRef{创 19:17}，即不可回头看，也不可停留在整个附近，而要逃到山上，免得被索多玛的异火追上，或因转回邪恶而凝结成盐柱；因为这是迟延的结果。\Emph{用苦菜}，因为按照天主旨意的生活是苦涩艰难的，尤其对初学者而言，高于享乐。因为虽然新的轭是容易的，担子是轻省的\BibleRef{玛 11:20}，如你所听说的，但这是由于希望和赏报，它们远比今生的艰辛更丰富。若非如此，谁不说福音比法律的规条更充满辛劳和麻烦？因为法律只禁止完成的罪恶行为，而我们连起因也被定罪，几乎如同行为一样。法律说：\Quote{不可奸淫}；但\Emph{你}甚至\Emph{不可贪恋}，通过好奇和热切的眼神点燃情欲。\Quote{不可杀人}，法律说；但你甚至不可还手，反而要把自己交给打你的人。福音比法律苦修多少！法律说：\Quote{不可发虚誓}；但\Emph{你}不可起誓，无论是大誓还是小誓，因为发誓是伪誓的父母。\Quote{不可房屋连房屋，田地连田地，欺压穷人}\BibleRef{依 5:8}；但你要甘心放弃甚至你正当的财产，为穷人赤身，好毫无牵挂地背起十字架\BibleRef{谷 10:21}，并被那看不见的财富所致富。\par

XVIII. 让无理性动物的腰松开解缓，因为它们没有能战胜快乐的理性恩赐（不用说它们也知道自然运动的限度）。但是，让你身上那作为情欲座位、圣经称之为\BibleQuote{嘶鸣}\BibleRef{耶 5:8}的部分，在扫除这可耻的情欲时，被节制的腰带所约束，这样你就能纯洁地吃逾越晚餐，既使你们在地上的肢体死去\BibleRef{哥 3:5}，也效法那隐修士、先驱和真理的伟大宣告者\Person{若翰}的腰带\BibleRef{玛 3:4}。我还知道另一种腰带，我指的是那军事的、雄壮的腰带，\Place{叙利亚}的\OriginalTerm{厄乌作尼人}{Euzoni}和某些\OriginalTerm{莫诺作尼人}{Monozoni}就是从它得名的。也正是在这一点上，天主在神谕中对\Person{约伯}说：\BibleQuote{不，但你要像男子汉束紧你的腰，并给出一个雄壮的回答。}\BibleRef{约 38:3} 圣\Person{达味}也以此夸耀自己以来自天主的力量束腰，并说天主自己以力量为衣，以权能束腰——当然是针对不敬虔的人——尽管也许有人更愿意将此视为祂权能丰沛及有如克制的宣示，正如祂以光明为衣披戴一样。因为谁能承受祂不受约束的权能和光明呢？我是否询问腰与真理有何共同之处？那么，对圣\Person{保禄}来说，这话是什么意思呢？\BibleQuote{因此要站稳，用真理束腰？}\BibleRef{弗 5:14} 也许默观是要约束私欲偏情，不让它被引向别处？因为那倾向于某一特定方向的爱，对其他快乐就不会有同样的力量。\par

XIX. 至于\Emph{鞋}，那将要触摸天主之脚踏过的圣地的人，要脱下它们，如\Person{梅瑟}在山上所做的一样\BibleRef{出 3:5}，免得他带去任何死亡之物；没有任何介于人和天主之间的东西。同样，如果有门徒被派遣去宣讲福音，让他本着哲学精神而没有过度地出发，因为他除了没有钱、没有杖、只有一件外衣之外，还必须赤脚\BibleRef{玛 10:9}，好叫那些传报和平福音及一切美好事物的人的脚显得美丽\BibleRef{依 52:7}。但那想要逃离\Place{埃及}及埃及事物的人，必须为安全起见穿上鞋，尤其因埃及多的是蝎子和蛇，以免被那窥伺脚跟的所伤\BibleRef{创 3:15}，而我们也奉命践踏它们\BibleRef{路 10:19}。至于\Emph{杖}及其意义，我的看法如下。我知道有一种杖是用来倚靠的，另一种属于牧者和导师，用来矫正人类的羊群。如今法律为你规定了倚靠的杖，免得当你听到天主之血、祂的苦难和祂的死亡时，你的心灵破碎；也免得你在护卫天主时被卷入异端；而是毫无羞耻、毫无疑惑地吃那肉、喝那血，如果你渴慕真正的生命，既不要不信祂关于肉的话，也不要因祂关于苦难的话而跌倒。倚靠这一点，站得坚定刚强，不被敌手动摇，也不被他们言论的动听所迷惑。站在你的高处；在\Place{耶路撒冷}的庭院里放你的脚；倚靠那磐石，使你在天主里的脚步不被摇动。\par

XX. 你说什么？祂已乐意使你从\Place{埃及}，那铁炉中出来；使你离开那国的偶像崇拜，并由\Person{梅瑟}及其立法和军事统治所领导。我给你一个建议，这建议不是我自己的，或者说，如果你灵性地看待这事，它其实很是我自己的。向埃及人借金器银器\BibleRef{出 11:2}；带着这些上路；用外邦人的财物，或者更可以说是你们自己的，在路上供应自己。有钱欠着你们，是你们奴役和制砖的工价；你们也要精明地要求赔偿；做一个诚实的强盗。你们在那里受苦，与泥土（即这麻烦而污秽的身体）搏斗，建造异邦而不安全的城市，其纪念随呼喊一声而消逝。那么，难道你们空手而出，没有工钱吗？但为什么你们要把他们通过邪恶所得、又将用更大邪恶花费的东西留给埃及人和你们敌人的权势呢？那不属于他们：他们夺取了它，并亵渎地把它当作战利品从祂那里拿走，而祂说：\BibleQuote{银子是我的，金子是我的}\BibleRef{哈 2:8}，\BibleQuote{我愿给谁就给谁。}昨日那是他们的，因为被允许如此；今日主取回来给了你们\BibleRef{玛 20:14}，为使你们好好而有益地使用它。让我们用不义的\OriginalTerm{玛门}{Mammon}为自己结交朋友\BibleRef{路 16:9}，这样当我们失败时，他们可以在审判的时候接待我们。\par

二一. 如果你是一个辣黑耳或一个肋阿，一个宗祖般伟大的灵魂，就偷取你父亲所能找到的任何神像；\BibleRef{创 31:19}但不可保留它们，而要毁灭它们；如果你是一个聪明的以色列人，就把它们带到应许之地，让那迫害者因失去它们而悲伤，并因被智胜而明白，他试图压迫和奴役比他更好的人乃是徒劳的。如果你这样做了，并且这样出了埃及，我深知你必被火柱和云柱昼夜引导。\BibleRef{出 13:20}旷野必为你驯服，海水分开；法郎必被淹没；\BibleRef{出 14:28}粮食必从天而降；盘石必成为泉源；阿玛肋克必被征服，不单靠武器，更靠义人的手形成祈祷和十字架无敌的凯旋；约但河必被截断；太阳必停住；月亮必被节制；\BibleRef{苏 3:15}城墙必不借器械而倒塌；黄蜂群必在你前面开路，为以色列预备道路，并阻止外邦人；以及此后和与此相关的经中所记载的一切其他事件（不赘述）都将由天主赐给你。这就是你今天庆祝的庆节；我愿你以如此方式庆祝为你而生并为你受苦的那位的生日与葬礼。这就是逾越节的奥迹；这就是法律所预表、并被基督所完成的奥迹，祂是文字的废除者，圣神的成全者，祂藉自己的苦难教导我们如何受苦，藉自己的荣耀赐我们与祂同享荣耀。\par

二二. 现在我们该查考另一个事实和道理，它被大多数人忽视，但依我看来很值得探究。那为我们倾流的血是献给谁的？为何倾流？我是指我们天主、大司祭和祭献那可贵而著名的血。我们曾被恶者拘留为奴，被卖于罪恶之下，以邪恶换取快乐。既然赎价只属于那拘留者，我问：这赎价是献给谁的？为了什么缘故？如果是献给恶者，真是可恶之极！如果强盗接受赎价，且不仅是出于天主的赎价，而是本身即是天主的赎价，并有如此显赫的代价作为他的暴政的报偿——为了这个代价，他本应彻底放我们自由才对。但如果是献给圣父，我首先要问：如何献上？因为压迫我们的并非祂；其次，圣父怎会喜悦祂独生子的血？祂连依撒格在其父献上时也不接受，而是改变了祭献，用一只公羊替代了人祭。难道不明显吗：圣父接受了祂，但既未要求也未索要祂；而是因为降生成人，因为人性必须藉天主的人性而被圣化，好使祂亲自拯救我们，战胜暴君，并藉其圣子的中介吸引我们归向自己——圣子也安排这一切以光荣圣父，祂显然在一切事上服从圣父。关于基督，我们已说了这些；但我们可说的更大一部分当以静默敬畏。但那铜蛇\BibleRef{户 21:9}被挂起作为被咬之蛇的解药，并非作为那为我们受苦者的预表，而是作为对比；它拯救了仰望它的人，不是因为他们相信它是活的，而是因为它被杀死了，并且与它一起杀死了那受它辖制的权势，而它本身也当受毁灭。我们从何而为它写下合宜的墓志铭？\BibleQuote{死亡！你的刺在哪里？坟墓！你的胜利在哪里？} 你被十字架推翻；你被生命赋予者所杀；你无气息、死亡、不动，即使你仍保持被高举在杆上的蛇的形状。\par

二三. 现在我们要参与一个仍是预像的逾越节；虽然它比旧的更清晰。因为那如今正在被认识的，永远是新的。我们应当学习那饮宴和享乐是什么，而祂则负责教导并传达圣言给祂的门徒。因为教导是食物，甚至是对那赐予食物者而言。那么，来吧，让我们参与法律，但以福音的方式，而非字面的方式；圆满地，而非不圆满地；永恒地，而非暂时地。让我们以天上的耶路撒冷、而非地上的耶路撒冷为我们的首；\BibleRef{希 12:22} 不是那如今被军队践踏的，\BibleRef{路 21:20} 而是那被天使光荣的。让我们不祭献牛犊，也不祭献长出角蹄的羔羊，其中许多部分没有生命和感觉；而是让我们在天上的祭台上，以天上的舞蹈，向天主祭献赞颂之祭；让我们拉开第一层帐幔，走近第二层，并窥视至圣所。我是否要说一件更伟大的事？让我们将\Emph{我们自己}祭献给天主；或者更确切地说，让我们每天都时时刻刻继续祭献。让我们为圣言的缘故接受一切。藉着苦难让我们模仿祂的苦难：藉着我们的血让我们崇敬祂的血：让我们欣然登上十字架。钉子虽极痛苦，却是甘甜的。因为为基督并与基督一起受苦，胜过与他人安逸度日。\par

XXIV. 如果你是一个\Person{基勒乃人西满}，\BibleRef{谷 15:21} 背起十字架来跟随。如果你与祂同钉十字架如同强盗，\BibleRef{路 23:42} 要像一个\Emph{忏悔的}强盗承认天主。如果祂为了你和你罪的缘故被列于罪犯之中，\BibleRef{依 53:12} 你要为祂而成为守法者。敬拜那位为你被悬挂的，即使你自己也被悬挂；从你的邪恶中获取一些益处；以你的死亡购买救恩；与耶稣一起进入乐园，\BibleRef{路 23:43} 如此你就能从你所坠落之处学习。\BibleRef{默 2:5} 默观那里的荣耀；让杀人者带着他的亵渎死在外面；如果你是一个\Person{阿黎玛特雅人若瑟}，\BibleRef{路 23:52} 从钉死祂的人那里求得身体，使那洁净世界的成为你自己的。\BibleRef{若一 1:7} 如果你是一个\Person{尼苛德摩}，夜间敬拜天主的人，以香料埋葬祂。\BibleRef{若 19:39} 如果你是一个\Person{玛利亚}，或另一个\Person{玛利亚}，或\Person{撒罗默}，或\Person{若翰纳}，清晨哭泣。首先要看见石头被移开，也许你会看见天使和耶稣本人。说话；听祂的声音。如果祂对你说：\Quote{不要摸我}，就远远站着；敬畏圣言，但不要悲伤；因为祂知道祂首先向谁显现。庆祝复活节；来帮助那首先堕落的\Person{厄娃}，那首先拥抱基督并使祂被门徒认识的那位。成为\Person{伯多禄}或\Person{若望}；奔向坟墓，一起跑，彼此竞争，在尊贵的赛跑中争胜。即使你在速度上被击败，也要以热忱赢得胜利；不是\Emph{探视}坟墓，而是\Emph{进去}。如果你像\Person{多默}，当门徒聚集时基督向他们显现而你被排除在外，当你看见祂时不要不信；如果你不信，就相信那些告诉你的人；如果你也不能相信他们，那么就信赖钉痕。如果祂下到地狱，\BibleRef{伯前 3:19} 与祂同下。在那里也学习认识基督的奥秘，这双重降下的眷顾目的是什么：是借着祂的显现绝对地拯救所有人，还是在那里也仅仅拯救那些相信的人。\par

XXV. 如果祂升到天堂，\BibleRef{路 24:51} 与祂一同上升。成为那些护送祂的天使之一，或那些迎接祂的天使之一。命令门户被举起，或被升高，以便它们能迎接祂，在祂受难后被高举。回答那些因祂携带身体和祂受难的标记（这些是祂降下时所没有的）而怀疑的人，他们因此询问：\Quote{这位光荣的君王是谁？} 祂是大能的主，如同在所有事情上祂从古至今所行和所行的一样，现在也为人类争战并得胜。对于问题的疑惑给予双重的回答。如果他们惊奇并像在\Person{依撒意亚}的戏剧中那样说：\Quote{这位从\Place{厄东}和属地之物而来的是谁？} 或 \Quote{那没有血或身体之人的衣服如何是红的，如同在榨酒池中踩踏的人？} \BibleRef{依 63:1} 展示那受苦的身体衣着的美丽，被苦难装饰，被神性荣耀，没有什么比它更可爱或更美丽。\par

XXVI. 对此，那些吹毛求疵的人，那些关于神性的尖刻论者，那些一切值得赞美之事的诽谤者，那些光的黑暗者，对于智慧不学无术（基督为他们白白死了），忘恩负义的受造物，\Person{恶者}的工作，会说些什么？你会把这恩惠变成对天主的侮辱吗？你会因此而轻视祂，因为祂为了你而谦卑自己，因为善牧\BibleRef{若 10:11}（祂为羊舍命）为了寻找迷失的羊，来到了你曾献祭的山上和丘陵上，\BibleRef{若 5:35} 找到了那迷失的；找到了就扛在肩上，\BibleRef{欧 4:13} 祂也曾在肩上背着木头；背了之后，带回到上面的生命中；带回来之后，列在那些从未迷失的羊中。祂点了一盏灯，\BibleRef{路 15:4} 祂自己的肉身，打扫了房屋，清除了世界的罪，寻找那被各种情欲覆盖的硬币，即君王的肖像，并在找到硬币时召集祂的朋友，即天使力量，使他们分享祂的喜乐，如同先前使他们分享祂降生成人的奥秘一样。那极其耀眼的光要跟随灯——前驱，圣言要跟随声音，新郎要跟随新郎的朋友，那为主准备了特殊子民、用水洁净他们以准备接受圣神的人。\BibleRef{玛 3:11} 你据此谴责天主吗？你因此认为祂较小，因为祂用毛巾束腰，洗门徒的脚，\BibleRef{若 13:4} 并表明谦卑是通往高举的最好道路；\BibleRef{玛 23:12} 因为祂为那弯向地面的灵魂而谦卑自己，以便将那被罪的重担压弯的与自己一同高举？你为什么不也控告祂与税吏一同吃饭，\BibleRef{谷 2:15} 在税吏的桌子上，并使税吏成为门徒，\BibleRef{路 15:2} 使祂也能有所得？什么所得？罪人的救恩。如果是这样，就必须责备医生弯腰照顾病人并忍受恶臭以赐予病人健康；也责备那人俯身靠近沟渠，以便按照法律拯救掉进去的牲畜。\par

XXVII. 祂被派遣，但按照祂的人性被派遣（因为祂具有两种本性），因祂按人性之律饥饿、口渴、疲倦、忧伤并哭泣。但即使祂也作为天主被派遣，那又如何呢？将这派遣视为圣父的美意，祂将自己的一切都归于圣父，既为尊崇永恒本原，又为避免显得是一位敌对的天主。因为一方面经上记着祂被出卖，另一方面又写着祂把自己交出；同样，经上既说祂被圣父复活并接升，又说祂凭自己的权能复活并升天。前者属于美意，后者属于祂自己的权威；但你们只盯住那些贬抑祂的事，却忽视那些显扬祂的事。例如，你们强调祂受苦，却不加上\Quote{是出于祂自己的意愿}。唉，如今圣言还在受苦！有人尊祂为天主，却将祂与圣父混同；有人因祂的肉体而贬抑祂，将祂与天主分离。祂会更恼怒谁——或者更宽恕谁呢？是那些错误地压缩祂的人，还是那些分裂祂的人？前者本应加以区分，后者本应使之合一，一是在数目上，一是在天主性上。你们因祂的肉体而跌倒吗？犹太人也如此。你们称祂为撒玛黎雅人（\BibleRef{若 8:48}）以及那些我不忍出口的恶名吗？连魔鬼也未如此，你这比魔鬼更不信、比犹太人更愚顽的人啊！犹太人承认\Quote{子}的名号表示同等的地位；魔鬼也知道驱逐它们的是天主，因为它们由自身的经历所证实。但你既不承认平等，也不宣认天主性。你若行过割礼并成了附魔者——把这事说得荒谬些——倒比未受割礼、身体健康却如此患病且不虔敬要好。至于我们与这些人的争战，就让他们何时醒悟（若他们愿意）便何时停止吧；否则，若他们执迷不悟，就暂且搁置，等他们保持现状时再说。无论如何，我们靠圣三的帮助为圣三争战，必无所畏惧。\par

XXVIII. 现在有必要将我们的讲论总结如下：我们被创造，是为得享幸福。我们被造时便得了幸福。我们被托付于乐园，是为享受生命。我们领受了诫命，是为因遵守而获得美名；并非天主不知道将要发生的事，而是祂立下了自由意志的法则。我们因被嫉妒而受骗。我们因违命而被逐出。我们禁食，是因拒绝禁食，被知善恶树胜过。因为那诫命是古老的，与我们同始，是一种灵魂的教育和对奢靡的约束，我们有理由服从它，为能藉遵守它而恢复那因不遵守而失去的。我们需要一个降生成人的天主，一个被置于死地的天主，好使我们得以生活。我们与祂同死，为得洁净；我们因与祂同死而与祂同复活；我们因与祂同复活而与祂同受光荣。\par

XXIX. 那时确有许多奇迹：天主被钉十字架；太阳昏暗后又重放光明——因为受造物理应与其创造者同受苦难；帐幔裂开；血和水从祂肋旁流出，一者出于人性，一者超乎人性；磐石因磐石而裂开；死人复活，作为众人最终复活的保证；坟墓内和坟墓后的征兆，无人能妥善歌颂；然而这些无一件能与我的救恩奇迹相比。几滴血重造了整个世界，对众人而言，就如凝乳剂之于牛奶，将我们聚合、压缩为一体。\par

XXX. 但是，逾越节啊，伟大、神圣、净化全世界的逾越节——因为我将如同对一位活者向你说话——天主的圣言、光明、生命、智慧、德能——因我喜爱你的一切名号——伟大理智的后裔、表达和印记；圣言，既被孕育又被默观为人，承载万有，以你德能的圣言束缚万有；请接纳这篇讲论，如今不是作为初果，也许是我奉献的完成，一份感恩，同时是一份恳求：愿我们除那些必要而神圣的照料（我们的生命在其中度过）之外，不受任何邪恶之苦；并求你制止身体对我们的暴政（主啊，你知道它多么巨大，如何压垮我），或你自身的判决，若我们要被你定罪。但若我们按愿望得释放，并被接入天上的帐幕，在那里我们或许也会在你的祭台上向你献上悦纳的祭品，归于圣父、圣言和圣神；因为荣耀、尊威和权能都归于你，万世无疆。阿们。\par

\SourceRecord{Translated by Charles Gordon Browne and James Edward Swallow.；Nicene and Post-Nicene Fathers, Second Series；Vol. 7.；Edited by Philip Schaff and Henry Wace.；Buffalo, NY: Christian Literature Publishing Co.,；1894.；<http://www.newadvent.org/fathers/310245.htm>.}
